% =====================================================================
%  Olasılık Teorisi ve Matematiksel İstatistik
%  EĞİTMEN ÇÖZÜM KILAVUZU
%  \alcozum görünür — tüm alıştırma çözümleri gösterilir
% =====================================================================
\documentclass[12pt, a4paper, twoside]{book}

% ---- Dil ve kodlama ----
\usepackage[utf8]{inputenc}
\usepackage[T1]{fontenc}
\usepackage[turkish]{babel}

% ---- Yazı tipi ----
\usepackage{lmodern}
\usepackage{microtype}

% ---- Sayfa düzeni ----
\usepackage[top=2.5cm, bottom=2.5cm, inner=3cm, outer=2cm]{geometry}
\usepackage{fancyhdr}
\pagestyle{fancy}
\fancyhf{}
\fancyhead[LE]{\small\leftmark}
\fancyhead[RO]{\small\rightmark}
\fancyfoot[C]{\thepage}

% Eğitmen sürümü uyarısı — her sayfada
\fancyhead[RE,LO]{\small\color{red!70!black}\bfseries EĞİTMEN SÜRÜMÜ}

% ---- Matematik ----
\usepackage{olasilik}

% ---- Pedagojik kutular (eğitmen sürümü: alcozum görünür) ----
\usepackage[cozumler]{maarif}

% ---- Diğer paketler ----
\usepackage{graphicx}
\usepackage{booktabs}
\usepackage{hyperref}
\hypersetup{
  colorlinks=true,
  linkcolor=blue!60!black,
  citecolor=green!50!black,
  urlcolor=blue!70!black,
  pdftitle={Olasılık Teorisi ve Matematiksel İstatistik — Eğitmen Çözüm Kılavuzu},
  pdfauthor={},
}
\usepackage{cleveref}

% ---- Başlık sayfası ----
\title{%
  {\Huge\bfseries Olasılık Teorisi ve}\\[4pt]
  {\Huge\bfseries Matematiksel İstatistik}\\[12pt]
  {\LARGE\color{red!70!black}\bfseries Eğitmen Çözüm Kılavuzu}\\[6pt]
  {\large\itshape Dağıtımı kısıtlıdır — yalnızca eğitmen kullanımına yöneliktir}%
}
\author{}
\date{}

% =====================================================================
\begin{document}

\frontmatter
\maketitle

{\color{red!70!black}\bfseries\large
Bu belge yalnızca eğitmen kullanımına yöneliktir.
Öğrencilerle paylaşılmamalıdır.
}
\vspace{1cm}

\tableofcontents

\mainmatter

% ---- PART I: Olasılık Teorisi ----
\part{Olasılık Teorisi}

% =====================================================================
%  BÖLÜM 0 — Ön Bilgiler: Kümeler, Fonksiyonlar ve Sayılabilirlik
%  Olasılık Teorisi ve Matematiksel İstatistik
%  Kaynaklar: Billingsley Ek A; Klenke Böl. 1; Durrett Ek A
% =====================================================================

\chapter{Ön Bilgiler: Kümeler, Fonksiyonlar ve Sayılabilirlik}

\bolumalinti{``Mathematicians always strive to confuse their audiences;
  where there is no confusion, there is no prestige.''}{Gian-Carlo Rota (1932--1999)}

\noindent\textbf{Ön Koşullar:} Lise düzeyinde kümeler ve fonksiyonlar kavramı.
Bir önceki bölüm yoktur; bu bölüm tüm kitabın temel taşıdır.

\noindent\textbf{Bu Bölüm Sonraki Bölümlerde Kullanılır:}
Bölüm~1 (ölçüm teorisi için doğrudan),
Bölüm~2 (olasılık uzayı yapısı),
Bölüm~3 (rasgele değişkenler).

\noindent\textbf{Tahmini Okuma Süresi:} $\approx 4$~saat.

\vspace{0.5\baselineskip}

% ─── KAZANIMLAR ────────────────────────────────────────────────────
\begin{kazanimlar}
Bu bölümün sonunda öğrenci:
\begin{itemize}[leftmargin=1.8em, itemsep=2pt]
  \item Küme cebiri işlemlerini (birleşim, kesişim, fark, tümleyen) hem
        sonlu hem sayılabilir hem de keyfi indis kümeleri üzerinde
        uygulayabilmeli ve De~Morgan yasalarını ispatlayabilmeli,
  \item Ters görüntü (\emph{preimage}) kavramını tanımlayabilmeli ve ters
        görüntünün birleşim, kesişim ve tümleme işlemlerini koruyduğunu
        ispatlayabilmeli,
  \item Sonlu, sayılabilir ve sayılamaz kümeleri ayırt edebilmeli;
        $\Q$'nun sayılabilir, $\R$'nin sayılamaz olduğunu ispatlayabilmeli,
  \item Gerçel sayılarda üst sınır (\emph{supremum}) ve alt sınır
        (\emph{infimum}) kavramlarını tanımlayabilmeli ve tamtümlük
        aksiyomunu ifade edebilmeli,
  \item Genişletilmiş gerçel sayı doğrusu $\Rbar$'yi
        ve üzerinde yapılan aritmetik kuralları açıklayabilmeli,
  \item Sayı dizileri için $\limsup$ ve $\liminf$'i tanımlayabilmeli ve
        hesaplayabilmeli,
  \item Küme dizileri için $\limsup$ ve $\liminf$'i yorumlayabilmeli
        (``sonsuz kez'' ve ``eninde sonunda'' olay dili ile ilişkilendirebilmeli).
\end{itemize}
\end{kazanimlar}

% ─── MOTİVASYON ────────────────────────────────────────────────────
\begin{motivasyon}
Olasılık teorisinin kalbi \emph{ölçüm teorisidir} (measure theory) ve ölçüm
teorisi tamamen küme işlemleri üzerine kuruludur. Bir olayın olasılığını
tanımlamak için, önce ``hangi kümelere olasılık atayabiliriz?'' sorusunu
yanıtlamamız gerekir. Bu sorunun yanıtı $\sigma$-cebiri kavramına götürür;
$\sigma$-cebirleri ise sayılabilir birleşim ve kesişim gibi işlemleri
barındırdıklarından, sayılabilirlik kavramını doğrudan olasılık teorisinin
merkezine taşır.

\medskip

\textbf{Köprü kurma.} Lise matematiğinde sonlu küme işlemlerini (Venn
şemaları, De~Morgan yasaları) öğrendiniz. Bu bölümde o bilgiyi \emph{sonsuz
küme aileleri} üzerine genişleteceğiz. Böylece Bölüm~1'de, sonsuz çok olayı
kapsayan bir $\sigma$-cebirini kolaylıkla tanımlayabilecek ve Kolmogorov'un
aksiyomlarını Bölüm~2'de tam anlamıyla anlayabileceksiniz.

\medskip

Bu bölümde şunları öğreneceğiz:
küme cebiri ile rahat çalışmayı ($\S0.1$),
fonksiyonların ters görüntü özelliklerini ($\S0.2$),
sayılabilirlik teorisini ($\S0.3$),
gerçel sayı sisteminin tamtümlük özelliğini ($\S0.4$)
ve sayı/küme dizilerinde $\limsup$/$\liminf$ kavramlarını ($\S0.5$).
\end{motivasyon}


% =====================================================================
\section{Küme Cebiri}
% =====================================================================

\begin{kesif}[Sonsuz birleşim ne demektir?]
Sonlu iki kümenin birleşimini $A \cup B$ olarak biliyorsunuz. Peki
$A_1 \cup A_2 \cup A_3 \cup \cdots$ ifadesini nasıl yazarsınız?
Düşüncelerinizi aşağıya not edin; ardından Tanım~\ref{def:kumecebiri}'yi
okuyun.
\ogrencialani[2.5cm]
\end{kesif}

\begin{tanimgerekce}
Sonsuz birleşimi $\bigcup_{n=1}^\infty A_n$ olarak yazmak sadece notasyon
kolaylığı değildir. Olasılık teorisinde olayların ``en az birinin gerçekleşmesi''
durumu doğal olarak sonsuz birleşim gerektirir. Örneğin bir madeni paranın
sonsuza dek atılması deneyi düşünüldüğünde, ``en az bir kez tura gelme'' olayı
sonsuz birleşimle ifade edilir.
\end{tanimgerekce}

\begin{tanim}[Temel Küme İşlemleri]\label{def:kumecebiri}
Bir $\Omega$ evren kümesi ve $\Omega$'nın alt kümeleri $A, B, (A_n)_{n \geq 1}$
verilsin. Aşağıdaki işlemleri tanımlıyoruz:
\begin{align*}
  \text{Tümleyen:}\quad
    & A^c \coloneqq \{\omega \in \Omega : \omega \notin A\}, \\[2pt]
  \text{Fark:}\quad
    & A \setminus B \coloneqq A \cap B^c
      = \{\omega \in \Omega : \omega \in A,\ \omega \notin B\}, \\[2pt]
  \text{Simetrik fark:}\quad
    & A \triangle B \coloneqq (A \setminus B) \cup (B \setminus A), \\[2pt]
  \text{Sayılabilir birleşim:}\quad
    & \bigcup_{n=1}^\infty A_n
      \coloneqq \{\omega \in \Omega : \omega \in A_n \text{ bazı } n \in \N \text{ için}\}, \\[2pt]
  \text{Sayılabilir kesişim:}\quad
    & \bigcap_{n=1}^\infty A_n
      \coloneqq \{\omega \in \Omega : \omega \in A_n \text{ her } n \in \N \text{ için}\}.
\end{align*}
İki küme \textbf{ayrık} \emph{(disjoint)} adını alır; eğer $A \cap B = \emptyset$.
Bir $(A_n)$ kümeler ailesi \textbf{ikili ayrık} \emph{(pairwise disjoint)} adını alır
eğer $A_i \cap A_j = \emptyset$, $i \neq j$ için.
\end{tanim}

\begin{kritiksorgu}[Keyfi indis kümeleri]
Yukarıdaki tanım $n \in \N$ için yazıldı. Peki $I$ herhangi bir indis
kümesi olsaydı, yani $\bigcup_{i \in I} A_i$, tanım nasıl değişirdi? $I$
sayılamaz olursa ne olur?
\end{kritiksorgu}

Küme cebiri işlemlerinin en önemli özellikleri De~Morgan yasaları aracılığıyla
ifade edilir. Bu yasalar, ölçüm teorisinde $\sigma$-cebirlerinin tümleme altında
kapalı olduğunu ispat etmede kullanılacak temel araçtır.

\begin{teorem}[De~Morgan Yasaları]\label{thm:demorgan}
Bir $(A_n)_{n \geq 1}$ kümeler ailesi için:
\[
  \left(\bigcup_{n=1}^\infty A_n\right)^c = \bigcap_{n=1}^\infty A_n^c
  \qquad \text{ve} \qquad
  \left(\bigcap_{n=1}^\infty A_n\right)^c = \bigcup_{n=1}^\infty A_n^c.
\]
\end{teorem}

\begin{proof}
\textit{Fikir:} ``$A_1 \cup A_2 \cup \cdots$'nin tümleyeni'' demek ``hiçbir $A_n$'de
bulunmamak'' demektir; bu ise tüm $A_n^c$'lerin kesişimidir.

\medskip

\textbf{Adım 1.} İlk eşitliği kanıtlıyoruz. $\omega \in \left(\bigcup_n A_n\right)^c$
olsun. Bu, $\omega \notin \bigcup_n A_n$ demektir; yani \emph{hiçbir} $n$ için
$\omega \in A_n$ değildir. Dolayısıyla \emph{her} $n$ için $\omega \in A_n^c$,
yani $\omega \in \bigcap_n A_n^c$.

\textbf{Adım 2.} Tersine, $\omega \in \bigcap_n A_n^c$ olsun. Her $n$ için
$\omega \in A_n^c$, yani $\omega \notin A_n$. O hâlde $\omega \notin \bigcup_n A_n$,
yani $\omega \in \left(\bigcup_n A_n\right)^c$.

\medskip

\textbf{Adım 3.} İki içerme birden elde edildiğinden
$\left(\bigcup_n A_n\right)^c = \bigcap_n A_n^c$.
İkinci eşitlik, birincisini $A_n^c$ yerine uygulamakla ya da
$A \mapsto A^c$ simetrisiyle elde edilir.
\end{proof}

\begin{ornek}[De~Morgan yasası --- somut uygulama]\label{ex:demorgan}
$\Omega = \R$ ve $A_n = (0, n)$, $n \geq 1$ olsun.
O hâlde $\bigcup_{n=1}^\infty A_n = (0,\infty)$ ve
\[
  \left(\bigcup_{n=1}^\infty A_n\right)^c = (-\infty,0].
\]
Öte yandan $A_n^c = (-\infty,0] \cup [n,\infty)$, dolayısıyla
\[
  \bigcap_{n=1}^\infty A_n^c = (-\infty,0],
\]
çünkü $[n,\infty)$'ların kesişimi boştur. De~Morgan eşitliği doğrulanmış oldu.
\end{ornek}

Küme dizileri için monoton diziler özellikle önemlidir.

\begin{tanim}[Artan ve Azalan Küme Dizileri]\label{def:monotondizi}
$(A_n)_{n \geq 1}$ bir küme dizisi olsun.
\begin{itemize}[leftmargin=1.8em]
  \item Dizi \textbf{artan} \emph{(increasing)} adını alır; eğer $A_1 \subseteq A_2 \subseteq \cdots$.
        Bu durumda $A_n \nearrow$ gösterimi kullanılır ve
        $\lim_n A_n \coloneqq \bigcup_n A_n$ tanımlanır.
  \item Dizi \textbf{azalan} \emph{(decreasing)} adını alır; eğer $A_1 \supseteq A_2 \supseteq \cdots$.
        Bu durumda $A_n \searrow$ gösterimi kullanılır ve
        $\lim_n A_n \coloneqq \bigcap_n A_n$ tanımlanır.
\end{itemize}
\end{tanim}

\begin{mikroozet}[§0.1 özeti]
Temel küme işlemleri sonsuz indis kümeleri üzerine genişletildi.
De~Morgan yasaları tümleme ile birleşim/kesişim arasındaki dönüşümü sağlar.
Monoton küme dizilerinin limiti birleşim ya da kesişim olarak tanımlanır.
\end{mikroozet}


% =====================================================================
\section{Fonksiyonlar ve Ters Görüntü}
% =====================================================================

\begin{motivasyon}[§0.2]
Bir rasgele değişken, tanımı gereği ölçülebilir bir fonksiyondur.
``Ölçülebilir'' sıfatının ne anlama geldiğini anlamak için,
bir fonksiyonun \emph{ters görüntüsünün} küme işlemlerini nasıl
koruduğunu bilmek şarttır. Bu alt bölüm tam olarak bu teknik temel
için vardır.
\end{motivasyon}

\begin{tanim}[Ters Görüntü (Preimage)]\label{def:tergoriuntu}
$f: X \to Y$ bir fonksiyon ve $B \subseteq Y$ olsun.
$f$'nin $B$ altındaki \textbf{ters görüntüsü} \emph{(preimage, inverse image)}:
\[
  f^{-1}(B) \coloneqq \{x \in X : f(x) \in B\}.
\]
$f$'nin tersi mevcut olmasa bile $f^{-1}(B)$ her zaman tanımlıdır.
\end{tanim}

\begin{hataanalizi}[Ters görüntü ile ters fonksiyon]
$f^{-1}(B)$ gösterimi yalnızca $f$ bijektif olduğunda
ters fonksiyonu ifade eder. Genel durumda $f^{-1}(B)$ bir
\emph{küme} işlemidir ve $f$'nin bijektif olmasını gerektirmez.
Örneğin $f(x) = x^2$ ve $B = \{4\}$ ise $f^{-1}(B) = \{-2, 2\}$.
\end{hataanalizi}

Ters görüntünün olasılık teorisini bu kadar kullanışlı kılan temel özellik
şudur: \emph{ters görüntü tüm küme işlemlerini korur.}

\begin{teorem}[Ters Görüntünün Küme İşlemlerini Koruması]\label{thm:tergoriuntu}
$f: X \to Y$ ve $(B_n)_{n \geq 1} \subseteq \mathcal{P}(Y)$ olsun. O zaman:
\begin{enumerate}[label=(\roman*)]
  \item $f^{-1}\!\left(\bigcup_n B_n\right) = \bigcup_n f^{-1}(B_n)$,
  \item $f^{-1}\!\left(\bigcap_n B_n\right) = \bigcap_n f^{-1}(B_n)$,
  \item $f^{-1}(B^c) = \left(f^{-1}(B)\right)^c$ her $B \subseteq Y$ için.
\end{enumerate}
\end{teorem}

\begin{proof}
\textit{Fikir:} Ters görüntü, ``$f(x) \in$'' koşuluyla çalışır; birleşim ve
kesişim de ``bazı / her'' niceleyicileriyle çalışır. Bu iki yapı birebir örtüşür.

\medskip

\textbf{(i)}\; $x \in f^{-1}\!\left(\bigcup_n B_n\right)$
$\iff f(x) \in \bigcup_n B_n$
$\iff$ bazı $n$ için $f(x) \in B_n$
$\iff$ bazı $n$ için $x \in f^{-1}(B_n)$
$\iff x \in \bigcup_n f^{-1}(B_n)$.

\textbf{(ii)}\; $x \in f^{-1}\!\left(\bigcap_n B_n\right)$
$\iff f(x) \in \bigcap_n B_n$
$\iff$ her $n$ için $f(x) \in B_n$
$\iff$ her $n$ için $x \in f^{-1}(B_n)$
$\iff x \in \bigcap_n f^{-1}(B_n)$.

\textbf{(iii)}\; $x \in f^{-1}(B^c)$
$\iff f(x) \in B^c$
$\iff f(x) \notin B$
$\iff x \notin f^{-1}(B)$
$\iff x \in \left(f^{-1}(B)\right)^c$.
\end{proof}

\begin{notkutu}
Teorem~\ref{thm:tergoriuntu}'nun önemi şuradan gelir: Bölüm~1'de
$\sigma$-cebirlerini tanımladığımızda, bir fonksiyonun ``ölçülebilir'' olması tam
olarak, ters görüntüsünün $\sigma$-cebirini $\sigma$-cebirine taşıması anlamına
gelecektir. Yukarıdaki üç özellik bu taşımanın iyi tanımlı olduğunu garanti eder.
\end{notkutu}

\begin{tanim}[Bijektif Fonksiyon Türleri]\label{def:bijektif}
$f: X \to Y$ bir fonksiyon olsun.
\begin{itemize}[leftmargin=1.8em]
  \item $f$ \textbf{enjektif} \emph{(injective, one-to-one)} adını alır;
        eğer $f(x_1) = f(x_2) \Rightarrow x_1 = x_2$.
  \item $f$ \textbf{sürjektif} \emph{(surjective, onto)} adını alır;
        eğer her $y \in Y$ için bir $x \in X$ mevcutsa $f(x) = y$.
  \item $f$ \textbf{bijektif} \emph{(bijective)} adını alır;
        eğer hem enjektif hem sürjektifse.
        Bijektif fonksiyona aynı zamanda \textbf{bire-bir eşleşme}
        \emph{(one-to-one correspondence)} da denir.
\end{itemize}
\end{tanim}

\begin{kendinleifade}
Teorem~\ref{thm:tergoriuntu}'nun (iii) maddesini kendi cümlenizle ifade edin:
``$f^{-1}(B^c) = (f^{-1}(B))^c$ ifadesi ne anlama geliyor?''
Özellikle $B = Y$ ve $B = \emptyset$ özel durumlarını kontrol edin.
\ogrencialani[2cm]
\end{kendinleifade}

\begin{mikroozet}[§0.2 özeti]
Ters görüntü $f^{-1}(B)$ her fonksiyon için tanımlıdır ve bijektiflik
gerektirmez. Ters görüntü birleşim, kesişim ve tümlemeyi korur. Bu özellik
ölçülebilirlik kavramının temel taşıdır.
\end{mikroozet}


% =====================================================================
\section{Sonlu, Sayılabilir ve Sayılamaz Kümeler}
% =====================================================================

\begin{motivasyon}[§0.3]
``Sayılabilir'' kavramı sezgilerimizi yanıltabilir: $\Z$ sonsuz olmasına
karşın ``sayılabilir'', $\R$ ise sonsuz olup aynı zamanda ``sayılamaz''. Bu
fark, $\sigma$-cebirlerinde neden \emph{sayılabilir} birleşim kullandığımızı
(sayılamaz değil) açıklar. Sayılamaz birleşime izin versek, tüm singletonların
birleşimi $\R$'dir; ama bu birleşime Lebesgue ölçüsü atayamazdık.
\end{motivasyon}

\begin{tanim}[Sayılabilirlik]\label{def:sayilabilir}
Bir $A$ kümesi:
\begin{itemize}[leftmargin=1.8em]
  \item \textbf{sonlu} adını alır; eğer $A = \emptyset$ ya da $\{1,\ldots,n\}$
        ile $A$ arasında bir bijeksiyon varsa (bazı $n \in \N$ için).
  \item \textbf{sayılabilir sonsuz} \emph{(countably infinite)} adını alır;
        eğer $\N$ ile $A$ arasında bir bijeksiyon mevcutsa.
  \item \textbf{sayılabilir} \emph{(countable)} adını alır; eğer sonlu ya da
        sayılabilir sonsuzsa.
  \item \textbf{sayılamaz} \emph{(uncountable)} adını alır; eğer sayılabilir değilse.
\end{itemize}
\end{tanim}

\begin{ornek}[Sayılabilir sonsuz kümeler]\label{ex:sayilabilir}
$\Z$ sayılabilir sonsuzdur: $f: \N \to \Z$,
\[
  f(n) = \begin{cases} n/2 & n \text{ çift}, \\ -(n-1)/2 & n \text{ tek} \end{cases}
\]
bijeksiyonu $0,1,-1,2,-2,3,-3,\ldots$ sırasıyla saydırır.

$\N \times \N$ de sayılabilir sonsuzdur: Cantor'un çapraz saydırma düzeniyle
$(1,1),(1,2),(2,1),(1,3),(2,2),(3,1),\ldots$ sırasıyla sayılır.
\end{ornek}

\begin{teorem}[$\Q$ Sayılabilirdir]\label{thm:Qsayilabilir}
Rasyonel sayılar kümesi $\Q$ sayılabilirdir.
\end{teorem}

\begin{proof}
Her $q \in \Q^+$, $q = p/n$ ($p \in \N$, $n \in \N$, $\gcd(p,n) = 1$) biçiminde
tek türlü yazılır. $\Q^+ \subseteq \N \times \N$ gömülüdür ve $\N \times \N$
sayılabilirdir (Örnek~\ref{ex:sayilabilir}). $\Q = \Q^+ \cup \{0\} \cup (-\Q^+)$
üç sayılabilir kümenin birleşimidir; dolayısıyla sayılabilirdir.
\end{proof}

\begin{teorem}[$\R$ Sayılamazdır --- Cantor'un Köşegen Argümanı]
\label{thm:Rsayilamaz}
Gerçel sayılar kümesi $\R$ sayılamazdır.
Daha güçlü bir ifadeyle, $(0,1)$ aralığı bile sayılamazdır.
\end{teorem}

\begin{proof}
\textit{Fikir:} $(0,1)$'in sayılabilir olduğunu varsayarak çelişki türetiriz;
verilen herhangi bir sayımın ``dışarıda bıraktığı'' bir eleman üretiriz.

\medskip

\textbf{Çelişki varsayımı.} $(0,1)$'in sayılabilir olduğunu, yani
$(0,1) = \{x_1, x_2, x_3, \ldots\}$ olduğunu varsayalım.

\textbf{Adım 1.} Her $x_n$'i ondalık gösterimle yazın:
$x_n = 0.a_{n1}\,a_{n2}\,a_{n3}\cdots$, $a_{nk} \in \{0,1,\ldots,9\}$.
(Sona giden $9$'ların tekrarlanmasını önlemek için: $0.999\cdots = 1$ eşitliğini
göz önünde bulundurarak sonlanan ondalıkları tercih edin.)

\textbf{Adım 2.} Yeni bir $y = 0.b_1 b_2 b_3 \cdots$ sayısı inşa edin:
\[
  b_n \coloneqq
  \begin{cases}
    5, & a_{nn} \neq 5, \\
    6, & a_{nn} = 5.
  \end{cases}
\]

\textbf{Adım 3.} $y \in (0,1)$'dir. Öte yandan $y \neq x_n$ her $n$ için,
çünkü $y$ ile $x_n$'in $n$-inci ondalık basamakları farklıdır ($b_n \neq a_{nn}$).

Bu durum, $y$'nin listede yer almadığı anlamına gelir; $(0,1) = \{x_1, x_2, \ldots\}$
varsayımı çelişkilidir.
\end{proof}

\begin{karsiornek}[Sayılabilir kümenin tümleni sayılamaz olabilir]
$\Q$, $\R$ içinde sayılabilirdir; fakat $\Q^c = \R \setminus \Q$
(irrasyonel sayılar) sayılamazdır. Evreni sayılamaz olan bir kümenin
sayılabilir bir alt kümesinin tümleyeni sayılamaz olmak zorundadır.
\end{karsiornek}

\begin{teorem}[Sayılabilir Ailenin Birleşimi]\label{thm:sayilabilirbir}
Eğer her $n \in \N$ için $A_n$ sayılabilir ise,
$\bigcup_{n=1}^\infty A_n$ da sayılabilirdir.
\end{teorem}

\begin{proof}
Her $A_n$ sayılabilir olduğundan $A_n = \{a_{n1}, a_{n2}, a_{n3}, \ldots\}$
olarak yazılabilir. $\bigcup_n A_n$'ın elemanlarını Cantor'un çapraz tarama
düzeniyle sıralıyoruz:
\[
  a_{11},\; a_{12}, a_{21},\; a_{13}, a_{22}, a_{31},\; \ldots
\]
Bu, $\bigcup_n A_n$'dan $\N$'ye enjektif bir fonksiyon verir (tekrarları atlayın);
dolayısıyla $\bigcup_n A_n$ sayılabilirdir.
\end{proof}

\begin{notkutu}
Teorem~\ref{thm:sayilabilirbir}'in tersi doğru değildir: sayılamaz bir küme
sayılamaz çok sayılabilir parçaya bölünebilir. Örneğin $\R$ sayılamaz,
ama her tekil küme $\{x\}$ sonludur.
\end{notkutu}

% Disiplinlerarası bağlantı — TYMM gereksinimi
\begin{tcolorbox}[maarif@base, colframe=maKopru, colback=maKopru!4,
    colbacktitle=maKopru,
    title={Disiplinlerarası bağlantı --- Bilgisayar bilimi ve hesaplanabilirlik}]
Sayılabilirlik kavramı bilgisayar biliminde merkezi bir rol oynar.
Herhangi bir programlama dilinde yazılabilen tüm programlar kümesi sayılabilirdir
(her program sonlu bir karakter dizisidir). Buna karşın $\R \to \R$
fonksiyonları kümesi sayılamazdır. Dolayısıyla ``hesaplanamaz fonksiyonlar'' vardır;
bu Turing'in durma problemi ile doğrudan bağlantılıdır.
\end{tcolorbox}

\begin{mikroozet}[§0.3 özeti]
Sayılabilirlik, bir kümenin $\N$ ile eşleştirilebilmesi anlamına gelir.
$\Q$ sayılabilir, $\R$ sayılamazdır. Sayılabilir ailenin sayılabilir
birleşimi sayılabilirdir. $\sigma$-cebiri tanımında ``sayılabilir birleşim'' koşulu
bu ayrımı işlevsel kılar.
\end{mikroozet}


% =====================================================================
\section{Gerçel Sayılar ve Genişletilmiş Gerçel Sayı Doğrusu}
% =====================================================================

\begin{motivasyon}[§0.4]
Olasılık teorisinde beklentiler, seriler ve integraller kullanılır. Bunların
bazıları $+\infty$ ya da $-\infty$ değer alabilir. Örneğin bir Poisson rasgele
değişkeninin tüm momentleri sonludur; ama bazı ağır kuyruklu dağılımların birinci
momenti bile sonsuz olabilir. Bu yüzden gerçel sayıları $\pm\infty$ ile
genişletmek zorunludur.
\end{motivasyon}

\begin{tanim}[Supremum ve İnfimum]\label{def:supinf}
$A \subseteq \R$, $A \neq \emptyset$ olsun.
\begin{itemize}[leftmargin=1.8em]
  \item $M \in \R$'ye $A$'nın \textbf{üst sınırı} \emph{(upper bound)} denir;
        eğer $a \leq M$ her $a \in A$ için.
  \item $A$'nın \textbf{üst kenar değeri} \emph{(supremum)}, var olan tüm üst
        sınırların en küçüğüdür; $\sup A$ ile gösterilir.
  \item Benzer biçimde \textbf{alt sınır} \emph{(lower bound)} ve
        \textbf{alt kenar değeri} \emph{(infimum)} $\inf A$ tanımlanır.
\end{itemize}
$A$ yukarıdan sınırlı değilse $\sup A = +\infty$; aşağıdan sınırlı değilse
$\inf A = -\infty$ olarak tanımlanır. Boş küme için $\sup \emptyset = -\infty$
ve $\inf \emptyset = +\infty$ konvansiyonu kullanılır.
\end{tanim}

\begin{aksiyom}[Gerçel Sayıların Tamtümlüğü]
\label{ax:tamtumluk}
Yukarıdan sınırlı her boşalmayan $A \subseteq \R$ kümesinin bir
$\sup A \in \R$ üst kenar değeri vardır.
\end{aksiyom}

\begin{notkutu}
$\Q$ bu aksiyomu sağlamaz: $\{q \in \Q : q^2 < 2\}$ kümesinin
$\Q$ içinde supremumu yoktur ($\sqrt{2} \notin \Q$).
Tamtümlük aksiyomu, $\R$'yi $\Q$'dan ayıran temel özelliktir ve
analizin tüm önemli teoremlerinin (Bolzano-Weierstrass, Heine-Borel, vb.) temelidir.
\end{notkutu}

\begin{tanim}[Genişletilmiş Gerçel Sayı Doğrusu]\label{def:Rbar}
$\Rbar \coloneqq \R \cup \{-\infty, +\infty\}$ kümesi
aşağıdaki düzen ve aritmetik kurallarla donatılır:
\begin{gather*}
  -\infty < x < +\infty \quad \text{her } x \in \R \text{ için},\\[4pt]
  x + \infty = +\infty, \quad
  x - \infty = -\infty \quad (x \in \R), \\[2pt]
  x \cdot (+\infty) = \operatorname{sgn}(x) \cdot \infty \quad (x \neq 0), \\[2pt]
  0 \cdot (\pm\infty) = 0
  \quad \text{\emph{(olasılık teorisinde bu konvansiyon kullanılır!)}}, \\[2pt]
  (+\infty) + (-\infty) \quad \text{tanımsız.}
\end{gather*}
\end{tanim}

\begin{uyari}
$(+\infty) - (+\infty)$ ve $(+\infty) + (-\infty)$ tanımsızdır.
Bu ifadelerle karşılaşıldığında problem ifadesini yeniden düzenlemek gerekir.
Olasılık teorisinde beklenti hesaplarken bu durum ortaya çıkabilir
(Bölüm~4'te ele alınacaktır).
\end{uyari}

\begin{hataanalizi}[$0 \cdot \infty = 0$ konvansiyonunun önemi]
Olasılık teorisinde $0 \cdot \infty = 0$ konvansiyonu \emph{kasıtlı olarak}
seçilmiştir: sıfır olasılıklı bir olayın beklentiye katkısının sıfır olmasını
istiyoruz. Gerçel analizin bazı kitapları bu konvansiyonu kullanmaz; bu nedenle
farklı kaynaklardaki tanımları dikkatle karşılaştırın.
\end{hataanalizi}

\begin{ornek}[Supremum ve İnfimum hesabı]
\begin{enumerate}[label=(\alph*)]
  \item $A = \{1 - 1/n : n \in \N\} = \{0, 1/2, 2/3, 3/4, \ldots\}$
        için $\sup A = 1$ (ama $1 \notin A$) ve $\inf A = 0$ ($0 \in A$).
  \item $B = \{n \in \N\} = \{1,2,3,\ldots\}$ için $\inf B = 1 \in B$ ve
        $\sup B = +\infty$.
  \item $C = \emptyset$ için $\sup C = -\infty$ ve $\inf C = +\infty$
        (konvansiyonel tanım).
\end{enumerate}
\end{ornek}

\begin{mikroozet}[§0.4 özeti]
Tamtümlük aksiyomu, her sınırlı kümenin $\R$'de supremum/infimum
sahibi olduğunu garanti eder.
$\Rbar = \R \cup \{-\infty, +\infty\}$,
özellikle $0 \cdot \infty = 0$ konvansiyonuyla, ölçüm ve olasılık teorisinin
standart çalışma alanıdır.
\end{mikroozet}


% =====================================================================
\section{Üst Limit ve Alt Limit}
% =====================================================================

\begin{motivasyon}[§0.5]
Bir dizi yakınsamıyorsa bile $\limsup$ ve $\liminf$ her zaman tanımlıdır
($\Rbar$'da değerleri mevcuttur). Olasılık teorisinde Borel-Cantelli
lemmalarını (Bölüm~7) ifade etmek için küme dizilerinin $\limsup$'ını,
büyük sayılar kanununu anlamak için ise sayı dizilerinin $\limsup$'ını
sıklıkla kullanacağız.
\end{motivasyon}

\subsection{Sayı Dizilerinde Üst ve Alt Limit}

\begin{tanim}[$\limsup$ ve $\liminf$ --- Sayı Dizileri]\label{def:limsup_sayi}
$(a_n)_{n \geq 1} \subseteq \Rbar$ bir dizi olsun. Tanımlarız:
\[
  \limsup_{n \to \infty} a_n
  \coloneqq \inf_{n \geq 1} \sup_{k \geq n} a_k
  = \lim_{n\to\infty} \sup_{k \geq n} a_k,
\]
\[
  \liminf_{n \to \infty} a_n
  \coloneqq \sup_{n \geq 1} \inf_{k \geq n} a_k
  = \lim_{n\to\infty} \inf_{k \geq n} a_k.
\]
İki limit $\Rbar$'da her zaman mevcuttur; çünkü $s_n \coloneqq \sup_{k \geq n} a_k$
azalan, $t_n \coloneqq \inf_{k \geq n} a_k$ artan bir dizidir.
\end{tanim}

\begin{onerme}[$\limsup$ ve $\liminf$'in Temel Özellikleri]\label{prop:limsup_ozellik}
Her $(a_n) \subseteq \Rbar$ için:
\begin{enumerate}[label=(\roman*)]
  \item $\liminf_n a_n \leq \limsup_n a_n$,
  \item $\lim_n a_n$ mevcutsa (uzun anlamda $\Rbar$'da)
        $\lim_n a_n = \liminf_n a_n = \limsup_n a_n$.
\end{enumerate}
\end{onerme}

\begin{proof}
(i) $n$ sabit alındığında $\inf_{k \geq n} a_k \leq \sup_{k \geq n} a_k$.
$n \to \infty$ alındığında sol tarafın limiti $\liminf$, sağ tarafın limiti
$\limsup$'tır; eşitsizlik korunur.

(ii) $\lim_n a_n = L$ olsun. Her $\varepsilon > 0$ için yeterince büyük
$n$'de $|a_k - L| < \varepsilon$ (her $k \geq n$). Bu,
$\sup_{k \geq n} a_k \leq L + \varepsilon$ ve
$\inf_{k \geq n} a_k \geq L - \varepsilon$ verir; $\varepsilon \to 0$
ile sonuç elde edilir.
\end{proof}

\begin{ornek}[$\limsup$ ve $\liminf$ hesabı]
$a_n = (-1)^n$ için: $\sup_{k \geq n} a_k = 1$ ve $\inf_{k \geq n} a_k = -1$
tüm $n$ için, dolayısıyla
$\limsup_n a_n = 1$ ve $\liminf_n a_n = -1$.
Dizi yakınsamadığından $\lim_n a_n$ mevcut değil, ama $\limsup/\liminf$
her zaman tanımlıdır.
\end{ornek}

\subsection{Küme Dizilerinde Üst ve Alt Limit}

Bu alt bölüm Bölüm~7'deki Borel-Cantelli lemmalarının doğrudan ön koşuludur.

\begin{tanim}[$\limsup$ ve $\liminf$ --- Küme Dizileri]\label{def:limsup_kume}
$(A_n)_{n \geq 1}$ bir küme dizisi olsun. Tanımlarız:
\begin{align*}
  \limsup_{n\to\infty} A_n
  &\coloneqq \bigcap_{n=1}^\infty \bigcup_{k=n}^\infty A_k, \\[4pt]
  \liminf_{n\to\infty} A_n
  &\coloneqq \bigcup_{n=1}^\infty \bigcap_{k=n}^\infty A_k.
\end{align*}
\end{tanim}

\begin{tanimgerekce}
Bu tanımların sezgisel yorumları aşağıdaki önerme ile netleşir.
``Sonsuz kez'' ve ``eninde sonunda'' dili, olasılık teorisindeki olayların
geleceğe dair yorumunu doğrudan yansıtır.
\end{tanimgerekce}

\begin{onerme}[Küme $\limsup$ ve $\liminf$'in Yorumu]\label{prop:limsup_kume}
$(A_n)$ bir küme dizisi olsun. O zaman:
\begin{align*}
  \limsup_{n} A_n
  &= \{\omega : \omega \in A_n \text{ sonsuz sayıda } n \text{ için}\}, \\[4pt]
  \liminf_{n} A_n
  &= \{\omega : \omega \in A_n \text{ eninde sonunda her } n \text{ için}\}.
\end{align*}
Olasılık literatüründe bu kümeler sırasıyla ``s.s.~çok'' (\emph{infinitely often}, i.o.)
ve ``e.s.'' (\emph{eventually}) olay kümeleri olarak adlandırılır.
\end{onerme}

\begin{proof}
$\omega \in \limsup_n A_n = \bigcap_n \bigcup_{k \geq n} A_k$ olsun.
Bu, her $n$ için $\omega \in \bigcup_{k \geq n} A_k$ anlamına gelir; yani
her $n$ için bazı $k \geq n$ ile $\omega \in A_k$. Dolayısıyla sonsuz
sayıda $k$ için $\omega \in A_k$ (yoksa bir $n_0$ bulunurdu ki $k \geq n_0$
için $\omega \notin A_k$, bu $\omega \notin \bigcup_{k \geq n_0} A_k$,
çelişki).

Tersine, $\omega$ sonsuz sayıda $n$'de $A_n$'deyse, her $N$ için $N$'den büyük
bir $k$ ile $\omega \in A_k$ vardır; dolayısıyla $\omega \in \bigcup_{k \geq N} A_k$
her $N$ için, yani $\omega \in \bigcap_N \bigcup_{k \geq N} A_k$.

\medskip

$\liminf$ için: $\omega \in \liminf_n A_n = \bigcup_n \bigcap_{k \geq n} A_k$
iff bazı $n_0$ için her $k \geq n_0$'da $\omega \in A_k$,
yani $\omega$ eninde sonunda her $A_n$'dedir.
\end{proof}

\begin{teorem}[$\liminf \subseteq \limsup$]\label{thm:liminf_subseteq}
Her küme dizisi $(A_n)$ için:
\[
  \liminf_{n} A_n \subseteq \limsup_{n} A_n.
\]
Dizi monoton artanlarsa ya da monoton azalanlarsa, $\liminf = \limsup$.
\end{teorem}

\begin{proof}
``Eninde sonunda her $A_n$'de olmak'' kesinlikle ``sonsuz sayıda $A_n$'de olmak''ı
gerektirir; dolayısıyla $\liminf \subseteq \limsup$.

Dizi artanlarsa ($A_n \nearrow$): $\omega \in \limsup_n A_n$ olsun; sonsuz çok $n$
için $\omega \in A_n$, dolayısıyla bazı $n_0$ için $\omega \in A_{n_0}$, ve
$A_{n_0} \subseteq A_k$ her $k \geq n_0$ için. Yani $\omega \in \liminf_n A_n$.
Azalan dizi benzer argümanla ele alınır.
\end{proof}

\begin{ispatiskele}[Teorem~\ref{thm:liminf_subseteq} --- azalan dizi]
Dizi azalanlarsa ($A_n \searrow$) $\liminf = \limsup$ olduğunu kanıtlayın:

\iskeleadim{1}{$A_n \searrow$ olduğundan, her $n$ için $\bigcap_{k \geq n} A_k = \bosluk{A_n}$'deki en küçük eleman, yani $\bigcap_{k \geq n} A_k = \bigcap_{k=1}^\infty A_k$.}

\iskeleadim{2}{$\liminf_n A_n = \bigcup_n \bigcap_{k \geq n} A_k = \bosluk{\bigcap_{k=1}^\infty A_k}$.}

\iskeleadim{3}{Benzer argümanla $\limsup_n A_n = \bosluk{\bigcap_{k=1}^\infty A_k}$, yani iki küme eşittir.}
\end{ispatiskele}

\begin{ornek}[$\limsup$ ve $\liminf$ küme hesabı]\label{ex:limsupkume}
$\Omega = \R$ ve
\[
  A_n =
  \begin{cases}
    (0,1), & n \text{ tek}, \\
    (0,2), & n \text{ çift}.
  \end{cases}
\]
O zaman:
\[
  \limsup_n A_n = (0,2),
  \qquad
  \liminf_n A_n = (0,1).
\]
Çünkü her $x \in (0,2)$ sonsuz sayıda çift indisli $A_n$'de yer alır;
yalnızca $(0,1)$ kümesi eninde sonunda tüm $A_n$'lerde yer alır.
\end{ornek}

\begin{kendinleifade}
$(A_n)$ ile $B_n = A_n^c$ arasındaki $\limsup$ ve $\liminf$ ilişkisini
De~Morgan yasasını kullanarak türetin. Beklenen sonuç:
$\left(\limsup_n A_n\right)^c = \liminf_n A_n^c$.
\ogrencialani[2.5cm]
\end{kendinleifade}

\begin{mikroozet}[§0.5 özeti]
$\limsup_n A_n$ = sonsuz sık olay; $\liminf_n A_n$ = eninde sonunda
gerçekleşen olay. Her zaman $\liminf \subseteq \limsup$. Monoton dizilerde
ikisi eşittir. Bu dil Bölüm~7'deki Borel-Cantelli lemmalarının tam dilidir.
\end{mikroozet}


% =====================================================================
%  BÖLÜM SONU ÖZETİ
% =====================================================================
\section*{Bölüm Özeti}
\addcontentsline{toc}{section}{Bölüm Özeti}

\noindent
\begin{tcolorbox}[maarif@base, colframe=maOzet, colback=maOzet!4,
    colbacktitle=maOzet, title={Bölüm 0 --- Temel Sonuçlar}]
\begin{itemize}[leftmargin=1.8em, itemsep=3pt]
  \item \textbf{De~Morgan:}
        $\left(\bigcup_n A_n\right)^c = \bigcap_n A_n^c$ ve
        $\left(\bigcap_n A_n\right)^c = \bigcup_n A_n^c$.
  \item \textbf{Ters görüntü:} $f^{-1}$ birleşim, kesişim ve tümlemeyi korur
        (Teorem~\ref{thm:tergoriuntu}).
  \item \textbf{Sayılabilirlik:} $\Q$ sayılabilir; $\R$ sayılamaz (Cantor köşegen argümanı).
        Sayılabilir ailenin birleşimi sayılabilirdir.
  \item \textbf{Tamtümlük:} Sınırlı her boşalmayan $A \subseteq \R$ kümesinin
        $\sup A$ ve $\inf A$ değerleri $\R$'de mevcuttur.
  \item \textbf{$\Rbar$:} $0 \cdot (\pm\infty) = 0$ konvansiyonuyla kullanılır;
        $(+\infty)+(-\infty)$ tanımsız.
  \item \textbf{$\limsup / \liminf$ (kümeler):}
        $\limsup_n A_n$ = s.s.~çok; $\liminf_n A_n$ = eninde~sonunda.
        $\left(\limsup_n A_n\right)^c = \liminf_n A_n^c$.
\end{itemize}

\medskip\noindent\textbf{Notasyon özeti:}
\begin{itemize}[leftmargin=1.8em, itemsep=2pt, label={--}]
  \item $A^c$: tümleyen;\quad $A \triangle B$: simetrik fark;\quad $\powerset{\Omega}$: kuvvet kümesi
  \item $f^{-1}(B)$: ters görüntü
  \item $\sup A$, $\inf A$: üst/alt kenar değer;\quad
        $\Rbar = \R \cup \{-\infty, +\infty\}$
  \item $\limsup_n A_n = \bigcap_n \bigcup_{k \geq n} A_k$;\quad
        $\liminf_n A_n = \bigcup_n \bigcap_{k \geq n} A_k$
\end{itemize}
\end{tcolorbox}

% ─── ÖZ-YANSITMA ──────────────────────────────────────────────────
\begin{ozyansima}
Bu bölümü tamamladıktan sonra kendinize şu soruları sorun:
\begin{itemize}[leftmargin=1.8em, itemsep=3pt]
  \item De~Morgan yasasını kaynağa bakmadan ispatlayabilir misiniz?
  \item Ters görüntünün neden tümlemeyi koruduğunu, bir arkadaşınıza sözlü
        açıklayabilir misiniz?
  \item Cantor'un köşegen argümanında hangi adım çelişkiyi üretir?
  \item $\limsup_n A_n$ ile $\liminf_n A_n$ arasındaki farkı tek cümleyle
        ifade edebilir misiniz?
\end{itemize}
\end{ozyansima}

% ─── KÖPRÜ NOTU ───────────────────────────────────────────────────
\begin{kopru}
\textbf{Sonraki bölüm (Bölüm~1: Ölçüm Teorisinin Temelleri).}
Küme cebiri bilgisi doğrudan $\sigma$-cebiri tanımına gidecek;
sayılabilirlik orada ``sayılabilir birleşim altında kapalı'' koşulunda karşınıza
çıkacak. Ters görüntünün küme işlemlerini koruması ise ölçülebilir fonksiyon
tanımının kalbidir.

\medskip

\textbf{İleri okuma.}
Billingsley \cite{billingsley1995} Ek~A;
Klenke \cite{klenke2014} Bölüm~1.1;
Durrett \cite{durrett2019} Ek~A.1.
\end{kopru}


% =====================================================================
%  ALIŞTIRMALAR
% =====================================================================

\cozumlualistirmalar

% ─── Çözümlü Alıştırma 0.1 ───────────────────────────────────────
\begin{alistirma}[Al.0.1]{A}{De~Morgan yasası --- doğrudan uygulama}
$\Omega = \{1,2,3,4,5,6\}$ (zar atma örnek uzayı) ve $A_n = \{k \in \Omega : k \leq n+1\}$,
$n = 1,2,3$ olsun. Hesaplayın:
(a)~$\bigcup_{n=1}^{3} A_n$,\quad
(b)~$\left(\bigcup_{n=1}^{3} A_n\right)^c$,\quad
(c)~$\bigcap_{n=1}^{3} A_n^c$.
İki sonucun eşit olduğunu doğrulayın.
\end{alistirma}
\alcozum{%
(a) $A_1 = \{1,2\}$, $A_2 = \{1,2,3\}$, $A_3 = \{1,2,3,4\}$.
$\bigcup_{n=1}^{3} A_n = \{1,2,3,4\}$.

(b) $\left(\bigcup_n A_n\right)^c = \Omega \setminus \{1,2,3,4\} = \{5,6\}$.

(c) $A_1^c = \{3,4,5,6\}$, $A_2^c = \{4,5,6\}$, $A_3^c = \{5,6\}$.
$\bigcap_n A_n^c = \{5,6\}$.

İki küme eşittir; De~Morgan yasası doğrulanmış oldu. $\blacksquare$%
}

% ─── Çözümlü Alıştırma 0.2 ───────────────────────────────────────
\begin{alistirma}[Al.0.2]{B}{Ters görüntü --- aralıklar}
$f: \R \to \R$, $f(x) = x^2$ olsun.
(a)~$f^{-1}([0,4])$'ü hesaplayın.
(b)~$f^{-1}((-1,0))$'ı hesaplayın ve yorumlayın.
(c)~$B_1 = [1,4]$ ve $B_2 = [4,9]$ için $f^{-1}(B_1 \cap B_2)$'yi
    hem doğrudan hem de $f^{-1}(B_1) \cap f^{-1}(B_2)$ olarak hesaplayın;
    eşit olduklarını doğrulayın.
\end{alistirma}
\alcozum{%
(a) $f^{-1}([0,4]) = \{x : x^2 \in [0,4]\} = [-2,2]$.

(b) $f^{-1}((-1,0)) = \{x : x^2 \in (-1,0)\} = \emptyset$.
Yorum: $f(x) = x^2 \geq 0$ olduğundan görüntü $[0,\infty)$;
negatif değerlerin ters görüntüsü boştur.

(c) $B_1 \cap B_2 = \{4\}$, dolayısıyla $f^{-1}(\{4\}) = \{-2,2\}$.

$f^{-1}(B_1) = [-2,-1] \cup [1,2]$ ve $f^{-1}(B_2) = [-3,-2] \cup [2,3]$.
$f^{-1}(B_1) \cap f^{-1}(B_2) = \{-2,2\}$.

İki sonuç eşittir: Teorem~\ref{thm:tergoriuntu}(ii) doğrulandı. $\blacksquare$%
}

% ─── Çözümlü Alıştırma 0.3 ───────────────────────────────────────
\begin{alistirma}[Al.0.3]{B}{Küme $\limsup$ ve $\liminf$ --- azalan dizi}
$\Omega = [0,1]$ ve $A_n = [0, 1/n]$ olsun.
(a)~$\limsup_n A_n$'i hesaplayın.
(b)~$\liminf_n A_n$'i hesaplayın.
(c)~Dizi monoton mu? Cevabınızı Teorem~\ref{thm:liminf_subseteq} ile
    ilişkilendirin.
\end{alistirma}
\alcozum{%
(a) Dizi azalan: $A_{n+1} \subset A_n$.
$\bigcup_{k \geq n} A_k = A_n = [0,1/n]$ (en büyük eleman $A_n$).
$\limsup_n A_n = \bigcap_{n \geq 1} [0,1/n] = \{0\}$.

(b) $\bigcap_{k \geq n} A_k = \bigcap_{k \geq n} [0,1/k] = \{0\}$ her $n$ için.
$\liminf_n A_n = \bigcup_{n \geq 1} \{0\} = \{0\}$.

(c) Dizi azalan; Teorem~\ref{thm:liminf_subseteq}'ye göre
$\limsup = \liminf = \{0\}$, yani $\lim_n A_n = \{0\}$.
Sezgisel: aralıklar gittikçe küçülür, tek ortak nokta $\{0\}$ kalır. $\blacksquare$%
}

% ─── Çözümlü Alıştırma 0.4 ───────────────────────────────────────
\begin{alistirma}[Al.0.4]{A}{$\Z$'nin sayılabilirliği}
$\Z$'nin sayılabilir sonsuz olduğunu, açık bir bijeksiyon
$f: \N \to \Z$ yazarak kanıtlayın. ($\N = \{1,2,3,\ldots\}$.)
\end{alistirma}
\alcozum{%
$f(n) = \begin{cases} n/2 & n \text{ çift}, \\ -(n-1)/2 & n \text{ tek}\end{cases}$
tanımlayın. Değerler: $f(1)=0$, $f(2)=1$, $f(3)=-1$, $f(4)=2$, $f(5)=-2$, \ldots

\emph{Enjektif:} $f(m) = f(n)$ olsun. $m,n$ ikisi de çift ise $m/2 = n/2$, yani $m=n$.
İkisi de tek ise $(m-1)/2 = (n-1)/2$, yani $m=n$. Biri çift biri tek olamaz
(biri pozitif, diğeri sıfır ya da negatif verir; çelişki).

\emph{Sürjektif:} $z > 0$ ise $n = 2z$ çift, $f(2z) = z$; $z \leq 0$ ise
$n = 1-2z$ tek, $f(1-2z) = z$.

$f$ bijektiftir; dolayısıyla $\Z$ sayılabilir sonsuzdur. $\blacksquare$%
}

\alistirmalar

% ─── A Seviyesi ───────────────────────────────────────────────────
\begin{alistirma}[Al.0.5]{A}{Küme özdeşlikleri}
Aşağıdakilerin doğru mu yanlış mı olduğunu belirleyin. Doğruysa ispatlayın;
yanlışsa karşı örnek verin.
\begin{enumerate}[label=(\alph*)]
  \item $A \setminus (B \cup C) = (A \setminus B) \cap (A \setminus C)$
  \item $A \triangle B = B \triangle A$
  \item $A \cap (B \triangle C) = (A \cap B) \triangle (A \cap C)$
  \item $(A \triangle B) \triangle C = A \triangle (B \triangle C)$
\end{enumerate}
\end{alistirma}
\alcozum{%
(a) Doğru. $A \setminus (B \cup C) = A \cap (B \cup C)^c = A \cap (B^c \cap C^c)
= (A \cap B^c) \cap (A \cap C^c) = (A \setminus B) \cap (A \setminus C)$.

(b) Doğru. Simetrik fark tanımı gereği simetriktir:
$(A \setminus B) \cup (B \setminus A) = (B \setminus A) \cup (A \setminus B)$.

(c) Doğru. Her $\omega$ için dört durum kontrol edilir;
$A \cap (B \triangle C)$ ve $(A \cap B) \triangle (A \cap C)$
her iki ifade ``$A$'da, ama $B$ ve $C$'den tam olarak birinde'' olan
elemanlar kümesidir.

(d) Doğru (simetrik fark birleştirme yasasına uyar). İspat: her $\omega$
için $A \triangle B \triangle C$ tam olarak tek sayıda kümede bulunan
elemanları içerir. $\blacksquare$%
}

\begin{alistirma}[Al.0.6]{A}{Supremum ve İnfimum}
Aşağıdaki kümelerin $\sup$ ve $\inf$ değerlerini bulun; her değerin kümede
bulunup bulunmadığını belirtin.
\begin{enumerate}[label=(\alph*)]
  \item $A = \left\{\dfrac{(-1)^n}{n} : n \in \N\right\}$
  \item $B = \{x \in \Q : x^2 < 3\}$
  \item $C = \left\{1 - \dfrac{1}{2^n} : n \in \N_0\right\}$
        ($\N_0 = \{0,1,2,\ldots\}$)
\end{enumerate}
\end{alistirma}
\alcozum{%
(a) $A = \{-1, 1/2, -1/3, 1/4, \ldots\}$. $\sup A = 1/2 \in A$;
$\inf A = -1 \in A$.

(b) $B = (-\sqrt{3}, \sqrt{3}) \cap \Q$. $\sup B = \sqrt{3} \notin B$
(irrasyonel); $\inf B = -\sqrt{3} \notin B$.

(c) $C = \{0, 1/2, 3/4, 7/8, \ldots\}$. $\sup C = 1 \notin C$
(limit değer, küme içermez); $\inf C = 0 \in C$. $\blacksquare$%
}

\begin{alistirma}[Al.0.7]{A}{$\limsup$ ve $\liminf$ --- salınan dizi}
$(a_n)$ dizisi $a_n = \sin(n\pi/2)$, $n \geq 1$ olsun.
(a)~İlk sekiz terimi yazın.
(b)~$\limsup_n a_n$ ve $\liminf_n a_n$'i hesaplayın.
\end{alistirma}
\alcozum{%
(a) $1,\,0,\,-1,\,0,\,1,\,0,\,-1,\,0$.

(b) Dizi $\{1,0,-1,0,\ldots\}$ döngüsündedir. $\sup_{k \geq n} a_k = 1$ her
$n$ için. $\limsup_n a_n = 1$. Benzer biçimde $\liminf_n a_n = -1$. $\blacksquare$%
}

\begin{alistirma}[Al.0.8]{A}{Gösterge fonksiyonu ve simetrik fark}
$\gostf{A}(\omega) = 1$ eğer $\omega \in A$, aksi hâlde $0$ olsun.
Gösterin ki:
\[
  \gostf{A \triangle B} = \gostf{A} + \gostf{B} - 2\,\gostf{A \cap B}.
\]
\end{alistirma}
\alcozum{%
Dört durum:
\begin{itemize}
  \item $\omega \in A \cap B$: sol = 0, sağ = $1+1-2 = 0$. \checkmark
  \item $\omega \in A \setminus B$: sol = 1, sağ = $1+0-0 = 1$. \checkmark
  \item $\omega \in B \setminus A$: sol = 1, sağ = $0+1-0 = 1$. \checkmark
  \item $\omega \notin A \cup B$: sol = 0, sağ = $0+0-0 = 0$. \checkmark
\end{itemize}
$\blacksquare$%
}

% ─── B Seviyesi ───────────────────────────────────────────────────
\begin{alistirma}[Al.0.9]{B}{Sayılabilir birleşim --- uygulama}
(a)~$\N \times \N$'nin sayılabilir olduğunu ispatlayın.
(b)~$\Q$'nun sayılabilirliğini Teorem~\ref{thm:sayilabilirbir}'i kullanarak
    yeniden ispatlayın.
(c)~Herhangi bir $A \subseteq B$ için $A$ sayılabilir olduğunda $B$'nin
    sayılabilmesinin \emph{yeterli} koşulunu tartışın. ($B$ sayılabilir
    olmak zorunda mı?)
\end{alistirma}
\alcozum{%
(a) $f: \N \times \N \to \N$, $f(m,n) = 2^{m-1}(2n-1)$ bijeksiyondur
(her pozitif tam sayı tek türlü $2^a \cdot b$ biçiminde yazılır, $b$ tek).
Dolayısıyla $\N \times \N \cong \N$, sayılabilirdir.

(b) $\Q^+ \subseteq \N \times \N$ gömülebilir ($p/q \mapsto (p,q)$ enjektif).
$\Q = \Q^+ \cup \{0\} \cup (-\Q^+)$, Teorem~\ref{thm:sayilabilirbir} ile sayılabilir.

(c) $A \subseteq B$ ise $A$'nın sayılabilirliği $B$'yi zorunlu olarak
sayılabilir kılmaz. Örneğin $A = \{0\}$ ve $B = \R$: $A$ sonlu,
$B$ sayılamazdır. Ama $B \subseteq A$ ile $A$ sayılabilir ise $B$
sayılabilirdir (sayılabilir bir kümenin her alt kümesi sayılabilir). $\blacksquare$%
}

\begin{alistirma}[Al.0.10]{B}{Küme $\limsup$ ve $\liminf$ --- iki dizi}
$\Omega = \R$ olsun. Her dizi için $\limsup$ ve $\liminf$'i hesaplayın.
\begin{enumerate}[label=(\alph*)]
  \item $A_n = \left(-\dfrac{1}{n},\, 1 + \dfrac{1}{n}\right)$
  \item $A_n = (n, n+1)$
\end{enumerate}
\end{alistirma}
\alcozum{%
(a) Dizi azalan: $A_{n+1} \subset A_n$. $\bigcup_{k \geq n} A_k = A_n$ ve
$\bigcap_{k \geq n} A_k \to [0,1]$.
$\limsup_n A_n = \liminf_n A_n = [0,1]$.

(b) $\bigcup_{k \geq n} A_k = (n,\infty)$ ve $\bigcap_{k \geq n} A_k = \emptyset$.
$\limsup_n A_n = \bigcap_{n \geq 1}(n,\infty) = \emptyset$.
$\liminf_n A_n = \bigcup_{n \geq 1} \emptyset = \emptyset$.
Yorum: $(n,n+1)$ sağa kayar; hiçbir sabit nokta bu dizinin sonsuz
üyesinde değildir. $\blacksquare$%
}

\begin{alistirma}[Al.0.11]{B}{Ters görüntü --- mutlak değer fonksiyonu}
$f: \R \to \R$, $f(x) = |x|$ olsun.
\begin{enumerate}[label=(\alph*)]
  \item $f^{-1}([1,2])$
  \item $f^{-1}(\{0\})$
  \item $f^{-1}(\R \setminus [0,1])$ hem doğrudan hem de
        $\left(f^{-1}([0,1])\right)^c$ olarak hesaplayın; eşitliği
        Teorem~\ref{thm:tergoriuntu}(iii) ile ilişkilendirin.
\end{enumerate}
\end{alistirma}
\alcozum{%
(a) $f^{-1}([1,2]) = \{x : |x| \in [1,2]\} = [-2,-1] \cup [1,2]$.

(b) $f^{-1}(\{0\}) = \{0\}$.

(c) Doğrudan: $\R \setminus [0,1] = (-\infty,0) \cup (1,\infty)$;
$f^{-1} = (-\infty,-1) \cup (1,\infty)$ (çünkü $|x| > 1 \iff x > 1$ ya da $x < -1$).

$(f^{-1}([0,1]))^c = [-1,1]^c = (-\infty,-1) \cup (1,\infty)$.

İki sonuç eşittir; Teorem~\ref{thm:tergoriuntu}(iii) doğrulandı. $\blacksquare$%
}

\begin{alistirma}[Al.0.12]{B}{$\limsup$ tümleme özdeşliği --- ispat}
$(A_n)$ bir küme dizisi için $\left(\limsup_n A_n\right)^c = \liminf_n A_n^c$
olduğunu ispatlayın.
\end{alistirma}
\alcozum{%
De~Morgan yasasını art arda uygulayın:
\begin{align*}
  \left(\limsup_n A_n\right)^c
  &= \left(\bigcap_{n=1}^\infty \bigcup_{k \geq n} A_k\right)^c \\
  &= \bigcup_{n=1}^\infty \left(\bigcup_{k \geq n} A_k\right)^c \\
  &= \bigcup_{n=1}^\infty \bigcap_{k \geq n} A_k^c \\
  &= \liminf_n A_n^c. \qquad \blacksquare
\end{align*}
Sezgisel yorum: ``$\omega$ sonsuz sık $A_n$'de \emph{değil}'' $\iff$
``$\omega$ eninde sonunda her $A_n^c$'de.'' $\blacksquare$%
}

% ─── C Seviyesi ───────────────────────────────────────────────────
\begin{alistirma}[Al.0.13]{C}{Cantor kümesinin sayılamazlığı}
\textbf{Cantor kümesi} $\mathcal{C}$, $[0,1]$'den ardı ardına orta üçte bir
açık aralıklar çıkarılarak elde edilir: önce $(1/3,2/3)$, ardından
$(1/9,2/9)$ ve $(7/9,8/9)$, vb.

(a)~$\mathcal{C}$'nin Lebesgue ölçüsünün $0$ olduğunu gösterin
(çıkarılan açık aralıkların uzunluk toplamı $1$'dir).

(b)~$\mathcal{C}$, üçlü taban gösteriminde yalnızca $0$ ve $2$ basamağı
içeren sayıların kümesidir. Bu gözlemi kullanarak $\mathcal{C}$'nin
sayılamaz olduğunu kanıtlayın.
\end{alistirma}
\alcozum{%
\begin{ipucu}{1}
$n$-inci adımda $2^{n-1}$ adet, her biri $1/3^n$ uzunluklu aralık çıkarılır.
\end{ipucu}

(a) Toplam çıkarılan uzunluk:
$\sum_{n=1}^\infty 2^{n-1}/3^n = \tfrac{1}{2}\sum_{n=1}^\infty (2/3)^n
= \tfrac{1}{2} \cdot \tfrac{2/3}{1/3} = 1$.

(b) Her $x \in \mathcal{C}$, $x = \sum_{n=1}^\infty d_n/3^n$ ($d_n \in \{0,2\}$)
biçiminde yazılır. $g: \mathcal{C} \to [0,1]$, $g(x) = \sum_{n=1}^\infty (d_n/2)/2^n$
sürjektiftir (her ikili gösterim elde edilir). $[0,1]$ sayılamaz; sürjektif
görüntünün alt kümesi sayılamaz olamaz; dolayısıyla $\mathcal{C}$ sayılamazdır.
$\blacksquare$%
}

\begin{alistirma}[Al.0.14]{C}{Alt dizi ve $\limsup$/$\liminf$ ilişkisi}
$(A_n)$ bir küme dizisi ve $(A_{n_k})_{k \geq 1}$, $(A_n)$'nin bir alt dizisi
($n_1 < n_2 < \cdots$) olsun. Kanıtlayın ki:
\[
  \liminf_n A_n \subseteq \liminf_k A_{n_k}
  \subseteq \limsup_k A_{n_k} \subseteq \limsup_n A_n.
\]
\end{alistirma}
\alcozum{%
\begin{ipucu}{1}
$\limsup_k A_{n_k} \subseteq \limsup_n A_n$ için: $\omega$'nın sonsuz çok $A_{n_k}$'da
olması, sonsuz çok $A_n$'da olması anlamına gelir (alt dizi endeksleri de $\N$'dedir).
\end{ipucu}

\textbf{$\limsup_k A_{n_k} \subseteq \limsup_n A_n$:}
$\omega$ sonsuz çok $k$ için $A_{n_k}$'da $\Rightarrow$ sonsuz çok $n \in \{n_1,n_2,\ldots\}$
için $\omega \in A_n$ $\Rightarrow$ $\omega \in \limsup_n A_n$.

\textbf{$\liminf_n A_n \subseteq \liminf_k A_{n_k}$:}
$\omega \in \liminf_n A_n$ $\Rightarrow$ bazı $N$ için $k \geq N \Rightarrow \omega \in A_k$.
$n_k \geq k$ olduğundan, yeterince büyük $k$ için $n_k \geq N$ ve $\omega \in A_{n_k}$;
yani $\omega \in \liminf_k A_{n_k}$.

Ortadaki $\liminf \subseteq \limsup$ Teorem~\ref{thm:liminf_subseteq}'den gelir. $\blacksquare$%
}

\begin{alistirma}[Al.0.15]{C}{Simetrik fark ve ölçü}
$(X, \mathcal{A}, \mu)$ bir ölçüm uzayı, $A, B \in \mathcal{A}$ ve $\mu$ sonlu
ölçü olsun. Gösterin ki $\mu(A \triangle B) = \mu(A) + \mu(B) - 2\mu(A \cap B)$.
Bu ifade neden bir ``küme metriği'' için temel oluşturur?
\end{alistirma}
\alcozum{%
$A \triangle B = (A \setminus B) \cup (B \setminus A)$ ve bu iki küme ayrık.
\begin{align*}
  \mu(A \triangle B)
  &= \mu(A \setminus B) + \mu(B \setminus A) \\
  &= [\mu(A) - \mu(A \cap B)] + [\mu(B) - \mu(A \cap B)] \\
  &= \mu(A) + \mu(B) - 2\mu(A \cap B).
\end{align*}
Metrik: $d(A,B) = \mu(A \triangle B)$ tanımlanırsa, $d(A,A)=0$, $d(A,B)=d(B,A)$
ve üçgen eşitsizliği $A \triangle C \subseteq (A \triangle B) \cup (B \triangle C)$
özdeşliğinden gelir. Bu ``ölçü teorik metrik'', $L^1$ normunun gösterge
fonksiyonları üzerindeki karşılığıdır. $\blacksquare$%
}

% =====================================================================
%  BÖLÜM 1 — Ölçüm Teorisinin Temelleri
%  Kaynaklar: Billingsley Böl.2-3; Klenke Böl.1-6; Shiryaev II.§1-§3
% =====================================================================

\chapter{Ölçüm Teorisinin Temelleri}

\bolumalinti{``The theory of measure is the foundation on which
  the modern theory of probability rests.''}{Andrey Kolmogorov (1903--1987)}

\noindent\textbf{Ön Koşullar:} Bölüm~0 (küme cebiri, sayılabilirlik, $\limsup/\liminf$).

\noindent\textbf{Bu Bölüm Sonraki Bölümlerde Kullanılır:}
Bölüm~2 (olasılık ölçüsü),
Bölüm~3 (rasgele değişkenler = ölçülebilir fonksiyonlar),
Bölüm~4 (beklenti = Lebesgue integrali).

\noindent\textbf{Tahmini Okuma Süresi:} $\approx 10$~saat.

\vspace{0.5\baselineskip}

\begin{kazanimlar}
Bu bölümün sonunda öğrenci:
\begin{itemize}
  \item $\sigma$-cebiri tanımını verebilmeli, örnekler üretebilmeli ve
        verilen bir aile tarafından üretilen $\sigma$-cebirini betimleyebilmeli,
  \item Borel $\sigma$-cebirini $\mathcal{B}(\mathbb{R})$ tanımlayabilmeli ve
        açık kümelerin, kapalı kümelerin, $F_\sigma$ ile $G_\delta$ kümelerinin
        Borel olduğunu gösterebilmeli,
  \item Ölçü aksiyomlarını ifade edebilmeli; monotonluk, alt-toplamsallık ve
        aşağıdan/yukarıdan süreklilik özelliklerini ispatlayabilmeli,
  \item Carathéodory genişleme teoremini ifade edebilmeli ve Lebesgue ölçüsünün
        bu teoremden nasıl elde edildiğini açıklayabilmeli,
  \item Ölçülebilir fonksiyon kavramını tanımlayabilmeli; basit fonksiyon
        yaklaşımını ve Lebesgue integralinin yapısını açıklayabilmeli,
  \item Monoton Yakınsama Teoremi, Fatou Lemması ve Dominante Yakınsama
        Teoremini ispatlayabilmeli ve bunları integral hesaplarında uygulayabilmeli,
  \item Fubini-Tonelli teoremini ifade edebilmeli ve çift integrallerin sırasını
        değiştirme koşulunu açıklayabilmeli.
\end{itemize}
\end{kazanimlar}

\begin{motivasyon}
Olasılık teorisinin Kolmogorov tarafından 1933'te aksiyomatik biçimde
kurulmadan önce, ``olasılık'' kavramı kesin matematiksel temellere sahip
değildi. Hangi olaylara olasılık atanabilir? Sonsuz çok olayı kapsayan
deneyler nasıl ele alınır? Lebesgue integralinin beklentiyle ilişkisi nedir?

Tüm bu soruların yanıtı ölçüm teorisinde yatar. Bu bölümde önce ``olası
olayların ailesinin'' hangi yapıya sahip olması gerektiğini ($\sigma$-cebiri),
ardından bu yapıya sayı atayan fonksiyonun özelliklerini (ölçü) inceleyeceğiz.

\medskip

\textbf{Köprü kurma.} Riemann integralini biliyorsunuz: fonksiyonu parçalara böler,
dikdörtgen alanları toplarsınız. Lebesgue integrali tam tersi yönde çalışır:
fonksiyonun değer aralığını böler, her değer kümesinin ``ölçüsünü'' hesaplar.
Bu fark küçük görünse de, Lebesgue integrali sınır alma ile değiştirilebilirlik
gibi son derece güçlü teoremler sağlar; Riemann integrali bunu sağlayamaz.
\end{motivasyon}

\begin{tarihselnot}
Henri Lebesgue, 1901 yılında doktora tezinde yeni bir integral teorisi kurdu.
Émile Borel daha önce belirli küme ailelerine ``ölçü'' atamayı düşünmüştü;
Lebesgue bu fikri sistematikleştirdi. Carathéodory ise 1914'te dış ölçü
yoluyla genişleme teoremini formüle etti. Kolmogorov, 1933'te tüm bu yapıyı
olasılık teorisine uyguladı: olasılık uzayı $(\Omega, \mathcal{F}, \mathbb{P})$
modern matematiksel olasılığın temel nesnesi oldu.
\end{tarihselnot}


% =====================================================================
\section{$\sigma$-Cebirleri ve Borel Kümeler}
% =====================================================================

\begin{kesif}[$\sigma$-cebiri neden gerekli?]
Bir madeni para \emph{sonsuz kez} atılıyor. ``İlk tura $n$-inci atışta gelir''
olayını $A_n$ ile gösterin. ``En az bir kez tura gelir'' olayı nedir?
Cevabınızı aşağıya yazın.
\kesifcevap{``En az bir kez tura gelir'' = $\bigcup_{n=1}^{\infty} A_n$.
  Bu sonsuz birleşimi içeren her aile $\sigma$-cebiri olmak zorundadır;
  aksi takdirde bu olaya olasılık atamak tutarsızlığa yol açar.
  $\sigma$-cebirinin kapalılık koşulu buradan doğrudan çıkar.}
Olasılık tanımlamak istediğimiz olaylar ailesi bu tür işlemler altında kapalı
olmalıdır. Bu koşulun kesin matematiksel karşılığı $\sigma$-cebiriyidir.
\end{kesif}

\begin{tanimgerekce}
Neden sadece sonlu işlemler değil, \emph{sayılabilir} işlemler istiyoruz?
Çünkü sonsuz deneyler (bir madeni paranın sonsuz kez atılması, Poisson süreci)
doğal olarak sayılabilir sonsuz birleşimler üretir. Sayılamaz birleşimlere
izin versek, her tekil kümenin birleşimi tüm $\mathbb{R}$'yi verir; bu kümeler
üzerinde Lebesgue ölçüsü gibi makul bir ölçü tanımlamak imkânsız olur.
\end{tanimgerekce}

\begin{tanim}[$\sigma$-Cebiri]\label{def:sigmacebiri}
$\Omega$ boşalmayan bir küme olsun. $\mathcal{F} \subseteq \mathcal{P}(\Omega)$
alt kümeler ailesine \textbf{$\sigma$-cebiri} \emph{($\sigma$-algebra,
$\sigma$-field)} denir; eğer:
\begin{enumerate}[label=\textbf{(\Alph*)}]
  \item $\Omega \in \mathcal{F}$,
  \item $A \in \mathcal{F} \Rightarrow A^c \in \mathcal{F}$ \quad(tümleme altında kapalı),
  \item $A_1, A_2, \ldots \in \mathcal{F} \Rightarrow \bigcup_{n=1}^\infty A_n \in \mathcal{F}$
        \quad(sayılabilir birleşim altında kapalı).
\end{enumerate}
$(\Omega, \mathcal{F})$ çiftine \textbf{ölçülebilir uzay} \emph{(measurable space)} denir.
\end{tanim}

\begin{sonuc}[Türetilen Özellikler]\label{cor:sigmaturev}
$\mathcal{F}$ bir $\sigma$-cebiri ise:
\begin{enumerate}[label=(\roman*)]
  \item $\emptyset \in \mathcal{F}$,
  \item $A_1, A_2, \ldots \in \mathcal{F} \Rightarrow \bigcap_{n=1}^\infty A_n \in \mathcal{F}$,
  \item $A, B \in \mathcal{F} \Rightarrow A \setminus B \in \mathcal{F}$,
  \item $A, B \in \mathcal{F} \Rightarrow A \triangle B \in \mathcal{F}$.
\end{enumerate}
\end{sonuc}

\begin{proof}
(i) $\Omega \in \mathcal{F}$'den (A), tümlemeyle $\emptyset = \Omega^c \in \mathcal{F}$.
(ii) De~Morgan: $\bigcap_n A_n = \left(\bigcup_n A_n^c\right)^c$;
her $A_n^c \in \mathcal{F}$ (B'den), birleşim $\mathcal{F}$'de (C), tümleyen $\mathcal{F}$'de (B).
(iii) $A \setminus B = A \cap B^c$; (ii)'nin özel hali.
(iv) $(A \setminus B) \cup (B \setminus A)$; (iii) ve (C)'den.
\end{proof}

\begin{ornek}[$\sigma$-cebiri örnekleri]\label{ex:sigmaoernekler}
\begin{enumerate}[label=(\alph*)]
  \item \textbf{Trivial (önemsiz) $\sigma$-cebiri:} $\mathcal{F} = \{\emptyset, \Omega\}$.
        Her $\Omega$ üzerinde en küçük $\sigma$-cebiriyidir.
  \item \textbf{Kuvvet kümesi:} $\mathcal{F} = \mathcal{P}(\Omega)$.
        En büyük $\sigma$-cebiriyidir.
  \item $\Omega = \{1,2,3,4\}$, $\mathcal{F} = \{\emptyset, \{1,2\}, \{3,4\}, \Omega\}$.
        Doğrulayın: (A), (B), (C) sağlanıyor.
  \item $\mathcal{F} = \{A \subseteq \Omega : A \text{ sayılabilir ya da } A^c \text{ sayılabilir}\}$.
        $\Omega$'nın herhangi bir güçte olabileceği bu ``sayılabilir-kosayılabilir''
        $\sigma$-cebiri standart bir örnektir.
\end{enumerate}
\end{ornek}

\begin{kritiksorgu}[Sonlu aileler yeterli değil mi?]
Tanımda sayılabilir birleşimi sonlu birleşimle değiştirseydik, elde edeceğimiz
yapıya \textbf{cebir} \emph{(algebra)} denir. Bir cebir, $\sigma$-cebiri olmak
zorunda değildir: $\mathcal{F}_0 = \{A \subseteq [0,1] : A \text{ ya da } A^c
\text{ sonlu birleşim aralıktır}\}$ bir cebirdir ama $\sigma$-cebiri değildir.
\end{kritiksorgu}

\subsection{Üretilmiş $\sigma$-Cebiri}

\begin{tanim}[Üretilmiş $\sigma$-Cebiri]\label{def:sigmauretilensi}
$\mathcal{E} \subseteq \mathcal{P}(\Omega)$ bir aile olsun.
$\mathcal{E}$'yi içeren tüm $\sigma$-cebirlerin kesişimine
$\mathcal{E}$ tarafından \textbf{üretilen $\sigma$-cebiri} denir ve
$\sigma(\mathcal{E})$ ile gösterilir:
\[
  \sigma(\mathcal{E}) \coloneqq \bigcap \{\mathcal{F} : \mathcal{F} \text{ bir } \sigma\text{-cebiri}, \ \mathcal{E} \subseteq \mathcal{F}\}.
\]
\end{tanim}

\begin{onerme}[$\sigma(\mathcal{E})$'nin iyi tanımlılığı]\label{prop:sigmaiyitanim}
$\sigma(\mathcal{E})$:
\begin{enumerate}[label=(\roman*)]
  \item Bir $\sigma$-cebiriyidir.
  \item $\mathcal{E}$'yi içeren en küçük $\sigma$-cebiriyidir:
        $\mathcal{E} \subseteq \mathcal{G}$ ve $\mathcal{G}$ $\sigma$-cebiri ise
        $\sigma(\mathcal{E}) \subseteq \mathcal{G}$.
\end{enumerate}
\end{onerme}

\begin{proof}
(i) $\sigma$-cebiri ailesinin keyfi kesişimi $\sigma$-cebiriyidir:
(A), (B), (C) aksiyomlarının her biri kesişimde korunur.
(ii) $\mathcal{G}$ kesişimde yer alan ailelerden biridir; dolayısıyla
kesişim $\subseteq \mathcal{G}$.
\end{proof}

\subsection{Borel $\sigma$-Cebiri}

\begin{tanim}[Borel $\sigma$-Cebiri]\label{def:borel}
$\mathbb{R}$ üzerindeki \textbf{Borel $\sigma$-cebiri}
\[
  \mathcal{B}(\mathbb{R}) \coloneqq \sigma\bigl(\{(a,b) : a < b,\ a,b \in \mathbb{R}\}\bigr)
\]
açık aralıklar ailesi tarafından üretilen $\sigma$-cebiriyidir.
$\mathcal{B}(\mathbb{R})$'nin elemanlarına \textbf{Borel kümesi} denir.
\end{tanim}

\begin{teorem}[Borel $\sigma$-Cebirinin Eşdeğer Üreticileri]\label{thm:boreldenk}
$\mathcal{B}(\mathbb{R})$ aşağıdakilerin her biri tarafından üretilir:
\begin{enumerate}[label=(\roman*)]
  \item Açık aralıklar: $\{(a,b) : a<b\}$
  \item Kapalı aralıklar: $\{[a,b] : a<b\}$
  \item Yarı-açık aralıklar: $\{(a,b] : a<b\}$\ veya\ $\{[a,b) : a<b\}$
  \item $\mathbb{R}$'nin açık kümeleri
  \item $(-\infty, b]$ biçimindeki kümeler ($b \in \mathbb{R}$)
\end{enumerate}
\end{teorem}

\begin{proof}
(v) için göstereceğiz; diğerleri benzerdir.

$\mathcal{E} = \{(-\infty, b] : b \in \mathbb{R}\}$ olsun. Önce $\mathcal{E} \subseteq \mathcal{B}(\mathbb{R})$:
\[
  (-\infty, b] = \bigcap_{n=1}^\infty (-\infty, b + 1/n)
  = \bigcap_{n=1}^\infty \Bigl[(-n, b+1/n)\Bigr]^c \cup \ldots
\]
Daha doğrudan: $(-\infty, b] = \mathbb{R} \setminus (b, \infty)$ ve
$(b,\infty) = \bigcup_{n=1}^\infty (b, b+n)$ açık aralıkların birleşimi;
dolayısıyla Borel. Tümleni de Borel.

Tersine, $(a,b) = (-\infty, b) \setminus (-\infty, a]$ ve
$(-\infty, b) = \bigcup_{n=1}^\infty (-\infty, b - 1/n]$; dolayısıyla
$\sigma(\mathcal{E})$ açık aralıkları içerir, yani $\mathcal{B}(\mathbb{R}) \subseteq \sigma(\mathcal{E})$.
İki taraflı içerme ile $\sigma(\mathcal{E}) = \mathcal{B}(\mathbb{R})$.
\end{proof}

\begin{notkutu}
Borel kümeleri çok geniş bir ailedir. Açık, kapalı, $G_\delta$ (açıkların sayılabilir
kesişimi), $F_\sigma$ (kapalıların sayılabilir birleşimi), $G_{\delta\sigma}$, \ldots
hepsi Boreldir. Buna karşın Borel olmayan Lebesgue-ölçülebilir kümeler
mevcuttur (Lebesgue $\sigma$-cebirinin Borel $\sigma$-cebirinden geniş olduğunun kanıtı).
\end{notkutu}

\begin{disiplinlerarasibaglan}[Topoloji ve Fonksiyonel Analiz]
$\mathcal{B}(\mathbb{R})$, metrik uzaylarda sürekli fonksiyonların doğal
tanım alanıdır. Bir $f: \mathbb{R} \to \mathbb{R}$ fonksiyonu sürekli ise,
her açık kümenin ters görüntüsü açıktır; dolayısıyla her Borel kümesinin
ters görüntüsü Boreldir. Bu gözlem Bölüm~3'teki ölçülebilir fonksiyon
kavramının temel örneğini oluşturur.
\end{disiplinlerarasibaglan}

\begin{mikroozet}[§1.1 özeti]
$\sigma$-cebiri: tümleme ve sayılabilir birleşim altında kapalı küme ailesi.
$\sigma(\mathcal{E})$: $\mathcal{E}$'yi içeren en küçük $\sigma$-cebiri.
$\mathcal{B}(\mathbb{R})$: açık aralıkların ya da açık kümelerin ürettiği $\sigma$-cebiri.
\end{mikroozet}


% =====================================================================
\section{Ölçüler ve Olasılık Ölçüleri}
% =====================================================================

\begin{tanimgerekce}
Bir $\sigma$-cebirine sayı atayan fonksiyon neden $\sigma$-toplamsallık
gerektiriyor? Çünkü ``uzunluk'', ``alan'', ``hacim'' ve ``olasılık'' gibi
kavramların hepsinde \emph{ayrık parçaların toplamı bütüne eşittir.} Sonlu
toplamsallık tek başına yeterli değil: sayılabilir sonsuz birleşimlerde
süreklilik özelliği istiyoruz.
\end{tanimgerekce}

\begin{tanim}[Ölçü]\label{def:olcu}
$(\Omega, \mathcal{F})$ bir ölçülebilir uzay ve
$\mu: \mathcal{F} \to [0,+\infty]$ bir fonksiyon olsun.
$\mu$'ya \textbf{ölçü} \emph{(measure)} denir; eğer:
\begin{enumerate}[label=\textbf{(\Alph*)}]
  \item $\mu(\emptyset) = 0$,
  \item \textbf{$\sigma$-toplamsallık} \emph{($\sigma$-additivity)}: $(A_n)_{n \geq 1}$
        ikili ayrık ve $A_n \in \mathcal{F}$ ise
        \[
          \mu\!\left(\bigcup_{n=1}^\infty A_n\right) = \sum_{n=1}^\infty \mu(A_n).
        \]
\end{enumerate}
$(\Omega, \mathcal{F}, \mu)$ üçlüsüne \textbf{ölçüm uzayı} \emph{(measure space)} denir.
\end{tanim}

\begin{tanim}[Olasılık Ölçüsü]\label{def:olasilikolucu}
$\mathbb{P}: \mathcal{F} \to [0,1]$ bir ölçüdür ve $\mathbb{P}(\Omega) = 1$
ise $\mathbb{P}$'ye \textbf{olasılık ölçüsü} \emph{(probability measure)} denir.
$(\Omega, \mathcal{F}, \mathbb{P})$ üçlüsüne \textbf{olasılık uzayı} denir.
\end{tanim}

\begin{notkutu}
Olasılık ölçüsü bir ölçünün özel halidir: değerleri $[0,1]$'e kısıtlı ve
tüm uzayın ölçüsü $1$'dir. Bölüm~2'de olasılık uzayını ayrıntılı inceleyeceğiz.
\end{notkutu}

\begin{teorem}[Ölçünün Temel Özellikleri]\label{thm:olcuozellik}
$(\Omega, \mathcal{F}, \mu)$ bir ölçüm uzayı olsun. O zaman:
\begin{enumerate}[label=(\roman*)]
  \item \textbf{Sonlu toplamsallık:} $A_1, \ldots, A_n$ ikili ayrık ise
        $\mu\!\left(\bigcup_{k=1}^n A_k\right) = \sum_{k=1}^n \mu(A_k)$.
  \item \textbf{Monotonluk:} $A \subseteq B \Rightarrow \mu(A) \leq \mu(B)$.
  \item \textbf{Alt-toplamsallık:} $\mu\!\left(\bigcup_{n=1}^\infty A_n\right)
        \leq \sum_{n=1}^\infty \mu(A_n)$.
  \item \textbf{Aşağıdan süreklilik:} $A_n \nearrow A \Rightarrow \mu(A_n) \nearrow \mu(A)$.
  \item \textbf{Yukarıdan süreklilik:} $A_n \searrow A$ ve $\mu(A_1) < \infty$
        $\Rightarrow \mu(A_n) \searrow \mu(A)$.
\end{enumerate}
\end{teorem}

\begin{proof}
\textbf{(i)} $A_{n+1} = A_{n+2} = \cdots = \emptyset$ koyup $\sigma$-toplamsallığı uygulayın;
$\mu(\emptyset) = 0$ kullanın.

\textbf{(ii)} $B = A \cup (B \setminus A)$ ayrık birleşim; (i)'den
$\mu(B) = \mu(A) + \mu(B \setminus A) \geq \mu(A)$.

\textbf{(iii)} $B_n = A_n \setminus \bigcup_{k=1}^{n-1} A_k$ tanımlayın.
$(B_n)$ ikili ayrık, $\bigcup B_n = \bigcup A_n$ ve $B_n \subseteq A_n$.
$\sigma$-toplamsallık ve monotonluktan:
$\mu(\bigcup A_n) = \sum \mu(B_n) \leq \sum \mu(A_n)$.

\textbf{(iv)} $B_1 = A_1$, $B_n = A_n \setminus A_{n-1}$ ($n \geq 2$) koyun.
$(B_n)$ ikili ayrık ve $A_n = \bigcup_{k=1}^n B_k$, $A = \bigcup_{k=1}^\infty B_k$.
$\sigma$-toplamsallıktan:
\[
  \mu(A) = \sum_{k=1}^\infty \mu(B_k) = \lim_{n\to\infty} \sum_{k=1}^n \mu(B_k) = \lim_{n\to\infty} \mu(A_n).
\]

\textbf{(v)} $C_n = A_1 \setminus A_n$ uygulayın; $C_n \nearrow A_1 \setminus A$.
(iv)'ten $\mu(C_n) \nearrow \mu(A_1 \setminus A) = \mu(A_1) - \mu(A)$.
$\mu(A_n) = \mu(A_1) - \mu(C_n) \to \mu(A)$.
($\mu(A_1) < \infty$ koşulu farkın tanımlı olması için gereklidir.)
\end{proof}

\begin{hataanalizi}[Yukarıdan süreklilik için sonluluk koşulu]
$\mu(A_1) < \infty$ koşulu kaldırılamaz. Karşı örnek: $\Omega = \mathbb{N}$,
$\mu$ sayma ölçüsü ($\mu(\{k\}) = 1$ her $k$ için), $A_n = \{n, n+1, n+2, \ldots\}$.
O zaman $A_n \searrow \emptyset$ ama $\mu(A_n) = \infty$ her $n$ için; $\mu(\emptyset) = 0$'a yakınsamaz.
\end{hataanalizi}

\begin{ornek}[Sayma ölçüsü ve Dirac ölçüsü]\label{ex:saymadirac}
\begin{enumerate}[label=(\alph*)]
  \item \textbf{Sayma ölçüsü} \emph{(counting measure)}: $\Omega$ herhangi,
        $\mathcal{F} = \mathcal{P}(\Omega)$,
        $\mu(A) = |A|$ (eleman sayısı, sonsuz ise $+\infty$).

  \item \textbf{Dirac ölçüsü} \emph{(Dirac measure)}: $x_0 \in \Omega$ sabit,
        \[
          \delta_{x_0}(A) = \gostf{A}(x_0) =
          \begin{cases} 1 & x_0 \in A, \\ 0 & x_0 \notin A. \end{cases}
        \]
        $\delta_{x_0}$ bir olasılık ölçüsüdür. Tüm ölçü $x_0$ noktasında yoğunlaşmıştır.
\end{enumerate}
\end{ornek}

\begin{mikroozet}[§1.2 özeti]
Ölçü: $\mu(\emptyset)=0$ ve $\sigma$-toplamsallık. Olasılık ölçüsü: ölçü + $\mu(\Omega)=1$.
Temel özellikler: monotonluk, alt-toplamsallık, aşağıdan/yukarıdan süreklilik
(sonluluk koşuluyla). Örnekler: sayma ölçüsü, Dirac ölçüsü, Lebesgue ölçüsü.
\end{mikroozet}


% =====================================================================
\section{Dış Ölçü ve Carathéodory'nin Genişleme Teoremi}
% =====================================================================

\begin{motivasyon}[§1.3]
Lebesgue ölçüsünü nasıl kurarız? ``Uzunluk'' fonksiyonunu aralıklar üzerinde
tanımlamak kolay: $\lambda((a,b)) = b-a$. Ama bu tanımı tüm Borel kümelerine
nasıl genişletiriz? Carathéodory'nin yaklaşımı: önce tüm $\mathcal{P}(\mathbb{R})$
üzerinde bir ``dış ölçü'' tanımla; ardından ölçülebilir kümeleri seç.
\end{motivasyon}

\begin{tanim}[Dış Ölçü]\label{def:dismolcu}
$\mu^*: \mathcal{P}(\Omega) \to [0,+\infty]$ fonksiyonuna \textbf{dış ölçü}
\emph{(outer measure)} denir; eğer:
\begin{enumerate}[label=(\alph*)]
  \item $\mu^*(\emptyset) = 0$,
  \item $A \subseteq B \Rightarrow \mu^*(A) \leq \mu^*(B)$,
  \item $\mu^*\!\left(\bigcup_{n=1}^\infty A_n\right) \leq \sum_{n=1}^\infty \mu^*(A_n)$.
\end{enumerate}
\end{tanim}

\begin{tanim}[Carathéodory Ölçülebilirliği]\label{def:caratheodory}
$E \subseteq \Omega$, $\mu^*$'a göre \textbf{Carathéodory-ölçülebilir} adını alır;
eğer her $A \subseteq \Omega$ için
\[
  \mu^*(A) = \mu^*(A \cap E) + \mu^*(A \cap E^c).
\]
Carathéodory-ölçülebilir kümelerin ailesi $\mathcal{M}^*$ ile gösterilir.
\end{tanim}

\begin{teorem}[Carathéodory'nin Genişleme Teoremi]\label{thm:caratheodory}
$\mu^*$ bir dış ölçü olsun. O zaman:
\begin{enumerate}[label=(\roman*)]
  \item $\mathcal{M}^*$ bir $\sigma$-cebiriyidir.
  \item $\mu^*\!\restriction_{\mathcal{M}^*}$ bir ölçüdür.
\end{enumerate}
\end{teorem}

\begin{proof}
\textit{Fikir:} Carathéodory koşulu, $E$'nin her $A$'yı ``mükemmel biçimde''
iki parçaya böldüğünü söyler; bu özellik $\sigma$-cebiri aksiyomlarıyla uyumludur.

\medskip

\textbf{$\mathcal{M}^*$ cebirdir (sonlu kapatma).}

\textit{Tümleme:} $E \in \mathcal{M}^*$ ise $E^c \in \mathcal{M}^*$: tanım simetrik.

\textit{Sonlu birleşim:} $E, F \in \mathcal{M}^*$, $A \subseteq \Omega$ olsun.
$E$'nin Carathéodory koşulunu $A \cap F^c$'ye uygula:
\[
  \mu^*(A \cap F^c) = \mu^*(A \cap F^c \cap E) + \mu^*(A \cap F^c \cap E^c).
\]
$F$'nin koşulunu $A$'ya uygula: $\mu^*(A) = \mu^*(A \cap F) + \mu^*(A \cap F^c)$.
Toplarsak ve $A \cap (E \cup F)^c = A \cap E^c \cap F^c$ göz önünde tutulursa,
$\mu^*(A) = \mu^*(A \cap (E \cup F)) + \mu^*(A \cap (E \cup F)^c)$.

\medskip

\textbf{$\mathcal{M}^*$ sayılabilir birleşim altında kapalıdır ve $\sigma$-toplamsallık.}

$(E_n) \subseteq \mathcal{M}^*$ ikili ayrık, $S_n = \bigcup_{k=1}^n E_k$, $S = \bigcup_n E_n$.
Tümevarımla $S_n \in \mathcal{M}^*$. Her $A$ için:
\[
  \mu^*(A \cap S_n) = \sum_{k=1}^n \mu^*(A \cap E_k).
\]
(Bu eşitlik $E_k$'nin Carathéodory koşulunun $A \cap S_k$'ya uygulanmasıyla tümevarımla elde edilir.)
$S_n \in \mathcal{M}^*$ koşulunu $A$'ya uygula ve $n \to \infty$ al:
\[
  \mu^*(A) = \mu^*(A \cap S_n) + \mu^*(A \cap S_n^c)
  \geq \sum_{k=1}^n \mu^*(A \cap E_k) + \mu^*(A \cap S^c).
\]
$n \to \infty$: $\mu^*(A) \geq \sum_k \mu^*(A \cap E_k) + \mu^*(A \cap S^c)
\geq \mu^*(A \cap S) + \mu^*(A \cap S^c) \geq \mu^*(A)$.
Tüm eşitsizlikler eşitlik verir; $S \in \mathcal{M}^*$ ve $\sigma$-toplamsallık.
\end{proof}

\begin{mikroozet}[§1.3 özeti]
Dış ölçü tüm kümelerde tanımlı ama yalnızca $\sigma$-alt-toplamsaldır.
Carathéodory kriteri ölçülebilir kümeleri seçer; bu kümelerin oluşturduğu
aile bir $\sigma$-cebiriyidir ve dış ölçünün kısıtlaması tam bir ölçüdür.
\end{mikroozet}


% =====================================================================
\section{Lebesgue Ölçüsü}
% =====================================================================

\begin{tanim}[Lebesgue Dış Ölçüsü]\label{def:lebesgue_dis}
$A \subseteq \mathbb{R}$ için \textbf{Lebesgue dış ölçüsü}:
\[
  \lambda^*(A) \coloneqq \inf\left\{
    \sum_{n=1}^\infty (b_n - a_n) :
    A \subseteq \bigcup_{n=1}^\infty (a_n, b_n)
  \right\}.
\]
\end{tanim}

\begin{teorem}[Lebesgue Ölçüsünün Varlığı]\label{thm:lebesgue_varlik}
$\lambda^*$ bir dış ölçüdür ve $\lambda^*((a,b)) = b - a$ her $a < b$ için.
Carathéodory teoremiyle, $\lambda^*$'nin kısıtlaması
$\mathcal{L} \supseteq \mathcal{B}(\mathbb{R})$ üzerinde bir ölçü verir.
Bu ölçüye \textbf{Lebesgue ölçüsü} $\lambda$ denir.
\end{teorem}

\begin{proof}[İspat (özet)]
Dış ölçü aksiyomlarının doğruluğu $\lambda^*$ tanımından açıktır.
$\lambda^*((a,b)) = b-a$: $\geq$ yönü $(a,b)$'nin kendisiyle örtülmesinden;
$\leq$ yönü Heine-Borel teoremiyle (her açık örtünün sonlu alt örtüsü var)
sağlanır; detay için Billingsley~\cite{billingsley1995} §3 bakınız.
Borel kümelerinin Carathéodory-ölçülebilir olduğu doğrulanır.
\end{proof}

\begin{onerme}[Lebesgue Ölçüsünün Özellikleri]\label{prop:lebesgueozell}
Lebesgue ölçüsü $\lambda$:
\begin{enumerate}[label=(\roman*)]
  \item \textbf{Ötelemede değişmez:} $\lambda(A + x) = \lambda(A)$ her $x \in \mathbb{R}$,
        $A \in \mathcal{L}$ için.
  \item \textbf{Ölçeklemede:} $\lambda(cA) = |c|\,\lambda(A)$ her $c \in \mathbb{R}$.
  \item \textbf{$\sigma$-sonlu:} $\mathbb{R} = \bigcup_n (-n,n)$ ve $\lambda((-n,n)) = 2n < \infty$.
  \item \textbf{Tam:} $\lambda(N) = 0$ ve $A \subseteq N$ ise $A \in \mathcal{L}$ ve $\lambda(A) = 0$.
\end{enumerate}
\end{onerme}

\begin{karsiornek}[Vitali kümesi --- Lebesgue-ölçülemeyen küme]
Seçim aksiyomu kabul edilirse, $[0,1]$'de Lebesgue-ölçülemeyen kümeler
mevcuttur (\textbf{Vitali kümesi}). İnşaat: $[0,1]$'deki sayıları
$\sim$ ile $x \sim y \iff x - y \in \mathbb{Q}$ denk bağıntısıyla sınıflandırın.
Her denk sınıftan bir eleman seçerek $V$ oluşturun. $V$ ölçülebilir değildir:
ölçülebilir olsaydı ya $\lambda(V) = 0$ (ama $[0,1] \subseteq \bigcup_{q \in \mathbb{Q} \cap [-1,1]} (V+q)$,
alt-toplamsallıkla çelişir) ya da $\lambda(V) > 0$ (ama $\lambda([0,2]) \geq \sum_{q} \lambda(V+q)$,
ötelemede değişmezlikle $= \infty$; çelişki).
\end{karsiornek}

\begin{mikroozet}[§1.4 özeti]
Lebesgue ölçüsü aralıkların uzunluğunu tüm Borel (ve daha geniş Lebesgue)
kümelerine genişletir. Öteleme değişmezliği, tamlık ve $\sigma$-sonluluk
temel özellikleridir. Ölçüm teorisinin zorunlu kıldığı Seçim Aksiyomu,
ölçülemeyen kümelerin varlığını garanti eder.
\end{mikroozet}


% =====================================================================
\section{Ölçülebilir Fonksiyonlar}
% =====================================================================

\begin{tanimgerekce}
Rasgele değişken, sezgisel olarak ``rastlantıya bağlı sayısal bir sonuç''tur.
Ama hangi fonksiyonlar rasgele değişken olarak kabul edilebilir? Cevap:
olayları olaylara götürenler, yani ters görüntüsü $\sigma$-cebirini koruyanlar.
\end{tanimgerekce}

\begin{tanim}[Ölçülebilir Fonksiyon]\label{def:olculebilfon}
$(\Omega, \mathcal{F})$ ve $(S, \mathcal{G})$ ölçülebilir uzaylar,
$f: \Omega \to S$ bir fonksiyon olsun.
$f$'ye $\mathcal{F}/\mathcal{G}$-\textbf{ölçülebilir} denir;
eğer her $B \in \mathcal{G}$ için $f^{-1}(B) \in \mathcal{F}$.

Özel olarak $S = \mathbb{R}$, $\mathcal{G} = \mathcal{B}(\mathbb{R})$ ise
$f$ kısaca $\mathcal{F}$-\textbf{ölçülebilir} adını alır.
\end{tanim}

\begin{teorem}[Ölçülebilirlik Karakterizasyonları]\label{thm:olculebichar}
$f: (\Omega,\mathcal{F}) \to (\mathbb{R}, \mathcal{B}(\mathbb{R}))$ için
aşağıdakiler denktir:
\begin{enumerate}[label=(\roman*)]
  \item $f$ ölçülebilirdir.
  \item Her $b \in \mathbb{R}$ için $\{f \leq b\} \coloneqq \{\omega : f(\omega) \leq b\} \in \mathcal{F}$.
  \item Her $b \in \mathbb{R}$ için $\{f < b\} \in \mathcal{F}$.
  \item Her $b \in \mathbb{R}$ için $\{f \geq b\} \in \mathcal{F}$.
  \item Her $b \in \mathbb{R}$ için $\{f > b\} \in \mathcal{F}$.
\end{enumerate}
\end{teorem}

\begin{proof}
(i) $\Rightarrow$ (ii): $(-\infty,b] \in \mathcal{B}(\mathbb{R})$
(Teorem~\ref{thm:boreldenk}(v)), dolayısıyla $f^{-1}((-\infty,b]) \in \mathcal{F}$.

(ii) $\Rightarrow$ (i): $\mathcal{B}(\mathbb{R}) = \sigma(\{(-\infty,b] : b \in \mathbb{R}\})$.
Koşul (ii), $\{(-\infty,b]\}$ ailesinin ters görüntülerinin $\mathcal{F}$'de
olduğunu söyler. $\{B : f^{-1}(B) \in \mathcal{F}\}$ ailesinin bir $\sigma$-cebiri
olduğu görülür (Bölüm~0, Teorem~\ref{thm:tergoriuntu}); bu $\sigma$-cebiri
$\{(-\infty,b]\}$ ailesini içerdiğinden $\mathcal{B}(\mathbb{R})$'yi içerir.

(ii) $\Leftrightarrow$ (iii)---(v) benzer dönüşümlerle: $\{f < b\} = \bigcup_n \{f \leq b - 1/n\}$, vb.
\end{proof}

\begin{onerme}[Ölçülebilir Fonksiyonların Kararlılığı]\label{prop:olculebikarar}
\begin{enumerate}[label=(\roman*)]
  \item $f, g$ ölçülebilir, $c \in \mathbb{R}$ sabit $\Rightarrow$
        $f+g$, $fg$, $cf$, $|f|$, $f \vee g$, $f \wedge g$ ölçülebilir.
  \item $f_n$ ölçülebilir ($n \geq 1$) $\Rightarrow$
        $\sup_n f_n$, $\inf_n f_n$, $\limsup_n f_n$, $\liminf_n f_n$ ölçülebilir.
  \item $f_n \to f$ noktasal ise ve her $f_n$ ölçülebilir ise $f$ ölçülebilir.
\end{enumerate}
\end{onerme}

\begin{proof}
(i) $\{f + g > b\} = \bigcup_{q \in \mathbb{Q}} (\{f > q\} \cap \{g > b-q\})$;
rasyoneller sayılabilir, her küme $\mathcal{F}$'de. Geri kalanlar benzer.

(ii) $\{\sup_n f_n > b\} = \bigcup_n \{f_n > b\} \in \mathcal{F}$.
$\limsup_n f_n = \inf_n \sup_{k \geq n} f_k$; (ii)'nin tekrarlı uygulaması.

(iii) Noktasal limit $\limsup = \liminf$ olduğu durumdur; (ii)'den.
\end{proof}

\begin{tanim}[Basit Fonksiyon]\label{def:basit}
$f: \Omega \to \mathbb{R}$ \textbf{basit} \emph{(simple function)} adını alır; eğer
yalnızca sonlu sayıda değer alıyorsa. Standart gösterim:
\[
  f = \sum_{k=1}^n c_k \gostf{A_k},
  \quad c_k \in \mathbb{R},\quad A_k = f^{-1}(\{c_k\}) \in \mathcal{F},
  \quad A_k \text{ ikili ayrık.}
\]
\end{tanim}

\begin{teorem}[Basit Fonksiyon Yaklaşımı]\label{thm:basityaklasim}
$f: (\Omega, \mathcal{F}) \to [0,\infty]$ ölçülebilir ise, her $n \geq 1$ için
$0 \leq f_n \leq f_{n+1} \leq f$ ve $f_n \nearrow f$ noktasal olan ölçülebilir
basit fonksiyonlar $(f_n)$ mevcuttur.
\end{teorem}

\begin{proof}
$f_n(\omega) = \sum_{k=0}^{n 2^n - 1} \frac{k}{2^n} \gostf{\{k/2^n \leq f < (k+1)/2^n\}}(\omega) + n \gostf{\{f \geq n\}}(\omega)$
tanımlayın. $f_n$ ölçülebilir (her $A_k \in \mathcal{F}$), basit (sonlu değer alır),
artan ($f_n \leq f_{n+1}$), ve $f_n(\omega) \nearrow f(\omega)$ her $\omega$ için
($f(\omega) < \infty$ ise yakınsama açık; $f(\omega) = \infty$ ise $f_n(\omega) = n \to \infty$).
\end{proof}

\begin{mikroozet}[§1.5 özeti]
Ölçülebilir fonksiyon: ters görüntüsü $\sigma$-cebirini koruyan fonksiyon.
Karakterizasyon: $\{f \leq b\} \in \mathcal{F}$ her $b$ için yeterli ve gerekli.
Basit fonksiyonlar ölçülebilir fonksiyonları aşağıdan yaklaşır: integral
teorisinin temel aracı.
\end{mikroozet}


% =====================================================================
\section{Lebesgue İntegrali}
% =====================================================================

\begin{motivasyon}[§1.6]
Lebesgue integrali üç aşamada tanımlanır: önce basit fonksiyonlar için
(sezgisel), sonra negatif olmayan fonksiyonlar için (supremum yoluyla),
son olarak işaretli fonksiyonlar için ($f = f^+ - f^-$). Her aşama
bir öncekine dayanır.
\end{motivasyon}

\begin{tanim}[Lebesgue İntegrali --- Üç Aşama]\label{def:lebesgueintegral}
$(\Omega, \mathcal{F}, \mu)$ ölçüm uzayı.

\medskip

\textbf{Aşama 1.} $f = \sum_{k=1}^n c_k \gostf{A_k}$ ($c_k \geq 0$) basit:
\[
  \int f \, d\mu \coloneqq \sum_{k=1}^n c_k \, \mu(A_k) \in [0,\infty].
\]

\textbf{Aşama 2.} $f: \Omega \to [0,\infty]$ ölçülebilir:
\[
  \int f \, d\mu \coloneqq \sup\left\{ \int g \, d\mu :
    g \text{ basit, } 0 \leq g \leq f \right\}.
\]

\textbf{Aşama 3.} $f: \Omega \to \mathbb{R}$ ölçülebilir. $f^+ = \max(f,0)$,
$f^- = \max(-f,0)$. $f$ \textbf{$\mu$-integrallenebilir} adını alır; eğer
$\int f^+ \, d\mu < \infty$ ve $\int f^- \, d\mu < \infty$. O zaman:
\[
  \int f \, d\mu \coloneqq \int f^+ \, d\mu - \int f^- \, d\mu.
\]
$L^1(\mu) \coloneqq \{f \text{ ölçülebilir} : \int |f| \, d\mu < \infty\}$.
\end{tanim}

\begin{notkutu}
$A$ üzerinde integral: $\int_A f \, d\mu \coloneqq \int f \gostf{A} \, d\mu$.
Olasılık teorisinde $\mu = \mathbb{P}$: $\int f \, d\mathbb{P} = \mathbb{E}[f]$.
\end{notkutu}

\begin{teorem}[Lebesgue İntegralin Özellikleri]\label{thm:lebesgueozell}
$f, g \in L^1(\mu)$, $c \in \mathbb{R}$:
\begin{enumerate}[label=(\roman*)]
  \item \textbf{Doğrusallık:} $\int (cf + g) \, d\mu = c \int f \, d\mu + \int g \, d\mu$.
  \item \textbf{Monotonluk:} $f \leq g$ n.k. $\Rightarrow \int f \, d\mu \leq \int g \, d\mu$.
  \item \textbf{Modül eşitsizliği:} $\left|\int f \, d\mu\right| \leq \int |f| \, d\mu$.
  \item \textbf{Sıfır ölçülü küme:} $\mu(A) = 0 \Rightarrow \int_A f \, d\mu = 0$.
  \item \textbf{N.k. eşitlik:} $f = g$ n.k. $\Rightarrow \int f \, d\mu = \int g \, d\mu$.
\end{enumerate}
\end{teorem}

\begin{proof}
(i) Basit fonksiyonlar için doğrusallık tanımdan; genel durum Teorem~\ref{thm:basityaklasim}
ve MYT (§1.7) ile taşınır.
(ii) $g - f \geq 0$ n.k.; Aşama 2 tanımdan $\int(g-f) \geq 0$.
(iii) $|f| = f^+ + f^-$, üçgen eşitsizliği.
(iv), (v) Aşama 1 tanımdan; $\mu(A) = 0$ ise $c_k \mu(A_k \cap A) = 0$.
\end{proof}

\begin{mikroozet}[§1.6 özeti]
Lebesgue integrali: basit fonksiyonlar $\to$ negatif olmayan $\to$ işaretli.
$L^1(\mu)$: integrallenebilir fonksiyonlar. Doğrusallık, monotonluk, n.k.
eşitlik temel özellikleri.
\end{mikroozet}


% =====================================================================
\section{Monoton Yakınsama, Fatou ve Dominante Yakınsama}
% =====================================================================

\begin{motivasyon}[§1.7]
Ne zaman $\lim_n \int f_n \, d\mu = \int \lim_n f_n \, d\mu$ yazabiliriz?
Riemann teorisinde bu çok kısıtlıdır. Lebesgue teorisinin gücü burada ortaya
çıkar: üç büyük geçiş teoremi bu soruyu kapsamlı biçimde yanıtlar.
\end{motivasyon}

\begin{teorem}[Monoton Yakınsama Teoremi (MYT)]\label{thm:myt}
$f_n: (\Omega,\mathcal{F},\mu) \to [0,\infty]$ ölçülebilir ve $f_n \nearrow f$
noktasal ise:
\[
  \int f \, d\mu = \lim_{n\to\infty} \int f_n \, d\mu.
\]
\end{teorem}

\begin{proof}
\textit{Fikir:} $\int f_n$ artan ve $\leq \int f$; yeterliyse limit $\int f$'ye eşit.

\medskip

\textbf{Adım 1.} $f_n \leq f$ ve monotonluktan $\int f_n \leq \int f$; dolayısıyla
$L \coloneqq \lim_n \int f_n \leq \int f$.

\textbf{Adım 2.} $L \geq \int f$ göstereceğiz. $0 < c < 1$ ve $g \leq f$ basit
fonksiyon olsun. $E_n = \{f_n \geq c \cdot g\}$ tanımla; $f_n \nearrow f \geq g$'den
$E_n \nearrow \Omega$.

$\int f_n \, d\mu \geq \int_{E_n} f_n \, d\mu \geq c \int_{E_n} g \, d\mu$.

$g = \sum_k a_k \gostf{A_k}$ basit; aşağıdan süreklilikten (Teorem~\ref{thm:olcuozellik}(iv)):
$\int_{E_n} g \, d\mu = \sum_k a_k \mu(A_k \cap E_n) \to \sum_k a_k \mu(A_k) = \int g \, d\mu$.

Dolayısıyla $L \geq c \int g \, d\mu$. $c \to 1^-$ ve ardından
$g \nearrow f$ (Teorem~\ref{thm:basityaklasim}) ile $L \geq \int f$.
\end{proof}

\begin{sonuc}[MYT --- Seriler]\label{cor:mytseriler}
$f_n \geq 0$ ölçülebilir ise $\int \sum_{n=1}^\infty f_n \, d\mu = \sum_{n=1}^\infty \int f_n \, d\mu$.
\end{sonuc}

\begin{proof}
$g_N = \sum_{n=1}^N f_n \nearrow \sum_n f_n$; MYT uygula.
\end{proof}

\begin{teorem}[Fatou Lemması]\label{thm:fatou}
$f_n \geq 0$ ölçülebilir ise:
\[
  \int \liminf_{n\to\infty} f_n \, d\mu \leq \liminf_{n\to\infty} \int f_n \, d\mu.
\]
\end{teorem}

\begin{proof}
$g_n = \inf_{k \geq n} f_k$ tanımla; $g_n \leq f_n$ ve $g_n \nearrow \liminf_n f_n$.
MYT: $\int \liminf_n f_n = \lim_n \int g_n$.
$g_n \leq f_k$ her $k \geq n$ için; dolayısıyla
$\int g_n \leq \inf_{k \geq n} \int f_k$.
$\lim_n \int g_n \leq \lim_n \inf_{k \geq n} \int f_k = \liminf_n \int f_n$.
\end{proof}

\begin{hataanalizi}[Fatou'da eşitsizlik keskin olabilir]
Eşitlik her zaman sağlanmaz. $f_n = \gostf{[n,n+1]}$ alalım; $f_n \to 0$ noktasal,
$\int f_n = 1$ her $n$ için. $\int \liminf f_n = 0 < 1 = \liminf \int f_n$.
\end{hataanalizi}

\begin{teorem}[Dominante Yakınsama Teoremi (DYT)]\label{thm:dyt}
$f_n \to f$ n.k. ve $|f_n| \leq g$ n.k. (her $n$ için), $g \in L^1(\mu)$ ise:
\[
  \lim_{n\to\infty} \int f_n \, d\mu = \int f \, d\mu
  \quad \text{ve} \quad
  \lim_{n\to\infty} \int |f_n - f| \, d\mu = 0.
\]
\end{teorem}

\begin{proof}
\textit{Fikir:} $g + f_n \geq 0$ ve $g - f_n \geq 0$; Fatou uygulanabilir.

\medskip

$g + f_n \geq 0$; Fatou:
$\int (g + f) \leq \liminf_n \int (g + f_n) = \int g + \liminf_n \int f_n$.
$\int g < \infty$'dan $\int f \leq \liminf_n \int f_n$.

$g - f_n \geq 0$; Fatou:
$\int (g - f) \leq \liminf_n \int (g - f_n) = \int g - \limsup_n \int f_n$.
$\limsup_n \int f_n \leq \int f$.

İkisi birden: $\limsup_n \int f_n \leq \int f \leq \liminf_n \int f_n$,
yani $\lim_n \int f_n = \int f$.

$L^1$ yakınsaması: $|f_n - f| \leq 2g \in L^1$; noktasal $\to 0$; DYT yeniden uygulanır.
\end{proof}

\begin{ornek}[DYT uygulaması]\label{ex:dytuyg}
$f_n(x) = \dfrac{nx}{1 + n^2 x^2}$, $x \in [0,1]$, Lebesgue ölçüsüyle.
$f_n(x) \to 0$ her $x > 0$ için (ve $f_n(0) = 0$).
$|f_n(x)| \leq \frac{1}{2}$ (AM-GM: $nx \cdot 1 \leq (nx + 1/(nx))/2$... daha basit:
$f_n'(x) = 0$ çözümü $x = 1/n$; $f_n(1/n) = 1/2$). Dolayısıyla $g(x) = 1/2 \in L^1([0,1])$.
DYT: $\int_0^1 f_n \, dx \to \int_0^1 0 \, dx = 0$.
\end{ornek}

\begin{mikroozet}[§1.7 özeti]
MYT: artan + negatif olmayan $\Rightarrow$ limit içine girebilir.
Fatou: negatif olmayan $\Rightarrow$ $\int \liminf \leq \liminf \int$.
DYT: integrallenebilir dom. fonksiyon varsa noktasal yakınsama $\Rightarrow$ $L^1$ yakınsaması.
\end{mikroozet}


% =====================================================================
\section{Ürün Ölçüleri ve Fubini-Tonelli Teoremi}
% =====================================================================

\begin{motivasyon}[§1.8]
İki rasgele değişkenin ortak dağılımını hesaplamak, $\mathbb{R}^2$ üzerinde
bir integral almak gerektirir. Fubini teoremi bu iki boyutlu integrali
tek boyutlu integrallerin ardışık uygulamasına indirgeyerek hesaplamayı
mümkün kılar.
\end{motivasyon}

\begin{tanim}[Ürün $\sigma$-Cebiri]\label{def:urun_sigma}
$(\Omega_1, \mathcal{F}_1)$ ve $(\Omega_2, \mathcal{F}_2)$ ölçülebilir uzaylar.
\[
  \mathcal{F}_1 \otimes \mathcal{F}_2 \coloneqq
  \sigma\bigl(\{A_1 \times A_2 : A_1 \in \mathcal{F}_1, A_2 \in \mathcal{F}_2\}\bigr)
\]
$\Omega_1 \times \Omega_2$ üzerinde \textbf{ürün $\sigma$-cebiriyidir}.
$A_1 \times A_2$ biçimindeki kümelere \textbf{dikdörtgen (measurable rectangle)} denir.
\end{tanim}

\begin{onerme}
$\mathcal{B}(\mathbb{R}) \otimes \mathcal{B}(\mathbb{R}) = \mathcal{B}(\mathbb{R}^2)$.
\end{onerme}

\begin{proof}
$\mathbb{R}^2$'deki açık dikdörtgenler $(a_1,b_1) \times (a_2,b_2)$ hem
$\mathcal{B}(\mathbb{R}) \otimes \mathcal{B}(\mathbb{R})$'de hem de
$\mathcal{B}(\mathbb{R}^2)$'de; karşılıklı içerme verir.
\end{proof}

\begin{teorem}[Ürün Ölçüsünün Varlığı ve Tekliği]\label{thm:urunoelcu}
$(\Omega_1, \mathcal{F}_1, \mu_1)$ ve $(\Omega_2, \mathcal{F}_2, \mu_2)$ $\sigma$-sonlu
ölçüm uzayları olsun. $\Omega_1 \times \Omega_2$ üzerinde tek bir $\mu_1 \otimes \mu_2$
ölçüsü mevcuttur; öyle ki her dikdörtgen için:
\[
  (\mu_1 \otimes \mu_2)(A_1 \times A_2) = \mu_1(A_1) \cdot \mu_2(A_2).
\]
\end{teorem}

\begin{proof}[İspat (özet)]
Dikdörtgenler bir $\pi$-sistemi oluşturur. Carathéodory genişlemesi
$\mu_1 \otimes \mu_2$'yi inşa eder. $\sigma$-sonluluk tekliği garantilemek
için gereklidir (Dynkin'in $\pi$-$\lambda$ argümanı).
Ayrıntı için Klenke~\cite{klenke2014} §14 bakınız.
\end{proof}

\begin{teorem}[Tonelli Teoremi]\label{thm:tonelli}
$f: (\Omega_1 \times \Omega_2, \mathcal{F}_1 \otimes \mathcal{F}_2) \to [0,\infty]$
ölçülebilir olsun. O zaman:
\begin{enumerate}[label=(\roman*)]
  \item $\omega_1 \mapsto \int f(\omega_1, \omega_2) \, d\mu_2(\omega_2)$ ve
        $\omega_2 \mapsto \int f(\omega_1, \omega_2) \, d\mu_1(\omega_1)$
        ölçülebilir fonksiyonlardır.
  \item
        \[
          \int f \, d(\mu_1 \otimes \mu_2)
          = \int\!\!\int f(\omega_1,\omega_2) \, d\mu_2(\omega_2) \, d\mu_1(\omega_1)
          = \int\!\!\int f(\omega_1,\omega_2) \, d\mu_1(\omega_1) \, d\mu_2(\omega_2).
        \]
\end{enumerate}
\end{teorem}

\begin{teorem}[Fubini Teoremi]\label{thm:fubini}
$f \in L^1(\mu_1 \otimes \mu_2)$ olsun. O zaman:
\begin{enumerate}[label=(\roman*)]
  \item $\mu_1$-n.k. her $\omega_1$ için $\omega_2 \mapsto f(\omega_1, \omega_2) \in L^1(\mu_2)$.
  \item Tonelli teoremindeki tekrarlı integrasyon formülü geçerlidir.
\end{enumerate}
\end{teorem}

\begin{proof}[İspat (özet)]
Tonelli teoremini $f^+$ ve $f^-$'ya ayrı ayrı uygulayın.
$f \in L^1$ olduğundan $\int f^+ < \infty$ ve $\int f^- < \infty$;
$\omega_1$-neredeyse her yerde tekil integraller de sonludur.
\end{proof}

\begin{hataanalizi}[$L^1$ koşulu olmadan Fubini başarısız olabilir]
$f(x,y) = \frac{x^2 - y^2}{(x^2+y^2)^2}$ üzerinde $[0,1]^2$ Lebesgue ölçüsü.
$\int_0^1 \int_0^1 f \, dx \, dy = \pi/4$ ama $\int_0^1 \int_0^1 f \, dy \, dx = -\pi/4$.
$f \notin L^1([0,1]^2)$; Fubini koşulu sağlanmaz.
\end{hataanalizi}

\begin{ornek}[Fubini --- somut uygulama]
$(\mathbb{R}^2, \mathcal{B}(\mathbb{R}^2), \lambda^2)$ üzerinde
$f(x,y) = e^{-x^2 - y^2}$ integralini hesaplayın.
\[
  \int_{\mathbb{R}^2} e^{-x^2-y^2} \, d\lambda^2
  = \int_{-\infty}^\infty e^{-x^2} \, dx \cdot \int_{-\infty}^\infty e^{-y^2} \, dy
  = \left(\int_{-\infty}^\infty e^{-x^2} \, dx\right)^{\!2}.
\]
Kutupsal koordinatlarla $\int_{\mathbb{R}^2} e^{-r^2} r \, dr \, d\theta = \pi$;
dolayısıyla $\int_{-\infty}^\infty e^{-x^2} \, dx = \sqrt{\pi}$.
Bu, Gaussian integralinin en zarif hesabıdır.
\end{ornek}

\begin{mikroozet}[§1.8 özeti]
Ürün $\sigma$-cebiri: dikdörtgenler tarafından üretilen. Ürün ölçüsü dikdörtgenlerde
çarpımsal. Tonelli: negatif olmayan için sıra serbesttir. Fubini: $L^1$'de sıra serbesttir.
$L^1$ koşulu olmadan sıra değiştirme yanlış sonuç verebilir.
\end{mikroozet}


% =====================================================================
\section{$L^p$ Uzayları ve İntegral Eşitsizlikleri}
\label{sec:lp_uzaylari}
% =====================================================================

\begin{tanim}[$L^p$ Uzayı]
\label{def:lp_uzay}
$(\Omega, \mathcal{F}, \mu)$ ölçüm uzayı, $1 \leq p < \infty$.
\[
  L^p(\mu) = \{f \text{ ölçülebilir} : \|f\|_p := \left(\int |f|^p \, d\mu\right)^{1/p} < \infty\}.
\]
$L^\infty(\mu)$: $\|f\|_\infty = \inf\{M : \mu(\{|f| > M\}) = 0\}$ (esansiyel üst sınır).
Teknik olarak $L^p$ sınıfları: eşdeğer fonksiyonlar (n.k.\ eşit) aynı sınıfa.
\end{tanim}

\begin{teorem}[Hölder Eşitsizliği]
\label{thm:holder}
$1/p + 1/q = 1$ ($p,q > 1$), $f \in L^p$, $g \in L^q$. O zaman $fg \in L^1$ ve
\[
  \int |fg| \, d\mu \leq \|f\|_p \, \|g\|_q.
\]
Özel durum $p = q = 2$: Cauchy-Schwarz eşitsizliği $\int|fg| \leq \|f\|_2\|g\|_2$.
\end{teorem}

\begin{proof}
Young eşitsizliği: $ab \leq a^p/p + b^q/q$ ($a,b \geq 0$, $1/p+1/q=1$).
$a = |f|/\|f\|_p$, $b = |g|/\|g\|_q$ koy; $\mu$-integral al:
$\int |fg|/(\|f\|_p\|g\|_q) \leq \int |f|^p/(p\|f\|_p^p) + \int|g|^q/(q\|g\|_q^q) = 1/p + 1/q = 1$.
\end{proof}

\begin{teorem}[Minkowski Eşitsizliği]
\label{thm:minkowski}
$1 \leq p \leq \infty$, $f, g \in L^p$. O zaman $f + g \in L^p$ ve
\[
  \|f + g\|_p \leq \|f\|_p + \|g\|_p.
\]
\end{teorem}

\begin{proof}
$p=1$: üçgen eşitsizliği, $|f+g| \leq |f| + |g|$.
$1 < p < \infty$: $\|f+g\|_p^p = \int|f+g|^p \leq \int|f+g|^{p-1}|f| + \int|f+g|^{p-1}|g|$.
Her terime Hölder ($|f+g|^{p-1} \in L^q$ çünkü $(p-1)q = p$):
$\leq \|f+g\|_p^{p-1}(\|f\|_p + \|g\|_p)$. $\|f+g\|_p^{p-1}$'e bölme ile sonuç.
\end{proof}

\begin{teorem}[$L^p$ Uzayı Tamlığı — Riesz-Fischer]
\label{thm:riesz_fischer}
$L^p(\mu)$ ($1 \leq p \leq \infty$) Minkowski normuyla tam metrize edilebilir uzaydır
(Banach uzayı). $p = 2$: Hilbert uzayı (iç çarpım $\langle f,g\rangle = \int fg$).
\end{teorem}

\begin{teorem}[$L^p$ Dahil Olma İlişkisi]
\label{thm:lp_dahil}
$\mu(\Omega) < \infty$ ise $1 \leq p \leq q \leq \infty$ için $L^q \subseteq L^p$.
\end{teorem}

\begin{proof}
Hölder: $\|f\|_p^p = \int|f|^p \cdot 1 \leq \|f^p\|_{q/p} \cdot \|1\|_{q/(q-p)} = \|f\|_q^p \cdot \mu(\Omega)^{1-p/q}$.
\end{proof}

\begin{ornek}[$L^p$ Farklılıkları]
$\Omega = (0,1)$, $\lambda$ Lebesgue. $f(x) = x^{-\alpha}$ ($0 < \alpha < 1$):
$\int_0^1 x^{-\alpha p} \, dx < \infty \Leftrightarrow \alpha p < 1 \Leftrightarrow p < 1/\alpha$.
Yani $f \in L^p$ ancak $p < 1/\alpha$ için. Küçük $p$'ler için $f \in L^p$, büyük $p$'ler için değil.
\end{ornek}

\begin{mikroozet}
$L^p$: $p$-inci mutlak moment sonlu fonksiyonlar. Hölder: $\int|fg| \leq \|f\|_p\|g\|_q$.
Minkowski: $\|f+g\|_p \leq \|f\|_p + \|g\|_p$ (üçgen). $L^p$ Banach; $L^2$ Hilbert.
\end{mikroozet}

% =====================================================================
\section{Radon-Nikodym Teoremi}
\label{sec:radon_nikodym}
% =====================================================================

\begin{tanim}[Mutlak Süreklilik]
\label{def:mutlak_sureklilik}
$\mu$ ve $\nu$ aynı $(\Omega, \mathcal{F})$ üzerinde ölçüler.
$\nu$, $\mu$'ya göre \textbf{mutlak sürekli} ($\nu \ll \mu$) denir, eğer
her $A \in \mathcal{F}$ için $\mu(A) = 0 \Rightarrow \nu(A) = 0$.
\end{tanim}

\begin{ornek}[Mutlak Süreklilik]
$f \geq 0$, $\int f \, d\mu = 1$. $\nu(A) = \int_A f \, d\mu$: o zaman $\nu \ll \mu$.
$\mu(A) = 0 \Rightarrow \int_A f \, d\mu = 0 \Rightarrow \nu(A) = 0$.
\end{ornek}

\begin{teorem}[Radon-Nikodym Teoremi]
\label{thm:radon_nikodym}
$(\Omega, \mathcal{F})$ üzerinde $\sigma$-sonlu $\mu$ ve $\nu$ ölçüleri, $\nu \ll \mu$.
O zaman tek türlü ($\mu$-n.k.) ölçülebilir $f \geq 0$ ($\mu$-n.k.) fonksiyon vardır:
\[
  \nu(A) = \int_A f \, d\mu \quad \text{her } A \in \mathcal{F} \text{ için.}
\]
$f = d\nu/d\mu$: \textbf{Radon-Nikodym türevi} (yoğunluk / Nikodym türevi).
\end{teorem}

\begin{proof}[İspat Taslağı (Hilbert Uzayı Argümanı)]
$\rho = \mu + \nu$. $\ell(g) = \int g \, d\nu$ ($g \in L^2(\rho)$) sınırlı lineer fonksiyonel.
Riesz temsil: $\ell(g) = \int gh \, d\rho$ olan $h \in L^2(\rho)$ var.
$0 \leq h \leq 1$ $\rho$-n.k.\ gösterilebilir. $f = h/(1-h)$ ($h < 1$ üzerinde) $\mu$-n.k.\ tanımlı;
$\nu(A) = \int_A f \, d\mu$ doğrulanır. $f$ benzersizliği: iki RN türevi $f_1, f_2$
ise $\int_A(f_1-f_2)d\mu = 0$ her $A$'da; $f_1 = f_2$ $\mu$-n.k.
\end{proof}

\begin{onerme}[Zincir Kuralı]
$\nu \ll \mu \ll \lambda$ ise $d\nu/d\lambda = (d\nu/d\mu)(d\mu/d\lambda)$ ($\lambda$-n.k.).
\end{onerme}

\begin{ornek}[Olasılıkta Önemi]
$X$ rasgele değişken, $P_X$ görüntü ölçüsü (dağılımı). $P_X \ll \lambda$ ise
RN türevi $f = dP_X/d\lambda$: bu tam olarak $X$'in PDF'sidir.
\emph{Koşullu beklenti} $E[X \mid \mathcal{G}]$ Radon-Nikodym teoreminden doğar
(Bölüm~\ref{chp:beklenti}'de ayrıntılı işlendi).
\end{ornek}

\begin{disiplinlerarasibaglan}
\textbf{İstatistik:} Olabilirlik oranı (likelihood ratio) $d\nu/d\mu$ Radon-Nikodym türevidir;
Neyman-Pearson lemması ve yeterli istatistik teorisi bu çerçevede kurulur.\\
\textbf{Bayesçi çıkarım:} Sonsal dağılım, önsel üzerine olabilirlik oranı ile Radon-Nikodym güncellemesidir.\\
\textbf{Stokastik kalkülüs:} Girsanov teoremi, olasılık ölçüsü değiştirmeyi Radon-Nikodym teoremiyle gerçekleştirir.
\end{disiplinlerarasibaglan}

% =====================================================================
\section{Lebesgue-Stieltjes İntegrali ve Beklenti}
\label{sec:ls_integral}
% =====================================================================

\begin{tanim}[Lebesgue-Stieltjes ölçüsü]\label{def:ls_olcusu}
$F:\mathbb{R}\to[0,1]$ KDF (sağdan sürekli, artan, $F(-\infty)=0$, $F(+\infty)=1$). Carathéodory genişlemesiyle $\mu_F$ tanımlanır: $\mu_F((a,b])=F(b)-F(a)$. Bu \emph{Lebesgue-Stieltjes ölçüsüdür}.
\end{tanim}

\begin{teorem}[Beklenti olarak LS integrali]\label{thm:beklenti_ls}
$X$ rasgele değişken, $F_X$ KDF. Her borelci $g:\mathbb{R}\to\mathbb{R}$ için:
\[
  \Beklenti[g(X)] = \int_\Omega g(X(\omega))\,dP(\omega) = \int_\mathbb{R} g(x)\,d\mu_{F_X}(x)
  = \int_\mathbb{R} g(x)\,dF_X(x).
\]
Sürekli $F_X$'te son integral Riemann-Stieltjes'e indirgenir.
\end{teorem}

\begin{ornek}[Karma dağılım beklentisi]\label{ex:karma_beklenti}
$X$: $p=1/3$ olasılıkla $X=0$ (kesikli kısım), $(1-p)=2/3$ olasılıkla $\Uniform(0,1)$ (sürekli kısım).
$F_X(x)=\frac13\mathbf{1}[x\geq0]+\frac23 x\mathbf{1}[0\leq x<1]+\frac23\mathbf{1}[x\geq1]$.
$\Beklenti[X]=\frac13\cdot0+\frac23\int_0^1x\,dx=\frac23\cdot\frac12=\frac13$.
LS integrali karma durumlarda tek çatı altında birleştirir.
\end{ornek}

\begin{teorem}[Riemann $\leftrightarrow$ Lebesgue]\label{thm:riemann_lebesgue}
$f:[a,b]\to\mathbb{R}$ sınırlı. $f$ Riemann integrallenebilir $\iff$ $f$'nin süreksizlik noktaları $\lambda$-ölçüsü sıfır kümedir. Bu durumda Riemann ve Lebesgue integralleri eşittir.
\end{teorem}

\begin{hataanalizi}
$f(x)=\mathbf{1}_\mathbb{Q}(x)$ Dirichlet fonksiyonu: Riemann integrallenemez (her alt aralıkta rasyonel + irrasyonel var), ama Lebesgue integrali sıfırdır ($\lambda(\mathbb{Q})=0$). Bu, Lebesgue integralinin Riemann'ı kesinlikle kapsadığını gösterir.
\end{hataanalizi}

% =====================================================================
\section{Olasılık Ölçülerinin Zayıf Topolojisi}
\label{sec:zayif_topoloji}
% =====================================================================

\begin{tanim}[Olasılık ölçülerinin zayıf yakınsaması]\label{def:zayif_yakin}
$\mathcal{P}(\mathbb{R})$ olasılık ölçüleri uzayı. $\mu_n\overset{w}{\to}\mu$: her sınırlı, sürekli $f$ için $\int f\,d\mu_n\to\int f\,d\mu$. Bu \emph{zayıf yakınsama} (dağılımda yakınsama ile eşdeğer).
\end{tanim}

\begin{teorem}[Portmanteau teoremi]\label{thm:portmanteau}
$\mu_n\overset{w}{\to}\mu$ aşağıdaki koşulların herhangi biriyle eşdeğerdir:
\begin{enumerate}[(i)]
  \item Her kapalı $F$ için $\limsup\mu_n(F)\leq\mu(F)$.
  \item Her açık $G$ için $\liminf\mu_n(G)\geq\mu(G)$.
  \item Her $\mu$-sınır kümesi $A$ ($\mu(\partial A)=0$) için $\mu_n(A)\to\mu(A)$.
  \item Her sınırlı, Lipschitz $f$ için $\int f\,d\mu_n\to\int f\,d\mu$.
\end{enumerate}
\end{teorem}

\begin{ornek}[Portmanteau uygulaması]\label{ex:portmanteau_uygulama}
$X_n\dto X\sim\Normal(0,1)$. $A=(-\infty,1.96]$: $\mu(\partial A)=\mu(\{1.96\})=0$ ($\Normal$ sürekli). Portmanteau (iii): $P(X_n\leq1.96)\to\Phi(1.96)=0.975$. Normal tablolar bu güvenceyle kullanılabilir.
\end{ornek}

\begin{disiplinlerarasibaglan}
\textbf{Bayes İstatistiği:} Zayıf topoloji, büyük örneklem Bayes posteriorlarının (Bernstein-von Mises) incelenmesinde kullanılır — posteriorun frekansçı normaline zayıf yakınsaması.

\textbf{Optik Ağlar/Kuyruk Teorisi:} Servis süreleri dağılımları zayıf yakınsama çerçevesinde ele alınır; Portmanteau teoremi ağ kapasitesi sınırlarını garanti eder.
\end{disiplinlerarasibaglan}

\begin{mikroozet}
$\nu \ll \mu$: sıfır $\mu$-ölçülü kümeler sıfır $\nu$-ölçülü.
Radon-Nikodym: $\nu \ll \mu$ ve $\sigma$-sonlu $\Rightarrow$ $\nu(A) = \int_A f\,d\mu$ olan $f = d\nu/d\mu$ var.
$f$: koşullu beklenti ve PDF'nin ölçüm teorisi temeli.
LS integrali: karma dağılımları tek çatıda birleştirir; Riemann $\subset$ Lebesgue.
Portmanteau: zayıf yakınsama için 4 eşdeğer koşul; normal tablo kullanımının garantisi.
\end{mikroozet}

% =====================================================================
%  BÖLÜM SONU ÖZETİ
% =====================================================================
\section*{Bölüm Özeti}
\addcontentsline{toc}{section}{Bölüm Özeti}

\begin{bolumozeti}[Bölüm 1 --- Temel Sonuçlar]
\begin{itemize}
  \item \textbf{$\sigma$-cebiri:} (A)~$\Omega \in \mathcal{F}$;
        (B)~tümleme kapalı; (C)~sayılabilir birleşim kapalı.
        $\sigma(\mathcal{E})$: $\mathcal{E}$'yi içeren en küçük.
        $\mathcal{B}(\mathbb{R}) = \sigma(\text{açık aralıklar})$.
  \item \textbf{Ölçü:} $\mu(\emptyset)=0$ ve $\sigma$-toplamsallık.
        Monotonluk, alt-toplamsallık, aşağıdan/yukarıdan süreklilik
        (yukarıdan için $\mu(A_1)<\infty$).
  \item \textbf{Carathéodory:} Dış ölçü $\mu^*$ verilince, Carathéodory-ölçülebilir
        kümeler bir $\sigma$-cebiri oluşturur ve $\mu^*$ bu ailede ölçüdür.
  \item \textbf{Lebesgue:} $\lambda^*$ dış ölçüden $\lambda$ elde edilir;
        $\lambda((a,b))=b-a$, öteleme değişmez, tam ve $\sigma$-sonlu.
  \item \textbf{Ölçülebilir $f$:} $f^{-1}(B) \in \mathcal{F}$ her $B \in \mathcal{B}$.
        Eşdeğer: $\{f \leq b\} \in \mathcal{F}$ her $b$.
        Basit fonksiyonlar artan limit.
  \item \textbf{Lebesgue integrali:} Basit $\to$ $[0,\infty]$ $\to$ $\mathbb{R}$.
        $L^1(\mu)$: $\int|f|<\infty$. Doğrusal, monoton, n.k.'ya duyarsız.
  \item \textbf{MYT:} $f_n \nearrow f \geq 0 \Rightarrow \int f_n \to \int f$.
  \item \textbf{Fatou:} $f_n \geq 0 \Rightarrow \int\liminf f_n \leq \liminf\int f_n$.
  \item \textbf{DYT:} $f_n \to f$ n.k., $|f_n| \leq g \in L^1 \Rightarrow \int f_n \to \int f$.
  \item \textbf{Fubini-Tonelli:} $\sigma$-sonlu + ($f\geq 0$ ya da $f\in L^1$)
        $\Rightarrow$ tekrarlı integral sırası serbesttir.
\end{itemize}
\end{bolumozeti}

\begin{ozyansima}
\begin{itemize}
  \item Bir $\sigma$-cebiri örnekleyin ve üç aksiyomu teker teker doğrulayın.
  \item Neden $\mu(A_1) < \infty$ koşulu olmadan yukarıdan süreklilik çöker?
        Kendi örneğinizi üretin.
  \item MYT, Fatou ve DYT'yi karşılaştırın: hangi durumlarda hangisini kullanırsınız?
  \item Fubini'nin neden $L^1$ koşulu gerektirdiğini karşı örnekle açıklayın.
\end{itemize}
\end{ozyansima}

\begin{kopru}
\textbf{Sonraki bölüm (Bölüm~2: Olasılık Uzayları).}
Bu bölümde geliştirilen ölçüm teorisi altyapısı, Bölüm~2'de olasılık ölçüsü
$\mathbb{P}$ ile hayata geçirilecek. Koşullu olasılık, bağımsızlık ve Bayes
teoremi $(\Omega,\mathcal{F},\mathbb{P})$ dili içinde yeniden formüle edilecek.

\medskip

\textbf{İleri okuma.}
Billingsley~\cite{billingsley1995} Bölüm~1--3;
Klenke~\cite{klenke2014} Bölüm~1--6;
Shiryaev~\cite{shiryaev1996} Bölüm~II §1--3.
\end{kopru}


% =====================================================================
%  ALIŞTIRMALAR
% =====================================================================

\cozumlualistirmalar

\begin{alistirma}[Al.1.1]{A}{$\sigma$-cebiri doğrulama}
$\Omega = \mathbb{R}$ ve $\mathcal{F} = \{A \subseteq \mathbb{R} : A \text{ ya da } A^c
\text{ sayılabilirdir}\}$ ailesinin bir $\sigma$-cebiri olduğunu gösterin.
\end{alistirma}
\alcozum{%
\textbf{(A)} $\mathbb{R}$ sayılamaz; dolayısıyla $\mathbb{R}^c = \emptyset$ sayılabilir.
$\mathbb{R} \in \mathcal{F}$.

\textbf{(B)} $A \in \mathcal{F}$: $A$ sayılabilirse $A^c$'nin tümleni $A$ sayılabilir,
yani $A^c \in \mathcal{F}$; $A^c$ sayılabilirse $A^c \in \mathcal{F}$ doğrudan.

\textbf{(C)} $(A_n) \subseteq \mathcal{F}$, $A = \bigcup_n A_n$.
Eğer her $A_n$ sayılabilirse $A$ sayılabilir (Teorem~0.\ref{thm:sayilabilirbir});
$A \in \mathcal{F}$.
Eğer bazı $A_{n_0}$ sayılamaz ise $A_{n_0}^c$ sayılabilir ve
$A^c = \bigcap_n A_n^c \subseteq A_{n_0}^c$ sayılabilir; $A \in \mathcal{F}$.
$\blacksquare$%
}

\begin{alistirma}[Al.1.2]{B}{Aşağıdan süreklilik uygulaması}
$(\mathbb{R}, \mathcal{B}(\mathbb{R}), \lambda)$ Lebesgue ölçüm uzayı.
$A_n = (0, 1/n)$ için $\lambda(A_n)$'i hesaplayın ve
$\lambda\!\left(\bigcap_{n=1}^\infty A_n\right)$'yi yukarıdan süreklilik
teoremi ile bulun. Doğrudan hesapla doğrulayın.
\end{alistirma}
\alcozum{%
$\lambda(A_n) = 1/n$. $A_n \searrow \emptyset$ ($\bigcap_n (0,1/n) = \emptyset$;
eğer $x > 0$ ise $x \notin (0,1/n)$ yeterince büyük $n$ için; $x \leq 0$ ise hiç içermez).
$\lambda(A_1) = 1 < \infty$; yukarıdan süreklilik:
$\lambda\!\left(\bigcap_n A_n\right) = \lim_n \lambda(A_n) = \lim_n 1/n = 0$.
Doğrudan: $\lambda(\emptyset) = 0$. Tutarlı. $\blacksquare$%
}

\begin{alistirma}[Al.1.3]{B}{MYT uygulaması}
$f_n(x) = \gostf{[0,n]}(x) \cdot e^{-x}$ ($x \geq 0$), Lebesgue ölçüsüyle.
$\int_0^\infty f_n \, d\lambda$'yı hesaplayın ve MYT kullanarak
$\int_0^\infty e^{-x} \, d\lambda$'yı bulun.
\end{alistirma}
\alcozum{%
$\int_0^\infty f_n \, d\lambda = \int_0^n e^{-x} \, dx = 1 - e^{-n}$.
$f_n \nearrow e^{-x}$ noktasal ($f_n(x) = e^{-x}$ eğer $x \leq n$, $= 0$ eğer $x > n$).
MYT: $\int_0^\infty e^{-x} \, d\lambda = \lim_n (1 - e^{-n}) = 1$. $\blacksquare$%
}

\begin{alistirma}[Al.1.4]{A}{Basit fonksiyon inşası}
$f(x) = x^2 + 1$, $x \in [0,2]$ için Teorem~\ref{thm:basityaklasim}'in
$n=2$ adımındaki $f_2$'yi açıkça yazın ve $\|f - f_2\|_\infty$'u hesaplayın.
\end{alistirma}
\alcozum{%
$n=2$: bölümleme $k/4$, $k=0,1,\ldots,7$ ve $\{f \geq 2\}$.
$f([0,2]) = [1,5]$. $f_2 = \sum_{k=0}^{7} \frac{k}{4} \gostf{\{k/4 \leq f < (k+1)/4\}} + 2\gostf{\{f \geq 2\}}$.
$f$ artan olduğundan $\{k/4 \leq f < (k+1)/4\}$ bir aralık. En büyük hata:
$f$ ile $f_2$ arasındaki fark en fazla $1/4 = 1/2^n$. $\|f-f_2\|_\infty \leq 1/4$.
$\blacksquare$%
}

\alistirmalar

\begin{alistirma}[Al.1.5]{A}{$\sigma$-cebiri işlemleri}
$\mathcal{F}_1$ ve $\mathcal{F}_2$, $\Omega$ üzerinde iki $\sigma$-cebiri.
(a)~$\mathcal{F}_1 \cap \mathcal{F}_2$'nin $\sigma$-cebiri olduğunu gösterin.
(b)~$\mathcal{F}_1 \cup \mathcal{F}_2$'nin genel olarak $\sigma$-cebiri olmayabileceğini
    bir karşı örnekle gösterin.
\end{alistirma}
\alcozum{%
(a) $A \in \mathcal{F}_1 \cap \mathcal{F}_2 \Rightarrow A \in \mathcal{F}_1$ ve
$A \in \mathcal{F}_2$; her iki $\sigma$-cebiri tümleme ve sayılabilir birleşim altında
kapalı; kesişim de kapalı.

(b) $\Omega = \{1,2,3\}$, $\mathcal{F}_1 = \{\emptyset,\{1\},\{2,3\},\Omega\}$,
$\mathcal{F}_2 = \{\emptyset,\{2\},\{1,3\},\Omega\}$.
$\{1\}, \{2\} \in \mathcal{F}_1 \cup \mathcal{F}_2$ ama
$\{1\} \cup \{2\} = \{1,2\} \notin \mathcal{F}_1 \cup \mathcal{F}_2$. $\blacksquare$%
}

\begin{alistirma}[Al.1.6]{A}{Ölçü özellikleri}
$(\Omega, \mathcal{F}, \mu)$ ölçüm uzayı, $A, B \in \mathcal{F}$ ve $\mu(A), \mu(B) < \infty$.
Gösterin ki $\mu(A \cup B) = \mu(A) + \mu(B) - \mu(A \cap B)$.
\end{alistirma}
\alcozum{%
$A \cup B = (A \setminus B) \cup (A \cap B) \cup (B \setminus A)$ ikili ayrık:
$\mu(A \cup B) = \mu(A \setminus B) + \mu(A \cap B) + \mu(B \setminus A)$.
$\mu(A) = \mu(A \setminus B) + \mu(A \cap B)$; $\mu(B) = \mu(B \setminus A) + \mu(A \cap B)$.
Toplarsak: $\mu(A) + \mu(B) = \mu(A \cup B) + \mu(A \cap B)$. $\blacksquare$%
}

\begin{alistirma}[Al.1.7]{A}{Ölçülebilir fonksiyon}
$f: \mathbb{R} \to \mathbb{R}$, $f(x) = \sin x$ fonksiyonunun
$(\mathbb{R}, \mathcal{B}(\mathbb{R}))$ üzerinde ölçülebilir olduğunu gösterin.
\end{alistirma}
\alcozum{%
$\sin: \mathbb{R} \to \mathbb{R}$ sürekli; sürekli fonksiyonlar Borel ölçülebilirdir
(açık kümenin ters görüntüsü açık, dolayısıyla Borel). $\blacksquare$%
}

\begin{alistirma}[Al.1.8]{A}{Fatou lemması --- kesin örnek}
$f_n = n \gostf{(0,1/n)}$ ($n \geq 1$), $[0,1]$ üzerinde Lebesgue ölçüsüyle.
(a)~$\liminf_n f_n$'i bulun.
(b)~$\int \liminf_n f_n$ ile $\liminf_n \int f_n$'yi karşılaştırın.
\end{alistirma}
\alcozum{%
(a) Her $x > 0$ için yeterince büyük $n$'de $x \geq 1/n$; $f_n(x) = 0$.
$f_n(0) = 0$ her zaman. $\liminf_n f_n = 0$ n.k.

(b) $\int \liminf f_n = 0$.
$\int f_n = n \cdot \lambda((0,1/n)) = n \cdot 1/n = 1$.
$\liminf_n \int f_n = 1 > 0$.
Fatou eşitsizliği $0 \leq 1$; kesin eşitsizlik. $\blacksquare$%
}

\begin{alistirma}[Al.1.9]{B}{DYT ile integral sınırı}
$f_n(x) = \dfrac{\sin(x/n)}{x/n} \cdot e^{-|x|}$ ($x \neq 0$),
$f_n(0) = e^0 = 1$, Lebesgue ölçüsüyle $\mathbb{R}$ üzerinde.
(a)~$f_n(x) \to e^{-|x|}$ noktasal olarak gösterin.
(b)~Uygun dominant fonksiyonu bulun.
(c)~$\lim_n \int_{-\infty}^\infty f_n \, d\lambda$'yı hesaplayın.
\end{alistirma}
\alcozum{%
(a) $|\sin t / t| \leq 1$ ve $\sin t / t \to 1$ ($t \to 0$); $x/n \to 0$;
$f_n(x) \to 1 \cdot e^{-|x|} = e^{-|x|}$.

(b) $|f_n(x)| \leq e^{-|x|}$ çünkü $|\sin t/t| \leq 1$.
$g(x) = e^{-|x|} \in L^1(\mathbb{R})$: $\int_{-\infty}^\infty e^{-|x|} dx = 2$.

(c) DYT: $\lim_n \int f_n = \int e^{-|x|} \, d\lambda = 2$. $\blacksquare$%
}

\begin{alistirma}[Al.1.10]{B}{Fubini ve sıra değiştirme}
$f(x,y) = e^{-xy}$, $(x,y) \in [0,\infty) \times [0,\infty)$, Lebesgue ölçüsüyle.
$\int_1^\infty \left(\int_0^\infty e^{-xy} \, dy\right) dx$'i
önce $y$'ye göre integralleyerek, ardından $x$'e göre integralleyerek hesaplayın.
Fubini'nin hangi koşulunu kontrol etmelisiniz?
\end{alistirma}
\alcozum{%
İç integral: $\int_0^\infty e^{-xy} dy = 1/x$ ($x > 0$).
Dış integral: $\int_1^\infty 1/x \, dx = \infty$.

Önce $x$: $\int_1^\infty e^{-xy} dx = e^{-y}/y$ ($y > 0$).
$\int_0^\infty e^{-y}/y \, dy$: $y \to 0^+$ yakınında $1/y$ integrallenemeyen sapma.

Fubini koşulu: $\int_1^\infty \int_0^\infty |f| \, dy \, dx = \int_1^\infty 1/x \, dx = \infty$;
$f \notin L^1([1,\infty) \times [0,\infty))$. Fubini uygulanamaz; sıra fark yaratır. $\blacksquare$%
}

\begin{alistirma}[Al.1.11]{B}{Ölçü sürekliliği --- uygulama}
$\mu$ $(\Omega,\mathcal{F})$ üzerinde bir olasılık ölçüsü. $A_n \in \mathcal{F}$
ve $\sum_{n=1}^\infty \mu(A_n) < \infty$ ise, $\mu\!\left(\limsup_n A_n\right) = 0$
olduğunu ispatlayın.
\end{alistirma}
\alcozum{%
$B_n = \bigcup_{k \geq n} A_k$ olsun; $B_n \searrow \limsup_n A_n$.
Alt-toplamsallık: $\mu(B_n) \leq \sum_{k \geq n} \mu(A_k) \to 0$
($\sum \mu(A_k) < \infty$ kuyruk toplamı).
$\mu(B_1) \leq \sum_k \mu(A_k) < \infty$; yukarıdan süreklilik:
$\mu(\limsup_n A_n) = \lim_n \mu(B_n) = 0$. $\blacksquare$

\emph{Not:} Bu, Borel-Cantelli lemmasının birinci yarısıdır (Bölüm~7).%
}

\begin{alistirma}[Al.1.12]{C}{Lebesgue--Stieltjes ölçüsü}
$F: \mathbb{R} \to \mathbb{R}$ sağdan sürekli, artan bir fonksiyon olsun.
$\mu_F((a,b]) \coloneqq F(b) - F(a)$ ile tanımlanan set fonksiyonunun
yarı-halkalar üzerinde bir ön-ölçü (pre-measure) olduğunu gösterin ve
Carathéodory genişlemesiyle $\mathbb{R}$ üzerinde bir ölçü elde edildiğini açıklayın.
$F(x) = x$ durumunda Lebesgue ölçüsü geri gelir; bu nasıl görülür?
\end{alistirma}
\alcozum{%
\begin{ipucu}{1}
$\mu_F$ için $\sigma$-toplamsallık: $(a,b] = \bigcup_n (a_n, b_n]$ ikili ayrık ise
$F(b) - F(a) = \sum_n (F(b_n) - F(a_n))$. Sağdan süreklilik bu eşitlik için gereklidir.
\end{ipucu}

$\mu_F(\emptyset) = F(a) - F(a) = 0$. Yarı-halka $\{(a,b]\}$ üzerinde $\sigma$-toplamsallık:
$(a,b]$ ikili ayrık $(a_n,b_n]$ birleşimi ise, $F$'nin sağdan sürekliliği ve Heine-Borel
gibi argümanlarla (ayrıntı için Billingsley~\cite{billingsley1995} §12 bakınız)
$F(b)-F(a) = \sum_n (F(b_n)-F(a_n))$ kanıtlanır.

Carathéodory genişlemesi bu ön-ölçüyü $\mathcal{B}(\mathbb{R})$'yi içeren bir
$\sigma$-cebire taşır. $F(x)=x$: $\mu_F((a,b]) = b-a$; bu Lebesgue dış ölçüsünü
aralıklarda tanımlar. $\blacksquare$%
}

\begin{alistirma}[Al.1.13]{C}{$L^p$ uzayları}
$(\Omega,\mathcal{F},\mu)$ ölçüm uzayı, $1 \leq p < \infty$.
$L^p(\mu) = \{f \text{ ölçülebilir} : \int |f|^p \, d\mu < \infty\}$
tanımlanıyor; $\|f\|_p = \left(\int |f|^p \, d\mu\right)^{1/p}$.

(a)~Hölder eşitsizliğini ifade edin ve ispatlayın:
$\|fg\|_1 \leq \|f\|_p \|g\|_q$ ($1/p + 1/q = 1$).

(b)~Minkowski eşitsizliğini kullanarak $\|\cdot\|_p$'nin üçgen eşitsizliğini
    sağladığını gösterin (ispat özeti yeterli).

(c)~$L^p \subseteq L^1$ olması ne zaman garantilenir? Karşı örnek veya koşul.
\end{alistirma}
\alcozum{%
(a) Young eşitsizliği: $ab \leq a^p/p + b^q/q$ ($a,b \geq 0$).
$|f(\omega)| / \|f\|_p$ ve $|g(\omega)| / \|g\|_q$'ya uygula ve integra al;
$\int |fg| / (\|f\|_p \|g\|_q) \leq 1/p + 1/q = 1$. $\|fg\|_1 \leq \|f\|_p \|g\|_q$.

(b) $\|f+g\|_p^p = \int|f+g|^p \leq \int|f+g|^{p-1}(|f|+|g|)$;
Hölder'i $|f+g|^{p-1} \in L^q$'ya uygula; $\|f+g\|_p \leq \|f\|_p + \|g\|_p$.

(c) $\mu(\Omega) < \infty$ ise Hölder ile $\|f\|_1 \leq \mu(\Omega)^{1/q} \|f\|_p < \infty$.
Sonsuz ölçüde $L^p \not\subseteq L^1$: $f(x) = (1+x)^{-1/p}$ üzerinde $[0,\infty)$,
$\int f^p = \infty$... daha dikkatli: $f(x) = x^{-1/2} \gostf{(0,1]} \in L^1 \setminus L^2$
ve $g(x) = (1+x)^{-1} \in L^2(\mathbb{R}) \setminus L^1(\mathbb{R})$. $\blacksquare$%
}

\begin{alistirma}[Al.1.14]{B}{Radon-Nikodym Türevi}
$\Omega = \{1, 2, 3, 4\}$, $\mathcal{F} = \mathcal{P}(\Omega)$.
$\mu(\{k\}) = 1/4$ (uniform), $\nu(\{1\}) = 1/2$, $\nu(\{2\}) = 1/4$, $\nu(\{3\}) = 1/8$, $\nu(\{4\}) = 1/8$.
(a) $\nu \ll \mu$ olduğunu doğrulayın. (b) Radon-Nikodym türevi $f = d\nu/d\mu$'yu bulun.
(c) $\nu(A) = \int_A f \, d\mu$ olduğunu $A = \{1,3\}$ için doğrulayın.
\end{alistirma}
\alcozum{%
(a) $\mu(A) = 0 \Rightarrow A = \emptyset \Rightarrow \nu(A) = 0$. Evet, $\nu \ll \mu$.\\
(b) RN türevi: $f(\omega) = \nu(\{\omega\})/\mu(\{\omega\})$ (ayrık durumda).
$f(1) = (1/2)/(1/4) = 2$; $f(2) = 1$; $f(3) = 1/2$; $f(4) = 1/2$.\\
(c) $\nu(\{1,3\}) = 1/2 + 1/8 = 5/8$.
$\int_{\{1,3\}} f \, d\mu = f(1)\mu(\{1\}) + f(3)\mu(\{3\}) = 2\cdot(1/4) + (1/2)\cdot(1/4) = 1/2 + 1/8 = 5/8$. \checkmark}

\begin{alistirma}[Al.1.15]{B}{$L^p$ Norm Hesabı}
$f(x) = e^{-x}$, $x \geq 0$; Lebesgue ölçüsü.
(a) $\|f\|_p = (\int_0^\infty e^{-px} dx)^{1/p}$'yi hesaplayın.
(b) $p \to \infty$ limitinde $\|f\|_p$'nin ne olduğunu yorumlayın.
(c) Hölder eşitsizliğiyle $\int_0^\infty e^{-x}\cos x \, dx \leq \|f\|_2\|\cos\|_{L^2([0,\infty))}$
    üst sınırını verin (dikkat: $\|\cos\|_{L^2([0,\infty))} = \infty$; bu ne anlama gelir?).
\end{alistirma}
\alcozum{%
(a) $\|f\|_p^p = \int_0^\infty e^{-px} dx = 1/p$. $\|f\|_p = p^{-1/p}$.\\
(b) $\|f\|_\infty = \sup_{x \geq 0} |e^{-x}| = e^0 = 1$. $\lim_{p\to\infty} p^{-1/p} = 1$ (çünkü $\ln(p^{-1/p}) = -\ln p/p \to 0$). $L^\infty$ normu esansiyel supremum.\\
(c) $\cos \notin L^2([0,\infty))$ çünkü $\int_0^\infty \cos^2 x \, dx = \infty$ (yarı-periyot toplamı ıraksak). Hölder uygulanamaz — $g \in L^q$ gerektirir. Bu, Hölder'in koşulsuz uygulanamayacağını gösterir.}

% =====================================================================
%  BÖLÜM 2 — Olasılık Uzayları
%  Kaynaklar: Kolmogorov 1933; Feller Vol. I Böl. I-V; Gnedenko Böl. 1
% =====================================================================

\chapter{Olasılık Uzayları}

\bolumalinti{``The theory of probability as mathematical discipline can
  and should be developed from axioms in exactly the same way as Geometry
  and Algebra.''}{Andrey Kolmogorov, \emph{Grundbegriffe der Wahrscheinlichkeitsrechnung} (1933)}

\noindent\textbf{Ön Koşullar:} Bölüm~0 (küme cebiri), Bölüm~1 (ölçüm teorisi).

\noindent\textbf{Bu Bölüm Sonraki Bölümlerde Kullanılır:}
Bölüm~3 (rasgele değişkenler),
Bölüm~4 (beklenti),
Bölüm~7 (büyük sayılar).

\noindent\textbf{Tahmini Okuma Süresi:} $\approx 8$~saat.

\vspace{0.5\baselineskip}

\begin{kazanimlar}
Bu bölümün sonunda öğrenci:
\begin{itemize}
  \item Kolmogorov aksiyomlarını ifade edebilmeli ve bu aksiyomlardan
        olasılığın temel özelliklerini (monotonluk, alt-toplamsallık,
        Poincaré formülü) türetebilmeli,
  \item Koşullu olasılığı tanımlayabilmeli, bunun bir olasılık ölçüsü
        olduğunu gösterebilmeli ve Toplam Olasılık Teoremini ispatlayabilmeli,
  \item Bağımsızlık kavramını (olaylar, olay aileleri, $\sigma$-cebirleri)
        kesin biçimde tanımlayabilmeli; ikili bağımsızlık ile karşılıklı
        bağımsızlık arasındaki farkı karşı örnekle sergileyebilmeli,
  \item Bayes teoremini ispatlayabilmeli ve önsel/sonsal olasılık yorumunu
        açıklayabilmeli,
  \item Klasik (eşit olası) olasılık uzaylarında kombinatoryal hesaplar
        yapabilmeli: permütasyon, kombinasyon, doğum günü paradoksu.
\end{itemize}
\end{kazanimlar}

\begin{motivasyon}
Bölüm~1'de ölçüm teorisinin genel çerçevesini kurduk. Bu bölümde o çerçeveyi
olasılık için özelleştiriyoruz: ölçü fonksiyonu $\mu$ yerine \emph{olasılık
ölçüsü} $\Prob$ kullanacak, $\mu(\Omega)=1$ normalizasyonu tüm yorumu değiştirecek.

Bir \emph{deneyden} söz ettiğimizde, aklımızdaki matematiksel yapı üç nesneye
sahip: mümkün tüm sonuçların kümesi $\Omega$, iyi tanımlanmış olay kavramı
$\mathcal{F}$, ve olaylara olasılık atayan $\Prob$. Kolmogorov, 1933'te bu
üçlünün tam olarak bir olasılık ölçüm uzayı $(\Omega,\mathcal{F},\Prob)$
olduğunu gösterdi.

\medskip

\textbf{Köprü kurma.} Lise matematiğinde ``eşit olası sonuçlar'' varsayımıyla
$\Prob(A) = |A|/|\Omega|$ formülünü kullandınız. Bu bölümde hem bu klasik
yaklaşımı hem de sonsuz ya da sürekli örnek uzaylı deneyleri kapsayan genel
çerçeveyi öğreneceğiz.
\end{motivasyon}

\begin{tarihselnot}
Jacob Bernoulli 1713'te (ölümünden sonra yayınlanan \emph{Ars Conjectandi}),
Pierre-Simon Laplace 1812'de (\emph{Théorie Analytique des Probabilités})
olasılığı sayısal biçimde ele aldılar; fakat her ikisi de yalnızca sonlu
eşit-olası uzaylara dayanıyordu. 19.~yüzyılda sonsuz seri ve integrallerle
olasılık hesapları yapıldıkça tutarsızlıklar ortaya çıktı (\emph{Bertrand
paradoksu}, 1888). Kolmogorov 1933'te Lebesgue ölçüm teorisini olasılığa
uyguladı ve bu paradoksları çözdü: olasılık artık bir ölçüm uzayıdır.
\end{tarihselnot}


% =====================================================================
\section{Kolmogorov Aksiyomları ve Temel Özellikler}
% =====================================================================

\begin{tanim}[Olasılık Uzayı]\label{def:olasilikuzayi}
Üç bileşenden oluşan $(\Omega, \mathcal{F}, \Prob)$ üçlüsüne
\textbf{olasılık uzayı} \emph{(probability space)} denir; burada:
\begin{itemize}
  \item $\Omega$: \textbf{örnek uzay} \emph{(sample space)} — deneyin tüm
        olası sonuçlarının kümesi; elemanlarına $\omega$ denir ve
        \textbf{sonuç} \emph{(outcome, elementary event)} adını alır.
  \item $\mathcal{F}$: $\Omega$ üzerinde bir \textbf{$\sigma$-cebiri}
        (Tanım~\ref{def:sigmacebiri}) — elemanlarına \textbf{olay} \emph{(event)} denir.
  \item $\Prob: \mathcal{F} \to [0,1]$: \textbf{olasılık ölçüsü}
        \emph{(probability measure)} — Kolmogorov'un üç aksiyomunu sağlar:
\end{itemize}
\begin{enumerate}[label=\textbf{(K\arabic*)}]
  \item $\Prob(A) \geq 0$ her $A \in \mathcal{F}$ için,
  \item $\Prob(\Omega) = 1$,
  \item $A_1, A_2, \ldots \in \mathcal{F}$ ikili ayrık ise
        $\Prob\!\left(\bigcup_{n=1}^\infty A_n\right) = \sum_{n=1}^\infty \Prob(A_n)$.
\end{enumerate}
\end{tanim}

\begin{notkutu}
Bölüm~1, Tanım~\ref{def:olasilikolucu} ile bu tanım özdeştir: olasılık ölçüsü,
$\Prob(\Omega)=1$ şartını sağlayan bir ölçüdür. Aksiyom (K3),
$\sigma$-toplamsallığın olasılık diline tercümesidir.
\end{notkutu}

\begin{ornek}[Olasılık uzayı örnekleri]\label{ex:olasilikorn}
\begin{enumerate}[label=(\alph*)]
  \item \textbf{Adil zar:} $\Omega = \{1,2,3,4,5,6\}$,
        $\mathcal{F} = \mathcal{P}(\Omega)$,
        $\Prob(\{k\}) = 1/6$ her $k$ için. Sayma ölçüsünün
        $1/|\Omega|$ ile ölçeklenmesi.

  \item \textbf{Düzgün dağılım $[0,1]$ üzerinde:} $\Omega = [0,1]$,
        $\mathcal{F} = \mathcal{B}([0,1])$,
        $\Prob = \lambda\!\restriction_{[0,1]}$ (Lebesgue ölçüsü).
        $\Prob((a,b)) = b-a$ her $0 \leq a < b \leq 1$ için.

  \item \textbf{Sonsuz yazı-tura:} $\Omega = \{0,1\}^{\mathbb{N}}$
        (sonsuz ikili diziler), $\mathcal{F}$ silindir kümelerin ürettiği
        $\sigma$-cebiri. Kolmogorov'un uzatma teoremiyle $\Prob$ inşa edilir
        (Bölüm~7'de ele alınacak).
\end{enumerate}
\end{ornek}

\begin{teorem}[Olasılığın Temel Özellikleri]\label{thm:olasiliktemel}
$(\Omega, \mathcal{F}, \Prob)$ bir olasılık uzayı, $A, B, A_n \in \mathcal{F}$.
\begin{enumerate}[label=(\roman*)]
  \item $\Prob(\emptyset) = 0$.
  \item \textbf{Sonlu toplamsallık:} $A_1,\ldots,A_n$ ikili ayrık $\Rightarrow$
        $\Prob\!\bigl(\bigcup_{k=1}^n A_k\bigr) = \sum_{k=1}^n \Prob(A_k)$.
  \item \textbf{Tümleyen:} $\Prob(A^c) = 1 - \Prob(A)$.
  \item \textbf{Monotonluk:} $A \subseteq B \Rightarrow \Prob(A) \leq \Prob(B)$.
  \item \textbf{Fark:} $A \subseteq B \Rightarrow \Prob(B \setminus A) = \Prob(B) - \Prob(A)$.
  \item \textbf{Toplam:} $\Prob(A \cup B) = \Prob(A) + \Prob(B) - \Prob(A \cap B)$.
  \item \textbf{Alt-toplamsallık (Boole):}
        $\Prob\!\bigl(\bigcup_n A_n\bigr) \leq \sum_n \Prob(A_n)$.
  \item \textbf{Aşağıdan süreklilik:} $A_n \nearrow A \Rightarrow \Prob(A_n) \nearrow \Prob(A)$.
  \item \textbf{Yukarıdan süreklilik:} $A_n \searrow A \Rightarrow \Prob(A_n) \searrow \Prob(A)$.
\end{enumerate}
\end{teorem}

\begin{proof}
Tümü Bölüm~1, Teorem~\ref{thm:olcuozellik}'den $\mu = \Prob$ ile alınır.
Ek olarak: (iii) $A \cup A^c = \Omega$ ayrık; (K2)'den $1 = \Prob(A)+\Prob(A^c)$.
(vi) $A \cup B = A \cup (B \setminus A)$ ayrık; (ii) + (v).
Yukarıdan süreklilik için $\Prob(\Omega)=1<\infty$; Teorem~\ref{thm:olcuozellik}(v)
uygulanabilir.
\end{proof}

\begin{notkutu}
(ix) için $\Prob(A_1) < \infty$ yerine $\Prob(\Omega) = 1 < \infty$ her zaman
sağlandığından, olasılık uzayında yukarıdan süreklilik \emph{koşulsuz} geçerlidir.
Bu, genel ölçüm teorisine kıyasla önemli bir kolaylıktır.
\end{notkutu}

\begin{teorem}[Poincaré Dahil-Hariç Formülü]\label{thm:poincare}
Sonlu aile $A_1, \ldots, A_n \in \mathcal{F}$ için:
\[
  \Prob\!\left(\bigcup_{k=1}^n A_k\right)
  = \sum_{k} \Prob(A_k)
  - \sum_{k < l} \Prob(A_k \cap A_l)
  + \sum_{k < l < m} \Prob(A_k \cap A_l \cap A_m)
  - \cdots
  + (-1)^{n-1} \Prob(A_1 \cap \cdots \cap A_n).
\]
\end{teorem}

\begin{proof}
\textit{Fikir:} Bir $\omega \in \bigcup_k A_k$, tam olarak $r$ adet $A_k$'da
bulunsun. Bu $\omega$'nın formülün sağ tarafına katkısını hesaplayın.

\medskip

$\omega$'nın sol tarafa katkısı $1$'dir. Sağ tarafta: $\omega$, $\binom{r}{1}$
adet tekli kümede, $\binom{r}{2}$ adet ikili kesişimde, \ldots yer alır.
Katkı:
\[
  \binom{r}{1} - \binom{r}{2} + \binom{r}{3} - \cdots + (-1)^{r-1}\binom{r}{r}
  = 1 - (1-1)^r = 1.
\]
(Binom açılımından $\sum_{k=0}^r (-1)^k \binom{r}{k} = 0$.)
$\omega \notin \bigcup_k A_k$ ise her iki taraftaki katkı $0$'dır.
\end{proof}

\begin{ornek}[Poincaré --- iki ve üç olay]\label{ex:poincare}
\textbf{İki olay:}
$\Prob(A \cup B) = \Prob(A) + \Prob(B) - \Prob(A \cap B)$ (Teorem~\ref{thm:olasiliktemel}(vi)).

\textbf{Üç olay:}
$\Prob(A \cup B \cup C) = \Prob(A) + \Prob(B) + \Prob(C)
- \Prob(A \cap B) - \Prob(A \cap C) - \Prob(B \cap C)
+ \Prob(A \cap B \cap C)$.
\end{ornek}

\begin{mikroozet}[§2.1 özeti]
Olasılık uzayı $(\Omega,\mathcal{F},\Prob)$: örnek uzay + $\sigma$-cebiri (olaylar) +
olasılık ölçüsü. Aksiyomlardan: tümleme, monotonluk, süreklilik, Poincaré formülü.
Olasılık uzayında yukarıdan süreklilik koşulsuzdur.
\end{mikroozet}


% =====================================================================
\section{Neredeyse Kesin Olaylar ve Sıfır Olasılıklı Olaylar}
% =====================================================================

\begin{tanim}[Neredeyse Kesin Olay]\label{def:neredeyse_kesin}
$A \in \mathcal{F}$:
\begin{itemize}
  \item \textbf{neredeyse kesin (n.k.)} \emph{(almost surely, a.s.)} adını alır;
        eğer $\Prob(A) = 1$.
  \item \textbf{sıfır olasılıklı} \emph{(null event)} adını alır;
        eğer $\Prob(A) = 0$.
  \item $A^c$ sıfır olasılıklı ise $A$ neredeyse kesindir ve $A^c$'ye
        \textbf{ihmal edilebilir olay} \emph{(negligible event)} denir.
\end{itemize}
\end{tanim}

\begin{onerme}[Sayılabilir birleşim n.k.]\label{prop:nk_birlesim}
$A_n$ n.k. her $n$ için ise $\bigcap_{n=1}^\infty A_n$ n.k.'dir.
\end{onerme}

\begin{proof}
$\Prob(A_n^c) = 0$ her $n$. Alt-toplamsallık:
$\Prob\!\bigl(\bigcup_n A_n^c\bigr) \leq \sum_n \Prob(A_n^c) = 0$.
De~Morgan: $\left(\bigcap_n A_n\right)^c = \bigcup_n A_n^c$; $\Prob(\text{tümleyen}) = 0$;
yani $\Prob(\bigcap_n A_n) = 1$.
\end{proof}

\begin{hataanalizi}[Sıfır olasılık $\neq$ imkânsız]
$\Omega = [0,1]$, $\Prob = \lambda$. Herhangi bir $x \in [0,1]$ için
$\Prob(\{x\}) = \lambda(\{x\}) = 0$. Ama $\{x\}$ gerçekleşebilir
(o sonuç fiilen mümkündür); yalnızca Lebesgue ölçüsü sıfır atamıştır.
``Sıfır olasılık = imkânsız'' yerine ``sıfır olasılık = ihmal edilebilir'' demek daha doğrudur.
\end{hataanalizi}

\begin{tanim}[Tam Olasılık Uzayı]\label{def:tamuzay}
$(\Omega,\mathcal{F},\Prob)$ \textbf{tam} \emph{(complete)} adını alır; eğer
$\Prob(N)=0$ ve $A \subseteq N$ ise $A \in \mathcal{F}$ (ve $\Prob(A)=0$).
Lebesgue ölçüm uzayı tam; Borel kümeler üzerinde Lebesgue ölçüsü tam değil.
\end{tanim}


% =====================================================================
\section{Koşullu Olasılık}
% =====================================================================

\begin{kesif}[Koşullu olasılık sezgisi]
Adil bir zar atılıyor. ``Çift sayı geldi'' ($B$) bilgisiyle, ``6 geldi'' ($A$)
olayının olasılığı nedir? Sezginizi yazın, ardından Tanım~\ref{def:kos_olasilik}'i okuyun.
\ogrencialani[1.5cm]
\end{kesif}

\begin{tanimgerekce}
$B$ olduğu biliniyorsa örnek uzayımız $\Omega$ yerine $B$ olur. $A$'nın bu
küçülmüş uzaydaki payı $\Prob(A \cap B) / \Prob(B)$'dir. Payda normalizasyonu
($B$'nin ölçüsünü 1 yapıyor) sağlıyor.
\end{tanimgerekce}

\begin{tanim}[Koşullu Olasılık]\label{def:kos_olasilik}
$(\Omega,\mathcal{F},\Prob)$ olasılık uzayı, $B \in \mathcal{F}$, $\Prob(B) > 0$.
$B$'ye koşullu $A$'nın olasılığı:
\[
  \Prob(A \mid B) \coloneqq \frac{\Prob(A \cap B)}{\Prob(B)}, \quad A \in \mathcal{F}.
\]
\end{tanim}

\begin{teorem}[Koşullu Olasılık Bir Olasılık Ölçüsüdür]\label{thm:kos_olcu}
$\Prob(B) > 0$ ise $A \mapsto \Prob(A \mid B)$ bir olasılık ölçüsüdür;
yani $(\Omega, \mathcal{F}, \Prob(\cdot \mid B))$ bir olasılık uzayıdır.
\end{teorem}

\begin{proof}
(K1) $\Prob(A \cap B) \geq 0$; $\Prob(B) > 0$; dolayısıyla $\Prob(A \mid B) \geq 0$.

(K2) $\Prob(\Omega \mid B) = \Prob(\Omega \cap B)/\Prob(B) = \Prob(B)/\Prob(B) = 1$.

(K3) $(A_n)$ ikili ayrık; $(A_n \cap B)$ da ikili ayrık:
\[
  \Prob\!\left(\bigcup_n A_n \,\Big|\, B\right)
  = \frac{\Prob\!\left(\bigcup_n (A_n \cap B)\right)}{\Prob(B)}
  = \frac{\sum_n \Prob(A_n \cap B)}{\Prob(B)}
  = \sum_n \Prob(A_n \mid B).  \qed
\]
\end{proof}

\begin{teorem}[Toplam Olasılık Teoremi]\label{thm:toplam_olasilik}
$(B_n)_{n=1}^N$ ($N \leq \infty$) ikili ayrık, $\bigcup_n B_n = \Omega$,
$\Prob(B_n) > 0$ olsun. Her $A \in \mathcal{F}$ için:
\[
  \Prob(A) = \sum_{n=1}^N \Prob(A \mid B_n)\, \Prob(B_n).
\]
\end{teorem}

\begin{proof}
$A = \bigcup_n (A \cap B_n)$ ayrık birleşim; $\sigma$-toplamsallık:
$\Prob(A) = \sum_n \Prob(A \cap B_n) = \sum_n \Prob(A \mid B_n)\Prob(B_n)$.
\end{proof}

\begin{ornek}[Toplam Olasılık --- fabrika sorunu]\label{ex:fabrika}
Üç fabrika sırasıyla üretimin $\%50$, $\%30$, $\%20$'sini yapıyor ve
hatalı oranları sırasıyla $\%2$, $\%4$, $\%5$. Rastgele seçilen bir
ürünün hatalı olma olasılığı?

$B_1, B_2, B_3$: fabrikalara ait ürünler. $A$: hatalı.
\[
  \Prob(A) = 0.02 \times 0.5 + 0.04 \times 0.3 + 0.05 \times 0.2
           = 0.010 + 0.012 + 0.010 = 0.032.
\]
\end{ornek}

\begin{teorem}[Çarpım Kuralı]\label{thm:carpim}
$A_1, \ldots, A_n \in \mathcal{F}$, $\Prob(A_1 \cap \cdots \cap A_{n-1}) > 0$:
\[
  \Prob(A_1 \cap \cdots \cap A_n)
  = \Prob(A_1)\,\Prob(A_2 \mid A_1)\,\Prob(A_3 \mid A_1 \cap A_2)
    \cdots \Prob(A_n \mid A_1 \cap \cdots \cap A_{n-1}).
\]
\end{teorem}

\begin{proof}
Tümevarım: $n=2$ tanımdan. $n$'de geçerliyse $n+1$:
son çarpanı $\Prob(A_1 \cap \cdots \cap A_n) \cdot \Prob(A_{n+1} \mid A_1 \cap \cdots \cap A_n)$
ile genişlet.
\end{proof}

\begin{mikroozet}[§2.3 özeti]
$\Prob(A|B) = \Prob(A \cap B)/\Prob(B)$. Koşullu olasılık bir olasılık ölçüsüdür.
Toplam olasılık: bölüntü üzerinden toplam. Çarpım kuralı: zincirleme hesap.
\end{mikroozet}


% =====================================================================
\section{Bayes Teoremi}
% =====================================================================

\begin{teorem}[Bayes Teoremi]\label{thm:bayes}
$(B_n)$ tam bir bölüntü oluştursun ($\Prob(B_n)>0$, $\bigcup_n B_n = \Omega$,
ikili ayrık). Her $A \in \mathcal{F}$, $\Prob(A)>0$ için:
\[
  \Prob(B_k \mid A)
  = \frac{\Prob(A \mid B_k)\,\Prob(B_k)}{\sum_n \Prob(A \mid B_n)\,\Prob(B_n)}.
\]
\end{teorem}

\begin{proof}
Pay: $\Prob(A \mid B_k)\Prob(B_k) = \Prob(A \cap B_k)$.
Payda: Toplam Olasılık Teoremi (Teorem~\ref{thm:toplam_olasilik}):
$\sum_n \Prob(A \mid B_n)\Prob(B_n) = \Prob(A)$.
Oran: $\Prob(A \cap B_k)/\Prob(A) = \Prob(B_k \mid A)$.
\end{proof}

\begin{notkutu}
Bayes teoremindeki terminoloji:
\begin{itemize}
  \item $\Prob(B_k)$: \textbf{önsel olasılık} \emph{(prior probability)} — $A$ gözlemlenmeden önceki bilgi.
  \item $\Prob(A \mid B_k)$: \textbf{olabilirlik} \emph{(likelihood)} — $B_k$ doğruysa $A$'nın olasılığı.
  \item $\Prob(B_k \mid A)$: \textbf{sonsal olasılık} \emph{(posterior probability)} — $A$ gözlemlendikten sonra.
  \item $\Prob(A)$: \textbf{marjinal (normalleştirme sabiti)}.
\end{itemize}
\end{notkutu}

\begin{ornek}[Bayes --- tıbbi test]\label{ex:tibbi_test}
Hastalığın görülme sıklığı $\Prob(H) = 0.001$. Test duyarlılığı
$\Prob(T^+ \mid H) = 0.99$; yanlış pozitif $\Prob(T^+ \mid H^c) = 0.05$.
Pozitif test sonucunda hasta olma olasılığı?

\[
  \Prob(H \mid T^+)
  = \frac{0.99 \times 0.001}{0.99 \times 0.001 + 0.05 \times 0.999}
  = \frac{0.00099}{0.00099 + 0.04995}
  \approx 0.0194.
\]

Sonuç $\approx \%2$: Hastalık çok nadir olduğunda bile \emph{çok iyi} bir testin
pozitif sonucu gerçek hastalığı büyük olasılıkla değil, yanlış pozitifi gösterir.
\end{ornek}

\begin{hataanalizi}[Savcının yanılgısı (Prosecutor's fallacy)]
$\Prob(E \mid \text{Suçsuz})$ çok küçükse, $\Prob(\text{Suçsuz} \mid E)$'nin
de küçük olduğu \emph{sonucu çıkmaz}. Örnek~\ref{ex:tibbi_test}'te
$\Prob(T^+ \mid H^c) = 0.05$ küçüktür, ama $\Prob(H \mid T^+) \approx 0.02$
daha küçüktür; çünkü $\Prob(H)$ çok küçüktür. Olası suçluluk oranı (önsel)
hesaba katılmadan $\Prob(E \mid \text{Suçsuz})$'u mahkûmiyet gerekçesi yapmak
Bayes teoremini göz ardı etmektir.
\end{hataanalizi}

\begin{disiplinlerarasibaglan}[Makine öğrenmesi ve Naive Bayes]
Bayes teoremi makine öğrenmesinde metin sınıflandırmanın temelini oluşturur.
\emph{Naive Bayes sınıflandırıcısı}, her özelliğin (kelime) bağımsız olduğunu
varsayarak bir belgenin hangi sınıfa ait olduğunu sınıflandırır:
$\Prob(\text{sınıf} \mid \text{kelimeler}) \propto \Prob(\text{sınıf}) \prod_i \Prob(\text{kelime}_i \mid \text{sınıf})$.
``Naive'' (saf) adı bağımsızlık varsayımının kaba olmasından gelir; ama pratikte
şaşırtıcı biçimde iyi çalışır.
\end{disiplinlerarasibaglan}

\begin{mikroozet}[§2.4 özeti]
Bayes teoremi: sonsal $\propto$ olabilirlik $\times$ önsel.
Payda = Toplam Olasılık. Nadir hastalıklar, Savcı yanılgısı, Naive Bayes:
hepsi aynı formülün görünümleri.
\end{mikroozet}


% =====================================================================
\section{Bağımsızlık}
% =====================================================================

\begin{kesif}[Bağımsızlık sezgisi]
``$A$ gerçekleşmesinin $B$'ye etkisi yoktur'' ne demektir?
$\Prob(A \mid B) = \Prob(A)$ formülü bunu yakalar; ama $\Prob(B)=0$ durumunda
bu yazılamaz. Bundan bağımsız bir tanım olasını düşünün.
\ogrencialani[1.5cm]
\end{kesif}

\begin{tanim}[Bağımsız Olaylar]\label{def:bagimsiz_olay}
$A, B \in \mathcal{F}$ olayları \textbf{bağımsız} \emph{(independent)} adını alır;
eğer $\Prob(A \cap B) = \Prob(A)\,\Prob(B)$.

Sonlu aile $\{A_1,\ldots,A_n\}$ \textbf{karşılıklı bağımsız}
\emph{(mutually independent)} adını alır; eğer her $I \subseteq \{1,\ldots,n\}$,
$I \neq \emptyset$ için:
\[
  \Prob\!\left(\bigcap_{i \in I} A_i\right) = \prod_{i \in I} \Prob(A_i).
\]

Sonsuz aile $\{A_i : i \in I\}$ karşılıklı bağımsızdır; eğer her sonlu
$J \subseteq I$ için karşılıklı bağımsızlık koşulu sağlanıyorsa.
\end{tanim}

\begin{kritiksorgu}[İkili bağımsızlık $\neq$ karşılıklı bağımsızlık]
İkili bağımsızlık ($\Prob(A_i \cap A_j) = \Prob(A_i)\Prob(A_j)$ tüm $i \neq j$
için) karşılıklı bağımsızlığı gerektirmez. Aşağıdaki karşı örnek bunu açıklar.
\end{kritiksorgu}

\begin{karsiornek}[Bernstein'ın örneği]\label{ex:bernstein}
Dengeli bir dört yüzlü zar: yüzler $\{1,2,3,4\}$, her biri $1/4$ olasılıklı.
$A = \{1,2\}$, $B = \{1,3\}$, $C = \{1,4\}$ olsun.

$\Prob(A)=\Prob(B)=\Prob(C)=1/2$.
$\Prob(A \cap B) = \Prob(\{1\}) = 1/4 = \Prob(A)\Prob(B)$. Benzer: $A,B$; $A,C$; $B,C$ ikili bağımsız.

Ama $\Prob(A \cap B \cap C) = 1/4 \neq 1/8 = \Prob(A)\Prob(B)\Prob(C)$.

$A, B, C$ ikili bağımsız ama karşılıklı bağımsız değil!
\end{karsiornek}

\begin{tanim}[$\sigma$-Cebirlerinin Bağımsızlığı]\label{def:sigmababag}
$(\mathcal{G}_i)_{i \in I}$ $\mathcal{F}$'nin alt $\sigma$-cebirleri ailesine
\textbf{bağımsız} denir; eğer her sonlu $J \subseteq I$ ve her
$A_j \in \mathcal{G}_j$ ($j \in J$) için:
\[
  \Prob\!\left(\bigcap_{j \in J} A_j\right) = \prod_{j \in J} \Prob(A_j).
\]
\end{tanim}

\begin{onerme}[$\pi$-sistemi ile bağımsızlık kontrolü]\label{prop:pisys_bag}
$\mathcal{P}_i$ bir $\pi$-sistemi ($\cap$ altında kapalı) ve
$\sigma(\mathcal{P}_i) = \mathcal{G}_i$ olsun. Eğer $(\mathcal{P}_i)$ ailesi
bağımsız ise $(\mathcal{G}_i)$ de bağımsızdır.
\end{onerme}

\begin{proof}
Dynkin'in $\pi$-$\lambda$ teoremi: $\pi$-sistemlerin bağımsızlığı,
ürettikleri $\sigma$-cebirlerine uzar.
Ayrıntı için Billingsley~\cite{billingsley1995} §4 bakınız.
\end{proof}

\begin{onerme}[Bağımsızlığın Korunması]\label{prop:bag_korunma}
$A$ ve $B$ bağımsız ise $A$ ve $B^c$; $A^c$ ve $B$; $A^c$ ve $B^c$ de bağımsızdır.
\end{onerme}

\begin{proof}
$\Prob(A \cap B^c) = \Prob(A) - \Prob(A \cap B) = \Prob(A) - \Prob(A)\Prob(B)
= \Prob(A)(1 - \Prob(B)) = \Prob(A)\Prob(B^c)$.
Diğerleri benzer.
\end{proof}

\begin{ornek}[Bağımsız denemeler uzayı]\label{ex:bag_deneme}
Adil yazı-tura $n$ kez atılsın. $\Omega = \{Y,T\}^n$, $|\Omega|=2^n$,
$\Prob(\{\omega\}) = (1/2)^n$ her $\omega$ için.

$A_k = \{k\text{-inci atışta tura}\}$. $A_1, \ldots, A_n$ karşılıklı bağımsızdır:
$\Prob(A_{k_1} \cap \cdots \cap A_{k_r}) = (1/2)^r = \prod_j \Prob(A_{k_j})$.
\end{ornek}

\begin{mikroozet}[§2.5 özeti]
Bağımsızlık: $\Prob(A \cap B) = \Prob(A)\Prob(B)$.
Karşılıklı bağımsızlık: tüm alt aileler için çarpım formülü.
İkili $\not\Rightarrow$ karşılıklı (Bernstein örneği).
Tümleyen bağımsızlığı korur.
\end{mikroozet}


% =====================================================================
\section{Klasik Kombinatoryal Olasılık}
% =====================================================================

\begin{motivasyon}[§2.6]
Sonlu, eşit-olası $\Omega$ ile $\Prob(A) = |A|/|\Omega|$ formülü, olasılık
hesabını sayma problemine indirgeler. Bu alt bölüm önemli sayma tekniklerini
ve klasik paradoksları içeriyor.
\end{motivasyon}

\begin{tanim}[Klasik Olasılık Modeli]\label{def:klasik}
$\Omega$ sonlu, $\mathcal{F} = \mathcal{P}(\Omega)$,
$\Prob(A) = |A|/|\Omega|$ her $A \subseteq \Omega$ için.
\end{tanim}

\subsection{Temel Sayma İlkeleri}

\begin{teorem}[Çarpım İlkesi]\label{thm:carpim_ilkesi}
İlk seçim $m$, ikinci seçim $n$ yolda yapılabiliyorsa, sıralı çift sayısı $mn$'dir.
Genel: $k$ aşamalı seçimde $n_1 \cdots n_k$ olasılık.
\end{teorem}

\begin{tanim}[Permütasyon ve Kombinasyon]\label{def:perm_komb}
\begin{itemize}
  \item $n$ elemanlı kümeden $r$ elemanın \textbf{sıralı} seçimi sayısı:
        $P(n,r) = n(n-1)\cdots(n-r+1) = \dfrac{n!}{(n-r)!}$.
  \item $n$ elemanlı kümeden $r$ elemanın \textbf{sırasız} seçimi sayısı:
        $\binom{n}{r} = \dfrac{n!}{r!\,(n-r)!}$.
  \item $n$ elemanlı kümenin \textbf{tam sıralaması}: $n!$
\end{itemize}
\end{tanim}

\begin{teorem}[Binom Teoremi]\label{thm:binom}
$(x+y)^n = \sum_{k=0}^n \binom{n}{k} x^k y^{n-k}$.
\end{teorem}

\begin{proof}
$(x+y)^n$'de her çarpımı açınca $n$ faktörün her birinden $x$ ya da $y$ seçilir.
$x^k y^{n-k}$ terimini veren seçim sayısı $\binom{n}{k}$ (hangi $k$ faktörden $x$ alınacağı).
\end{proof}

\begin{ornek}[Doğum günü paradoksu]\label{ex:dogumgunu}
$n$ kişilik bir grupta en az iki kişinin aynı doğum gününe sahip olma
olasılığı ne zaman $1/2$'yi aşar?

$A^c$: hiçbir iki kişi aynı günde doğmamış. $|\Omega| = 365^n$,
$|A^c| = 365 \cdot 364 \cdots (365-n+1) = P(365,n)$.
\[
  \Prob(A^c) = \frac{365!/(365-n)!}{365^n} = \prod_{k=0}^{n-1}\frac{365-k}{365}.
\]
$n=23$: $\Prob(A^c) \approx 0.493$; $\Prob(A) \approx 0.507 > 1/2$.

\textbf{Paradoks neden şaşırtıcı?} Biz $n/365$'i düşünürüz; oysa olası çift
sayısı $\binom{n}{2} \approx n^2/2$'dir.
\end{ornek}

\begin{ornek}[Zindan sorunu --- Monty Hall]\label{ex:montyhall}
Üç kapının arkasında: bir araba, iki keçi. Bir kapı seçtiniz; sunucu (keçi
olan) başka bir kapıyı açıyor ve değiştirmenizi öneriyor. Değişmeli misiniz?

$B_k$: arabanın $k$'inci kapıda olması, $k=1,2,3$.
Siz 1'i seçtiniz; sunucu 3'ü açtı.

$\Prob(B_1 \mid \text{3 açıldı}) = \Prob(\text{3 açıldı} \mid B_1)\Prob(B_1) / \Prob(\text{3 açıldı})$.

$\Prob(\text{3 açıldı} \mid B_1) = 1/2$ (sunucu 2 veya 3'ü seçebilir),
$\Prob(\text{3 açıldı} \mid B_2) = 1$ (zorunlu: 3 açılır),
$\Prob(\text{3 açıldı} \mid B_3) = 0$ (3 açılamaz).
$\Prob(\text{3 açıldı}) = 1/2 \cdot 1/3 + 1 \cdot 1/3 = 1/2$.
$\Prob(B_1 \mid \text{3 açıldı}) = (1/6)/(1/2) = 1/3$.
$\Prob(B_2 \mid \text{3 açıldı}) = (1/3)/(1/2) = 2/3$.

\textbf{Değiştirin!} Değiştirme olasılığı $2/3 > 1/3$.
\end{ornek}

\begin{kendinleifade}
Monty Hall sorusunda sunucunun rastgele değil, \emph{her zaman keçi olan}
bir kapıyı açtığını not edin. Sunucu rastgele açsaydı sonuç değişir miydi?
\ogrencialani[2cm]
\end{kendinleifade}

\begin{ornek}[Hat çekme paradoksu]\label{ex:hatchekme}
$n$ kişi arasında $n$ hat dağıtılıyor; $k$ hat kısa ($k < n$). Herkes
rastgele bir hat çekiyor. İlk çeken mi, son çeken mi avantajlı?

$A_j$: $j$-inci kişi kısa hat çekiyor.
$\Prob(A_j) = k/n$ her $j$ için (simetri).

İlk çeken için: $k/n$ açıkça.
İkinci çeken için: $A_1^c$ koşulunda $k/(n-1)$;
$\Prob(A_2) = \Prob(A_2|A_1)\Prob(A_1) + \Prob(A_2|A_1^c)\Prob(A_1^c)$
$= (k-1)/(n-1) \cdot k/n + k/(n-1) \cdot (n-k)/n = k(n-1)/n(n-1) = k/n$.

Hiçbir sıra avantajlı değildir.
\end{ornek}

\begin{mikroozet}[§2.6 özeti]
Klasik model: eşit olası, $\Prob(A) = |A|/|\Omega|$. Çarpım ilkesi, permütasyon,
kombinasyon. Doğum günü: $n=23$'te $>1/2$. Monty Hall: Bayes uygulaması.
Hat çekme: sıranın önemi yok.
\end{mikroozet}


% =====================================================================
\section{Koşullu Bağımsızlık ve Bayesçi Ağlar}
\label{sec:kosullu_bagimsizlik}
% =====================================================================

\begin{tanim}[Koşullu Bağımsızlık]
\label{def:kosullu_bagimsizlik}
$A$ ve $B$ olayları $C$ verildiğinde \textbf{koşullu bağımsız}
(conditionally independent) denir ve $A \perp B \mid C$ gösterilir, eğer
\[
  \Prob(A \cap B \mid C) = \Prob(A \mid C)\,\Prob(B \mid C).
\]
\end{tanim}

\begin{ornek}[Koşullu Bağımsızlık $\neq$ Mutlak Bağımsızlık]
Bir kişi iki gün üst üste hasta olsun ($A$: 1. gün hasta, $B$: 2. gün hasta).
$C$ = "grip mevsiminde". $A$ ve $B$ koşulsuz bağımlı (grip yakalanmak sürer).
$C$ verildiğinde $A$ ve $B$ koşullu bağımsız olabilir: gribi olan kişide
birinci ve ikinci günün durumları hastalığın seyrine bağlı, ama $C$'yi bilirsek
korelasyon büyük ölçüde $C$'den kaynaklanıyor.

Ters yön: $A \perp B$ ama $A \not\perp B \mid C$.
Örnek: $A = \{\text{Türkiye'de yağmur}\}$, $B = \{\text{Japonya'da yağmur}\}$, $C = \{\text{El Niño}\}$.
$A \perp B$ (bağımsız ülkeler) ama $A \not\perp B \mid C$ (El Niño her ikisini etkiler).
\end{ornek}

\begin{teorem}[Koşullu Bağımsızlık Özellikleri]
\label{thm:kosullu_bagimsizlik}
$A \perp B \mid C$ ise:
\begin{enumerate}[(i)]
  \item $A^c \perp B \mid C$ (tümleme kapalı).
  \item $\Prob(A \cap B \cap C) = \Prob(A \mid C)\Prob(B \mid C)\Prob(C)$ (zincir kuralı ile).
  \item $A \perp B \mid C$ ve $A \perp D \mid C$ ise $A \perp (B \cup D) \mid C$ (belirli koşullarla).
\end{enumerate}
\end{teorem}

\begin{proof}[(i)]
$\Prob(A^c \cap B \mid C) = \Prob(B \mid C) - \Prob(A \cap B \mid C)
= \Prob(B \mid C) - \Prob(A \mid C)\Prob(B \mid C) = (1-\Prob(A \mid C))\Prob(B \mid C)
= \Prob(A^c \mid C)\Prob(B \mid C)$.
\end{proof}

\begin{tanim}[Olasılıklı Grafik Model — Bayesçi Ağ]
\label{def:bayesci_ag}
Rasgele değişkenler $X_1, \ldots, X_n$; bir yönlü asiklik graf (DAG).
Her $X_i$, ebeveynleri $\mathrm{Pa}(X_i)$ koşulunda diğer atalardan bağımsız:
\[
  X_i \perp \{X_j : j \notin \mathrm{Pa}(X_i), j \text{ torun değil}\} \mid \mathrm{Pa}(X_i).
\]
Ortak dağılım: $\Prob(X_1 = x_1, \ldots, X_n = x_n) = \prod_{i=1}^n \Prob(X_i = x_i \mid \mathrm{Pa}(X_i))$.
\end{tanim}

\begin{ornek}[Basit Bayesçi Ağ]
$D$ (hastalık) $\to$ $S$ (semptom) $\to$ $T$ (test).
$S \perp T \mid D$ değil, ama $D \perp T \mid S$ (test sadece semptomla ilgili).
Ortak: $\Prob(D, S, T) = \Prob(D)\Prob(S \mid D)\Prob(T \mid S)$.
$\Prob(D \mid T)$: Bayes teoremi ile hesaplanır.
\end{ornek}

\begin{disiplinlerarasibaglan}
\textbf{Yapay zeka:} Bayesçi ağlar, belirsizlik altında akıl yürütmenin
temel matematiksel çerçevesidir; teşhis sistemleri, konuşma tanıma ve robotik karar sistemlerinde kullanılır.\\
\textbf{Epidemiyoloji:} Bulaşıcı hastalık modelleri DAG yapısıyla koşullu bağımsızlıkları kodlar.\\
\textbf{Makine öğrenmesi:} Naive Bayes sınıflandırıcısı, özellikler sınıf koşulunda bağımsız varsayımına dayanır.
\end{disiplinlerarasibaglan}

% =====================================================================
\section{Bilgi Teorisi Temelleri: Entropi ve Karşılıklı Bilgi}
\label{sec:entropi}
% =====================================================================

\begin{tanim}[Shannon entropisi]\label{def:entropi}
Sonlu $\Omega=\{x_1,\ldots,x_n\}$, $p_i=P(X=x_i)$. $X$'in \emph{Shannon entropisi} (\emph{entropy}):
\[
  H(X) = -\sum_{i=1}^n p_i \log_2 p_i,\quad (0\log 0:=0)
\]
bit cinsinden. $H(X)\geq0$; $H(X)=0\iff X$ deterministik; $H(X)\leq\log_2 n$ (maksimum: düzgün dağılım).
\end{tanim}

\begin{teorem}[Maksimum entropi]\label{thm:maks_entropi}
$P(X=x_i)=1/n$ için $H(X)=\log_2 n$; bu $n$-durumlu herhangi bir dağılımın maksimum entropisidir.
\end{teorem}

\begin{proof}
Jensen eşitsizliği: $H(p)=-\sum p_i\log p_i\leq\log_2(\sum p_i/p_i \cdot n^{-1})^{-1}$ ... Doğrudan: Gibbs eşitsizliği, $-\sum p_i\log p_i\leq-\sum p_i\log(1/n)=\log n$, eşitlik $p_i=1/n$.
\end{proof}

\begin{tanim}[Koşullu entropi ve karşılıklı bilgi]\label{def:karsılikli_bilgi}
\begin{align*}
  H(X|Y) &= -\sum_{x,y}P(X=x,Y=y)\log P(X=x|Y=y) \quad \text{(koşullu entropi)}\\
  I(X;Y) &= H(X)-H(X|Y) = H(Y)-H(Y|X) \quad \text{(karşılıklı bilgi)}\\
  D_{\mathrm{KL}}(P\|Q) &= \sum_x P(x)\log\frac{P(x)}{Q(x)} \geq 0 \quad \text{(KL iraksaması)}
\end{align*}
$X\perp Y\iff I(X;Y)=0$.
\end{tanim}

\begin{teorem}[Veri işleme eşitsizliği]\label{thm:veri_isleme}
$X\to Y\to Z$ Markov zinciri ise $I(X;Z)\leq I(X;Y)$: yeterli istatistikler bilgiyi korur, diğer işlemler azaltır.
\end{teorem}

\begin{ornek}[KL iraksaması ile olabilirlik]\label{ex:kl_olasilik}
$P$ gerçek dağılım, $Q_\theta$ model. $D_{\mathrm{KL}}(P\|Q_\theta)=H(P,Q_\theta)-H(P)$; çapraz entropi minimize etmek KL minimize etmekle eşdeğer. MLE = empirik dağılıma en küçük KL olan model.
\end{ornek}

% =====================================================================
\section{Olasılık Uzaylarının Genişlemesi: Sonsuz Dizi Uzayları}
\label{sec:sonsuz_uzay}
% =====================================================================

\begin{tanim}[Ürün uzayı]\label{def:urun_uzayi}
$(\Omega_n,\mathcal{F}_n,P_n)$ olasılık uzayları. Sonsuz ürün $\Omega=\prod_{n=1}^\infty\Omega_n$, $\mathcal{F}=\bigotimes_{n=1}^\infty\mathcal{F}_n$ (silindir kümeler tarafından üretilen $\sigma$-cebir).
\end{tanim}

\begin{teorem}[Kolmogorov genişleme teoremi]\label{thm:kolmogorov_genisleme}
Her tutarlı sonlu boyutlu dağılım ailesi $\{P_{n_1,\ldots,n_k}\}$ için, bu dağılımları marjinal olarak veren tek bir olasılık ölçüsü $P$ sonsuz ürün uzayında mevcuttur.
\end{teorem}

\begin{ornek}[İ.i.d.\ dizi uzayı]\label{ex:iid_uzay}
$X_1,X_2,\ldots\iid F$. Kolmogorov teoremi bu dizinin bir olasılık uzayında var olduğunu garanti eder. Bu, GBSK ve MLT ispatlarının gizli temelidir: sonsuz bir dizi tanımlamak için $(\Omega,\mathcal{F},P)$'nin var olduğunu kabul ederiz.
\end{ornek}

\begin{ispatiskele}
\textbf{Kolmogorov genişleme teoremi ispatı taslağı}
\begin{enumerate}
  \item Sonlu boyutlu dağılımların tutarlılık koşulunu yaz (marjinalizasyon kapalı).
  \item Silindir kümeleri üzerinde $P$'yi tüm tutarlı dağılımları kullanarak tanımla.
  \item Carathéodory genişleme teoremini uygulayarak $\sigma$-cebirde $P$'yi genişlet.
  \item $\sigma$-toplamsallığını (sayılabilir toplanabilirlik) topolojik argümanla doğrula.
\end{enumerate}
\end{ispatiskele}

\begin{mikroozet}
Koşullu bağımsızlık: $\Prob(A\cap B|C) = \Prob(A|C)\Prob(B|C)$.
Koşullu bağımsızlık mutlak bağımsızlığı ne gerektirir ne de garantiler.
Bayesçi ağlar: DAG + koşullu olasılık tabloları $\Rightarrow$ ortak dağılım.
Entropi: belirsizlik ölçüsü; KL iraksaması $\geq0$; MLE = minimum KL.
Kolmogorov genişlemesi: sonsuz i.i.d.\ dizinin varlığı garanti edilir.
\end{mikroozet}

% =====================================================================
%  BÖLÜM SONU ÖZETİ
% =====================================================================
\section*{Bölüm Özeti}
\addcontentsline{toc}{section}{Bölüm Özeti}

\begin{bolumozeti}[Bölüm 2 --- Temel Sonuçlar]
\begin{itemize}
  \item \textbf{Olasılık uzayı:} $(\Omega,\mathcal{F},\Prob)$; (K1) $\Prob \geq 0$;
        (K2) $\Prob(\Omega)=1$; (K3) $\sigma$-toplamsallık.
  \item \textbf{Poincaré:} $\Prob(\bigcup_k A_k) = \sum|I|(-1)^{|I|-1}\Prob(\bigcap_{i\in I}A_i)$.
  \item \textbf{Koşullu olasılık:} $\Prob(A|B) = \Prob(A \cap B)/\Prob(B)$;
        bu bir olasılık ölçüsü.
  \item \textbf{Toplam olasılık:} $\Prob(A) = \sum_n \Prob(A|B_n)\Prob(B_n)$.
  \item \textbf{Bayes:} sonsal $= $ olabilirlik $\times$ önsel / marjinal.
  \item \textbf{Bağımsızlık:} $\Prob(A \cap B) = \Prob(A)\Prob(B)$.
        Karşılıklı bağımsızlık $\neq$ ikili bağımsızlık (Bernstein).
        Tümleyen bağımsızlığı korur.
  \item \textbf{Klasik model:} $\Prob(A) = |A|/|\Omega|$;
        doğum günü $n=23$'te çarpıcı.
\end{itemize}

\medskip\noindent\textbf{Notasyon özeti:}
\begin{itemize}
  \item $\Prob(A|B)$: koşullu olasılık;\quad $\Prob(B_k|A)$: sonsal
  \item $\Prob(B_k)$: önsel;\quad $\Prob(A|B_k)$: olabilirlik
  \item $\binom{n}{k}$: kombinasyon;\quad $n!$: permütasyon
\end{itemize}
\end{bolumozeti}

\begin{ozyansima}
\begin{itemize}
  \item Kolmogorov aksiyomlarından tüm temel özellikleri türetebilir misiniz
        (tümleme, monotonluk, Poincaré)?
  \item Koşullu olasılığın neden bir olasılık ölçüsü olduğunu üç aksiyomu
        kontrol ederek açıklayabilir misiniz?
  \item Tıbbi test örneğini farklı bir hastalık sıklığı ($1/100$) ile yeniden
        hesaplayın; sonuç nasıl değişti?
  \item Bernstein'ın karşı örneğinde neden karşılıklı bağımsızlık bozuluyor?
\end{itemize}
\end{ozyansima}

\begin{kopru}
\textbf{Sonraki bölüm (Bölüm~3: Rasgele Değişkenler ve Dağılımlar).}
Olasılık uzayı altyapısı artık hazır. Rasgele değişken, $(\Omega,\mathcal{F})$'den
$(\mathbb{R},\mathcal{B}(\mathbb{R}))$'ye ölçülebilir fonksiyon olarak
tanımlanacak; dağılım fonksiyonu bu fonksiyonun $\Prob$ altındaki görüntüsü
olacak.

\medskip

\textbf{İleri okuma.}
Kolmogorov~\cite{kolmogorov1933} (orijinal kaynak, 77 sayfa);
Feller~\cite{feller1968vol1} Bölüm~I--V;
Gnedenko~\cite{gnedenko1962} Bölüm~1--3.
\end{kopru}


% =====================================================================
%  ALIŞTIRMALAR
% =====================================================================

\cozumlualistirmalar

\begin{alistirma}[Al.2.1]{A}{Poincaré formülü --- üç olay}
$\Prob(A) = 0.5$, $\Prob(B) = 0.4$, $\Prob(C) = 0.3$,
$\Prob(A \cap B) = 0.2$, $\Prob(A \cap C) = 0.15$,
$\Prob(B \cap C) = 0.1$, $\Prob(A \cap B \cap C) = 0.05$.
$\Prob(A \cup B \cup C)$'yi hesaplayın.
\end{alistirma}
\alcozum{%
Poincaré (Teorem~\ref{thm:poincare}):
\begin{align*}
  \Prob(A \cup B \cup C)
  &= 0.5 + 0.4 + 0.3 - 0.2 - 0.15 - 0.1 + 0.05 = 0.80.
\end{align*}
$\blacksquare$%
}

\begin{alistirma}[Al.2.2]{B}{Bayes teoremi --- fabrika}
Örnek~\ref{ex:fabrika}'ya devam: Seçilen ürün hatalı çıktı.
Hangi fabrikadan geldiğinin sonsal olasılıklarını hesaplayın.
\end{alistirma}
\alcozum{%
$\Prob(A) = 0.032$ (Örnek~\ref{ex:fabrika}'dan).
\[
  \Prob(B_1 \mid A) = \frac{0.02 \times 0.5}{0.032} = \frac{0.010}{0.032} \approx 0.3125.
\]
\[
  \Prob(B_2 \mid A) = \frac{0.04 \times 0.3}{0.032} = \frac{0.012}{0.032} = 0.375.
\]
\[
  \Prob(B_3 \mid A) = \frac{0.05 \times 0.2}{0.032} = \frac{0.010}{0.032} \approx 0.3125.
\]
Toplam: $0.3125 + 0.375 + 0.3125 = 1$. \checkmark
Hatalı ürün en olası fabrika 2'den geliyor. $\blacksquare$%
}

\begin{alistirma}[Al.2.3]{A}{Bağımsızlık doğrulama}
Adil zar: $\Omega = \{1,2,3,4,5,6\}$, $\Prob(\{k\}) = 1/6$.
$A = \{1,2,3,4\}$, $B = \{3,4,5,6\}$, $C = \{1,3,5\}$.
(a) $A$ ve $B$ bağımsız mı?
(b) $A$ ve $C$ bağımsız mı?
(c) $B$ ve $C$ bağımsız mı?
(d) $A, B, C$ karşılıklı bağımsız mı?
\end{alistirma}
\alcozum{%
$\Prob(A) = 4/6 = 2/3$, $\Prob(B) = 4/6 = 2/3$, $\Prob(C) = 3/6 = 1/2$.

(a) $A \cap B = \{3,4\}$; $\Prob(A \cap B) = 2/6 = 1/3$.
$\Prob(A)\Prob(B) = 4/9 \neq 1/3$. Bağımsız \emph{değil}.

(b) $A \cap C = \{1,3\}$; $\Prob(A \cap C) = 2/6 = 1/3$.
$\Prob(A)\Prob(C) = 2/3 \cdot 1/2 = 1/3$. Bağımsız. \checkmark

(c) $B \cap C = \{3,5\}$; $\Prob(B \cap C) = 1/3$.
$\Prob(B)\Prob(C) = 1/3$. Bağımsız. \checkmark

(d) $A \cap B \cap C = \{3\}$; $\Prob = 1/6$.
$\Prob(A)\Prob(B)\Prob(C) = 2/3 \cdot 2/3 \cdot 1/2 = 2/9 \neq 1/6$.
Karşılıklı bağımsız \emph{değil}. $\blacksquare$%
}

\begin{alistirma}[Al.2.4]{B}{Doğum günü --- genelleştirilmiş}
$n = 50$ kişilik bir grupta aynı doğum gününe sahip iki kişinin olmama
olasılığını hesaplayın (logaritma kullanın). $\Prob(A)$ nedir?
\end{alistirma}
\alcozum{%
$\Prob(A^c) = \prod_{k=0}^{49}\frac{365-k}{365}$.
$\ln \Prob(A^c) = \sum_{k=0}^{49} \ln(1 - k/365)$.
$\ln(1-k/365) \approx -k/365$ için kaba tahmin:
$\sum_{k=0}^{49} k/365 = \binom{50}{2}/365 = 1225/365 \approx 3.356$.
$\Prob(A^c) \approx e^{-3.356} \approx 0.035$.
$\Prob(A) \approx 0.965$.

Daha kesin: $\Prob(A^c) \approx 0.0296$; $\Prob(A) \approx 0.970$. $\blacksquare$%
}

\begin{alistirma}[Al.2.4b]{B}{Koşullu Olasılık Paradoksu: Monty Hall}
Monty Hall problemi: 3 kapı, birinde araba diğer ikisinde keçi. Yarışmacı kapı 1'i seçer. Sunucu, keçili bir kapıyı açar (kapı 3). Değiştirmeli mi?
(a) $A$: kapı 1'de araba; $B$: sunucu kapı 3'ü açar; Bayes teoremi ile $P(A|B)$'yi hesaplayın.
(b) $P(\text{araba}|\text{değiştirir})$'yi bulun.
(c) Sezgisiz insanların neden yanıldığını olasılık teorisi açısından açıklayın.
\end{alistirma}
\alcozum{%
(a) Toplam olasılık: $P(B) = P(B|A)P(A)+P(B|A^c)P(A^c)$.
$P(A)=1/3$; $P(A^c)=2/3$.
$P(B|A)=1/2$ (sunucu 2 veya 3'ü eşit olasılıkla açar; biz kapı 3'ü koşullandırdık).
$P(B|A^c)=1$ eğer araba kapı 2'deyse sunucu sadece kapı 3'ü açabilir; $P(B|A^c=\text{kapı 2})=1$, $P(B|A^c=\text{kapı 3})=0$.

Daha dikkatli: $A_k$: araba kapı $k$'da.
$P(A_1)=P(A_2)=P(A_3)=1/3$.
Sunucu kapı 3'ü açar: $P(B|A_1)=1/2$, $P(B|A_2)=1$, $P(B|A_3)=0$.
$P(B)=\frac{1}{3}\cdot\frac{1}{2}+\frac{1}{3}\cdot1+\frac{1}{3}\cdot0=\frac{1}{2}$.
$P(A_1|B)=\frac{(1/2)(1/3)}{1/2}=\frac{1}{3}$; $P(A_2|B)=\frac{1\cdot(1/3)}{1/2}=\frac{2}{3}$.

(b) Değiştirir $\to$ kapı 2 seçer: $P(\text{araba})=P(A_2|B)=2/3$. Değiştirmez: $1/3$. Değiştirmek iki kat daha iyi!

(c) Yanılgı: Sunucunun ek bilgi verdiği göz ardı ediliyor. Sunucunun kapı açması rastgele değil — keçi içeren kapıyı bilinçli açar. Bu bilgi kapı 2'nin olasılığını günceller. İnsanlar genellikle sunucunun hamlesi bağımsız diye düşünür ($P=1/3$ değişmez sanır); Bayes teoremi bu bağımlılığı doğru biçimde yakalar.}

\begin{alistirma}[Al.2.4c]{C}{Olasılık Ölçüsünün Süreklilik Özellikleri}
$(\Omega,\mathcal{F},\Prob)$ olasılık uzayı.
(a) $A_n\uparrow A$ (artan limit) için $\Prob(A_n)\to\Prob(A)$ olduğunu kanıtlayın.
(b) $B_n\downarrow B$ (azalan limit) için $\Prob(B_n)\to\Prob(B)$ olduğunu kanıtlayın.
(c) $\Prob(\limsup_n A_n)$'ı yorumlayın ve Borel-Cantelli ile bağlantısını kurun.
\end{alistirma}
\alcozum{%
(a) \textbf{Arttan süreklilik}: $C_1=A_1$, $C_n=A_n\setminus A_{n-1}$ (ayrık) olsun. $A=\bigcup_n C_n$ (ayrık birleşim).
$\Prob(A)=\sum_{k=1}^\infty\Prob(C_k)=\lim_{n\to\infty}\sum_{k=1}^n\Prob(C_k)=\lim_{n\to\infty}\Prob(A_n)$.
(Sayılabilir toplamsellik ve dizi yakınsaması birlikte.)

(b) \textbf{Azalan süreklilik}: $D_n=B_1\setminus B_n\uparrow B_1\setminus B$.
(a)'dan $\Prob(D_n)\to\Prob(B_1\setminus B)$.
$\Prob(D_n)=\Prob(B_1)-\Prob(B_n)$ (tek yönlü kapsama); $\Prob(B_n)\to\Prob(B)$. $\checkmark$

(c) $\limsup_n A_n=\{A_n\text{ sonsuz sık gerçekleşir}\}=\bigcap_{n=1}^\infty\bigcup_{k=n}^\infty A_k$.
Borel-Cantelli 1: $\sum\Prob(A_n)<\infty\Rightarrow\Prob(\limsup A_n)=0$ (sonsuz sık gerçekleşmez).
Borel-Cantelli 2: $A_n$ bağımsız, $\sum\Prob(A_n)=\infty\Rightarrow\Prob(\limsup A_n)=1$ (sonsuz sık gerçekleşir).
Bu iki yönlü sonuç, neredeyse kesin yakınsama kriterleri için temel araçtır.}

\alistirmalar

\begin{alistirma}[Al.2.5]{A}{Olasılık aksiyomlarından türetme}
$(\Omega,\mathcal{F},\Prob)$ olasılık uzayı, $A, B \in \mathcal{F}$.
Yalnızca Kolmogorov aksiyomlarını (K1)--(K3) kullanarak gösterin:
\begin{enumerate}[label=(\alph*)]
  \item $\Prob(A) \leq 1$.
  \item $\Prob(A^c) = 1 - \Prob(A)$.
  \item $A \subseteq B \Rightarrow \Prob(A) \leq \Prob(B)$.
\end{enumerate}
\end{alistirma}
\alcozum{%
(a) $\Omega = A \cup A^c$ ayrık; (K3) + (K1): $1 = \Prob(A)+\Prob(A^c) \geq \Prob(A)$.
(b) Aynı ayrışım: $\Prob(A^c) = 1 - \Prob(A)$.
(c) $B = A \cup (B \setminus A)$ ayrık; $\Prob(B) = \Prob(A)+\Prob(B \setminus A) \geq \Prob(A)$. $\blacksquare$%
}

\begin{alistirma}[Al.2.6]{A}{Koşullu olasılık zinciri}
Bir kutuda 3 kırmızı, 4 mavi top var. Geri koymadan iki top çekiliyor.
(a) İlk kırmızı, ikinci mavi çekme olasılığı.
(b) Her ikisi de mavi çekme olasılığı.
(c) En az bir kırmızı çekme olasılığı.
\end{alistirma}
\alcozum{%
(a) $\Prob(K_1 \cap M_2) = \Prob(K_1)\Prob(M_2|K_1) = \frac{3}{7} \cdot \frac{4}{6} = \frac{12}{42} = \frac{2}{7}$.

(b) $\Prob(M_1 \cap M_2) = \frac{4}{7} \cdot \frac{3}{6} = \frac{12}{42} = \frac{2}{7}$.

(c) $\Prob(\text{en az bir kırmızı}) = 1 - \Prob(M_1 \cap M_2) = 1 - 2/7 = 5/7$. $\blacksquare$%
}

\begin{alistirma}[Al.2.7]{A}{Boole eşitsizliği}
$n$ olay için Boole eşitsizliğini doğrulayın:
$\Prob\!\bigl(\bigcap_{k=1}^n A_k\bigr) \geq 1 - \sum_{k=1}^n \Prob(A_k^c)$.
Bu eşitsizliğin adı Bonferroni eşitsizliğidir.
\end{alistirma}
\alcozum{%
$\Prob\!\bigl(\bigcap_k A_k\bigr) = 1 - \Prob\!\bigl(\bigcup_k A_k^c\bigr)$
(De~Morgan ve tümleyen kuralı).
Alt-toplamsallık: $\Prob(\bigcup_k A_k^c) \leq \sum_k \Prob(A_k^c)$.
Dolayısıyla $\Prob(\bigcap_k A_k) \geq 1 - \sum_k \Prob(A_k^c)$. $\blacksquare$%
}

\begin{alistirma}[Al.2.8]{B}{Koşullu bağımsızlık}
$A$ ve $B$, $C$'ye koşullu bağımsız adını alır; eğer
$\Prob(A \cap B \mid C) = \Prob(A \mid C)\Prob(B \mid C)$.
(a) $A$ ve $B$ koşulsuz bağımsızsa, $C$'ye koşullu bağımsız olması gerekir mi?
    Karşı örnek verin.
(b) $A$ ve $B$, $C$'ye koşullu bağımsızsa, koşulsuz bağımsız olması gerekir mi?
    Karşı örnek verin.
\end{alistirma}
\alcozum{%
(a) Hayır. İki adil zar: $A =$ \{1. zar çift\}, $B =$ \{2. zar çift\}, $C = \{$toplam 7$\}$.
$A,B$ bağımsız. Ama $\Prob(A|C) = 3/6 = 1/2$ ve $\Prob(B|C) = 3/6 = 1/2$;
$\Prob(A \cap B \mid C) = \Prob(\{(2,5),(4,3),(6,1)\} \mid C) = 0$
(çift+çift=çift, toplam 7 olamaz). $\neq 1/4$. Koşullu bağımsız değil.

(b) Hayır. Örnek: $\Omega=\{1,2,3,4\}$ düzgün; $A=\{1,2\}$, $B=\{2,3\}$, $C=\{1,3\}$.
$\Prob(A)=\Prob(B)=1/2$, $\Prob(A\cap B)=1/4=\Prob(A)\Prob(B)$: bağımsız.
$\Prob(A\cap B\mid C)=\Prob(\emptyset\mid C)=0$,
$\Prob(A|C)\Prob(B|C)=\frac{1}{2}\cdot\frac{1}{2}=\frac{1}{4}\neq 0$:
koşullu bağımsız değil. $\blacksquare$%
}

\begin{alistirma}[Al.2.9]{B}{Toplam olasılık --- iki aşamalı deney}
Bir torba 5 beyaz, 3 siyah top içeriyor. Geri koymadan iki top çekiliyor.
$A_k =$ \{$k$-inci top beyaz\}.
(a) $\Prob(A_1)$, $\Prob(A_2)$, $\Prob(A_1 \cap A_2)$'yi hesaplayın.
(b) $\Prob(A_2)$'yi Toplam Olasılık Teoremi ile doğrulayın.
(c) $\Prob(A_1 \mid A_2)$ ve $\Prob(A_2 \mid A_1)$'i hesaplayın.
\end{alistirma}
\alcozum{%
(a) $\Prob(A_1) = 5/8$.
$\Prob(A_2) = \Prob(A_2|A_1)\Prob(A_1)+\Prob(A_2|A_1^c)\Prob(A_1^c)$
$= \frac{4}{7}\cdot\frac{5}{8}+\frac{5}{7}\cdot\frac{3}{8} = \frac{20}{56}+\frac{15}{56}=\frac{35}{56}=\frac{5}{8}$.
$\Prob(A_1 \cap A_2) = \frac{5}{8}\cdot\frac{4}{7} = \frac{20}{56} = \frac{5}{14}$.

(b) Yukarıdaki hesap Toplam Olasılık Teoremi'nin uygulaması.

(c) $\Prob(A_1|A_2) = \Prob(A_1\cap A_2)/\Prob(A_2) = (5/14)/(5/8) = 8/14 = 4/7$.
$\Prob(A_2|A_1) = \Prob(A_1\cap A_2)/\Prob(A_1) = (5/14)/(5/8) = 4/7$. $\blacksquare$%
}

\begin{alistirma}[Al.2.10]{B}{Hat çekme genelleştirilmiş}
$n$ kişi, $m$ şanslı hat ($m < n$). Hiç kimse önceki seçimi görmeden
kendi hatını çekiyor (tüm sıralamalar eşit olası).
$j$-inci kişinin şanslı hat çekme olasılığının $j$'den bağımsız olduğunu
tümevarımla ispatlayın.
\end{alistirma}
\alcozum{%
$\Prob(A_j) = m/n$ her $j$ için.

$j=1$: $\Prob(A_1) = m/n$ açıkça.

$j$'de doğru varsayalım; $j+1$:
$\Prob(A_{j+1}) = \Prob(A_{j+1}|A_j)\Prob(A_j) + \Prob(A_{j+1}|A_j^c)\Prob(A_j^c)$
$= \frac{m-1}{n-1}\cdot\frac{m}{n} + \frac{m}{n-1}\cdot\frac{n-m}{n}$
$= \frac{m(m-1) + m(n-m)}{n(n-1)} = \frac{m(n-1)}{n(n-1)} = \frac{m}{n}$.
$\blacksquare$%
}

\begin{alistirma}[Al.2.11]{C}{$\pi$-$\lambda$ teoremi ve bağımsızlık}
$\mathcal{P}_1 = \{A_1 \times A_2 : A_i \in \mathcal{F}_i\}$ ikili ayrık
dikdörtgenler ailesi; $\mathcal{P}_1$ bir $\pi$-sistemdir.
$\Prob_1$ ve $\Prob_2$ aynı $\pi$-sistem üzerinde aynı değerleri alıyorsa,
$\sigma(\mathcal{P}_1)$ üzerinde eşit olduklarını Dynkin'in $\pi$-$\lambda$
teoremini kullanarak ispatlayın. Bu teoremin ürün ölçüsünün tekliği ile
bağlantısını açıklayın.
\end{alistirma}
\alcozum{%
\begin{ipucu}{1}
Dynkin'in $\pi$-$\lambda$ teoremi: $\mathcal{P}$ bir $\pi$-sistem (sonlu kesişim
altında kapalı) ve $\lambda(\mathcal{P}) = \mathcal{P}$ içeren en küçük
Dynkin sistemi ise $\lambda(\mathcal{P}) = \sigma(\mathcal{P})$.
\end{ipucu}

$D = \{B \in \sigma(\mathcal{P}_1) : \Prob_1(B) = \Prob_2(B)\}$ tanımlayın.
$D$ bir Dynkin sistemidir: (i) $\Omega \in D$; (ii) $A \subseteq B$, $A,B \in D
\Rightarrow B \setminus A \in D$; (iii) $A_n \nearrow A$, $A_n \in D \Rightarrow A \in D$.

$\mathcal{P}_1 \subseteq D$ (verilen koşul).
Dynkin teoreminden $\sigma(\mathcal{P}_1) \subseteq D$, yani $\Prob_1 = \Prob_2$ üzerinde.

Ürün ölçüsü: $\mu_1 \otimes \mu_2$ dikdörtgenlerde
$\mu_1(A_1)\mu_2(A_2)$'yi üretir; bu değerler başka bir ölçüyü belirliyorsa
her ikisi $\sigma(\mathcal{P}_1)$ üzerinde eşittir; yani ürün ölçüsü
$\sigma$-sonlu durumlarda tekdir. $\blacksquare$%
}

\begin{alistirma}[Al.2.12]{C}{Sonsuz bağımsız denemeler}
$(\Omega, \mathcal{F}, \Prob)$ üzerinde $A_1, A_2, \ldots$ bağımsız,
$\Prob(A_n) = p \in (0,1)$ her $n$ için.
(a) $\Prob(A_n \text{ s.s. çok})$ — yani $\Prob(\limsup_n A_n)$ — nedir?
(b) $\Prob(A_n \text{ eninde sonunda hep gerçekleşir})$ nedir?
(c) Cevaplarınızı $p$ değerine göre yorumlayın.
\end{alistirma}
\alcozum{%
(a) $\sum_n \Prob(A_n) = \sum_n p = \infty$ (sonsuz toplam).
Borel-Cantelli II (Bölüm~7'de kanıtlanacak): $A_n$ bağımsız ve $\sum \Prob(A_n) = \infty$
ise $\Prob(\limsup_n A_n) = 1$.

(b) $\liminf_n A_n$: eninde sonunda hep gerçekleşir.
$B_n = \bigcap_{k \geq n} A_k^c$; $\Prob(B_n) = \prod_{k \geq n}(1-p) = 0$
(sonsuz çarpım, $p>0$). $\Prob(\liminf_n A_n) = 1 - \Prob(\limsup_n A_n^c)$.
$\sum_n \Prob(A_n^c) = \infty$; $A_n^c$ da bağımsız; Borel-Cantelli II:
$\Prob(\limsup_n A_n^c) = 1$.
$\Prob(\liminf_n A_n) = 0$.

(c) Her $p \in (0,1)$ için: sonsuz bağımsız denemede her olay sonsuz sık
gerçekleşir (a.s.), ama hiçbir zaman sonunda \emph{sürekli} gerçekleşmez. $\blacksquare$%
}

\begin{alistirma}[Al.2.13]{A}{Koşullu Bağımsızlık}
Sağlık veri tabanında: $D$ = hastada kanser var, $S$ = semptom pozitif, $T$ = test pozitif.
Verilen: $\Prob(D) = 0.01$, $\Prob(S \mid D) = 0.8$, $\Prob(S \mid D^c) = 0.05$,
$\Prob(T \mid S) = 0.9$, $\Prob(T \mid S^c) = 0.1$.
Bayesçi ağ: $D \to S \to T$.
(a) $\Prob(S)$'yi hesaplayın. (b) $\Prob(T)$'yi hesaplayın. (c) $\Prob(D \mid T)$'yi hesaplayın.
\end{alistirma}
\alcozum{%
(a) $\Prob(S) = 0.8\cdot0.01 + 0.05\cdot0.99 = 0.008 + 0.0495 = 0.0575$.\\
(b) $\Prob(T) = 0.9\cdot\Prob(S) + 0.1\cdot\Prob(S^c) = 0.9(0.0575) + 0.1(0.9425) = 0.05175 + 0.09425 = 0.146$.\\
(c) $\Prob(D \mid T) = \Prob(T \mid D)\Prob(D)/\Prob(T)$.
$\Prob(T \mid D) = 0.9\cdot0.8 + 0.1\cdot0.2 = 0.74$.
$\Prob(D \mid T) = 0.74\cdot0.01/0.146 \approx 0.051$. Yüksek test duyarlılığına rağmen düşük: nadir hastalık etkisi.}

\begin{alistirma}[Al.2.14]{B}{Koşullu Bağımsızlık vs Mutlak Bağımsızlık}
$\Omega = \{1,2,3,4\}$, eşit olası.
$A = \{1,2\}$, $B = \{1,3\}$, $C = \{1,4\}$.
(a) $A$ ve $B$ bağımsız mı? (b) $A$ ve $B$, $C$ koşulunda bağımsız mı?
(c) $A$ ve $C$ bağımsız mı? $A$ ve $C$, $B$ koşulunda?
\end{alistirma}
\alcozum{%
$\Prob(A) = \Prob(B) = \Prob(C) = 1/2$.
$\Prob(A\cap B) = \Prob(\{1\}) = 1/4 = \Prob(A)\Prob(B)$: (a) evet, $A \perp B$.\\
(b) $\Prob(C) = 1/2$; $\Prob(A\cap B\cap C) = \Prob(\{1\}) = 1/4$.
$\Prob(A\cap B \mid C) = \Prob(A\cap B\cap C)/\Prob(C) = (1/4)/(1/2) = 1/2$.
$\Prob(A\mid C)\Prob(B\mid C) = (1/2)(1/2) = 1/4 \neq 1/2$.
Koşullu bağımsız değil!\\
(c) $A \perp C$: $\Prob(A\cap C) = 1/4 = \Prob(A)\Prob(C)$. Evet.
$\Prob(A\cap C\mid B) = \Prob(A\cap B\cap C)/\Prob(B) = (1/4)/(1/2) = 1/2 \neq \Prob(A|B)\Prob(C|B) = (1/2)(1/2)$.
$A \not\perp C \mid B$.}

\begin{alistirma}[Al.2.15]{B}{Bayesçi Ağ — Ortak Dağılım}
$X_1$, $X_2 \mid X_1$, $X_3 \mid X_2$ zinciri (her ikili normal):
$X_1 \sim \Normal(0,1)$; $X_2 \mid X_1 = x_1 \sim \Normal(x_1, 1)$; $X_3 \mid X_2 = x_2 \sim \Normal(x_2, 1)$.
(a) $E[X_2]$ ve $\mathrm{Var}(X_2)$'yi hesaplayın.
(b) $X_1$ ve $X_3$'ün koşulsuz korelasyonunu bulun.
(c) $X_3 \mid X_2 = 0$ ile $X_3 \mid X_1 = 0$'ı karşılaştırın.
\end{alistirma}
\alcozum{%
(a) $E[X_2] = E[E[X_2|X_1]] = E[X_1] = 0$.
$\mathrm{Var}(X_2) = E[\mathrm{Var}(X_2|X_1)] + \mathrm{Var}(E[X_2|X_1]) = 1 + \mathrm{Var}(X_1) = 1+1 = 2$.\\
(b) Aynı şekilde: $E[X_3] = 0$, $\mathrm{Var}(X_3) = 3$.
$\mathrm{Cov}(X_1,X_3) = E[X_1 X_3]$. Kule beklentisi: $E[X_1 X_3] = E[X_1 E[X_3|X_1,X_2]]$... zincirde: $E[X_3|X_2] = X_2$; $E[X_1 X_3] = E[X_1 X_2] = E[X_1^2] = 1$.
$\mathrm{Cor}(X_1,X_3) = 1/\sqrt{1\cdot3} = 1/\sqrt{3} \approx 0.577$.\\
(c) $X_3|X_2=0 \sim \Normal(0,1)$. $X_3|X_1=0$: $X_2|X_1=0 \sim \Normal(0,1)$, sonra $X_3|X_2 \sim \Normal(X_2,1)$; marjinal $X_3|X_1=0 \sim \Normal(0,2)$.
Daha geniş: $X_1$ bilgisi $X_2$ üzerinden dolaylı.}

% =====================================================================
%  BÖLÜM 3: Rasgele Değişkenler ve Dağılımlar
%  Kaynak: Feller Vol.I-II; Shiryaev; Billingsley; Klenke
% =====================================================================
\chapter{Rasgele Değişkenler ve Dağılımlar}
\label{chp:rasgele_degiskenler}

\bolumalinti{%
  Olasılık teorisi, belirsizliği sayıya dönüştürmenin bilimidir.
  Rasgele değişken bu dönüşümün ta kendisidir.%
}{Andrei Kolmogorov}

% ---- Kazanımlar (TYMM) ----
\begin{kazanimlar}
Bu bölümün sonunda öğrenci:
\begin{itemize}[leftmargin=1.8em]
  \item Rasgele değişkeni ölçülebilir fonksiyon olarak tanımlayabilir ve
        bu tanımın neden "şans sonucu sayı üretme" sezgisini yakaladığını açıklayabilir.
  \item Kümülatif dağılım fonksiyonu (KDF) kavramını tanımlayabilir,
        karakteristik özelliklerini ispatlayabilir ve KDF'den olasılık hesabı yapabilir.
  \item Kesikli, mutlak sürekli ve karma dağılımlar arasındaki farkı
        Lebesgue ayrışımı çerçevesinde ifade edebilir.
  \item Ortak dağılım, marjinal dağılım ve koşullu dağılım kavramlarını
        tanımlayabilir; bağımsız rasgele değişkenlerin ayırt edici özelliklerini kullanabilir.
  \item Birebir dönüşüm yöntemi ve Jacobian formülünü tek ve çok boyutlu
        durumlarda uygulayabilir; evrişim (konvolüsyon) yöntemiyle toplam dağılımını bulabilir.
\end{itemize}
\end{kazanimlar}

% ---- Motivasyon ----
\begin{motivasyon}
Bölüm~\ref{chp:olasilik_uzaylari}'da bir olasılık uzayı $(\Omega, \mathcal{F}, P)$ inşa ettik.
Ancak $\Omega$ kümesi çoğu zaman soyut bir uzaydır: yazı-tura atışında
$\Omega = \{Y, T\}$, bir deneyin tüm sonuçları kümesinde ise karmaşık nesneler barındırabilir.
Pratikte ilgilendiğimiz şey genellikle \emph{sayısal} sonuçlardır:
"kaç tura geldi?", "bekleme süresi ne kadar?", "ölçülen sıcaklık nedir?"

\textbf{Rasgele değişken} (random variable), $\Omega$'yı reel sayı eksenine taşıyan
ölçülebilir bir fonksiyondur. Bu taşıma sayesinde karmaşık olasılık uzaylarını
\emph{gerçel sayılar üzerindeki dağılımlarla} inceleyebiliriz; asıl ağırlık
$(\Omega, \mathcal{F}, P)$'den $(\mathbb{R}, \mathcal{B}(\mathbb{R}), P_X)$'e geçer.

\medskip
\textbf{Köprü.} Bölüm~\ref{chp:olcum_teorisi}'de ölçülebilir fonksiyon ve
Lebesgue integrali; Bölüm~\ref{chp:olasilik_uzaylari}'da Kolmogorov aksiyomları
ve bağımsızlık kavramı bu bölümün temelini oluşturur.
Bölüm~\ref{chp:beklenti}'de bu dağılımlar üzerinden beklenti ve momentler hesaplanacak;
Bölüm~\ref{chp:dagilim_aileleri}'de somut dağılım aileleri incelenecektir.
\end{motivasyon}

% =====================================================================
\section{Rasgele Değişken Kavramı}
\label{sec:rd_kavram}
% =====================================================================

\begin{kesif}
Yazı-tura atışında $\Omega = \{Y, T\}$ olsun.
"Tura sayısı" bir fonksiyon tanımlar: $X(Y) = 0$, $X(T) = 1$.
Bu fonksiyonun hangi özelliği onu "rasgele değişken" yapar?
Eğer $A = \{1\} \subset \mathbb{R}$ ise $\{X \in A\} = \{T\}$ olayı
olasılık uzayında tanımlı olmalı — yani $\mathcal{F}$'ye ait olmalıdır.
Peki ya $A = (0.5, 1.5)$ veya $A = \emptyset$ için?
Tüm Borel kümelerine bu özelliği istiyorsak ne elde ederiz?
\end{kesif}

\begin{tanimgerekce}
Neden "$f^{-1}(B) \in \mathcal{F}$ her $B \in \mathcal{B}(\mathbb{R})$ için" koşulu?
Çünkü $P(X \in B)$ ifadesi ancak $\{X \in B\} = X^{-1}(B)$ bir olay olduğunda,
yani $\mathcal{F}$'de olduğunda, anlam taşır.
Ölçülebilirlik tam da bu hesabı mümkün kılan koşuldur.
\end{tanimgerekce}

\begin{tanim}[Rasgele Değişken]
\label{def:rasgele_degisken}
$(\Omega, \mathcal{F}, P)$ bir olasılık uzayı olsun.
$X : \Omega \to \mathbb{R}$ fonksiyonu, her $B \in \mathcal{B}(\mathbb{R})$ için
\[
  X^{-1}(B) = \{\omega \in \Omega : X(\omega) \in B\} \in \mathcal{F}
\]
koşulunu sağlıyorsa $X$'e \textbf{rasgele değişken} (random variable) denir.
Başka bir deyişle $X$, $(\Omega, \mathcal{F})$'den $(\mathbb{R}, \mathcal{B}(\mathbb{R}))$'ye
\textbf{ölçülebilir bir fonksiyondur}.
\end{tanim}

\begin{onerme}[Eşdeğer Karakterizasyon]
\label{prop:rd_char}
$X : \Omega \to \mathbb{R}$ için aşağıdaki koşullar birbirine denktir:
\begin{enumerate}[(i)]
  \item $X$ rasgele değişkendir.
  \item Her $c \in \mathbb{R}$ için $\{X \leq c\} \in \mathcal{F}$.
  \item Her $c \in \mathbb{R}$ için $\{X < c\} \in \mathcal{F}$.
  \item Her $c \in \mathbb{R}$ için $\{X > c\} \in \mathcal{F}$.
  \item Her $c \in \mathbb{R}$ için $\{X \geq c\} \in \mathcal{F}$.
\end{enumerate}
\end{onerme}

\begin{proof}
Bölüm~\ref{chp:olcum_teorisi}'deki Teorem~\ref{thm:olculebichar}'in doğrudan
bir uygulamasıdır; $\pi$-sistemi argümanı $\{(-\infty, c] : c \in \mathbb{R}\}$'nin
$\mathcal{B}(\mathbb{R})$'yi ürettiğini kullanır.
$(i) \Rightarrow (ii)$: $(-\infty, c] \in \mathcal{B}(\mathbb{R})$ olduğundan $X^{-1}((-\infty, c]) \in \mathcal{F}$.
$(ii) \Rightarrow (iii)$: $\{X < c\} = \bigcup_{n=1}^\infty \{X \leq c - 1/n\}$,
sayılabilir birleşim $\mathcal{F}$'de kalır.
$(iii) \Leftrightarrow (iv)$: $\{X < c\} = \{X \geq c\}^c$.
$(ii) \Leftrightarrow (v)$: $\{X \geq c\} = \{X < c\}^c$.
$(ii) \Rightarrow (i)$: $\pi$-sistemi $\{(-\infty, c]\}$ üzerinde
$X^{-1}$ $\mathcal{F}$'ye düşüyor; $\mathcal{B}(\mathbb{R})$'nin üretilmesi bu yeterli.
\end{proof}

\begin{kritiksorgu}
"$X$ sürekli olması rasgele değişken olması için yeterli midir?"
Hayır: süreklilik analitik bir kavramdır, ölçülebilirliği garanti etmez —
ancak Borel-ölçülebilir bir $\sigma$-cebir düzeni altında sürekli fonksiyonlar
her zaman ölçülebilirdir (dolayısıyla rasgele değişkendir).
Vitali kümesinin gösterge fonksiyonu, $(\mathbb{R}, \mathcal{B}(\mathbb{R}))$'de
ölçülebilir olmayan bir fonksiyon örneğidir (ama bu durum uygulamada sorun yaratmaz).
\end{kritiksorgu}

\begin{onerme}[Cebirsel Kararlılık]
\label{prop:rd_cebir}
$X, Y$ rasgele değişkenler ve $c \in \mathbb{R}$ sabit olsun. O zaman
$X + Y$, $cX$, $XY$, $\max(X, Y)$, $\min(X, Y)$ ve $|X|$ de rasgele değişkenlerdir.
Ayrıca $Y \neq 0$ n.k.\ ise $X/Y$ rasgele değişkendir.
$\{X_n\}$ rasgele değişken dizisi ise $\sup_n X_n$, $\inf_n X_n$,
$\limsup_{n} X_n$, $\liminf_n X_n$ ölçülebilir fonksiyonlardır
(sonlu değerleri yoksa $\overline{\mathbb{R}}$'ye değer alırlar).
\end{onerme}

\begin{proof}
$\{X + Y < c\} = \bigcup_{q \in \mathbb{Q}} (\{X < q\} \cap \{Y < c-q\}) \in \mathcal{F}$
(rasyonel sayılar üzerinden sayılabilir birleşim).
$\{XY > c\}$ için $XY = \tfrac{1}{4}[(X+Y)^2 - (X-Y)^2]$ açılımı kullanılır.
$\sup_n$: $\{\sup_n X_n > c\} = \bigcup_n \{X_n > c\} \in \mathcal{F}$.
Diğerleri benzer.
\end{proof}

\begin{tarihselnot}
"Rasgele değişken" terimi bir değişken değil, fonksiyondur — tarihsel adlandırma yanıltıcıdır.
Kolmogorov 1933'te olasılığı ölçüm teorisi üzerine oturtmadan önce,
rasgele değişkenler daha çok sezgisel olarak tanımlanıyordu.
Modern tanım, Lebesgue'nin ölçüm teorisi ile Borel'in $\sigma$-cebirleri
kavramlarının olasılığa uygulanmasıyla ortaya çıktı.
\end{tarihselnot}

\begin{hocanotu}
Bundan sonra "$X$ rasgele değişken" derken her zaman
$(\Omega, \mathcal{F}, P)$ olasılık uzayından $(\mathbb{R}, \mathcal{B}(\mathbb{R}))$'ye
ölçülebilir fonksiyonu kastediyoruz. Vektör değerli rasgele değişkenler
(Bölüm~\ref{sec:rd_ortak}) aynı mantıkla $\mathbb{R}^n$'e değer alan ölçülebilir
fonksiyonlar olarak tanımlanır.
\end{hocanotu}

\begin{kendinleifade}
"Rasgele değişken" neden bir "değişken" değil, bir "fonksiyondur"?
Hangi uzaydan hangi uzaya gider? Ölçülebilirlik koşulu ne anlama gelir?
Bir arkadaşınıza matematiksiz bir dille açıklayın.
\ogrencialani[2.5cm]
\end{kendinleifade}

\begin{mikroozet}
Rasgele değişken $X : (\Omega, \mathcal{F}) \to (\mathbb{R}, \mathcal{B}(\mathbb{R}))$
ölçülebilir bir fonksiyondur; bu $\{X \leq c\} \in \mathcal{F}$ koşuluyla
eşdeğerdir. Rasgele değişkenler cebirsel işlemler ve limit altında kapalıdır.
\end{mikroozet}

% =====================================================================
\section{Dağılım Fonksiyonu}
\label{sec:dagilim_fonksiyonu}
% =====================================================================

\begin{kesif}
$X$ bir rasgele değişken. $P(X \leq 1.5)$, $P(X \leq 3.7)$ gibi değerleri
biliyorsak $X$ hakkında ne kadar bilgi sahibi oluruz?
Eğer tüm $c \in \mathbb{R}$ için $P(X \leq c)$ değerlerini biliyorsak,
bu yeterli midir — yani bu değerler $X$'in dağılımını tamamen belirler mi?
\end{kesif}

\begin{tanim}[Kümülatif Dağılım Fonksiyonu]
\label{def:kdf}
$X$ bir rasgele değişken olsun. \textbf{Kümülatif dağılım fonksiyonu} (KDF)
(cumulative distribution function, CDF)
\[
  F_X(x) = P(X \leq x), \qquad x \in \mathbb{R}
\]
ile tanımlanır.
\end{tanim}

\begin{teorem}[KDF'nin Karakteristik Özellikleri]
\label{thm:kdf_ozellik}
Her KDF aşağıdaki özellikleri sağlar:
\begin{enumerate}[(i)]
  \item \textbf{Monoton artmayan değil:} $x \leq y \Rightarrow F(x) \leq F(y)$.
  \item \textbf{Sağdan sürekli:} $\lim_{y \downarrow x} F(y) = F(x)$.
  \item \textbf{Sınır davranışı:} $\lim_{x \to -\infty} F(x) = 0$ ve
        $\lim_{x \to +\infty} F(x) = 1$.
\end{enumerate}
Tersine, bu üç özelliği sağlayan her $F : \mathbb{R} \to [0,1]$ fonksiyonu
bir rasgele değişkenin KDF'sidir.
\end{teorem}

\begin{proof}
\textit{(i):} $x \leq y$ ise $\{X \leq x\} \subseteq \{X \leq y\}$, dolayısıyla
$F(x) = P(X \leq x) \leq P(X \leq y) = F(y)$.

\textit{(ii):} $y_n \downarrow x$ olsun. O zaman $\{X \leq y_n\} \downarrow \{X \leq x\}$
(azalan olaylar dizisi, $\bigcap_n \{X \leq y_n\} = \{X \leq x\}$).
Ölçünün yukarıdan sürekliliği (Bölüm~\ref{chp:olcum_teorisi}, Teorem~\ref{thm:olcuozellik}):
\[
  \lim_{y \downarrow x} F(y) = \lim_{n \to \infty} P(X \leq y_n) = P\!\left(\bigcap_n \{X \leq y_n\}\right) = P(X \leq x) = F(x).
\]

\textit{(iii):} $x_n \to -\infty$ ise $\{X \leq x_n\} \downarrow \emptyset$, dolayısıyla
$F(x_n) \to P(\emptyset) = 0$. Benzer şekilde $x_n \to +\infty$ ise
$\{X \leq x_n\} \uparrow \Omega$, dolayısıyla $F(x_n) \to P(\Omega) = 1$.

\textit{Tersine:} $(i)$--$(iii)$'ü sağlayan $F$ verildiğinde,
$\mathbb{R}$ üzerindeki Lebesgue-Stieltjes ölçüsü inşası yoluyla
(Carathéodory genişleme teoremi) $P((a,b]) = F(b) - F(a)$ sağlayan
bir olasılık ölçüsü $P$ inşa edilebilir; bu ölçünün üreteceği KDF tam olarak $F$'dir.
\end{proof}

\begin{onerme}[KDF'den Olasılık Hesabı]
\label{prop:kdf_hesap}
$F = F_X$ KDF'si olan $X$ için:
\begin{align*}
  P(X > x)   &= 1 - F(x), \\
  P(a < X \leq b) &= F(b) - F(a), \\
  P(X = x)   &= F(x) - F(x^-) \quad \text{(sol limit: } F(x^-) = \lim_{y \uparrow x} F(y)\text{)}, \\
  P(a \leq X \leq b) &= F(b) - F(a^-), \\
  P(a < X < b) &= F(b^-) - F(a).
\end{align*}
\end{onerme}

\begin{proof}
$P(a < X \leq b) = P(X \leq b) - P(X \leq a) = F(b) - F(a)$.
$P(X = x) = P(X \leq x) - P(X < x) = F(x) - \lim_{y \uparrow x} F(y) = F(x) - F(x^-)$.
Diğerleri bu iki formülden türetilir.
\end{proof}

\begin{hataanalizi}
"$P(X < x) = F(x)$" hatası: Sol açık aralık için $P(X < x) = F(x^-)$
olup $F$'nin $x$'te süreksiz olduğu noktalarda $F(x^-) < F(x)$.
Sağdan sürekli tanım $P(X \leq x) = F(x)$'i verir; $P(X < x)$ farklıdır.
\end{hataanalizi}

\begin{kendinleifade}
KDF'nin üç temel özelliğini (monoton, sağdan sürekli, sınırlar)
bir yabancı dil örneğiyle canlandırın: örneğin "dil testi skoru"nun
KDF'si nasıl görünürdü? Sürekli mi olurdu, yoksa sıçramalı mı?
\ogrencialani[2.5cm]
\end{kendinleifade}

\begin{mikroozet}
$F_X(x) = P(X \leq x)$. KDF: monoton, sağdan sürekli, $0 \to 1$ sınırları.
$P(X = x) = F(x) - F(x^-)$; süreklilik noktalarında bu sıfırdır.
\end{mikroozet}

% =====================================================================
\section{Kesikli Dağılımlar}
\label{sec:kesikli}
% =====================================================================

\begin{tanim}[Kesikli Rasgele Değişken]
\label{def:kesikli}
$X$ rasgele değişkeni \textbf{kesikli} (discrete) dır, eğer
sayılabilir bir $S \subset \mathbb{R}$ kümesi için $P(X \in S) = 1$ ise.
$S$'nin elemanlarına $X$'in \textbf{olası değerleri} denir.
\textbf{Olasılık kütle fonksiyonu} (probability mass function, PMF):
\[
  p_X(x) = P(X = x), \qquad x \in S.
\]
\end{tanim}

\begin{teorem}[PMF-KDF İlişkisi]
\label{thm:pmf_kdf}
$X$ kesikli rasgele değişken, $S = \{x_1, x_2, \ldots\}$ ($x_1 < x_2 < \cdots$) olsun. O zaman
\[
  F_X(x) = \sum_{x_i \leq x} p_X(x_i)
\]
ve $F_X$, her $x_i$'de büyüklüğü $p_X(x_i)$ olan sıçramayla
sağdan sürekli bir basamak fonksiyonudur.
\end{teorem}

\begin{proof}
$\{X \leq x\} = \bigsqcup_{x_i \leq x} \{X = x_i\}$ ayrık birleşimidir.
Sayılabilir toplamlılık:
$F_X(x) = P(X \leq x) = \sum_{x_i \leq x} P(X = x_i) = \sum_{x_i \leq x} p_X(x_i)$.
$y \downarrow x$ alındığında $\{x_i \leq y\}$ kümesi $y$'nin küçük aralıklarda
değişimi boyunca $\{x_i \leq x\}$'e yakınsadığından sağdan süreklilik sağlanır.
\end{proof}

\begin{ornek}[Binom]
\label{ex:binom_pmf}
$X \sim \Binom(n, p)$: $n$ bağımsız Bernoulli denemesinde başarı sayısı.
\[
  p_X(k) = \binom{n}{k} p^k (1-p)^{n-k}, \quad k = 0, 1, \ldots, n.
\]
$\sum_{k=0}^n p_X(k) = (p + (1-p))^n = 1$ (Binom teoremi).
\end{ornek}

% =====================================================================
\section{Mutlak Sürekli Dağılımlar ve Yoğunluk Fonksiyonu}
\label{sec:surekli}
% =====================================================================

\begin{kesif}
Bir terazi ibresi [0,1] aralığında rasgele bir değere işaret ediyor.
$P(X = 0.5) = ?$ — sonsuz çok nokta var, her birinin olasılığı sıfır olmalı.
Peki $P(0.3 \leq X \leq 0.7) = ?$ Bu sıfır değil.
Nasıl bir kavramla olasılıkları hesaplayabiliriz?
\end{kesif}

\begin{tanimgerekce}
Radon-Nikodym teoremi der ki: $P_X \ll \lambda$ (Lebesgue ölçüsüne göre
mutlak sürekli) ise $P_X(B) = \int_B f(x) \, d\lambda(x)$ yazan bir $f \geq 0$ var.
Bu $f$ "yoğunluk"tur — noktasal olasılık değil, uzunluk başına olasılık yoğunluğu.
\end{tanimgerekce}

\begin{tanim}[Mutlak Sürekli Rasgele Değişken]
\label{def:surekli_rd}
$X$ rasgele değişkenin \textbf{olasılık yoğunluk fonksiyonu} (probability density function, PDF)
$f_X : \mathbb{R} \to [0, \infty)$ mevcutsa ve her $B \in \mathcal{B}(\mathbb{R})$ için
\[
  P(X \in B) = \int_B f_X(x) \, dx
\]
sağlanıyorsa $X$'e \textbf{mutlak sürekli} (absolutely continuous) rasgele değişken denir.
Özellikle $F_X(x) = \int_{-\infty}^x f_X(t) \, dt$ ve $\int_{-\infty}^{+\infty} f_X(x) \, dx = 1$.
\end{tanim}

\begin{teorem}[PDF-KDF İlişkisi]
\label{thm:pdf_kdf}
$X$ mutlak sürekli rasgele değişken ve $f_X$ PDF'si olsun. O zaman:
\begin{enumerate}[(i)]
  \item $F_X(x) = \int_{-\infty}^x f_X(t) \, dt$, dolayısıyla $F_X$ süreklidir.
  \item $f_X$ sürekli olduğu her $x$ noktasında $F_X'(x) = f_X(x)$.
  \item Her $x$ için $P(X = x) = 0$.
  \item $f_X \geq 0$ n.k.\ ve $\int_{\mathbb{R}} f_X = 1$.
\end{enumerate}
\end{teorem}

\begin{proof}
\textit{(i)}: Tanımdan; $(-\infty, x] \in \mathcal{B}(\mathbb{R})$ seçilir.
\textit{(ii)}: Lebesgue farklılaştırma teoremi; $f_X$ sürekli ise
$\frac{d}{dx} \int_{-\infty}^x f_X(t) \, dt = f_X(x)$.
\textit{(iii)}: $P(X = x) = F_X(x) - F_X(x^-) = 0$ ($F_X$ sürekli olduğundan).
\textit{(iv)}: $f_X \geq 0$ a.e., $\int f_X = P(X \in \mathbb{R}) = 1$.
\end{proof}

\begin{hataanalizi}
"$f_X(x) = P(X = x)$" en yaygın hatadır. Doğrusu: $f_X(x) \, dx \approx P(x < X \leq x + dx)$
yani $f_X(x)$, küçük aralığa düşen \emph{olasılık yoğunluğudur}, bizzat olasılık değildir.
$f_X(x) > 1$ olabilir (örneğin $X \sim \mathrm{Uniform}(0, 0.1)$ ise $f_X(x) = 10$).
\end{hataanalizi}

\begin{karsiornek}
$f_X(x) = 10 \cdot \mathbf{1}_{[0, 0.1]}(x)$: $f_X(x) = 10 > 1$ ama yine de geçerli bir PDF.
$\int_0^{0.1} 10 \, dx = 1$. "Yoğunluk 1'i geçemez" iddiası yanlıştır.
\end{karsiornek}

\begin{ornek}[Normal Dağılım]
$X \sim \Normal(\mu, \sigma^2)$: PDF
\[
  f_X(x) = \frac{1}{\sigma\sqrt{2\pi}} \exp\!\left(-\frac{(x-\mu)^2}{2\sigma^2}\right).
\]
$\int_{-\infty}^{+\infty} f_X(x) \, dx = 1$ doğrulaması: Bölüm~\ref{chp:olcum_teorisi}'deki
Gauss integrali hesabına bakınız (Fubini-Tonelli).
\end{ornek}

\begin{mikroozet}
PDF: $P(X \in B) = \int_B f_X \, dx$; $f_X \geq 0$, $\int f_X = 1$.
$f_X(x)$ bir olasılık değil, yoğunluktur; $f_X > 1$ olabilir.
Her tekil noktanın olasılığı sıfırdır.
\end{mikroozet}

% =====================================================================
\section{Karma Dağılımlar ve Lebesgue Ayrışımı}
\label{sec:karma}
% =====================================================================

\begin{tanim}[Karma Dağılım]
\label{def:karma}
$X$ ne kesikli ne de mutlak sürekli olan bir rasgele değişken ise,
KDF'si $F_X = p \cdot F_{\text{kes}} + (1-p) \cdot F_{\text{sür}}$
biçiminde ($0 < p < 1$) \textbf{karma dağılım} (mixed distribution) olarak
yazılabilir; burada $F_{\text{kes}}$ kesikli, $F_{\text{sür}}$ mutlak sürekli
bir KDF'dir.
\end{tanim}

\begin{onerme}[Lebesgue Ayrışım Teoremi — Özel Durum]
\label{prop:lebesgue_ayrisim}
Her $F : \mathbb{R} \to [0,1]$ KDF'si şu biçimde tek türlü yazılabilir:
\[
  F = \alpha_1 F_{\text{saf-kes}} + \alpha_2 F_{\text{mut-sür}} + \alpha_3 F_{\text{tekil-sür}},
\]
$\alpha_1 + \alpha_2 + \alpha_3 = 1$, $\alpha_i \geq 0$; burada
$F_{\text{tekil-sür}}$ sürekli ama türevi neredeyse her yerde sıfır olan
(Lebesgue ölçüsüne göre mutlak sürekli olmayan) bir KDF'dir.
Cantor fonksiyonu (şeytan merdiveni) $F_{\text{tekil-sür}}$'ye örnektir.
\end{onerme}

\begin{hocanotu}
Pratikte tekil-sürekli bileşen nadiren karşımıza çıkar; uygulamalar
büyük ölçüde $\alpha_3 = 0$ durumunda (karma = kesikli + mutlak sürekli) çalışır.
Ancak tam teorik çerçeve için bu bileşeni bilmek önemlidir.
\end{hocanotu}

\begin{ornek}[Karma Dağılım]
Bir sigorta poliçesinde:
$P(X = 0) = 0.3$ (hasar yok), $X > 0$ durumunda $X \sim \mathrm{Exp}(\lambda)$.
KDF: $F_X(x) = 0$ ($x < 0$); $F_X(0) = 0.3$; $F_X(x) = 0.3 + 0.7(1 - e^{-\lambda x})$ ($x > 0$).
Ne tamamen kesikli (sürekli bir parçası var) ne tamamen sürekli (0'da sıçrama var).
\end{ornek}

% =====================================================================
\section{Birden Fazla Rasgele Değişken: Ortak Dağılımlar}
\label{sec:rd_ortak}
% =====================================================================

\begin{tanim}[Rasgele Vektör ve Ortak KDF]
\label{def:ortak_kdf}
$X_1, \ldots, X_n : \Omega \to \mathbb{R}$ rasgele değişkenler olsun.
$(X_1, \ldots, X_n) : \Omega \to \mathbb{R}^n$ \textbf{rasgele vektördür}
(ölçülebilirlik: $\mathcal{F}$'den $\mathcal{B}(\mathbb{R}^n)$'ye).
\textbf{Ortak KDF} (joint CDF):
\[
  F_{X_1,\ldots,X_n}(x_1, \ldots, x_n) = P(X_1 \leq x_1, \ldots, X_n \leq x_n).
\]
İki değişkenli durumda:
\begin{itemize}
  \item \textbf{Ortak PMF} (kesikli): $p(x,y) = P(X=x, Y=y)$.
  \item \textbf{Ortak PDF} (mutlak sürekli): $f_{X,Y}(x,y) \geq 0$,
        $P((X,Y) \in B) = \iint_B f_{X,Y}(x,y) \, dx \, dy$.
\end{itemize}
\end{tanim}

% =====================================================================
\section{Marjinal ve Koşullu Dağılımlar}
\label{sec:marjinal_kosullu}
% =====================================================================

\begin{tanim}[Marjinal Dağılım]
\label{def:marjinal}
$(X, Y)$ ortak dağılımlı rasgele vektör olsun.
$X$'in \textbf{marjinal dağılımı}:
\begin{itemize}
  \item Kesikli: $p_X(x) = \sum_y p(x, y)$.
  \item Sürekli: $f_X(x) = \int_{-\infty}^{+\infty} f_{X,Y}(x,y) \, dy$.
\end{itemize}
\end{tanim}

\begin{kritiksorgu}
"Aynı marjinallere sahip iki rasgele vektörün ortak dağılımı aynı mıdır?"
Hayır: Marjinaller ortak dağılımı belirlemez.
Örnek: $X \sim \mathrm{Uniform}(-1,1)$, $Y = X$ ile $Y = -X$ seçimleri
aynı marjinalleri verir ama farklı ortak dağılımlar üretir.
Bu gerçek, istatistiksel bağımlılık teorisinin (kopula, korelasyon) özüdür.
\end{kritiksorgu}

\begin{tanim}[Koşullu Dağılım]
\label{def:kosullu_dagilim}
\textbf{Kesikli durum:} $p_Y(y) > 0$ ise
\[
  p_{X|Y}(x \mid y) = P(X = x \mid Y = y) = \frac{p(x,y)}{p_Y(y)}.
\]
\textbf{Sürekli durum:} $f_Y(y) > 0$ ise
\[
  f_{X|Y}(x \mid y) = \frac{f_{X,Y}(x,y)}{f_Y(y)}.
\]
$f_{X|Y}(\cdot \mid y)$, $y$ sabit tutulduğunda $x$'in bir PDF'sidir:
$\int f_{X|Y}(x \mid y) \, dx = 1$.
\end{tanim}

\begin{teorem}[Toplam Olasılık — Sürekli Versiyon]
\label{thm:toplam_surekli}
$(X, Y)$ ortak sürekli dağılımlı olsun. O zaman
\[
  f_X(x) = \int_{-\infty}^{+\infty} f_{X|Y}(x \mid y) \, f_Y(y) \, dy.
\]
\end{teorem}

\begin{proof}
$f_X(x) = \int f_{X,Y}(x,y) \, dy = \int \frac{f_{X,Y}(x,y)}{f_Y(y)} \cdot f_Y(y) \, dy
= \int f_{X|Y}(x \mid y) \, f_Y(y) \, dy$.
\end{proof}

\begin{ornek}[Koşullu Dağılım Hesabı]
$f_{X,Y}(x,y) = 6x$ ($0 < x < y < 1$), sıfır başka yerde.
Marjinal: $f_Y(y) = \int_0^y 6x \, dx = 3y^2$.
Koşullu: $f_{X|Y}(x \mid y) = \frac{6x}{3y^2} = \frac{2x}{y^2}$, $0 < x < y$.
\end{ornek}

% =====================================================================
\section{Bağımsız Rasgele Değişkenler}
\label{sec:rd_bagimsizlik}
% =====================================================================

\begin{tanim}[Rasgele Değişken Bağımsızlığı]
\label{def:rd_bagimsiz}
$X_1, \ldots, X_n$ rasgele değişkenler \textbf{bağımsızdır} (independent)
ancak ve ancak her $B_1, \ldots, B_n \in \mathcal{B}(\mathbb{R})$ için
\[
  P(X_1 \in B_1, \ldots, X_n \in B_n) = \prod_{i=1}^n P(X_i \in B_i).
\]
Eşdeğer olarak $\sigma(X_1), \ldots, \sigma(X_n)$ $\sigma$-cebirleri bağımsızdır.
\end{tanim}

\begin{teorem}[Bağımsızlığın Eşdeğer Koşulları]
\label{thm:bagimsiz_denk}
$(X, Y)$ rasgele vektör için aşağıdakiler eşdeğerdir:
\begin{enumerate}[(i)]
  \item $X$ ve $Y$ bağımsızdır.
  \item Her $x, y$ için $F_{X,Y}(x,y) = F_X(x) \cdot F_Y(y)$.
  \item \textit{(Kesikli:)} Her $x, y$ için $p(x,y) = p_X(x) \cdot p_Y(y)$.
  \item \textit{(Sürekli:)} $f_{X,Y}(x,y) = f_X(x) \cdot f_Y(y)$ n.k.
  \item $f_{X|Y}(x \mid y) = f_X(x)$ n.k.\ (koşullu dağılım marjinale eşittir).
\end{enumerate}
\end{teorem}

\begin{proof}
$(i) \Leftrightarrow (ii)$: $B_i = (-\infty, x_i]$ seçilirse $(i)$'den $(ii)$ çıkar;
$\pi$-sistemi argümanıyla tersi de kurulur.
$(ii) \Leftrightarrow (iv)$: Sürekli durumda $F_{X,Y}(x,y) = \int_{-\infty}^x \int_{-\infty}^y f(s,t) \, dt \, ds$;
bu $F_X(x) F_Y(y)$'ye eşit olması $f = f_X \cdot f_Y$ a.e.\ ile eşdeğerdir
(Fubini).
$(iv) \Leftrightarrow (v)$: $f_{X|Y}(x \mid y) = f(x,y)/f_Y(y) = f_X(x)$.
\end{proof}

\begin{hataanalizi}
"Korelasyonsuz $\Rightarrow$ bağımsız" hatası: $\mathrm{Cor}(X, Y) = 0$
bağımsızlık için yeterli değildir.
Örnek: $X \sim \mathrm{Uniform}(-1,1)$, $Y = X^2$. O zaman $E[XY] = E[X^3] = 0 = E[X]E[Y]$,
yani $\mathrm{Cov}(X, Y) = 0$, ama $Y$ tamamen $X$'e bağımlıdır.
Yalnızca normal dağılım için $\mathrm{Cor} = 0 \Leftrightarrow$ bağımsızlık.
\end{hataanalizi}

\begin{mikroozet}
Bağımsızlık: $F_{X,Y} = F_X \cdot F_Y$ (veya PDF çarpımlanması).
Bağımsızlık $\Rightarrow$ korelasyonsuzluk; tersi genel olarak yanlış.
\end{mikroozet}

% =====================================================================
\section{Rasgele Değişkenlerin Fonksiyonları}
\label{sec:rd_donusum}
% =====================================================================

\begin{motivasyon}[Neden Dönüşüm?]
$X$ bilinen dağılımlı ise $Y = g(X)$'in dağılımı nedir?
Bu soru istatistiksel hesaplarda sürekli karşımıza çıkar:
örneğin $X \sim \Normal(0,1)$ ise $Y = X^2$'nin dağılımı chi-kare,
$Y = e^X$ ise lognormal dağılım verir.
\end{motivasyon}

\subsection{Tekil Boyutlu Dönüşüm}
\label{subsec:donusum_1d}

\begin{teorem}[Birebir Dönüşüm Yöntemi]
\label{thm:donusum_1d}
$X$ sürekli rasgele değişken, PDF'si $f_X$. $g : \mathbb{R} \to \mathbb{R}$
monoton, türevlenebilir ve birebir olsun, $Y = g(X)$.
$h = g^{-1}$ (ters fonksiyon) olmak üzere:
\[
  f_Y(y) = f_X(h(y)) \cdot \left|\frac{dh}{dy}\right|.
\]
\end{teorem}

\begin{proof}
\textbf{$g$ artan} olsun. O zaman
\[
  F_Y(y) = P(Y \leq y) = P(g(X) \leq y) = P(X \leq g^{-1}(y)) = F_X(h(y)).
\]
Türev alınca: $f_Y(y) = f_X(h(y)) \cdot h'(y) = f_X(h(y)) \cdot |h'(y)|$.

\textbf{$g$ azalan} olsun. O zaman
\[
  F_Y(y) = P(Y \leq y) = P(g(X) \leq y) = P(X \geq h(y)) = 1 - F_X(h(y)).
\]
Türev: $f_Y(y) = -f_X(h(y)) \cdot h'(y) = f_X(h(y)) \cdot |h'(y)|$ ($h' < 0$ olduğundan).
\end{proof}

\begin{ornek}[Lognormal]
$X \sim \Normal(\mu, \sigma^2)$, $Y = e^X$ (lognormal dağılım).
$h(y) = \ln y$, $|h'(y)| = 1/y$.
\[
  f_Y(y) = \frac{1}{y\sigma\sqrt{2\pi}} \exp\!\left(-\frac{(\ln y - \mu)^2}{2\sigma^2}\right), \quad y > 0.
\]
\end{ornek}

\begin{onerme}[Genel Monoton Olmayan Dönüşüm]
\label{prop:genel_donusum}
$g$ birebir olmayabilir. $g^{-1}(y) = \{x : g(x) = y\}$ sonlu küme ise:
\[
  f_Y(y) = \sum_{x : g(x) = y} f_X(x) \cdot \left|\frac{dg}{dx}\bigg|_{x}\right|^{-1}.
\]
\end{onerme}

\begin{ornek}[$Y = X^2$, $X \sim \Normal(0,1)$]
$g(x) = x^2$, $g^{-1}(y) = \{\sqrt{y}, -\sqrt{y}\}$ ($y > 0$).
$|dg/dx|^{-1} = 1/(2|x|)$.
\[
  f_Y(y) = f_X(\sqrt{y}) \cdot \frac{1}{2\sqrt{y}} + f_X(-\sqrt{y}) \cdot \frac{1}{2\sqrt{y}}
         = \frac{1}{\sqrt{2\pi y}} e^{-y/2}, \quad y > 0.
\]
Bu chi-kare dağılımının ($\chi^2_1$) PDF'sidir.
\end{ornek}

\subsection{Çok Boyutlu Dönüşüm ve Jacobian}
\label{subsec:donusum_jacobian}

\begin{teorem}[Jacobian Formülü]
\label{thm:jacobian}
$(X_1, \ldots, X_n)$ ortak sürekli dağılımlı; PDF'si $f_{X_1,\ldots,X_n}$.
$g : \mathbb{R}^n \to \mathbb{R}^n$ birebir ve türevlenebilir olsun,
$(Y_1, \ldots, Y_n) = g(X_1, \ldots, X_n)$.
Ters dönüşüm $h = g^{-1}$ ve Jacobian matrisi
\[
  J = \frac{\partial(x_1,\ldots,x_n)}{\partial(y_1,\ldots,y_n)}
    = \left(\frac{\partial h_i}{\partial y_j}\right)_{n \times n}
\]
olmak üzere:
\[
  f_{Y_1,\ldots,Y_n}(y_1,\ldots,y_n) = f_{X_1,\ldots,X_n}(h_1(y),\ldots,h_n(y)) \cdot |\det J|.
\]
\end{teorem}

\begin{proof}
Çok değişkenli değişken değiştirme formülü:
$P((Y_1,\ldots,Y_n) \in B) = \int_B f_Y \, dy = \int_{g^{-1}(B)} f_X \, dx$.
$dx = |\det J| \, dy$ (Jacobian dönüşümü); ikisini eşitleyince $f_Y$ elde edilir.
\end{proof}

\begin{ornek}[Polar Dönüşüm]
$(X, Y)$ bağımsız, her ikisi $\Normal(0,1)$.
$R = \sqrt{X^2+Y^2}$, $\Theta = \arctan(Y/X)$.
Ters: $x = r\cos\theta$, $y = r\sin\theta$.
$\det J = r$. Ortak PDF $f_{R,\Theta}(r, \theta) = f_{X,Y}(r\cos\theta, r\sin\theta) \cdot r$.
\[
  f_{X,Y}(x,y) = \frac{1}{2\pi} e^{-(x^2+y^2)/2}
  \implies f_{R,\Theta}(r,\theta) = r e^{-r^2/2} \cdot \frac{1}{2\pi},
\]
$r > 0$, $\theta \in [0, 2\pi)$.
Marjinaller: $\Theta \sim \mathrm{Uniform}(0, 2\pi)$ ve $R^2 \sim \chi^2_2$ (Rayleigh dağılımı).
Bu Box-Muller dönüşümünün temelidir.
\end{ornek}

\subsection{Toplam Dağılımı: Evrişim}
\label{subsec:evrisim}

\begin{tanim}[Evrişim (Konvolüsyon)]
\label{def:evrisim}
$X, Y$ bağımsız, ortak sürekli dağılımlı; $Z = X + Y$ olsun. O zaman
\[
  f_Z(z) = (f_X * f_Y)(z) = \int_{-\infty}^{+\infty} f_X(z - y) \, f_Y(y) \, dy.
\]
Kesikli durumda: $p_Z(z) = \sum_y p_X(z-y) \, p_Y(y)$.
\end{tanim}

\begin{proof}
$F_Z(z) = P(X + Y \leq z) = \int_{-\infty}^{+\infty} P(X \leq z - y) \, f_Y(y) \, dy
= \int_{-\infty}^{+\infty} F_X(z-y) \, f_Y(y) \, dy$.
$z$'ye göre türev: Leibniz kuralıyla $f_Z(z) = \int f_X(z-y) f_Y(y) \, dy$.
\end{proof}

\begin{ornek}[Normal + Normal = Normal]
$X \sim \Normal(\mu_1, \sigma_1^2)$, $Y \sim \Normal(\mu_2, \sigma_2^2)$ bağımsız.
Evrişim hesabı: $Z = X + Y \sim \Normal(\mu_1+\mu_2, \sigma_1^2+\sigma_2^2)$.
Hesap, Gauss integrallerini tamamlama yöntemiyle yapılır;
Bölüm~\ref{chp:beklenti}'deki karakteristik fonksiyon yöntemi çok daha hızlıdır.
\end{ornek}

% ---- Disiplinlerarası Bağlantı (TYMM) ----
\begin{kopru}[Disiplinlerarası — Finans ve Mühendislik]
\textbf{Finans:} Hisse senedi getirisinin lognormal modellenmesi (Black-Scholes),
rasgele değişken dönüşüm teorisinin doğrudan uygulamasıdır.\\
\textbf{Mühendislik:} Haberleşmede sinyal+gürültü modeli $Z = S + N$,
$N \sim \Normal(0, \sigma^2)$: evrişim yoluyla alıcı tasarımı.\\
\textbf{Fizik:} Kuantum mekaniğinde konum ve momentum dağılımları
PDF kavramıyla doğrudan ifade edilir; Fourier dönüşümü ile karakteristik
fonksiyon arasındaki bağlantı Heisenberg belirsizlik ilkesini yansıtır.
\end{kopru}

% ---- Öz-Yansıtma ----
\begin{ozyansima}
\begin{itemize}[leftmargin=1.8em]
  \item "Rasgele değişken bir sayı mı, yoksa fonksiyon mu?" — cevabınız ve gerekçeniz?
  \item KDF'nin hangi özelliği sizi en çok şaşırttı?
  \item Sürekli durumda $P(X = x) = 0$ gerçeği sezgiye aykırı geliyor mu?
        Bunu fizik veya mühendislikten bir örnekle nasıl uzlaştırırsınız?
  \item Bağımsızlık ile korelasyonsuzluk arasındaki fark neden önemlidir?
\end{itemize}
\end{ozyansima}

% ---- Köprü Notu ----
\begin{kopru}
\textbf{Önceki bölüm:} Bölüm~\ref{chp:olasilik_uzaylari} — Kolmogorov aksiyomları,
bağımsızlık olaylar düzeyinde; burada rasgele değişkenlere taşındı.\\
\textbf{Sonraki bölüm:} Bölüm~\ref{chp:beklenti} — Dağılımlar verildiğinde
$E[X]$, $\mathrm{Var}(X)$, momentler; Lebesgue integrali olarak beklenti.\\
\textbf{Dijital katman:} Dönüşüm formüllerinin sayısal doğrulaması ve
Monte Carlo simülasyonu \texttt{bolum03\_donusumler.ipynb} defterinde.
\defterbaglantisi{bolum03\_donusumler.ipynb}{%
  Jacobian dönüşümü, evrişim ve KDF simülasyonu}
\end{kopru}

% =====================================================================
\section{Çok Boyutlu Normal Dağılım}
\label{sec:cok_boyutlu_normal}
% =====================================================================

\begin{motivasyon}[Neden Çok Boyutlu Normal?]
Tek boyutlu Normal dağılım, pek çok bağımsız etkinin toplamının limit dağılımı
(MLT) olarak doğal biçimde ortaya çıkar.
Çok boyutlu uzantısı, bağımlı Normal rasgele değişkenlerin tam ailesini verir
ve lineer dönüşümler altında kapalıdır.
İstatistik, makine öğrenmesi ve sinyal işlemede varsayılan "default" dağılımdır.
\end{motivasyon}

\begin{tanim}[Çok Boyutlu Normal Dağılım]
\label{def:cok_boyutlu_normal}
$\mathbf{X} = (X_1, \ldots, X_k)^\top$ rasgele vektör.
$\mathbf{X} \sim \Normal_k(\boldsymbol{\mu}, \boldsymbol{\Sigma})$
(\textbf{$k$-boyutlu Normal}, multivariate normal) denir eğer:
\begin{enumerate}[(i)]
  \item $\boldsymbol{\mu} = E[\mathbf{X}] \in \mathbb{R}^k$ ortalama vektör.
  \item $\boldsymbol{\Sigma} = \mathrm{Cov}(\mathbf{X}) \in \mathbb{R}^{k\times k}$ simetrik pozitif yarı-belirli kovaryans matrisi.
  \item $\boldsymbol{\Sigma}$ pozitif belirli ise ortak PDF:
  \[
    f(\mathbf{x}) = \frac{1}{(2\pi)^{k/2}|\boldsymbol{\Sigma}|^{1/2}}
    \exp\!\left(-\frac{1}{2}(\mathbf{x}-\boldsymbol{\mu})^\top \boldsymbol{\Sigma}^{-1}(\mathbf{x}-\boldsymbol{\mu})\right).
  \]
\end{enumerate}
Eşdeğer tanım: $\mathbf{X} = \boldsymbol{\mu} + A\mathbf{Z}$, $\mathbf{Z} \sim \Normal_m(\mathbf{0}, I_m)$ bağımsız,
$\boldsymbol{\Sigma} = AA^\top$.
\end{tanim}

\begin{teorem}[Çok Boyutlu Normal'in Temel Özellikleri]
\label{thm:cb_normal}
$\mathbf{X} \sim \Normal_k(\boldsymbol{\mu}, \boldsymbol{\Sigma})$:
\begin{enumerate}[(i)]
  \item \textbf{Marjinaller:} $X_i \sim \Normal(\mu_i, \sigma_{ii})$ (marjinallar normal).
  \item \textbf{Lineer dönüşüm:} $A\mathbf{X} + \mathbf{b} \sim \Normal(A\boldsymbol{\mu}+\mathbf{b}, A\boldsymbol{\Sigma}A^\top)$.
  \item \textbf{Bağımsızlık:} $X_i$ ve $X_j$ bağımsız $\Leftrightarrow$ $\sigma_{ij} = 0$ (Normal için korelasyonsuzluk = bağımsızlık!).
  \item \textbf{Koşullu dağılım:} $\mathbf{X}$'i $(\mathbf{X}_1, \mathbf{X}_2)$ şeklinde parçalayın.
  $\mathbf{X}_1 \mid \mathbf{X}_2 = \mathbf{x}_2 \sim \Normal(\boldsymbol{\mu}_{1|2},\, \boldsymbol{\Sigma}_{11|2})$:
  \begin{align*}
    \boldsymbol{\mu}_{1|2} &= \boldsymbol{\mu}_1 + \boldsymbol{\Sigma}_{12}\boldsymbol{\Sigma}_{22}^{-1}(\mathbf{x}_2 - \boldsymbol{\mu}_2),\\
    \boldsymbol{\Sigma}_{11|2} &= \boldsymbol{\Sigma}_{11} - \boldsymbol{\Sigma}_{12}\boldsymbol{\Sigma}_{22}^{-1}\boldsymbol{\Sigma}_{21}.
  \end{align*}
\end{enumerate}
\end{teorem}

\begin{proof}[(i) ve (ii)]
(i) Marjinal: $X_i = \mathbf{e}_i^\top \mathbf{X}$ ($\mathbf{e}_i$ birim vektör); (ii) özel durumu.\\
(ii) KF: $\varphi_{A\mathbf{X}+\mathbf{b}}(\mathbf{t}) = e^{i\mathbf{t}^\top \mathbf{b}} \varphi_{\mathbf{X}}(A^\top \mathbf{t})$.
$\varphi_{\mathbf{X}}(\mathbf{t}) = e^{i\mathbf{t}^\top\boldsymbol{\mu} - \mathbf{t}^\top\boldsymbol{\Sigma}\mathbf{t}/2}$ (Normal KF).
$\varphi_{A\mathbf{X}+\mathbf{b}}(\mathbf{t}) = e^{i\mathbf{t}^\top(A\boldsymbol{\mu}+\mathbf{b}) - \mathbf{t}^\top A\boldsymbol{\Sigma}A^\top\mathbf{t}/2}$: $\Normal(A\boldsymbol{\mu}+\mathbf{b}, A\boldsymbol{\Sigma}A^\top)$ KF'si.
\end{proof}

\begin{ornek}[İki Boyutlu Normal]
$k=2$: $\boldsymbol{\mu} = (\mu_1, \mu_2)^\top$,
$\boldsymbol{\Sigma} = \begin{pmatrix}\sigma_1^2 & \rho\sigma_1\sigma_2 \\ \rho\sigma_1\sigma_2 & \sigma_2^2\end{pmatrix}$.
PDF:
\[
  f(x_1,x_2) = \frac{1}{2\pi\sigma_1\sigma_2\sqrt{1-\rho^2}}
  \exp\!\left(-\frac{1}{2(1-\rho^2)}\left[\frac{(x_1-\mu_1)^2}{\sigma_1^2}
  - \frac{2\rho(x_1-\mu_1)(x_2-\mu_2)}{\sigma_1\sigma_2}
  + \frac{(x_2-\mu_2)^2}{\sigma_2^2}\right]\right).
\]
$X_1 \mid X_2 = x_2 \sim \Normal\!\left(\mu_1 + \rho\frac{\sigma_1}{\sigma_2}(x_2-\mu_2),\,
\sigma_1^2(1-\rho^2)\right)$.
\end{ornek}

\begin{hataanalizi}
"Marjinalleri Normal olan rasgele vektör çok boyutlu Normal'dir."
Yanlış: Korrelasyon yapısı yanlış olabilir. Örnek: $Z \sim \Normal(0,1)$,
$W = Z \cdot \mathrm{sgn}(Z)$; hem $Z$ hem $W$ normal marjinale sahip
ama $(Z, W)$ çift boyutlu normal değil (ortak PDF oval konturlar çizmiyor).
Çok boyutlu Normal tanımı KF veya affine dönüşüm koşuluyla yapılmalı.
\end{hataanalizi}

\begin{mikroozet}
$\mathbf{X}\sim\Normal_k(\boldsymbol{\mu},\boldsymbol{\Sigma})$: marjinaller normal; lineer dönüşümler kapalı;
Normal'de korelasyonsuzluk = bağımsızlık; koşullu dağılım da normaldir.
\end{mikroozet}

% =====================================================================
\section{Kuantil Fonksiyon ve Olasılık İntegral Dönüşümü}
\label{sec:kuantil_fonksiyon}
% =====================================================================

\begin{tanim}[Kuantil Fonksiyon]
\label{def:kuantil_fonksiyon}
$F$ KDF için \textbf{kuantil fonksiyon} (quantile function), genelleştirilmiş ters:
\[
  F^{-1}(u) = \inf\{x : F(x) \geq u\}, \quad u \in (0,1).
\]
$u$'nun $p$. yüzdelik (\textit{percentile}): $F^{-1}(p)$.
$F^{-1}(1/2)$: medyan. $F^{-1}(1/4), F^{-1}(3/4)$: çeyrekler.
\end{tanim}

\begin{teorem}[Olasılık İntegral Dönüşümü]
\label{thm:oit}
\textbf{İleri:} $X$ sürekli rasgele değişken, $F_X$ KDF. $F_X(X) \sim \Uniform(0,1)$.\\
\textbf{Geri:} $U \sim \Uniform(0,1)$. $X = F^{-1}(U)$ rasgele değişkeni $F$ KDF'ye sahip.
\end{teorem}

\begin{proof}
\textbf{İleri:} $P(F_X(X) \leq u) = P(X \leq F_X^{-1}(u)) = F_X(F_X^{-1}(u)) = u$.
(Son eşitlik: $F_X$ sürekli ise $F_X \circ F_X^{-1} = \mathrm{id}$.)\\
\textbf{Geri:} $P(F^{-1}(U) \leq x) = P(U \leq F(x)) = F(x)$ ($U$ uniform).
\end{proof}

\begin{ornek}[Simülasyon Uygulaması]
$\Exp(\lambda)$ dağılımından örneklem üretmek:
$F^{-1}(u) = -\ln(1-u)/\lambda$.
$U \sim \Uniform(0,1)$: $X = -\ln(1-U)/\lambda \sim \Exp(\lambda)$.
Bu standart simülasyon tekniğinin (Inverse Transform Method) matematiksel temelidir.
Normal: $\Phi^{-1}(U)$ kapalı form yok; sayısal yöntemler ($\Phi^{-1} \approx$ probit fonksiyon) kullanılır.
\end{ornek}

\begin{teorem}[Kuantil-Kuantil İlişkisi]
\label{thm:qq_iliskisi}
$X \sim F_X$, $Y \sim F_Y$. $Y \overset{d}{=} F_Y^{-1}(F_X(X))$.
\emph{QQ-plot} prensibi: $X$ örneklemi ile $F_Y^{-1}$ karşılaştırılarak $F_Y$ uyumu test edilir.
\end{teorem}

\begin{mikroozet}
$F^{-1}(u) = \inf\{x: F(x)\geq u\}$: kuantil fonksiyon.
OID: $F(X)\sim\Uniform$; $F^{-1}(U)\sim F$.
Simülasyon: her sürekli dağılım $\Uniform(0,1)$'den üretilir.
\end{mikroozet}

% =====================================================================
%  ALIŞTIRMALAR
% =====================================================================

\cozumlualistirmalar

\begin{alistirma}[Al.3.1]{A}{Ölçülebilirlik Kontrolü}
$\Omega = \{1, 2, 3, 4\}$, $\mathcal{F} = \{\emptyset, \{1,2\}, \{3,4\}, \Omega\}$.
$X(\omega) = \omega$ ve $Y(\omega) = \mathbf{1}_{\{1,2\}}(\omega)$ fonksiyonlarından
hangisi rasgele değişkendir? Gerekçelendirin.
\end{alistirma}
\alcozum{%
$Y$: $\{Y \leq 0\} = \{3,4\} \in \mathcal{F}$; $\{Y \leq 1\} = \Omega \in \mathcal{F}$;
yani tüm alt kümeler $\mathcal{F}$'de. $Y$ rasgele değişkendir.\\
$X$: $\{X \leq 1\} = \{1\} \notin \mathcal{F}$. $X$ rasgele değişken değildir.}

\begin{alistirma}[Al.3.2]{A}{KDF'den Olasılık Hesabı}
$X$ rasgele değişkenin KDF'si:
\[
  F(x) = \begin{cases} 0 & x < 0 \\ x/2 & 0 \leq x < 1 \\ 1/2 + (x-1)/4 & 1 \leq x < 3 \\ 1 & x \geq 3 \end{cases}
\]
$P(X = 0)$, $P(X = 1)$, $P(0.5 < X \leq 2)$ ve $P(X = 3)$'ü bulun.
\end{alistirma}
\alcozum{%
$P(X = 0) = F(0) - F(0^-) = 0 - 0 = 0$ ($F$ 0'da sürekli).\\
$P(X = 1) = F(1) - F(1^-) = (1/2+0) - (1/2) = 0$ (sürekli).\\
$P(0.5 < X \leq 2) = F(2) - F(0.5) = (1/2 + 1/4) - 1/4 = 3/4 - 1/4 = 1/2$.\\
$P(X = 3) = F(3) - F(3^-) = 1 - (1/2 + 2/4) = 1 - 1 = 0$.
Bu KDF tamamen sürekli; dağılım mutlak süreklidir (karma bile değil).}

\begin{alistirma}[Al.3.3]{B}{Koşullu Dağılım}
$(X,Y)$ ortak PDF'si $f(x,y) = c\,xy$, $0 < x < 1$, $0 < y < 2$.
(a) $c$'yi bulun. (b) $f_{X|Y}(x \mid y)$'i bulun.
(c) $P(X > 0.5 \mid Y = 1)$'i hesaplayın.
\end{alistirma}
\alcozum{%
(a) $\int_0^1 \int_0^2 cxy \, dy \, dx = c \int_0^1 x \, dx \int_0^2 y \, dy = c \cdot \frac{1}{2} \cdot 2 = c$. $c = 1$.\\
(b) $f_Y(y) = \int_0^1 xy \, dx = y/2$. $f_{X|Y}(x \mid y) = \frac{xy}{y/2} = 2x$, $0 < x < 1$. ($y$'den bağımsız!)\\
(c) $P(X > 0.5 \mid Y = 1) = \int_{0.5}^1 2x \, dx = [x^2]_{0.5}^1 = 1 - 0.25 = 0.75$.}

\begin{alistirma}[Al.3.4]{B}{Jacobian Dönüşümü}
$X \sim \mathrm{Uniform}(0,1)$. $Y = -\ln X$ dönüşümü ile $Y$'nin dağılımını bulun.
\end{alistirma}
\alcozum{%
$g(x) = -\ln x$ azalan; ters $h(y) = e^{-y}$, $|h'(y)| = e^{-y}$.\\
$f_X(x) = 1$ ($0 < x < 1$), $x$ yerine $e^{-y}$ yazılırsa:
$f_Y(y) = f_X(e^{-y}) \cdot e^{-y} = 1 \cdot e^{-y} = e^{-y}$, $y > 0$.\\
$Y \sim \mathrm{Exp}(1)$. Bu quantile dönüşümü yönteminin temelidir:
herhangi bir dağılımdan örnek üretmek için KDF'nin tersini kullanmak.}

% ---

\alistirmalar

\begin{alistirma}[Al.3.5]{A}{KDF Özellikleri}
$F(x) = \frac{1}{1+e^{-x}}$ (lojistik fonksiyon) bir KDF midir?
Üç temel özelliği kontrol edin. Bu KDF mutlak sürekli bir dağılımı mı temsil eder?
PDF'yi bulun.
\end{alistirma}
\alcozum{%
Monoton artan: $F'(x) = e^{-x}/(1+e^{-x})^2 > 0$. $\checkmark$\\
Sağdan sürekli: $F$ her yerde türevlenebilir, dolayısıyla sürekli. $\checkmark$\\
Sınırlar: $F(-\infty) = 0$, $F(+\infty) = 1$. $\checkmark$\\
PDF: $f(x) = F'(x) = e^{-x}/(1+e^{-x})^2$. Bu lojistik dağılımdır.}

\begin{alistirma}[Al.3.6]{A}{PMF Doğrulama}
$p(k) = c \cdot (1/2)^k$, $k = 1, 2, 3, \ldots$
(a) $c$'yi bulun. (b) $P(X > 3)$'ü hesaplayın. (c) $P(X \text{ çift})$'i bulun.
\end{alistirma}
\alcozum{%
(a) $\sum_{k=1}^\infty c(1/2)^k = c \cdot (1/2)/(1-1/2) = c = 1$.\\
(b) $P(X > 3) = \sum_{k=4}^\infty (1/2)^k = (1/2)^4/(1-1/2) = 1/8$.\\
(c) $P(\text{çift}) = \sum_{k=1}^\infty (1/2)^{2k} = (1/4)/(1-1/4) = 1/3$.}

\begin{alistirma}[Al.3.7]{A}{Marjinal Dağılımlar}
$(X, Y)$ ortak PMF: $p(0,0)=0.1$, $p(0,1)=0.2$, $p(1,0)=0.3$, $p(1,1)=0.4$.
Marjinal PMF'leri $p_X$ ve $p_Y$'yi bulun. $X$ ve $Y$ bağımsız mıdır?
\end{alistirma}
\alcozum{%
$p_X(0) = 0.3$, $p_X(1) = 0.7$. $p_Y(0) = 0.4$, $p_Y(1) = 0.6$.\\
Bağımsız mı? $p_X(0) \cdot p_Y(0) = 0.3 \cdot 0.4 = 0.12 \neq 0.1 = p(0,0)$. Hayır.}

\begin{alistirma}[Al.3.8]{B}{Evrişim}
$X, Y$ bağımsız, $X \sim \mathrm{Uniform}(0,1)$, $Y \sim \mathrm{Uniform}(0,1)$.
$Z = X + Y$'nin PDF'sini evrişim yöntemiyle bulun ve grafiğini çizin.
\end{alistirma}
\alcozum{%
$f_Z(z) = \int_0^1 f_X(z-y) f_Y(y) \, dy = \int_0^1 \mathbf{1}_{[0,1]}(z-y) \, dy$.\\
$0 \leq z \leq 1$: $z-y \in [0,1]$ için $y \in [z-1, z] \cap [0,1] = [0, z]$; $f_Z(z) = z$.\\
$1 < z \leq 2$: $y \in [z-1, 1]$; $f_Z(z) = 2-z$.\\
Üçgen dağılım: $f_Z(z) = z$ (artan), $2-z$ (azalan). Merkezi Limit Teoremi'nin habercisi.}

\begin{alistirma}[Al.3.9]{B}{Dönüşüm: Karekök}
$X \sim \mathrm{Exp}(\lambda)$. $Y = \sqrt{X}$'in PDF'sini bulun.
$E[Y]$'yi $f_Y$ üzerinden ifade edin (Bölüm~\ref{chp:beklenti}'de hesaplanacak).
\end{alistirma}
\alcozum{%
$g(x) = \sqrt{x}$ artan; $h(y) = y^2$, $|h'(y)| = 2y$.\\
$f_X(x) = \lambda e^{-\lambda x}$ ($x > 0$).\\
$f_Y(y) = \lambda e^{-\lambda y^2} \cdot 2y = 2\lambda y e^{-\lambda y^2}$, $y > 0$.\\
Bu Rayleigh dağılımıdır (parametre $1/\sqrt{2\lambda}$).}

\begin{alistirma}[Al.3.10]{B}{Bağımsızlık Kontrolü}
$(X,Y)$ PDF'si $f(x,y) = e^{-x-y}$, $x > 0$, $y > 0$. $X$ ve $Y$ bağımsız mıdır?
Aynı soruyu $f(x,y) = 2e^{-x-y}$, $0 < y < x < \infty$ için yanıtlayın.
\end{alistirma}
\alcozum{%
Birinci: $f(x,y) = e^{-x} \cdot e^{-y} = f_X(x) \cdot f_Y(y)$. Evet, bağımsız.\\
İkinci: $f_X(x) = \int_0^x 2e^{-x-y} dy = 2e^{-x}(1-e^{-x})$. $f_Y(y) = \int_y^\infty 2e^{-x-y} dx = 2e^{-2y}$.
$f_X(x) \cdot f_Y(y) \neq f(x,y)$. Hayır, bağımsız değil.}

\begin{alistirma}[Al.3.11]{B}{Çok Boyutlu Jacobian}
$(X, Y)$ bağımsız $\Normal(0,1)$. $U = X + Y$, $V = X - Y$ dönüşümü ile
$(U, V)$'nin ortak dağılımını bulun. $U$ ve $V$ bağımsız mıdır?
\end{alistirma}
\alcozum{%
Ters: $x = (u+v)/2$, $y = (u-v)/2$. $J = \partial(x,y)/\partial(u,v)$:
$\partial x/\partial u = 1/2$, $\partial x/\partial v = 1/2$,
$\partial y/\partial u = 1/2$, $\partial y/\partial v = -1/2$.
$|\det J| = |{-1/4-1/4}| = 1/2$.\\
$f_{U,V}(u,v) = f_X((u+v)/2) f_Y((u-v)/2) \cdot 1/2
= \frac{1}{2\pi} e^{-u^2/4} e^{-v^2/4} \cdot \frac{1}{2}
= \frac{1}{4\pi} e^{-(u^2+v^2)/4}$.\\
$U \sim \Normal(0, 2)$, $V \sim \Normal(0, 2)$ ve çarpımlanıyor: bağımsızlar.}

\begin{alistirma}[Al.3.12]{C}{İspat: Bağımsız Değişkenlerin Toplamının KDF'si}
$X, Y$ bağımsız rasgele değişkenler, $Z = X + Y$ olsun.
(a) $F_Z(z) = \int_{-\infty}^{+\infty} F_X(z - y) \, dF_Y(y)$ olduğunu ispatlayın
(Lebesgue-Stieltjes integrali).
(b) Mutlak sürekli durumda bu formülün evrişim formülüne nasıl indirgeneceğini gösterin.
\end{alistirma}
\alcozum{%
(a) Bağımsızlıktan:
$F_Z(z) = P(X+Y \leq z) = E[\mathbf{1}_{X+Y \leq z}]
= E_Y[E_X[\mathbf{1}_{X \leq z-Y}]] = \int P(X \leq z-y) \, dF_Y(y) = \int F_X(z-y) \, dF_Y(y)$.\\
(b) Sürekli $Y$: $dF_Y(y) = f_Y(y) \, dy$. O zaman $F_Z(z) = \int F_X(z-y) f_Y(y) \, dy$.
$z$'ye göre türev (Leibniz): $f_Z(z) = \int f_X(z-y) f_Y(y) \, dy = (f_X * f_Y)(z)$.}

\begin{alistirma}[Al.3.13]{B}{Olasılık İntegral Dönüşümü}
$X \sim \Gama(\alpha, \beta)$. $U = F_X(X)$ dönüşümünü tanımlayın.
(a) $U$'nun dağılımını bulun.
(b) Geri dönüşüm: $U \sim \Uniform(0,1)$ verildiğinde $\Gama(2,1)$ dağılımı
nasıl simüle edersiniz? (İpucu: Üstel dağılım üzerinden)
(c) OID'nin QQ-plot oluşturmadaki rolünü açıklayın.
\end{alistirma}
\alcozum{%
(a) OID (Teorem~\ref{thm:oit}): $F_X(X) \sim \Uniform(0,1)$. Sürekli dağılım için geçerli; $\Gama$ sürekli. $U \sim \Uniform(0,1)$.\\
(b) $\Gama(2,1) = \Exp(1)+\Exp(1)$ (iki bağımsız Üstel toplamı).
$U_1, U_2 \sim \Uniform(0,1)$ bağımsız.
$X = -\ln(1-U_1) - \ln(1-U_2) = -\ln((1-U_1)(1-U_2))$.
Bu $\Gama(2,1)$ dağılımını verir.\\
(c) Veri $\{X_i\}$ için $\{F_0^{-1}(i/(n+1))\}$ teorik kuantiller ile sıralı örneklem karşılaştırılır.
$F_0 = F_X$ ise noktalar doğrusal dizilir.}

\begin{alistirma}[Al.3.14]{B}{İki Boyutlu Normal}
$(X_1, X_2) \sim \Normal_2\!\left((0,0)^\top, \begin{pmatrix}1 & \rho\\\rho & 1\end{pmatrix}\right)$.
(a) $X_1 + X_2$ ve $X_1 - X_2$'nin dağılımını bulun.
(b) $X_1 + X_2$ ve $X_1 - X_2$ bağımsız mıdır?
(c) $X_1 \mid X_2 = 1$'in dağılımını verin.
\end{alistirma}
\alcozum{%
(a) Lineer dönüşüm: $A = \begin{pmatrix}1&1\\1&-1\end{pmatrix}$, $\boldsymbol{\Sigma} = \begin{pmatrix}1&\rho\\\rho&1\end{pmatrix}$.
$A\boldsymbol{\Sigma}A^\top = \begin{pmatrix}2+2\rho & 0 \\ 0 & 2-2\rho\end{pmatrix}$.
$X_1+X_2 \sim \Normal(0, 2+2\rho)$; $X_1-X_2 \sim \Normal(0, 2-2\rho)$.\\
(b) Kovaryans sıfır (matris köşegen) $\Rightarrow$ Normal için bağımsız. Evet, bağımsız.\\
(c) $X_1|X_2=1 \sim \Normal(\rho, 1-\rho^2)$.}

\begin{alistirma}[Al.3.15]{C}{Kuantil Fonksiyon — Karma Dağılım}
$X$ karma dağılım: $P(X = 0) = 1/2$, $P(X > 0) = 1/2$ ve $X \mid X>0 \sim \Exp(1)$.
(a) $F_X(x)$'i her $x$ için yazın.
(b) $F_X^{-1}(u)$'yu $u \in (0,1)$ için hesaplayın.
(c) $U \sim \Uniform(0,1)$ kullanarak $X$ dağılımından nasıl örneklem alırsınız?
\end{alistirma}
\alcozum{%
(a) $x < 0$: $F(x) = 0$. $x = 0$: $F(0) = 1/2$. $x > 0$: $F(x) = 1/2 + (1/2)(1-e^{-x})$.\\
(b) $u \leq 1/2$: $F^{-1}(u) = 0$. $u > 1/2$: $F(x) = u \Rightarrow 1-e^{-x} = 2u-1 \Rightarrow x = -\ln(2-2u)$.
$F^{-1}(u) = -\ln(2(1-u))$ ($u > 1/2$).\\
(c) $U \sim \Uniform(0,1)$: $U \leq 1/2$ ise $X = 0$; $U > 1/2$ ise $X = -\ln(2(1-U))$.}

% =====================================================================
\endinput

% =====================================================================
%  BÖLÜM 4: Beklenti ve Momentler
%  Kaynak: Feller; Billingsley; Klenke; Shiryaev
% =====================================================================
\chapter{Beklenti ve Momentler}
\label{chp:beklenti}

\bolumalinti{%
  Beklenti, belirsizliğin ağırlık merkezidir.%
}{Patrick Billingsley, \textit{Probability and Measure}}

% ---- Kazanımlar (TYMM) ----
\begin{kazanimlar}
Bu bölümün sonunda öğrenci:
\begin{itemize}[leftmargin=1.8em]
  \item Beklentiyi Lebesgue integrali olarak tanımlayabilir ve
        bu tanımın kesikli/sürekli formüllere nasıl indirgeneceğini gösterebilir.
  \item Beklentinin doğrusallık, monotonluk ve önemli eşitsizliklerini
        (Jensen, Markov, Chebyshev) ispatlayabilir ve uygulayabilir.
  \item Varyans, kovaryans ve korelasyon kavramlarını hesaplayabilir;
        bağımsızlıkla ilişkilerini açıklayabilir.
  \item Moment üreten fonksiyon (MÜF) ve karakteristik fonksiyonun (KF)
        tanımını, yakınsama bölgesini ve momentlerle ilişkisini açıklayabilir.
  \item Koşullu beklentiyi $L^2$ projeksiyonu olarak yorumlayabilir,
        temel özelliklerini ve yineleme (kule) özelliğini ispatlayabilir.
\end{itemize}
\end{kazanimlar}

% ---- Motivasyon ----
\begin{motivasyon}
Bir rasgele değişkenin dağılımı tam bilgiyi içerir; ama çoğu zaman bu kadar bilgiye
ihtiyaç duymayız. Dağılımı birkaç sayıyla özetlemek istiyoruz: "ortalama değer nedir?"
ve "ne kadar saçılım var?" Beklenti ve varyans bu iki soruyu yanıtlar.

Daha incelikli özetler için momentler, çarpıklık (skewness) ve basıklık (kurtosis)
devreye girer. Tüm momentleri birden kodlayan moment üreten fonksiyon ve
karakteristik fonksiyon ise dağılım teorisinin (Böl.~\ref{chp:dagilim_aileleri})
ve Merkezi Limit Teoremi'nin (Böl.~\ref{chp:merkezi_limit}) temel araçlarıdır.

\medskip
\textbf{Köprü.} Bölüm~\ref{chp:olcum_teorisi}'deki Lebesgue integrali (MYT, Fatou, DYT)
burada olasılık dili ile yeniden okunur.
Bölüm~\ref{chp:rasgele_degiskenler}'deki PDF ve PMF kavramları, beklenti hesabını
somutlaştırır. Bölüm~\ref{chp:buyuk_sayilar}'da Büyük Sayılar Kanunu bu bölümün
sonuçlarına dayanacaktır.
\end{motivasyon}

% =====================================================================
\section{Beklenti: Lebesgue İntegrali Olarak}
\label{sec:beklenti_tanim}
% =====================================================================

\begin{kesif}
Altı yüzlü zarın beklenen değerini klasik yöntemle hesaplayalım:
$\frac{1}{6}(1+2+3+4+5+6) = 3.5$.
Bu hesap $\sum_x x \cdot P(X=x)$ formülünü izliyor.
Peki ya sürekli dağılımlar? Ya da karma dağılımlar?
Tüm bu durumları kapsayan tek bir tanım mümkün müdür?
\end{kesif}

\begin{tanimgerekce}
Bölüm~\ref{chp:olcum_teorisi}'de basit fonksiyonlarla başlayıp
non-negatif ve ardından genel fonksiyonlara geçen üç aşamalı Lebesgue
integral inşasını hatırlayın. Beklenti, olasılık ölçüsüne göre bu
integralin ta kendisidir — hem kesikli hem sürekli durumu tek çatı altında toplar.
\end{tanimgerekce}

\begin{tanim}[Beklenti]
\label{def:beklenti}
$(\Omega, \mathcal{F}, P)$ üzerinde $X$ rasgele değişken olsun.
$\int_\Omega |X| \, dP < \infty$ ise $X$'in \textbf{beklentisi} (expected value)
\[
  E[X] = \int_\Omega X(\omega) \, dP(\omega)
\]
olarak tanımlanır; $E[X]$ Bölüm~\ref{chp:olcum_teorisi}'deki Lebesgue integralidir.
\end{tanim}

\begin{teorem}[Beklentinin Somut Formülleri]
\label{thm:beklenti_formul}
$X$ rasgele değişken ve $E[|X|] < \infty$ olsun.
\begin{enumerate}[(i)]
  \item \textbf{Kesikli ($S = \{x_1, x_2, \ldots\}$):}
        $E[X] = \sum_k x_k \, P(X = x_k)$.
  \item \textbf{Mutlak sürekli (PDF $f_X$):}
        $E[X] = \int_{-\infty}^{+\infty} x \, f_X(x) \, dx$.
  \item \textbf{Bilinç değişimi (transfer teoremi):}
        $g : \mathbb{R} \to \mathbb{R}$ ölçülebilir ve $E[|g(X)|] < \infty$ ise
        \[
          E[g(X)] = \int_\Omega g(X(\omega)) \, dP(\omega) = \int_{\mathbb{R}} g(x) \, dP_X(x),
        \]
        yani sürekli durumda $E[g(X)] = \int g(x) f_X(x) \, dx$.
\end{enumerate}
\end{teorem}

\begin{proof}
\textit{(i):} $X = \sum_k x_k \mathbf{1}_{\{X = x_k\}}$ ayrışımı; Lebesgue integralinin doğrusallığı.\\
\textit{(ii):} $P_X(B) = \int_B f_X \, d\lambda$ (mutlak süreklilik);
$\int x \, dP_X(x) = \int x f_X(x) \, d\lambda(x)$.\\
\textit{(iii):} $h = g \circ X$ bileşkesi ölçülebilir; ölçü değişim teoremi (push-forward):
$\int_\Omega h \, dP = \int_{\mathbb{R}} g \, dP_X$.
\end{proof}

\begin{ornek}[Geometrik Dağılım]
$X \sim \mathrm{Geom}(p)$: $P(X = k) = (1-p)^{k-1} p$, $k \geq 1$.
\[
  E[X] = \sum_{k=1}^\infty k(1-p)^{k-1} p = p \cdot \frac{d}{dp}\!\left[-\sum_{k=1}^\infty (1-p)^k\right]
  = p \cdot \frac{d}{dp}\!\left[-\frac{1-p}{p}\right] = \frac{1}{p}.
\]
\end{ornek}

\begin{ornek}[Üstel Dağılım]
$X \sim \mathrm{Exp}(\lambda)$: $f_X(x) = \lambda e^{-\lambda x}$, $x > 0$.
\[
  E[X] = \int_0^\infty x \lambda e^{-\lambda x} \, dx
  = \left[-x e^{-\lambda x}\right]_0^\infty + \int_0^\infty e^{-\lambda x} \, dx
  = 0 + \frac{1}{\lambda} = \frac{1}{\lambda}.
\]
\end{ornek}

\begin{hataanalizi}
"Her rasgele değişkenin beklentisi vardır" — yanlış.
Cauchy dağılımı ($f(x) = \frac{1}{\pi(1+x^2)}$) için $\int_{-\infty}^{+\infty} |x| f(x) \, dx = \infty$;
beklenti tanımsızdır. St.\ Petersburg paradoksu (bir madeni para oyununda beklenti sonsuz
ama oyunu oynamak istediğiniz ücret sonludur) da bu olgunun pratik yansımasıdır.
\end{hataanalizi}

\begin{onerme}[Kuyruk Formülü]
\label{prop:kuyruk_beklenti}
$X \geq 0$ rasgele değişken için
\[
  E[X] = \int_0^\infty P(X > t) \, dt.
\]
Daha genel: $E[X^r] = r \int_0^\infty t^{r-1} P(X > t) \, dt$ ($r > 0$).
\end{onerme}

\begin{proof}
Fubini:
\[
  \int_0^\infty P(X > t) \, dt = \int_0^\infty \int_\Omega \mathbf{1}_{X(\omega) > t} \, dP(\omega) \, dt
  = \int_\Omega \int_0^{X(\omega)} dt \, dP(\omega) = \int_\Omega X \, dP = E[X].
\]
\end{proof}

\begin{mikroozet}
$E[X] = \int X \, dP$; kesikli: $\sum x_k p_k$; sürekli: $\int x f(x) \, dx$.
Her rasgele değişkenin beklentisi olmayabilir ($E[|X|] < \infty$ gereklidir).
\end{mikroozet}

% =====================================================================
\section{Beklentinin Temel Özellikleri}
\label{sec:beklenti_ozellik}
% =====================================================================

\begin{teorem}[Beklentinin Temel Özellikleri]
\label{thm:beklenti_ozellik}
$X, Y$ rasgele değişkenler, $E[|X|], E[|Y|] < \infty$, $a, b \in \mathbb{R}$:
\begin{enumerate}[(i)]
  \item \textbf{Doğrusallık:} $E[aX + bY] = aE[X] + bE[Y]$.
  \item \textbf{Monotonluk:} $X \leq Y$ n.k.\ ise $E[X] \leq E[Y]$.
  \item \textbf{$|E[X]| \leq E[|X|]$} (mutlak değer eşitsizliği).
  \item \textbf{Sabit:} $E[c] = c$.
  \item \textbf{Bağımsızlık:} $X, Y$ bağımsız ise $E[XY] = E[X] \cdot E[Y]$
        (her iki beklenti de mevcutsa).
\end{enumerate}
\end{teorem}

\begin{proof}
\textit{(i)}: Lebesgue integralinin doğrusallığı.\\
\textit{(ii)}: $X \leq Y$ n.k.\ $\Rightarrow$ $Y - X \geq 0$ n.k.\ $\Rightarrow$ $E[Y-X] \geq 0$.\\
\textit{(iii)}: $-|X| \leq X \leq |X|$; monotonluktan $|E[X]| \leq E[|X|]$.\\
\textit{(v)}: Sürekli durumda Fubini:
$E[XY] = \iint xy \, f_X(x) f_Y(y) \, dx \, dy = \left(\int x f_X(x) \, dx\right)\!\left(\int y f_Y(y) \, dy\right)
= E[X] \cdot E[Y]$.
\end{proof}

\begin{teorem}[Markov Eşitsizliği]
\label{thm:markov}
$X \geq 0$ rasgele değişken ve $a > 0$ için
\[
  P(X \geq a) \leq \frac{E[X]}{a}.
\]
\end{teorem}

\begin{proof}
$X \geq X \cdot \mathbf{1}_{X \geq a} \geq a \cdot \mathbf{1}_{X \geq a}$.
Beklenti alınca: $E[X] \geq a \cdot P(X \geq a)$. Düzenlenir.
\end{proof}

\begin{teorem}[Chebyshev Eşitsizliği]
\label{thm:chebyshev}
$X$ rasgele değişken, $E[X] = \mu$, $\mathrm{Var}(X) = \sigma^2 < \infty$. Her $k > 0$ için
\[
  P(|X - \mu| \geq k\sigma) \leq \frac{1}{k^2}.
\]
Eşdeğer form: $P(|X - \mu| \geq \varepsilon) \leq \frac{\sigma^2}{\varepsilon^2}$.
\end{teorem}

\begin{proof}
$(X-\mu)^2 \geq 0$ rasgele değişkenine $a = \varepsilon^2$ ile Markov:
$P((X-\mu)^2 \geq \varepsilon^2) \leq E[(X-\mu)^2]/\varepsilon^2 = \sigma^2/\varepsilon^2$.
$P((X-\mu)^2 \geq \varepsilon^2) = P(|X-\mu| \geq \varepsilon)$.
\end{proof}

\begin{teorem}[Jensen Eşitsizliği]
\label{thm:jensen}
$\varphi : \mathbb{R} \to \mathbb{R}$ konveks fonksiyon ve $E[|X|] < \infty$ ise
\[
  \varphi(E[X]) \leq E[\varphi(X)].
\]
\end{teorem}

\begin{proof}
$\varphi$ konveks $\Rightarrow$ her $\mu = E[X]$ için bir $c \in \mathbb{R}$ var:
$\varphi(x) \geq \varphi(\mu) + c(x - \mu)$ (destek doğrusu).
Beklenti alınca: $E[\varphi(X)] \geq \varphi(\mu) + c(E[X] - \mu) = \varphi(\mu) = \varphi(E[X])$.
\end{proof}

\begin{karsiornek}[Jensen Sıkılığı]
$\varphi$ katı konveks ve $X$ sabit değil ise eşitsizlik kesin katıdır:
$\varphi(E[X]) < E[\varphi(X)]$.
Örnek: $\varphi(x) = x^2$, $X$ sabit olmayan $\Rightarrow$ $(E[X])^2 < E[X^2]$,
yani $\mathrm{Var}(X) = E[X^2] - (E[X])^2 > 0$.
\end{karsiornek}

\begin{ornek}[Jensen Uygulamaları]
\begin{itemize}
  \item $\varphi(x) = e^x$ konveks: $e^{E[X]} \leq E[e^X]$.
  \item $\varphi(x) = -\ln x$ konveks ($x>0$): $-\ln E[X] \leq E[-\ln X]$,
        yani $\ln E[X] \geq E[\ln X]$ (AM-GM eşitsizliğinin olasılıksal formu).
  \item $\varphi(x) = |x|^p$ ($p \geq 1$) konveks: $|E[X]|^p \leq E[|X|^p]$ (güç eşitsizliği).
\end{itemize}
\end{ornek}

\begin{mikroozet}
Doğrusallık + Monotonluk + Jensen: beklentinin üç temel aracı.
Markov: $P(X \geq a) \leq E[X]/a$. Chebyshev: $P(|X-\mu| \geq k\sigma) \leq 1/k^2$.
\end{mikroozet}

% =====================================================================
\section{Varyans ve Standart Sapma}
\label{sec:varyans}
% =====================================================================

\begin{tanim}[Varyans ve Standart Sapma]
\label{def:varyans}
$E[X^2] < \infty$ ise $X$'in \textbf{varyansı} (variance):
\[
  \mathrm{Var}(X) = E[(X - E[X])^2] = E[X^2] - (E[X])^2.
\]
\textbf{Standart sapma:} $\sigma_X = \sqrt{\mathrm{Var}(X)}$.
\end{tanim}

\begin{proof}[Hesaplama Formülünün İspat]
$E[(X - \mu)^2] = E[X^2 - 2\mu X + \mu^2]
= E[X^2] - 2\mu E[X] + \mu^2 = E[X^2] - \mu^2$.
\end{proof}

\begin{onerme}[Varyansın Özellikleri]
\label{prop:varyans_ozellik}
\begin{enumerate}[(i)]
  \item $\mathrm{Var}(X) \geq 0$; $\mathrm{Var}(X) = 0 \Leftrightarrow X = \mu$ n.k.
  \item $\mathrm{Var}(aX + b) = a^2 \mathrm{Var}(X)$.
  \item $X, Y$ bağımsız ise $\mathrm{Var}(X + Y) = \mathrm{Var}(X) + \mathrm{Var}(Y)$.
  \item Genel: $\mathrm{Var}(X + Y) = \mathrm{Var}(X) + \mathrm{Var}(Y) + 2\mathrm{Cov}(X, Y)$.
\end{enumerate}
\end{onerme}

\begin{proof}
\textit{(ii)}: $E[(aX+b-aE[X]-b)^2] = a^2 E[(X-E[X])^2]$.\\
\textit{(iii)}: Bağımsız $\Rightarrow$ $\mathrm{Cov}(X,Y) = 0$ (aşağıda ispat edilecek).\\
\textit{(iv)}: $\mathrm{Var}(X+Y) = E[(X+Y-\mu_X-\mu_Y)^2]$; açılımda $E[(X-\mu_X)(Y-\mu_Y)] = \mathrm{Cov}(X,Y)$.
\end{proof}

\begin{ornek}[Binom Dağılımı]
$X \sim \Binom(n, p)$: $X = \sum_{i=1}^n X_i$ ($X_i \sim \mathrm{Bernoulli}(p)$ bağımsız).
$E[X_i] = p$, $\mathrm{Var}(X_i) = p(1-p)$.
Doğrusallık ve bağımsızlıktan: $E[X] = np$, $\mathrm{Var}(X) = np(1-p)$.
\end{ornek}

% =====================================================================
\section{Kovaryans ve Korelasyon}
\label{sec:kovaryans}
% =====================================================================

\begin{tanim}[Kovaryans ve Korelasyon]
\label{def:kovaryans}
$E[X^2], E[Y^2] < \infty$ ise
\begin{align*}
  \mathrm{Cov}(X, Y) &= E[(X - E[X])(Y - E[Y])] = E[XY] - E[X]E[Y], \\
  \mathrm{Cor}(X, Y) &= \rho_{XY} = \frac{\mathrm{Cov}(X,Y)}{\sqrt{\mathrm{Var}(X)\,\mathrm{Var}(Y)}}.
\end{align*}
\end{tanim}

\begin{teorem}[Korelasyonun Sınırları]
\label{thm:cauchy_schwarz}
$-1 \leq \mathrm{Cor}(X, Y) \leq 1$.
$|\rho_{XY}| = 1$ ancak ve ancak $Y = aX + b$ (n.k.) ise ($a \neq 0$).
\end{teorem}

\begin{proof}
\textbf{Cauchy-Schwarz eşitsizliği:} $(\mathrm{Cov}(X,Y))^2 \leq \mathrm{Var}(X) \cdot \mathrm{Var}(Y)$.
İspat: $f(t) = \mathrm{Var}(tX + Y) = t^2 \mathrm{Var}(X) + 2t\,\mathrm{Cov}(X,Y) + \mathrm{Var}(Y) \geq 0$
(diskriminant $\leq 0$): $4(\mathrm{Cov}(X,Y))^2 - 4\mathrm{Var}(X)\mathrm{Var}(Y) \leq 0$.
Eşitlik $f(t^*) = 0$ için; bu $t^*X + Y = 0$ n.k.\ anlamına gelir, yani $Y = -t^*X$ lineer ilişkisi.
\end{proof}

\begin{hataanalizi}
"$\rho = 0$ ise $X, Y$ bağımsızdır" — yanlış.
Bölüm~\ref{sec:rd_bagimsizlik}'daki $X \sim U(-1,1)$, $Y = X^2$ örneği:
$\mathrm{Cov}(X,Y) = E[X^3] - E[X]E[X^2] = 0 - 0 = 0$ ama bağımsız değiller.
Korelasyon yalnızca doğrusal ilişkiyi ölçer.
\end{hataanalizi}

% =====================================================================
\section{Momentler ve Merkezi Momentler}
\label{sec:momentler}
% =====================================================================

\begin{tanim}[Momentler]
\label{def:moment}
$E[|X|^r] < \infty$ ise ($r > 0$):
\begin{itemize}
  \item \textbf{$r$'inci moment:} $\mu_r' = E[X^r]$.
  \item \textbf{$r$'inci merkezi moment:} $\mu_r = E[(X-\mu)^r]$ ($\mu = E[X]$).
  \item \textbf{Standardize $r$'inci moment:} $\tilde{\mu}_r = \mu_r / \sigma^r$.
\end{itemize}
Önemli özel durumlar: $\mu_1' = E[X]$ (birinci moment), $\mu_2 = \mathrm{Var}(X)$.
\end{tanim}

\begin{tanim}[Çarpıklık ve Basıklık]
\label{def:carpiklik}
$\sigma > 0$ ve $E[|X|^4] < \infty$ olmak üzere:
\begin{align*}
  \text{Çarpıklık (skewness):} &\quad \gamma_1 = \tilde{\mu}_3 = \frac{E[(X-\mu)^3]}{\sigma^3}. \\
  \text{Basıklık (kurtosis):} &\quad \gamma_2 = \tilde{\mu}_4 - 3 = \frac{E[(X-\mu)^4]}{\sigma^4} - 3.
\end{align*}
$\gamma_2 = 0$ normal dağılımın referans değeridir (ekses basıklık).
$\gamma_2 > 0$: ağır kuyruk (leptokurtik); $\gamma_2 < 0$: hafif kuyruk (platikurtik).
\end{tanim}

\begin{ornek}[Normal Dağılımın Momentleri]
$X \sim \Normal(\mu, \sigma^2)$:
Simetri $\Rightarrow$ tüm tek merkezi momentler sıfır ($\gamma_1 = 0$).
$E[(X-\mu)^4] = 3\sigma^4$ $\Rightarrow$ $\gamma_2 = 0$.
Genel: $E[(X-\mu)^{2k}] = (2k-1)!! \cdot \sigma^{2k}$ (çift momentler formülü).
\end{ornek}

% =====================================================================
\section{Moment Üreten Fonksiyon}
\label{sec:muf}
% =====================================================================

\begin{kesif}
Bir rasgele değişkenin tüm momentlerini ($E[X], E[X^2], E[X^3], \ldots$) tek bir
fonksiyona kodlayabilir miyiz? Ve bu fonksiyon dağılımı tam olarak belirler mi?
\end{kesif}

\begin{tanim}[Moment Üreten Fonksiyon]
\label{def:muf}
$X$ rasgele değişkeninin \textbf{moment üreten fonksiyonu} (MÜF, moment generating function):
\[
  M_X(t) = E[e^{tX}], \qquad t \in \mathbb{R},
\]
eğer bu beklenti $t$'nin bir $(-h, h)$ ($h > 0$) komşuluğunda sonlu ise.
Bu koşul sağlanmıyorsa MÜF tanımsız sayılır.
\end{tanim}

\begin{teorem}[MÜF'ten Momentler]
\label{thm:muf_moment}
$M_X(t)$ açık bir $(- h, h)$ aralığında sonlu ise:
\begin{enumerate}[(i)]
  \item $E[|X|^r] < \infty$ her $r \geq 1$ için.
  \item $M_X(t) = \sum_{r=0}^\infty \frac{E[X^r]}{r!} t^r$ ($|t| < h$).
  \item $E[X^r] = M_X^{(r)}(0) = \left.\frac{d^r}{dt^r} M_X(t)\right|_{t=0}$.
\end{enumerate}
\end{teorem}

\begin{proof}
\textit{(ii)--(iii):} $e^{tX} = \sum_{r=0}^\infty \frac{(tX)^r}{r!}$.
$|t| < h$ için DCT (Bölüm~\ref{chp:olcum_teorisi}) uygulanabilir;
seri ve integral sıraları değiştirilebilir:
\[
  M_X(t) = E\!\left[\sum_{r=0}^\infty \frac{(tX)^r}{r!}\right]
  = \sum_{r=0}^\infty \frac{E[X^r]}{r!} t^r.
\]
$t = 0$'da $r$'inci türev: $M_X^{(r)}(0) = E[X^r]$.
\end{proof}

\begin{teorem}[MÜF Dağılımı Belirler]
\label{thm:muf_tek}
$M_X(t) = M_Y(t)$ açık bir aralıkta ise $X \stackrel{d}{=} Y$ (dağılımca eşit).
\end{teorem}

Bu teoremin ispatı için moment probleminin çözülebilirliği (Hamburger, Stieltjes)
ve Laplace dönüşümünün tek-türlülüğü gerekmektedir; ayrıntı için Billingsley~\cite{Billingsley1995}
veya Klenke~\cite{Klenke2014}'e bakılabilir.

\begin{teorem}[MÜF ve Bağımsızlık]
\label{thm:muf_bagimsiz}
$X, Y$ bağımsız ise $M_{X+Y}(t) = M_X(t) \cdot M_Y(t)$.
\end{teorem}

\begin{proof}
$M_{X+Y}(t) = E[e^{t(X+Y)}] = E[e^{tX} e^{tY}] = E[e^{tX}] \cdot E[e^{tY}] = M_X(t) M_Y(t)$.
(Bağımsızlık: $E[fg] = E[f]E[g]$.)
\end{proof}

\begin{ornek}[Normal Dağılımın MÜF'ü]
$X \sim \Normal(\mu, \sigma^2)$:
\[
  M_X(t) = \exp\!\left(\mu t + \tfrac{1}{2}\sigma^2 t^2\right), \qquad t \in \mathbb{R}.
\]
\textit{İspat:} $M_X(t) = \int_{-\infty}^{+\infty} e^{tx} \frac{1}{\sigma\sqrt{2\pi}} e^{-(x-\mu)^2/(2\sigma^2)} dx$.
Üstel ifadeler birleştirilir:
$tx - \frac{(x-\mu)^2}{2\sigma^2} = -\frac{(x-\mu-\sigma^2 t)^2}{2\sigma^2} + \mu t + \frac{\sigma^2 t^2}{2}$.
Gauss integrali 1'e tamamlanır; $M_X(t) = e^{\mu t + \sigma^2 t^2/2}$.\\
Bağımsız $X \sim \Normal(\mu_1, \sigma_1^2)$, $Y \sim \Normal(\mu_2, \sigma_2^2)$:
$M_{X+Y}(t) = e^{(\mu_1+\mu_2)t + (\sigma_1^2+\sigma_2^2)t^2/2}$,
yani $X + Y \sim \Normal(\mu_1+\mu_2, \sigma_1^2+\sigma_2^2)$.
\end{ornek}

\begin{hataanalizi}
"Her rasgele değişkenin MÜF'ü vardır" — yanlış.
Cauchy dağılımı için $E[e^{tX}] = \infty$ her $t \neq 0$ için.
Lognormal dağılımı için tüm momentler mevcut ama MÜF sonsuz ($t > 0$).
MÜF yokluğu problemi, karakteristik fonksiyon kullanılarak aşılır (bir sonraki kesim).
\end{hataanalizi}

\begin{mikroozet}
$M_X(t) = E[e^{tX}]$; $E[X^r] = M_X^{(r)}(0)$.
Bağımsız $\Rightarrow$ $M_{X+Y} = M_X \cdot M_Y$.
MÜF her dağılım için var olmayabilir.
\end{mikroozet}

% =====================================================================
\section{Karakteristik Fonksiyon}
\label{sec:kf}
% =====================================================================

\begin{tanimgerekce}
MÜF $E[e^{tX}]$ ($t$ gerçel) her zaman mevcut değildir.
$e^{itX}$ ($i = \sqrt{-1}$, $t$ gerçel) ise $|e^{itX}| = 1$ olduğundan
her zaman integrallenebilir: $E[|e^{itX}|] = 1 < \infty$.
Kompleks üsseli beklenti her dağılım için tanımlıdır.
\end{tanimgerekce}

\begin{tanim}[Karakteristik Fonksiyon]
\label{def:kf}
$X$ rasgele değişkeninin \textbf{karakteristik fonksiyonu} (KF, characteristic function):
\[
  \varphi_X(t) = E[e^{itX}] = E[\cos(tX)] + i \, E[\sin(tX)], \qquad t \in \mathbb{R}.
\]
\end{tanim}

\begin{teorem}[Karakteristik Fonksiyonun Temel Özellikleri]
\label{thm:kf_ozellik}
\begin{enumerate}[(i)]
  \item $\varphi_X(0) = 1$, $|\varphi_X(t)| \leq 1$.
  \item $\varphi_X$ düzgün süreklidir (uniformly continuous).
  \item $\varphi_X$ dağılımı tek türlü belirler (ters dönüşüm teoremi).
  \item $E[|X|^r] < \infty$ ise $\varphi_X^{(r)}(0) = i^r E[X^r]$.
  \item $X, Y$ bağımsız ise $\varphi_{X+Y}(t) = \varphi_X(t) \cdot \varphi_Y(t)$.
\end{enumerate}
\end{teorem}

\begin{proof}
\textit{(i)}: $|e^{itX}| = 1$, dolayısıyla $|\varphi_X(t)| = |E[e^{itX}]| \leq E[|e^{itX}|] = 1$.\\
\textit{(ii)}: $|\varphi_X(t+h) - \varphi_X(t)| = |E[e^{itX}(e^{ihX}-1)]| \leq E[|e^{ihX}-1|]$.
$h \to 0$: $|e^{ihX}-1| \leq 2$, DCT ile $E[|e^{ihX}-1|] \to 0$.\\
\textit{(v)}: MÜF için aynı argüman; $|e^{itX}|=1$ dolayısıyla sonluluk kaygısı yok.
\end{proof}

\begin{teorem}[Ters Dönüşüm]
\label{thm:kf_ters}
$X$ mutlak sürekli, $\int |\varphi_X(t)| \, dt < \infty$ ise
\[
  f_X(x) = \frac{1}{2\pi} \int_{-\infty}^{+\infty} e^{-itx} \varphi_X(t) \, dt.
\]
\end{teorem}

Bu, Fourier dönüşümünün ters formülüdür; gerçel analize (Plancherel teoremi) dayanır.

\begin{ornek}[Normal Dağılımın KF'si]
$X \sim \Normal(\mu, \sigma^2)$:
\[
  \varphi_X(t) = \exp\!\left(i\mu t - \tfrac{1}{2}\sigma^2 t^2\right).
\]
MÜF'te $t \to it$ ikamesi yapar gibi: $M_X(it) = e^{i\mu t - \sigma^2 t^2/2}$.
(Gerçek ispat için analitik sürdürüm veya doğrudan Gauss hesabı gereklidir.)
\end{ornek}

\begin{onerme}[Dağılımla Yakınsama ve KF]
\label{prop:kf_yaklasim}
$X_n \xrightarrow{d} X$ ancak ve ancak her $t \in \mathbb{R}$ için
$\varphi_{X_n}(t) \to \varphi_X(t)$ ise (Lévy süreklilik teoremi).
\end{onerme}

Bu önerme, Bölüm~\ref{chp:merkezi_limit}'teki Merkezi Limit Teoremi'nin kanıtında
kilit rol oynayacaktır.

\begin{mikroozet}
$\varphi_X(t) = E[e^{itX}]$: her dağılım için mevcut, düzgün sürekli, $|\varphi| \leq 1$.
Bağımsız toplam: $\varphi_{X+Y} = \varphi_X \cdot \varphi_Y$.
Dağılımla yakınsama $\Leftrightarrow$ KF'lerin noktasal yakınsaması (Lévy).
\end{mikroozet}

% =====================================================================
\section{Kümülan Üreten Fonksiyon}
\label{sec:kgf}
% =====================================================================

\begin{tanim}[Kümülan Üreten Fonksiyon]
\label{def:kgf}
$M_X(t)$ mevcut ise \textbf{kümülan üreten fonksiyon} (cumulant generating function):
\[
  K_X(t) = \ln M_X(t) = \sum_{r=1}^\infty \frac{\kappa_r}{r!} t^r.
\]
$\kappa_r = K_X^{(r)}(0)$: $r$'inci \textbf{kümülan} (cumulant).
\end{tanim}

\begin{onerme}[Önemli Kümanlar]
\label{prop:kumulanlar}
\begin{align*}
  \kappa_1 &= E[X] = \mu, \\
  \kappa_2 &= \mathrm{Var}(X) = \sigma^2, \\
  \kappa_3 &= E[(X-\mu)^3] \quad \text{(çarpıklıkla orantılı)}, \\
  \kappa_4 &= E[(X-\mu)^4] - 3\sigma^4 \quad \text{(ekses basıklıkla orantılı)}.
\end{align*}
Bağımsız $X, Y$ için $K_{X+Y}(t) = K_X(t) + K_Y(t)$, dolayısıyla
kümanlar toplanır: $\kappa_r(X+Y) = \kappa_r(X) + \kappa_r(Y)$.
\end{onerme}

\begin{proof}
$K_X(t) = \ln M_X(t)$; $K_X'(0) = M_X'(0)/M_X(0) = E[X]/1 = \mu$.
$K_X''(0) = E[X^2] - (E[X])^2 = \mathrm{Var}(X)$.
Bağımsız toplam: $\ln M_{X+Y} = \ln M_X + \ln M_Y$.
\end{proof}

\begin{ornek}[Normal Dağılımın Kümanları]
$X \sim \Normal(\mu, \sigma^2)$: $K_X(t) = \mu t + \sigma^2 t^2/2$.
Sadece $\kappa_1 = \mu$ ve $\kappa_2 = \sigma^2$ sıfırdan farklı;
$r \geq 3$ için $\kappa_r = 0$.
Bu, normal dağılımın kümülan diliyle karakterizasyonudur.
\end{ornek}

% =====================================================================
\section{Koşullu Beklenti}
\label{sec:kosullu_beklenti}
% =====================================================================

\begin{kesif}
$Y = 2X + \varepsilon$ ($\varepsilon \perp X$) modelini düşünün.
$X$'in değerini bildiğimizde $Y$ için en iyi tahminimiz nedir?
Sezgisel olarak $E[Y \mid X = x] = 2x$ diyoruz.
Peki ya $X$ yerine bir olay $\mathcal{G}$ verildiğinde?
"En iyi tahmin" kavramını nasıl soyutlarız?
\end{kesif}

\begin{tanim}[Koşullu Beklenti — Elemanter]
\label{def:kosullu_beklenti_elem}
\textbf{Kesikli:} $P(Y = y_j) > 0$ olayları üretilen $\mathcal{G} = \sigma(Y)$ için
\[
  E[X \mid Y = y_j] = \frac{E[X \cdot \mathbf{1}_{\{Y=y_j\}}]}{P(Y = y_j)}.
\]
\textbf{Sürekli:} $f_Y(y) > 0$ ise $E[X \mid Y = y] = \int x \, f_{X|Y}(x \mid y) \, dx$.
\end{tanim}

\begin{tanimgerekce}
Genel $\sigma$-cebiri için yukarıdaki formüller yetersiz kalır (sıfır olasılıklı koşullamada anlam kaybolur).
Radon-Nikodym teoremi, koşullu beklentiyi "en iyi $\mathcal{G}$-ölçülebilir tahmin"
olarak $L^2$-projeksiyon ile tanımlamamızı sağlar.
\end{tanimgerekce}

\begin{tanim}[Koşullu Beklenti — Genel]
\label{def:kosullu_beklenti}
$(\Omega, \mathcal{F}, P)$ üzerinde $X \in L^1$ ve $\mathcal{G} \subseteq \mathcal{F}$
alt $\sigma$-cebiri verilsin. $E[X \mid \mathcal{G}]$, şu iki koşulu sağlayan
(n.k.\ tek türlü) rasgele değişkendir:
\begin{enumerate}[(i)]
  \item $E[X \mid \mathcal{G}]$ \textbf{$\mathcal{G}$-ölçülebilir}.
  \item Her $G \in \mathcal{G}$ için $\int_G E[X \mid \mathcal{G}] \, dP = \int_G X \, dP$.
\end{enumerate}
Özel olarak $\mathcal{G} = \sigma(Y)$ için $E[X \mid Y] := E[X \mid \sigma(Y)]$.
\end{tanim}

\begin{proof}[Varlık ve Tek Türlülük]
Tek türlülük: $Z_1, Z_2$ her ikisi de koşulları sağlasın; $G = \{Z_1 > Z_2\} \in \mathcal{G}$:
$\int_G (Z_1 - Z_2) \, dP = 0$ ve $Z_1 - Z_2 > 0$ burada; yani $P(G) = 0$.
Benzer şekilde $\{Z_1 < Z_2\}$ için; dolayısıyla $Z_1 = Z_2$ n.k.\\
Varlık: $\nu(G) = \int_G X \, dP$ ($G \in \mathcal{G}$) $\sigma$-sonlu imzalı ölçü;
$P|_{\mathcal{G}}$'ye mutlak sürekli. Radon-Nikodym: $\nu = Z \cdot P|_{\mathcal{G}}$
yazan $\mathcal{G}$-ölçülebilir $Z$ var; $Z = E[X \mid \mathcal{G}]$.
\end{proof}

\begin{teorem}[Koşullu Beklentinin Temel Özellikleri]
\label{thm:kosullu_beklenti}
$X, Y \in L^1$, $\mathcal{G} \subseteq \mathcal{H} \subseteq \mathcal{F}$ alt $\sigma$-cebirleri.
\begin{enumerate}[(i)]
  \item \textbf{Doğrusallık:} $E[aX + bY \mid \mathcal{G}] = aE[X\mid\mathcal{G}] + bE[Y\mid\mathcal{G}]$.
  \item \textbf{Kule (yineleme) özelliği:} $E[E[X \mid \mathcal{H}] \mid \mathcal{G}] = E[X \mid \mathcal{G}]$.
        Özellikle $E[E[X \mid Y]] = E[X]$ (toplam beklenti formülü).
  \item \textbf{$\mathcal{G}$-ölçülebilir çarpan:} $Z$ $\mathcal{G}$-ölçülebilir ise
        $E[ZX \mid \mathcal{G}] = Z \cdot E[X \mid \mathcal{G}]$.
  \item \textbf{Bağımsızlık:} $X$ $\mathcal{G}$'den bağımsız ise $E[X \mid \mathcal{G}] = E[X]$.
  \item \textbf{Monotonluk:} $X \leq Y$ n.k.\ ise $E[X \mid \mathcal{G}] \leq E[Y \mid \mathcal{G}]$ n.k.
  \item \textbf{Jensen:} $\varphi$ konveks ise $\varphi(E[X \mid \mathcal{G}]) \leq E[\varphi(X) \mid \mathcal{G}]$.
\end{enumerate}
\end{teorem}

\begin{proof}
\textit{(ii) Kule özelliği:} $\mathcal{G} \subseteq \mathcal{H}$.
Her $G \in \mathcal{G} \subseteq \mathcal{H}$ için tanım koşulu (ii):
$\int_G E[X\mid\mathcal{H}] \, dP = \int_G X \, dP$.
Ama $E[X\mid\mathcal{H}]$'nın $\mathcal{G}$-koşullu beklentisi için de bu değer kullanılabilir;
dolayısıyla $E[E[X\mid\mathcal{H}]\mid\mathcal{G}]$ tanım koşullarını $E[X\mid\mathcal{G}]$ ile aynı şekilde sağlar.
Tek türlülükten eşitlik.\\
\textit{(iii):} $G \in \mathcal{G}$: $\int_G ZE[X\mid\mathcal{G}] \, dP = \int_G Z \cdot E[X\mid\mathcal{G}] \, dP$.
$Z$ basamak fonksiyonu ise doğrudan; genel $Z$ için yaklaşım.
\end{proof}

\begin{onerme}[$L^2$ Projeksiyon Yorumu]
\label{prop:l2_projeksiyon}
$X \in L^2(\Omega, \mathcal{F}, P)$ için $E[X \mid \mathcal{G}]$,
$L^2(\Omega, \mathcal{G}, P)$ alt uzayına $X$'in dik projeksiyonudur:
\[
  E[X \mid \mathcal{G}] = \arg\min_{Z \in L^2(\mathcal{G})} E[(X - Z)^2].
\]
\end{onerme}

\begin{proof}
$Z \in L^2(\mathcal{G})$ için:
$E[(X-Z)^2] = E[(X - E[X\mid\mathcal{G}] + E[X\mid\mathcal{G}] - Z)^2]$.
Açılımda çapraz terim $E[(X - E[X\mid\mathcal{G}])(E[X\mid\mathcal{G}] - Z)]$:
$(E[X\mid\mathcal{G}] - Z)$ $\mathcal{G}$-ölçülebilir; (iii) özelliğiyle çapraz terim sıfır.
Dolayısıyla $E[(X-Z)^2] = E[(X-E[X\mid\mathcal{G}])^2] + E[(E[X\mid\mathcal{G}]-Z)^2]$,
sağ tarafın minimumu $Z = E[X \mid \mathcal{G}]$'de sağlanır.
\end{proof}

\begin{ornek}[Regresyon]
$E[Y \mid X] = g(X)$; burada $g(x) = E[Y \mid X = x]$.
$E[Y \mid X]$ koşullu beklentisi, $Y$'yi $X$ bilgisiyle tahmin etmede
en küçük ortalama karesel hata veren tahminci (MMSE tahmin edici) olup
"doğrusal olmayan regresyon" kavramıyla özdeşleşir.
Doğrusal kısıt altında ($g(X) = a + bX$): $b = \mathrm{Cov}(X,Y)/\mathrm{Var}(X)$.
\end{ornek}

\begin{kendinleifade}
Koşullu beklentinin "bilgiyi güncellemek" yorumunu kendi cümlelerinizle açıklayın.
"Kule özelliği" ne anlama gelir? Somut bir örnek üretin.
\ogrencialani[2.5cm]
\end{kendinleifade}

% ---- Disiplinlerarası Bağlantı (TYMM) ----
\begin{kopru}[Disiplinlerarası — Ekonomi, Fizik, Bilgisayar Bilimi]
\textbf{Ekonomi:} Beklenen fayda teorisi (von Neumann-Morgenstern), karar teorisinin
matematiksel temelidir; Jensen eşitsizliği risk kaçınımını açıklar ($u$ içbükey $\Rightarrow$ $u(E[X]) > E[u(X)]$).\\
\textbf{Makine öğrenmesi:} Koşullu beklenti $E[Y|X]$, gözetimli öğrenmede kayıp fonksiyonu
minimize eden tahmin edicinin kendisidir; nöral ağlar bu fonksiyonu yaklaşık öğrenir.\\
\textbf{Kuantum mekaniği:} Observable $A$ operatörünün beklenti değeri $\langle A \rangle = \langle \psi | A | \psi \rangle$,
istatistiksel beklentiyle yapısal olarak örtüşür.
\end{kopru}

% ---- Öz-Yansıtma ----
\begin{ozyansima}
\begin{itemize}[leftmargin=1.8em]
  \item Beklentinin neden Lebesgue integrali olarak tanımlanmasının gerekli olduğunu açıklayabilir misiniz?
  \item Jensen eşitsizliği sizce en güçlü hangi uygulamada karşınıza çıktı?
  \item MÜF ile KF arasındaki temel fark nedir, ve hangisi daha "evrensel"dir?
  \item Koşullu beklentiyi $L^2$ projeksiyonu olarak görmek, sezginizi nasıl değiştirdi?
\end{itemize}
\end{ozyansima}

% ---- Köprü Notu ----
\begin{kopru}
\textbf{Önceki bölüm:} Bölüm~\ref{chp:rasgele_degiskenler} — dağılım fonksiyonu ve PDF;
burada bunlar Lebesgue integrali ile beklentiye dönüştü.\\
\textbf{Sonraki bölüm:} Bölüm~\ref{chp:dagilim_aileleri} — somut dağılımların
beklenti, varyans, MÜF hesapları bu bölümün araçlarıyla yapılacak.\\
\textbf{İleride:} Bölüm~\ref{chp:buyuk_sayilar}'da Chebyshev ile ZBSK,
Bölüm~\ref{chp:merkezi_limit}'te KF ile MLT kanıtı.
\defterbaglantisi{bolum04\_beklenti.ipynb}{%
  MÜF hesabı, Jensen görselleştirmesi, koşullu beklenti simülasyonu}
\end{kopru}

% =====================================================================
\section{Olasılık Eşitsizlikleri}
\label{sec:olasilik_esitsizlikleri}
% =====================================================================

\begin{motivasyon}[Neden Eşitsizlikler?]
Bir dağılımı tam olarak bilmeden $P(|X - \mu| > \varepsilon)$ hakkında ne söyleyebiliriz?
Bu soru yalnızca momentlere dayalı sınırlarla yanıtlanabilir.
Markov, Chebyshev ve Jensen eşitsizlikleri bu sınırların sistematik kaynağıdır;
Büyük Sayılar Kanunu ve Büyük Sapmalar teorisinin altyapısını oluştururlar.
\end{motivasyon}

\begin{teorem}[Markov Eşitsizliği]
\label{thm:markov_beklenti}
$X \geq 0$ ve $E[X] < \infty$. Her $a > 0$ için
\[
  P(X \geq a) \leq \frac{E[X]}{a}.
\]
\end{teorem}

\begin{proof}
$X \geq a\,\mathbf{1}_{\{X \geq a\}}$ (çünkü $X \geq 0$).
Beklenti alınca: $E[X] \geq a\,P(X \geq a)$. Düzenlenince sonuç.
\end{proof}

\begin{teorem}[Chebyshev Eşitsizliği — Genelleştirilmiş]
\label{thm:chebyshev_genellestirilmis}
$\varphi : \mathbb{R} \to [0, \infty)$ simetrik ve artan ($|x|$'de) fonksiyon,
$\varphi(a) > 0$. O zaman
\[
  P(|X - \mu| \geq a) \leq \frac{E[\varphi(X-\mu)]}{\varphi(a)}.
\]
Özel durum $\varphi(x) = x^2$: $P(|X-\mu|\geq a) \leq \mathrm{Var}(X)/a^2$.
\end{teorem}

\begin{proof}
$\varphi(|X-\mu|) \geq \varphi(a)\,\mathbf{1}_{\{|X-\mu| \geq a\}}$ ($\varphi$ artan).
Markov: $E[\varphi(|X-\mu|)] \geq \varphi(a)\,P(|X-\mu|\geq a)$.
\end{proof}

\begin{teorem}[Tek Taraflı Chebyshev (Cantelli Eşitsizliği)]
\label{thm:cantelli}
$E[X] = \mu$, $\mathrm{Var}(X) = \sigma^2 < \infty$. Her $\lambda > 0$ için
\[
  P(X - \mu \geq \lambda) \leq \frac{\sigma^2}{\sigma^2 + \lambda^2}.
\]
\end{teorem}

\begin{proof}
Her $t > 0$ için $P(X-\mu \geq \lambda) = P(X-\mu+t \geq \lambda+t)$.
Markov ($X-\mu+t \geq 0$ değil; bunun yerine $(X-\mu+t)^2$ ile):
$P((X-\mu+t)^2 \geq (\lambda+t)^2) \leq E[(X-\mu+t)^2]/(\lambda+t)^2 = (\sigma^2+t^2)/(\lambda+t)^2$.
$t = \sigma^2/\lambda$'da minimize et: $(\sigma^2 + \sigma^4/\lambda^2)/(\lambda + \sigma^2/\lambda)^2 = \sigma^2/(\sigma^2+\lambda^2)$.
\end{proof}

\begin{ornek}[Eşitsizliklerin Karşılaştırması]
$X \sim \Normal(0,1)$, $a = 3$.
\begin{itemize}
  \item Markov ($|X|$ ile): $P(|X| \geq 3) \leq E[|X|]/3 = \sqrt{2/\pi}/3 \approx 0.266$.
  \item Chebyshev: $P(|X| \geq 3) \leq 1/9 \approx 0.111$.
  \item Cantelli (tek taraf): $P(X \geq 3) \leq 1/(1+9) = 0.1$.
  \item Gerçek değer: $P(|X| \geq 3) = 2(1-\Phi(3)) \approx 0.0027$.
\end{itemize}
Eşitsizlikler "kötü" ama dağılım bilgisi olmadan geçerlidir.
\end{ornek}

\begin{teorem}[Hoeffding Eşitsizliği]
\label{thm:hoeffding}
$X_1, \ldots, X_n$ bağımsız, $a_i \leq X_i \leq b_i$ n.k., $E[X_i] = \mu_i$.
Her $t > 0$ için
\[
  P\!\left(\sum_{i=1}^n (X_i - \mu_i) \geq t\right) \leq \exp\!\left(-\frac{2t^2}{\sum_{i=1}^n(b_i-a_i)^2}\right).
\]
\end{teorem}

\begin{proof}[İspat Taslağı]
Hoeffding lemması: $E[X_i] = 0$, $a_i \leq X_i \leq b_i$ ise
$E[e^{sX_i}] \leq e^{s^2(b_i-a_i)^2/8}$.
(Konvekslik: $e^{sx} \leq \frac{b-x}{b-a}e^{sa} + \frac{x-a}{b-a}e^{sb}$;
beklenti al, $f(u) = -su + e^u$'nun Taylor açılımıyla sınırla.)
Markov + MÜF çarpımı: $P(\sum(X_i-\mu_i)\geq t) \leq e^{-st}\prod e^{s^2(b_i-a_i)^2/8}$.
$s = 4t/\sum(b_i-a_i)^2$'de minimize et: istenen sınır.
\end{proof}

\begin{ornek}[Makine Öğrenmesi — Genelleme Hatası]
$n$ bağımsız veri; her $X_i \in [0,1]$, $E[X_i] = \mu$.
Hoeffding: $P(\bar{X}_n - \mu \geq \varepsilon) \leq e^{-2n\varepsilon^2}$.
$n = 1000$, $\varepsilon = 0.05$: $P \leq e^{-5} \approx 0.007$.
PAC öğrenmesinde bu tür sınırlar "garantili öğrenme" teorisinin temelidir.
\end{ornek}

\begin{mikroozet}
Markov: $P(X \geq a) \leq E[X]/a$ ($X \geq 0$).
Chebyshev: $P(|X-\mu|\geq a) \leq \sigma^2/a^2$.
Cantelli: tek taraflı, $\sigma^2/(\sigma^2+\lambda^2)$.
Hoeffding: sınırlı değişkenler için üstel sınır.
\end{mikroozet}

% =====================================================================
\section{Değişim Katsayısı ve Standartlaştırma}
\label{sec:degisim_katsayisi}
% =====================================================================

\begin{tanim}[Değişim Katsayısı]
\label{def:degisim_katsayisi}
$\mu = E[X] \neq 0$ ve $\sigma = \sqrt{\mathrm{Var}(X)}$ için
\textbf{değişim katsayısı} (coefficient of variation):
\[
  \mathrm{CV}(X) = \frac{\sigma}{|\mu|}.
\]
Boyutsuz göreli değişkenlik ölçüsü; farklı birim veya ölçekteki verileri karşılaştırır.
\end{tanim}

\begin{ornek}[CV Karşılaştırması]
İki yatırım:
$X$: $\mu_X = 100$, $\sigma_X = 15$ $\Rightarrow$ $\mathrm{CV}(X) = 0.15$.
$Y$: $\mu_Y = 1000$, $\sigma_Y = 120$ $\Rightarrow$ $\mathrm{CV}(Y) = 0.12$.
Mutlak değişkenlik ($\sigma$) bakımından $Y$ daha değişken;
göreli değişkenlik (CV) bakımından $X$ daha riskli.
\end{ornek}

\begin{tanim}[Çarpıklık ve Basıklık]
\label{def:carpiklik_basiklik}
$\mu_k' = E[(X-\mu)^k]$ $k$'ıncı merkezi moment.
\[
  \text{Çarpıklık (skewness):} \quad \gamma_1 = \frac{\mu_3'}{\sigma^3}.
\]
\[
  \text{Basıklık (kurtosis):} \quad \gamma_2 = \frac{\mu_4'}{\sigma^4} - 3.
\]
$\gamma_2 = 0$: Normal dağılım. $\gamma_2 > 0$: ağır kuyruk (leptokurtik).
$\gamma_2 < 0$: hafif kuyruk (platikurtik).
\end{tanim}

\begin{teorem}[Normal Dağılımın Momentleri]
\label{thm:normal_momentler}
$X \sim \Normal(0,1)$: tüm tek momentler sıfır;
$E[X^{2k}] = (2k-1)!! = (2k)!/(2^k k!)$.
Özellikle: $E[X^4] = 3$ $\Rightarrow$ $\gamma_2 = 3 - 3 = 0$.
\end{teorem}

\begin{proof}
$E[X^{2k}] = \int_{-\infty}^\infty x^{2k} \phi(x) dx$.
Parçalı integral: $u = x^{2k-1}$, $dv = x\phi(x)dx = -\phi'(x)dx$; $v = -\phi(x)$.
$\int x^{2k}\phi(x)dx = [-x^{2k-1}\phi(x)]_{-\infty}^\infty + (2k-1)\int x^{2k-2}\phi(x)dx = (2k-1)E[X^{2k-2}]$.
Yinelemeli: $E[X^{2k}] = (2k-1)(2k-3)\cdots 1 = (2k-1)!!$.
\end{proof}

\begin{ornek}[Çarpıklık Yorumu]
$X \sim \Exp(\lambda)$: $\mu_3' = 2/\lambda^3$, $\sigma^3 = 1/\lambda^3$, $\gamma_1 = 2 > 0$: sağa çarpık.
$X \sim \Beta(2,5)$: $\gamma_1 = 2(5-2)\sqrt{(2+5+1)/((2\cdot5)\cdot 1)} / (2+5+2) > 0$: sağa çarpık.
Gelir dağılımları tipik olarak pozitif çarpık: kalabalık düşük-orta, seyrek yüksek gelirler.
\end{ornek}

% =====================================================================
\section{L$^p$ Uzayları ve Beklenti Eşitsizlikleri}
\label{sec:lp_uzaylari}
% =====================================================================

\begin{tanim}[$L^p$ normu]\label{def:lp_norm}
$p\geq1$. Rasgele değişken $X$ için
\[
  \|X\|_p = (\Beklenti[|X|^p])^{1/p}.
\]
$L^p(\Omega,\mathcal{F},P)=\{X:\|X\|_p<\infty\}$: $p$'inci mutlak moment sonlu olan rasgele değişkenler sınıfı.
\end{tanim}

\begin{teorem}[Lyapunov eşitsizliği]\label{thm:lyapunov_esit}
$1\leq r<s$: $\|X\|_r\leq\|X\|_s$.
Yani $L^s\subset L^r$ ve $r < s$ için $r$'inci moment $s$'inci momentten daha sık sonludur.
\end{teorem}

\begin{proof}
Jensen: $\phi(x)=x^{s/r}$ dışbükey ($s>r$). $\Beklenti[|X|^r]^{s/r}\leq\Beklenti[|X|^s]$. Her iki tarafın $r/s$'inci kuvveti: $\Beklenti[|X|^r]^{1}\leq\Beklenti[|X|^s]^{r/s}$, yani $\|X\|_r^r\leq\|X\|_s^r$, bu da $\|X\|_r\leq\|X\|_s$.
\end{proof}

\begin{teorem}[Hölder eşitsizliği]\label{thm:holder}
$1/p+1/q=1$, $p,q\geq1$:
\[
  \Beklenti[|XY|] \leq \|X\|_p\cdot\|Y\|_q.
\]
Özel durum $p=q=2$: Cauchy-Schwarz $\Beklenti[|XY|]\leq\sqrt{\Beklenti[X^2]\Beklenti[Y^2]}$.
\end{teorem}

\begin{proof}[Proof (Young ile)]
Young eşitsizliği: $ab\leq a^p/p+b^q/q$. $a=|X|/\|X\|_p$, $b=|Y|/\|Y\|_q$: $\frac{|XY|}{\|X\|_p\|Y\|_q}\leq\frac{|X|^p}{p\|X\|_p^p}+\frac{|Y|^q}{q\|Y\|_q^q}$. Beklenti alınca sağ taraf $1/p+1/q=1$.
\end{proof}

\begin{teorem}[Minkowski eşitsizliği]\label{thm:minkowski}
$p\geq1$: $\|X+Y\|_p\leq\|X\|_p+\|Y\|_p$ (üçgen eşitsizliği).
\end{teorem}

\begin{ornek}[$L^p$ dahil olma ilişkisi]\label{ex:lp_dahil}
Sonlu $\Omega=\{1,\ldots,n\}$ (her eleman eşit olasılıklı). $\|X\|_p=(\frac1n\sum|x_i|^p)^{1/p}$: RMS$(p=2)>$ortalama$(p=1)$ — karesel ortalamanın aritmetik ortalamayı aşması. $\|X\|_\infty=\max|x_i|$ = supremum norm.
\end{ornek}

% =====================================================================
\section{Koşullu Beklentinin Geometrisi}
\label{sec:kb_geometri}
% =====================================================================

\begin{tanim}[Projeksiyon yorumu]\label{def:projeksiyon}
$L^2(\Omega,\mathcal{F},P)$ Hilbert uzayı, iç çarpım $\langle X,Y\rangle=\Beklenti[XY]$. $\Beklenti[X|\mathcal{G}]$ (alt $\sigma$-cebir $\mathcal{G}$ üzerine koşullu beklenti), $X$'in $L^2(\Omega,\mathcal{G},P)$ altuzayına \emph{ortogonal projeksiyonudur}:
\[
  \Beklenti[\Beklenti[X|\mathcal{G}]Y] = \Beklenti[XY] \quad \forall Y\in L^2(\Omega,\mathcal{G},P).
\]
\end{tanim}

\begin{teorem}[Projeksiyon karakterizasyonu]\label{thm:projeksiyon}
$Z=\Beklenti[X|\mathcal{G}]$ iki koşulu sağlayan tek $\mathcal{G}$-ölçülebilir rasgele değişkendir:
\begin{enumerate}[(i)]
  \item $\Beklenti[ZY]=\Beklenti[XY]$ her $\mathcal{G}$-ölçülebilir $Y\in L^2$ için.
  \item $Z$, $\Beklenti[X-Z)^2]$'yi minimize eden — MMSE tahmincisi.
\end{enumerate}
\end{teorem}

\begin{proof}[Proof (taslak)]
(ii): $\Beklenti[(X-Z')^2]\geq\Beklenti[(X-Z)^2]$ her $\mathcal{G}$-ölçülebilir $Z'$ için. Genişlet: $\Beklenti[(X-Z')^2]=\Beklenti[(X-Z)^2]+\Beklenti[(Z-Z')^2]\geq\Beklenti[(X-Z)^2]$; orta terim sıfır (ortogonallik).
\end{proof}

\begin{ornek}[Lineer MMSE]\label{ex:lineer_mmse}
$(X,Y)$ ortak $L^2$. Lineer tahminci $\hat X = aY+b$ üzerinde MMSE:
$a^*=\Cov(X,Y)/\Var(Y)$, $b^*=E[X]-a^*E[Y]$.
$(X,Y)$ birlikte normal ise bu tahmin edici optimal (doğrusal olmayan gerekmez) ve $\hat X=\Beklenti[X|Y]$.
\end{ornek}

\begin{disiplinlerarasibaglan}
\textbf{Sinyal İşleme:} Wiener filtresi, gürültülü sinyalden MMSE tahmincisidir — koşullu beklentinin projeksiyon yorumunun mühendislik uygulaması.

\textbf{Ekonometri:} OLS (En Küçük Kareler) tahmincisi, $Y$'yi doğrusal $\mathbf{x}$ fonksiyonlarına projeksiyon olarak yorumlanır; $E[Y|\mathbf{X}=\mathbf{x}]$ doğrusal ise OLS MMSE elde eder.
\end{disiplinlerarasibaglan}

\begin{mikroozet}
CV: boyutsuz göreli değişkenlik; $\sigma/|\mu|$.
Çarpıklık: $\mu_3'/\sigma^3$; yönsel asimetri.
Basıklık fazlası: $\mu_4'/\sigma^4 - 3$; Normal'e göre kuyruk ağırlığı.
$L^p$ uzayları: Lyapunov ($\|X\|_r\leq\|X\|_s$, $r<s$), Hölder, Minkowski.
Koşullu beklenti = ortogonal projeksiyon $\Rightarrow$ MMSE tahmincisi.
\end{mikroozet}

% =====================================================================
%  ALIŞTIRMALAR
% =====================================================================

\cozumlualistirmalar

\begin{alistirma}[Al.4.1]{A}{Beklenti Hesabı}
$X$ PMF'si $p(k) = (1/2)^k$, $k = 1, 2, 3, \ldots$ olan kesikli rasgele değişken.
$E[X]$ ve $E[X^2]$'yi hesaplayın. Ardından $\mathrm{Var}(X)$'i bulun.
\end{alistirma}
\alcozum{%
$E[X] = \sum_{k=1}^\infty k (1/2)^k$. $\sum_{k=1}^\infty k r^k = r/(1-r)^2$ ($|r|<1$) ile $r=1/2$:
$E[X] = (1/2)/(1/2)^2 = 2$.\\
$E[X^2] = \sum_{k=1}^\infty k^2 (1/2)^k$. $\sum k^2 r^k = r(1+r)/(1-r)^3$: $E[X^2] = (1/2)(3/2)/(1/2)^3 = 6$.\\
$\mathrm{Var}(X) = 6 - 4 = 2$.}

\begin{alistirma}[Al.4.2]{A}{Jensen Eşitsizliği}
$X > 0$ rasgele değişken, $E[X] = 2$.
(a) $E[\ln X]$'i en fazla kaç yapabilirsiniz? Üst sınırı bulun.
(b) $E[1/X]$'in alt sınırını Jensen ile bulun.
\end{alistirma}
\alcozum{%
(a) $\varphi(x) = \ln x$ içbükey ($\varphi'' < 0$), yani $-\ln x$ konveks:
$-\ln E[X] \leq E[-\ln X]$, yani $E[\ln X] \leq \ln E[X] = \ln 2$.\\
(b) $\varphi(x) = 1/x$ konveks ($x>0$): Jensen $\Rightarrow$ $1/E[X] \leq E[1/X]$, yani $E[1/X] \geq 1/2$.}

\begin{alistirma}[Al.4.3]{B}{MÜF ile Moment Bulma}
$X \sim \mathrm{Pois}(\lambda)$: PMF $p(k) = e^{-\lambda}\lambda^k/k!$, $k=0,1,2,\ldots$
(a) $M_X(t) = e^{\lambda(e^t-1)}$ olduğunu gösterin.
(b) $M_X'(0)$ ve $M_X''(0)$ ile $E[X]$ ve $E[X^2]$'yi bulun.
(c) $\mathrm{Var}(X)$'i hesaplayın.
\end{alistirma}
\alcozum{%
(a) $M_X(t) = \sum_{k=0}^\infty e^{tk} e^{-\lambda} \lambda^k/k! = e^{-\lambda} \sum_{k=0}^\infty (\lambda e^t)^k/k!
= e^{-\lambda} e^{\lambda e^t} = e^{\lambda(e^t-1)}$.\\
(b) $M_X'(t) = \lambda e^t \cdot e^{\lambda(e^t-1)}$; $M_X'(0) = \lambda \cdot 1 = \lambda = E[X]$.
$M_X''(t) = (\lambda e^t + \lambda^2 e^{2t}) e^{\lambda(e^t-1)}$;
$M_X''(0) = \lambda + \lambda^2 = E[X^2]$.\\
(c) $\mathrm{Var}(X) = E[X^2] - (E[X])^2 = \lambda + \lambda^2 - \lambda^2 = \lambda$.
Poisson: beklenti = varyans = $\lambda$.}

\begin{alistirma}[Al.4.4]{B}{Koşullu Beklenti}
$X \sim \mathrm{Uniform}(0,1)$, $Y \mid X = x \sim \mathrm{Uniform}(0, x)$.
(a) $E[Y \mid X = x]$'i bulun.
(b) $E[Y]$'yi kule özelliğiyle hesaplayın.
(c) $\mathrm{Var}(Y)$'yi $\mathrm{Var}(Y) = E[\mathrm{Var}(Y\mid X)] + \mathrm{Var}(E[Y\mid X])$ ile bulun.
\end{alistirma}
\alcozum{%
(a) $E[Y \mid X=x] = x/2$.\\
(b) $E[Y] = E[E[Y\mid X]] = E[X/2] = (1/2) \cdot 1/2 = 1/4$.\\
(c) $\mathrm{Var}(Y\mid X=x) = x^2/12$.
$E[\mathrm{Var}(Y\mid X)] = E[X^2/12] = (1/3)/12 = 1/36$.
$\mathrm{Var}(E[Y\mid X]) = \mathrm{Var}(X/2) = (1/4)\mathrm{Var}(X) = (1/4)(1/12) = 1/48$.
$\mathrm{Var}(Y) = 1/36 + 1/48 = 4/144 + 3/144 = 7/144$.}

\begin{alistirma}[Al.4.4b]{B}{Karakteristik Fonksiyon ve Bağımsızlık}
$X,Y$ bağımsız rasgele değişkenler; $\varphi_X(t)=E[e^{itX}]$, $\varphi_Y(t)=E[e^{itY}]$.
(a) $\varphi_{X+Y}(t)=\varphi_X(t)\varphi_Y(t)$ olduğunu gösterin.
(b) $X,Y\iid\Normal(0,1)$: $\varphi_X(t)=e^{-t^2/2}$ kullanarak $Z=X+Y$ dağılımını bulun.
(c) $X\sim\Pois(\lambda_1)$, $Y\sim\Pois(\lambda_2)$ bağımsız: $Z=X+Y$ dağılımını KF yoluyla belirleyin.
\end{alistirma}
\alcozum{%
(a) $\varphi_{X+Y}(t)=E[e^{it(X+Y)}]=E[e^{itX}e^{itY}]=E[e^{itX}]E[e^{itY}]=\varphi_X(t)\varphi_Y(t)$.
(Bağımsızlık beklentinin çarpımını garanti eder.)

(b) $\varphi_Z(t)=\varphi_X(t)\varphi_Y(t)=e^{-t^2/2}\cdot e^{-t^2/2}=e^{-t^2}=e^{-(\sqrt{2})^2t^2/2}$.
Bu $\Normal(0,2)$'nin KF'sidir. Dolayısıyla $Z\sim\Normal(0,2)$.

(c) $\varphi_{\Pois(\lambda)}(t)=\exp(\lambda(e^{it}-1))$.
$\varphi_Z(t)=\exp(\lambda_1(e^{it}-1))\cdot\exp(\lambda_2(e^{it}-1))=\exp((\lambda_1+\lambda_2)(e^{it}-1))$.
Bu $\Pois(\lambda_1+\lambda_2)$'nin KF'sidir. $Z\sim\Pois(\lambda_1+\lambda_2)$. $\checkmark$}

\begin{alistirma}[Al.4.4c]{C}{Toplam Kuralı ile Ortak Dağılım Analizi}
$(X,Y)$ ortak yoğunluğu $f(x,y)=6xy^2$, $0<x<1$, $0<y<1$.
(a) Marjinal yoğunlukları $f_X(x)$ ve $f_Y(y)$'yi bulun.
(b) $X$ ve $Y$ bağımsız mı?
(c) $\Cov(X,Y)$, $\Cor(X,Y)$'yi hesaplayın.
(d) $E[Y\mid X=0.5]$'i bulun. Kule özelliğiyle $E[Y]$'yi doğrulayın.
\end{alistirma}
\alcozum{%
(a) $f_X(x)=\int_0^1 6xy^2\,dy=6x\cdot[y^3/3]_0^1=2x$, $x\in(0,1)$.\\
$f_Y(y)=\int_0^1 6xy^2\,dx=6y^2\cdot[x^2/2]_0^1=3y^2$, $y\in(0,1)$.

(b) $f(x,y)=6xy^2=(2x)(3y^2)=f_X(x)f_Y(y)$. Evet, $X$ ve $Y$ \textbf{bağımsız}.

(c) Bağımsız olduklarından $\Cov(X,Y)=0$ ve $\Cor(X,Y)=0$.
Doğrulama: $E[X]=\int_0^1 2x^2dx=2/3$; $E[X^2]=\int_0^1 2x^3dx=1/2$; $\Var(X)=1/2-4/9=1/18$.\\
$E[Y]=\int_0^1 3y^3dy=3/4$; $E[Y^2]=\int_0^1 3y^4dy=3/5$; $\Var(Y)=3/5-9/16=48/80-45/80=3/80$.\\
$E[XY]=E[X]E[Y]=(2/3)(3/4)=1/2$. $\Cov(X,Y)=E[XY]-E[X]E[Y]=1/2-1/2=0$. $\checkmark$

(d) $f_{Y|X}(y|x)=f(x,y)/f_X(x)=6xy^2/(2x)=3y^2$ — $X$'den bağımsız! $E[Y|X=0.5]=E[Y]=\int_0^1 3y^3dy=3/4$.\\
Kule: $E[E[Y|X]]=E[3/4]=3/4=E[Y]$. $\checkmark$}

% ---

\alistirmalar

\begin{alistirma}[Al.4.5]{A}{Beklenti Özellikleri}
$X$ rasgele değişken, $E[X] = 3$, $E[X^2] = 13$.
$E[(2X-1)^2]$ ve $\mathrm{Var}(3X+5)$'i hesaplayın.
\end{alistirma}
\alcozum{%
$E[(2X-1)^2] = 4E[X^2] - 4E[X] + 1 = 52 - 12 + 1 = 41$.
$\mathrm{Var}(X) = 13-9=4$. $\mathrm{Var}(3X+5) = 9\cdot 4 = 36$.}

\begin{alistirma}[Al.4.6]{A}{Chebyshev Uygulaması}
$X$ rasgele değişken, $E[X] = 50$, $\mathrm{Var}(X) = 25$.
(a) $P(|X - 50| \geq 10)$'un Chebyshev üst sınırını bulun.
(b) $P(40 \leq X \leq 60) \geq ?$ alt sınırını verin.
\end{alistirma}
\alcozum{%
(a) $\varepsilon = 10$: $P(|X-50|\geq 10) \leq 25/100 = 1/4$.
(b) $P(40 \leq X \leq 60) = 1 - P(|X-50| \geq 10) \geq 3/4$.}

\begin{alistirma}[Al.4.7]{B}{Markov Eşitsizliği ve Kesinlik}
$X \geq 0$, $E[X] = \mu$.
(a) $P(X \geq k\mu) \leq 1/k$ olduğunu gösterin.
(b) $P(X \geq 2\mu) = 1/2$ sağlayan bir dağılım inşa edin (Markov sıkılığı).
\end{alistirma}
\alcozum{%
(a) $a = k\mu$ ile Markov: $P(X \geq k\mu) \leq E[X]/(k\mu) = 1/k$.\\
(b) $P(X=0) = 1/2$, $P(X = 2\mu) = 1/2$: $E[X] = \mu$, $P(X \geq 2\mu) = 1/2$. Eşitlik sağlanır.}

\begin{alistirma}[Al.4.8]{B}{Kovaryans Matrisi}
$(X, Y, Z)$: $\mathrm{Var}(X) = 1$, $\mathrm{Var}(Y) = 4$, $\mathrm{Var}(Z) = 9$;
$\mathrm{Cov}(X,Y) = 1$, $\mathrm{Cov}(X,Z) = -1$, $\mathrm{Cov}(Y,Z) = 2$.
$\mathrm{Var}(X + Y + Z)$'i hesaplayın.
\end{alistirma}
\alcozum{%
$\mathrm{Var}(X+Y+Z) = \mathrm{Var}(X) + \mathrm{Var}(Y) + \mathrm{Var}(Z) + 2\mathrm{Cov}(X,Y) + 2\mathrm{Cov}(X,Z) + 2\mathrm{Cov}(Y,Z)$
$= 1 + 4 + 9 + 2(1) + 2(-1) + 2(2) = 14 + 2 - 2 + 4 = 18$.}

\begin{alistirma}[Al.4.9]{B}{Karakteristik Fonksiyon}
$X \sim \mathrm{Bernoulli}(p)$.
(a) $\varphi_X(t) = 1 - p + pe^{it}$ olduğunu gösterin.
(b) $X_1, \ldots, X_n$ bağımsız $\mathrm{Bernoulli}(p)$; $S_n = \sum X_i$.
$\varphi_{S_n}(t)$'yi bulun ve bunun hangi dağılıma ait olduğunu tanıyın.
\end{alistirma}
\alcozum{%
(a) $\varphi_X(t) = E[e^{itX}] = (1-p)e^{0} + pe^{it} = 1-p+pe^{it}$.\\
(b) Bağımsız çarpım: $\varphi_{S_n}(t) = (1-p+pe^{it})^n = \varphi_{\mathrm{Binom}(n,p)}(t)$.
Bu binom dağılımının KF'sidir; $S_n \sim \Binom(n,p)$.}

\begin{alistirma}[Al.4.10]{B}{Kule Özelliği}
$N \sim \mathrm{Pois}(\lambda)$ ve $X_1, X_2, \ldots$ bağımsız, $E[X_i] = \mu$, $\mathrm{Var}(X_i) = \sigma^2$.
$S_N = \sum_{i=1}^N X_i$ (toplam öğe sayısı rastlantısal).
(a) $E[S_N]$'i bulun. (b) $\mathrm{Var}(S_N)$'i bulun.
\end{alistirma}
\alcozum{%
(a) $E[S_N] = E[E[S_N \mid N]] = E[N\mu] = \mu E[N] = \mu\lambda$.\\
(b) Varyans ayrışım formülü:
$\mathrm{Var}(S_N) = E[\mathrm{Var}(S_N \mid N)] + \mathrm{Var}(E[S_N \mid N])
= E[N\sigma^2] + \mathrm{Var}(N\mu) = \sigma^2\lambda + \mu^2\lambda = \lambda(\sigma^2+\mu^2)$.}

\begin{alistirma}[Al.4.11]{C}{İspat: Cauchy-Schwarz Genel Formu}
$X, Y \in L^2(\Omega, \mathcal{F}, P)$ için $(E[XY])^2 \leq E[X^2] E[Y^2]$ olduğunu ispatlayın.
Eşitlik ne zaman sağlanır?
\end{alistirma}
\alcozum{%
Her $t \in \mathbb{R}$ için $f(t) = E[(tX+Y)^2] = t^2E[X^2] + 2tE[XY] + E[Y^2] \geq 0$.
Bu kuadratik $t$'de negatif olamaz; diskriminant $\leq 0$:
$4(E[XY])^2 - 4E[X^2]E[Y^2] \leq 0$.
Eşitlik: $f(t^*) = 0$ yani $t^*X + Y = 0$ n.k., yani $Y = cX$ n.k.\ (lineer bağımlılık).}

\begin{alistirma}[Al.4.12]{C}{Varyans Ayrışım Formülü (Kule Varyansı)}
$E[X^2] < \infty$ ve $\mathcal{G} \subseteq \mathcal{F}$ alt $\sigma$-cebiri olsun.
$\mathrm{Var}(X) = E[\mathrm{Var}(X \mid \mathcal{G})] + \mathrm{Var}(E[X \mid \mathcal{G}])$
olduğunu ispatlayın.
\end{alistirma}
\alcozum{%
$\mathrm{Var}(X) = E[X^2] - (E[X])^2$.
$E[\mathrm{Var}(X\mid\mathcal{G})] = E[E[X^2\mid\mathcal{G}] - (E[X\mid\mathcal{G}])^2]
= E[X^2] - E[(E[X\mid\mathcal{G}])^2]$.
$\mathrm{Var}(E[X\mid\mathcal{G}]) = E[(E[X\mid\mathcal{G}])^2] - (E[E[X\mid\mathcal{G}]])^2
= E[(E[X\mid\mathcal{G}])^2] - (E[X])^2$.
Toplam: $E[X^2] - (E[X])^2 = \mathrm{Var}(X)$. \checkmark}

\begin{alistirma}[Al.4.13]{A}{Chebyshev ve Cantelli Karşılaştırması}
$X$ rasgele değişken, $E[X] = 2$, $\mathrm{Var}(X) = 1$.
(a) Chebyshev ile $P(|X - 2| \geq 2)$ için üst sınır verin.
(b) Cantelli ile $P(X - 2 \geq 2)$ için üst sınır verin.
(c) Sonuçları karşılaştırın.
\end{alistirma}
\alcozum{%
(a) Chebyshev: $P(|X-2|\geq2) \leq \mathrm{Var}(X)/4 = 1/4$.\\
(b) Cantelli: $P(X-2\geq2) \leq \sigma^2/(\sigma^2+\lambda^2) = 1/(1+4) = 1/5$.\\
(c) Cantelli tek taraflı sınır veriyor ve daha keskin: $1/5 < 1/4$.
Chebyshev iki taraflı; tam $P(X-2\geq2)$ için Cantelli daha iyi.}

\begin{alistirma}[Al.4.14]{B}{Hoeffding Uygulaması}
$X_1, \ldots, X_{100}$ bağımsız, $X_i \in [0,1]$, $E[X_i] = 0.5$.
(a) Hoeffding eşitsizliği ile $P(\bar{X}_{100} \geq 0.6)$ için üst sınır verin.
(b) Chebyshev ile (varyanslı $\sigma^2 = 1/12$ alarak) karşılaştırın.
(c) $P \leq 0.01$ için gereken minimum $n$'yi Hoeffding ile bulun ($\varepsilon = 0.1$).
\end{alistirma}
\alcozum{%
(a) Hoeffding: $P(\bar{X}_n - 0.5 \geq 0.1) \leq e^{-2n(0.1)^2/1} = e^{-2(100)(0.01)} = e^{-2} \approx 0.135$.\\
(b) Chebyshev: $P(|\bar{X}_{100}-0.5|\geq0.1) \leq \sigma^2/(n\varepsilon^2) = (1/12)/(100\cdot0.01) = 1/12 \approx 0.083$.
Bu örnekte Chebyshev daha iyi — Hoeffding (sınırlı) avantajı daha büyük $n$'de belirginleşir.\\
(c) $e^{-2n\varepsilon^2} \leq 0.01$: $-2n(0.01) \leq \ln(0.01) = -4.605$; $n \geq 230$.}

\begin{alistirma}[Al.4.15]{B}{Çarpıklık ve Basıklık}
$X \sim \Exp(1)$: $E[X] = 1$, $\mathrm{Var}(X) = 1$.
(a) $E[X^3]$ ve $E[X^4]$'ü hesaplayın ($E[X^k] = k!$ Üstel için).
(b) Çarpıklık $\gamma_1 = (E[(X-\mu)^3])/\sigma^3$'ü hesaplayın.
(c) Basıklık fazlasını $\gamma_2 = E[(X-\mu)^4]/\sigma^4 - 3$ hesaplayın.
\end{alistirma}
\alcozum{%
(a) $E[X^3] = 3! = 6$; $E[X^4] = 4! = 24$.\\
(b) $\mu_3' = E[(X-1)^3] = E[X^3 - 3X^2 + 3X - 1] = 6 - 3(2) + 3(1) - 1 = 6-6+3-1 = 2$.
$\gamma_1 = 2/1^3 = 2$: kuvvetli sağ çarpıklık.\\
(c) $\mu_4' = E[(X-1)^4] = E[X^4 - 4X^3 + 6X^2 - 4X + 1] = 24 - 24 + 12 - 4 + 1 = 9$.
$\gamma_2 = 9/1 - 3 = 6$: kuvvetli ağır kuyruk (Normal'den 6 kat fazla).}

% =====================================================================
\endinput

% =====================================================================
%  BÖLÜM 5: Önemli Dağılım Aileleri
%  Kaynak: Feller Vol.I-II; Hogg & Craig; Casella & Berger; Klenke
% =====================================================================
\chapter{Önemli Dağılım Aileleri}
\label{chp:dagilim_aileleri}

\bolumalinti{%
  Doğa, birkaç temel dağılımın sonsuz çeşitlemesinden ibarettir.%
}{William Feller, \textit{An Introduction to Probability Theory and Its Applications}}

% ---- Kazanımlar (TYMM) ----
\begin{kazanimlar}
Bu bölümün sonunda öğrenci:
\begin{itemize}[leftmargin=1.8em]
  \item Temel kesikli dağılımları (Bernoulli, Binom, Poisson, Geometrik,
        Negatif Binom, Hipergeometrik, Çok Terimli) tanımlayabilir;
        beklenti, varyans ve MÜF'lerini hesaplayabilir.
  \item Temel sürekli dağılımları (Uniform, Normal, Üstel, Gama, Beta,
        Cauchy, Lognormal, $\chi^2$, $t$, $F$) tanımlayabilir; türetme ilişkilerini açıklayabilir.
  \item Üstel aile kavramını tanımlayabilir; doğal parametre ve yeterli istatistik
        kavramlarıyla bağlantısını kurabilir.
  \item Yer-ölçek ailesi kavramını açıklayabilir; standartlaştırma işlemini uygulayabilir.
  \item Dağılımlar arasındaki türetme ve limit ilişkilerini (Poisson limiti,
        normal yaklaşım, $\chi^2$ toplamı) ispatlayabilir.
\end{itemize}
\end{kazanimlar}

% ---- Motivasyon ----
\begin{motivasyon}
Bölüm~\ref{chp:rasgele_degiskenler} ve~\ref{chp:beklenti}'de dağılım fonksiyonu
ve beklenti soyut düzeyde incelendi. Şimdi pratikte en sık karşılaşılan dağılım
ailelerini somutlaştırıyoruz. Bu dağılımlar teorik güzelliklerinin yanı sıra
istatistiksel modellemenin (Kısım~\ref{part:istatistik}) temel yapı taşlarıdır.

Her dağılım için standart bilgi şeması: PMF/PDF — parametreler — $E[X]$, $\mathrm{Var}(X)$ —
MÜF/KF — özellikler — tipik kullanım alanı.

\medskip
\textbf{Köprü.} Bölüm~\ref{chp:beklenti}'deki MÜF ve KF araçları burada her dağılım
için hesaplanacak; Bölüm~\ref{chp:buyuk_sayilar} ve~\ref{chp:merkezi_limit}'te
bu dağılımların asimptotik davranışı incelenecektir.
\end{motivasyon}

% =====================================================================
\section{Kesikli Dağılımlar}
\label{sec:kesikli_dagilimlar}
% =====================================================================

\subsection{Bernoulli ve Binom Dağılımı}
\label{subsec:bernoulli_binom}

\begin{tanim}[Bernoulli Dağılımı]
\label{def:bernoulli}
$0 < p < 1$ olsun. $X \sim \mathrm{Bernoulli}(p)$:
\[
  P(X = 1) = p, \quad P(X = 0) = 1 - p =: q.
\]
\textbf{Beklenti:} $E[X] = p$. \quad
\textbf{Varyans:} $\mathrm{Var}(X) = pq$. \quad
\textbf{MÜF:} $M_X(t) = q + pe^t$.
\end{tanim}

\begin{tanim}[Binom Dağılımı]
\label{def:binom}
$X \sim \Binom(n, p)$: $n$ bağımsız $\mathrm{Bernoulli}(p)$ denemesinde başarı sayısı.
\[
  P(X = k) = \binom{n}{k} p^k q^{n-k}, \quad k = 0, 1, \ldots, n.
\]
\textbf{Beklenti:} $E[X] = np$. \quad
\textbf{Varyans:} $\mathrm{Var}(X) = npq$. \quad
\textbf{MÜF:} $M_X(t) = (q + pe^t)^n$.
\end{tanim}

\begin{proof}[MÜF İspat]
$X = X_1 + \cdots + X_n$ ($X_i$ bağımsız Bernoulli).
$M_X(t) = \prod_{i=1}^n M_{X_i}(t) = (q+pe^t)^n$.
\end{proof}

\begin{teorem}[Binom Yeniden Üretkenliği]
\label{thm:binom_yeniden}
$X \sim \Binom(n, p)$, $Y \sim \Binom(m, p)$ bağımsız $\Rightarrow$
$X + Y \sim \Binom(n+m, p)$.
\end{teorem}

\begin{proof}
$M_{X+Y}(t) = (q+pe^t)^n(q+pe^t)^m = (q+pe^t)^{n+m}$.
\end{proof}

\begin{onerme}[Binom Mod ve Medyan]
$\lfloor (n+1)p \rfloor$ binom dağılımının modudur (birden fazla mod olabilir).
$n = 10$, $p = 0.3$: mod $= \lfloor 11 \cdot 0.3 \rfloor = 3$.
\end{onerme}

\subsection{Poisson Dağılımı}
\label{subsec:poisson}

\begin{tanim}[Poisson Dağılımı]
\label{def:poisson}
$\lambda > 0$ olsun. $X \sim \Pois(\lambda)$:
\[
  P(X = k) = \frac{\lambda^k e^{-\lambda}}{k!}, \quad k = 0, 1, 2, \ldots
\]
\textbf{Beklenti:} $E[X] = \lambda$. \quad
\textbf{Varyans:} $\mathrm{Var}(X) = \lambda$. \quad
\textbf{MÜF:} $M_X(t) = e^{\lambda(e^t - 1)}$.
\end{tanim}

\begin{proof}[Beklenti ve MÜF]
$E[X] = \sum_{k=0}^\infty k \frac{\lambda^k e^{-\lambda}}{k!}
= \lambda e^{-\lambda} \sum_{k=1}^\infty \frac{\lambda^{k-1}}{(k-1)!} = \lambda$.
$M_X(t) = \sum_{k=0}^\infty e^{tk} \frac{\lambda^k e^{-\lambda}}{k!}
= e^{-\lambda} \sum_{k=0}^\infty \frac{(\lambda e^t)^k}{k!} = e^{-\lambda} e^{\lambda e^t} = e^{\lambda(e^t-1)}$.
\end{proof}

\begin{teorem}[Binom'dan Poisson Limiti]
\label{thm:binom_poisson}
$n \to \infty$, $p = p_n \to 0$, $np_n \to \lambda > 0$ iken
$\Binom(n, p_n) \xrightarrow{d} \Pois(\lambda)$.
\end{teorem}

\begin{proof}
$P(X_n = k) = \binom{n}{k} p_n^k (1-p_n)^{n-k}$.
$np_n \to \lambda$, yani $p_n = \lambda_n/n$ ($\lambda_n = np_n \to \lambda$).
\begin{align*}
  \binom{n}{k} p_n^k (1-p_n)^{n-k}
  &= \frac{n(n-1)\cdots(n-k+1)}{k!} \cdot \frac{\lambda_n^k}{n^k} \cdot \left(1 - \frac{\lambda_n}{n}\right)^{n-k} \\
  &\to \frac{1}{k!} \cdot \lambda^k \cdot e^{-\lambda} = \frac{\lambda^k e^{-\lambda}}{k!}.
\end{align*}
(Her $k$ için: ilk çarpan $\to 1$, $(1-\lambda_n/n)^n \to e^{-\lambda}$, $(1-\lambda_n/n)^{-k} \to 1$.)
\end{proof}

\begin{onerme}[Poisson Yeniden Üretkenliği]
$X \sim \Pois(\lambda)$, $Y \sim \Pois(\mu)$ bağımsız $\Rightarrow$ $X+Y \sim \Pois(\lambda+\mu)$.
\end{onerme}

\begin{proof}
$M_{X+Y}(t) = e^{\lambda(e^t-1)} \cdot e^{\mu(e^t-1)} = e^{(\lambda+\mu)(e^t-1)}$.
\end{proof}

\begin{ornek}[Poisson Süreci]
Birim zamanda ortalama $\lambda$ müşteri gelen bir bankada,
$k$ müşteri gelme olasılığı $P(X=k) = e^{-\lambda}\lambda^k/k!$.
$\lambda = 3$: $P(X = 0) = e^{-3} \approx 0.050$, $P(X \leq 2) \approx 0.423$.
\end{ornek}

\subsection{Geometrik ve Negatif Binom}
\label{subsec:geometrik}

\begin{tanim}[Geometrik Dağılım]
\label{def:geometrik}
$X \sim \mathrm{Geom}(p)$: ilk başarıya kadar gereken deneme sayısı.
\[
  P(X = k) = q^{k-1} p, \quad k = 1, 2, 3, \ldots
\]
\textbf{Beklenti:} $E[X] = 1/p$. \quad \textbf{Varyans:} $\mathrm{Var}(X) = q/p^2$.
\textbf{MÜF:} $M_X(t) = pe^t / (1 - qe^t)$, $t < -\ln q$.
\end{tanim}

\begin{onerme}[Bellek Yoksunluğu]
\label{prop:bellek_yoksunlugu}
$X \sim \mathrm{Geom}(p)$ için her $m, n \geq 1$'de
$P(X > m+n \mid X > m) = P(X > n)$.
Kesikli dağılımlar arasında bellek yoksunluğu (memorylessness) özelliğine sahip
tek dağılım geometrik dağılımdır.
\end{onerme}

\begin{proof}
$P(X > k) = q^k$.
$P(X > m+n \mid X > m) = \frac{P(X > m+n)}{P(X > m)} = \frac{q^{m+n}}{q^m} = q^n = P(X > n)$.
Tersine: bu özelliği sağlayan tek kesikli dağılımın geometrik olduğu,
$P(X > k) = P(X > 1)^k$ çarpım özelliğiyle gösterilir.
\end{proof}

\begin{tanim}[Negatif Binom Dağılımı]
\label{def:neg_binom}
$X \sim \mathrm{NegBin}(r, p)$: $r$'inci başarıya kadar gereken deneme sayısı.
\[
  P(X = k) = \binom{k-1}{r-1} p^r q^{k-r}, \quad k = r, r+1, \ldots
\]
\textbf{Beklenti:} $E[X] = r/p$. \quad
\textbf{Varyans:} $\mathrm{Var}(X) = rq/p^2$.\\
$\mathrm{NegBin}(1, p) = \mathrm{Geom}(p)$.
\end{tanim}

\subsection{Hipergeometrik Dağılım}
\label{subsec:hiper}

\begin{tanim}[Hipergeometrik Dağılım]
\label{def:hipergeometrik}
$N$ elemanlı popülasyonda $K$ "başarılı" eleman var; $n$ elemanlık örneklem çekilsin
(yerine koymadan). $X$: örneklemdeki başarılı sayısı.
\[
  P(X = k) = \frac{\binom{K}{k}\binom{N-K}{n-k}}{\binom{N}{n}},
  \quad \max(0, n+K-N) \leq k \leq \min(n, K).
\]
\textbf{Beklenti:} $E[X] = nK/N$. \quad
\textbf{Varyans:} $\mathrm{Var}(X) = n \frac{K}{N}\frac{N-K}{N}\frac{N-n}{N-1}$.
\end{tanim}

\begin{onerme}[Binom Yaklaşımı]
$N \to \infty$, $K/N \to p$ iken $\mathrm{Hiper}(N, K, n) \to \Binom(n, p)$.
Yani büyük popülasyonda yerine koymadan örneklem ≈ yerine koyarak örneklem.
\end{onerme}

\subsection{Çok Terimli Dağılım}
\label{subsec:multinomial}

\begin{tanim}[Çok Terimli Dağılım]
\label{def:multinomial}
$n$ bağımsız deneme; her deneyde $k$ kategoriden biri $p_1, \ldots, p_k$ olasılıklarla seçilsin
($\sum p_i = 1$). $X_i$: $i$'inci kategorinin deneme sayısı.
$(X_1, \ldots, X_k) \sim \mathrm{Multinomial}(n; p_1, \ldots, p_k)$:
\[
  P(X_1 = n_1, \ldots, X_k = n_k) = \frac{n!}{n_1! \cdots n_k!} p_1^{n_1} \cdots p_k^{n_k},
  \quad \sum n_i = n.
\]
\textbf{Beklenti:} $E[X_i] = np_i$. \quad
\textbf{Varyans:} $\mathrm{Var}(X_i) = np_i(1-p_i)$. \quad
\textbf{Kovaryans:} $\mathrm{Cov}(X_i, X_j) = -np_i p_j$ ($i \neq j$).
\end{tanim}

% =====================================================================
\section{Sürekli Dağılımlar}
\label{sec:surekli_dagilimlar}
% =====================================================================

\subsection{Düzgün (Uniform) Dağılım}
\label{subsec:uniform}

\begin{tanim}[Uniform Dağılım]
\label{def:uniform}
$X \sim \Uniform(a, b)$ ($a < b$):
\[
  f_X(x) = \frac{1}{b-a}, \quad a < x < b.
\]
\textbf{Beklenti:} $E[X] = (a+b)/2$. \quad
\textbf{Varyans:} $\mathrm{Var}(X) = (b-a)^2/12$. \quad
\textbf{MÜF:} $M_X(t) = \frac{e^{tb}-e^{ta}}{t(b-a)}$ ($t \neq 0$).
\end{tanim}

\begin{onerme}[Quantile Dönüşümü]
\label{prop:quantile}
$U \sim \Uniform(0,1)$ ve $F$ herhangi bir KDF ise $X = F^{-1}(U) \sim F$.
Bu simülasyonun temelidir: uniform rassal sayılardan herhangi bir dağılım üretilebilir.
\end{onerme}

\begin{proof}
$P(F^{-1}(U) \leq x) = P(U \leq F(x)) = F(x)$ ($U$ uniform olduğundan).
\end{proof}

\subsection{Normal Dağılım}
\label{subsec:normal}

\begin{tanim}[Normal (Gauss) Dağılımı]
\label{def:normal}
$X \sim \Normal(\mu, \sigma^2)$ ($\sigma > 0$):
\[
  f_X(x) = \frac{1}{\sigma\sqrt{2\pi}} \exp\!\left(-\frac{(x-\mu)^2}{2\sigma^2}\right), \quad x \in \mathbb{R}.
\]
\textbf{Beklenti:} $E[X] = \mu$. \quad
\textbf{Varyans:} $\mathrm{Var}(X) = \sigma^2$. \quad
\textbf{MÜF:} $M_X(t) = e^{\mu t + \sigma^2 t^2/2}$.\\
$Z = (X-\mu)/\sigma \sim \Normal(0,1)$ \textbf{standart normal}.
\end{tanim}

\begin{onerme}[Standart Normal Özellikleri]
$Z \sim \Normal(0,1)$. PDF: $\phi(z) = (2\pi)^{-1/2} e^{-z^2/2}$. KDF: $\Phi(z) = \int_{-\infty}^z \phi(t)\,dt$.
\begin{itemize}
  \item $\phi(-z) = \phi(z)$ (simetri); $\Phi(-z) = 1 - \Phi(z)$.
  \item $E[Z^{2k}] = (2k-1)!! = 1 \cdot 3 \cdots (2k-1)$; $E[Z^{2k+1}] = 0$.
  \item $E[Z^4] = 3$: basıklık 0 (referans).
\end{itemize}
\end{onerme}

\begin{teorem}[Normal Yeniden Üretkenliği]
\label{thm:normal_yeniden}
$X_i \sim \Normal(\mu_i, \sigma_i^2)$ bağımsız ($i=1,\ldots,n$) $\Rightarrow$
$\sum_{i=1}^n a_i X_i \sim \Normal\!\left(\sum a_i\mu_i,\ \sum a_i^2\sigma_i^2\right)$.
\end{teorem}

\begin{proof}
MÜF çarpım: $M_{\sum a_i X_i}(t) = \prod_i e^{\mu_i(a_i t) + \sigma_i^2(a_i t)^2/2}
= \exp\!\bigl(\sum a_i\mu_i \cdot t + \tfrac{t^2}{2}\sum a_i^2\sigma_i^2\bigr)$.
\end{proof}

\begin{tanim}[Çok Değişkenli Normal Dağılım]
\label{def:cok_normal}
$\mathbf{X} = (X_1, \ldots, X_n)^\top \sim \Normal_n(\boldsymbol{\mu}, \boldsymbol{\Sigma})$:
$\boldsymbol{\mu} \in \mathbb{R}^n$, $\boldsymbol{\Sigma}$ pozitif yarı-tanımlı $n\times n$ matris.
$\boldsymbol{\Sigma}$ pozitif tanımlı ise ortak PDF:
\[
  f(\mathbf{x}) = \frac{1}{(2\pi)^{n/2}|\boldsymbol{\Sigma}|^{1/2}}
  \exp\!\left(-\tfrac{1}{2}(\mathbf{x}-\boldsymbol{\mu})^\top \boldsymbol{\Sigma}^{-1}(\mathbf{x}-\boldsymbol{\mu})\right).
\]
$E[\mathbf{X}] = \boldsymbol{\mu}$, $\mathrm{Cov}(\mathbf{X}) = \boldsymbol{\Sigma}$.
\end{tanim}

\begin{teorem}[Çok Değişkenli Normal: Bağımsızlık]
\label{thm:cok_normal_bagimsiz}
$\mathbf{X} \sim \Normal_n(\boldsymbol{\mu}, \boldsymbol{\Sigma})$ ise
bileşenler çiftler arası korelasyonsuzluğu $\Leftrightarrow$ bağımsızlıktır:
$\Sigma_{ij} = 0$ ($i \neq j$) $\Leftrightarrow$ $X_i \perp X_j$.
\end{teorem}

\begin{proof}
$\boldsymbol{\Sigma}$ köşegen ise ortak PDF $\prod_i f_{X_i}(x_i)$ çarpımına ayrışır.
\end{proof}

\begin{hataanalizi}
Bu teorem genel dağılımlar için yanlıştır! Yalnızca çok değişkenli normal için
korelasyonsuzluk ile bağımsızlık eşdeğerdir. Genel dağılımlarda korelasyonsuz
ama bağımlı rasgele değişkenler inşa etmek mümkündür (bkz.\ Bölüm~\ref{sec:kovaryans}).
\end{hataanalizi}

\subsection{Üstel ve Gama Dağılımı}
\label{subsec:gama}

\begin{tanim}[Üstel Dağılım]
\label{def:exp}
$X \sim \Exp(\lambda)$ ($\lambda > 0$):
\[
  f_X(x) = \lambda e^{-\lambda x}, \quad x > 0.
\]
\textbf{Beklenti:} $E[X] = 1/\lambda$. \quad
\textbf{Varyans:} $\mathrm{Var}(X) = 1/\lambda^2$. \quad
\textbf{MÜF:} $M_X(t) = \lambda/(\lambda - t)$, $t < \lambda$.
\end{tanim}

\begin{onerme}[Bellek Yoksunluğu — Sürekli]
$X \sim \Exp(\lambda)$ için her $s, t > 0$'da $P(X > s+t \mid X > s) = P(X > t)$.
Sürekli dağılımlar arasında bellek yoksunluğu özelliğine sahip tek dağılım üsteldir.
\end{onerme}

\begin{proof}
$P(X > t) = e^{-\lambda t}$.
$P(X > s+t \mid X > s) = e^{-\lambda(s+t)}/e^{-\lambda s} = e^{-\lambda t} = P(X > t)$.
Tersine: $P(X > s+t) = P(X > s)P(X > t)$ ise $G(t) := P(X>t)$ için
$G(s+t) = G(s)G(t)$ ve süreklilik $G(t) = e^{-\lambda t}$ verir.
\end{proof}

\begin{tanim}[Gama Dağılımı]
\label{def:gama}
$X \sim \Gama(\alpha, \beta)$ ($\alpha, \beta > 0$):
\[
  f_X(x) = \frac{\beta^\alpha}{\Gamma(\alpha)} x^{\alpha-1} e^{-\beta x}, \quad x > 0.
\]
$\Gamma(\alpha) = \int_0^\infty t^{\alpha-1} e^{-t} \, dt$ (Gama fonksiyonu).
\textbf{Beklenti:} $E[X] = \alpha/\beta$. \quad
\textbf{Varyans:} $\mathrm{Var}(X) = \alpha/\beta^2$. \quad
\textbf{MÜF:} $M_X(t) = (\beta/(\beta-t))^\alpha$, $t < \beta$.\\
Özel durumlar: $\Gama(1, \beta) = \Exp(\beta)$; $\Gama(n/2, 1/2) = \chi^2_n$ (aşağıda).
\end{tanim}

\begin{proof}[Beklenti]
$E[X] = \frac{\beta^\alpha}{\Gamma(\alpha)} \int_0^\infty x^\alpha e^{-\beta x} dx
= \frac{\beta^\alpha}{\Gamma(\alpha)} \cdot \frac{\Gamma(\alpha+1)}{\beta^{\alpha+1}}
= \frac{\alpha}{\beta}$.
($\Gamma(\alpha+1) = \alpha\Gamma(\alpha)$ yinelemeli bağıntısı kullanıldı.)
\end{proof}

\begin{teorem}[Gama Yeniden Üretkenliği]
$X_i \sim \Gama(\alpha_i, \beta)$ bağımsız ($\beta$ aynı) $\Rightarrow$
$\sum X_i \sim \Gama(\sum \alpha_i, \beta)$.
\end{teorem}

\begin{proof}
$M_{\sum X_i}(t) = \prod_i (\beta/(\beta-t))^{\alpha_i} = (\beta/(\beta-t))^{\sum\alpha_i}$.
\end{proof}

\subsection{Beta Dağılımı}
\label{subsec:beta}

\begin{tanim}[Beta Dağılımı]
\label{def:beta}
$X \sim \Beta(\alpha, \beta)$ ($\alpha, \beta > 0$):
\[
  f_X(x) = \frac{x^{\alpha-1}(1-x)^{\beta-1}}{B(\alpha,\beta)}, \quad 0 < x < 1.
\]
$B(\alpha, \beta) = \Gamma(\alpha)\Gamma(\beta)/\Gamma(\alpha+\beta)$ (Beta fonksiyonu).\\
\textbf{Beklenti:} $E[X] = \alpha/(\alpha+\beta)$.\\
\textbf{Varyans:} $\mathrm{Var}(X) = \alpha\beta / [(\alpha+\beta)^2(\alpha+\beta+1)]$.\\
Özel durum: $\Beta(1,1) = \Uniform(0,1)$.
\end{tanim}

\begin{onerme}[Beta'nın Gama ile Bağlantısı]
$X \sim \Gama(\alpha, 1)$, $Y \sim \Gama(\beta, 1)$ bağımsız $\Rightarrow$
$U = X/(X+Y) \sim \Beta(\alpha, \beta)$, bağımsız olarak $X+Y \sim \Gama(\alpha+\beta, 1)$.
\end{onerme}

\begin{proof}
$(U, V) = (X/(X+Y), X+Y)$ dönüşümü; Jacobian $= v$;
$f_{U,V}(u,v) = f_{X,Y}(uv, v(1-u)) \cdot v$ hesaplanır;
$U$ ve $V$'nin bağımsız Beta ve Gama olduğu görülür.
\end{proof}

\subsection{Cauchy Dağılımı}
\label{subsec:cauchy}

\begin{tanim}[Cauchy Dağılımı]
\label{def:cauchy}
$X \sim \mathrm{Cauchy}(\mu, \sigma)$ ($\sigma > 0$):
\[
  f_X(x) = \frac{1}{\pi\sigma\left[1 + \left(\frac{x-\mu}{\sigma}\right)^2\right]}, \quad x \in \mathbb{R}.
\]
\textbf{Beklenti mevcut değildir} ($\int |x| f(x) dx = \infty$).\\
\textbf{KF:} $\varphi_X(t) = e^{i\mu t - \sigma|t|}$ (Laplace çekirdeği).\\
$\mathrm{Cauchy}(0,1)$: $Z_1/Z_2$ ($Z_1, Z_2$ bağımsız standart normal).
\end{tanim}

\begin{karsiornek}[Cauchy — Büyük Sayılar Kanunu Çöküyor]
$X_1, X_2, \ldots$ bağımsız $\mathrm{Cauchy}(0,1)$ ise
$\bar{X}_n = n^{-1}\sum X_i$ de $\mathrm{Cauchy}(0,1)$ dağılımlıdır —
ortalama $n$ büüdükçe stabilize olmaz. Büyük Sayılar Kanunu beklentinin
varlığını gerektirir; Cauchy bunu ihlal eder.
\end{karsiornek}

\subsection{Lognormal Dağılım}
\label{subsec:lognormal}

\begin{tanim}[Lognormal Dağılım]
\label{def:lognormal}
$Y = e^X$, $X \sim \Normal(\mu, \sigma^2)$ ise $Y \sim \mathrm{Lognormal}(\mu, \sigma^2)$:
\[
  f_Y(y) = \frac{1}{y\sigma\sqrt{2\pi}} \exp\!\left(-\frac{(\ln y - \mu)^2}{2\sigma^2}\right), \quad y > 0.
\]
\textbf{Beklenti:} $E[Y] = e^{\mu + \sigma^2/2}$.\\
\textbf{Varyans:} $\mathrm{Var}(Y) = e^{2\mu+\sigma^2}(e^{\sigma^2}-1)$.\\
Tüm momentler mevcuttur ama MÜF yoktur ($E[e^{tY}] = \infty$ her $t > 0$ için).
\end{tanim}

\begin{ornek}[Finans Uygulaması]
Black-Scholes modelinde hisse senedi fiyatı $S_t = S_0 e^{(\mu-\sigma^2/2)t + \sigma W_t}$
lognormal dağılım izler. $\ln(S_t/S_0) \sim \Normal((\mu-\sigma^2/2)t, \sigma^2 t)$.
\end{ornek}

\subsection{Chi-Kare, \texorpdfstring{$t$}{t} ve \texorpdfstring{$F$}{F} Dağılımları}
\label{subsec:chi_t_f}

\begin{tanim}[Chi-Kare Dağılımı]
\label{def:chi_kare}
$Z_1, \ldots, Z_n$ bağımsız $\Normal(0,1)$ ise
\[
  \chi^2_n = \sum_{i=1}^n Z_i^2 \sim \Gama(n/2,\, 1/2).
\]
\textbf{Beklenti:} $E[\chi^2_n] = n$. \quad
\textbf{Varyans:} $\mathrm{Var}(\chi^2_n) = 2n$. \quad
\textbf{MÜF:} $M(t) = (1-2t)^{-n/2}$, $t < 1/2$.
\end{tanim}

\begin{proof}
$Z_i^2 \sim \Gama(1/2, 1/2)$ (Bölüm~\ref{sec:rd_donusum}'deki $Y=X^2$ dönüşümü).
Gama yeniden üretkenliği: $\sum Z_i^2 \sim \Gama(n/2, 1/2)$.
\end{proof}

\begin{tanim}[$t$ Dağılımı (Student)]
\label{def:t_dagilim}
$Z \sim \Normal(0,1)$, $V \sim \chi^2_n$ bağımsız. $T = Z/\sqrt{V/n} \sim t_n$ (Student-$t$).
PDF:
\[
  f_T(t) = \frac{\Gamma((n+1)/2)}{\sqrt{n\pi}\,\Gamma(n/2)}\left(1+\frac{t^2}{n}\right)^{-(n+1)/2}, \quad t \in \mathbb{R}.
\]
$n = 1$: $t_1 = \mathrm{Cauchy}(0,1)$.
$n \to \infty$: $t_n \xrightarrow{d} \Normal(0,1)$ (Böl.~\ref{chp:merkezi_limit}).\\
\textbf{Beklenti:} $E[T] = 0$ ($n > 1$). \quad \textbf{Varyans:} $\mathrm{Var}(T) = n/(n-2)$ ($n > 2$).
\end{tanim}

\begin{tanim}[$F$ Dağılımı]
\label{def:f_dagilim}
$U \sim \chi^2_m$, $V \sim \chi^2_n$ bağımsız.
$W = (U/m)/(V/n) \sim F(m, n)$.\\
\textbf{Beklenti:} $E[W] = n/(n-2)$ ($n > 2$).\\
Bağlantı: $T_n^2 \sim F(1, n)$.
\end{tanim}

\begin{hocanotu}
$t_n$ ve $F(m,n)$ dağılımları istatistiksel hipotez testinin (Bölüm~\ref{chp:hipotez_testi})
temel test istatistikleridir. Normal popülasyondan çekilen örneklem ortalaması
$\bar{X}$ ve örneklem varyansı $S^2$ ile bunlar türetilir (Böl.~\ref{chp:istatistik_temelleri}).
\end{hocanotu}

% =====================================================================
\section{Üstel Aile}
\label{sec:ustel_aile}
% =====================================================================

\begin{kesif}
Binom, Poisson, Normal, Gama, Beta — bu dağılımların PDF/PMF'leri çok farklı görünüyor.
Ama hepsinin $\exp(\ldots)$ biçiminde yazılabildiğini fark ettiniz mi?
Bu yapı tesadüf mü, yoksa derin bir ortak özelliği mi yansıtıyor?
\end{kesif}

\begin{tanim}[Tek Parametreli Üstel Aile]
\label{def:ustel_aile}
$\theta$'ya bağlı bir dağılım ailesi, eğer PDF/PMF'si
\[
  f(x \mid \theta) = h(x) \exp\!\bigl(\eta(\theta) T(x) - B(\theta)\bigr)
\]
biçiminde yazılabiliyorsa \textbf{(tek parametreli) üstel aile} (exponential family) üyesidir.
Burada:
\begin{itemize}
  \item $\eta(\theta)$: \textbf{doğal parametre} (natural parameter),
  \item $T(x)$: \textbf{yeterli istatistik} (sufficient statistic),
  \item $B(\theta)$: normalleştirme sabiti (log-bölüşüm fonksiyonu),
  \item $h(x)$: $\theta$'dan bağımsız taban ölçüsü.
\end{itemize}
\textbf{Çok parametreli genelleştirme:}
$f(x \mid \boldsymbol{\theta}) = h(x) \exp\!\bigl(\sum_{j=1}^k \eta_j(\boldsymbol{\theta}) T_j(x) - B(\boldsymbol{\theta})\bigr)$.
\end{tanim}

\begin{ornek}[Üstel Aile Örnekleri]
\begin{itemize}
  \item \textbf{Normal} ($\sigma^2$ biliniyor): $\eta(\mu) = \mu/\sigma^2$, $T(x) = x$, $B(\mu) = \mu^2/(2\sigma^2)$.
  \item \textbf{Poisson}: $\eta(\lambda) = \ln\lambda$, $T(k) = k$, $B(\lambda) = \lambda$, $h(k) = 1/k!$.
  \item \textbf{Binom} ($n$ sabit): $\eta(p) = \ln(p/q)$ (log-odds), $T(k) = k$.
  \item \textbf{Gama} ($\beta$ biliniyor): $\eta(\alpha) = \alpha-1$, $T(x) = \ln x$.
\end{itemize}
\end{ornek}

\begin{teorem}[Log-Bölüşüm Fonksiyonunun Türevleri]
\label{thm:log_bolüşüm}
Doğal parametre formu $f(x \mid \eta) = h(x)\exp(\eta T(x) - B(\eta))$ için:
\[
  E_\eta[T(X)] = B'(\eta), \qquad \mathrm{Var}_\eta(T(X)) = B''(\eta).
\]
\end{teorem}

\begin{proof}
$1 = \int h(x) e^{\eta T(x) - B(\eta)} dx$ her $\eta$ için.
$\eta$'ya göre türev: $0 = \int T(x) h(x) e^{\eta T - B} dx - B'(\eta) \cdot 1$,
yani $B'(\eta) = E[T(X)]$.
İkinci türev: $B''(\eta) = E[T(X)^2] - (E[T(X)])^2 = \mathrm{Var}(T(X))$.
\end{proof}

\begin{onerme}[Üstel Ailenin İstatistik Teorisindeki Önemi]
Üstel aile üyesi dağılımlar için $\sum_{i=1}^n T(X_i)$ yeterli istatistiktir
(Neyman faktorizasyon teoremi, Böl.~\ref{chp:nokta_tahmin}).
Bu dağılımlar UMVUE tahminlerine (Rao-Blackwell, Böl.~\ref{chp:nokta_tahmin}) sahiptir.
\end{onerme}

% =====================================================================
\section{Yer-Ölçek Aileleri}
\label{sec:yer_olcek}
% =====================================================================

\begin{tanim}[Yer-Ölçek Ailesi]
\label{def:yer_olcek}
$f_0$ bir "standart" PDF olsun. \textbf{Yer-ölçek ailesi}:
\[
  f(x \mid \mu, \sigma) = \frac{1}{\sigma} f_0\!\left(\frac{x-\mu}{\sigma}\right),
  \quad \mu \in \mathbb{R},\ \sigma > 0.
\]
$\mu$: yer parametresi (konum, location); $\sigma$: ölçek parametresi (scale).
\end{tanim}

\begin{ornek}[Yer-Ölçek Aileleri]
\begin{itemize}
  \item Normal: $f_0(z) = \phi(z)$; $\Normal(\mu, \sigma^2)$ yer-ölçek ailesi.
  \item Üstel (ölçek): $f_0(z) = e^{-z}$; $\Exp(\lambda) = \mathrm{Ölçek}(1/\lambda)$ ailesi.
  \item Cauchy: $f_0(z) = [\pi(1+z^2)]^{-1}$; $\mathrm{Cauchy}(\mu,\sigma)$ yer-ölçek.
  \item Logistik: $f_0(z) = e^{-z}/(1+e^{-z})^2$; lojistik regresyon temeli.
\end{itemize}
\end{ornek}

\begin{onerme}[Standartlaştırma]
$(X - \mu)/\sigma$ "standart" dağılıma sahip; yani yer-ölçek ailelerinde
tek bir "standart" tablo yeterlidir: $F(x \mid \mu, \sigma) = F_0((x-\mu)/\sigma)$.
\end{onerme}

% ---- Disiplinlerarası Bağlantı (TYMM) ----
\begin{kopru}[Disiplinlerarası — Biyoloji, Fizik, Mühendislik]
\textbf{Biyoloji:} Popülasyon büyüklüğü, hücre bölünme süreleri üstel ve gama
dağılımlarıyla modellenir; lognormal dağılım protein konsantrasyonlarını açıklar.\\
\textbf{Fizik:} Maxwell-Boltzmann hız dağılımı gama ailesine aittir;
termal gürültü normal dağılımdan kaynaklanır (Merkezi Limit Teoremi).\\
\textbf{Mühendislik/Güvenilirlik:} Weibull dağılımı (Gama genellemesi)
ürün ömrü modellemesinin standardıdır; üstel dağılım "bellek yoksunluğu"
sayesinde sıra kuyruk modellerinin temelidir.
\end{kopru}

% ---- Öz-Yansıtma ----
\begin{ozyansima}
\begin{itemize}[leftmargin=1.8em]
  \item Binom ile Poisson arasındaki "limit" ilişkisini kendi cümlelerinizle açıklayın.
        Hangi koşullar altında Poisson yaklaşımı makul bir seçimdir?
  \item Bellek yoksunluğu özelliği fiziksel olarak ne anlama gelir?
        Hangi gerçek hayat fenomenlerini modelleyebilir, hangilerini modelleyemez?
  \item Üstel aile kavramı neden istatistik teorisi açısından önemlidir?
  \item $t$ ve $F$ dağılımları neden "örneklem dağılımları" olarak anılır?
\end{itemize}
\end{ozyansima}

% ---- Köprü Notu ----
\begin{kopru}
\textbf{Önceki bölüm:} Bölüm~\ref{chp:beklenti} — beklenti, MÜF, KF araçları
buradaki her dağılım için somutlaştırıldı.\\
\textbf{Sonraki bölüm:} Bölüm~\ref{chp:buyuk_sayilar} — büyük sayılar kanunlarında
bu dağılımların asimptotik davranışı; Bölüm~\ref{chp:merkezi_limit}'te normal
dağılıma yakınsama.\\
\textbf{İstatistik kısmı:} Böl.~\ref{chp:nokta_tahmin}'de üstel aile ve yeterli
istatistik; Böl.~\ref{chp:hipotez_testi}'nde $\chi^2$, $t$, $F$ test istatistikleri.
\defterbaglantisi{bolum05\_dagilimlar.ipynb}{%
  Dağılım grafikleri, parametre duyarlılığı, simülasyon karşılaştırması}
\end{kopru}

% =====================================================================
%  ALIŞTIRMALAR
% =====================================================================

\cozumlualistirmalar

\begin{alistirma}[Al.5.1]{A}{Poisson Hesabı}
Birim saatte ortalama 5 arıza oluşan bir sistemde:
(a) Tam 3 arıza olasılığı. (b) En az 1 arıza olasılığı. (c) En fazla 2 arıza olasılığı.
\end{alistirma}
\alcozum{%
$X \sim \Pois(5)$.\\
(a) $P(X=3) = e^{-5}5^3/3! = 125e^{-5}/6 \approx 0.1404$.\\
(b) $P(X \geq 1) = 1 - P(X=0) = 1 - e^{-5} \approx 0.9933$.\\
(c) $P(X \leq 2) = e^{-5}(1 + 5 + 25/2) = 18.5e^{-5} \approx 0.1247$.}

\begin{alistirma}[Al.5.2]{A}{Normal Dağılım}
$X \sim \Normal(70, 100)$ (not: sınav puanı, $\sigma = 10$).
(a) $P(X > 85)$. (b) $P(60 < X < 80)$. (c) 90. yüzdeliği bulun.
\end{alistirma}
\alcozum{%
$Z = (X-70)/10 \sim \Normal(0,1)$.\\
(a) $P(X>85) = P(Z > 1.5) = 1 - \Phi(1.5) \approx 1 - 0.9332 = 0.0668$.\\
(b) $P(60<X<80) = P(-1 < Z < 1) = 2\Phi(1) - 1 \approx 2(0.8413) - 1 = 0.6827$.\\
(c) $\Phi(z_{0.9}) = 0.9 \Rightarrow z_{0.9} \approx 1.282$; $x_{0.9} = 70 + 10(1.282) = 82.82$.}

\begin{alistirma}[Al.5.3]{B}{Gama Dağılımı ve Beklenti}
$X \sim \Gama(3, 2)$.
(a) $E[X]$, $\mathrm{Var}(X)$ ve $M_X(t)$'yi yazın.
(b) $P(X > 3)$'ü sayısal ifade edin (yarı analitik form yeterli).
(c) $X_1, X_2 \sim \Gama(3, 2)$ bağımsız ise $X_1 + X_2$'nin dağılımını bulun.
\end{alistirma}
\alcozum{%
(a) $E[X] = 3/2$, $\mathrm{Var}(X) = 3/4$, $M_X(t) = (2/(2-t))^3$ ($t < 2$).\\
(b) $P(X>3) = \int_3^\infty \frac{2^3}{\Gamma(3)} x^2 e^{-2x} dx = 1 - F_{\Gama(3,2)}(3)$
(tamamlanmamış gama fonksiyonu).\\
(c) $X_1+X_2 \sim \Gama(6, 2)$.}

\begin{alistirma}[Al.5.4]{B}{Üstel Ailede Log-Bölüşüm}
$X \sim \Pois(\lambda)$'yı doğal parametre $\eta = \ln\lambda$ ile üstel aile formunda yazın.
$B(\eta)$'yı bulun ve $B'(\eta) = E[T(X)]$ formülüyle $E[X] = \lambda$'yı doğrulayın.
\end{alistirma}
\alcozum{%
$P(X=k) = \frac{\lambda^k e^{-\lambda}}{k!} = \frac{1}{k!} \exp(k\ln\lambda - \lambda)
= \frac{1}{k!} \exp(\eta k - e^\eta)$.\\
$T(k) = k$, $h(k) = 1/k!$, $B(\eta) = e^\eta = \lambda$.\\
$B'(\eta) = e^\eta = \lambda = E[T(X)] = E[X]$. \checkmark}

% ---

\alistirmalar

\begin{alistirma}[Al.5.5]{A}{Binom Olasılık}
10 soru içeren çoktan seçmeli testte her sorunun 4 şıkkı var; öğrenci tamamen rastgele cevaplıyor.
(a) Tam 3 doğru cevap olasılığı. (b) En az 5 doğru cevap olasılığı.
(c) $E[X]$ ve $\mathrm{Var}(X)$.
\end{alistirma}
\alcozum{%
$X \sim \Binom(10, 1/4)$.
(a) $\binom{10}{3}(1/4)^3(3/4)^7 \approx 0.2503$.
(b) $P(X\geq5) = \sum_{k=5}^{10}\binom{10}{k}(1/4)^k(3/4)^{10-k} \approx 0.0781$.
(c) $E[X] = 2.5$, $\mathrm{Var}(X) = 1.875$.}

\begin{alistirma}[Al.5.6]{A}{Geometrik Uygulama}
Bir üretim hattında her ürün $p = 0.05$ olasılıkla bozuk.
(a) İlk bozuk ürünün 10. üründe çıkma olasılığı.
(b) İlk bozuk üründen önce en az 20 ürün üretilme olasılığı.
(c) Kaçıncı üründe ilk bozuğun çıkması beklenir?
\end{alistirma}
\alcozum{%
$X \sim \mathrm{Geom}(0.05)$.
(a) $P(X=10) = (0.95)^9(0.05) \approx 0.0315$.
(b) $P(X>20) = (0.95)^{20} \approx 0.3585$.
(c) $E[X] = 1/0.05 = 20$.}

\begin{alistirma}[Al.5.7]{B}{Beta Dağılımı}
$X \sim \Beta(2, 3)$.
(a) $E[X]$ ve $\mathrm{Var}(X)$'i hesaplayın.
(b) $P(X > 0.5)$'i analitik olarak bulun.
(c) Bu dağılım hangi pratik durumu modelleyebilir?
\end{alistirma}
\alcozum{%
(a) $E[X] = 2/5 = 0.4$. $\mathrm{Var}(X) = 2\cdot3/(25\cdot6) = 6/150 = 1/25$.\\
(b) $f(x) = 12x(1-x)^2$.
$P(X>0.5) = 12\int_{0.5}^1 x(1-x)^2 dx = 12[x^2/2 - 2x^3/3 + x^4/4]_{0.5}^1
= 12[(1/2-2/3+1/4)-(1/8-1/12+1/16)] = 5/16$.\\
(c) $[0,1]$'de değer alan oranlar: dönüşüm hata oranı, bir proje tamamlanma yüzdesi, vb.}

\begin{alistirma}[Al.5.8]{B}{$\chi^2$ ve $t$ Dağılımı}
$X_1, \ldots, X_5$ bağımsız $\Normal(0,1)$.
(a) $V = \sum X_i^2$'nin dağılımını yazın; $E[V]$ ve $\mathrm{Var}(V)$'yi bulun.
(b) $Z \sim \Normal(0,1)$ bu $X_i$'lerden bağımsız. $T = Z/\sqrt{V/5}$'in dağılımını yazın.
(c) $T^2$'nin dağılımını yazın.
\end{alistirma}
\alcozum{%
(a) $V \sim \chi^2_5 = \Gama(5/2, 1/2)$. $E[V] = 5$, $\mathrm{Var}(V) = 10$.\\
(b) $T \sim t_5$.\\
(c) $T^2 \sim F(1, 5)$.}

\begin{alistirma}[Al.5.9]{B}{Poisson-Gama Karışımı}
$\Lambda \sim \Gama(\alpha, \beta)$ ve $X \mid \Lambda = \lambda \sim \Pois(\lambda)$.
$X$'in koşulsuz dağılımının negatif binom olduğunu gösterin.
\end{alistirma}
\alcozum{%
$P(X=k) = \int_0^\infty \frac{\lambda^k e^{-\lambda}}{k!} \cdot \frac{\beta^\alpha}{\Gamma(\alpha)}\lambda^{\alpha-1}e^{-\beta\lambda} d\lambda
= \frac{\beta^\alpha}{k!\,\Gamma(\alpha)} \int_0^\infty \lambda^{k+\alpha-1} e^{-(\beta+1)\lambda} d\lambda
= \frac{\beta^\alpha}{k!\,\Gamma(\alpha)} \cdot \frac{\Gamma(k+\alpha)}{(\beta+1)^{k+\alpha}}
= \binom{k+\alpha-1}{k} \left(\frac{\beta}{\beta+1}\right)^\alpha \left(\frac{1}{\beta+1}\right)^k$.
Bu $\mathrm{NegBin}(\alpha, \beta/(\beta+1))$ dağılımıdır.}

\begin{alistirma}[Al.5.10]{B}{Yer-Ölçek Dönüşümü}
$X \sim \mathrm{Cauchy}(0,1)$ ise $Y = \mu + \sigma X$ ($\sigma > 0$) için
$Y \sim \mathrm{Cauchy}(\mu, \sigma)$ olduğunu gösterin.
Cauchy dağılımı neden bir yer-ölçek ailesidir?
\end{alistirma}
\alcozum{%
$f_X(x) = 1/(\pi(1+x^2))$. Jacobian dönüşümü: $x = (y-\mu)/\sigma$, $|dx/dy| = 1/\sigma$.
$f_Y(y) = \frac{1}{\pi(1+((y-\mu)/\sigma)^2)} \cdot \frac{1}{\sigma}
= \frac{1}{\pi\sigma(1+((y-\mu)/\sigma)^2)}$. Bu $\mathrm{Cauchy}(\mu,\sigma)$ PDF'sidir.
Yer-ölçek ailesi: standart form $f_0(z) = 1/(\pi(1+z^2))$; $f(y|\mu,\sigma) = f_0((y-\mu)/\sigma)/\sigma$.}

\begin{alistirma}[Al.5.11]{C}{Normal Dağılımın Üstel Aile Formu}
$X \sim \Normal(\mu, \sigma^2)$'yi iki parametreli üstel ailede yazın:
$\boldsymbol{\eta} = (\eta_1, \eta_2)$, $\mathbf{T}(x) = (T_1(x), T_2(x))$ ile.
$B(\boldsymbol{\eta})$'yı bulun ve $\partial B/\partial\eta_1 = E[T_1(X)]$,
$\partial B/\partial\eta_2 = E[T_2(X)]$ formüllerini doğrulayın.
\end{alistirma}
\alcozum{%
$f(x|\mu,\sigma^2) = \frac{1}{\sqrt{2\pi}\sigma}\exp\!\left(-\frac{x^2}{2\sigma^2} + \frac{\mu}{\sigma^2}x - \frac{\mu^2}{2\sigma^2}\right)$.
Doğal parametreler: $\eta_1 = \mu/\sigma^2$, $\eta_2 = -1/(2\sigma^2)$ (yani $\sigma^2 = -1/(2\eta_2)$, $\mu = \eta_1/(-2\eta_2)$).
$T_1(x) = x$, $T_2(x) = x^2$.
$B(\boldsymbol{\eta}) = -\eta_1^2/(4\eta_2) - \frac{1}{2}\ln(-2\eta_2) + \frac{1}{2}\ln(2\pi)$.
$\partial B/\partial\eta_1 = -\eta_1/(2\eta_2) = \mu/\sigma^2 \cdot \sigma^2 = \mu = E[X]$. \checkmark
$\partial B/\partial\eta_2 = \eta_1^2/(4\eta_2^2) - 1/(2\eta_2) = \mu^2 + \sigma^2 = E[X^2]$. \checkmark}

\begin{alistirma}[Al.5.12]{C}{Binom $\to$ Poisson: Tam İspat}
Teorem~\ref{thm:binom_poisson}'in ispatında kullanılan iki limiti kesin biçimde ispatlayın:
(a) $n(n-1)\cdots(n-k+1)/n^k \to 1$.
(b) $(1 - \lambda_n/n)^n \to e^{-\lambda}$ ($\lambda_n \to \lambda$).
\end{alistirma}
\alcozum{%
(a) $\frac{n(n-1)\cdots(n-k+1)}{n^k} = 1 \cdot (1-1/n) \cdots (1-(k-1)/n) \to 1^k = 1$ (sabit $k$, $n\to\infty$).\\
(b) $\ln(1-\lambda_n/n)^n = n\ln(1-\lambda_n/n)$. Taylor: $\ln(1-x) = -x - x^2/2 - \cdots$,
$n \cdot (-\lambda_n/n + O(1/n^2)) = -\lambda_n + O(1/n) \to -\lambda$.
Dolayısıyla $(1-\lambda_n/n)^n \to e^{-\lambda}$.}

% =====================================================================
\endinput

% =====================================================================
%  BÖLÜM 6: Dönüşümler ve Dağılım Teorisi
%  Kaynak: Hogg & Craig; Casella & Berger; Feller Vol.II; Klenke
% =====================================================================
\chapter{Dönüşümler ve Dağılım Teorisi}
\label{chp:donusumler}

\bolumalinti{%
  Bir rasgele değişkenin dönüşümü, dağılım uzayındaki bir yolculuktur.%
}{George Casella \& Roger Berger, \textit{Statistical Inference}}

% ---- Kazanımlar (TYMM) ----
\begin{kazanimlar}
Bu bölümün sonunda öğrenci:
\begin{itemize}[leftmargin=1.8em]
  \item Tekil ve çok boyutlu dönüşümlerde Jacobian formülünü ustalıkla uygulayabilir;
        birebir olmayan dönüşümleri parçalara bölerek ele alabilir.
  \item MÜF ve karakteristik fonksiyon yöntemiyle dağılım tanımlayabilir;
        bu yöntemin Jacobian yöntemine göre avantajlarını açıklayabilir.
  \item Evrişim (konvolüsyon) integralini hesaplayabilir; Gama, Normal ve Poisson
        toplamlarının dağılımını türetebilir.
  \item Sıra istatistiklerini tanımlayabilir; tekil ve ortak PDF'lerini türetip
        hesaplayabilir; örneklem minimum, maksimum ve medyanını bulabilir.
  \item Delta yöntemini tek ve çok boyutlu durumda uygulayabilir.
\end{itemize}
\end{kazanimlar}

% ---- Motivasyon ----
\begin{motivasyon}
Bölüm~\ref{chp:rasgele_degiskenler}'de dönüşüm yönteminin temellerini gördük.
Bu bölümde araç kutusunu genişletiyoruz. Çoğu zaman Jacobian hesabı yerine
MÜF ya da KF yöntemi çok daha hızlı sonuç verir: $X+Y$'nin dağılımını bulmak için
evrişim integralini hesaplamak yerine $M_{X+Y}(t) = M_X(t)M_Y(t)$ özdeşliğini
kullanmak yeterlidir. Sıra istatistikleri ise güvenilirlik mühendisliği ve
çevresel istatistiğin temel aracıdır: "en zayıf halka ne zaman kopar?"

\medskip
\textbf{Köprü.} Bölüm~\ref{chp:rasgele_degiskenler}'deki Jacobian ve evrişim,
Bölüm~\ref{chp:beklenti}'deki MÜF/KF bu bölümde birleştiriliyor.
Bölüm~\ref{chp:merkezi_limit}'teki Delta yöntemi bu bölümün doğal devamıdır.
\end{motivasyon}

% =====================================================================
\section{Tekil Dönüşümler: Derinlemesine}
\label{sec:tekil_donusum}
% =====================================================================

Bölüm~\ref{sec:rd_donusum}'da birebir dönüşüm formülünü ve genel monoton-olmayan
durumu gördük. Burada bu tekniği pekiştiriyor, sınır durumlarını ve pratikte
sık karşılaşılan kalıpları inceliyoruz.

\begin{teorem}[Tekil Dönüşüm — Kapsamlı Form]
\label{thm:tekil_donusum}
$X$ sürekli rasgele değişken, PDF $f_X$. $g : \mathcal{X} \to \mathcal{Y}$
($\mathcal{X}$ ve $\mathcal{Y}$ açık aralıklar birleşimi). $\mathcal{X} = \bigsqcup_k A_k$
ayrışımında $g|_{A_k}$ her parça üzerinde birebir, türevlenebilir olsun.
$h_k = (g|_{A_k})^{-1}$ ters fonksiyon. O zaman $Y = g(X)$'in PDF'si:
\[
  f_Y(y) = \sum_k f_X(h_k(y)) \cdot |h_k'(y)| \cdot \mathbf{1}_{y \in g(A_k)}.
\]
\end{teorem}

\begin{ornek}[$Y = |X|$, $X \sim \Normal(0,1)$]
$g(x) = |x|$: $A_1 = (-\infty, 0)$, $A_2 = (0, \infty)$.
$h_1(y) = -y$, $h_2(y) = y$; $|h_1'| = |h_2'| = 1$.
$f_Y(y) = \phi(-y) + \phi(y) = 2\phi(y)$, $y > 0$ (yarı-normal dağılım).
$P(|Z| > 1.96) = 2(1 - \Phi(1.96)) \approx 0.05$ — iki-kuyruklu normal tablosu.
\end{ornek}

\begin{ornek}[$Y = X^2$, $X$ simetrik]
$X$ simetrik (çift PDF) olsun. $A_1 = (-\infty,0)$, $A_2 = (0,\infty)$.
$h_1(y) = -\sqrt{y}$, $h_2(y) = \sqrt{y}$; $|h_{1,2}'(y)| = 1/(2\sqrt{y})$.
$f_Y(y) = [f_X(\sqrt{y}) + f_X(-\sqrt{y})]/(2\sqrt{y}) = f_X(\sqrt{y})/\sqrt{y}$ (simetri).
$Z \sim \Normal(0,1)$: $f_{Z^2}(y) = \phi(\sqrt{y})/\sqrt{y} = (2\pi y)^{-1/2}e^{-y/2}$
$= \Gama(1/2, 1/2)$ PDF'si $= \chi^2_1$ PDF'si.
\end{ornek}

\begin{ispatiskele}[Quantile Dönüşümü — İspat İskeleti]
$U \sim \Uniform(0,1)$, $F$ herhangi KDF, $X = F^{-1}(U)$.
\iskeleadim{1}{$P(X \leq x) = P(F^{-1}(U) \leq x)$ ifadesini yazın.}
\iskeleadim{2}{$F^{-1}(u) \leq x \Leftrightarrow u \leq F(x)$ eşdeğerliğini kullanın ($F$ sağdan sürekli).}
\iskeleadim{3}{$P(U \leq F(x)) = \bosluk{F(x)}$ (uniform KDF).}
\iskeleadim{4}{Sonuç: $F_X(x) = F(x)$, yani $X \sim F$.}
\end{ispatiskele}

% =====================================================================
\section{Çok Boyutlu Dönüşümler}
\label{sec:cok_boyut_donusum}
% =====================================================================

Bölüm~\ref{subsec:donusum_jacobian}'da Jacobian formülünü verdik.
Burada sistematik uygulama tekniklerini ve önemli örnekleri işliyoruz.

\begin{onerme}[Jacobian Hesabında Adım Adım Yöntem]
\label{prop:jacobian_adim}
$(X_1, \ldots, X_n) \to (Y_1, \ldots, Y_n)$ dönüşümü için:
\begin{enumerate}[(1)]
  \item Ters dönüşümü bulun: $x_i = h_i(y_1,\ldots,y_n)$.
  \item Jacobian matrisini oluşturun: $J_{ij} = \partial h_i/\partial y_j$.
  \item $|\det J|$'yi hesaplayın.
  \item Tanım bölgesini dönüştürün: $\mathcal{X} \to \mathcal{Y}$.
  \item $f_Y(\mathbf{y}) = f_X(h(\mathbf{y})) \cdot |\det J|$ ile PDF'yi yazın.
\end{enumerate}
\end{onerme}

\begin{ornek}[Beta-Gama İlişkisi — Tam Türetme]
\label{ex:beta_gama_turetme}
$U \sim \Gama(\alpha, 1)$, $V \sim \Gama(\beta, 1)$ bağımsız.
$X = U/(U+V)$, $W = U+V$ dönüşümü.

\textbf{Ters:} $u = xw$, $v = w(1-x)$; $0 < x < 1$, $w > 0$.

\textbf{Jacobian:}
$J = \begin{pmatrix} w & x \\ -w & 1-x \end{pmatrix}$; $|\det J| = w(1-x) + xw = w$.

\textbf{Ortak PDF:}
\begin{align*}
  f_{X,W}(x,w) &= f_U(xw)\,f_V(w(1-x))\cdot w \\
  &= \frac{(xw)^{\alpha-1}e^{-xw}}{\Gamma(\alpha)} \cdot \frac{(w(1-x))^{\beta-1}e^{-w(1-x)}}{\Gamma(\beta)} \cdot w \\
  &= \frac{x^{\alpha-1}(1-x)^{\beta-1}}{B(\alpha,\beta)} \cdot \frac{w^{\alpha+\beta-1}e^{-w}}{\Gamma(\alpha+\beta)}.
\end{align*}

$X$ ve $W$ bağımsız; $X \sim \Beta(\alpha,\beta)$, $W \sim \Gama(\alpha+\beta, 1)$.
\end{ornek}

\begin{ornek}[$t$ Dağılımının Türetilmesi]
\label{ex:t_turetme}
$Z \sim \Normal(0,1)$, $V \sim \chi^2_n$ bağımsız. $T = Z/\sqrt{V/n}$, $W = V$.

\textbf{Ters:} $z = t\sqrt{w/n}$, $v = w$.
\textbf{Jacobian:} $|\det J| = \sqrt{w/n}$.

\textbf{Ortak PDF:}
$f_{T,W}(t,w) = \phi(t\sqrt{w/n}) \cdot f_{\chi^2_n}(w) \cdot \sqrt{w/n}$.

$W$'ye göre integre et ($w > 0$):
\[
  f_T(t) = \frac{\Gamma((n+1)/2)}{\sqrt{n\pi}\,\Gamma(n/2)} \left(1 + \frac{t^2}{n}\right)^{-(n+1)/2}.
\]
(Gama integrali $\int_0^\infty w^{(n-1)/2} e^{-w(1+t^2/n)/2} dw = \Gamma((n+1)/2) \cdot [2/(1+t^2/n)]^{(n+1)/2}$
kullanılır.)
\end{ornek}

\begin{ornek}[$F$ Dağılımının Türetilmesi]
$U \sim \chi^2_m$, $V \sim \chi^2_n$ bağımsız. $W = (U/m)/(V/n)$ için:
$f_W(w) = \frac{\Gamma((m+n)/2)}{\Gamma(m/2)\Gamma(n/2)} \left(\frac{m}{n}\right)^{m/2}
w^{m/2-1}\left(1 + \frac{m}{n}w\right)^{-(m+n)/2}$, $w > 0$.
Bu $F(m,n)$ dağılımının PDF'sidir.
\end{ornek}

% =====================================================================
\section{MÜF ve KF Yöntemiyle Dağılım Belirleme}
\label{sec:muf_yontemi}
% =====================================================================

\begin{motivasyon}[Jacobian Yerine MÜF]
$X+Y$'nin PDF'sini bulmak için ya evrişim integrali ya da MÜF çarpımı kullanılır.
MÜF yöntemi genellikle çok daha hızlıdır: toplam PDF'si yerine MÜF bulunur,
bu standart bir dağılıma karşılık geliyorsa PDF'yi yazmaya gerek kalmaz.
\end{motivasyon}

\begin{teorem}[MÜF ile Dağılım Tanımlama]
\label{thm:muf_tanim}
$Y = g(X_1, \ldots, X_n)$ ve $M_Y(t) = E[e^{tY}]$ hesaplanabiliyorsa:
$M_Y(t)$ bilinen bir dağılımın MÜF'üne eşit ise $Y$ o dağılımı izler
(MÜF tek türlü dağılım belirler, Teorem~\ref{thm:muf_tek}).
\end{teorem}

\begin{ornek}[Normal Toplamı — MÜF ile]
$X_i \sim \Normal(\mu_i, \sigma_i^2)$ bağımsız ($i=1,\ldots,n$), $a_i \in \mathbb{R}$.
\[
  M_{\sum a_i X_i}(t) = \prod_i M_{X_i}(a_i t) = \prod_i e^{\mu_i a_i t + \sigma_i^2 a_i^2 t^2/2}
  = \exp\!\left(\sum a_i\mu_i \cdot t + \frac{t^2}{2}\sum a_i^2\sigma_i^2\right).
\]
Bu $\Normal(\sum a_i\mu_i, \sum a_i^2\sigma_i^2)$ MÜF'üdür. Jacobian veya evrişim gerekmez.
\end{ornek}

\begin{ornek}[Gama Toplamı — MÜF ile]
$X_i \sim \Gama(\alpha_i, \beta)$ bağımsız ($\beta$ aynı):
$M_{\sum X_i}(t) = \prod (\beta/(\beta-t))^{\alpha_i} = (\beta/(\beta-t))^{\sum\alpha_i}$.
$\Rightarrow$ $\sum X_i \sim \Gama(\sum\alpha_i, \beta)$.

Özel durum: $Z_i \sim \Normal(0,1)$ bağımsız $\Rightarrow$ $Z_i^2 \sim \chi^2_1 = \Gama(1/2,1/2)$
$\Rightarrow$ $\sum_{i=1}^n Z_i^2 \sim \Gama(n/2,1/2) = \chi^2_n$.
\end{ornek}

\begin{ornek}[Poisson Toplamı — KF ile]
$X_i \sim \Pois(\lambda_i)$ bağımsız:
$\varphi_{\sum X_i}(t) = \prod e^{\lambda_i(e^{it}-1)} = \exp\!\bigl((\sum\lambda_i)(e^{it}-1)\bigr)$.
$\Rightarrow$ $\sum X_i \sim \Pois(\sum\lambda_i)$.
\end{ornek}

\begin{teorem}[KF ile Dağılımla Yakınsama]
\label{thm:kf_yaklinsama}
$\varphi_{X_n}(t) \to \varphi_X(t)$ her $t$'de ise $X_n \xrightarrow{d} X$
(Lévy süreklilik teoremi; ispat Bölüm~\ref{chp:merkezi_limit}'te kullanılacak).
\end{teorem}

% =====================================================================
\section{Evrişim}
\label{sec:evrisim}
% =====================================================================

Bölüm~\ref{subsec:evrisim}'de evrişim formülünü tanımladık. Burada hesaplama
tekniklerini ve önemli örnekleri derinleştiriyoruz.

\begin{teorem}[Evrişim Özellikleri]
\label{thm:evrisim_ozellik}
$f, g, h$ PDF'ler olmak üzere:
\begin{enumerate}[(i)]
  \item \textbf{Değişmelilik:} $f * g = g * f$.
  \item \textbf{Birleşmelilik:} $(f * g) * h = f * (g * h)$.
  \item \textbf{KF çarpımı:} $\varphi_{X+Y}(t) = \varphi_X(t)\,\varphi_Y(t)$.
  \item \textbf{MÜF çarpımı:} $M_{X+Y}(t) = M_X(t)\,M_Y(t)$ (bağımsız $X,Y$).
\end{enumerate}
\end{teorem}

\begin{ornek}[Üçgen Dağılım — Evrişim]
$X, Y \sim \Uniform(0,1)$ bağımsız, $Z = X + Y$:
\[
  f_Z(z) = \int_0^1 \mathbf{1}_{[0,1]}(z-y)\,dy
  = \begin{cases} z & 0 \leq z \leq 1 \\ 2-z & 1 < z \leq 2 \end{cases}.
\]
Üçgen dağılım (Bölüm~\ref{sec:buyuk_sayilar}'daki Al.7'nin açılımı).
$n$ bağımsız $\Uniform(0,1)$ toplamı: Irwin-Hall dağılımı.
$n \to \infty$: Merkezi Limit Teoremi'ne hazırlık.
\end{ornek}

\begin{ornek}[Normal Evrişim — İntegral Yöntemi]
$X \sim \Normal(\mu_1, \sigma_1^2)$, $Y \sim \Normal(\mu_2, \sigma_2^2)$ bağımsız.
Evrişim:
\[
  f_{X+Y}(z) = \int_{-\infty}^{+\infty} \frac{1}{\sigma_1\sqrt{2\pi}}e^{-(z-y-\mu_1)^2/(2\sigma_1^2)}
  \cdot \frac{1}{\sigma_2\sqrt{2\pi}}e^{-(y-\mu_2)^2/(2\sigma_2^2)} dy.
\]
Üstel ifade birleştirilip $y$'deki kareyi tamamlama:
$(z-y-\mu_1)^2/\sigma_1^2 + (y-\mu_2)^2/\sigma_2^2$'nin $y$ üzerinde integral
$\sigma_1^2+\sigma_2^2$ varyanslı Gauss integraline döner.
Sonuç: $X+Y \sim \Normal(\mu_1+\mu_2, \sigma_1^2+\sigma_2^2)$.
(MÜF yöntemi bu hesabı beş satıra indirger.)
\end{ornek}

\begin{onerme}[Cauchy Toplamı]
$X, Y \sim \mathrm{Cauchy}(0,1)$ bağımsız $\Rightarrow$ $X+Y \sim \mathrm{Cauchy}(0,2)$.
\end{onerme}

\begin{proof}
KF: $\varphi_{X+Y}(t) = e^{-|t|} \cdot e^{-|t|} = e^{-2|t|}$ — bu $\mathrm{Cauchy}(0,2)$'nin KF'sidir.
\end{proof}

\begin{kritiksorgu}
Cauchy toplamı neden $\mathrm{Cauchy}(0,2)$, "beklentiden" elde ettiğimiz $\mathrm{Cauchy}(0,1)$ değil?
Cauchy'nin beklentisi mevcut değildir — standartlaştırma ($n$'e bölme) dağılımı stabilize etmez.
Büyük Sayılar Kanunu çöker; Merkezi Limit Teoremi uygulanmaz.
\end{kritiksorgu}

% =====================================================================
\section{Sıra İstatistikleri}
\label{sec:sira_istatistikleri}
% =====================================================================

\begin{tanim}[Sıra İstatistikleri]
\label{def:sira_istatistik}
$X_1, \ldots, X_n$ i.i.d., ortak PDF $f$, KDF $F$.
Küçükten büyüğe sıralanmış değerler:
\[
  X_{(1)} \leq X_{(2)} \leq \cdots \leq X_{(n)}.
\]
$X_{(k)}$: $k$'ıncı \textbf{sıra istatistiği} (order statistic).
$X_{(1)} = \min$, $X_{(n)} = \max$, $X_{(\lceil n/2 \rceil)}$: medyan.
\end{tanim}

\begin{teorem}[Tekil Sıra İstatistiğinin PDF'si]
\label{thm:sira_pdf}
$k$'ıncı sıra istatistiğinin PDF'si:
\[
  f_{X_{(k)}}(x) = \frac{n!}{(k-1)!(n-k)!} [F(x)]^{k-1}[1-F(x)]^{n-k} f(x).
\]
\end{teorem}

\begin{proof}
Kombinatoryal argüman: $X_{(k)} \in (x, x+dx)$ olması için
$k-1$ gözlem $(-\infty, x)$'de, 1 gözlem $(x, x+dx)$'de, $n-k$ gözlem $(x+dx, \infty)$'da olmalı.
Çok terimli katsayı $\frac{n!}{(k-1)!\,1!\,(n-k)!}$ ile olasılık:
$[F(x)]^{k-1} \cdot f(x)\,dx \cdot [1-F(x)]^{n-k}$.
$dx$'e bölerek PDF elde edilir.
\end{proof}

\begin{ornek}[Minimum ve Maksimum]
\textbf{Minimum} ($k=1$):
$f_{X_{(1)}}(x) = n[1-F(x)]^{n-1}f(x)$.
$F_{X_{(1)}}(x) = 1 - [1-F(x)]^n$.

\textbf{Maksimum} ($k=n$):
$f_{X_{(n)}}(x) = n[F(x)]^{n-1}f(x)$.
$F_{X_{(n)}}(x) = [F(x)]^n$.

Bu formüller hemen doğrulanabilir:
$P(X_{(1)} > x) = P(\text{tüm } X_i > x) = [1-F(x)]^n$ (bağımsızlık).
$P(X_{(n)} \leq x) = P(\text{tüm } X_i \leq x) = [F(x)]^n$.
\end{ornek}

\begin{teorem}[Ortak PDF]
\label{thm:sira_ortak_pdf}
$(X_{(1)}, \ldots, X_{(n)})$'nin ortak PDF'si:
\[
  f_{X_{(1)},\ldots,X_{(n)}}(x_1, \ldots, x_n) = n! \prod_{i=1}^n f(x_i), \quad x_1 < x_2 < \cdots < x_n.
\]
\end{teorem}

\begin{proof}
$n!$ adet permütasyon; her biri $(x_1, \ldots, x_n)$ sıralamasını $x_{\sigma(1)} < \cdots < x_{\sigma(n)}$
ile üretir. Bağımsızlıktan her permütasyonun katkısı $\prod f(x_i)$; toplam $n! \prod f(x_i)$.
\end{proof}

\begin{teorem}[İkili Marjinal PDF]
\label{thm:sira_ikili}
$j < k$ için $(X_{(j)}, X_{(k)})$'nın ortak PDF'si:
\[
  f_{X_{(j)},X_{(k)}}(u,v) = \frac{n!}{(j-1)!(k-j-1)!(n-k)!}
  [F(u)]^{j-1}[F(v)-F(u)]^{k-j-1}[1-F(v)]^{n-k}f(u)f(v),
\]
$u < v$.
\end{teorem}

\begin{ornek}[$\Uniform(0,1)$ Sıra İstatistikleri]
$X_1, \ldots, X_n \sim \Uniform(0,1)$ i.i.d.; $f(x) = 1$, $F(x) = x$ ($0<x<1$).

$k$'ıncı sıra istatistiği:
$f_{X_{(k)}}(x) = \frac{n!}{(k-1)!(n-k)!} x^{k-1}(1-x)^{n-k} = \frac{x^{k-1}(1-x)^{n-k}}{B(k, n-k+1)}$.

$X_{(k)} \sim \Beta(k, n-k+1)$.
$E[X_{(k)}] = k/(n+1)$: $k$'ıncı sıra istatistiğinin beklentisi $[0,1]$'i eşit böler.
\end{ornek}

\begin{ornek}[Üstel Sıra İstatistikleri — Yarışma Süresi]
$X_i \sim \Exp(\lambda)$ bağımsız (cihaz ömürleri). $X_{(1)} = \min X_i$:
$P(X_{(1)} > t) = [P(X_i > t)]^n = e^{-n\lambda t}$.
$X_{(1)} \sim \Exp(n\lambda)$: minimum de üstel, oran $n$ kat büyük.

Sıra arası boşluklar: $D_k = X_{(k)} - X_{(k-1)}$ bağımsız,
$D_k \sim \Exp((n-k+1)\lambda)$ (eksponensiyal gecikme özelliği).
Bu, güvenilirlik mühendisliğinde yaygın kullanılan "rekabetçi hayatta kalma" modelidir.
\end{ornek}

\begin{hataanalizi}
"$X_{(k)}$ ile $X_{(j)}$ bağımsızdır" — yanlış.
Ortak PDF $f_{X_{(j)},X_{(k)}}(u,v) \neq f_{X_{(j)}}(u) f_{X_{(k)}}(v)$
(bölge $u < v$ koşulu bağımlılık yaratır).
Sıra istatistikleri $\mathbf{bağımlıdır}$; ikili ve yüksek dereceli dağılımlar ayrıca hesaplanmalıdır.
\end{hataanalizi}

\begin{mikroozet}
$f_{X_{(k)}}(x) = \frac{n!}{(k-1)!(n-k)!}F^{k-1}(1-F)^{n-k}f$.
Uniform: $X_{(k)} \sim \Beta(k, n-k+1)$.
Üstel: minimum de üstel; aralıklar bağımsız üstel.
Sıra istatistikleri genel olarak bağımlıdır.
\end{mikroozet}

% =====================================================================
\section{Delta Yöntemi}
\label{sec:delta_yontemi}
% =====================================================================

\begin{motivasyon}[Neden Delta Yöntemi?]
$\bar{X}_n \xrightarrow{d} \Normal(\mu, \sigma^2/n)$ — bu Merkezi Limit Teoremi'nin (bir sonraki
bölüm) sonucudur. Peki $g(\bar{X}_n)$ nasıl dağılır? $g$ doğrusal değilse doğrudan
dönüşüm uygulanamaz; ancak Taylor açılımı bu boşluğu kapatır.
\end{motivasyon}

\begin{teorem}[Delta Yöntemi — Tek Boyutlu]
\label{thm:delta_yontemi}
$\sqrt{n}(X_n - \mu) \xrightarrow{d} \Normal(0, \sigma^2)$ ve $g$ $\mu$'da türevlenebilir,
$g'(\mu) \neq 0$ ise:
\[
  \sqrt{n}(g(X_n) - g(\mu)) \xrightarrow{d} \Normal(0,\, \sigma^2 [g'(\mu)]^2).
\]
\end{teorem}

\begin{proof}
Taylor: $g(X_n) = g(\mu) + g'(\mu)(X_n-\mu) + R_n$ ($R_n = o(X_n-\mu)$).
$\sqrt{n}(g(X_n)-g(\mu)) = g'(\mu)\sqrt{n}(X_n-\mu) + \sqrt{n}R_n$.
$X_n \xrightarrow{P} \mu$ ($\sqrt{n}$-normalite bunu içerir) $\Rightarrow$ $\sqrt{n}R_n \xrightarrow{P} 0$.
Slutsky (Teorem~\ref{thm:slutsky}): $g'(\mu)\sqrt{n}(X_n-\mu) \xrightarrow{d} g'(\mu)\Normal(0,\sigma^2)
= \Normal(0, \sigma^2[g'(\mu)]^2)$.
\end{proof}

\begin{ornek}[Delta Yöntemi — Logaritma]
$X_i \sim \mathrm{Exp}(\lambda)$, $n$ büyük. MLT: $\sqrt{n}(\bar{X}_n - 1/\lambda) \xrightarrow{d} \Normal(0, 1/\lambda^2)$.
$g(x) = \ln x$, $g'(x) = 1/x$, $g'(1/\lambda) = \lambda$.
Delta yöntemi:
\[
  \sqrt{n}(\ln\bar{X}_n - \ln(1/\lambda)) \xrightarrow{d} \Normal(0, \lambda^2 \cdot 1/\lambda^2) = \Normal(0,1).
\]
\end{ornek}

\begin{teorem}[Çok Boyutlu Delta Yöntemi]
\label{thm:delta_cok}
$\sqrt{n}(\mathbf{X}_n - \boldsymbol{\mu}) \xrightarrow{d} \Normal_k(\mathbf{0}, \boldsymbol{\Sigma})$
ve $g : \mathbb{R}^k \to \mathbb{R}$ $\boldsymbol{\mu}$'da türevlenebilir ise:
\[
  \sqrt{n}(g(\mathbf{X}_n) - g(\boldsymbol{\mu})) \xrightarrow{d}
  \Normal\!\left(0,\, \nabla g(\boldsymbol{\mu})^\top \boldsymbol{\Sigma}\, \nabla g(\boldsymbol{\mu})\right).
\]
\end{teorem}

\begin{ornek}[Oran Tahmincisi]
$(X_i, Y_i)$ i.i.d., $E[X] = \mu_X$, $E[Y] = \mu_Y > 0$. $R_n = \bar{X}_n/\bar{Y}_n$.
$g(x,y) = x/y$; $\nabla g = (1/\mu_Y,\, -\mu_X/\mu_Y^2)^\top$.
Çok boyutlu delta:
\[
  \sqrt{n}(R_n - \mu_X/\mu_Y) \xrightarrow{d} \Normal\!\left(0,\,
  \frac{\sigma_X^2}{\mu_Y^2} - \frac{2\mu_X\sigma_{XY}}{\mu_Y^3} + \frac{\mu_X^2\sigma_Y^2}{\mu_Y^4}
  \right).
\]
\end{ornek}

% =====================================================================
\section{Karakteristik Fonksiyon: Tersçevirme ve Benzersizlik}
\label{sec:kf_ters}
% =====================================================================

\begin{motivasyon}[KF Neden Bu Kadar Güçlü?]
Moment üreten fonksiyon her dağılım için mevcut olmayabilir (Cauchy).
Karakteristik fonksiyon $\varphi_X(t) = E[e^{itX}]$ her dağılım için vardır.
Dahası, KF dağılımı \emph{tam olarak} belirler: iki rasgele değişkenin KF'leri özdeşse
dağılımları da özdeştir. Tersçevirme formülü, KF'den PDF'ye geçiş sağlar.
\end{motivasyon}

\begin{teorem}[KF Benzersizlik Teoremi]
\label{thm:kf_benzersizlik}
$X$ ve $Y$ rasgele değişkenler. Eğer $\varphi_X(t) = \varphi_Y(t)$ tüm $t \in \mathbb{R}$ için
geçerliyse $X \overset{d}{=} Y$ (aynı dağılım).
\end{teorem}

\begin{proof}[İspat Taslağı]
$F_X = F_Y$ olduğu yeter. Herhangi $a < b$ için
\[
  P(a < X \leq b) = \lim_{T\to\infty} \frac{1}{2\pi} \int_{-T}^T \frac{e^{-ita}-e^{-itb}}{it}\varphi_X(t)\,dt,
\]
(Lévy tersçevirme formülü). Aynı formül $F_Y$ için de geçerli; $\varphi_X = \varphi_Y$ koşuluyla
$F_X(b)-F_X(a) = F_Y(b)-F_Y(a)$ her $a < b$ için. $F_X = F_Y$ çıkar.
\end{proof}

\begin{teorem}[Tersçevirme Formülü]
\label{thm:kf_terscevir}
$X$ sürekli rasgele değişken, $\int_{-\infty}^{+\infty} |\varphi_X(t)|\,dt < \infty$ ise
\[
  f_X(x) = \frac{1}{2\pi} \int_{-\infty}^{+\infty} e^{-itx}\,\varphi_X(t)\,dt.
\]
Yani KF, PDF'nin Fourier dönüşümüdür; tersçevirme Fourier ters dönüşümüdür.
\end{teorem}

\begin{ornek}[Normal KF'sinden PDF'ye]
$\varphi_X(t) = e^{i\mu t - \sigma^2 t^2/2}$ (Normal KF).
$\int |\varphi_X(t)|\,dt = \int e^{-\sigma^2 t^2/2}\,dt = \sqrt{2\pi}/\sigma < \infty$.
Tersçevirme:
\[
  f_X(x) = \frac{1}{2\pi}\int_{-\infty}^{+\infty} e^{-itx} e^{i\mu t - \sigma^2 t^2/2}\,dt
  = \frac{1}{\sigma\sqrt{2\pi}} e^{-(x-\mu)^2/(2\sigma^2)}.
\]
(Gauss integralini tamamlayarak $t$'deki kareyi tamamlama kullanılır.)
\end{ornek}

\begin{teorem}[Lévy Süreklilik Teoremi]
\label{thm:levy_sureklilik}
$\varphi_{X_n}(t) \to \varphi(t)$ her $t$'de, ve $\varphi$ $t=0$'da sürekli ise
$X_n \xrightarrow{d} X$ burada $X$'in KF'si $\varphi$'dir.
\end{teorem}

\begin{proof}[İspat Notu]
Temel fikir: KDF dizisinin görece kompaktlığı (Helly seçim teoremi) + KF yakınsamasının
limit ölçümünü tam belirlemesi. Tam ispat için Billingsley~\cite{Billingsley1995}, §26.
Bu teorem MLT'nin KF ile ispatının kalbini oluşturur (Bölüm~\ref{chp:merkezi_limit}).
\end{proof}

\begin{hataanalizi}
"$\varphi_{X_n}(t) \to \varphi(t)$ her $t$'de, dolayısıyla $X_n \xrightarrow{d} X$."
Bu çıkarım Lévy Sürekliliği'ni atlıyor: $\varphi$'nin $t=0$'da sürekli (yani geçerli bir KF) olması
gerekir. $\varphi_n(t) = e^{-n|t|}$ ise limit $\varphi(0)=1$, $\varphi(t)=0$ ($t\neq0$) —
bu bir KF değildir. Sonuç: $X_n \to \infty$ dağılıma (mass escape), limit dağılım mevcut değil.
\end{hataanalizi}

\begin{mikroozet}
KF her dağılım için vardır; benzersiz olarak belirler.
Tersçevirme: $f_X = (2\pi)^{-1}\int e^{-itx}\varphi_X(t)\,dt$ ($|\varphi|$ integrallenebilir ise).
Lévy Sürekliliği: KF yakınsaması $\Leftrightarrow$ dağılımla yakınsama.
\end{mikroozet}

% =====================================================================
\section{İkinci Dereceden Delta Yöntemi}
\label{sec:delta2}
% =====================================================================

\begin{motivasyon}[Birinci Türev Sıfır Olduğunda]
Delta yöntemi $g'(\mu) \neq 0$ gerektirir. $g'(\mu) = 0$ ise $\sqrt{n}(g(\bar{X}_n) - g(\mu))$
sıfıra yakınsar: limit dağılım dejenere. Doğru ölçek $n$'dir; limit dağılım normal değil,
ki-kare ile ilgilidir.
\end{motivasyon}

\begin{teorem}[İkinci Dereceden Delta Yöntemi]
\label{thm:delta2}
$\sqrt{n}(X_n - \mu) \xrightarrow{d} \Normal(0, \sigma^2)$, $g'(\mu) = 0$,
$g''(\mu) \neq 0$ ise:
\[
  n(g(X_n) - g(\mu)) \xrightarrow{d} \frac{\sigma^2 g''(\mu)}{2} \chi^2_1.
\]
\end{teorem}

\begin{proof}
İkinci derecede Taylor:
$g(X_n) - g(\mu) = g'(\mu)(X_n-\mu) + \tfrac{g''(\mu)}{2}(X_n-\mu)^2 + o((X_n-\mu)^2)$.
$g'(\mu) = 0$:
$n(g(X_n)-g(\mu)) = \tfrac{g''(\mu)}{2}\cdot n(X_n-\mu)^2 + o_P(1)$.
$\sqrt{n}(X_n-\mu) \xrightarrow{d} \Normal(0,\sigma^2)$;
sürekli haritalama: $n(X_n-\mu)^2 \xrightarrow{d} \sigma^2 Z^2$ ($Z \sim \Normal(0,1)$).
$Z^2 \sim \chi^2_1$: $n(g(X_n)-g(\mu)) \xrightarrow{d} \frac{\sigma^2 g''(\mu)}{2}\chi^2_1$.
\end{proof}

\begin{ornek}[Varyans Tahmincisi]
$X_i$ i.i.d., $E[X_i] = \mu$, $\mathrm{Var}(X_i) = \sigma^2$, $E[(X_i-\mu)^4] = \mu_4 < \infty$.
$g(m) = (m - \mu)^2$ fonksiyonu, $g'(\mu) = 0$, $g''(\mu) = 2$.

$n(\bar{X}_n - \mu)^2 \xrightarrow{d} \sigma^2 \chi^2_1$.

Örneğin $\sigma^2 = 1$: $n(\bar{X}_n-\mu)^2 \xrightarrow{d} \chi^2_1$.
Bu, varyans-örneklem büyüklüğü ilişkisinin ikinci dereceden hassas analizini sağlar.
\end{ornek}

\begin{ornek}[Sıfır Noktasında Sinüs Dönüşümü]
$\bar{X}_n \xrightarrow{P} 0$; $g(x) = \sin(x)$. $g'(0) = 1 \neq 0$, normal delta yöntemi geçerli.
Ama $h(x) = 1 - \cos(x)$: $h'(0) = 0$, $h''(0) = 1$.
$n(1 - \cos(\bar{X}_n)) \xrightarrow{d} \frac{\sigma^2}{2}\chi^2_1$.
\end{ornek}

% =====================================================================
\section{Olasılık İntegral Dönüşümü ve Kuantil Dönüşümü}
\label{sec:prob_integral}
% =====================================================================

\begin{teorem}[Olasılık İntegral Dönüşümü (OİD)]\label{thm:oid}
$X$ sürekli rasgele değişken, KDF $F$ kesinlikle artan ve sürekli. O zaman
\[
  U = F(X) \sim \Uniform(0,1).
\]
\end{teorem}

\begin{proof}
$P(U \leq u) = P(F(X)\leq u) = P(X\leq F^{-1}(u)) = F(F^{-1}(u)) = u$, $u\in(0,1)$. Bu, $\Uniform(0,1)$ KDF'sidir.
\end{proof}

\begin{teorem}[Kuantil dönüşümü]\label{thm:kuantil}
$U\sim\Uniform(0,1)$ ve $F$ herhangi bir KDF, $F^{-1}$ genelleştirilmiş tersi. O zaman $X=F^{-1}(U)$ dağılımı $F$'dir.
\end{teorem}

\begin{proof}
$P(X\leq x)=P(F^{-1}(U)\leq x)=P(U\leq F(x))=F(x)$.
\end{proof}

\begin{ornek}[Simülasyon yöntemi]\label{ex:simulasyon}
$\Exp(\lambda)$ simülasyonu: $F(x)=1-e^{-\lambda x}$, $F^{-1}(u)=-\log(1-u)/\lambda$. $U\sim\Uniform(0,1)$ üret; $X=-\log U/\lambda$ hesapla ($1-U\overset{d}{=}U$).
Box-Muller ($\Normal(0,1)$): $U_1,U_2\iid\Uniform(0,1)$:
\[
  Z_1=\sqrt{-2\log U_1}\cos(2\pi U_2),\quad Z_2=\sqrt{-2\log U_1}\sin(2\pi U_2)\iid\Normal(0,1).
\]
\end{ornek}

\begin{hataanalizi}
OİD yalnızca sürekli dağılımlar için geçerlidir. Kesikli $X$ için $F(X)$ $\Uniform(0,1)$'e dağılmaz; bunun yerine $U\in(F(X^-),F(X)]$ aralığında üretilmiş $\Uniform$ ile çalışmak gerekir.
\end{hataanalizi}

% =====================================================================
\section{Çarpışma Transformasyonları ve Simetri}
\label{sec:simetri_donusumler}
% =====================================================================

\begin{tanim}[Simetrik dağılım]\label{def:simetrik}
$X$ dağılımı $0$ çevresinde simetriktir: $X\overset{d}{=}-X$.
\end{tanim}

\begin{teorem}[Mutlak değer dönüşümü]\label{thm:mutlak}
$X$, $0$ çevresinde simetrik, $f$ yoğunluğu mevcut. $Y=|X|$'in yoğunluğu:
\[
  f_Y(y) = 2f(y),\quad y>0.
\]
\end{teorem}

\begin{proof}
$P(|X|\leq y)=P(-y\leq X\leq y)=F(y)-F(-y)$. Türev: $f_Y(y)=f(y)+f(-y)=2f(y)$ (simetri).
\end{proof}

\begin{ornek}[Yarı-normal]\label{ex:yari_normal}
$Z\sim\Normal(0,1)$: $|Z|$ yarı-normal (\emph{half-normal}), $f_{|Z|}(y)=2\phi(y)$, $y>0$.
$\Beklenti[|Z|]=\sqrt{2/\pi}$, $\Var(|Z|)=1-2/\pi$.
\end{ornek}

\begin{teorem}[İşaret-büyüklük bağımsızlığı]\label{thm:isaret_buyukluk}
$X$ $0$ çevresinde simetrik. $S=\mathrm{sgn}(X)$ ve $M=|X|$ bağımsızdır; $S\sim\mathrm{Bernoulli}(1/2)$.
\end{teorem}

\begin{proof}
$P(S=1,M\leq m)=P(0<X\leq m)=F(m)-1/2=\tfrac12 P(M\leq m)$. Benzer şekilde $P(S=-1,M\leq m)=\tfrac12 P(M\leq m)$: bağımsızlık.
\end{proof}

\begin{disiplinlerarasibaglan}
\textbf{Finansal Matematik:} OİD, Monte Carlo risk simülasyonunun temelidir. Sklar teoremi: her ortak dağılım $F(\mathbf{x})=C(F_1(x_1),\ldots,F_p(x_p))$ biçiminde yazılır; $C$ \emph{copula}'dır — CDO kredi riski modellemesinde kullanılır.

\textbf{Kalite Mühendisliği:} Yarı-normal dağılım, Taguchi yönteminde tek taraflı üretim toleranslarını modellemek için kullanılır.
\end{disiplinlerarasibaglan}

\begin{mikroozet}
$g'(\mu) = 0$: standart delta yöntemi çöker.
İkinci derecede Delta: $n(g(X_n)-g(\mu)) \to \frac{\sigma^2 g''(\mu)}{2}\chi^2_1$.
Limit dağılım ki-kare; ölçek $\sqrt{n}$ yerine $n$.
OİD: $F(X)\sim\Uniform(0,1)$ — simülasyon ve copula modellemesinin temeli.
Simetrik dağılım: işaret ve büyüklük bağımsız; yarı-normal $\Beklenti[|Z|]=\sqrt{2/\pi}$.
\end{mikroozet}

% ---- Disiplinlerarası Bağlantı (TYMM) ----
\begin{kopru}[Disiplinlerarası — Mühendislik ve Veri Bilimi]
\textbf{Güvenilirlik mühendisliği:} Paralel sistemlerde hayatta kalma süresi $X_{(n)}$,
seri sistemlerde $X_{(1)}$. "$n$ bileşenden $k$'si çalışırsa sistem çalışır"
modeli, $k$'ıncı sıra istatistiğinin dağılımına dayanır.\\
\textbf{Veri bilimi:} Örneklem medyanı, aşırı değerlere dirençli bir merkez ölçüsüdür;
sıra istatistiklerinden türetilen yüzdelikler ($Q_1$, $Q_3$) kutu grafiğinin temelidir.\\
\textbf{Ekonometri:} Delta yöntemi, doğrusal olmayan parametre dönüşümlerinde
güven aralığı hesabında standart araçtır (oran tahmincisi, elastisite).
\end{kopru}

% ---- Öz-Yansıtma ----
\begin{ozyansima}
\begin{itemize}[leftmargin=1.8em]
  \item MÜF yönteminin Jacobian yöntemine göre avantajı nedir? Hangi durumda
        Jacobian daha kullanışlıdır?
  \item Sıra istatistiklerinin neden bağımlı olduğunu kendi cümlelerinizle açıklayın.
  \item Delta yönteminin temel fikri nedir? Taylor açılımının bu bağlamda nasıl çalıştığını anlatın.
  \item Üstel dağılımlı cihazlarda sıra arası boşlukların bağımsız olması sezgiye uygun mu?
\end{itemize}
\end{ozyansima}

% ---- Köprü Notu ----
\begin{kopru}
\textbf{Önceki bölüm:} Bölüm~\ref{chp:dagilim_aileleri} — $t$, $F$, $\chi^2$ dağılımları
burada türetildi; Bölüm~\ref{chp:beklenti}'deki MÜF/KF araçları evrişimde kullanıldı.\\
\textbf{Sonraki bölüm:} Bölüm~\ref{chp:merkezi_limit} — Delta yöntemi ve KF
yöntemi Merkezi Limit Teoremi'nin ispatında merkezi rol oynayacak.\\
\textbf{İstatistik kısmı:} Sıra istatistikleri Böl.~\ref{chp:istatistik_temelleri}'de
yüzdelik tahmincisi ve Böl.~\ref{chp:hipotez_testi}'nde parametresiz testlerde kullanılır.
\defterbaglantisi{bolum06\_donusumler.ipynb}{%
  Jacobian görselleştirmesi, sıra istatistikleri simülasyonu, delta yöntemi uygulaması}
\end{kopru}

% =====================================================================
%  ALIŞTIRMALAR
% =====================================================================

\cozumlualistirmalar

\begin{alistirma}[Al.6.1]{A}{Tekil Dönüşüm}
$X \sim \mathrm{Exp}(1)$. $Y = \sqrt{X}$ için:
(a) $f_Y(y)$'yi Jacobian yöntemiyle bulun.
(b) $E[Y]$'yi $f_Y$ üzerinden hesaplayın ($\Gamma(3/2) = \sqrt{\pi}/2$).
\end{alistirma}
\alcozum{%
(a) $g(x) = \sqrt{x}$, ters $h(y) = y^2$, $|h'(y)| = 2y$.
$f_Y(y) = e^{-y^2} \cdot 2y$, $y > 0$ (Rayleigh dağılımı, Bölüm~\ref{chp:rasgele_degiskenler}'de de elde edilmişti).\\
(b) $E[Y] = \int_0^\infty y \cdot 2y e^{-y^2} dy = 2\int_0^\infty y^2 e^{-y^2} dy$.
$u = y^2$: $= \int_0^\infty u^{1/2} e^{-u} du = \Gamma(3/2) = \sqrt{\pi}/2$.}

\begin{alistirma}[Al.6.2]{A}{Sıra İstatistikleri}
$X_1, X_2, X_3 \sim \Uniform(0,1)$ bağımsız.
(a) $X_{(1)}$ ve $X_{(3)}$'ün PDF'lerini yazın.
(b) $E[X_{(1)}]$ ve $E[X_{(3)}]$'ü hesaplayın.
(c) $P(X_{(2)} > 0.5)$'i bulun.
\end{alistirma}
\alcozum{%
(a) $f_{X_{(1)}}(x) = 3(1-x)^2$, $f_{X_{(3)}}(x) = 3x^2$, $0<x<1$.\\
(b) $E[X_{(1)}] = 1/(3+1) = 1/4$. $E[X_{(3)}] = 3/4$.\\
(c) $X_{(2)} \sim \Beta(2,2)$: $f_{X_{(2)}}(x) = 6x(1-x)$.
$P(X_{(2)}>0.5) = \int_{0.5}^1 6x(1-x)dx = 6[x^2/2-x^3/3]_{0.5}^1
= 6[(1/2-1/3)-(1/8-1/24)] = 6[1/6 - 1/12] = 6 \cdot 1/12 = 1/2$. (Simetri!)}

\begin{alistirma}[Al.6.3]{B}{MÜF Yöntemi}
$X_1, \ldots, X_n \sim \Pois(\lambda)$ bağımsız. $S_n = \sum X_i$.
(a) $M_{S_n}(t)$'yi bulun.
(b) $S_n$'nin dağılımını belirtin.
(c) $\bar{X}_n = S_n/n$'nin MÜF'ünü yazın ve $n=100$, $\lambda=2$ için
$P(\bar{X}_{100} \geq 2.5)$'e Chernoff sınırı verin.
\end{alistirma}
\alcozum{%
(a) $M_{S_n}(t) = \prod M_{X_i}(t) = [e^{\lambda(e^t-1)}]^n = e^{n\lambda(e^t-1)}$.\\
(b) $S_n \sim \Pois(n\lambda)$.\\
(c) $M_{\bar{X}_n}(t) = E[e^{t\sum X_i/n}] = M_{S_n}(t/n) = e^{n\lambda(e^{t/n}-1)}$.
Chernoff: $\Lambda^*(a) = \sup_t(ta - n\lambda(e^{t/n}-1))/n$. Optimize: $t^* = n\ln(a/\lambda)$.
$\Lambda^*(2.5) = 2.5\ln(2.5/2) - 2(2.5/2 - 1) = 2.5\ln(1.25) - 0.5 \approx 0.558 - 0.5 = 0.058$.
$P(\bar{X}_{100}\geq2.5) \leq e^{-100\cdot0.058} = e^{-5.8} \approx 0.003$.}

\begin{alistirma}[Al.6.4]{B}{Delta Yöntemi}
$X_1, \ldots, X_n$ i.i.d., $E[X_i] = \mu > 0$, $\mathrm{Var}(X_i) = \sigma^2$.
MLT (bir sonraki bölüm): $\sqrt{n}(\bar{X}_n - \mu) \xrightarrow{d} \Normal(0,\sigma^2)$.
Delta yöntemi ile $\sqrt{n}(1/\bar{X}_n - 1/\mu)$'nun limit dağılımını bulun.
\end{alistirma}
\alcozum{%
$g(x) = 1/x$, $g'(x) = -1/x^2$, $g'(\mu) = -1/\mu^2$.
Delta yöntemi: $\sqrt{n}(g(\bar{X}_n)-g(\mu)) \xrightarrow{d} \Normal(0, \sigma^2/\mu^4)$.
$\sqrt{n}(1/\bar{X}_n - 1/\mu) \xrightarrow{d} \Normal(0, \sigma^2/\mu^4)$.}

\begin{alistirma}[Al.6.4b]{B}{İki Boyutlu Jacobian — Polar Dönüşüm}
$(X,Y)\iid\Normal(0,1)$ bağımsız. $R=\sqrt{X^2+Y^2}$, $\Theta=\arctan(Y/X)$ polar dönüşümü.
(a) $(R,\Theta)$'nın birleşik yoğunluğunu Jacobian yöntemiyle bulun.
(b) $R$ ve $\Theta$'nın marjinal dağılımlarını yazın. $R$'nin adı ne?
(c) Bu sonucu Box-Muller dönüşümüyle ilişkilendirin.
\end{alistirma}
\alcozum{%
(a) Ters dönüşüm: $X=R\cos\Theta$, $Y=R\sin\Theta$. Jacobian:
\[
  J = \begin{vmatrix}\cos\Theta & -R\sin\Theta \\ \sin\Theta & R\cos\Theta\end{vmatrix} = R.
\]
Birleşik yoğunluk: $f_{X,Y}(x,y)=\frac{1}{2\pi}e^{-(x^2+y^2)/2}$.
$f_{R,\Theta}(r,\theta) = f_{X,Y}(r\cos\theta, r\sin\theta)\cdot|J|^{-1} = \frac{1}{2\pi}e^{-r^2/2}\cdot r$, $r>0$, $\theta\in[0,2\pi)$.

(b) $f_\Theta(\theta)=\frac{1}{2\pi}$ (Uniform$[0,2\pi)$); $f_R(r)=re^{-r^2/2}$ ($r>0$): Rayleigh dağılımı. $R^2\sim\chi^2_2=\mathrm{Exp}(1/2)$.

(c) Box-Muller: $U_1,U_2\iid\Uniform[0,1]$'den $X=\sqrt{-2\log U_1}\cos(2\pi U_2)$, $Y=\sqrt{-2\log U_1}\sin(2\pi U_2)$ bağımsız standart normaldir. Burada $R=\sqrt{-2\log U_1}\sim$ Rayleigh ve $\Theta=2\pi U_2\sim\Uniform[0,2\pi)$; bu tam olarak $(R,\Theta)$'nın dağılımıdır.}

\begin{alistirma}[Al.6.4c]{C}{Konvolüsyon ve Toplam Dağılımı}
$X\sim\Gama(\alpha,1)$, $Y\sim\Gama(\beta,1)$ bağımsız.
(a) MÜF yöntemiyle $Z=X+Y$'nin dağılımını bulun.
(b) $W=X/(X+Y)$'nin dağılımını konvolüsyon + Jacobian ile türetin.
(c) $\alpha=\beta=1$ durumunda (yani $X,Y\iid\mathrm{Exp}(1)$): $W$'nin dağılımını yazın ve sezgisel yorumunu verin.
\end{alistirma}
\alcozum{%
(a) $M_X(t)=(1-t)^{-\alpha}$, $M_Y(t)=(1-t)^{-\beta}$ (t<1). $M_Z(t)=(1-t)^{-(\alpha+\beta)}$; dolayısıyla $Z\sim\Gama(\alpha+\beta,1)$.

(b) Değişken dönüşümü: $(X,Y)\mapsto(Z=X+Y,W=X/(X+Y))$. Ters: $X=WZ$, $Y=(1-W)Z$. Jacobian: $|J|=Z$. Birleşik: $f_{Z,W}(z,w)=f_X(wz)f_Y((1-w)z)\cdot z$.
\[
  f_X(x)f_Y(y) = \frac{x^{\alpha-1}e^{-x}}{\Gamma(\alpha)}\cdot\frac{y^{\beta-1}e^{-y}}{\Gamma(\beta)}.
\]
$f_{Z,W}(z,w)=\frac{(wz)^{\alpha-1}((1-w)z)^{\beta-1}e^{-z}}{\Gamma(\alpha)\Gamma(\beta)}\cdot z = \frac{z^{\alpha+\beta-1}e^{-z}}{\Gamma(\alpha+\beta)}\cdot\frac{w^{\alpha-1}(1-w)^{\beta-1}}{B(\alpha,\beta)}$.

Marjinal $w$: $f_W(w)=\frac{w^{\alpha-1}(1-w)^{\beta-1}}{B(\alpha,\beta)}$; yani $W\sim\Beta(\alpha,\beta)$.

(c) $\alpha=\beta=1$: $W\sim\Beta(1,1)=\Uniform[0,1]$. Yorum: İki bağımsız Üstel sürenin toplamındaki birincinin oranı, tüm $[0,1]$'de eşit olasılıklı. Yenileme teorisi: bekleme kesiri üniform.}

% ---

\alistirmalar

\begin{alistirma}[Al.6.5]{A}{Dönüşüm Uygulaması}
$X \sim \Normal(0,1)$. $Y = e^X$ (lognormal).
(a) $F_Y(y)$'yi $\Phi$ cinsinden yazın. (b) $f_Y(y)$'yi türetin. (c) $P(1 < Y < e)$'yi bulun.
\end{alistirma}
\alcozum{%
(a) $F_Y(y) = P(e^X \leq y) = P(X \leq \ln y) = \Phi(\ln y)$.
(b) $f_Y(y) = \phi(\ln y)/y = e^{-(\ln y)^2/2}/(y\sqrt{2\pi})$, $y>0$.
(c) $P(1<Y<e) = P(0<X<1) = \Phi(1)-\Phi(0) \approx 0.8413 - 0.5 = 0.3413$.}

\begin{alistirma}[Al.6.6]{A}{Minimum Dağılımı}
$X_1,\ldots,X_5 \sim \mathrm{Exp}(2)$ bağımsız. $X_{(1)} = \min X_i$.
(a) $X_{(1)}$'in dağılımını bulun. (b) $P(X_{(1)} > 0.3)$'ü hesaplayın.
(c) $E[X_{(1)}]$'i bulun.
\end{alistirma}
\alcozum{%
(a) $X_{(1)} \sim \mathrm{Exp}(5\cdot2) = \mathrm{Exp}(10)$.
(b) $P(X_{(1)}>0.3) = e^{-10\cdot0.3} = e^{-3} \approx 0.050$.
(c) $E[X_{(1)}] = 1/10 = 0.1$.}

\begin{alistirma}[Al.6.7]{B}{Çok Boyutlu Jacobian}
$X, Y$ bağımsız $\Normal(0,1)$. $U = X+Y$, $V = X-Y$.
(a) $(U,V)$'nin ortak PDF'sini bulun.
(b) $U$ ve $V$ bağımsız mı? Dağılımlarını belirtin.
(c) $U$ ve $V$'nin $\mathrm{Cov}(U,V)$'sini hesaplayın.
\end{alistirma}
\alcozum{%
(a) Ters: $x=(u+v)/2$, $y=(u-v)/2$. $|\det J| = 1/2$.
$f_{U,V}(u,v) = \phi((u+v)/2)\phi((u-v)/2)/2 = \frac{1}{4\pi}e^{-(u^2+v^2)/4}$.
(b) $f_{U,V} = f_U(u)f_V(v)$ ile $U \sim \Normal(0,2)$, $V \sim \Normal(0,2)$, bağımsız.
(c) $\mathrm{Cov}(U,V) = \mathrm{Cov}(X+Y,X-Y) = \mathrm{Var}(X)-\mathrm{Var}(Y) = 1-1 = 0$. (Bağımsızlıkla tutarlı.)}

\begin{alistirma}[Al.6.8]{B}{Sıra İstatistikleri — Beklenti}
$X_1,\ldots,X_n \sim \Uniform(0,\theta)$. $X_{(n)} = \max X_i$ ile $\theta$'yı tahmin edin.
(a) $X_{(n)}$'nin PDF'sini bulun. (b) $E[X_{(n)}]$'i hesaplayın.
(c) $E[X_{(n)}] \neq \theta$: yansız tahminci nedir?
\end{alistirma}
\alcozum{%
(a) $F(x) = x/\theta$, $f_{X_{(n)}}(x) = n(x/\theta)^{n-1}/\theta = nx^{n-1}/\theta^n$.
(b) $E[X_{(n)}] = \int_0^\theta x \cdot nx^{n-1}/\theta^n dx = n\theta/(n+1)$.
(c) $\hat\theta = (n+1)/n \cdot X_{(n)}$ yansız tahminci: $E[\hat\theta] = \theta$.}

\begin{alistirma}[Al.6.9]{C}{Evrişim: Gama Toplamı}
$X \sim \Gama(\alpha, \beta)$, $Y \sim \Gama(\gamma, \beta)$ bağımsız ($\beta$ aynı).
Evrişim integraliyle (MÜF kullanmadan) $X+Y \sim \Gama(\alpha+\gamma, \beta)$
olduğunu ispatlayın.
\end{alistirma}
\alcozum{%
$f_{X+Y}(z) = \int_0^z f_X(t)f_Y(z-t)dt
= \int_0^z \frac{\beta^\alpha t^{\alpha-1}e^{-\beta t}}{\Gamma(\alpha)} \cdot \frac{\beta^\gamma(z-t)^{\gamma-1}e^{-\beta(z-t)}}{\Gamma(\gamma)} dt$.
$= \frac{\beta^{\alpha+\gamma}e^{-\beta z}}{\Gamma(\alpha)\Gamma(\gamma)} \int_0^z t^{\alpha-1}(z-t)^{\gamma-1}dt$.
$u = t/z$: $\int_0^z t^{\alpha-1}(z-t)^{\gamma-1}dt = z^{\alpha+\gamma-1}B(\alpha,\gamma)
= z^{\alpha+\gamma-1}\Gamma(\alpha)\Gamma(\gamma)/\Gamma(\alpha+\gamma)$.
$f_{X+Y}(z) = \frac{\beta^{\alpha+\gamma}}{\Gamma(\alpha+\gamma)}z^{\alpha+\gamma-1}e^{-\beta z}$
— bu $\Gama(\alpha+\gamma,\beta)$ PDF'sidir.}

\begin{alistirma}[Al.6.10]{B}{KF Tersçevirme}
$X$ rasgele değişken, $\varphi_X(t) = e^{-|t|}$ (Cauchy KF).
(a) $\int_{-\infty}^{+\infty} |\varphi_X(t)|\,dt$'nin varlığını kontrol edin.
(b) Tersçevirme formülünü kullanarak $f_X(x)$'i hesaplayın.
(c) Sonucun hangi standart dağılıma karşılık geldiğini belirtin.
\end{alistirma}
\alcozum{%
(a) $\int_{-\infty}^\infty |e^{-|t|}|\,dt = 2\int_0^\infty e^{-t}\,dt = 2 < \infty$. Formül uygulanabilir.\\
(b) $f_X(x) = \frac{1}{2\pi}\int_{-\infty}^{\infty} e^{-itx}e^{-|t|}\,dt
= \frac{1}{2\pi}\left[\int_0^\infty e^{t(-1-ix)}\,dt + \int_0^\infty e^{t(-1+ix)}\,dt\right]$
$= \frac{1}{2\pi}\left[\frac{1}{1+ix} + \frac{1}{1-ix}\right] = \frac{1}{2\pi}\cdot\frac{2}{1+x^2} = \frac{1}{\pi(1+x^2)}$.\\
(c) $f_X(x) = \frac{1}{\pi(1+x^2)}$ — standart Cauchy$(0,1)$ PDF'si.}

\begin{alistirma}[Al.6.11]{B}{İkinci Dereceden Delta Yöntemi}
$X_1, \ldots, X_n$ i.i.d., $E[X_i] = 0$, $\mathrm{Var}(X_i) = 1$.
$g(x) = x^2 + 3x$.
(a) $g'(0)$ ve $g''(0)$'ı hesaplayın.
(b) $\sqrt{n}(g(\bar{X}_n) - g(0))$ için limit dağılımı verin.
(c) Eğer $h(x) = x^2$ olsaydı, $n(h(\bar{X}_n) - h(0))$ için limit dağılımını verin.
\end{alistirma}
\alcozum{%
$g(x) = x^2 + 3x$; $g'(x) = 2x+3$; $g'(0) = 3 \neq 0$; $g''(0) = 2$.\\
(a) $g'(0) = 3$, $g''(0) = 2$.\\
(b) Standart delta (çünkü $g'(0)=3\neq0$): $\sqrt{n}(g(\bar{X}_n)-0) \xrightarrow{d} \Normal(0, 1\cdot 9) = \Normal(0,9)$.\\
(c) $h(x)=x^2$; $h'(0)=0$, $h''(0)=2$. İkinci dereceden delta:
$n(h(\bar{X}_n)-0) \xrightarrow{d} \frac{1\cdot2}{2}\chi^2_1 = \chi^2_1$.
Yani $n\bar{X}_n^2 \xrightarrow{d} \chi^2_1$ — bu MLT'nin doğal bir sonucudur.}

% =====================================================================
\endinput

% =====================================================================
%  BÖLÜM 7: Büyük Sayılar Kanunları
%  Kaynak: Feller Vol.I-II; Billingsley; Durrett; Klenke
% =====================================================================
\chapter{Büyük Sayılar Kanunları}
\label{chp:buyuk_sayilar}

\bolumalinti{%
  Büyük sayılar kanunu olasılığın en derin gerçeğidir:\\
  belirsizlik çoğunlukla sönümlenir.%
}{Richard Durrett, \textit{Probability: Theory and Examples}}

% ---- Kazanımlar (TYMM) ----
\begin{kazanimlar}
Bu bölümün sonunda öğrenci:
\begin{itemize}
  \item Neredeyse kesin, olasılıkla, dağılımla ve $L^p$ yakınsamalarını tanımlayabilir;
        aralarındaki hiyerarşik ilişkileri ispatlayabilir.
  \item Borel-Cantelli lemmalarını (her ikisini) ispatlayabilir ve uygulayabilir.
  \item Zayıf Büyük Sayılar Kanunu'nu (ZBSK) Chebyshev yöntemiyle ispatlayabilir;
        Helly seçim teoremi ile genel versiyonunu açıklayabilir.
  \item Güçlü Büyük Sayılar Kanunu'nu (GBSK) Kolmogorov yaklaşımıyla ispatlayabilir.
  \item Kolmogorov'un 0-1 Kanunu'nu ifade edip ispatlayabilir;
        "erillik" ve olasılık-0 veya 1 sonuçlarını tanıyabilir.
  \item Büyük sapma üst sınırlarını (Chernoff sınırı) türetebilir ve uygulayabilir.
\end{itemize}
\end{kazanimlar}

% ---- Motivasyon ----
\begin{motivasyon}
Bir para atışında yazı gelme oranı $n$ büyüdükçe $1/2$'ye yaklaşır — bu sezgi,
matematiksel olarak tam olarak ne anlama gelir? "Yaklaşmak" birden fazla biçimde
tanımlanabilir: \emph{her} örnek yolu için mi, yoksa çoğunlukla mı?

Büyük Sayılar Kanunları bu soruyu yanıtlar. Zayıf versiyon (ZBSK) olasılıkla
yakınsamayı, güçlü versiyon (GBSK) ise neredeyse kesin yakınsamayı garantiler.
Aradaki fark ince ama derin: GBSK, "uzun vadede ortalama gerçek beklentiye eşit
olur" ifadesinin en kesin matematiksel anlatımıdır.

\medskip
\textbf{Köprü.} Bölüm~\ref{chp:beklenti}'deki Chebyshev eşitsizliği ZBSK'nın
kalbi; Bölüm~\ref{chp:olcum_teorisi}'deki Monoton Yakınsama Teoremi GBSK'nın
temelidir. Bölüm~\ref{chp:merkezi_limit}'te yakınsama \emph{hızı} araştırılacak.
\end{motivasyon}

% =====================================================================
\section{Yakınsama Türleri}
\label{sec:yaklinsama_turleri}
% =====================================================================

\begin{motivasyon}[Neden Birden Fazla Yakınsama Türü?]
$X_n \to X$ demek nedir? Dizinin her terimi bir rasgele değişken.
Analiz derslerinin "nokta yakınsaması" (her $\omega$ için) çok kuvvetli;
"düzgün yakınsama" daha da kuvvetli. Olasılık teorisi bu ikisinin yanı sıra
daha zayıf ama yeterince güçlü iki kavram sunar: olasılıkla ve dağılımla yakınsama.
\end{motivasyon}

\begin{tanim}[Neredeyse Kesin Yakınsama]
\label{def:nk_yaklinsama}
$\{X_n\}$ rasgele değişkenler dizisi. $X_n \xrightarrow{\text{n.k.}} X$
(\textbf{neredeyse kesin yakınsama}, almost sure convergence) denir, eğer
\[
  P\!\left(\left\{\omega : \lim_{n\to\infty} X_n(\omega) = X(\omega)\right\}\right) = 1.
\]
Eşdeğer: $P(|X_n - X| > \varepsilon \text{ sonsuz çok } n \text{ için}) = 0$ her $\varepsilon > 0$'da.
\end{tanim}

\begin{tanim}[Olasılıkla Yakınsama]
\label{def:p_yaklinsama}
$X_n \xrightarrow{P} X$ (\textbf{olasılıkla yakınsama}, convergence in probability) denir, eğer
her $\varepsilon > 0$ için
\[
  \lim_{n\to\infty} P(|X_n - X| > \varepsilon) = 0.
\]
\end{tanim}

\begin{tanim}[Dağılımla Yakınsama]
\label{def:d_yaklinsama}
$X_n \xrightarrow{d} X$ (\textbf{dağılımla yakınsama}, convergence in distribution) denir, eğer
$F_X$ sürekli olan her $x$'te
\[
  \lim_{n\to\infty} F_{X_n}(x) = F_X(x).
\]
\end{tanim}

\begin{tanim}[$L^p$ Yakınsaması]
\label{def:lp_yaklinsama}
$p \geq 1$ ve $E[|X_n|^p], E[|X|^p] < \infty$ ise $X_n \xrightarrow{L^p} X$ denir, eğer
\[
  \lim_{n\to\infty} E[|X_n - X|^p] = 0.
\]
\end{tanim}

% =====================================================================
\section{Yakınsama Türleri Arasındaki İlişkiler}
\label{sec:yaklinsama_iliskileri}
% =====================================================================

\begin{teorem}[Yakınsama Hiyerarşisi]
\label{thm:yaklinsama_hiyerarsi}
Aşağıdaki çıkarımlar geçerlidir:
\[
  X_n \xrightarrow{\text{n.k.}} X \;\Rightarrow\; X_n \xrightarrow{P} X \;\Rightarrow\; X_n \xrightarrow{d} X,
\]
\[
  X_n \xrightarrow{L^p} X \;\Rightarrow\; X_n \xrightarrow{P} X.
\]
Hiçbir ok ters yönde genel olarak geçerli değildir.
\end{teorem}

\begin{proof}
\textbf{n.k.\ $\Rightarrow$ P:} Sabit $\varepsilon > 0$ için
$\limsup_n \{|X_n - X| > \varepsilon\} \subseteq \{\omega : X_n(\omega) \not\to X(\omega)\}$.
n.k.\ yakınsama: sağ tarafın olasılığı sıfır. Ölçünün monotonluğu:
$P(|X_n-X| > \varepsilon) \leq P(\bigcup_{k\geq n}\{|X_k-X|>\varepsilon\})$,
bu da $n\to\infty$'da sıfıra gider (azalan olaylar dizisi, Bölüm~\ref{chp:olcum_teorisi}).

\textbf{P $\Rightarrow$ d:} $\varepsilon > 0$ ve $x$ süreklilik noktası alın.
$F_{X_n}(x) = P(X_n \leq x) \leq P(X \leq x+\varepsilon) + P(|X_n - X| > \varepsilon)$.
$n\to\infty$: $\limsup F_{X_n}(x) \leq F_X(x+\varepsilon)$.
Benzer şekilde alt sınır kurulur; $\varepsilon\to 0$ ile $F_{X_n}(x) \to F_X(x)$.

\textbf{$L^p$ $\Rightarrow$ P:} Markov ($|X_n-X|^p \geq \varepsilon^p \mathbf{1}_{|X_n-X|>\varepsilon}$):
$P(|X_n-X|>\varepsilon) \leq E[|X_n-X|^p]/\varepsilon^p \to 0$.
\end{proof}

\begin{karsiornek}[Okların Terslenmediğine Örnekler]
\label{ex:yaklinsama_karsiornek}

\textbf{P $\not\Rightarrow$ n.k.:} $\Omega = [0,1]$, $P = \lambda$. "Gezgin blok":
$X_n = \mathbf{1}_{[(k-1)/m, k/m]}$ ($n = m(m-1)/2 + k$, $1 \leq k \leq m$).
$P(|X_n| > 0) = 1/m \to 0$ ama her $\omega$ için sonsuz çok $n$'de $X_n(\omega) = 1$.

\textbf{d $\not\Rightarrow$ P:} $X_n \equiv X \sim \Normal(0,1)$ ama $Y = -X \sim \Normal(0,1)$.
$X_n \xrightarrow{d} Y$ ama $X_n - Y = 2X$ sıfıra olasılıkla yakınsamaz.
\end{karsiornek}

\begin{teorem}[Slutsky Teoremi]
\label{thm:slutsky}
$X_n \xrightarrow{d} X$ ve $Y_n \xrightarrow{P} c$ (sabit) ise:
\begin{enumerate}[(i)]
  \item $X_n + Y_n \xrightarrow{d} X + c$,
  \item $X_n Y_n \xrightarrow{d} cX$,
  \item $X_n / Y_n \xrightarrow{d} X/c$ ($c \neq 0$).
\end{enumerate}
\end{teorem}

\begin{proof}
\textit{(i):} $P(X_n + Y_n \leq x) \leq P(X_n \leq x - c + \varepsilon) + P(|Y_n - c| > \varepsilon)$.
$n\to\infty$: $\limsup \leq F_X(x-c+\varepsilon)$; $\varepsilon\to 0$: $F_{X+c}(x)$.
Alt sınır benzer şekilde; eşitlik sağlanır.
\end{proof}

\begin{teorem}[Sürekli Haritalama Teoremi]
\label{thm:surekli_haritalama}
$g : \mathbb{R} \to \mathbb{R}$ sürekli ve $X_n \xrightarrow{d} X$ (veya $\xrightarrow{P}$ veya $\xrightarrow{\text{n.k.}}$)
ise $g(X_n) \xrightarrow{d} g(X)$ (veya $\xrightarrow{P}$, $\xrightarrow{\text{n.k.}}$ sırasıyla).
\end{teorem}

\begin{proof}
n.k.\ durumu: $X_n(\omega) \to X(\omega)$ ve $g$ sürekli $\Rightarrow$ $g(X_n(\omega)) \to g(X(\omega))$
($P = 1$ kümesi üzerinde).
Dağılımla yakınsama için Portmanteau teoremi kullanılır (ayrıntı: Billingsley~\cite{Billingsley1995}).
\end{proof}

\begin{mikroozet}
Kuvvetli sıralaması: n.k.\ $\Rightarrow$ P $\Rightarrow$ d; $L^p$ $\Rightarrow$ P.
Tersleri genel olarak yanlış. Slutsky: $X_n \xrightarrow{d} X$, $Y_n \xrightarrow{P} c \Rightarrow X_n+Y_n \xrightarrow{d} X+c$.
\end{mikroozet}

% =====================================================================
\section{Borel-Cantelli Lemmaları}
\label{sec:borel_cantelli}
% =====================================================================

\begin{kesif}
$\{A_n\}$ olaylar dizisi. "$A_n$ sonsuz çok gerçekleşir" olayını
Bölüm~\ref{chp:onbilgiler}'deki limsup dili ile nasıl ifade ederiz?
Eğer $\sum P(A_n) < \infty$ ise bu olayın gerçekleşme olasılığı hakkında
ne söyleyebiliriz?
\kesifcevap{``$A_n$ sonsuz çok gerçekleşir'' $= \limsup A_n = \bigcap_{n=1}^\infty \bigcup_{k=n}^\infty A_k$.
  Borel-Cantelli I: $\sum P(A_n)<\infty \Rightarrow P(\limsup A_n)=0$.
  Sezgi: toplam olasılık sonludur, dolayısıyla sonsuz kez gerçekleşme imkânsızdır (n.k.).}
\end{kesif}

\begin{tanim}[Sonsuz Sık Olay]
$\{A_n\}$ olaylar dizisi için $\limsup_n A_n = \bigcap_{n=1}^\infty \bigcup_{k=n}^\infty A_k$.
Bu olay "\textbf{$A_n$ sonsuz çok gerçekleşir}" (i.o.: infinitely often) diye okunur.
\end{tanim}

\begin{teorem}[Borel-Cantelli Birinci Lemması]
\label{thm:bc1}
$\{A_n\} \subseteq \mathcal{F}$ ve $\sum_{n=1}^\infty P(A_n) < \infty$ ise
$P(\limsup_n A_n) = 0$.
\end{teorem}

\begin{proof}
$B_n = \bigcup_{k=n}^\infty A_k$ için monoton azalma $B_n \downarrow \limsup_n A_n$.
\[
  P(\limsup_n A_n) = \lim_{n\to\infty} P(B_n) \leq \lim_{n\to\infty} \sum_{k=n}^\infty P(A_k) = 0,
\]
çünkü $\sum P(A_n) < \infty$ olduğundan kuyruk toplamı sıfıra gider.
\end{proof}

\begin{teorem}[Borel-Cantelli İkinci Lemması]
\label{thm:bc2}
$\{A_n\}$ \textbf{bağımsız} olaylar ve $\sum_{n=1}^\infty P(A_n) = \infty$ ise
$P(\limsup_n A_n) = 1$.
\end{teorem}

\begin{proof}
$P(\limsup_n A_n)^c = P(\liminf_n A_n^c)$.
$P\!\left(\bigcap_{k=n}^m A_k^c\right) = \prod_{k=n}^m (1-P(A_k))$ (bağımsızlık).
$1 - x \leq e^{-x}$: $\prod_{k=n}^m (1-P(A_k)) \leq \exp\!\left(-\sum_{k=n}^m P(A_k)\right) \to 0$
($\sum P(A_k) = \infty$ olduğundan). Yani $P(\bigcap_{k=n}^\infty A_k^c) = 0$ her $n$ için.
Sayılabilir birleşim: $P(\liminf A_n^c) = P(\bigcup_n \bigcap_{k\geq n} A_k^c) = 0$.
\end{proof}

\begin{hataanalizi}
İkinci lemma için bağımsızlık şartı çıkarılamaz.
Örnek: $A_n = A$ sabit olay, $P(A) = 1/2$, $\sum P(A_n) = \infty$.
Ama $P(\limsup A_n) = P(A) = 1/2 \neq 1$.
Bağımsızlık olmadan $\sum P(A_n) = \infty$ yalnızca üst sınır verir, kesin sonuç değil.
\end{hataanalizi}

\begin{ornek}[Para Atışı]
$A_n = \{n\text{'inci atışta arka}\ k\ \text{kez art arda}\}$.
$P(A_n) = 1/2^k$. $\sum P(A_n) = \infty$ ve $\{A_n\}$ bağımsız.
İkinci Borel-Cantelli: herhangi bir $k$ için art arda $k$ arka gelme n.k.\ sonsuz kez gerçekleşir.
\end{ornek}

% =====================================================================
\section{Zayıf Büyük Sayılar Kanunu}
\label{sec:zbsk}
% =====================================================================

\begin{tanim}[Örneklem Ortalaması]
$X_1, X_2, \ldots$ rasgele değişkenler; $\bar{X}_n = \frac{1}{n}\sum_{i=1}^n X_i$.
\end{tanim}

\begin{teorem}[Chebyshev ZBSK]
\label{thm:chebyshev_zbsk}
$X_1, X_2, \ldots$ çiftler arası korelasyonsuz (pairwise uncorrelated),
$E[X_i] = \mu$ ve $\mathrm{Var}(X_i) \leq M < \infty$ ise
$\bar{X}_n \xrightarrow{P} \mu$.
\end{teorem}

\begin{proof}
$E[\bar{X}_n] = \mu$. Korelasyonsuzluktan:
$\mathrm{Var}(\bar{X}_n) = \frac{1}{n^2}\sum_{i=1}^n \mathrm{Var}(X_i) \leq \frac{M}{n}$.
Chebyshev: $P(|\bar{X}_n - \mu| > \varepsilon) \leq \mathrm{Var}(\bar{X}_n)/\varepsilon^2 \leq M/(n\varepsilon^2) \to 0$.
\end{proof}

\begin{teorem}[Genel ZBSK (Kolmogorov)]
\label{thm:zbsk_genel}
$X_1, X_2, \ldots$ \textbf{i.i.d.}, $E[|X_1|] < \infty$, $E[X_1] = \mu$ ise
$\bar{X}_n \xrightarrow{P} \mu$.
\end{teorem}

\begin{proof}[Kırpma Yöntemi]
$\mu = 0$ (genel durum ötelemeli indirgenir). Sabit $M > 0$ için kırpılmış rasgele değişken:
$Y_i = X_i \mathbf{1}_{|X_i| \leq M}$.
$E[Y_i] = E[X_i \mathbf{1}_{|X_i|\leq M}] \to E[X_i] = 0$ ($M\to\infty$, DCT).
$\mathrm{Var}(Y_i) \leq E[Y_i^2] \leq ME[|Y_i|] \leq M \cdot E[|X_1|] < \infty$.

$\bar{X}_n - \bar{Y}_n = \frac{1}{n}\sum(X_i - Y_i) = \frac{1}{n}\sum X_i \mathbf{1}_{|X_i|>M}$.
Markov: $P(|\bar{X}_n - \bar{Y}_n| > \varepsilon) \leq \frac{1}{\varepsilon}E[|X_1|\mathbf{1}_{|X_1|>M}] \to 0$
($M\to\infty$, DCT).
$\bar{Y}_n \xrightarrow{P} E[Y_1] \approx 0$ (Chebyshev, $M$ büyük); ikisini birleştirince $\bar{X}_n \xrightarrow{P} 0$.
\end{proof}

\begin{ornek}[Monte Carlo Entegrasyonu]
$g : [0,1] \to \mathbb{R}$ integrallenebilir. $U_1, \ldots, U_n \sim \Uniform(0,1)$ bağımsız.
$\frac{1}{n}\sum g(U_i) \xrightarrow{P} E[g(U)] = \int_0^1 g(x)\,dx$.
ZBSK, Monte Carlo yönteminin matematiksel garantisidir.
\end{ornek}

\begin{kritiksorgu}
ZBSK şunu söyler: $P(|\bar{X}_n - \mu| > \varepsilon) \to 0$.
Bu, $\bar{X}_n$'nin hiçbir zaman $\mu$'dan uzak olmadığı anlamına mı gelir?
Hayır — yalnızca büyük sapmalar giderek daha düşük olasılıklı hale gelir.
$\bar{X}_n$'nin $\mu$'dan uzak olduğu anların kümesi boş olmayabilir;
ama bu anların "oranı" sıfıra yaklaşır.
\end{kritiksorgu}

% =====================================================================
\section{Güçlü Büyük Sayılar Kanunu}
\label{sec:gbsk}
% =====================================================================

\begin{teorem}[Kolmogorov GBSK]
\label{thm:gbsk}
$X_1, X_2, \ldots$ \textbf{i.i.d.}, $E[|X_1|] < \infty$, $E[X_1] = \mu$ ise
$\bar{X}_n \xrightarrow{\text{n.k.}} \mu$.
\end{teorem}

\begin{proof}
$\mu = 0$ varsayalım. Dört-moment argümanı (i.i.d., sonlu dördüncü moment varsayarak):

\textbf{Adım 1 (Dördüncü Moment Durumu):} $E[X_1^4] < \infty$ olsun.
\[
  E[(\bar{X}_n)^4] = \frac{1}{n^4} E\!\left[\left(\sum X_i\right)^4\right].
\]
$\sum X_i^4$ terim beklentisi $nE[X_1^4]$; çapraz terimler bağımsızlıkla yok olur;
$E[(X_i X_j)^2]$ terimleri $n(n-1)/2$ adet, her biri $E[X_1^2]^2$:
\[
  E[(\bar{X}_n)^4] = \frac{nE[X_1^4] + 3n(n-1)(E[X_1^2])^2}{n^4} \leq \frac{C}{n^2}.
\]
Markov ile $\varepsilon > 0$: $P(|\bar{X}_n| > \varepsilon) \leq C/(n^2\varepsilon^4)$.
$\sum_n P(|\bar{X}_n| > \varepsilon) \leq C/\varepsilon^4 \sum 1/n^2 < \infty$.
Borel-Cantelli 1: $P(|\bar{X}_n| > \varepsilon \text{ i.o.}) = 0$ her $\varepsilon > 0$ için.
Yani $\bar{X}_n \to 0$ n.k.

\textbf{Adım 2 (Genel Durum — Kolmogorov Kırpma):}
$E[|X_1|] < \infty$ için kırpılmış dizi $Y_n = X_n \mathbf{1}_{|X_n|\leq n}$ kullanılır.
Kolmogorov ölçütü: $\sum \mathrm{Var}(Y_n)/n^2 < \infty$ (Böl.~\ref{chp:olcum_teorisi}'deki
Kronecker lemması). Ayrıntılar Billingsley~\cite{Billingsley1995}, §22 veya Durrett~\cite{Durrett2019}, §2.4'te.
\end{proof}

\begin{teorem}[GBSK $\leftrightarrow$ ZBSK İlişkisi]
\label{thm:gbsk_zbsk}
GBSK, ZBSK'yi içerir: n.k.\ yakınsama $\Rightarrow$ olasılıkla yakınsama
(Teorem~\ref{thm:yaklinsama_hiyerarsi}). Tersi genel olarak yanlış.
\end{teorem}

\begin{ornek}[GBSK Uygulaması — Monte Carlo Hassasiyeti]
$g : [0,1] \to \mathbb{R}$ integrallenebilir. GBSK:
$\frac{1}{n}\sum_{i=1}^n g(U_i) \xrightarrow{\text{n.k.}} \int_0^1 g(x)\,dx$.
Bu, simülasyon çıktısının yalnızca "büyük ihtimalle" değil,
neredeyse tüm örnek yollarında kesinlikle doğru değere yakınsadığını garanti eder.
\end{ornek}

\begin{karsiornek}[GBSK'nın Beklenti Koşulunu Kaldıramazsınız]
$X_i \sim \mathrm{Cauchy}(0,1)$ bağımsız: $E[|X_1|] = \infty$.
$\bar{X}_n$, Cauchy(0,1) dağılımını korur — $0$'a yakınsamaz.
$E[|X_1|] < \infty$ koşulu vazgeçilmezdir.
\end{karsiornek}

\begin{mikroozet}
ZBSK: $\bar{X}_n \xrightarrow{P} \mu$ ($E[|X_1|] < \infty$, i.i.d.).
GBSK: $\bar{X}_n \xrightarrow{\text{n.k.}} \mu$ (aynı koşul, daha güçlü sonuç).
\end{mikroozet}

% =====================================================================
\section{Kolmogorov'un 0-1 Kanunu}
\label{sec:01_kanunu}
% =====================================================================

\begin{tanim}[Kuyruk $\sigma$-Cebiri]
\label{def:kuyruk_sigma}
$\{X_n\}$ rasgele değişkenler. $\mathcal{F}_n = \sigma(X_1, \ldots, X_n)$,
$\mathcal{T}_n = \sigma(X_{n+1}, X_{n+2}, \ldots)$. \textbf{Kuyruk $\sigma$-cebiri}:
\[
  \mathcal{T} = \bigcap_{n=1}^\infty \mathcal{T}_n.
\]
$\mathcal{T}$'nin olayları "sonlu sayıda $X_i$'yi değiştirmekten etkilenmeyen" olaylardır.
Örnekler: $\{\lim_n \bar{X}_n \text{ var}\}$, $\{\limsup_n X_n > c\}$, $\{\sum X_n \text{ yakınsar}\}$.
\end{tanim}

\begin{teorem}[Kolmogorov'un 0-1 Kanunu]
\label{thm:01_kanunu}
$X_1, X_2, \ldots$ bağımsız rasgele değişkenler ise kuyruk $\sigma$-cebirindeki
her $A \in \mathcal{T}$ olayı için $P(A) \in \{0, 1\}$.
\end{teorem}

\begin{proof}
\textbf{Adım 1:} Her $A \in \mathcal{T}$ ve her $n$ için $A \in \mathcal{T}_n = \sigma(X_{n+1}, X_{n+2}, \ldots)$,
dolayısıyla $A$, $\mathcal{F}_n = \sigma(X_1, \ldots, X_n)$'den bağımsızdır.

\textbf{Adım 2:} $\pi$-sistemi argümanı. $\mathcal{F}_\infty = \sigma(\bigcup_n \mathcal{F}_n)$.
Adım 1'deki bağımsızlık, $A$'nın $\mathcal{F}_\infty$'dan da bağımsız olduğunu gösterir
($\pi$-sistemi genişleme: $\bigcup_n \mathcal{F}_n$ bir $\pi$-sistemi; bağımsızlık üzerinde
$\lambda$-sistemi tutarlılığı ile Dynkin genişler).

\textbf{Adım 3:} $A \in \mathcal{T} \subseteq \mathcal{F}_\infty$ (her kuyruk olayı $\mathcal{F}_\infty$'dadır).
$A$, $\mathcal{F}_\infty$'dan bağımsız ve $A \in \mathcal{F}_\infty$:
$P(A) = P(A \cap A) = P(A) \cdot P(A)$.
Yani $P(A)^2 = P(A)$; çözüm: $P(A) \in \{0, 1\}$.
\end{proof}

\begin{ornek}[0-1 Kanunu Uygulamaları]
$\{X_n\}$ bağımsız rasgele değişkenler:
\begin{itemize}
  \item $P(\lim_n \bar{X}_n \text{ var}) \in \{0,1\}$. (GBSK bu olasılığın 1 olduğunu söyler.)
  \item $P(\sum X_n \text{ yakınsar}) \in \{0,1\}$.
  \item $P(\limsup_n X_n = +\infty) \in \{0,1\}$.
  \item $P(\limsup_n X_n/\sqrt{n\ln\ln n} = \sqrt{2}) \in \{0,1\}$. (Yinelenmiş logaritma kanunu.)
\end{itemize}
\end{ornek}

\begin{kritiksorgu}
0-1 Kanunu "olasılık ya 0 ya da 1" diyor. Bu deterministik mi?
Hayır — bu olaylar gerçek anlamda belirsiz olabilir; ancak birden fazla bağımsız
deneyin asimptotik davranışını yöneten olaylar, her örnek yolunu etkilese de
bütünleşik olarak keskin sınıflandırılabilir.
\end{kritiksorgu}

% =====================================================================
\section{Büyük Sapmalar: Giriş}
\label{sec:buyuk_sapmalar}
% =====================================================================

\begin{motivasyon}[Neden Büyük Sapmalar?]
ZBSK: $P(|\bar{X}_n - \mu| > \varepsilon) \to 0$. Chebyshev bize $O(1/n)$ hız verir.
Ama gerçek azalma çoğunlukla \emph{üstel}: $P(\bar{X}_n - \mu > \varepsilon) \approx e^{-nI(\varepsilon)}$.
Bu üstel azalma hızını veren $I(\varepsilon)$ fonksiyonu "oran fonksiyonu"dur.
\end{motivasyon}

\begin{teorem}[Chernoff Sınırı]
\label{thm:chernoff}
$X_1, \ldots, X_n$ i.i.d., $M(t) = E[e^{tX_1}]$ mevcut. Her $a > \mu = E[X_1]$ için
\[
  P(\bar{X}_n \geq a) \leq e^{-n \sup_{t \geq 0}(ta - \ln M(t))} = e^{-n \Lambda^*(a)},
\]
burada $\Lambda^*(a) = \sup_{t \geq 0}(ta - \ln M(t)) > 0$: \textbf{oran fonksiyonu}.
\end{teorem}

\begin{proof}
Her $t > 0$ için Markov ($e^{t\sum X_i} \geq e^{tna}$ olayından):
\[
  P(\bar{X}_n \geq a) = P\!\left(e^{t\sum X_i} \geq e^{tna}\right)
  \leq \frac{E[e^{t\sum X_i}]}{e^{tna}} = \frac{M(t)^n}{e^{tna}} = e^{n(\ln M(t) - ta)}.
\]
$t$'ye göre minimize etmek: $\inf_{t>0} e^{n(\ln M(t)-ta)} = e^{-n\sup_t(ta - \ln M(t))} = e^{-n\Lambda^*(a)}$.
\end{proof}

\begin{ornek}[Binom için Chernoff]
$X_i \sim \mathrm{Bernoulli}(p)$, $a > p$.
$\ln M(t) = \ln(q + pe^t)$.
$\Lambda^*(a) = a\ln(a/p) + (1-a)\ln((1-a)/q)$ (KL ıraksaması $\mathrm{Binom}$'dan $\mathrm{Binom}$'a).
$P(\bar{X}_n \geq a) \leq \exp(-n\,\mathrm{KL}(\mathrm{Ber}(a)\|\mathrm{Ber}(p)))$.
\end{ornek}

\begin{hocanotu}
Büyük sapmalar teorisi (Cramér teoremi, Varadhan'ın genel teorisi) bu kitabın
kapsamını aşar; ancak Chernoff sınırı istatistik ve algoritmik uygulamalarda
(öğrenme teorisi, randomize algoritmalar) sıkça kullanılır.
Derinlemesine okuma için Dembo-Zeitouni~\cite{Dembo1998}.
\end{hocanotu}

% ---- Disiplinlerarası Bağlantı (TYMM) ----
\begin{kopru}[Disiplinlerarası — Ekonomi, Fizik, Bilgisayar Bilimi]
\textbf{Ekonomi:} Büyük Sayılar Kanunu, ankete dayalı istatistiksel tahminlerin
(seçim anketi, tüketici fiyat endeksi) hangi koşulda güvenilir olduğunu açıklar:
gözlem sayısı $n$ büyüdükçe örneklem ortalaması nüfus ortalamasına yaklaşır.\\
\textbf{Fizik:} Ergodik teori (zamana göre ortalama = uzaya göre ortalama)
BSK'nın dinamik sistemler dilidir; termal dengedeki gaz molekülleri bu prensiple modellenir.\\
\textbf{Bilgisayar bilimi:} Randomize algoritmaların analizi BSK'ya dayanır:
$n$ tekrarlı bağımsız deneme, ortalama davranışı $O(1/\sqrt{n})$ hassasiyetle tahmin eder.
\end{kopru}

% ---- Öz-Yansıtma ----
\begin{ozyansima}
\begin{itemize}
  \item ZBSK ile GBSK arasındaki matematiksel fark nedir? Pratikte bu fark önemli mi?
  \item Borel-Cantelli lemmalarından birinin bağımsızlık koşulunu gerektirmesi,
        diğerinin gerektirmemesi hakkında ne düşünüyorsunuz?
  \item 0-1 Kanunu size "belirsizliğin çift değerli" olduğunu söylüyor;
        sezginizle bu nasıl bağdaşıyor?
  \item Chernoff sınırının Chebyshev'den daha güçlü olduğunu bir örnekle gösterin.
\end{itemize}
\end{ozyansima}

% ---- Köprü Notu ----
\begin{kopru}
\textbf{Önceki bölümler:} Bölüm~\ref{chp:beklenti} (Chebyshev, DCT) ve
Bölüm~\ref{chp:olcum_teorisi} (MYT) bu bölümün ispat araçlarını sağladı.\\
\textbf{Sonraki bölüm:} Bölüm~\ref{chp:donusumler} — Dönüşümler;
ardından Bölüm~\ref{chp:merkezi_limit} — BSK'nın öte boyutu:
$\bar{X}_n$ ne kadar hızlı yaklaşır? Cevap: $\sqrt{n}$ ölçeğinde Merkezi Limit Teoremi.\\
\textbf{İstatistik kısmı:} Bölüm~\ref{chp:istatistik_temelleri}'de büyük örneklem
tutarlılığı BSK çerçevesinde yeniden ele alınacak.
\defterbaglantisi{bolum07\_buyuk\_sayilar.ipynb}{%
  BSK simülasyonu, yakınsama hızı görselleştirmesi, Chernoff sınırı karşılaştırması}
\end{kopru}

% =====================================================================
\section{Yinelenmiş Logaritma Kanunu}
\label{sec:yil_kanunu}
% =====================================================================

\begin{motivasyon}[Büyük Sayılar Ötesinde]
GBSK: $\bar{X}_n \to \mu$ n.k. Ama ne \emph{kadar} hızlı yaklaşıyor?
$\sqrt{n}$ ölçeğinde Merkezi Limit Teoremi normal dağılım verir.
Daha ince soru: n.k.\ her örnek yolunun salınım büyüklüğü nedir?
Yinelenmiş Logaritma Kanunu (YLK) bu soruyu tam olarak yanıtlar.
\end{motivasyon}

\begin{teorem}[Kolmogorov'un Yinelenmiş Logaritma Kanunu]
\label{thm:ylk}
$X_1, X_2, \ldots$ i.i.d., $E[X_1] = 0$, $E[X_1^2] = \sigma^2 \in (0,\infty)$.
$S_n = X_1 + \cdots + X_n$ ise
\[
  P\!\left(\limsup_{n\to\infty} \frac{S_n}{\sqrt{2\sigma^2 n \ln\ln n}} = 1\right) = 1.
\]
Eşdeğer olarak: n.k.\ her örnek yolunda
$\limsup_{n\to\infty} S_n/\sqrt{2\sigma^2 n\ln\ln n} = 1$
ve $\liminf_{n\to\infty} S_n/\sqrt{2\sigma^2 n\ln\ln n} = -1$.
\end{teorem}

\begin{proof}[İspat Taslağı]
İki yönlü ispat gerekir.

\textbf{Üst sınır ($\limsup \leq 1$ n.k.):}
Sabit $b > 1$. Geometrik alt dizi $n_k = \lfloor b^k \rfloor$ seçin.
$A_k = \{S_{n_k} > (1+\varepsilon)\sqrt{2\sigma^2 n_k \ln\ln n_k}\}$.
Olağan sapma tahmini ile $P(A_k) \leq C/k^{1+\varepsilon}$.
$\sum P(A_k) < \infty$: Borel-Cantelli 1 ile $P(A_k \text{ i.o.}) = 0$.
Alt diziler arası maksimum kontrolü tüm $n$'ye genişler.

\textbf{Alt sınır ($\limsup \geq 1$ n.k.):}
Uygun $\varepsilon_k \to 0$ ile olaylar
$B_k = \{S_{n_{k+1}} - S_{n_k} > (1-\varepsilon)\sqrt{2\sigma^2(n_{k+1}-n_k)\ln\ln n_{k+1}}\}$
tanımlanır. $\{B_k\}$ asimptotik bağımsız; $\sum P(B_k) = \infty$.
Bağımsız Borel-Cantelli: $P(B_k \text{ i.o.}) = 1$.
Tam ispat: Hartman-Wintner (1941), bkz. Billingsley~\cite{Billingsley1995}, §12.
\end{proof}

\begin{ornek}[YLK'nın Yorumu]
$X_i \sim \Normal(0,1)$, $\sigma^2 = 1$. YLK:
\[
  \limsup_{n\to\infty} \frac{S_n}{\sqrt{2n\ln\ln n}} = 1 \quad \text{n.k.}
\]
$n = 10^6$: $\sqrt{2 \cdot 10^6 \cdot \ln\ln 10^6} = \sqrt{2\times10^6 \times 2.97} \approx 2439$.
$S_n$ sonsuz çok kez $\approx 2439$'a ulaşır, ama asla $2439 \cdot (1+\varepsilon)$'ı
sonsuz kez geçmez.
\end{ornek}

\begin{hataanalizi}
YLK ve MLT'yi karıştırmak yaygın hatadır.
MLT: $S_n/\sqrt{n\sigma^2} \xrightarrow{d} \Normal(0,1)$ — dağılımsal limit.
YLK: $\limsup S_n/\sqrt{2n\sigma^2\ln\ln n} = 1$ n.k. — yol bazında keskin sınır.
YLK, MLT'den \emph{daha güçlü} bir sonuçtur; her bir örnek yolu hakkında bilgi verir.
\end{hataanalizi}

\begin{mikroozet}
YLK: $S_n$'nin n.k.\ salınım büyüklüğü tam olarak $\sqrt{2\sigma^2 n \ln\ln n}$ ölçeğindedir.
GBSK ($S_n = o(n)$) ve MLT ($S_n = O(\sqrt{n})$) arasına düşer: $S_n = O(\sqrt{n\ln\ln n})$ n.k.
\end{mikroozet}

% =====================================================================
\section{Yenileme Teorisi Girişi}
\label{sec:yenileme_teorisi}
% =====================================================================

\begin{tanim}[Yenileme Süreci]
\label{def:yenileme_sureci}
$X_1, X_2, \ldots$ i.i.d., $X_i > 0$ n.k., $E[X_1] = \mu < \infty$.
$X_i$: arası zamanlar. $S_0 = 0$, $S_n = X_1 + \cdots + X_n$.
\textbf{Yenileme sayacı}: $N(t) = \max\{n \geq 0 : S_n \leq t\}$.
$\{N(t), t \geq 0\}$: \textbf{yenileme süreci} (renewal process).
\end{tanim}

\begin{teorem}[Temel Yenileme Teoremi]
\label{thm:temel_yenileme}
$E[X_1] = \mu \in (0,\infty)$ ise
\[
  \frac{N(t)}{t} \xrightarrow{\text{n.k.}} \frac{1}{\mu} \quad (t \to \infty).
\]
Daha kesin: $E[N(t)]/t \to 1/\mu$ ve $\mathrm{Var}(N(t))/t \to \sigma^2/\mu^3$
($\sigma^2 = \mathrm{Var}(X_1)$).
\end{teorem}

\begin{proof}
$N(t) \geq n \Leftrightarrow S_n \leq t$. Dolayısıyla $S_{N(t)} \leq t < S_{N(t)+1}$.
Her iki tarafı $N(t)$'ye böl:
\[
  \frac{S_{N(t)}}{N(t)} \leq \frac{t}{N(t)} < \frac{S_{N(t)+1}}{N(t)}.
\]
GBSK: $S_{N(t)}/N(t) \xrightarrow{\text{n.k.}} \mu$ ($N(t) \to \infty$ çünkü $\mu < \infty$).
Sıkıştırma: $t/N(t) \xrightarrow{\text{n.k.}} \mu$, yani $N(t)/t \xrightarrow{\text{n.k.}} 1/\mu$.
\end{proof}

\begin{ornek}[Makine Bakım Modeli]
Bir makine $X_i \sim \Exp(\lambda)$ ömrü sonunda yenileniyor.
$E[X_i] = 1/\lambda$. Temel Yenileme: $N(t)/t \to \lambda$ n.k.
$t$ birim zamanda ortalama $\lambda t$ yenileme yapılır.
\end{ornek}

\begin{teorem}[Blackwell'in Yenileme Teoremi]
\label{thm:blackwell}
$X_i$ aritmetik-olmayan (non-arithmetic) dağılımlı ise: her $h > 0$ için
\[
  E[N(t+h) - N(t)] \to \frac{h}{\mu} \quad (t \to \infty).
\]
Yani uzun vadede $[t, t+h]$ aralığındaki beklenen yenileme sayısı $h/\mu$'ya yaklaşır:
sistem "kararlı duruma" (steady state) girer.
\end{teorem}

\begin{disiplinlerarasibaglan}
\textbf{Güvenilirlik Mühendisliği:} Yenileme süreci, ekipman ömürleri ve
bakım planlamasının temel modelidir: $N(t)$ arızaların kümülatif sayısıdır.\\
\textbf{Kuyruk Teorisi:} M/G/1 kuyruğunda servis süreleri yenileme sürecini oluşturur;
Temel Yenileme Teoremi ortalama bekleme süresi hesaplamasında kullanılır.\\
\textbf{Ekonomi:} Stok yönetiminde sipariş dönemleri yenileme yapısı taşır.
\end{disiplinlerarasibaglan}

% =====================================================================
\section{Durağan Diziler için Ergodik Teorem}
\label{sec:ergodik_teorem}
% =====================================================================

\begin{tanim}[Durağan Dizi]
\label{def:duragan_dizi}
$\{X_n\}$ rasgele değişkenler dizisi \textbf{(geniş anlamda) durağan}
(stationary) denir, eğer her $k, n$ ve $h$ için
$(X_k, X_{k+1}, \ldots, X_{k+n})$ ile $(X_{k+h}, X_{k+h+1}, \ldots, X_{k+h+n})$
aynı dağılıma sahipse.

Özel durumlar: i.i.d.\ diziler, durağan Markov zincirleri.
\end{tanim}

\begin{tanim}[Kaydırma Dönüşümü ve Ergodiklik]
\label{def:ergodiklik}
$\Omega$'da \textbf{kaydırma dönüşümü} $\theta: \theta(\omega)_n = \omega_{n+1}$.
Dizi $\theta$-uyumlu: $X_n(\omega) = X_0(\theta^n \omega)$.
$A \in \mathcal{F}$ \textbf{değişmez} (invariant) olay: $\theta^{-1}A = A$.
$P$ ölçümü \textbf{ergodik}: her değişmez olay $A$ için $P(A) \in \{0,1\}$.
\end{tanim}

\begin{teorem}[Birkhoff'un Ergodik Teoremi]
\label{thm:birkhoff}
$\{X_n\}$ durağan ve ergodik (ölçüm $P$ ergodik), $E[|X_0|] < \infty$. O zaman
\[
  \frac{1}{n}\sum_{k=0}^{n-1} X_k \xrightarrow{\text{n.k.}} E[X_0].
\]
\end{teorem}

\begin{proof}[İspat Notu]
Bağımsız diziler için (i.i.d.) bu Kolmogorov GBSK'dır.
Genel ergodik durum için ispat daha derindir: Maximal Ergodik Lemma (Maximal Lemma)
ve Garsia'nın kombinatoryal argümanı kullanılır.
Referans: Durrett~\cite{Durrett2019}, §7.2; Billingsley~\cite{Billingsley1995}, §24.
\end{proof}

\begin{ornek}[Ergodik Markov Zinciri]
$\{Z_n\}$ sonlu durum uzaylı, geri dönüşümlü ve pozitif tekrarlı Markov zinciri;
$\pi$ durağan dağılımı. $f: S \to \mathbb{R}$ her duruma bir değer atar.
Birkhoff: $\frac{1}{n}\sum_{k=0}^{n-1} f(Z_k) \xrightarrow{\text{n.k.}} \sum_{s\in S} f(s)\pi(s) = E_\pi[f]$.

Yani "zamansal ortalama = uzamsal ortalama" — klasik ergodiklik eşitliği.
\end{ornek}

\begin{kritiksorgu}
Ergodik teorem i.i.d.\ koşulunu gevşetiyor. Hangi koşulun gerektiğini tartışın:
bağımsızlık mı, durağanlık mı, yoksa ergodiklik mi? Ergodikliğin 0-1 Kanunu
ile ilişkisini açıklayın.
\end{kritiksorgu}

% =====================================================================
\section{Yavaş Değişen Fonksiyonlar ve Düzenli Değişim}
\label{sec:duzenli_degisim}
% =====================================================================

\begin{tanim}[Yavaş değişen fonksiyon]\label{def:yavas_degisen}
$L:(0,\infty)\to(0,\infty)$ fonksiyonu \emph{yavaş değişen} (\emph{slowly varying}) denir, her $c>0$ için
\[
  \lim_{x\to\infty}\frac{L(cx)}{L(x)}=1.
\]
\end{tanim}

\begin{ornek}[Yavaş değişen fonksiyon örnekleri]\label{ex:yavas_ornekler}
\begin{enumerate}
  \item Sabitler: $L(x)=c>0$.
  \item $L(x)=\log x$, $L(x)=(\log\log x)^2$, $L(x)=\log(1+x^2)$.
  \item İterated logarithm: $L(x)=\log_k(x)$ ($k$'ıncı yinelenmiş logaritma).
\end{enumerate}
\emph{Düzenli değişen}: $R(x)=x^\alpha L(x)$ ($L$ yavaş değişen, $\alpha\in\mathbb{R}$ indeks).
\end{ornek}

\begin{teorem}[Karamata temsil teoremi]\label{thm:karamata}
$L$ yavaş değişen $\iff$ $L(x)=c(x)\exp\!\left(\int_1^x \frac{\varepsilon(t)}{t}dt\right)$,
burada $c(x)\to c>0$ ve $\varepsilon(t)\to0$.
\end{teorem}

\begin{teorem}[Potter sınırları]\label{thm:potter}
$L$ yavaş değişen. Her $\delta>0$ için $x_0>0$ mevcut:
$x,y\geq x_0$ ise
\[
  \frac{L(y)}{L(x)} \leq 2\max\!\left\{\left(\frac{y}{x}\right)^\delta, \left(\frac{x}{y}\right)^\delta\right\}.
\]
\end{teorem}

\begin{ornek}[Ağır kuyruklu dağılımlar]\label{ex:agir_kuyruk}
$P(X>x)\sim x^{-\alpha}L(x)$ ($\alpha>0$, $L$ yavaş değişen): \emph{düzenli değişen kuyruk}.
\begin{itemize}
  \item $\alpha>2$: $E[X^2]<\infty$ — standart MLT geçerli.
  \item $1<\alpha\leq2$: $E[X]<\infty$ ama $E[X^2]=\infty$ — kararlı dağılımlara yakınsama, ölçek $n^{1/\alpha}$.
  \item $0<\alpha\leq1$: $E[X]=\infty$ — büyük sayılar kanunu çöker.
\end{itemize}
Pareto$(\alpha)$, Cauchy $(\alpha=1)$, Student-$t_\nu$ ($\nu=\alpha$) bu sınıftadır.
\end{ornek}

\begin{hataanalizi}
"Büyük sayılar kanunu her zaman geçerlidir" yanlıştır. $E[|X|]=\infty$ durumunda ZBSK ve GBSK çöker: Cauchy dağılımında örneklem ortalaması hiçbir değere yakınsamaz, her örneklem için farklı bir değer verir. Bu, finansta "ağır kuyruklu" kayıpların klasik istatistikle modellenememesinin matematiksel gerekçesidir.
\end{hataanalizi}

% =====================================================================
\section{Büyük Sapmalar İçin Üst Sınırlar}
\label{sec:buyuk_sapma_sinir}
% =====================================================================

\begin{teorem}[Markov eşitsizliği — hatırlatma]\label{thm:markov_hatirlatma}
$X\geq0$, $a>0$: $P(X\geq a)\leq E[X]/a$.
\end{teorem}

\begin{teorem}[Chebyshev eşitsizliği]\label{thm:chebyshev_hatirlatma}
$E[X]=\mu$, $\Var(X)=\sigma^2<\infty$, $k>0$:
$P(|X-\mu|\geq k\sigma)\leq 1/k^2$.
\end{teorem}

\begin{teorem}[Chernoff sınırı]\label{thm:chernoff}
$X_1,\ldots,X_n\iid$, MGF $M(\lambda)=E[e^{\lambda X}]$ $\lambda_0>0$ için sonlu. $S_n=\sum_{i=1}^n X_i$. Her $t>E[X]$ için:
\[
  P(S_n\geq nt) \leq e^{-n\sup_{\lambda\geq0}[\lambda t - \log M(\lambda)]}.
\]
Üstel $[\lambda t - \log M(\lambda)]$ büyük sapma hızıdır (\emph{rate function} $I(t)$).
\end{teorem}

\begin{proof}
$P(S_n\geq nt)=P(e^{\lambda S_n}\geq e^{n\lambda t})\leq e^{-n\lambda t}E[e^{\lambda S_n}]=e^{-n\lambda t}[M(\lambda)]^n=e^{-n[\lambda t-\log M(\lambda)]}$.
$\lambda\geq0$'da supremum alınır.
\end{proof}

\begin{ornek}[Hoeffding eşitsizliği]\label{ex:hoeffding}
$X_i\in[a_i,b_i]$ sınırlı bağımsız, $E[X_i]=0$. $P(S_n\geq t)\leq\exp\!\left(-\frac{2t^2}{\sum(b_i-a_i)^2}\right)$.
Özellikle $a_i=0$, $b_i=1$, $t=n\varepsilon$: $P(\bar X_n\geq\varepsilon)\leq e^{-2n\varepsilon^2}$.
\end{ornek}

\begin{disiplinlerarasibaglan}
\textbf{Makine Öğrenmesi — PAC öğrenimi:} Hoeffding eşitsizliği, bir öğrenme algoritmasının eğitim hatasının test hatasına $\varepsilon$ içinde yakınsadığını yüksek olasılıkla garanti eden PAC (\emph{probably approximately correct}) öğrenimi sınırının temelidir.

\textbf{Finansal Risk:} Düzenli değişen kuyruklar (Pareto, $t_\nu$) piyasa çöküşlerini modellemek için kullanılır; Basle III sermaye yeterliliği hesaplamalarında VaR (değer-risk) bu dağılımlara dayanır.
\end{disiplinlerarasibaglan}

\begin{mikroozet}
Durağan + ergodik dizilerde $\frac{1}{n}\sum_{k=0}^{n-1}X_k \to E[X_0]$ n.k.
Bu, GBSK'nın i.i.d.\ ötesine genellemesidir; Markov zincirleri ve zaman serisi
analizinin matematiksel temelidir.
Düzenli değişen kuyruklar ($P(X>x)\sim x^{-\alpha}L(x)$): $\alpha\leq1$'de BSK çöker.
Chernoff/Hoeffding: üstel büyük sapma sınırları; makine öğrenmesinde PAC sınırlarının temeli.
\end{mikroozet}

% =====================================================================
%  ALIŞTIRMALAR
% =====================================================================

\cozumlualistirmalar

\begin{alistirma}[Al.7.1]{A}{Yakınsama Türlerini Tanımak}
$X_n = n^{-1/2}Z$ ($Z \sim \Normal(0,1)$, sabit). Gösterin:
(a) $X_n \xrightarrow{\text{n.k.}} 0$. (b) $X_n \xrightarrow{P} 0$. (c) $X_n \xrightarrow{d} 0$.
Bu örnek hangi tipik yanılgıyı düzeltiyor?
\end{alistirma}
\alcozum{%
$X_n(\omega) = Z(\omega)/\sqrt{n} \to 0$ her $\omega$'da (sabit $Z$, deterministik yakınsama).
(a)--(c) üçü de trivyal olarak sağlanır.\\
Düzeltilen yanılgı: Burada $X_n$ "aynı $Z$'yi ölçekliyor"; dizi gerçekten bağımsız değil.
Bağımsız $X_n$'lerin neredeyse kesin yakınsaması için Borel-Cantelli gerekir.}

\begin{alistirma}[Al.7.2]{A}{Borel-Cantelli Uygulaması}
$X_1, X_2, \ldots$ bağımsız, $P(X_n > n) = 1/n^2$.
$P(X_n > n \text{ sonsuz çok } n \text{ için}) = ?$
\end{alistirma}
\alcozum{%
$\sum_{n=1}^\infty P(X_n > n) = \sum 1/n^2 = \pi^2/6 < \infty$.
Borel-Cantelli 1: $P(X_n > n \text{ i.o.}) = 0$.
Yani yalnızca sonlu sayıda $n$ için $X_n > n$ gerçekleşir.}

\begin{alistirma}[Al.7.3]{B}{ZBSK ve GBSK Karşılaştırması}
$X_1, X_2, \ldots$ i.i.d.\ $\mathrm{Exp}(1)$. $\bar{X}_n$'nin hangi değere yakınsadığını
belirtin. (a) Chebyshev ZBSK ile $P(|\bar{X}_n - 1| > \varepsilon)$ için $n$ sayısına bağlı
üst sınır verin. (b) Bu yakınsamanın GBSK anlamında da geçerli olduğunu söyleyin.
\end{alistirma}
\alcozum{%
$E[X_1] = 1$, $\mathrm{Var}(X_1) = 1$.\\
(a) Chebyshev: $P(|\bar{X}_n - 1| > \varepsilon) \leq 1/(n\varepsilon^2)$.
$\varepsilon = 0.1$ için $n = 100$ yeterli: $P(\ldots) \leq 1$.
(İyi bir üst sınır değil ama metodoloji doğru; Chernoff çok daha iyi.)\\
(b) $E[|X_1|] = 1 < \infty$, i.i.d.; GBSK: $\bar{X}_n \xrightarrow{\text{n.k.}} 1$.}

\begin{alistirma}[Al.7.4]{B}{Chernoff Sınırı}
$X_i \sim \mathrm{Bernoulli}(0.5)$ bağımsız, $n = 100$, $a = 0.6$.
(a) Chebyshev ile $P(\bar{X}_{100} \geq 0.6)$ için üst sınır verin.
(b) Chernoff ile üst sınır verin ve karşılaştırın.
\end{alistirma}
\alcozum{%
(a) Chebyshev: $P(\bar{X}_n \geq 0.6) \leq P(|\bar{X}_n-0.5|\geq 0.1) \leq \mathrm{Var}(\bar{X}_n)/(0.1)^2
= (0.25/100)/0.01 = 0.25$.\\
(b) Chernoff: $\Lambda^*(0.6) = 0.6\ln(0.6/0.5)+0.4\ln(0.4/0.5) = 0.6\ln(1.2)+0.4\ln(0.8)$.
$\approx 0.6(0.1823)+0.4(-0.2231) = 0.1094-0.0892 = 0.0202$.
$P(\bar{X}_{100}\geq 0.6)\leq e^{-100\times0.0202} = e^{-2.02} \approx 0.133$.
Chebyshev: 0.25; Chernoff: 0.133. Gerçek değer $\approx 0.028$ (Chernoff daha yakın).}

\begin{alistirma}[Al.7.4b]{B}{Bağımsızlık Olmadan BSK — Ergodik Teorem}
$(X_n)_{n\geq1}$ \emph{durağan ergodik} dizi; $E[|X_1|]<\infty$, $\mu=E[X_1]$.
(a) Birkhoff ergodik teoremi ne söyler? Güçlü BSK ile nasıl kıyaslanır?
(b) Markov zinciri $X_n$ durağan dağılım $\pi$ ile ergodikse $\frac{1}{n}\sum_{k=1}^n f(X_k)\to\sum_x f(x)\pi(x)$ n.k. Bunu simülasyon bağlamında yorumlayın.
(c) $X_1,X_2,\ldots$ i.i.d.\ ama simetrik Cauchy dağılımlı. $\bar X_n$ ne yapar? Neden BSK çalışmaz?
\end{alistirma}
\alcozum{%
(a) \textbf{Birkhoff ergodik teoremi}: Durağan ergodik $(X_n)$ için $\bar X_n=\frac{1}{n}\sum_{k=1}^n X_k\xrightarrow{\text{n.k.}}\mu=E[X_1]$. Güçlü BSK, i.i.d.\ varsayımı altında ergodikliğin özel halidir; Birkhoff bağımlı ama durağan-ergodik dizilere geneller.

(b) Markov Monte Carlo (MCMC) temeli: Zinciri durağan dağılımdan başlatın (veya yeterince uzun ısın sonra) ve $\frac{1}{n}\sum f(X_k)$ ile $E_\pi[f]$'yi tahmin edin. Ergodiklik garantisi yakınsamayı sağlar; bağımsız örnekleme yerine geçer.

(c) Cauchy için $E[|X_1|]=\infty$ (beklenti var olmaz). Güçlü BSK'nın koşulu $E[|X_1|]<\infty$ ihlal edilir. Gerçekte $\bar X_n$ Cauchy($0,1$) dağılımlı kalır: $n$ büyüse de dağılım daralmaz. Büyük sayılar kanunları geçerli değil.}

\begin{alistirma}[Al.7.4c]{C}{Yavaş Değişen Fonksiyonlar ve Kuyruğun Ağırlığı}
$L$ yavaş değişen (slowly varying) demek: her $c>0$ için $L(cx)/L(x)\to1$ ($x\to\infty$).
(a) $L(x)=\log x$ ve $L(x)=(\log x)^2$'nin yavaş değişen olduğunu doğrulayın.
(b) $P(X>x)=L(x)/x^\alpha$ ($\alpha>0$) ise $X$'in ortalaması ne zaman sonlu?
(c) $\alpha=1$, $L(x)=1$ durumunda (Cauchy benzeri), $S_n/n\to?$ Limitin ne olduğunu belirtin.
\end{alistirma}
\alcozum{%
(a) $L(x)=\log x$: $\frac{\log(cx)}{\log x}=\frac{\log c+\log x}{\log x}=1+\frac{\log c}{\log x}\to1$. $\checkmark$ Benzer şekilde $(\log x)^2$ için de aynı argüman.

(b) $E[X]=\int_0^\infty P(X>x)dx=\int_0^1 dx+\int_1^\infty \frac{L(x)}{x^\alpha}dx$. İkinci integral $\int_1^\infty x^{-\alpha}dx$ ile karşılaştırılır (Potter lemması: $L(x)\leq Cx^\varepsilon$ yeterince büyük $x$'te). Sonlu olması için $\alpha>1$ gerekir ($\alpha=1$ sınır: $\log$ düzeltmesiyle ince analiz).

(c) $\alpha=1$, $P(X>x)\sim1/x$ ($L=1$): Cauchy dağılımı modeli. Stabil dağılım teorisi: $S_n/(n\log n)^{1/\alpha}$ için uygun normalleştirme gerekir; $S_n/n$ hiçbir sonlu limite yakınsamaz ($\bar X_n$ Cauchy kalır). BSK ve CLT her ikisi de çöker.}

% ---

\alistirmalar

\begin{alistirma}[Al.7.5]{A}{Yakınsama Türleri}
Her biri için örnek verin ya da imkânsız olduğunu gösterin:
(a) $X_n \xrightarrow{P} X$ ama $X_n \not\xrightarrow{\text{n.k.}} X$.
(b) $X_n \xrightarrow{d} c$ (sabit) ama $X_n \not\xrightarrow{P} c$.
\end{alistirma}
\alcozum{%
(a) Gezgin blok (Karşı Örnek~\ref{ex:yaklinsama_karsiornek}).\\
(b) İmkânsız: $X_n \xrightarrow{d} c$ (dejenere) $\Leftrightarrow$ $X_n \xrightarrow{P} c$.
Çünkü $P(|X_n-c|>\varepsilon) = P(X_n \leq c-\varepsilon) + P(X_n > c+\varepsilon) \to 0+0=0$.}

\begin{alistirma}[Al.7.6]{A}{0-1 Kanunu Uygulaması}
$X_1, X_2, \ldots$ bağımsız. Aşağıdaki olayların olasılığını belirtin:
(a) $A = \{\bar{X}_n \to 0\}$.
(b) $B = \{\sum_{n=1}^\infty X_n/n \text{ yakınsar}\}$.
(c) $C = \{\limsup_n X_n < \infty\}$.
\end{alistirma}
\alcozum{%
Hepsi kuyruk $\sigma$-cebirinde. 0-1 Kanunu: $P(A), P(B), P(C) \in \{0,1\}$.
Belirli dağılıma bağlı olmadan bu sonuç kesin — spesifik değer ayrıca hesaplanmalı.}

\begin{alistirma}[Al.7.7]{B}{Borel-Cantelli — Kuvvetli Kanun}
$X_1, X_2, \ldots$ bağımsız, $E[X_i] = \mu$, $\mathrm{Var}(X_i) = \sigma^2 < \infty$.
Borel-Cantelli 1'i kullanarak $\bar{X}_{n^2} \xrightarrow{\text{n.k.}} \mu$ olduğunu ispatlayın.
($n^2$'nin kare indeksler olduğuna dikkat: $\bar{X}$ dizisinin alt dizisi.)
\end{alistirma}
\alcozum{%
$\varepsilon > 0$. Chebyshev: $P(|\bar{X}_{n^2}-\mu|>\varepsilon) \leq \sigma^2/(n^2\varepsilon^2)$.
$\sum_n 1/n^2 < \infty$: $\sum_n P(|\bar{X}_{n^2}-\mu|>\varepsilon) < \infty$.
Borel-Cantelli 1: $P(|\bar{X}_{n^2}-\mu|>\varepsilon \text{ i.o.}) = 0$.
Her $\varepsilon>0$ için geçerli: $\bar{X}_{n^2} \xrightarrow{\text{n.k.}} \mu$.}

\begin{alistirma}[Al.7.8]{B}{Slutsky Uygulaması}
$X_n \xrightarrow{d} \Normal(0,1)$ ve $S_n^2 \xrightarrow{P} \sigma^2 > 0$ (örnek varyans).
$T_n = X_n / S_n$'nin limit dağılımını bulun.
\end{alistirma}
\alcozum{%
$Y_n = S_n \xrightarrow{P} \sigma$; $g(y) = 1/y$ sürekli ($y\neq0$).
Sürekli haritalama: $1/S_n \xrightarrow{P} 1/\sigma$.
Slutsky çarpım: $T_n = X_n \cdot (1/S_n) \xrightarrow{d} \Normal(0,1) \cdot (1/\sigma) = \Normal(0,1/\sigma^2)$.}

\begin{alistirma}[Al.7.9]{C}{GBSK — Kolmogorov Yeterli Koşulu}
$\{X_n\}$ bağımsız (mutlaka i.i.d.\ değil), $E[X_n] = 0$,
$\sum_{n=1}^\infty \mathrm{Var}(X_n)/n^2 < \infty$.
Kronecker lemmasını kullanarak $\bar{X}_n \xrightarrow{\text{n.k.}} 0$ olduğunu gösterin.
(\textit{İpucu:} $S_n = \sum_{k=1}^n X_k/k$'nin n.k.\ yakınsadığını Borel-Cantelli ile kurun;
ardından Kronecker: $S_n \to S$ ve $S_n = \sum X_k/k$ ise $\frac{1}{n}\sum_{k=1}^n X_k \to 0$.)
\end{alistirma}
\alcozum{%
\begin{ipucu}{1}
$Y_k = X_k/k$ tanımlayın; $\sum E[Y_k^2] = \sum \mathrm{Var}(X_k)/k^2 < \infty$.
$L^2$ yakınsaması + martingal argümanı veya Borel-Cantelli ile $\sum Y_k$ n.k.\ yakınsar.
\end{ipucu}
\textbf{Kronecker Lemması:} $\sum_{k=1}^\infty a_k/k$ yakınsarsa $\frac{1}{n}\sum_{k=1}^n a_k \to 0$.
$a_k = X_k$ uygulandığında: $\sum X_k/k$ yakınsıyor (yukarıdaki ipucu) $\Rightarrow$ $\bar{X}_n \to 0$ n.k.}

\begin{alistirma}[Al.7.10]{B}{Yinelenmiş Logaritma Kanunu — Uygulama}
$X_i \sim \Normal(0,1)$ i.i.d., $S_n = \sum_{i=1}^n X_i$.
(a) YLK'ya göre $\limsup_{n\to\infty} S_n/\sqrt{2n\ln\ln n}$'nin değerini verin.
(b) $n = 10^4$ için $\sqrt{2n\ln\ln n}$'yi sayısal olarak hesaplayın.
(c) "GBSK $\bar{X}_n \to 0$" ile "YLK $S_n/\sqrt{2n\ln\ln n} \to 1$" ifadelerinin
    çelişmediğini açıklayın.
\end{alistirma}
\alcozum{%
(a) YLK (Teorem~\ref{thm:ylk}): $\sigma^2 = 1$ olduğundan $\limsup = 1$ n.k.\\
(b) $n = 10^4$: $\ln n = \ln 10^4 = 4\ln 10 \approx 9.21$; $\ln\ln n = \ln 9.21 \approx 2.22$.
$\sqrt{2 \times 10^4 \times 2.22} \approx \sqrt{44400} \approx 210.7$.
Yani $S_{10^4}$ sonsuz kez $\approx 210.7$'ye ulaşır.\\
(c) $\bar{X}_n = S_n/n \to 0$ (GBSK); $S_n/\sqrt{2n\ln\ln n} \to 1$ (YLK).
$n^{-1} = o(n^{-1/2}(\ln\ln n)^{-1/2})$ olduğundan iki ifade farklı ölçek dilleridir.
$\bar{X}_n \to 0$ yavaş yavaş gerçekleşir; YLK tam ne kadar yavaş olduğunu söyler.}

\begin{alistirma}[Al.7.11]{B}{Yenileme Teorisi}
Bir makine parçası $X_i \sim \Exp(2)$ dağılımına sahip ömürler geçiriyor ($i = 1, 2, \ldots$).
Parça bozulduğunda hemen yenisiyle değiştiriliyor. $N(t)$ yenileme sayısı.
(a) Temel Yenileme Teoremi ile $\lim_{t\to\infty} N(t)/t$'yi bulun.
(b) Blackwell teoremini kullanarak $[100, 101]$ aralığındaki beklenen yenileme sayısını yaklaşık hesaplayın.
(c) $N(t)$'nin standart sapmasının $\sim \sqrt{\sigma^2 t/\mu^3}$ ölçeğinde olduğunu yorumlayın
    ($\sigma^2 = \mathrm{Var}(X_1) = 1/4$, $\mu = 1/2$).
\end{alistirma}
\alcozum{%
$\mu = E[X_1] = 1/2$, $\sigma^2 = \mathrm{Var}(X_1) = 1/4$ ($\Exp(2)$ için).\\
(a) $N(t)/t \xrightarrow{\text{n.k.}} 1/\mu = 2$. Uzun vadede dakikada 2 yenileme.\\
(b) Blackwell: $E[N(101)-N(100)] \approx h/\mu = 1/(1/2) = 2$.
$[100,101]$ aralığında beklenen 2 yenileme.\\
(c) $\sigma^2 t/\mu^3 = (1/4) \cdot t/(1/8) = 2t$.
$\mathrm{SD}(N(t)) \approx \sqrt{2t}$: büyük $t$'de $N(t)$ etrafında $\sqrt{2t}$ ölçeğinde salınır.}

\begin{alistirma}[Al.7.12]{C}{Ergodik Markov Zinciri}
$Z_n$ iki durumlı Markov zinciri, $S = \{0, 1\}$, geçiş matrisi
$P = \begin{pmatrix}1-\alpha & \alpha \\ \beta & 1-\beta\end{pmatrix}$
($0 < \alpha, \beta < 1$).
(a) Durağan dağılım $\pi$'yi bulun.
(b) $f(0) = 0$, $f(1) = 1$ ise Birkhoff Ergodik Teoremi'ne göre
    $\frac{1}{n}\sum_{k=0}^{n-1} f(Z_k)$'nın sınırını verin.
(c) Bu sonucu BSK ile karşılaştırın: bağımsızlık gerekiyor mu?
\end{alistirma}
\alcozum{%
(a) $\pi P = \pi$: $\pi_0(1-\alpha) + \pi_1 \beta = \pi_0$ ve $\pi_0 + \pi_1 = 1$.
$\pi_0 \alpha = \pi_1 \beta \Rightarrow \pi_1/\pi_0 = \alpha/\beta$.
$\pi_0 = \beta/(\alpha+\beta)$, $\pi_1 = \alpha/(\alpha+\beta)$.\\
(b) Zincir geri dönüşümlü ve pozitif tekrarlı (sonlu durum, $\alpha,\beta > 0$) $\Rightarrow$ ergodik.
Birkhoff: $\frac{1}{n}\sum f(Z_k) \xrightarrow{\text{n.k.}} E_\pi[f] = \pi_1 = \alpha/(\alpha+\beta)$.\\
(c) BSK: bağımsızlık + i.i.d.\ gerekir. Ergodik teorem: yalnızca durağanlık + ergodiklik.
Markov zinciri bağımlı olduğu hâlde "zamansal ortalama = uzamsal ortalama" geçerlidir.
Bağımsızlık, ergodikliğin özel bir durumudur.}

% =====================================================================
\endinput

% =====================================================================
%  BÖLÜM 8: Merkezi Limit Teoremi ve Asimptotik Teori
%  Kaynak: Feller Vol.II; Billingsley; Durrett; Klenke; Gnedenko
% =====================================================================
\chapter{Merkezi Limit Teoremi ve Asimptotik Teori}
\label{chp:merkezi_limit}

\bolumalinti{%
  Merkezi Limit Teoremi, olasılık teorisinin en güzel ve en derin sonuçlarından biridir.
  Normal dağılım, sanki evrenin derinliklerine işlenmiş gibidir.%
}{William Feller, \textit{An Introduction to Probability Theory and Its Applications}, Cilt II}

% ---- Kazanımlar (TYMM) ----
\begin{kazanimlar}
Bu bölümün sonunda öğrenci:
\begin{itemize}[leftmargin=1.8em]
  \item Merkezi Limit Teoremi'ni (MLT) ifade edip karakteristik fonksiyon yöntemiyle ispatlayabilir;
        De Moivre-Laplace özel durumunu tarihsel bağlamda açıklayabilir.
  \item Lindeberg koşulunu tanımlayabilir; bağımsız ama özdeş dağılımlı olmayan
        değişkenler için MLT'nin geçerli olduğu koşulları açıklayabilir.
  \item Çok boyutlu MLT'yi ifade edip kovaryans yapısıyla ilişkilendirebilir.
  \item Berry-Esseen teoremini kullanarak yakınsama hızına sayısal üst sınır verebilir.
  \item Yinelenmiş logaritma kanununun MLT ile ilişkisini sezgisel düzeyde açıklayabilir.
  \item Delta yöntemini büyük örneklem istatistiğine uygulayabilir;
        normal yaklaşım kalitesini değerlendirebilir.
\end{itemize}
\end{kazanimlar}

% ---- Motivasyon ----
\begin{motivasyon}
Büyük Sayılar Kanunu der ki: $\bar{X}_n \to \mu$. Soru şu: bu yakınsama \emph{ne kadar hızlı}
gerçekleşir ve $\bar{X}_n$ ile $\mu$ arasındaki fark \emph{nasıl dağılır}?

Merkezi Limit Teoremi yanıtlar: $\sqrt{n}(\bar{X}_n - \mu)/\sigma$ dağılımı,
$X_i$'lerin dağılımından \emph{bağımsız} olarak Normal(0,1)'e yakınsar.
Bu olağanüstü evrensellik, normal dağılımın olasılık ve istatistiğin merkezine
yerleşmesinin matematiksel gerekçesidir.

\medskip
\textbf{Köprü.} Bölüm~\ref{chp:beklenti}'deki karakteristik fonksiyon,
Bölüm~\ref{chp:donusumler}'deki delta yöntemi ve Bölüm~\ref{chp:buyuk_sayilar}'daki
Slutsky teoremi bu bölümün araçlarıdır.
Bölüm~\ref{chp:istatistik_temelleri}'nde büyük örneklem istatistiği bu sonuçlara dayanır.
\end{motivasyon}

% =====================================================================
\section{De Moivre-Laplace Teoremi}
\label{sec:demoivre_laplace}
% =====================================================================

\begin{tarihselnot}
Abraham de Moivre (1733) binom dağılımının büyük $n$ için normale yakınsadığını
fark etti. Pierre-Simon Laplace (1812) bunu daha geniş bir çerçevede ifade etti.
Carl Friedrich Gauss ise normal dağılımı hata teorisi bağlamında bağımsız olarak geliştirdi.
Genel i.i.d.\ versiyonu Aleksandr Lyapunov (1901) ve Jarl Waldemar Lindeberg (1922) tarafından
ispatlandı.
\end{tarihselnot}

\begin{teorem}[De Moivre-Laplace]
\label{thm:demoivre_laplace}
$X_i \sim \mathrm{Bernoulli}(p)$ bağımsız, $S_n = \sum_{i=1}^n X_i \sim \Binom(n,p)$.
$\mu = np$, $\sigma^2 = npq$. O zaman her $a < b$'de
\[
  P\!\left(a \leq \frac{S_n - np}{\sqrt{npq}} \leq b\right) \to \Phi(b) - \Phi(a) \quad (n\to\infty).
\]
\end{teorem}

\begin{proof}[İskelet İspat — Stirling Yaklaşımı]
Stirling formülü: $k! \approx \sqrt{2\pi k}(k/e)^k$.
$P(S_n = k)$ için $\binom{n}{k}p^kq^{n-k}$'yi Stirling ile açarak
$k = np + x\sqrt{npq}$ ($x$ sabit) için bu olasılığın
$\phi(x)/\sqrt{npq}$'ya yakınsadığı gösterilir (yerel limit teoremi).
Yerel limit toplamından dağılımla yakınsama elde edilir.
\end{proof}

\begin{ornek}[Binom Normal Yaklaşımı]
$n = 100$, $p = 0.4$: $\mu = 40$, $\sigma = \sqrt{24} \approx 4.90$.
$P(35 \leq S_{100} \leq 45) \approx \Phi\!\left(\frac{45.5-40}{4.90}\right) - \Phi\!\left(\frac{34.5-40}{4.90}\right)$
$= \Phi(1.12) - \Phi(-1.12) \approx 0.7372$.
(Süreklilik düzeltmesi: ayrık $\to$ sürekli, sınırlara $\pm 0.5$ eklenir.)
\end{ornek}

% =====================================================================
\section{Genel Merkezi Limit Teoremi}
\label{sec:genel_mlt}
% =====================================================================

\begin{teorem}[Merkezi Limit Teoremi — i.i.d.\ Durum]
\label{thm:mlt}
$X_1, X_2, \ldots$ i.i.d., $E[X_1] = \mu$, $\mathrm{Var}(X_1) = \sigma^2 \in (0,\infty)$. O zaman
\[
  Z_n = \frac{\bar{X}_n - \mu}{\sigma/\sqrt{n}} = \frac{\sum_{i=1}^n (X_i - \mu)}{\sigma\sqrt{n}}
  \xrightarrow{d} \Normal(0,1).
\]
\end{teorem}

\begin{proof}[Karakteristik Fonksiyon İspati]
$Y_i = (X_i - \mu)/\sigma$ tanımı: $E[Y_i] = 0$, $E[Y_i^2] = 1$.
$\varphi_Y(t) := E[e^{itY_1}]$ KF'si. Taylor açılımı ($E[Y_1]=0$, $E[Y_1^2]=1$):
\[
  \varphi_Y(t) = 1 - \frac{t^2}{2} + o(t^2) \quad (t\to 0).
\]
$Z_n = \frac{1}{\sqrt{n}}\sum_{i=1}^n Y_i$; bağımsızlık + KF çarpım:
\[
  \varphi_{Z_n}(t) = \left[\varphi_Y\!\left(\frac{t}{\sqrt{n}}\right)\right]^n
  = \left[1 - \frac{t^2}{2n} + o\!\left(\frac{1}{n}\right)\right]^n.
\]
$n\to\infty$ limiti: $\left(1 - \frac{t^2}{2n} + o(1/n)\right)^n \to e^{-t^2/2}$.
(Temel limit: $(1+a_n/n)^n \to e^a$ eğer $na_n \to a$.)
$e^{-t^2/2}$: standart normalin KF'si.
Lévy süreklilik teoremi (Teorem~\ref{prop:kf_yaklinsama}): $Z_n \xrightarrow{d} \Normal(0,1)$.
\end{proof}

\begin{kritiksorgu}
MLT hangi koşulları \emph{gerektirmez}?
\begin{itemize}
  \item $X_i$ normal olmak \emph{zorunda değil} (herhangi bir dağılım).
  \item $E[X_i^3] < \infty$ \emph{gerekmez} — yalnızca $\sigma^2 < \infty$ yeterli.
  \item $X_i$'lerin sürekli olması \emph{gerekmez} (kesikli dağılımlar da MLT'ye tabidir).
\end{itemize}
Ne gerekir: i.i.d. + sonlu varyans. Bunlar ortadan kalkarsa MLT bozulabilir
(Cauchy: varyans yok; ağır kuyruklar: stabil dağılımlara yakınsama).
\end{kritiksorgu}

\begin{onerme}[Yararlı Eşdeğer Formlar]
\label{prop:mlt_form}
MLT'den türetilen ifadeler:
\begin{align*}
  \sqrt{n}(\bar{X}_n - \mu) &\xrightarrow{d} \Normal(0, \sigma^2), \\
  \frac{S_n - n\mu}{\sigma\sqrt{n}} &\xrightarrow{d} \Normal(0,1), \\
  P\!\left(\bar{X}_n \leq \mu + \frac{z_\alpha \sigma}{\sqrt{n}}\right) &\to \Phi(z_\alpha) = 1-\alpha.
\end{align*}
\end{onerme}

% =====================================================================
\section{Lindeberg-Feller Koşulları}
\label{sec:lindeberg}
% =====================================================================

\begin{kesif}
i.i.d.\ koşulunu kaldırırsak ne olur? $X_i$'lerin farklı dağılımları olsa bile
toplamları normale yakınsar mı? Hangi koşul "hiçbir terim baskın değil"
sezgisini matematiksel olarak yakalar?
\end{kesif}

\begin{tanim}[Lindeberg Koşulu]
\label{def:lindeberg}
$\{X_{ni}\}$ bağımsız, $E[X_{ni}] = 0$, $\sigma_{ni}^2 = \mathrm{Var}(X_{ni})$,
$s_n^2 = \sum_{i=1}^n \sigma_{ni}^2$. \textbf{Lindeberg koşulu}: her $\varepsilon > 0$ için
\[
  L_n(\varepsilon) = \frac{1}{s_n^2} \sum_{i=1}^n E\!\left[X_{ni}^2 \mathbf{1}_{|X_{ni}|>\varepsilon s_n}\right] \to 0.
\]
\end{tanim}

\begin{teorem}[Lindeberg-Feller MLT]
\label{thm:lindeberg_feller}
Lindeberg koşulu sağlanıyorsa
$\frac{1}{s_n}\sum_{i=1}^n X_{ni} \xrightarrow{d} \Normal(0,1)$.
Tersine (Feller): Lindeberg koşulu, $\max_i \sigma_{ni}/s_n \to 0$ ("ufaklık koşulu")
altında MLT için zorunludur.
\end{teorem}

\begin{proof}[Ana Fikir]
Lindeberg koşulu, her terimin katkısının $s_n$'e kıyasla küçük olmasını garantiler.
KF ispatında $\varphi_{X_{ni}/s_n}(t) = 1 - t^2\sigma_{ni}^2/(2s_n^2) + R_{ni}(t)$;
hata terimi $|R_{ni}(t)| \leq E[X_{ni}^2\mathbf{1}_{|X_{ni}|>\varepsilon s_n}] \cdot t^2/s_n^2$.
Lindeberg koşulu ile hataların toplamı sıfıra gider;
$\prod \varphi_{X_{ni}/s_n}(t) \to e^{-t^2/2}$ sonucuna varılır.
Ayrıntı: Billingsley~\cite{Billingsley1995}, §27.
\end{proof}

\begin{onerme}[Lyapunov Koşulu $\Rightarrow$ Lindeberg]
\label{prop:lyapunov}
$\sum_{i=1}^n E[|X_{ni}|^{2+\delta}] / s_n^{2+\delta} \to 0$ (bazı $\delta > 0$)
$\Rightarrow$ Lindeberg koşulu $\Rightarrow$ MLT.
Lyapunov koşulu daha hesaplanabilir ama Lindeberg'den daha kuvvetlidir.
\end{onerme}

\begin{ornek}[Lindeberg Koşulunun Önemi]
$n$ işçinin maaş artışı: $X_i \sim \Normal(0, \sigma_i^2)$, $\sigma_i^2 = i$.
$s_n^2 = \sum i = n(n+1)/2$.
Lyapunov ($\delta=1$): $\sum E[|X_i|^3]/s_n^3 \propto \sum i^{3/2}/n^{3/2} \to 0$.
Toplam büyük olduğunda bile MLT geçerli; Lindeberg sağlanır.
\end{ornek}

% =====================================================================
\section{Çok Boyutlu Merkezi Limit Teoremi}
\label{sec:cok_boyut_mlt}
% =====================================================================

\begin{teorem}[Çok Boyutlu MLT]
\label{thm:cok_boyut_mlt}
$\mathbf{X}_1, \mathbf{X}_2, \ldots$ i.i.d.\ $\mathbb{R}^k$ değerli,
$E[\mathbf{X}_i] = \boldsymbol{\mu}$, kovaryans matrisi $\boldsymbol{\Sigma}$ (pozitif yarı-tanımlı). O zaman
\[
  \sqrt{n}(\bar{\mathbf{X}}_n - \boldsymbol{\mu}) \xrightarrow{d} \Normal_k(\mathbf{0},\, \boldsymbol{\Sigma}).
\]
\end{teorem}

\begin{proof}
\textbf{Cramér-Wold cihazı:} $\mathbf{Y}_n \xrightarrow{d} \mathbf{Y}$ ($\mathbb{R}^k$'de) ancak ve ancak
her $\mathbf{a} \in \mathbb{R}^k$ için $\mathbf{a}^\top \mathbf{Y}_n \xrightarrow{d} \mathbf{a}^\top \mathbf{Y}$.

$\mathbf{a}^\top \sqrt{n}(\bar{\mathbf{X}}_n - \boldsymbol{\mu}) = \sqrt{n}(\overline{\mathbf{a}^\top\mathbf{X}_n} - \mathbf{a}^\top\boldsymbol{\mu})$.

$\mathbf{a}^\top\mathbf{X}_i$: i.i.d., ortalama $0$, varyans $\mathbf{a}^\top\boldsymbol{\Sigma}\mathbf{a}$.
Tek boyutlu MLT: $\xrightarrow{d} \Normal(0, \mathbf{a}^\top\boldsymbol{\Sigma}\mathbf{a})$.
Bu $\Normal_k(\mathbf{0},\boldsymbol{\Sigma})$'nın doğrusal kombinasyonudur; Cramér-Wold tamamlar.
\end{proof}

\begin{ornek}[Çok Boyutlu Normal Yaklaşım]
$(X_i, Y_i)$ i.i.d., $E[\mathbf{X}] = (\mu_X, \mu_Y)^\top$,
$\boldsymbol{\Sigma} = \begin{pmatrix}\sigma_X^2 & \rho\sigma_X\sigma_Y \\ \rho\sigma_X\sigma_Y & \sigma_Y^2\end{pmatrix}$.
Büyük $n$: $(\bar{X}_n, \bar{Y}_n)$ yaklaşık $\Normal_2\!\left((\mu_X,\mu_Y)^\top,\, \boldsymbol{\Sigma}/n\right)$.
Bu, iki boyutlu güven bölgelerinin (elips) temelidir.
\end{ornek}

% =====================================================================
\section{Berry-Esseen Teoremi}
\label{sec:berry_esseen}
% =====================================================================

\begin{motivasyon}[MLT'nin Hızı Nedir?]
MLT $Z_n \xrightarrow{d} \Normal(0,1)$ diyor. Ama $n = 30$ iken ne kadar yakınız?
Berry-Esseen teoremi hata büyüklüğüne sayısal üst sınır koyar.
\end{motivasyon}

\begin{teorem}[Berry-Esseen]
\label{thm:berry_esseen}
$X_1, \ldots, X_n$ i.i.d., $E[X_1] = 0$, $E[X_1^2] = \sigma^2 > 0$,
$E[|X_1|^3] = \rho < \infty$. O zaman her $x \in \mathbb{R}$ için
\[
  \left|F_{Z_n}(x) - \Phi(x)\right| \leq \frac{C\,\rho}{\sigma^3\sqrt{n}},
\]
burada $C \leq 0.4748$ (en iyi bilinen sabit; Shevtsova 2010).
\end{teorem}

İspatın ana fikri: KF'ler farkına Fourier ters dönüşümü uygulanması;
"yumuşatma" (smoothing) lemması aracılığıyla KDF farkına bağlanması.
Ayrıntı için Feller Vol.II~\cite{Feller1971}, Bölüm XVI.5 veya Durrett~\cite{Durrett2019}, §3.4.

\begin{ornek}[Berry-Esseen Uygulaması]
$X_i \sim \mathrm{Bernoulli}(0.5)$: $\sigma^2 = 0.25$, $\rho = E[|X_1|^3] = 0.5$.
$n = 36$: hata $\leq C \cdot 0.5/(0.25^{3/2}\cdot 6) = C \cdot 0.5/(0.125\cdot 6) \approx 0.663C \approx 0.315$.
$n = 100$: hata $\leq C\cdot 0.5/(0.125\cdot 10) \approx 0.190$.
$n = 1000$: hata $\leq 0.060$: 3 ondalık basamak için $n \approx 1000$ önerilir.
\end{ornek}

\begin{hataanalizi}
"$n \geq 30$ ise MLT yeterince iyi yaklaşır" — çok kaba bir kuraldır.
Berry-Esseen: hata üçüncü mutlak moment ile orantılı.
Çarpık dağılımlarda ($\rho$ büyük) $n = 30$ yetersizdir;
çok çarpık dağılımlar (lognormal, Pareto) için $n = 1000$ bile tartışmalıdır.
\end{hataanalizi}

% =====================================================================
\section{Poisson Yakınsaması}
\label{sec:poisson_yaklinsama}
% =====================================================================

\begin{teorem}[Le Cam Teoremi]
\label{thm:le_cam}
$X_1, \ldots, X_n$ \emph{bağımsız} (özdeş dağılımlı olmak zorunda değil),
$X_i \sim \mathrm{Bernoulli}(p_i)$, $S_n = \sum X_i$.
$\lambda_n = \sum p_i$. Toplam varyasyon mesafesi:
\[
  d_{TV}(S_n, \Pois(\lambda_n)) \leq 2\sum_{i=1}^n p_i^2 = 2\sum p_i^2.
\]
\end{teorem}

\begin{proof}[Ana Fikir — Stein-Chen Yöntemi]
Poisson operatörü: $\mathcal{A}f(k) = \lambda f(k+1) - kf(k)$.
$E[\mathcal{A}f(S_n)]$ ile $E[\mathcal{A}f(\Pois(\lambda))]$ arasındaki fark,
$\sum p_i^2$ terimlerinden beslenir.
Ayrıntı: Arratia-Goldstein-Gordon (1989) veya Durrett~\cite{Durrett2019}, §3.6.
\end{proof}

\begin{ornek}[Le Cam Uygulaması]
$n = 200$, $p_i = 1/200$ homojen: $\lambda = 1$, $\sum p_i^2 = 200 \cdot (1/200)^2 = 1/200$.
$d_{TV} \leq 2/200 = 0.01$: Poisson yaklaşımı $S_{200}$ için çok iyidir.
Bu, Bölüm~\ref{subsec:poisson}'daki binom-Poisson limitinin nicel versiyonudur.
\end{ornek}

% =====================================================================
\section{Yinelenmiş Logaritma Kanunu}
\label{sec:yinelenmis_log}
% =====================================================================

\begin{motivasyon}[MLT'nin Ötesi]
MLT: $S_n/({\sigma\sqrt{n}}) \xrightarrow{d} \Normal(0,1)$. Yani $S_n \approx \sigma\sqrt{n}$ büyüklüğünde.
Ama $S_n/(\sigma\sqrt{n})$'nin üstel rastlanımsal salınımı ne kadar büyük olabilir?
Sezgi: bazen $k\sigma\sqrt{n}$'ye ($k > 1$) çıkabilir; bu "bazen"in kesin sınırı nedir?
\end{motivasyon}

\begin{teorem}[Yinelenmiş Logaritma Kanunu (YLK)]
\label{thm:ylk}
$X_1, X_2, \ldots$ i.i.d., $E[X_1] = 0$, $\mathrm{Var}(X_1) = \sigma^2 \in (0,\infty)$,
$S_n = \sum_{i=1}^n X_i$. O zaman
\[
  \limsup_{n\to\infty} \frac{S_n}{\sigma\sqrt{2n\ln\ln n}} = 1 \quad \text{n.k.}
\]
ve $\liminf_{n\to\infty} S_n/(\sigma\sqrt{2n\ln\ln n}) = -1$ n.k.
\end{teorem}

İspat oldukça teknik olup Kolmogorov'un orijinal 1929 kanıtına, daha sonra
Hartman-Wintner (1941) genel versiyonuna dayanır.
Temel fikir: uygun olayların bağımsızlığı + Borel-Cantelli lemmalarının dikkatli uygulanması.

\begin{onerme}[YLK ve MLT Karşılaştırması]
\begin{center}
\begin{tabular}{lcc}
\toprule
 & MLT & YLK \\
\midrule
Normalize eden & $\sigma\sqrt{n}$ & $\sigma\sqrt{2n\ln\ln n}$ \\
Limit türü & Dağılımla ($\Normal(0,1)$) & n.k.\ ($\limsup = 1$) \\
Anlam & Tipik davranış & Ekstrem salınım \\
\bottomrule
\end{tabular}
\end{center}
\end{onerme}

\begin{ornek}[YLK Yorumu]
$X_i \sim \Normal(0,1)$ bağımsız. $S_n$ neredeyse kesinlikle
$\pm\sqrt{2n\ln\ln n}$ bandının içinde kalır; bu bandı sonsuz kez tam olarak dokunur
ama aşmaz (kesin üst sınır).
$n = 10^6$: bant genişliği $\approx \sqrt{2 \times 10^6 \times \ln\ln(10^6)} \approx \sqrt{2\times10^6\times2.93} \approx 2420$.
MLT bandı $\sqrt{n} \approx 1000$; YLK bant $\approx 2.4$ kat daha geniş.
\end{ornek}

% =====================================================================
\section{Büyük Örneklemde Normal Yaklaşımlar}
\label{sec:buyuk_orneklem}
% =====================================================================

\begin{onerme}[Örneklem Ortalaması — Güven Aralığı]
\label{prop:guven_araligi}
$X_1, \ldots, X_n$ i.i.d., $\mu$ bilinmiyor, $\sigma^2$ biliniyor. MLT'den:
\[
  P\!\left(\bar{X}_n - z_{\alpha/2}\frac{\sigma}{\sqrt{n}} \leq \mu \leq \bar{X}_n + z_{\alpha/2}\frac{\sigma}{\sqrt{n}}\right) \to 1-\alpha.
\]
$z_{0.025} = 1.96$: yaklaşık $\%95$ güven aralığı $\bar{X}_n \pm 1.96\sigma/\sqrt{n}$.
\end{onerme}

\begin{onerme}[Örneklem Oranı]
\label{prop:oran_guven}
$X_i \sim \mathrm{Bernoulli}(p)$, $\hat{p} = \bar{X}_n$. MLT:
$\sqrt{n}(\hat{p}-p)/\sqrt{p(1-p)} \xrightarrow{d} \Normal(0,1)$.
$p$ bilinmiyorsa $\sqrt{\hat{p}(1-\hat{p})}$ ile değiştirin (Slutsky):
yaklaşık $\%95$ GA: $\hat{p} \pm 1.96\sqrt{\hat{p}(1-\hat{p})/n}$.
\end{onerme}

\begin{proof}
$\hat{p}$ yerine $p$ yazılan versiyonu MLT'nin doğrudan uygulaması.
$\sqrt{\hat{p}(1-\hat{p})} \xrightarrow{P} \sqrt{p(1-p)}$ (BSK + sürekli haritalama).
Slutsky: bölüm dağılımla $\Normal(0,1)$'e yakınsar.
\end{proof}

\begin{teorem}[Delta Yöntemi — MLT Uygulaması]
\label{thm:delta_mlt}
$\sqrt{n}(\bar{X}_n - \mu) \xrightarrow{d} \Normal(0,\sigma^2)$ ve $g$, $\mu$'da türevlenebilir ise
(Teorem~\ref{thm:delta_yontemi} ile birlikte):
\[
  \sqrt{n}(g(\bar{X}_n) - g(\mu)) \xrightarrow{d} \Normal(0,\, \sigma^2[g'(\mu)]^2).
\]
Bu, Bölüm~\ref{chp:nokta_tahmin}'deki asimptotik güven aralıklarının temelidir.
\end{teorem}

\begin{ornek}[Varyans Stabilize Eden Dönüşüm]
$X_i \sim \mathrm{Bernoulli}(p)$: $\sigma^2 = p(1-p)$, $p$'ye bağlı.
$g(p) = \arcsin(\sqrt{p})$: $g'(p) = 1/(2\sqrt{p(1-p)})$.
Delta: $[g'(p)]^2 \sigma^2 = p(1-p)/(4p(1-p)) = 1/4$ — sabit!
$\sqrt{n}(\arcsin\sqrt{\hat{p}} - \arcsin\sqrt{p}) \xrightarrow{d} \Normal(0, 1/4)$.
Varyans artık $p$'den bağımsız: güven aralıkları her $p$ için eşit genişlikte.
\end{ornek}

\begin{onerme}[Poisson Gözlemleri için Normal Yaklaşım]
$X \sim \Pois(\lambda)$, $\lambda$ büyük:
$(\sqrt{X} - \sqrt{\lambda}) \xrightarrow{d} \Normal(0, 1/4)$ (karekök dönüşümü).
$P(X \leq k) \approx \Phi\!\left(\frac{k+0.5-\lambda}{\sqrt{\lambda}}\right)$ (süreklilik düzeltmesi).
\end{onerme}

% ---- Disiplinlerarası Bağlantı (TYMM) ----
\begin{kopru}[Disiplinlerarası — Fizik, Sosyal Bilimler, Biyoistatistik]
\textbf{Fizik:} Termal dengedeki gaz moleküllerinin Maxwell-Boltzmann hız dağılımı,
çok sayıda bağımsız çarpışmanın MLT sonucu olarak Gaussian'a yakınsar.
Hata teorisi (Gauss): ölçüm hatalarının bağımsız küçük bileşenlerden oluştuğu
varsayımı altında toplam hata normal dağılır.\\
\textbf{Sosyal bilimler:} Büyük örneklem istatistiği (anket, seçim tahminleri, A/B testi)
%\%95 güven aralığı MLT'ye dayanır. "Hata payı $\pm 3\%$" beyanı normal yaklaşımın uygulamasıdır.\\
\textbf{Biyoistatistik:} Klinik denemelerde $t$-testi, MLT'nin küçük örneklemlerde
doğrudan uygulaması değil, $t$ dağılımının kullanımıdır; büyük $n$'de $t_n \to \Normal(0,1)$.
\end{kopru}

% ---- Öz-Yansıtma ----
\begin{ozyansima}
\begin{itemize}[leftmargin=1.8em]
  \item MLT'nin "evrenselliği" size sezgisel olarak neden doğru geliyor (ya da gelmiyor)?
  \item Berry-Esseen teoremi "$n \geq 30$ yeterli" kuralını nasıl sorgular?
  \item Lindeberg koşulu intuisyon olarak ne anlama gelir?
        "Hiçbir terim baskın değil" yorumunu bir örnekle somutlaştırın.
  \item YLK ile MLT arasındaki farkı bir grafikle anlatın (sözlü olarak tarif edin).
\end{itemize}
\end{ozyansima}

% ---- Köprü Notu ----
\begin{kopru}
\textbf{Önceki bölüm:} Bölüm~\ref{chp:buyuk_sayilar} BSK'yı verdi (yakınsama garantisi);
bu bölüm yakınsama \emph{hızını} ve \emph{dağılımını} belirledi.\\
\textbf{Sonraki (opsiyonel):} Bölüm~\ref{chp:martingaller} — koşullu beklenti
ve martingallar; martingal MLT.\\
\textbf{İstatistik kısmı:} Bölüm~\ref{chp:istatistik_temelleri}'nde asimptotik
dağılım teorisi; Bölüm~\ref{chp:nokta_tahmin}'de MLE'nin asimptotik normalliği
($\sqrt{n}(\hat\theta_{MLE}-\theta) \xrightarrow{d} \Normal(0, I(\theta)^{-1})$).
\defterbaglantisi{bolum08\_mlt.ipynb}{%
  MLT simülasyonu (çeşitli dağılımlar), Berry-Esseen görselleştirmesi, YLK yolu}
\end{kopru}

% =====================================================================
%  ALIŞTIRMALAR
% =====================================================================

\cozumlualistirmalar

\begin{alistirma}[Al.8.1]{A}{MLT Uygulaması}
Bir sigorta şirketinin yıllık hasar başına ödeme miktarı $E[X] = 1000$ TL,
$\sigma = 200$ TL. $n = 100$ hasar için toplam ödemenin $105\,000$ TL'yi aşma
olasılığını MLT ile tahmin edin.
\end{alistirma}
\alcozum{%
$S_{100}$: $E[S_{100}] = 100\,000$, $\mathrm{Var}(S_{100}) = 100 \times 200^2 = 4\,000\,000$,
$\sigma_{S_{100}} = 2000$.
$P(S_{100} > 105\,000) = P\!\left(Z > \frac{105000-100000}{2000}\right) = P(Z > 2.5) = 1-\Phi(2.5) \approx 0.0062$.}

\begin{alistirma}[Al.8.2]{A}{De Moivre-Laplace}
$S_{400} \sim \Binom(400, 0.4)$.
(a) $P(150 \leq S_{400} \leq 170)$'i MLT ile hesaplayın (süreklilik düzeltmesiyle).
(b) $E[S_{400}]$ ve $\sigma_{S_{400}}$'ü hesaplayın.
\end{alistirma}
\alcozum{%
(b) $E[S_{400}]=160$, $\sigma=\sqrt{400\cdot0.4\cdot0.6}=\sqrt{96}\approx9.798$.\\
(a) Süreklilik düzeltmesiyle: $P(149.5\leq S\leq170.5)$.
$z_1=(149.5-160)/9.798\approx-1.072$, $z_2=(170.5-160)/9.798\approx1.072$.
$P\approx\Phi(1.072)-\Phi(-1.072)=2\Phi(1.072)-1\approx2(0.8582)-1=0.7164$.}

\begin{alistirma}[Al.8.3]{B}{Berry-Esseen Sınırı}
$X_i \sim \mathrm{Exp}(1)$ bağımsız: $\mu=1$, $\sigma^2=1$, $E[|X_1-1|^3]$'ü hesaplayın.
Berry-Esseen ile $n=50$ için maksimum KDF hatasına üst sınır verin.
\end{alistirma}
\alcozum{%
$Y_i = X_i-1$ için $E[|Y_i|^3]$.
$E[|X-1|^3] = \int_0^\infty |x-1|^3 e^{-x}dx = \int_0^1(1-x)^3e^{-x}dx + \int_1^\infty(x-1)^3e^{-x}dx$.
Entegrasyon by parts (veya mgf türevi): $E[|X-1|^3] = 2$ (Exp(1) için).
$\sigma^3 = 1$, $\rho = 2$.
Berry-Esseen: hata $\leq C\cdot2/(\sqrt{50})\approx 0.4748\cdot2/7.071\approx0.134$.
$n=50$, üstel dağılım için maksimum KDF hatası yaklaşık $0.134$.}

\begin{alistirma}[Al.8.4]{B}{Delta Yöntemi + MLT}
$X_1,\ldots,X_n$ i.i.d.\ $\mathrm{Pois}(\lambda)$, $\bar{X}_n$ örneklem ortalaması.
(a) MLT: $\sqrt{n}(\bar{X}_n-\lambda)$'nın limit dağılımını yazın.
(b) $g(x) = \sqrt{x}$ dönüşümü ile $\sqrt{n}(\sqrt{\bar{X}_n}-\sqrt{\lambda})$'nın limit dağılımını bulun.
(c) (b)'nin sonucu neden uygulamada faydalıdır?
\end{alistirma}
\alcozum{%
(a) $\mathrm{Var}(X_i)=\lambda$: $\sqrt{n}(\bar{X}_n-\lambda)\xrightarrow{d}\Normal(0,\lambda)$.\\
(b) $g'(x)=1/(2\sqrt{x})$, $g'(\lambda)=1/(2\sqrt{\lambda})$.
$\sqrt{n}(\sqrt{\bar{X}_n}-\sqrt{\lambda})\xrightarrow{d}\Normal(0,\lambda/(4\lambda))=\Normal(0,1/4)$.\\
(c) Varyans $1/4$ artık $\lambda$'ya bağlı değil: güven aralığı $\sqrt{\hat\lambda}\pm z_{\alpha/2}/(2\sqrt{n})$.
Bilinmeyen parametre ile güven aralığı genişliğini kontrol etmek mümkün.}

% ---

\alistirmalar

\begin{alistirma}[Al.8.5]{A}{Örneklem Oranı}
$n=400$ kişilik ankette $180$ kişi "evet" dedi. $\hat{p}=0.45$.
(a) $p$ için yaklaşık $\%95$ güven aralığını hesaplayın.
(b) $p=0.5$ hipotezini $\alpha=0.05$ için test edin ($z$-testi).
\end{alistirma}
\alcozum{%
(a) $\hat{p}=0.45$, $\sqrt{\hat{p}(1-\hat{p})/n}=\sqrt{0.2475/400}=\sqrt{0.000619}\approx0.02488$.
GA: $0.45\pm1.96\times0.02488 = (0.401, 0.499)$.
(b) $z=(0.45-0.5)/\sqrt{0.5\times0.5/400}=(-0.05)/0.025=-2.0$.
$|z|=2.0>1.96$: $H_0$ reddedilir, $\alpha=0.05$.}

\begin{alistirma}[Al.8.6]{A}{Poisson Normal Yaklaşımı}
$X \sim \Pois(25)$. MLT ile $P(X \leq 30)$ için normal yaklaşım verin.
Süreklilik düzeltmesini kullanın.
\end{alistirma}
\alcozum{%
$E[X]=25$, $\sigma=5$.
$P(X\leq30)\approx P\!\left(Z\leq\frac{30.5-25}{5}\right)=\Phi(1.10)\approx0.8643$.
(Gerçek değer $\approx0.870$; süreklilik düzeltmesi olmadan $\Phi(1.0)\approx0.841$, daha az doğru.)}

\begin{alistirma}[Al.8.7]{B}{Lindeberg Koşulu}
$X_{ni} = \pm 1/\sqrt{n}$ eşit olasılıkla ($i = 1,\ldots,n$), bağımsız.
(a) $E[X_{ni}]=0$, $\sigma_{ni}^2=1/n$, $s_n^2=1$ olduğunu gösterin.
(b) Lindeberg koşulunun sağlandığını kontrol edin.
(c) $\sum X_{ni}$'nin limit dağılımını belirtin.
\end{alistirma}
\alcozum{%
(a) $\sigma_{ni}^2 = E[X_{ni}^2]=1/n$; $s_n^2=n\cdot(1/n)=1$. \\
(b) $L_n(\varepsilon) = \sum_i E[X_{ni}^2\mathbf{1}_{|X_{ni}|>\varepsilon s_n}]/s_n^2$.
$|X_{ni}|=1/\sqrt{n}$; $\varepsilon s_n=\varepsilon$.
$n$'yeterince büyüyünce $1/\sqrt{n}<\varepsilon$: indikatör sıfır.
$L_n(\varepsilon)=0$ (tüm büyük $n$). Sağlandı.\\
(c) Lindeberg-Feller: $\sum X_{ni}\xrightarrow{d}\Normal(0,1)$.}

\begin{alistirma}[Al.8.8]{B}{Çok Boyutlu MLT}
$(X_i, Y_i)$ i.i.d., $E[X]=E[Y]=0$, $\mathrm{Var}(X)=\mathrm{Var}(Y)=1$, $\mathrm{Cov}(X,Y)=\rho$.
$(\bar{X}_n, \bar{Y}_n)$'nin büyük örneklem dağılımını yazın.
$U_n = \bar{X}_n + \bar{Y}_n$'nin yaklaşık dağılımını bulun.
\end{alistirma}
\alcozum{%
$\boldsymbol{\Sigma} = \begin{pmatrix}1&\rho\\\rho&1\end{pmatrix}$.
$\sqrt{n}(\bar{X}_n,\bar{Y}_n)^\top \xrightarrow{d} \Normal_2(\mathbf{0},\boldsymbol{\Sigma})$.\\
$U_n = \bar{X}_n+\bar{Y}_n$: $\mathbf{a}=(1,1)^\top$.
$\mathrm{Var}(U_n)=\mathbf{a}^\top\boldsymbol{\Sigma}\mathbf{a}/n=(1+1+2\rho)/n=2(1+\rho)/n$.
$\sqrt{n}\,U_n\xrightarrow{d}\Normal(0,2(1+\rho))$.}

\begin{alistirma}[Al.8.9]{C}{MLT İspat Detayı}
$(1-t^2/(2n))^n \to e^{-t^2/2}$ limitini $n\to\infty$ için kesin biçimde ispatlayın.
Daha genel olarak: $a_n \to 0$ ve $na_n \to a$ iken $(1+a_n)^n \to e^a$ olduğunu gösterin.
\end{alistirma}
\alcozum{%
$\ln(1+a_n)^n = n\ln(1+a_n)$.
$\ln(1+x) = x - x^2/2 + O(x^3)$ ($x\to0$).
$n\ln(1+a_n) = n(a_n - a_n^2/2 + O(a_n^3))$.
$na_n \to a$; $na_n^2 = (na_n)\cdot a_n \to a\cdot0 = 0$ ($a_n\to0$).
$nO(a_n^3) = na_n \cdot a_n^2 \to 0$.
$\Rightarrow n\ln(1+a_n) \to a$.
$\Rightarrow (1+a_n)^n = e^{n\ln(1+a_n)} \to e^a$.
Özel durum: $a_n = -t^2/(2n)$, $na_n = -t^2/2 \to -t^2/2$.
$(1-t^2/(2n))^n \to e^{-t^2/2}$. \checkmark}

% =====================================================================
\endinput

% % =====================================================================
%  BÖLÜM 9: Koşullu Beklenti ve Martingaller  [OPSİYONEL / İLERİ]
%  Kaynak: Williams; Durrett Böl.4; Klenke Böl.9-12; Shiryaev Böl.VII
% =====================================================================
\chapter{Koşullu Beklenti ve Martingaller}
\label{chp:martingaller}

\bolumalinti{%
  Martingal, dürüst bir oyunun matematiksel modelidir.%
}{Joseph Leo Doob, \textit{Stochastic Processes}}

% ---- Kazanımlar (TYMM) ----
\begin{kazanimlar}
Bu bölümün sonunda öğrenci:
\begin{itemize}[leftmargin=1.8em]
  \item Koşullu beklentinin Radon-Nikodym tanımını ve temel özelliklerini açıklayabilir.
  \item Martingal, süpmartingal ve altmartingal kavramlarını tanımlayabilir; örnekler verebilir.
  \item Durma zamanı ve isteğe bağlı durma teoremini ifade edebilir.
  \item Doob'un ayrışım ve yakınsama teoremlerini kullanabilir.
  \item Martingal yakınsamasının GBSK ile bağlantısını açıklayabilir.
\end{itemize}
\end{kazanimlar}

\begin{motivasyon}
Bölüm~\ref{sec:kosullu_beklenti}'de koşullu beklentiyi $L^2$ projeksiyon olarak tanımladık.
Bu bölümde \emph{bilgi akışı} boyutunu ekliyoruz: zamanla biriken gözlemler $\mathcal{F}_0 \subseteq \mathcal{F}_1 \subseteq \cdots$
bir filtrasyon oluşturur. Martingal, "dürüst oyun" sezgisini bu dil ile yakalar:
geçmişi bilince gelecekteki beklenti bugünkü değere eşit.

\medskip
\textbf{Köprü.} Bölüm~\ref{sec:kosullu_beklenti}'deki koşullu beklenti ve
Bölüm~\ref{chp:buyuk_sayilar}'daki Borel-Cantelli bu bölümün araçlarıdır.
Martingal yakınsaması, Güçlü BSK'nın zarif bir genellemesidir.
\end{motivasyon}

% =====================================================================
\section{Filtrasyon ve Uyum}
\label{sec:filtrasyon}
% =====================================================================

\begin{tanim}[Filtrasyon]
\label{def:filtrasyon}
$(\Omega,\mathcal{F},P)$ olasılık uzayı üzerinde \textbf{filtrasyon}
$\mathbb{F} = (\mathcal{F}_n)_{n\geq0}$; $\mathcal{F}_0 \subseteq \mathcal{F}_1 \subseteq \cdots \subseteq \mathcal{F}$
artan $\sigma$-cebirleri dizisi. $\mathcal{F}_n$, zaman $n$'deki bilgiyi temsil eder.

Dizi $(X_n)_{n\geq0}$ filtrasyona \textbf{uyumlu} (adapted) dır:
her $n$ için $X_n$ $\mathcal{F}_n$-ölçülebilir.
\end{tanim}

\begin{tanim}[Doğal Filtrasyon]
$X_0, X_1, \ldots$ dizisi için \textbf{doğal filtrasyon}:
$\mathcal{F}_n = \sigma(X_0, X_1, \ldots, X_n)$ — "ilk $n+1$ gözlemden üretilen bilgi".
\end{tanim}

% =====================================================================
\section{Martingal Tanımı ve Örnekler}
\label{sec:martingal_tanim}
% =====================================================================

\begin{tanim}[Martingal]
\label{def:martingal}
$(\mathbb{F},P)$-uyumlu ve $E[|X_n|]<\infty$ olan $(X_n)$ dizisi \textbf{martingal}dır, eğer
\[
  E[X_{n+1} \mid \mathcal{F}_n] = X_n \quad \text{n.k., her } n\geq0.
\]
\begin{itemize}
  \item \textbf{Süpmartingal (supermartingale):} $E[X_{n+1}\mid\mathcal{F}_n] \leq X_n$.
  \item \textbf{Altmartingal (submartingale):} $E[X_{n+1}\mid\mathcal{F}_n] \geq X_n$.
\end{itemize}
Sezgi: martingal $\equiv$ dürüst oyun; süpmartingal $\equiv$ oyuncuya karşı oyun.
\end{tanim}

\begin{ornek}[Temel Martingal Örnekleri]
\begin{enumerate}[(1)]
  \item \textbf{Kümülatif toplam:} $X_i$ bağımsız, $E[X_i]=0$; $S_n = \sum_{i=1}^n X_i$ martingal.
        $E[S_{n+1}\mid\mathcal{F}_n] = S_n + E[X_{n+1}\mid\mathcal{F}_n] = S_n + 0 = S_n$.

  \item \textbf{Ürün martingali:} $X_i$ bağımsız, $E[X_i]=1$; $M_n = \prod_{i=1}^n X_i$ martingal.

  \item \textbf{Koşullu beklenti:} $Y \in L^1$; $M_n = E[Y\mid\mathcal{F}_n]$ martingal
        (kule özelliği: $E[M_{n+1}\mid\mathcal{F}_n] = E[E[Y\mid\mathcal{F}_{n+1}]\mid\mathcal{F}_n]
        = E[Y\mid\mathcal{F}_n] = M_n$).

  \item \textbf{Doğrusal kazanç:} $M_n = S_n^2 - n\sigma^2$ ($X_i$ bağımsız, $E[X_i]=0$, $\mathrm{Var}(X_i)=\sigma^2$) martingal.
\end{enumerate}
\end{ornek}

\begin{onerme}[Jensen — Martingalden Altmartingale]
$M_n$ martingal ve $\varphi$ konveks ise $\varphi(M_n)$ altmartingal:
$E[\varphi(M_{n+1})\mid\mathcal{F}_n] \geq \varphi(E[M_{n+1}\mid\mathcal{F}_n]) = \varphi(M_n)$.
Özel: $|M_n|$ altmartingal ($\varphi = |\cdot|$).
\end{onerme}

% =====================================================================
\section{Durma Zamanları}
\label{sec:durma_zamani}
% =====================================================================

\begin{tanim}[Durma Zamanı]
\label{def:durma_zamani}
$\tau : \Omega \to \{0,1,2,\ldots\} \cup \{\infty\}$ rasgele değişkeni
\textbf{durma zamanı} (stopping time) dır, eğer her $n$ için
$\{\tau \leq n\} \in \mathcal{F}_n$ ise (durma kararı geleceği gerektirmez).
\end{tanim}

\begin{ornek}[Durma Zamanı Örnekleri]
\begin{itemize}
  \item İlk geçiş zamanı: $\tau_a = \inf\{n \geq 0 : S_n \geq a\}$ durma zamanıdır.
  \item Sabit $\tau = n_0$: her sabit zaman durma zamanıdır.
  \item "Son $n$ için $X_n > a$": bu gelecek bilgisi gerektirir — durma zamanı \emph{değildir}.
\end{itemize}
\end{ornek}

% =====================================================================
\section{İsteğe Bağlı Durma Teoremi}
\label{sec:ibd_teoremi}
% =====================================================================

\begin{teorem}[Doob'un İsteğe Bağlı Durma Teoremi]
\label{thm:ibd}
$(M_n)$ martingal, $\tau$ durma zamanı. Aşağıdaki koşullardan biri sağlanırsa
$E[M_\tau] = E[M_0]$:
\begin{enumerate}[(i)]
  \item $\tau$ sınırlı: $\tau \leq N$ sabit $N$ için.
  \item $\tau < \infty$ n.k.\ ve $|M_n|$ düzgün integrallenebilir.
  \item $E[\tau] < \infty$ ve $|M_{n+1}-M_n| \leq C$ sabit $C$.
\end{enumerate}
\end{teorem}

\begin{proof}[Sınırlı Durum İspat]
$\tau \leq N$. $M_\tau = M_N - \sum_{k=\tau}^{N-1}(M_{k+1}-M_k) \cdot \mathbf{1}_{\tau \leq k}$...
Alternatif: dürztaması takası.
$E[M_\tau] = \sum_{n=0}^N E[M_n \mathbf{1}_{\tau=n}] = \sum_{n=0}^N E[E[M_N\mid\mathcal{F}_n]\mathbf{1}_{\tau=n}]
= E[M_N] = E[M_0]$.
($\{\tau=n\} \in \mathcal{F}_n$; martingal özelliği ve kule.)
\end{proof}

\begin{ornek}[Gamblerın İflası]
Simetrik yürüyüş: $S_n = \sum_{i=1}^n \xi_i$ ($\xi_i = \pm1$ eşit olasılıkla).
$\tau = \inf\{n: S_n = a \text{ veya } S_n = -b\}$ ($a,b > 0$).
$M_n = S_n$ martingal; IBD: $E[S_\tau] = S_0 = 0$.
$E[S_\tau] = a \cdot P(S_\tau=a) - b \cdot P(S_\tau=-b) = 0$.
$P(S_\tau=a) = b/(a+b)$, $P(S_\tau=-b) = a/(a+b)$.
Başlangıç $S_0 = 0$: $a$ liraya ulaşma olasılığı $= b/(a+b)$.
\end{ornek}

% =====================================================================
\section{Doob'un Yakınsama Teoremi}
\label{sec:doob_yaklinsama}
% =====================================================================

\begin{tanim}[Yukarı Geçiş Sayısı]
$[a,b]$ aralığı için $(X_n)$'nin $n$ adımda $[a,b]$'yi \emph{aşağıdan yukarıya}
kaç kez geçtiği sayısına $U_n(a,b)$ denir (upcrossing number).
\end{tanim}

\begin{lemma}[Doob Yukarı Geçiş Lemması]
\label{lem:doob_upcrossing}
$(X_n)$ altmartingal ve $a < b$. O zaman:
$(b-a) E[U_n(a,b)] \leq E[(X_n-a)^+] - E[(X_0-a)^+]$.
\end{lemma}

\begin{teorem}[Doob Martingal Yakınsama Teoremi]
\label{thm:doob_yaklinsama}
$(X_n)$ altmartingal, $\sup_n E[X_n^+] < \infty$. O zaman
$X_\infty = \lim_{n\to\infty} X_n$ n.k.\ var ve $E[|X_\infty|] < \infty$.
\end{teorem}

\begin{proof}
Her $a < b$ için yukarı geçiş lemması:
$(b-a)E[U_\infty(a,b)] \leq \sup_n E[(X_n-a)^+] < \infty$.
$P(U_\infty(a,b) = \infty) = 0$ her rasyonel $a < b$ için.
$P(\liminf X_n < a < b < \limsup X_n) = 0$ tüm rasyonel çiftlerde.
Sayılabilir birleşim: $P(\liminf X_n < \limsup X_n) = 0$,
yani $X_n$ her yerde limit sahibi (sonsuz olabilir).
$\sup E[X_n^+] < \infty$ $\Rightarrow$ $E[X_\infty^+] < \infty$; monoton sınır $\Rightarrow$ sonlu.
\end{proof}

\begin{onerme}[GBSK — Martingal İspatı]
$X_i$ i.i.d., $E[|X_1|] < \infty$. $S_n = \sum X_i$.
$M_n = E[S_N/N \mid \mathcal{F}_n]$ (geri-martingal argümanıyla) $\to E[X_1]$ n.k.
Bu, GBSK'nın martingal teorisi aracılığıyla zarif alternatif ispatıdır.
\end{onerme}

% =====================================================================
\section{$L^2$ Martingalleri}
\label{sec:l2_martingal}
% =====================================================================

\begin{teorem}[$L^2$ Martingal Yakınsaması]
\label{thm:l2_martingal}
$(M_n)$ martingal, $\sup_n E[M_n^2] < \infty$. O zaman
$M_\infty$ mevcut, $E[M_\infty^2] < \infty$, ve $M_n \xrightarrow{L^2} M_\infty$.
\end{teorem}

\begin{proof}
$M_n^2$ altmartingal (Jensen); $\sup E[M_n^2] < \infty$ $\Rightarrow$ Doob yakınsama teoremi.
$L^2$ yakınsama: Doob'un eşitsizliği $E[\sup_n M_n^2] \leq 4\sup E[M_n^2]$;
DCT ile $L^2$ yakınsama.
\end{proof}

\begin{ornek}[Rassal Yürüyüş — $L^2$ Martingal]
$S_n = \sum_{i=1}^n \xi_i$ ($E[\xi_i]=0$, $E[\xi_i^2]=\sigma^2$).
$M_n = S_n/\sqrt{n}$: $E[M_n^2] = \sigma^2$ sınırlı.
Ama $M_n/\sqrt{\ln n} \to 0$ n.k.\ (YLK'nın sonucu),
$M_n$ ise $\Normal(0,\sigma^2)$'e dağılımla yakınsar (MLT).
$L^2$ sınırlılık n.k.\ yakınsamayı garanti etmez (sınır değişkenlik gösterebilir).
\end{ornek}

% ---- Öz-Yansıtma ----
\begin{ozyansima}
\begin{itemize}[leftmargin=1.8em]
  \item "Martingal = dürüst oyun" metaforunu iki farklı örnekle somutlaştırın.
  \item IBD teoremi neden durma zamanının sınırlı olmasını (veya başka bir koşul) gerektiriyor?
  \item GBSK'nın martingal ispatı neden daha "zarif"tir?
\end{itemize}
\end{ozyansima}

% ---- Köprü Notu ----
\begin{kopru}
\textbf{Önceki:} Bölüm~\ref{sec:kosullu_beklenti} koşullu beklenti; Böl.~\ref{chp:buyuk_sayilar} BSK.\\
\textbf{Sonraki (opsiyonel):} Bölüm~\ref{chp:stokastik} — Stokastik Süreçler.\\
\textbf{Uygulamalar:} Finansta martingal fiyatlama; sıralı analizde isteğe bağlı durma.
\defterbaglantisi{bolum09\_martingaller.ipynb}{Martingal simülasyonu, gamblerın iflası}
\end{kopru}

% =====================================================================
%  ALIŞTIRMALAR
% =====================================================================

\cozumlualistirmalar

\begin{alistirma}[Al.9.1]{A}{Martingal Doğrulama}
$X_i$ bağımsız, $E[X_i]=1$, $M_n = \prod_{i=1}^n X_i$.
$E[M_{n+1}\mid\mathcal{F}_n] = M_n$ olduğunu gösterin.
\end{alistirma}
\alcozum{%
$E[M_{n+1}\mid\mathcal{F}_n] = E[M_n X_{n+1}\mid\mathcal{F}_n] = M_n E[X_{n+1}\mid\mathcal{F}_n]
= M_n E[X_{n+1}] = M_n \cdot 1 = M_n$.
(Bağımsızlık: $X_{n+1}$ $\mathcal{F}_n$'den bağımsız.)}

\begin{alistirma}[Al.9.2]{B}{Gamblerın İflası — Varyans Martingali}
$S_n$ simetrik rassal yürüyüş ($\xi_i = \pm1$, $\sigma^2=1$). $M_n = S_n^2 - n$.
(a) $M_n$'nin martingal olduğunu gösterin.
(b) $\tau = \inf\{n: S_n = a\}$ ($a > 0$) için $E[\tau] = a^2$ olduğunu IBD ile bulun.
\end{alistirma}
\alcozum{%
(a) $E[M_{n+1}\mid\mathcal{F}_n] = E[(S_n+\xi_{n+1})^2 - (n+1)\mid\mathcal{F}_n]
= S_n^2 + 2S_n E[\xi_{n+1}] + E[\xi_{n+1}^2] - n - 1 = S_n^2 - n = M_n$. $\checkmark$\\
(b) IBD ($\tau$ sınırsız ama Koşul (iii) çalışır — artışlar $\leq 1+1+1$):
$E[M_\tau] = E[S_\tau^2 - \tau] = E[M_0] = 0$.
$S_\tau = a$ (emici bariyer): $a^2 - E[\tau] = 0$. $E[\tau] = a^2$.}

% ---

\alistirmalar

\begin{alistirma}[Al.9.3]{A}{Altmartingal}
$(X_n)$ martingal. $\varphi(x) = x^2$ için $\varphi(X_n)$'nin altmartingal
olduğunu Jensen'i kullanarak gösterin.
\end{alistirma}
\alcozum{%
$E[X_{n+1}^2\mid\mathcal{F}_n] \geq (E[X_{n+1}\mid\mathcal{F}_n])^2 = X_n^2$. (Jensen, $\varphi$ konveks.)}

\begin{alistirma}[Al.9.4]{B}{Doob Yakınsama Uygulaması}
$X_i \geq 0$ bağımsız, $E[X_i] = 1$. $M_n = \prod_{i=1}^n X_i$.
$\sup_n E[M_n] = 1 < \infty$ ise $M_n$ n.k.\ yakınsar.
$M_\infty = 0$ n.k.\ olduğunu BSK ile gösterin.
\end{alistirma}
\alcozum{%
$\ln M_n = \sum \ln X_i$. BSK: $(\ln M_n)/n \to E[\ln X_1]$.
Jensen: $E[\ln X_1] \leq \ln E[X_1] = 0$; genellikle $< 0$.
$\ln M_n / n \to c < 0$ $\Rightarrow$ $M_n = e^{n \cdot c+o(n)} \to 0$ n.k.}

% =====================================================================
\endinput

% %% ============================================================
%% BÖLÜM 10 — Stokastik Süreçlere Giriş  [OPSİYONEL]
%% Olasılık Teorisi ve Matematiksel İstatistik
%% ============================================================
\chapter{Stokastik Süreçlere Giriş}
\label{chp:stokastik}

\bolumalinti{%
  The theory of stochastic processes is really a theory of families of random variables
  indexed by a parameter set.
}{Joseph Doob}

\begin{kazanimlar}
Bu bölümün sonunda okuyucu:
\begin{itemize}
  \item Stokastik süreç, örneklem yolu, durumlar uzayı kavramlarını tanımlar.
  \item Markov zincirlerini tanımlayıcı özellikleriyle açıklar; geçiş matrislerini ve Chapman-Kolmogorov denklemini uygular.
  \item Süreç sınıflarını (geri dönüşlü, geçici, ergodik, dönemsel) belirler ve durağan dağılımı hesaplar.
  \item Poisson sürecini aksiyomlardan türetir; süreçler arası ilişkileri tanımlar.
  \item Brownian hareketi tanımlar ve temel özelliklerini (bağımsız artışlar, süreklilik, Gaussian dağılım) listeler.
\end{itemize}
\end{kazanimlar}

\begin{motivasyon}
Tek bir rasgele değişken statik bir belirsizliği modellerken, \emph{stokastik süreç} zamanla evrilen belirsizliği modelleyen rasgele değişken ailesidir. Hisse senedi fiyatları, kuyruklar, epidemi yayılımı, iklim değişkenliği — bunların tümü stokastik süreçtir.

\medskip
\noindent\textbf{Köprü:} Bölüm~\ref{chp:martingaller}'daki filtrasyon ve uyum kavramları doğrudan burada kullanılacak. Markov zincirlerindeki geçişler aslında $\sigma$-cebirlerinin monoton dizisi anlamında "bilgi artışı" kavramıdır.
\end{motivasyon}

%% ============================================================
\section{Temel Tanımlar}
\label{sec:stokastik_temel}
%% ============================================================

\begin{tanim}[Stokastik süreç]\label{def:stokastik_surec}
$(\Omega,\mathcal{F},\Prob)$ olasılık uzayı ve $T$ indeks kümesi olsun. \textbf{Stokastik süreç} \emph{(stochastic process)}, $\{X_t : t\in T\}$ rasgele değişken ailesidir. $T\subseteq\mathbb{Z}$ ise \textbf{ayrık zamanlı}, $T\subseteq[0,\infty)$ ise \textbf{sürekli zamanlı} denir. $S = X_t(\Omega)$ kümesine \textbf{durum uzayı} \emph{(state space)} denir.
\end{tanim}

\begin{tanim}[Örneklem yolu]\label{def:orneklem_yolu}
Sabit $\omega\in\Omega$ için $t\mapsto X_t(\omega)$ fonksiyonuna \textbf{örneklem yolu} \emph{(sample path)} veya \textbf{yörünge} \emph{(trajectory)} denir.
\end{tanim}

\begin{ornek}[Basit rasgele yürüyüş]\label{ex:basit_yuruyus}
$\xi_1,\xi_2,\ldots\iid$ ile $\Prob(\xi_i=+1)=p$, $\Prob(\xi_i=-1)=1-p$. $S_0=0$, $S_n=\sum_{k=1}^n\xi_k$. $\{S_n:n\geq0\}$ tam sayı değerli ayrık zamanlı stokastik süreçtir.
\end{ornek}

%% ============================================================
\section{Markov Zincirleri}
\label{sec:markov_zincirleri}
%% ============================================================

\begin{kesif}
Basit rasgele yürüyüşte $S_{n+1}$'in dağılımı, tüm geçmiş $S_0,S_1,\ldots,S_n$ yerine yalnızca $S_n$'e bağlıdır. Bu "geçmişten bağımsızlık" özelliği Markov zincirlerin özüdür.
\end{kesif}

\begin{tanim}[Markov zinciri]\label{def:markov_zinciri}
Ayrık durum uzaylı ($S$ sayılabilir) ayrık zamanlı süreç $\{X_n\}$, aşağıdaki \textbf{Markov özelliğini} \emph{(Markov property)} sağlıyorsa Markov zinciri adını alır:
\[
  \Prob(X_{n+1}=j \mid X_0=i_0,\ldots,X_{n-1}=i_{n-1},X_n=i)
  = \Prob(X_{n+1}=j\mid X_n=i).
\]
\end{tanim}

\begin{tanim}[Geçiş olasılıkları]\label{def:gecis_olasiliklari}
\textbf{Tek adım geçiş olasılığı}: $p_{ij} = \Prob(X_{n+1}=j\mid X_n=i)$.
\textbf{Geçiş matrisi}: $P = (p_{ij})_{i,j\in S}$; $p_{ij}\geq0$, $\sum_j p_{ij}=1$ (stokastik matris).
\end{tanim}

\begin{teorem}[Chapman-Kolmogorov]\label{thm:chapman_kolmogorov}
$p_{ij}^{(m+n)} = \sum_{k\in S} p_{ik}^{(m)}p_{kj}^{(n)}$, yani $P^{(m+n)} = P^{(m)}P^{(n)} = P^{m+n}$.
\end{teorem}

\begin{proof}
Koşullu olasılığı $X_m$ üzerinden toplayalım:
\[
  p_{ij}^{(m+n)} = \Prob(X_{m+n}=j\mid X_0=i)
  = \sum_k \Prob(X_{m+n}=j\mid X_m=k)\Prob(X_m=k\mid X_0=i)
  = \sum_k p_{kj}^{(n)}p_{ik}^{(m)}.
\]
\end{proof}

\subsection{Sınıf Yapısı}

\begin{tanim}[Erişilebilirlik ve sınıf]\label{def:erisim}
$j$ durumu $i$'den \textbf{erişilebilir} \emph{(accessible)}: $p_{ij}^{(n)}>0$ olan $n\geq0$ varsa; $i\leftrightarrow j$ ise ikisi \textbf{iletişimli} \emph{(communicating)}. İletişimli durumlar kümesi \textbf{sınıf} \emph{(class)} oluşturur. Her durum bir sınıfta. Tüm durumlar tek sınıf ise zincir \textbf{geri dönülemez} \emph{(irreducible)}.
\end{tanim}

\begin{tanim}[Geri dönüşlü ve geçici]\label{def:geri_donusluluk}
Durum $i$ için $f_i = \Prob(\exists n\geq1: X_n=i\mid X_0=i)$ geri dönüş olasılığı.
\begin{itemize}
  \item $f_i=1$: \textbf{geri dönüşlü} \emph{(recurrent)}.
  \item $f_i<1$: \textbf{geçici} \emph{(transient)}.
\end{itemize}
\end{tanim}

\begin{teorem}[Geri dönüşlülük kriteri]\label{thm:geri_donus_krit}
$\sum_{n=0}^\infty p_{ii}^{(n)} = \infty \iff i$ geri dönüşlü.
\end{teorem}

\begin{proof}
$R = \sum_{n=1}^\infty \mathbf{1}[X_n=i]$ geri dönüş sayısı. $\Beklenti_i[R]=\sum_{n=1}^\infty p_{ii}^{(n)}$. Geometrik dağılım argümanıyla $\Beklenti_i[R]=f_i/(1-f_i)$: $f_i=1$ iff $\Beklenti_i[R]=\infty$ iff seri ıraksar.
\end{proof}

\begin{tanim}[Periyot]\label{def:periyot}
$d(i) = \gcd\{n\geq1: p_{ii}^{(n)}>0\}$ durum $i$'nin \textbf{periyodu} \emph{(period)}. $d(i)=1$ ise \textbf{aperiyodik} \emph{(aperiodic)}.
\end{tanim}

\subsection{Durağan Dağılım}

\begin{tanim}[Durağan dağılım]\label{def:duraganlik}
Olasılık vektörü $\pi = (\pi_j)_{j\in S}$ ($\pi_j\geq0$, $\sum\pi_j=1$), $\pi P=\pi$ koşulunu sağlıyorsa \textbf{durağan dağılım} \emph{(stationary distribution)} denir.
\end{tanim}

\begin{teorem}[Ergodik teorem]\label{thm:ergodik}
$\{X_n\}$ geri dönülemez, geri dönüşlü, aperiyodik Markov zinciri olsun. O zaman:
\begin{enumerate}[(a)]
  \item Yegâne durağan dağılım $\pi$ vardır: $\pi_j = 1/m_j$ ($m_j$ = ortalama geri dönüş süresi).
  \item Her $i,j$: $p_{ij}^{(n)}\to\pi_j$ ($n\to\infty$).
  \item Her $f:S\to\mathbb{R}$ sınırlı için zaman ortalaması olasılık ortalamasına yakınsar:
  \[
    \frac{1}{n}\sum_{k=0}^{n-1}f(X_k) \asto \sum_j f(j)\pi_j.
  \]
\end{enumerate}
\end{teorem}

\begin{ispatiskele}
Geri dönüş teoremi, yenileme teorisine dayanır. $m_j=\Beklenti_j[T_j]$ (ilk geri dönüş zamanı). Ergodik limit, büyük sayılar kanunu analog. Tam ispat: Norris (1997) §1.8.
\end{ispatiskele}

\begin{ornek}[İki durumlu zincir]\label{ex:iki_durumlu}
$S=\{0,1\}$, $P=\begin{pmatrix}1-\alpha&\alpha\\\beta&1-\beta\end{pmatrix}$. Durağan dağılım: $\pi_0=\beta/(\alpha+\beta)$, $\pi_1=\alpha/(\alpha+\beta)$. $\alpha,\beta>0$ ise ergodik.
\end{ornek}

\begin{ornek}[Gambler's ruin — Markov zinciri]\label{ex:gambler_markov}
$S=\{0,1,\ldots,N\}$; $p_{i,i+1}=p$, $p_{i,i-1}=q=1-p$, sınır durumları absorbe edici. Bu Markov zincirinde $0$ ve $N$ geri dönüşlü, diğerleri geçici. $i$'den başlayıp $N$'ye ulaşma olasılığı $\rho_i$:
\[
  \rho_i = \frac{1-(q/p)^i}{1-(q/p)^N} \quad (p\neq q), \qquad \rho_i = i/N \quad (p=q).
\]
\end{ornek}

%% ============================================================
\section{Poisson Süreci}
\label{sec:poisson_sureci}
%% ============================================================

\begin{tanim}[Sayma süreci]\label{def:sayma_sureci}
$N(t)$ ($t\geq0$) tam sayı değerli süreç; $N(0)=0$ ve $t\mapsto N(t)$ n.k.\ sağdan sürekli, artan ise \textbf{sayma süreci} \emph{(counting process)}.
\end{tanim}

\begin{tanim}[Poisson süreci]\label{def:poisson_sureci}
Sayma süreci $\{N(t)\}$, aşağıdakileri sağlıyorsa $\lambda>0$ hızlı \textbf{Poisson süreci} \emph{(Poisson process)} adını alır:
\begin{enumerate}[(i)]
  \item $N(0)=0$.
  \item \textbf{Bağımsız artışlar}: $0\leq s<t$, $N(t)-N(s)$ önceki artışlardan bağımsız.
  \item \textbf{Durağan artışlar}: $N(t+s)-N(s) \overset{d}{=} N(t)$.
  \item $\Prob(N(h)=1)=\lambda h+o(h)$, $\Prob(N(h)\geq2)=o(h)$.
\end{enumerate}
\end{tanim}

\begin{teorem}[Poisson dağılımı]\label{thm:poisson_dagilim}
$\{N(t)\}$ Poisson süreci ise $N(t)\sim\Pois(\lambda t)$.
\end{teorem}

\begin{proof}
$p_n(t)=\Prob(N(t)=n)$ olsun. İleriye denklem: bağımsız artışlar ve aksiyom (iv) ile
\[
  p_n(t+h) = p_n(t)(1-\lambda h) + p_{n-1}(t)\lambda h + o(h).
\]
$h\to0$: $p_n'(t) = -\lambda p_n(t)+\lambda p_{n-1}(t)$, $p_0(0)=1$, $p_n(0)=0$ ($n\geq1$). Çözüm: $p_n(t)=e^{-\lambda t}(\lambda t)^n/n!$.
\end{proof}

\begin{teorem}[Varış araları]\label{thm:varis_aralari}
$T_k$ ($k$-inci varış arası) i.i.d.\ $\Exp(\lambda)$.
\end{teorem}

\begin{proof}
$\Prob(T_1>t) = \Prob(N(t)=0)=e^{-\lambda t}$; dolayısıyla $T_1\sim\Exp(\lambda)$. Markov özelliği (bağımsız artışlar) $T_2,T_3,\ldots$'yı da bağımsız $\Exp(\lambda)$ yapar.
\end{proof}

\begin{teorem}[İnceleme — süperpozisyon]\label{thm:poisson_ozellikler}
$\{N_1(t)\}$ ve $\{N_2(t)\}$ bağımsız, sırasıyla $\lambda_1$ ve $\lambda_2$ hızlı Poisson süreçleri.
\begin{enumerate}[(a)]
  \item \textbf{Süperpozisyon:} $N(t)=N_1(t)+N_2(t)$ ise $(\lambda_1+\lambda_2)$ hızlı Poisson.
  \item \textbf{İnceleme:} Her olayı $p$ olasılıkla tut. Tutulan $\{M(t)\}$ ise $p\lambda$ hızlı Poisson.
  \item \textbf{Koşullu:} $N(t)=n$ verildiğinde $n$ varış zamanı $[0,t]$ üzerinde i.i.d.\ $\Uniform[0,t]$.
\end{enumerate}
\end{teorem}

\begin{proof}[Proof (a)]
Karakteristik fonksiyonla: $N_1(t)+N_2(t)$ bağımsız Poisson toplamı $\Pois((\lambda_1+\lambda_2)t)$.
\end{proof}

%% ============================================================
\section{Brownian Hareket}
\label{sec:brownian}
%% ============================================================

\begin{tarihselnot}
Robert Brown (1827) su altındaki polenlerin rastgele hareketini gözlemledi. Einstein (1905) diffüzyon denklemini türetti. Wiener (1923) matematiksel temeli kurdu — bu nedenle \emph{Wiener süreci} de denir.
\end{tarihselnot}

\begin{tanim}[Standart Brownian hareket]\label{def:brownian}
$\{W(t): t\geq0\}$ süreci aşağıdakileri sağlıyorsa \textbf{standart Brownian hareket} \emph{(standard Brownian motion / Wiener process)} adını alır:
\begin{enumerate}[(i)]
  \item $W(0)=0$ n.k.
  \item \textbf{Bağımsız artışlar:} $0\leq s<t$, $W(t)-W(s)$ öncesinden bağımsız.
  \item \textbf{Gaussian artışlar:} $W(t)-W(s)\sim\Normal(0,t-s)$.
  \item \textbf{Sürekli yollar:} $t\mapsto W(t)$ n.k.\ sürekli.
\end{enumerate}
\end{tanim}

\begin{teorem}[Brownian hareketin temel özellikleri]\label{thm:brownian_ozellik}
$\{W(t)\}$ standart Brownian hareket. O zaman:
\begin{enumerate}[(a)]
  \item $W(t)\sim\Normal(0,t)$.
  \item $\Cov(W(s),W(t))=\min(s,t)$.
  \item Yollar n.k.\ hiçbir noktada türevlenemez.
  \item $W$ bir martingaldir.
\end{enumerate}
\end{teorem}

\begin{proof}
(a) $W(t)=W(t)-W(0)\sim\Normal(0,t)$ tanımdan.

(b) $s\leq t$: $\Cov(W(s),W(t))=\Cov(W(s),W(s)+(W(t)-W(s)))=\Var(W(s))=s$ (bağımsız artışlar).

(c) Yol türevsizliği: Paley-Wiener-Zygmund teoremi (ileri konu, referans: Mörters-Peres).

(d) $\Beklenti[W(t)\mid\mathcal{F}_s]=W(s)$: $W(t)-W(s)\perp\mathcal{F}_s$, $\Beklenti[W(t)-W(s)]=0$.
\end{proof}

\begin{ornek}[Geometrik Brownian hareket]\label{ex:geometrik_brownian}
$S(t) = S_0\exp\!\bigl[(\mu-\tfrac{1}{2}\sigma^2)t + \sigma W(t)\bigr]$ formülü Black-Scholes modelinde hisse senedi fiyatını verir. $\log S(t)\sim\Normal(\log S_0+(\mu-\sigma^2/2)t,\,\sigma^2 t)$.
\end{ornek}

\begin{disiplinlerarasibaglan}
\textbf{Fizik:} Brownian hareket ısı difüzyonunu modelleyen $\partial u/\partial t = \frac{1}{2}\partial^2 u/\partial x^2$ denkleminin olasılıksal çözümüdür.
\textbf{Finans:} Black-Scholes opsiyon fiyatlaması geometrik Brownian harekete dayanır.
\textbf{Biyoloji:} Protein difüzyonu, nöral gürültü modelleri.
\end{disiplinlerarasibaglan}

%% ============================================================

\begin{mikroozet}
\begin{itemize}
  \item Markov zinciri: gelecek yalnızca şimdiye bağlı; Chapman-Kolmogorov $P^{m+n}=P^m P^n$.
  \item Durağan dağılım: $\pi P=\pi$; ergodik zincirde yegâne ve $p_{ij}^{(n)}\to\pi_j$.
  \item Poisson süreci: bağımsız+durağan artışlar; $N(t)\sim\Pois(\lambda t)$; varış araları $\Exp(\lambda)$.
  \item Brownian hareket: Gaussian bağımsız artışlar, sürekli türevsiz yollar, martingal.
\end{itemize}
\end{mikroozet}

\begin{ozyansima}
Markov özelliği "geçmişi unutan" sistemleri ne kadar iyi modelliyor? Hangisi daha zengin: Markov zincirleri mi yoksa martingaller mi? Bu ikisinin kesişimi nedir?
\end{ozyansima}

\begin{kopru}
\textbf{Nereden geldik:} Bölüm~\ref{chp:martingaller}'daki Doob yakınsama teoremi, geri dönüşlü Markov zincirlerin ergodik teoreminin özel durumudur.

\textbf{Nereye gidiyoruz:} Part II'de (Böl.~\ref{chp:guven_araliklari} vd.) bu süreçlerin parametrelerini veriden tahmin etmeye odaklanacağız. Markov zinciri Monte Carlo (MCMC) Bayesçi istatistikte doğrudan kullanılır.
\end{kopru}

%% ============================================================
%% ALISTIRMALAR
%% ============================================================

\cozumlualistirmalar

\begin{alistirma}[Al.10.1]{A}{Durağan dağılım}
Üç durumlu Markov zinciri $P=\begin{pmatrix}0.5&0.4&0.1\\0.3&0.5&0.2\\0.2&0.3&0.5\end{pmatrix}$. Durağan dağılımı bul.
\end{alistirma}
\alcozum{%
$\pi P=\pi$, $\sum\pi_i=1$. Denklem sistemi: $0.5\pi_0+0.3\pi_1+0.2\pi_2=\pi_0$ vb. Çözüm: $\pi=(0.303,0.404,0.293)$ (yaklaşık; tam çözüm için lineer denklem sistemi).%
}

\begin{alistirma}[Al.10.2]{B}{Poisson süperpozisyon}
Bir hastane acil birimine birbirinden bağımsız iki kaynaktan hasta geliyor: 112 araçlarından $\lambda_1=3$/saat, bireysel başvurulardan $\lambda_2=2$/saat. Bir saatte 10 veya daha fazla hasta gelme olasılığı nedir?
\end{alistirma}
\alcozum{%
Toplam süreç $\Pois(5)$. $\Prob(N(1)\geq10)=1-\Prob(N(1)\leq9)=1-\sum_{k=0}^9 e^{-5}5^k/k!\approx1-0.9682=0.0318$.%
}

\alistirmalar

\begin{alistirma}[Al.10.3]{A}{Brownian varyans}
$W(t)$ standart Brownian hareket. $\Beklenti[W(3)^2]$, $\Cov(W(2),W(5))$ ve $\Cor(W(2),W(5))$'i hesapla.
\end{alistirma}
\alcozum{%
$\Beklenti[W(3)^2]=\Var(W(3))=3$. $\Cov(W(2),W(5))=\min(2,5)=2$. $\Cor=2/\sqrt{2\cdot5}=2/\sqrt{10}$.%
}

\begin{alistirma}[Al.10.4]{B}{İki durumlu ergodiklik}
İki durumlu Markov zinciri $P=\begin{pmatrix}1-\alpha&\alpha\\\beta&1-\beta\end{pmatrix}$ ($\alpha,\beta\in(0,1)$). (a) Durağan dağılımı bul. (b) $p_{00}^{(n)}$'in kapalı formülünü türet. (c) $n\to\infty$ limitinin durağan dağılıma yakınsadığını doğrula.
\end{alistirma}
\alcozum{%
(a) $\pi_0=\beta/(\alpha+\beta)$, $\pi_1=\alpha/(\alpha+\beta)$. (b) Özdeğerler $1$ ve $1-\alpha-\beta$: $p_{00}^{(n)}=\pi_0+\pi_1(1-\alpha-\beta)^n$. (c) $|1-\alpha-\beta|<1$ olduğundan ikinci terim $\to0$; $p_{00}^{(n)}\to\pi_0$.%
}

\begin{alistirma}[Al.10.5]{C}{Gambler's ruin — genel}
$p\neq q$ için Gambler's Ruin'de $i$'den başlayıp $N$'ye ulaşma olasılığı $\rho_i=\frac{1-(q/p)^i}{1-(q/p)^N}$ olduğunu fark denklemiyle ispat et. Ortalama oyun uzunluğunu $i=N/2$, $p=q=1/2$ özel durumunda hesapla.
\end{alistirma}
\alcozum{%
Fark denklemi: $\rho_i=p\rho_{i+1}+q\rho_{i-1}$, $\rho_0=0$, $\rho_N=1$. Karakteristik denklem $pr^2-r+q=0$; kökler $1$ ve $q/p$. Genel çözüm $\rho_i=A+B(q/p)^i$; sınır koşullarıyla $\rho_i$ bulunur. $p=q=1/2$, $i=N/2$: ortalama süre $= i(N-i) = N^2/4$.%
}


% ---- PART II: Matematiksel İstatistik ----
\part{Matematiksel İstatistik}

% =====================================================================
%  BÖLÜM 11: İstatistiğin Temelleri — Model ve Örneklem
%  Kaynak: Casella & Berger; Hogg & Craig; Lehmann & Casella
% =====================================================================
\chapter{İstatistiğin Temelleri: Model ve Örneklem}
\label{chp:istatistik_temelleri}

\bolumalinti{%
  İstatistik, gözlemlerin ötesinde olanı anlamaya çalışır;
  bu anlamda her istatistiksel çıkarım bir inançtır.%
}{Erich Leo Lehmann, \textit{Theory of Point Estimation}}

% ---- Kazanımlar (TYMM) ----
\begin{kazanimlar}
Bu bölümün sonunda öğrenci:
\begin{itemize}[leftmargin=1.8em]
  \item İstatistiksel model kavramını parametrik ve parametrik olmayan çerçevede tanımlayabilir;
        istatistiksel karar, kayıp fonksiyonu ve risk kavramlarını açıklayabilir.
  \item Örneklem, istatistik ve örneklem dağılımı kavramlarını tanımlayabilir;
        $\bar{X}_n$ ve $S^2$'nin dağılımlarını türetebilir.
  \item Yeterli istatistiği tanımlayabilir; Fisher-Neyman çarpanlaştırma teoremini
        ispatlayabilir ve uygulamalarda yeterli istatistik bulabilir.
  \item Tamamlayıcı ve en az yeterli istatistik kavramlarını tanımlayabilir;
        bunları UMVUE teorisindeki rolleriyle ilişkilendirebilir.
\end{itemize}
\end{kazanimlar}

% ---- Motivasyon ----
\begin{motivasyon}
Part~I'de olasılığın matematiksel temellerini inşa ettik: ölçüm teorisi,
rasgele değişkenler, beklenti, dağılımlar, büyük sayılar ve limit teoremler.
Şimdi bu araçları ters yönde kullanacağız: \emph{gözlemlerden} hareketle
bilinmeyen parametreler hakkında çıkarım yapmak.

İstatistik dili şöyle çalışır: gerçeklik, bilinmeyen parametre $\theta$ tarafından
belirlenen bir dağılım $P_\theta$'dur. Biz yalnızca bu dağılımdan çekilen
$n$ gözlemi görürüz ve $\theta$'yı tahmin etmeye ya da test etmeye çalışırız.

\medskip
\textbf{Köprü.} Bölüm~\ref{chp:dagilim_aileleri}'deki üstel aile,
Bölüm~\ref{chp:beklenti}'deki koşullu beklenti ve
Bölüm~\ref{chp:merkezi_limit}'teki asimptotik teori bu bölümün doğrudan ön koşullarıdır.
\end{motivasyon}

% =====================================================================
\section{İstatistiksel Model Kavramı}
\label{sec:istatistik_model}
% =====================================================================

\begin{tanim}[İstatistiksel Model]
\label{def:istatistik_model}
\textbf{İstatistiksel model} bir aile $\mathcal{P} = \{P_\theta : \theta \in \Theta\}$'dır;
burada her $P_\theta$ $(\mathcal{X}, \mathcal{B})$ üzerinde bir olasılık ölçüsü,
$\Theta \subseteq \mathbb{R}^k$ parametre uzayıdır.

\begin{itemize}
  \item \textbf{Parametrik model:} $\Theta$ sonlu boyutlu ($\Theta \subseteq \mathbb{R}^k$, $k$ sabit).
        Örnek: $\Normal(\mu, \sigma^2)$ ailesi, $\Theta = \mathbb{R}\times\mathbb{R}_{>0}$.
  \item \textbf{Parametrik olmayan model:} $\Theta$ sonsuz boyutlu.
        Örnek: "tüm sürekli dağılımlar" ailesi.
  \item \textbf{Yarı-parametrik model:} Parametrik bileşen + parametrik olmayan bileşen.
        Örnek: Cox orantılı tehlike modeli.
\end{itemize}
\end{tanim}

\begin{tanim}[Tanımlanabilirlik]
\label{def:tanimlanabilirlik}
Model $\{P_\theta\}$ \textbf{tanımlanabilir} (identifiable) dır, eğer
$\theta \neq \theta'$ $\Rightarrow$ $P_\theta \neq P_{\theta'}$ ise.
Tanımlanabilirlik olmadan $\theta$'yı veri ile ayırt etmek imkânsızdır.
\end{tanim}

\begin{ornek}[Tanımlanabilir ve Tanımlanamaz Modeller]
\textbf{Tanımlanabilir:} $\Normal(\mu, \sigma^2)$: farklı $(\mu,\sigma^2)$ çiftleri
farklı dağılımlar üretir.\\
\textbf{Tanımlanamaz (örnek):} $\Normal(\mu_1+\mu_2, \sigma^2)$:
$\mu_1=1, \mu_2=2$ ile $\mu_1=0, \mu_2=3$ aynı dağılımı verir.
Bu model $\mu_1$ ve $\mu_2$'yi ayrı ayrı tanımlayamaz.
\end{ornek}

% =====================================================================
\section{İstatistiksel Kararlar ve Kayıp Fonksiyonu}
\label{sec:karar_teorisi}
% =====================================================================

\begin{tanim}[Karar Çerçevesi]
\label{def:karar}
\textbf{Karar teorisi} çerçevesi:
\begin{itemize}
  \item $\Theta$: parametre uzayı (gerçeklik kümesi),
  \item $\mathcal{D}$: karar uzayı (verilen yanıtlar kümesi),
  \item $L(\theta, d)$: \textbf{kayıp fonksiyonu} (loss function) —
        gerçek parametre $\theta$ iken $d$ kararı almanın maliyeti,
  \item $\delta(\mathbf{X})$: \textbf{karar kuralı} — gözlemden karara giden fonksiyon.
\end{itemize}
\textbf{Risk fonksiyonu:} $R(\theta, \delta) = E_\theta[L(\theta, \delta(\mathbf{X}))]$.
\end{tanim}

\begin{ornek}[Yaygın Kayıp Fonksiyonları]
\begin{itemize}
  \item \textbf{Karesel kayıp (squared error loss):} $L(\theta, d) = (\theta - d)^2$.
        Risk $= E[(\theta - \delta(\mathbf{X}))^2]$ = ortalama karesel hata (OKH, MSE).
  \item \textbf{Mutlak değer kaybı:} $L(\theta, d) = |\theta - d|$.
        Optimal karar: koşullu medyan.
  \item \textbf{0-1 kaybı (sınıflandırma):} $L(\theta, d) = \mathbf{1}_{d\neq\theta}$.
        Risk = hata olasılığı.
\end{itemize}
\end{ornek}

\begin{tanim}[Yansızlık ve OKH Ayrışımı]
\label{def:yansizlik_okh}
$\delta(\mathbf{X})$, $\theta$'nın \textbf{yansız tahmin edicisi} (unbiased estimator):
$E_\theta[\delta(\mathbf{X})] = \theta$ her $\theta \in \Theta$.
Karesel kayıp altında:
\[
  \mathrm{OKH}(\delta, \theta) = \mathrm{Var}_\theta(\delta) + [\mathrm{Yan}(\delta, \theta)]^2,
\]
burada $\mathrm{Yan}(\delta,\theta) = E_\theta[\delta] - \theta$. Yansız: $\mathrm{OKH} = \mathrm{Var}$.
\end{tanim}

\begin{proof}[OKH Ayrışımı]
$E[(\delta-\theta)^2] = E[(\delta - E[\delta] + E[\delta]-\theta)^2]
= \mathrm{Var}(\delta) + (E[\delta]-\theta)^2$ (çapraz terim sıfır).
\end{proof}

% =====================================================================
\section{Örneklem ve İstatistik Kavramı}
\label{sec:orneklem}
% =====================================================================

\begin{tanim}[Rasgele Örneklem]
\label{def:rasgele_orneklem}
$X_1, \ldots, X_n$ i.i.d.\ $P_\theta$ dağılımlı rasgele değişkenler;
bu dizi $\theta$'lı dağılımdan alınan \textbf{rasgele örneklem} (random sample) dir.
Veri $\mathbf{x} = (x_1,\ldots,x_n)$ bu dizinin bir gerçekleşmesidir.
\end{tanim}

\begin{tanim}[İstatistik]
\label{def:istatistik}
$T = T(X_1,\ldots,X_n) : \mathcal{X}^n \to \mathbb{R}^k$ ölçülebilir bir fonksiyon;
$\theta$'ya \emph{doğrudan} bağlı olmayan (parametreyi bilmeden hesaplanabilen) böyle
bir fonksiyona \textbf{istatistik} (statistic) denir.

Temel istatistikler:
\begin{align*}
  \bar{X}_n &= \frac{1}{n}\sum_{i=1}^n X_i \quad \text{(örneklem ortalaması)}, \\
  S^2 &= \frac{1}{n-1}\sum_{i=1}^n (X_i - \bar{X}_n)^2 \quad \text{(örneklem varyansı)}, \\
  M_k &= \frac{1}{n}\sum_{i=1}^n X_i^k \quad \text{($k$'ıncı örneklem momenti)}.
\end{align*}
\end{tanim}

% =====================================================================
\section{Örneklem Dağılımları}
\label{sec:orneklem_dagilim}
% =====================================================================

\subsection{Örneklem Ortalaması ve Varyansının Dağılımı}

\begin{teorem}[Normal Popülasyondan Örneklem]
\label{thm:normal_orneklem}
$X_1,\ldots,X_n$ i.i.d.\ $\Normal(\mu,\sigma^2)$. O zaman:
\begin{enumerate}[(i)]
  \item $\bar{X}_n \sim \Normal(\mu,\, \sigma^2/n)$.
  \item $\frac{(n-1)S^2}{\sigma^2} \sim \chi^2_{n-1}$.
  \item $\bar{X}_n$ ve $S^2$ \textbf{bağımsızdır}.
  \item $T = \frac{\bar{X}_n - \mu}{S/\sqrt{n}} \sim t_{n-1}$.
\end{enumerate}
\end{teorem}

\begin{proof}
\textit{(i):} $\bar{X}_n = n^{-1}\sum X_i$; bağımsız normalların doğrusal bileşimi
Bölüm~\ref{thm:normal_yeniden}'den $\Normal(\mu, \sigma^2/n)$.

\textit{(ii) ve (iii):} Ortogonal dönüşüm argümanı.
$(X_1,\ldots,X_n)$ birim vektörlerle ifade edilirse, $\bar{X}_n$
$\mathbf{1}/\sqrt{n}$ yönüne, $\sum(X_i-\bar{X}_n)^2$ ise buna dik olan $n-1$ boyutlu
alt uzaya karşılık gelir. Ortogonal dönüşüm normalliği korur ve bağımsızlığı üretir.
Dik bileşen: $n-1$ bağımsız standart normal karesinin toplamı $= \chi^2_{n-1}$.

Alternatif ispat: $Y_1 = \sqrt{n}\bar{X}$, $Y_i = X_{i-1}-X_i$ ($i=2,\ldots,n$) dönüşümü;
Jacobian hesabı; Fubini ile bağımsızlık.

\textit{(iv):} $Z = \sqrt{n}(\bar{X}_n-\mu)/\sigma \sim \Normal(0,1)$;
$V = (n-1)S^2/\sigma^2 \sim \chi^2_{n-1}$ bağımsız.
$T = Z/\sqrt{V/(n-1)} \sim t_{n-1}$ (tanım gereği).
\end{proof}

\begin{hataanalizi}
"$S^2 = \frac{1}{n}\sum(X_i-\bar{X}_n)^2$" örneklem varyansı için yaygın bir notasyon.
Ancak bu yanlı tahmin edicidir: $E[S_n^2] = (n-1)\sigma^2/n \neq \sigma^2$.
Kitabımızda $S^2 = \frac{1}{n-1}\sum(X_i-\bar{X}_n)^2$ (Bessel düzeltmesi) kullanılır:
$E[S^2] = \sigma^2$.
\end{hataanalizi}

\begin{teorem}[Örneklem Ortalamasının Büyük Örneklem Dağılımı]
\label{thm:buyuk_orneklem_dagilim}
$X_1,\ldots,X_n$ i.i.d., $E[X_i]=\mu$, $\mathrm{Var}(X_i)=\sigma^2<\infty$.
Dağılımı ne olursa olsun: $\bar{X}_n \xrightarrow{d} \Normal(\mu, \sigma^2/n)$ (MLT).
$S^2 \xrightarrow{P} \sigma^2$ (BSK + sürekli haritalama).
$T_n = \sqrt{n}(\bar{X}_n-\mu)/S \xrightarrow{d} \Normal(0,1)$ (Slutsky).
\end{teorem}

% =====================================================================
\section{Yeterli İstatistik}
\label{sec:yeterli_istatistik}
% =====================================================================

\begin{kesif}
$n$ gözlemden elde edilen tüm bilgi için $T(X_1,\ldots,X_n)$'yi hesaplamak yeterli mi?
Yani $T$'yi bilince, orijinal veriyi $\theta$ hakkında daha fazla bilgi için kullanabilir miyiz?
Yanıt "hayır" ise $T$ yeterli istatistik adını alır.
\end{kesif}

\begin{tanimgerekce}
$X_1,\ldots,X_n$ birlikte $\theta$ hakkında tüm bilgiyi içerir.
$T$ "yeterli"yse, $T$'nin değeri bilindiğinde verinin $\theta$'dan bağımsız bir kalan dağılımı olur.
Başka deyişle $T$, $\theta$ bilgisini "özetlemiş"tir ve orijinal veriden geriye kalan
rasgelelik $\theta$ ile ilgisizdir.
\end{tanimgerekce}

\begin{tanim}[Yeterli İstatistik]
\label{def:yeterli}
$T = T(\mathbf{X})$ istatistiği $\theta$ için \textbf{yeterli} (sufficient) dır, eğer
$\mathbf{X}$'in koşullu dağılımı $T = t$ verildiğinde $\theta$'dan bağımsız ise, yani
$P_\theta(\mathbf{X} \in A \mid T = t)$ her $t$ için $\theta$'ya bağlı değilse.
\end{tanim}

\begin{teorem}[Fisher-Neyman Çarpanlaştırma Teoremi]
\label{thm:fisher_neyman}
$\mathbf{X}$ ortak PDF/PMF'si $f(\mathbf{x} \mid \theta)$ olan rasgele vektör.
$T(\mathbf{X})$, $\theta$ için yeterlidir ancak ve ancak
\[
  f(\mathbf{x} \mid \theta) = g(T(\mathbf{x}), \theta) \cdot h(\mathbf{x})
\]
ayrışımı mevcut ise; burada $g$ yalnızca $T(\mathbf{x})$ ve $\theta$'ya, $h$ yalnızca $\mathbf{x}$'e bağlıdır.
\end{teorem}

\begin{proof}
\textbf{Kesikli durum ($\Leftarrow$):}
$P_\theta(\mathbf{X}=\mathbf{x} \mid T(\mathbf{X})=t) = \frac{P_\theta(\mathbf{X}=\mathbf{x},\, T(\mathbf{X})=t)}{P_\theta(T(\mathbf{X})=t)}$.
$T(\mathbf{x})=t$ ise pay $= f(\mathbf{x}|\theta) = g(t,\theta)h(\mathbf{x})$.
Payda $= \sum_{\mathbf{x}': T(\mathbf{x}')=t} g(t,\theta)h(\mathbf{x}') = g(t,\theta)\sum_{\mathbf{x}':T=t}h(\mathbf{x}')$.
Oran $= h(\mathbf{x})/\sum_{\mathbf{x}':T=t}h(\mathbf{x}')$ — $\theta$'dan bağımsız. $\Rightarrow$ $T$ yeterli.

\textbf{($\Rightarrow$):} $T$ yeterli olsun.
$f(\mathbf{x}|\theta) = P_\theta(\mathbf{X}=\mathbf{x}, T=T(\mathbf{x}))
= P_\theta(\mathbf{X}=\mathbf{x}|T=T(\mathbf{x})) \cdot P_\theta(T=T(\mathbf{x}))$.
$P_\theta(\mathbf{X}=\mathbf{x}|T=T(\mathbf{x})) =: h(\mathbf{x})$ ($\theta$'dan bağımsız).
$P_\theta(T=T(\mathbf{x})) =: g(T(\mathbf{x}),\theta)$.
Ayrışım sağlandı.
\end{proof}

\begin{ornek}[Binom için Yeterli İstatistik]
$X_1,\ldots,X_n$ i.i.d.\ $\mathrm{Bernoulli}(p)$.
$f(\mathbf{x}|p) = \prod p^{x_i}(1-p)^{1-x_i} = p^{\sum x_i}(1-p)^{n-\sum x_i}$.
$g(T,p) = p^T(1-p)^{n-T}$ ($T = \sum x_i$), $h(\mathbf{x}) = 1$.
Fisher-Neyman: $T = \sum X_i$ yeterli.
\end{ornek}

\begin{ornek}[Normal için Yeterli İstatistik]
$X_i \sim \Normal(\mu, \sigma^2)$, $\theta = (\mu, \sigma^2)$.
$f(\mathbf{x}|\theta) = (2\pi\sigma^2)^{-n/2}\exp\!\left(-\frac{\sum(x_i-\mu)^2}{2\sigma^2}\right)$.
$\sum(x_i-\mu)^2 = \sum x_i^2 - 2\mu\sum x_i + n\mu^2$.
$g(T_1,T_2,\theta) = (2\pi\sigma^2)^{-n/2}\exp\!\left(\frac{2\mu T_1 - T_2 - n\mu^2}{2\sigma^2}\right)$,
$h(\mathbf{x})=1$; $T_1 = \sum x_i$, $T_2 = \sum x_i^2$.
$(\sum X_i, \sum X_i^2)$ — dolayısıyla $(\bar{X}_n, S^2)$ — yeterli.
\end{ornek}

\begin{ornek}[Üstel Dağılım]
$X_i \sim \mathrm{Exp}(\lambda)$: $f(\mathbf{x}|\lambda) = \lambda^n e^{-\lambda\sum x_i}$.
$T = \sum X_i$ yeterli; $g(T,\lambda) = \lambda^n e^{-\lambda T}$, $h = 1$.
\end{ornek}

\begin{teorem}[Üstel Ailede Yeterli İstatistik]
\label{thm:ustel_yeterli}
Tek parametreli üstel aile $f(x|\theta) = h(x)\exp(\eta(\theta)T(x) - B(\theta))$.
$n$ gözlemde: $f(\mathbf{x}|\theta) = \prod h(x_i) \cdot \exp\!\bigl(\eta(\theta)\sum T(x_i) - nB(\theta)\bigr)$.
Fisher-Neyman: $\sum_{i=1}^n T(X_i)$ yeterli istatistiktir.
\end{teorem}

\begin{mikroozet}
Yeterlilik $\equiv$ çarpanlaştırma: $f = g(T,\theta)\cdot h(\mathbf{x})$.
Üstel ailede her zaman doğal yeterli istatistik: $\sum T(X_i)$.
\end{mikroozet}

% =====================================================================
\section{Tamamlayıcı İstatistik}
\label{sec:tamamlayici}
% =====================================================================

\begin{tanim}[Tamamlayıcı İstatistik]
\label{def:tamamlayici}
$T$ istatistiği $\theta$ için \textbf{tamamlayıcı} (complete) dır, eğer
her ölçülebilir $g : \mathbb{R} \to \mathbb{R}$ için
\[
  E_\theta[g(T)] = 0 \quad \forall\,\theta \in \Theta \implies g(T) = 0 \quad \text{n.k.\ (her } \theta\text{'da)}.
\]
\end{tanim}

\begin{onerme}[Tamamlayıcılığın Önemi]
Tamamlayıcı bir istatistik varsa, $T$ üzerinden ifade edilen her yansız tahmin edici tek türlüdür.
Bu, UMVUE (en küçük varyanslı yansız tahmin edici) teorisinin kilit adımıdır
(Bölüm~\ref{chp:nokta_tahmin}, Lehmann-Scheffé teoremi).
\end{onerme}

\begin{teorem}[Üstel Ailede Tamamlayıcılık]
\label{thm:ustel_tam}
Tam rankli (tam parametreli) tek parametreli üstel ailede yeterli istatistik
$T = \sum T(X_i)$ tamamlayıcıdır.
\end{teorem}

\begin{proof}
$E_\theta[g(T)] = 0$ tüm $\theta$'da koşulu, Laplace ya da Moment dönüşümü argümanıyla
$g = 0$ a.e.\ verir. Üstel aile ailesiyle MÜF biricikliği kullanılır.
Ayrıntı: Lehmann-Casella~\cite{Lehmann1998}, Teorem 1.6.22.
\end{proof}

\begin{ornek}[Binom Tamamlayıcılığı]
$T = \sum X_i \sim \Binom(n,p)$. $E_p[g(T)] = 0$ tüm $p \in (0,1)$:
$\sum_{k=0}^n g(k)\binom{n}{k}p^k(1-p)^{n-k} = 0$.
Bu $p$ polinomu — sıfır polinom ancak $g(k) = 0$ tüm $k$'larda.
$T$ tamamlayıcı.
\end{ornek}

% =====================================================================
\section{En Az Yeterli İstatistik}
\label{sec:en_az_yeterli}
% =====================================================================

\begin{tanim}[En Az Yeterli İstatistik]
\label{def:en_az_yeterli}
$T$ \textbf{en az yeterli} (minimal sufficient) dır, eğer:
\begin{enumerate}[(i)]
  \item $T$ yeterlidir.
  \item Her yeterli $U$ için, $U$'nun bir fonksiyonu olarak $T$ ifade edilebilir:
        $T = h(U)$ (ölçülebilir $h$).
\end{enumerate}
En az yeterli istatistik, veriyi mümkün olan en yoğun biçimde özetler.
\end{tanim}

\begin{teorem}[Lehmann-Scheffé En Az Yeterlilik Kriteri]
\label{thm:ls_en_az}
$f(\mathbf{x}|\theta)/f(\mathbf{y}|\theta)$ oranının $\theta$'dan bağımsız olması
$\mathbf{x}$ ve $\mathbf{y}$'nin aynı $T$ değerini vermesiyle eşdeğer ise,
$T$ en az yeterlidir.
\end{teorem}

\begin{ornek}[Normal — En Az Yeterlilik]
$X_i \sim \Normal(\mu, \sigma^2)$.
$f(\mathbf{x}|\theta)/f(\mathbf{y}|\theta) = \exp\!\bigl(\frac{\sum x_i^2-\sum y_i^2}{-2\sigma^2}
+ \frac{\mu(\sum x_i - \sum y_i)}{\sigma^2}\bigr)$.
Bu $\theta$'dan bağımsız $\Leftrightarrow$ $\sum x_i = \sum y_i$ ve $\sum x_i^2 = \sum y_i^2$.
$(\sum X_i, \sum X_i^2)$ en az yeterli (= tamamlayıcı yeterli $\Rightarrow$ UMVUE mevcut).
\end{ornek}

\begin{hocanotu}
Tamamlayıcı yeterli istatistik (complete sufficient statistic) her zaman en az yeterlidir.
Tersine, en az yeterli istatistik her zaman tamamlayıcı olmak zorunda değildir —
ama üstel aileler için ikisi çakışır.
\end{hocanotu}

% ---- Disiplinlerarası Bağlantı (TYMM) ----
\begin{kopru}[Disiplinlerarası — Veri Bilimi, Biyomedikal Araştırma, Ekonometri]
\textbf{Veri bilimi:} "Hangi özellikler (features) modeli tam olarak özetler?" sorusu,
yeterli istatistik kavramının makine öğrenmesi karşılığıdır.
Sıkıştırma (compression) teorisi, yeterlilik ile entropi arasındaki bağı kurar.\\
\textbf{Klinik araştırma:} İlaç klinik denemelerinde birincil uç nokta (primary endpoint)
seçimi bir istatistik seçim problemidir; yeterli ve duyarlı bir istatistik
çalışmanın güç (power) hesabını doğrudan etkiler.\\
\textbf{Ekonometri:} Araç değişken (instrumental variable) ve panel veri modellerinde
parametre tanımlanabilirliği, bu bölümdeki kavramlarla doğrudan ilişkilidir.
\end{kopru}

% ---- Öz-Yansıtma ----
\begin{ozyansima}
\begin{itemize}[leftmargin=1.8em]
  \item "Parametrik model" ve "parametrik olmayan model" arasındaki farkı bir örnekle açıklayın.
  \item Yeterli istatistiğin "veri özetleme" yorumunu ve pratikte neden önemli olduğunu anlatın.
  \item $(\bar{X}_n, S^2)$'nin $\sigma^2$ bilinen $\Normal(\mu,\sigma^2)$'de yeterli istatistik olup
        olmadığını tartışın. Tek bilinmeyen $\mu$ ise ne değişir?
  \item OKH ayrışımı $= \mathrm{Var} + \mathrm{Yan}^2$ istatistiksel bir "uzlaşı"yı nasıl yansıtır?
\end{itemize}
\end{ozyansima}

% ---- Köprü Notu ----
\begin{kopru}
\textbf{Önceki kısım:} Part~I'deki dağılım teorisi (Böl.~\ref{chp:dagilim_aileleri}),
MLT (Böl.~\ref{chp:merkezi_limit}) ve koşullu beklenti (Böl.~\ref{chp:beklenti})
bu bölümde doğrudan kullanıldı.\\
\textbf{Sonraki bölüm:} Bölüm~\ref{chp:nokta_tahmin} — yeterli istatistik + tamamlayıcılık
$\to$ UMVUE; Fisher bilgisi + Cramér-Rao $\to$ verimlilik alt sınırı; MLE.\\
\textbf{Bölüm~\ref{chp:guven_araliklari}:} Pivot yöntemi örneklem dağılımlarına dayanır.
\defterbaglantisi{bolum11\_istatistik\_temelleri.ipynb}{%
  Örneklem dağılımı simülasyonu, yeterli istatistik görselleştirmesi}
\end{kopru}

% =====================================================================
\section{Sıra İstatistikleri}
\label{sec:sira_istatistik}
% =====================================================================

\begin{tanim}[Sıra istatistikleri]\label{def:sira_ist}
$X_1,\ldots,X_n\iid F$ sürekli dağılım. Sıralı değerler $X_{(1)}\leq X_{(2)}\leq\cdots\leq X_{(n)}$ \textbf{sıra istatistikleri} \emph{(order statistics)}, toplu vektörü $(X_{(1)},\ldots,X_{(n)})$ \textbf{sıralanmış örneklem} adını alır.
\end{tanim}

\begin{teorem}[Sıra istatistiği yoğunluğu]\label{thm:sira_yogunluk}
$X_{(k)}$'nın yoğunluğu:
\[
  f_{X_{(k)}}(x) = \frac{n!}{(k-1)!(n-k)!}[F(x)]^{k-1}[1-F(x)]^{n-k}f(x).
\]
$X_{(1)}$ için $k=1$: $f_{X_{(1)}}(x) = n[1-F(x)]^{n-1}f(x)$.
$X_{(n)}$ için $k=n$: $f_{X_{(n)}}(x) = n[F(x)]^{n-1}f(x)$.
\end{teorem}

\begin{proof}
Kombinatorik argüman: $P(X_{(k)}\in[x,x+dx])$ = $\binom{n}{1}\binom{n-1}{k-1}$ yolu ile $k-1$ gözlem $<x$, bir gözlem $\in[x,x+dx]$, $n-k$ gözlem $>x$.
\end{proof}

\begin{teorem}[Min-max için özel formüller]\label{thm:min_max}
$X_i\iid F$ bağımsız.
\begin{enumerate}[(a)]
  \item $F_{X_{(1)}}(x) = 1-[1-F(x)]^n$ (minimum KDF).
  \item $F_{X_{(n)}}(x) = [F(x)]^n$ (maksimum KDF).
  \item Aralık (range): $R_n=X_{(n)}-X_{(1)}$.
\end{enumerate}
\end{teorem}

\begin{ornek}[Uniform sıra istatistikleri]\label{ex:uniform_sira}
$X_i\iid\Uniform[0,1]$. $X_{(k)}\sim\mathrm{Beta}(k,n-k+1)$; $E[X_{(k)}]=k/(n+1)$. Dolayısıyla sıra istatistiklerinin beklentileri $[0,1]$'i eşit aralıklara böler.
\end{ornek}

% =====================================================================
\section{Ampirik Dağılım Fonksiyonu}
\label{sec:ampirik_dagilim}
% =====================================================================

\begin{tanim}[Ampirik dağılım fonksiyonu]\label{def:edf}
$X_1,\ldots,X_n\iid F$. \textbf{Ampirik dağılım fonksiyonu} \emph{(empirical distribution function, EDF)}:
\[
  F_n(x) = \frac{1}{n}\sum_{i=1}^n\mathbf{1}[X_i\leq x].
\]
$F_n$ sıra istatistiklerinin basamak fonksiyonu: $F_n(X_{(k)}) = k/n$.
\end{tanim}

\begin{teorem}[Glivenko-Cantelli]\label{thm:glivenko_cantelli}
$\|F_n - F\|_\infty = \sup_x|F_n(x)-F(x)|\asto0$.
\end{teorem}

\begin{proof}[Proof (taslak)]
Her sabit $x$ için $F_n(x)\asto F(x)$ (GBSK). Süreklilik + sonlu sayıda bölüm argümanıyla süpremum kontrolü sağlanır. Detay: Bölüm~\ref{chp:buyuk_sayilar}.
\end{proof}

\begin{teorem}[Kolmogorov-Smirnov]\label{thm:ks_istatistik}
$F$ sürekli altında $\sqrt{n}\|F_n-F\|_\infty\dto\|\mathbb{B}\|_\infty$ (Bölüm~\ref{chp:asimptotik}, Donsker teoremi). $H_0:F=F_0$ testi:
\[
  D_n = \sqrt{n}\sup_x|F_n(x)-F_0(x)|.
\]
$\alpha=0.05$ kritik değeri $\approx1.36$.
\end{teorem}

% =====================================================================
\section{Neyman-Fisher Çarpanlaştırma — Daha Fazla Örnek}
\label{sec:carpanlastirma_ornekler}
% =====================================================================

\begin{ornek}[İki parametreli Normal — $\mu$ bilinmiyor, $\sigma^2$ sabit]\label{ex:normal_mu_yeterli}
$X_i\iid\Normal(\mu,\sigma_0^2)$. Fisher-Neyman:
\[
  f(\mathbf{x}|\mu) = (2\pi\sigma_0^2)^{-n/2}\exp\!\left(-\frac{\sum x_i^2}{2\sigma_0^2}+\frac{\mu\sum x_i}{\sigma_0^2}-\frac{n\mu^2}{2\sigma_0^2}\right).
\]
$g(T,\mu)=\exp(\mu T/\sigma_0^2-n\mu^2/(2\sigma_0^2))$ ($T=\sum X_i$), $h=$ sabit çarpan. $\bar{X}_n$ yeterli.
\end{ornek}

\begin{ornek}[Çift parametreli üstel aile]\label{ex:cift_param_ustel}
Normal$(\mu,\sigma^2)$ her iki bilinmeyen: $\eta_1(\theta)=\mu/\sigma^2$, $\eta_2(\theta)=-1/(2\sigma^2)$. Doğal yeterli istatistik: $(\sum X_i, \sum X_i^2)$; tam rankli — tamamlayıcı.
\end{ornek}

\begin{teorem}[Yeterli istatistik ve bağımsızlık]\label{thm:yeterli_bagimsizlik}
$T$ yeterli ve $V$ herhangi bir istatistik. Eğer $\Prob_\theta(V\leq v\mid T=t)$ $\theta$'dan bağımsız ise, $V$'nin dağılımı yalnızca $T$'nin değerine bağlıdır. Özellikle $V$ bir pivot seçilirse çıkarım $T$ aracılığıyla yapılır.
\end{teorem}

% =====================================================================
\section{Bayes Çıkarımına Giriş}
\label{sec:bayes_giris}
% =====================================================================

\begin{kesif}
Frekansçı istatistikte $\theta$ sabit (bilinmeyen) bir değerdir; Bayesçi istatistikte ise $\theta$ da rasgele bir değişkendir ve ona dair önsel (prior) inancımızı bir olasılık dağılımı olarak ifade ederiz.
\end{kesif}

\begin{tanim}[Bayes teoremi — istatistik]\label{def:bayes_istatistik}
$\theta$'nın önsel dağılımı $\pi(\theta)$, veri olabilirliği $f(\mathbf{x}|\theta)$. \textbf{Sonsal dağılım} \emph{(posterior distribution)}:
\[
  \pi(\theta|\mathbf{x}) = \frac{f(\mathbf{x}|\theta)\pi(\theta)}{m(\mathbf{x})}, \quad m(\mathbf{x}) = \int f(\mathbf{x}|\theta)\pi(\theta)\,d\theta.
\]
$m(\mathbf{x})$ marjinal olabilirlik \emph{(marginal likelihood)} veya kanıt \emph{(evidence)}.
\end{tanim}

\begin{tanim}[Konjugat önsel]\label{def:konjugat_onsel}
$\pi(\theta)$ \textbf{konjugat önsel} \emph{(conjugate prior)}: sonsal $\pi(\theta|\mathbf{x})$ önsel ile aynı ailedendir. Hesap kolaylığı sağlar.
\begin{itemize}
  \item $\Binom(n,p)$ için $p\sim\Beta(\alpha,\beta)$: sonsal $\Beta(\alpha+\sum x_i, \beta+n-\sum x_i)$.
  \item $\Pois(\lambda)$ için $\lambda\sim\Gama(\alpha,\beta)$: sonsal $\Gama(\alpha+n\bar x, \beta+n)$.
  \item $\Normal(\mu,\sigma^2)$ ($\sigma^2$ bilinen) için $\mu\sim\Normal(\mu_0,\tau^2)$: sonsal $\Normal(\mu_1,\tau_1^2)$.
\end{itemize}
\end{tanim}

\begin{teorem}[Normal konjugat güncelleme]\label{thm:normal_konjugat}
$X_1,\ldots,X_n\iid\Normal(\mu,\sigma^2)$ ($\sigma^2$ bilinen); $\mu\sim\Normal(\mu_0,\tau^2)$. Sonsal:
\[
  \mu|\mathbf{X}\sim\Normal(\mu_1,\tau_1^2), \quad
  \mu_1 = \frac{\mu_0/\tau^2 + n\bar{X}/\sigma^2}{1/\tau^2+n/\sigma^2}, \quad
  \tau_1^2 = \frac{1}{1/\tau^2+n/\sigma^2}.
\]
\end{teorem}

\begin{proof}
$\pi(\mu|\mathbf{x})\propto\pi(\mu)f(\mathbf{x}|\mu)$. İkisi de Gaussian; çarpım da Gaussian (karede lineer). Üsteki $\exp[-(\mu-\mu_1)^2/(2\tau_1^2)+\text{sabit}]$ biçimine getirilince $\mu_1$ ve $\tau_1^2$ elde edilir.
\end{proof}

\begin{ornek}[Bayesçi ağırlıklı ortalama]
$\mu_1 = w_0\mu_0 + w_1\bar{X}$ ($w_0=(\sigma^2/n)/(\tau^2+\sigma^2/n)$, $w_1=\tau^2/(\tau^2+\sigma^2/n)$). Sonsal ortalama, önsel ortalama ile örneklem ortalamasının ağırlıklı ortalamasıdır. $n\to\infty$: $\mu_1\to\bar{X}$ (veri öneli baskılar); $\tau\to\infty$: $\mu_1\to\bar{X}$ (bilgisiz önsel).
\end{ornek}

\begin{teorem}[Bayes tahmincisi ve OKH]\label{thm:bayes_tahmin}
Karesel kayıp altında Bayes tahmincisi sonsal ortalama: $\hat\theta_{\text{Bayes}}=E[\theta|\mathbf{X}]$.

Mutlak değer kaybında Bayes tahmincisi sonsal medyan.
\end{teorem}

\begin{proof}
$E_{\theta|\mathbf{x}}[(\theta-d)^2]$ minimizasyonu: türev alıp sıfıra eşitle; $d^*=E[\theta|\mathbf{x}]$.
\end{proof}

\begin{kritiksorgu}
Frekansçı ve Bayesçi yaklaşım ne zaman aynı sonucu verir? Büyük örneklemde (Bernstein-von Mises teoremi) sonsal $\Normal(\hat\theta_{\mathrm{MLE}}, I(\hat\theta)^{-1}/n)$'ye yakınsar — Bayes ve frekansçı güven aralıkları örtüşür. Küçük örneklemde önsel seçim kritik önem taşır.
\end{kritiksorgu}

% =====================================================================
\section{Yeterli İstatistik ve Bayes}
\label{sec:yeterli_bayes}
% =====================================================================

\begin{teorem}[Yeterlilik ve Bayes]\label{thm:yeterlilik_bayes}
$T$ yeterli istatistik ise $\pi(\theta|\mathbf{x}) = \pi(\theta|T(\mathbf{x}))$. Yani Bayesçi güncelleme tamamen $T$ üzerinden gerçekleşir; orijinal verinin tüm $\theta$-ilgili bilgisi $T$'de toplanmıştır.
\end{teorem}

\begin{proof}
$\pi(\theta|\mathbf{x})\propto f(\mathbf{x}|\theta)\pi(\theta) = g(T(\mathbf{x}),\theta)h(\mathbf{x})\pi(\theta)$. $h(\mathbf{x})$ $\theta$'dan bağımsız; normalleştirmede iptal olur. Dolayısıyla $\pi(\theta|\mathbf{x})\propto g(T,\theta)\pi(\theta)=\pi(\theta|T)$.
\end{proof}

\begin{ornek}[Bayes ve yeterlilik — Bernoulli]
$X_i\iid\mathrm{Ber}(p)$; $T=\sum X_i$ yeterli. Sonsal $\pi(p|T=t)=\pi(p|\mathbf{x})$ — 100 gözlem yerine sadece $t=\sum x_i$ yeterli. Hiçbir ek bilgi kaybolmadı.
\end{ornek}

\begin{disiplinlerarasibaglan}
\textbf{Makine öğrenmesi:} Bayesçi sinir ağları (BNN) belirsizlik tahmini için sonsal dağılımı öğrenir. \textbf{Tıbbi teşhis:} Ön-test/son-test olasılığı (Bayes güncellemesi) tanısal testlerin yorumunda temel. \textbf{Sinyal işleme:} Kalman filtresi Bayesçi güncellemeler zinciridir (Gaussian konjugat).
\end{disiplinlerarasibaglan}

% =====================================================================
%  ALIŞTIRMALAR
% =====================================================================

\cozumlualistirmalar

\begin{alistirma}[Al.11.1]{A}{Örneklem Ortalaması ve Varyansı}
$X_1,\ldots,X_5$ i.i.d.\ $\Normal(3, 4)$ ($\sigma^2=4$).
(a) $\bar{X}_5$'in dağılımını yazın.
(b) $P(\bar{X}_5 > 4)$'ü hesaplayın.
(c) $(n-1)S^2/\sigma^2$'nin dağılımını yazın ve $E[S^2]$'yi bulun.
\end{alistirma}
\alcozum{%
(a) $\bar{X}_5 \sim \Normal(3, 4/5)$; $\sigma_{\bar{X}} = 2/\sqrt{5}$.\\
(b) $P(\bar{X}_5>4)=P(Z>(4-3)/(2/\sqrt{5}))=P(Z>\sqrt{5}/2)=P(Z>1.118)=1-\Phi(1.118)\approx0.132$.\\
(c) $4S^2/4 = S^2 \sim \chi^2_4/1$: $(n-1)S^2/\sigma^2 = 4S^2/4 \sim \chi^2_4$.
$E[S^2] = \sigma^2 = 4$.}

\begin{alistirma}[Al.11.2]{A}{Fisher-Neyman — Poisson}
$X_1,\ldots,X_n$ i.i.d.\ $\Pois(\lambda)$.
(a) Ortak PMF'yi yazın. (b) $T = \sum X_i$'nin yeterli istatistik olduğunu F-N teoremi ile gösterin.
(c) $T \sim \Pois(n\lambda)$ olduğunu belirtin.
\end{alistirma}
\alcozum{%
(a) $f(\mathbf{x}|\lambda) = \prod \frac{\lambda^{x_i}e^{-\lambda}}{x_i!} = \frac{\lambda^{\sum x_i}e^{-n\lambda}}{\prod x_i!}$.\\
(b) $g(T,\lambda) = \lambda^T e^{-n\lambda}$, $h(\mathbf{x}) = 1/\prod x_i!$ — ayrışım var.
F-N: $T = \sum X_i$ yeterli.\\
(c) Poisson yeniden üretkenliği: $T \sim \Pois(n\lambda)$.}

\begin{alistirma}[Al.11.3]{B}{OKH Hesabı}
$X_1,\ldots,X_n$ i.i.d.\ $\Normal(\mu,\sigma^2)$.
$\tilde{S}^2 = \frac{1}{n}\sum(X_i-\bar{X})^2$ (yanlı tahmin edici) için:
(a) $E[\tilde{S}^2]$ ve $\mathrm{Var}(\tilde{S}^2)$'yi bulun.
(b) $\mathrm{OKH}(\tilde{S}^2, \sigma^2)$'yi hesaplayın.
(c) $\mathrm{OKH}(\tilde{S}^2,\sigma^2)$ ile $\mathrm{OKH}(S^2,\sigma^2)$'yi karşılaştırın.
\end{alistirma}
\alcozum{%
(a) $\tilde{S}^2 = \frac{n-1}{n}S^2$; $E[\tilde{S}^2] = \frac{n-1}{n}\sigma^2$.
$\mathrm{Var}(\tilde{S}^2) = \left(\frac{n-1}{n}\right)^2\mathrm{Var}(S^2) = \left(\frac{n-1}{n}\right)^2 \frac{2\sigma^4}{n-1} = \frac{2(n-1)\sigma^4}{n^2}$.\\
(b) $\mathrm{Yan} = -\sigma^2/n$; $\mathrm{OKH}(\tilde{S}^2) = \frac{2(n-1)\sigma^4}{n^2}+\frac{\sigma^4}{n^2} = \frac{(2n-1)\sigma^4}{n^2}$.\\
(c) $\mathrm{OKH}(S^2) = \mathrm{Var}(S^2) = \frac{2\sigma^4}{n-1}$.
$n \geq 3$'te $\tilde{S}^2$ daha küçük OKH'ye sahip (küçük yanlılık kazancı).}

\begin{alistirma}[Al.11.4]{B}{Tamamlayıcılık — Uniform}
$X_1,\ldots,X_n$ i.i.d.\ $\Uniform(0,\theta)$.
(a) $T = X_{(n)} = \max X_i$'nin yeterli olduğunu F-N ile gösterin.
(b) $E_\theta[g(T)] = 0$ tüm $\theta > 0$'da koşulundan $g \equiv 0$ sonucunu türetin.
(c) $T$ tamamlayıcı mı?
\end{alistirma}
\alcozum{%
(a) $f(\mathbf{x}|\theta) = \theta^{-n}\mathbf{1}_{x_{(n)}\leq\theta}$;
$g(T,\theta) = \theta^{-n}\mathbf{1}_{T\leq\theta}$, $h=1$. F-N: $T$ yeterli.\\
(b) $f_{X_{(n)}}(t|\theta) = nt^{n-1}/\theta^n$ ($0<t<\theta$).
$E[g(T)] = \int_0^\theta g(t)\frac{nt^{n-1}}{\theta^n}dt = 0$.
$\frac{d}{d\theta}[\theta^n \cdot 0] = 0 = ng(\theta)\theta^{n-1}$, yani $g(\theta)=0$ tüm $\theta>0$.
(c) Evet, tamamlayıcı.}

% ---

\alistirmalar

\begin{alistirma}[Al.11.5]{A}{Örneklem $t$ Dağılımı}
$X_1,\ldots,X_{16}$ i.i.d.\ $\Normal(\mu, 9)$.
$T = \sqrt{16}(\bar{X}_{16}-\mu)/S$'nin dağılımını yazın.
$P(T > 2.131)$'i $t$ tablosundan söyleyin.
\end{alistirma}
\alcozum{%
$T \sim t_{15}$.
$t_{15,0.025} \approx 2.131$: $P(T > 2.131) = 0.025$.}

\begin{alistirma}[Al.11.6]{A}{Yeterli İstatistik — Üstel Aile}
$X_i \sim \Gama(\alpha, \beta)$ ($\alpha$ biliniyor, $\beta$ bilinmiyor).
(a) Ortak PDF'yi üstel aile formunda yazın. (b) Yeterli istatistiği bulun.
\end{alistirma}
\alcozum{%
(a) $f(\mathbf{x}|\beta) = \prod \frac{\beta^\alpha}{\Gamma(\alpha)}x_i^{\alpha-1}e^{-\beta x_i}
= \left(\frac{\beta^\alpha}{\Gamma(\alpha)}\right)^n \prod x_i^{\alpha-1} \cdot e^{-\beta\sum x_i}$.
(b) $g(T,\beta) = (\beta^\alpha/\Gamma(\alpha))^n e^{-\beta T}$, $h = \prod x_i^{\alpha-1}$.
$T = \sum X_i$ yeterli.}

\begin{alistirma}[Al.11.7]{B}{Bağımsızlık — $\bar{X}$ ve $S^2$}
$X_1, X_2, X_3$ i.i.d.\ $\Normal(\mu, \sigma^2)$.
$(Y_1, Y_2, Y_3) = (X_1+X_2+X_3, X_1-X_2, X_1+X_2-2X_3)/\sqrt{6}$ dönüşümünü yapın.
(a) $Y_i$'lerin $\Normal(?, ?)$ dağılımını bulun.
(b) $Y_1$ ile $(Y_2, Y_3)$'ün bağımsız olduğunu gösterin.
(c) $\bar{X}$ ve $S^2$'nin bağımsızlığını bu sonuçtan çıkarın.
\end{alistirma}
\alcozum{%
(a) $Y_1 = (X_1+X_2+X_3)/\sqrt{3}\cdot\sqrt{3}/\sqrt{6}\cdot\sqrt{6}$ — uygun normalize ile
$Y_1 = \sqrt{3}\bar{X}$: $\Normal(\sqrt{3}\mu, \sigma^2)$.
$Y_2 = (X_1-X_2)/\sqrt{2}$: $\Normal(0,\sigma^2)$. $Y_3 = (X_1+X_2-2X_3)/\sqrt{6}$: $\Normal(0,\sigma^2)$.\\
(b) $(Y_1,Y_2,Y_3)$ bağımsız normal: kovaryans matrisi köşegen (ortogonal dönüşüm).
$\mathrm{Cov}(Y_1,Y_2) = 0$; normal için $\Rightarrow$ bağımsız.\\
(c) $\bar{X} = Y_1/\sqrt{3}$; $S^2 = (Y_2^2+Y_3^2)/2\sigma^2 \cdot \sigma^2/1$
($(Y_2^2+Y_3^2)/\sigma^2 \sim \chi^2_2$). $Y_1 \perp (Y_2,Y_3)$ $\Rightarrow$ $\bar{X} \perp S^2$.}

\begin{alistirma}[Al.11.8]{B}{En Az Yeterlilik — Cauchy}
$X_1,\ldots,X_n$ i.i.d.\ $\mathrm{Cauchy}(\theta, 1)$.
$f(\mathbf{x}|\theta)/f(\mathbf{y}|\theta)$ oranını yazın.
Bu oranın $\theta$'dan bağımsız olması için hangi koşul gerekir?
En az yeterli istatistik nedir?
\end{alistirma}
\alcozum{%
$f(\mathbf{x}|\theta) = \prod \frac{1}{\pi(1+(x_i-\theta)^2)}$.
Oran: $\prod \frac{1+(y_i-\theta)^2}{1+(x_i-\theta)^2}$ — bu $\theta$'dan bağımsız olması için
$(x_1,\ldots,x_n)$ ve $(y_1,\ldots,y_n)$'nin aynı sıra istatistiklerine sahip olması gerekir
(çünkü Cauchy simetrik; terimlerin $\theta$-bağımlılığı birbirini ancak permütasyonla götürür).
En az yeterli istatistik: $(X_{(1)},\ldots,X_{(n)})$ — tam sıralı vektör.
Bu, Cauchy'de boyut azaltılmış bir özet olmamasını gösterir.}

\begin{alistirma}[Al.11.9]{A}{Sıra istatistiği dağılımı}
$X_1,\ldots,X_5\iid\Uniform[0,1]$. (a) $X_{(3)}$'ün yoğunluğunu yaz. (b) $E[X_{(3)}]$'ü bul. (c) $\Prob(X_{(5)}>0.9)$'u hesapla.
\end{alistirma}
\alcozum{%
(a) $f_{X_{(3)}}(x)=5!/(2!\cdot2!)x^2(1-x)^2\cdot1=30x^2(1-x)^2$ ($0<x<1$). (b) $E[X_{(3)}]=3/(5+1)=1/2$ (Uniform sıra istatistiği beklentisi $k/(n+1)$). (c) $F_{X_{(5)}}(x)=x^5$; $\Prob(X_{(5)}>0.9)=1-0.9^5=1-0.5905=0.4095$.%
}

\begin{alistirma}[Al.11.10]{B}{Ampirik dağılım ve KS testi}
$n=5$ gözlem: $1.2, 2.5, 0.8, 3.1, 1.7$. $H_0: F=\mathrm{Exp}(1)$ (yani $F_0(x)=1-e^{-x}$) için KS istatistiğini hesapla.
\end{alistirma}
\alcozum{%
Sıralı: $0.8,1.2,1.7,2.5,3.1$. $F_n(x_k)=k/5$. $F_0$: $F_0(0.8)=0.551$, $F_0(1.2)=0.699$, $F_0(1.7)=0.817$, $F_0(2.5)=0.918$, $F_0(3.1)=0.955$. $|F_n-F_0|$: $|0.2-0.551|=0.351$, $|0.4-0.699|=0.299$, $|0.6-0.817|=0.217$, $|0.8-0.918|=0.118$, $|1-0.955|=0.045$. $D_5=0.351$; $\sqrt{5}\cdot0.351=0.785<1.36$: $H_0$ reddedilemez.%
}

\begin{alistirma}[Al.11.11]{B}{Min için $n$ belirleme}
$X_i\iid\Exp(\lambda)$. $P(X_{(1)}>1)=0.9$ olmasını istiyorsunuz. $n$ kaç olmalı?
\end{alistirma}
\alcozum{%
$X_{(1)}\sim\Exp(n\lambda)$ (minimum üstel i.i.d. bağımsızsa). $P(X_{(1)}>1)=e^{-n\lambda}=0.9$; $n=-\log(0.9)/\lambda$. $\lambda=1$: $n=-\log(0.9)\approx0.105$; yani $n=1$ yetmez, tam tam değil — düz: $e^{-n}=0.9$; $n=0.105$ gerçek sayı. Yani tek gözlemde $\lambda=0.105$ ya da $n=1$, $\lambda=0.105$. Pratik: $\lambda=1$ ise $P(X_{(1)}>1)=e^{-n}=0.9$ için $n=\lceil-\log0.9\rceil=1$.%
}

% =====================================================================
\endinput

%% ============================================================
%% BÖLÜM 12 — Nokta Tahmini
%% Olasılık Teorisi ve Matematiksel İstatistik
%% ============================================================
\chapter{Nokta Tahmini}
\label{chp:nokta_tahmin}

\bolumalinti{%
  It is better to be roughly right than precisely wrong.
}{John Maynard Keynes}

\begin{kazanimlar}
Bu bölümün sonunda okuyucu:
\begin{itemize}
  \item Tahmin edici kavramını tanımlar; yansızlık, verimlilik ve tutarlılık özelliklerini ispat ederek karşılaştırır.
  \item Fisher bilgi matrisini hesaplar ve Cramér-Rao alt sınırını uygular.
  \item Rao-Blackwell teoremini ve Lehmann-Scheffé teoremini kullanarak UMVUE inşa eder.
  \item En çok olabilirlik yöntemini uygular; MLE'nin invaryant, tutarlı ve asimptotik normal olduğunu açıklar.
  \item Momentler yöntemiyle tahmin edici türetir.
  \item Konjugat önsel dağılımlar kullanarak Bayes tahmin edicisini hesaplar ve güvenilir aralık kurar.
\end{itemize}
\end{kazanimlar}

\begin{motivasyon}
Önceki bölümlerde bir dağılımın olasılık yapısını teorik olarak inceledik. Şimdi tersine gidiyoruz: elimizde veriden oluşan bir örneklem var, gizli parametre $\theta$'yı bu veriden tahmin etmek istiyoruz. Bu soru modern istatistiğin merkezindedir.

\medskip
\noindent\textbf{Köprü:} Bölüm~\ref{chp:istatistik_temelleri}'de tanımladığımız yeterli istatistik kavramı burada kritik rol oynayacak. Bir tahmin edicinin yeterliliği, Rao-Blackwell teoremi aracılığıyla verimliliği artırmak için kullanılacak.
\end{motivasyon}

%% ============================================================
\section{Tahmin Edici Kavramı}
\label{sec:tahmin_edici}
%% ============================================================

\begin{kesif}
$X_1,\ldots,X_n \sim \Normal(\mu,\sigma^2)$ i.i.d.\ olsun. $\mu$'yü nasıl tahmin ederiz? Doğal yanıt $\overline{X}_n$'dir. Peki neden bu iyi bir seçim? Başka seçenekler var mı?
\end{kesif}

\begin{tanim}[Tahmin edici]\label{def:tahmin_edici}
$(X_1,\ldots,X_n)$ bir örneklem olsun. $\theta \in \Theta$ parametresinin \textbf{tahmin edicisi} \emph{(estimator)}, veriye dayanan $W = W(X_1,\ldots,X_n)$ istatistiğidir. Gözlemlenen $x_1,\ldots,x_n$ için $w = W(x_1,\ldots,x_n)$ değerine \textbf{tahmin} \emph{(estimate)} denir.
\end{tanim}

\begin{tanimgerekce}
Tahmin edici bir rasgele değişken; tahmin ise onun gözlemlenen gerçekleşmesidir. Bu ayrım çok önemli: $W = \overline{X}_n$ bir rasgele değişken, $w = \bar{x}_n = 3{,}14$ ise bir sayıdır.
\end{tanimgerekce}

\subsection{Yanlılık ve Yansızlık}

\begin{tanim}[Yanlılık]\label{def:yanlilik}
$W$ istatistiğinin $\theta$ için \textbf{yanlılığı} \emph{(bias)}:
\[
  \mathrm{Yanlılık}(W,\theta) = \Beklenti_\theta[W] - \theta.
\]
$\Beklenti_\theta[W] = \theta$ ise $W$'ye \textbf{yansız tahmin edici} \emph{(unbiased estimator)} denir.
\end{tanim}

\begin{teorem}[Standart örneklem istatistikleri]\label{thm:orn_istatistik}
$X_1,\ldots,X_n \iid F$ ile $\Beklenti[X_i]=\mu$, $\Var(X_i)=\sigma^2 < \infty$ olsun. O zaman:
\begin{enumerate}[(a)]
  \item $\overline{X}_n = \frac{1}{n}\sum_{i=1}^n X_i$ yansız: $\Beklenti[\overline{X}_n] = \mu$.
  \item $S^2 = \frac{1}{n-1}\sum_{i=1}^n (X_i - \overline{X}_n)^2$ yansız: $\Beklenti[S^2] = \sigma^2$.
  \item $\hat{\sigma}^2 = \frac{1}{n}\sum_{i=1}^n(X_i-\overline{X}_n)^2$ yanlı: $\Beklenti[\hat{\sigma}^2] = \frac{n-1}{n}\sigma^2$.
\end{enumerate}
\end{teorem}

\begin{proof}
(a) Beklentinin doğrusallığından doğrudan.

(b) Özdeşlik kullanılır:
\[
  \sum_{i=1}^n (X_i - \overline{X}_n)^2 = \sum_{i=1}^n X_i^2 - n\overline{X}_n^2.
\]
Her iki tarafın beklentisini alalım. $\Beklenti[X_i^2] = \sigma^2 + \mu^2$ ve $\Beklenti[\overline{X}_n^2] = \Var(\overline{X}_n) + \mu^2 = \sigma^2/n + \mu^2$ olduğundan
\[
  \Beklenti\!\left[\sum_{i=1}^n(X_i-\overline{X}_n)^2\right] = n(\sigma^2+\mu^2) - n\!\left(\frac{\sigma^2}{n}+\mu^2\right) = (n-1)\sigma^2.
\]
$S^2$'nin böleni $n-1$ olduğundan $\Beklenti[S^2]=\sigma^2$.

(c) $\hat{\sigma}^2 = \frac{n-1}{n}S^2$ olduğundan $\Beklenti[\hat{\sigma}^2] = \frac{n-1}{n}\sigma^2$.
\end{proof}

\begin{kritiksorgu}
$S^2$'de neden $n-1$ kullanıyoruz? Özgürlük derecesi argümanı: $\sum(X_i-\overline{X}_n)=0$ kısıtı, $n$ sayıdan bağımsız olanları $n-1$'e indirger.
\end{kritiksorgu}

\subsection{Ortalama Karesel Hata}

\begin{tanim}[OKH]\label{def:okh}
$W$ tahmin edicisinin $\theta$ için \textbf{ortalama karesel hatası} \emph{(mean squared error)}:
\[
  \mathrm{OKH}(W,\theta) = \Beklenti_\theta\!\bigl[(W-\theta)^2\bigr].
\]
\end{tanim}

\begin{teorem}[OKH ayrışımı]\label{thm:okh_ayrisim}
$\mathrm{OKH}(W,\theta) = \Var_\theta(W) + [\mathrm{Yanlılık}(W,\theta)]^2$.
\end{teorem}

\begin{proof}
$b = \Beklenti_\theta[W]-\theta$ yanlılık olsun. O zaman
\[
  (W-\theta)^2 = \bigl[(W-\Beklenti_\theta W) + b\bigr]^2
  = (W-\Beklenti_\theta W)^2 + 2b(W-\Beklenti_\theta W) + b^2.
\]
Beklenti alındığında orta terim $2b\cdot\Beklenti[W-\Beklenti W]=0$ olur.
\end{proof}

\begin{hataanalizi}
Yansızlık her zaman en iyi ölçüt değildir. Yanlı ama küçük varyanslı bir tahmin edici, yansız ama büyük varyanslı birinden daha küçük OKH'a sahip olabilir.
\end{hataanalizi}

\subsection{Tutarlılık}

\begin{tanim}[Tutarlılık]\label{def:tutarlilik}
Tahmin edici dizisi $\{W_n\}$ \textbf{tutarlı} \emph{(consistent)} ise: her $\theta$ ve $\varepsilon>0$ için
\[
  \Prob_\theta(|W_n - \theta|>\varepsilon) \to 0 \quad (n\to\infty).
\]
Yani $W_n \pto \theta$.
\end{tanim}

\begin{teorem}[Yeterli koşul]\label{thm:tutarli_yeterli}
$\Beklenti_\theta[W_n]\to\theta$ ve $\Var_\theta(W_n)\to 0$ ise $W_n$ tutarlıdır.
\end{teorem}

\begin{proof}
Chebyshev eşitsizliğinden: $\Prob(|W_n-\theta|>\varepsilon) \leq \mathrm{OKH}(W_n,\theta)/\varepsilon^2 \to 0$.
\end{proof}

%% ============================================================
\section{Cramér-Rao Alt Sınırı}
\label{sec:cramer_rao}
%% ============================================================

\begin{kesif}
Yansız tahmin ediciler arasında en küçük varyanslısını bulmak istiyoruz. Varyansın ne kadar küçük olabileceğine dair evrensel bir alt sınır var mıdır?
\end{kesif}

\subsection{Fisher Bilgi Miktarı}

\begin{tanim}[Fisher bilgisi]\label{def:fisher_bilgisi}
$f(x;\theta)$ yoğunluğu (ya da olasılık kitlesi) düzgün regularity koşullarını sağlasın. \textbf{Fisher bilgi miktarı} \emph{(Fisher information)}:
\[
  \mathcal{I}(\theta) = \Beklenti_\theta\!\left[\left(\frac{\partial}{\partial\theta}\log f(X;\theta)\right)^{\!2}\right].
\]
\end{tanim}

\begin{teorem}[Fisher bilgisinin hesabı]\label{thm:fisher_hesap}
Regularity koşulları altında:
\begin{enumerate}[(a)]
  \item $\Beklenti_\theta\!\left[\dfrac{\partial}{\partial\theta}\log f(X;\theta)\right] = 0$.
  \item $\mathcal{I}(\theta) = -\Beklenti_\theta\!\left[\dfrac{\partial^2}{\partial\theta^2}\log f(X;\theta)\right]$.
  \item $n$ i.i.d.\ gözlem için: $\mathcal{I}_n(\theta) = n\,\mathcal{I}_1(\theta)$.
\end{enumerate}
\end{teorem}

\begin{proof}
(a) $f(x;\theta)$'yı integre eden özdeliği $\theta$'ya göre türetelim:
\[
  0 = \frac{\partial}{\partial\theta}\int f(x;\theta)\,dx
  = \int \frac{\partial f}{\partial\theta}\,dx
  = \int \frac{\partial\log f}{\partial\theta} f(x;\theta)\,dx
  = \Beklenti_\theta\!\left[\frac{\partial\log f}{\partial\theta}\right].
\]
(b) İkinci türev için aynı işlemi tekrarlayalım:
\[
  0 = \frac{\partial^2}{\partial\theta^2}\int f\,dx
  = \int\!\left[\frac{\partial^2\log f}{\partial\theta^2} + \left(\frac{\partial\log f}{\partial\theta}\right)^{\!2}\right]f\,dx.
\]
Yeniden düzenlenince $\mathcal{I}(\theta) = -\Beklenti[\partial^2\log f/\partial\theta^2]$ elde edilir.

(c) İ.i.d.\ durumda $\log f(\mathbf{x};\theta) = \sum_i \log f(x_i;\theta)$; türevler bağımsız toplam olduğundan varyans $n$ ile ölçeklenir.
\end{proof}

\begin{ornek}[Normal dağılım Fisher bilgisi]
$X\sim\Normal(\mu,\sigma^2)$ sabit $\sigma^2$ ile. $\theta=\mu$ için:
\[
  \log f(x;\mu) = -\tfrac{1}{2}\log(2\pi\sigma^2) - \frac{(x-\mu)^2}{2\sigma^2}.
\]
$\partial\log f/\partial\mu = (x-\mu)/\sigma^2$; $\partial^2\log f/\partial\mu^2 = -1/\sigma^2$. Dolayısıyla $\mathcal{I}(\mu) = 1/\sigma^2$.
\end{ornek}

\subsection{Cramér-Rao Eşitsizliği}

\begin{teorem}[Cramér-Rao]\label{thm:cramer_rao}
$X_1,\ldots,X_n \iid f(x;\theta)$ ve $W = W(\mathbf{X})$ yansız tahmin edici olsun. Regularity koşulları altında:
\[
  \Var_\theta(W) \geq \frac{1}{n\,\mathcal{I}(\theta)}.
\]
\end{teorem}

\begin{proof}
Kovaryans-Cauchy-Schwarz argümanı kullanılır. $U = \partial\log f(\mathbf{X};\theta)/\partial\theta$ skor fonksiyonu olsun. Yansızlık koşulu $\int W(\mathbf{x})f(\mathbf{x};\theta)\,d\mathbf{x}=\theta$'yı $\theta$'ya göre türetelim:
\[
  1 = \int W(\mathbf{x})\frac{\partial\log f}{\partial\theta}f(\mathbf{x};\theta)\,d\mathbf{x}
  = \Beklenti_\theta[W\cdot U] = \Cov_\theta(W,U),
\]
çünkü $\Beklenti[U]=0$. Cauchy-Schwarz:
\[
  1 = [\Cov(W,U)]^2 \leq \Var(W)\cdot\Var(U) = \Var(W)\cdot n\mathcal{I}(\theta).
\]
Dolayısıyla $\Var(W)\geq 1/(n\mathcal{I}(\theta))$.
\end{proof}

\begin{tanim}[Verimlilik]\label{def:verimlilik}
Cramér-Rao alt sınırına ulaşan yansız tahmin ediciye \textbf{verimli tahmin edici} \emph{(efficient estimator)} denir.
\end{tanim}

\begin{teorem}[Üstel aile verimliliği]\label{thm:ustel_verimli}
Tek parametreli üstel ailede $f(x;\theta)=h(x)\exp[\eta(\theta)T(x)-B(\theta)]$; $T(\mathbf{X}) = \sum T(X_i)/n$ her $\theta$ için Cramér-Rao sınırına ulaşır (verimlidir).
\end{teorem}

\begin{ispatiskele}
Üstel ailede $\partial\log f/\partial\theta = \eta'(\theta)T(x) - B'(\theta)$; bu skor $T(x)$'in doğrusal fonksiyonu. Cramér-Rao'da eşitlik Cauchy-Schwarz'da eşitlikle, o da $W = aU+b$ doğrusallığıyla eşdeğer. Sonuç: $T$ verimlidir.
\end{ispatiskele}

%% ============================================================
\section{UMVUE: Düzgün Minimum Varyanslı Yansız Tahmin Edici}
\label{sec:umvue}
%% ============================================================

\begin{tanim}[UMVUE]\label{def:umvue}
$W^*$ yansız tahmin edici, her yansız $W$ ve her $\theta\in\Theta$ için $\Var_\theta(W^*)\leq\Var_\theta(W)$ sağlıyorsa \textbf{UMVUE} \emph{(uniformly minimum variance unbiased estimator)} adını alır.
\end{tanim}

\subsection{Rao-Blackwell Teoremi}

\begin{teorem}[Rao-Blackwell]\label{thm:rao_blackwell}
$W$ yansız tahmin edici ve $T$ yeterli istatistik olsun. $W^* = \Beklenti[W\mid T]$ tanımlayalım. O zaman:
\begin{enumerate}[(a)]
  \item $W^*$ yansız: $\Beklenti[W^*] = \theta$.
  \item $\Var(W^*)\leq\Var(W)$.
  \item Eşitlik yalnızca $W = W^*(T)$ (neredeyse kesin) olduğunda geçerli.
\end{enumerate}
\end{teorem}

\begin{proof}
(a) Kule özelliği: $\Beklenti[W^*] = \Beklenti[\Beklenti[W\mid T]] = \Beklenti[W] = \theta$.

(b) Varyans ayrışımı: $\Var(W) = \Beklenti[\Var(W\mid T)] + \Var(\Beklenti[W\mid T]) = \Beklenti[\Var(W\mid T)] + \Var(W^*) \geq \Var(W^*)$.
\end{proof}

\begin{kritiksorgu}
Rao-Blackwell sadece yeterli istatistiğe koşullanmak yeterli mi, UMVUE garanti eder mi? Hayır — $T$ tamamlayıcı olmalı. Tamamlayıcılık olmazsa $W^*$ yeterli fakat UMVUE değil olabilir.
\end{kritiksorgu}

\subsection{Lehmann-Scheffé Teoremi}

\begin{teorem}[Lehmann-Scheffé]\label{thm:lehmann_scheffe}
$T$ tamamlayıcı yeterli istatistik ve $g(T)$ yansız tahmin edici olsun. O zaman $g(T)$ UMVUE'dür ve yegânedir (n.k.).
\end{teorem}

\begin{proof}
Yegânlik: $g_1(T)$ ve $g_2(T)$ her ikisi de yansız ise $\Beklenti_\theta[g_1(T)-g_2(T)]=0$ her $\theta$ için. Tamamlayıcılıktan $g_1(T)=g_2(T)$ n.k.

UMVUE: $W$ başka herhangi yansız tahmin edici olsun. $W^* = \Beklenti[W\mid T] = g(T)$ (yeterlilikten, $g$ sabit). Rao-Blackwell'den $\Var(g(T)) \leq \Var(W)$.
\end{proof}

\begin{ornek}[Poisson UMVUE]\label{ex:poisson_umvue}
$X_1,\ldots,X_n\iid\Pois(\lambda)$. $T=\sum X_i \sim \Pois(n\lambda)$ tamamlayıcı yeterli. $\Beklenti[\bar{X}_n]=\lambda$ ve $\bar{X}_n = T/n$ bir $T$'nin fonksiyonu. Lehmann-Scheffé'den $\bar{X}_n$ UMVUE'dür.

Peki $\Prob(X_1=0)=e^{-\lambda}$'yı tahmin et. $g(T)=\bigl(\frac{n-1}{n}\bigr)^T$ yansız ve $T$'nin fonksiyonu, dolayısıyla UMVUE'dür (hesabı alıştırma).
\end{ornek}

%% ============================================================
\section{En Çok Olabilirlik Tahmini}
\label{sec:mle}
%% ============================================================

\begin{tanim}[En çok olabilirlik tahmin edicisi]\label{def:mle}
$\mathbf{x} = (x_1,\ldots,x_n)$ gözlemine karşılık \textbf{olabilirlik fonksiyonu}:
\[
  L(\theta;\mathbf{x}) = \prod_{i=1}^n f(x_i;\theta).
\]
$L(\hat{\theta};\mathbf{x}) = \sup_{\theta\in\Theta} L(\theta;\mathbf{x})$ sağlayan $\hat{\theta} = \hat{\theta}(\mathbf{x})$'ye \textbf{en çok olabilirlik tahmini} \emph{(maximum likelihood estimate, MLE)} denir; $\MLE{\theta}$ yazılır.
\end{tanim}

\begin{tanimgerekce}
Pratikte log-olabilirlik $\ell(\theta) = \log L(\theta) = \sum\log f(x_i;\theta)$ maksimize edilir. Log dönüşümü maksimumu değiştirmez ama türev hesabını basitleştirir.
\end{tanimgerekce}

\begin{ornek}[Normal MLE]\label{ex:normal_mle}
$X_1,\ldots,X_n\iid\Normal(\mu,\sigma^2)$, $\theta=(\mu,\sigma^2)$.
\[
  \ell(\mu,\sigma^2) = -\frac{n}{2}\log(2\pi\sigma^2) - \frac{1}{2\sigma^2}\sum(x_i-\mu)^2.
\]
$\partial\ell/\partial\mu = 0 \Rightarrow \MLE{\mu}=\bar{x}$.
$\partial\ell/\partial\sigma^2 = 0 \Rightarrow \MLE{\sigma^2} = \frac{1}{n}\sum(x_i-\bar{x})^2 = \hat{\sigma}^2$.

Not: $\MLE{\sigma^2}$ yanlıdır ($1/n$ böleni), ama MLE doğal olarak bu sonucu verir.
\end{ornek}

\subsection{MLE'nin Özellikleri}

\begin{teorem}[İnvaryant]\label{thm:mle_invaryant}
$\MLE{\theta}$ var ve $g:\Theta\to\mathbb{R}$ bire-bir ise $\widehat{g(\theta)} = g(\MLE{\theta})$.
\end{teorem}

\begin{proof}
$g$ bire-bir olduğundan $\eta=g(\theta)$, $\theta=g^{-1}(\eta)$. $L^*(\eta;\mathbf{x}) = L(g^{-1}(\eta);\mathbf{x})$'yi maksimize eden $\hat\eta = g(\MLE\theta)$.
\end{proof}

\begin{teorem}[Tutarlılık]\label{thm:mle_tutarlilik}
Regularity koşulları altında $\MLE{\theta} \pto \theta_0$ (gerçek parametre).
\end{teorem}

\begin{ispatiskele}
\textbf{Fikir:} $\frac{1}{n}\ell(\theta) = \frac{1}{n}\sum\log f(X_i;\theta) \asto \Beklenti_{\theta_0}[\log f(X;\theta)]$ (GBSK). Bu beklenti $\theta=\theta_0$'da maksimum alır (Gibbs eşitsizliği / KL ıraksama $\geq 0$). Dolayısıyla maksimum noktası $\theta_0$'a yakınsar.
\end{ispatiskele}

\begin{teorem}[Asimptotik normallik]\label{thm:mle_asimptotik}
Regularity koşulları altında:
\[
  \sqrt{n}\bigl(\MLE{\theta} - \theta\bigr) \dto \Normal\!\left(0,\frac{1}{\mathcal{I}(\theta)}\right).
\]
Yani MLE asimptotik olarak verimlidir.
\end{teorem}

\begin{ispatiskele}
\textbf{Fikir:} Skor denklemi $\ell'(\MLE\theta)=0$'ı $\theta_0$ çevresinde Taylor açılımı:
\[
  0 = \ell'(\theta_0) + (\MLE\theta-\theta_0)\ell''(\theta^*).
\]
$\frac{1}{\sqrt{n}}\ell'(\theta_0)\dto\Normal(0,\mathcal{I}(\theta))$ (MLT) ve $\frac{1}{n}\ell''(\theta^*)\asto -\mathcal{I}(\theta)$ (GBSK). Slütseva: $\sqrt{n}(\MLE\theta-\theta_0)\dto\Normal(0,1/\mathcal{I}(\theta))$.
\end{ispatiskele}

\begin{hataanalizi}
MLE her zaman var olmak zorunda değildir, benzersiz olmak zorunda değildir, tutarlı olmak zorunda değildir (regularity ihlalinde). Uniform$[0,\theta]$ için $\MLE\theta=X_{(n)}$ tutarlı; ama bu durumda sınır değer problemi olduğundan regularity bozulur.
\end{hataanalizi}

%% ============================================================
\section{Momentler Yöntemi}
\label{sec:momentler_yontemi}
%% ============================================================

\begin{tanim}[Momentler yöntemi]\label{def:momentler_yontemi}
$k$. popülasyon momenti $\mu_k'(\theta) = \Beklenti_\theta[X^k]$, $k$. örneklem momenti $m_k = \frac{1}{n}\sum X_i^k$ olsun. \textbf{Momentler yöntemi tahmin edicisi} \emph{(method of moments estimator)}, $\mu_k'(\hat{\theta}) = m_k$ denklem sisteminin çözümüdür.
\end{tanim}

\begin{ornek}[Gama momentler]\label{ex:gama_momentler}
$X\sim\Gama(\alpha,\beta)$ ise $\Beklenti[X]=\alpha\beta$, $\Var(X)=\alpha\beta^2$.
\[
  \hat\alpha = \frac{\bar{X}^2}{S^2}, \qquad \hat\beta = \frac{S^2}{\bar{X}}.
\]
\end{ornek}

\begin{teorem}[Momentler tahmin edicisinin tutarlılığı]
Regularity koşulları altında momentler yöntemi tahmin edicisi tutarlıdır.
\end{teorem}

\begin{proof}
GBSK'dan $m_k \asto \mu_k'(\theta_0)$. $g(\mu_1',\ldots,\mu_p')=\theta$ düzgün (sürekli türevlenebilir) ise Slütseva teoreminden $\hat\theta \asto \theta_0$.
\end{proof}

%% ============================================================
\section{Bayes Tahmini}
\label{sec:bayes_tahmin}
%% ============================================================

\begin{kesif}
MLE'de $\theta$ sabit ama bilinmeyen parametre idi. Bayes yaklaşımında $\theta$'yı da bir rasgele değişken olarak modelleriz: tahmin öncesinde $\theta$ hakkında bir \emph{önsel} (\emph{prior}) dağılımımız var, veriyi gördükten sonra \emph{sonsal} (\emph{posterior}) dağılıma geçiyoruz.
\end{kesif}

\begin{tanim}[Önsel ve sonsal dağılım]\label{def:onsel_sonsal}
$\pi(\theta)$ önsel yoğunluk ve $f(\mathbf{x}\mid\theta)$ olabilirlik olsun. \textbf{Sonsal dağılım}:
\[
  \pi(\theta\mid\mathbf{x}) = \frac{f(\mathbf{x}\mid\theta)\,\pi(\theta)}{\int f(\mathbf{x}\mid\theta')\pi(\theta')\,d\theta'} \propto f(\mathbf{x}\mid\theta)\,\pi(\theta).
\]
\end{tanim}

\begin{tanim}[Konjugat önsel]\label{def:konjugat_onsel}
$\pi(\theta)$ \textbf{konjugat önsel} \emph{(conjugate prior)} ise, sonsal dağılım önsel ile aynı aile içinde kalır.
\end{tanim}

\subsection{Temel Konjugat Çiftler}

\begin{teorem}[Beta-Binom konjugasyonu]\label{thm:beta_binom}
$X\mid p \sim \Binom(n,p)$ ve $p\sim\Beta(\alpha,\beta)$ ise:
\[
  p\mid\mathbf{x} \sim \Beta\!\left(\alpha + \sum x_i,\; \beta + n - \sum x_i\right).
\]
Bayes tahmini (karesel kayıp altında):
\[
  \hat{p}_{\text{Bayes}} = \Beklenti[p\mid\mathbf{x}] = \frac{\alpha+\sum x_i}{\alpha+\beta+n}.
\]
\end{teorem}

\begin{proof}
$f(x\mid p)\propto p^x(1-p)^{n-x}$ ve $\pi(p)\propto p^{\alpha-1}(1-p)^{\beta-1}$. Sonsal $\propto p^{\alpha+\sum x_i-1}(1-p)^{\beta+n-\sum x_i-1}$, bu $\Beta(\alpha+\sum x_i, \beta+n-\sum x_i)$.
\end{proof}

\begin{teorem}[Normal-Normal konjugasyonu]\label{thm:normal_normal}
$X_1,\ldots,X_n\mid\mu \iid \Normal(\mu,\sigma^2)$ (bilinen $\sigma^2$) ve $\mu\sim\Normal(\mu_0,\tau^2)$ ise:
\[
  \mu\mid\mathbf{x} \sim \Normal\!\left(\mu_n, \sigma_n^2\right)
\]
ile
\[
  \frac{1}{\sigma_n^2} = \frac{n}{\sigma^2} + \frac{1}{\tau^2}, \qquad
  \mu_n = \sigma_n^2\!\left(\frac{n\bar{x}}{\sigma^2} + \frac{\mu_0}{\tau^2}\right).
\]
\end{teorem}

\begin{proof}
$\log\pi(\mu\mid\mathbf{x}) \propto -\frac{n}{2\sigma^2}\sum(x_i-\mu)^2 - \frac{1}{2\tau^2}(\mu-\mu_0)^2$. $\mu$'ya göre karesel; kareyi tamamla: $-\frac{1}{2\sigma_n^2}(\mu-\mu_n)^2$.
\end{proof}

\begin{teorem}[Gama-Poisson konjugasyonu]\label{thm:gama_poisson}
$X_1,\ldots,X_n\mid\lambda\iid\Pois(\lambda)$ ve $\lambda\sim\Gama(\alpha,\beta)$ ise:
\[
  \lambda\mid\mathbf{x} \sim \Gama\!\left(\alpha+\sum x_i,\; \frac{\beta}{n\beta+1}\right).
\]
\end{teorem}

\begin{proof}
$f(\mathbf{x}\mid\lambda)\propto\lambda^{\sum x_i}e^{-n\lambda}$ ve $\pi(\lambda)\propto\lambda^{\alpha-1}e^{-\lambda/\beta}$. Sonsal $\propto\lambda^{\alpha+\sum x_i-1}e^{-\lambda(n+1/\beta)}$, bu $\Gama$.
\end{proof}

\subsection{Güvenilir Aralık}

\begin{tanim}[Güvenilir aralık]\label{def:guvenilir_aralik}
$(1-\alpha)$ kapsama olasılıklı \textbf{güvenilir aralık} \emph{(credible interval)} $[a,b]$:
\[
  \Prob(\theta\in[a,b]\mid\mathbf{x}) = \int_a^b \pi(\theta\mid\mathbf{x})\,d\theta = 1-\alpha.
\]
\textbf{HPD (highest posterior density)} aralığı: en kısa güvenilir aralık; sonsal yoğunluğun $c$ eşik değerini aşan kümesi.
\end{tanim}

\begin{ornek}[Beta sonsal güvenilir aralık]\label{ex:beta_credible}
$n=20$ deneme, $\sum x_i=14$ başarı; $p\sim\Beta(1,1)$ (düzgün önsel). Sonsal $p\mid\mathbf{x}\sim\Beta(15,7)$. \%95 HPD aralığı Beta sonsal dağılımının \%2.5--\%97.5 kantilleri: sayısal olarak $[0{,}515, 0{,}913]$.
\end{ornek}

\begin{kendinleifade}
Beta-Binom Bayes tahmin edicisini MLE olan $\bar{X}_n$ ile karşılaştır. $\alpha+\beta$ "hayali örneklem büyüklüğü" yorumunu açıkla. $n\to\infty$ limitinde ne olur?
\end{kendinleifade}

%% ============================================================
\section{Minimax Tahmin ve Kabul Edilebilirlik}
\label{sec:minimax}
%% ============================================================

\begin{tanim}[Minimax tahmin edici]\label{def:minimax}
Tüm $\delta$ karar kuralları arasında $\delta^*$ \textbf{minimax}:
\[
  \sup_\theta R(\theta,\delta^*) = \inf_\delta\sup_\theta R(\theta,\delta).
\]
Worst-case riski minimize eder; $\theta$ hakkında bilgisiz (en kötümser) karar.
\end{tanim}

\begin{tanim}[Kabul edilemezlik]\label{def:kabul_edilemez}
$\delta$ karar kuralı \textbf{kabul edilemez} \emph{(inadmissible)} eğer başka bir $\delta'$ ile $R(\theta,\delta')\leq R(\theta,\delta)$ tüm $\theta$'da ve en az bir $\theta^*$'da $R(\theta^*,\delta')<R(\theta^*,\delta)$. Aksi hâlde \textbf{kabul edilebilir} \emph{(admissible)}.
\end{tanim}

\begin{teorem}[Stein paradoksu]\label{thm:stein}
$\mathbf{X}\sim\Normal_p(\boldsymbol\mu,I_p)$, $p\geq3$, karesel kayıp. O zaman $\bar{\mathbf{X}}$ (yansız MLE) kabul edilemezdir: \textbf{James-Stein tahmini}
\[
  \hat{\boldsymbol\mu}^{JS} = \left(1 - \frac{p-2}{\|\mathbf{X}\|^2}\right)\mathbf{X}
\]
her $\boldsymbol\mu$ için daha küçük OKH'a sahip.
\end{teorem}

\begin{ispatiskele}
$\Beklenti[\|\hat{\boldsymbol\mu}^{JS}-\boldsymbol\mu\|^2]=p-(p-2)^2\Beklenti[1/\|\mathbf{X}\|^2]<p=\mathrm{OKH}(\mathbf{X})$. Detay: Efron-Morris (1977). Paradoks: $p=1,2$'de yansız tahmin kabul edilebilir; $p\geq3$'te değil.
\end{ispatiskele}

%% ============================================================
\section{Skor Fonksiyonu ve Fisher Bilgisi — İleri}
\label{sec:fisher_ileri}
%% ============================================================

\begin{teorem}[Fisher bilgisi birleşimi]\label{thm:fisher_birlesen}
$X_1,\ldots,X_n\iid f(x;\theta)$. Birleşik Fisher bilgisi $\mathcal{I}_n(\theta)=n\mathcal{I}(\theta)$.
\end{teorem}

\begin{proof}
Bağımsız gözlemlerden skor toplamı: $\ell_n'(\theta)=\sum_{i=1}^n\partial\log f(X_i;\theta)/\partial\theta$. $\Var(\ell_n')=n\Var(\partial\log f/\partial\theta)=n\mathcal{I}(\theta)$.
\end{proof}

\begin{teorem}[Çok parametreli Cramér-Rao]\label{thm:cr_cok}
$\theta\in\mathbb{R}^p$, Fisher bilgi matrisi $\mathbf{I}(\theta)=[I_{jk}]$ ile $I_{jk}=\Beklenti[-\partial^2\log f/\partial\theta_j\partial\theta_k]$. Her yansız $W$:
\[
  \Cov_\theta(W) \succeq \mathbf{I}(\theta)^{-1}
\]
(matris sıralaması anlamında).
\end{teorem}

\begin{proof}[Proof (taslak)]
Tek boyutlu Cauchy-Schwarz'ın matris versiyonu. $\Cov(W,\ell'(\theta))=I_p$ (normal denklem); Cauchy-Schwarz: $\Cov(W)\cdot\Cov(\ell')\succeq[\Cov(W,\ell')]^2$. $\Cov(\ell')=\mathbf{I}(\theta)$.
\end{proof}

\begin{ornek}[Normal Fisher bilgi matrisi]\label{ex:normal_fisher}
$X\sim\Normal(\mu,\sigma^2)$, $\theta=(\mu,\sigma^2)$. $\log f = -\frac12\log\sigma^2 - \frac{(x-\mu)^2}{2\sigma^2}+$sabit. Fisher matrisi:
\[
  \mathbf{I}(\theta) = \begin{pmatrix}1/\sigma^2 & 0 \\ 0 & 1/(2\sigma^4)\end{pmatrix}.
\]
Cramér-Rao: $\Var(\hat\mu)\geq\sigma^2/n$, $\Var(\hat\sigma^2)\geq2\sigma^4/n$ — yansız MLE her ikisini sağlar.
\end{ornek}

%% ============================================================

\begin{mikroozet}
\begin{itemize}
  \item Yansızlık + minimum varyans = UMVUE; Cramér-Rao alt sınırı ulaşılabilir sınır koyar.
  \item Yeterli + tamamlayıcı $\Rightarrow$ Rao-Blackwell + Lehmann-Scheffé $\Rightarrow$ UMVUE.
  \item MLE: kolay hesap, invaryant, tutarlı, asimptotik verimli; ama yanlı olabilir.
  \item Bayes: önsel + olabilirlik $\to$ sonsal; konjugat önsel analitik çözüm verir; güvenilir aralık Bayesçi belirsizlik ölçüsü.
  \item Minimax: worst-case optimal; Stein paradoksu $p\geq3$'te yansız MLE'nin kabul edilemez olduğunu gösterir.
\end{itemize}
\end{mikroozet}

\begin{ozyansima}
Hangi durumlarda MLE, UMVUE ile çakışır? Üstel ailede bu her zaman geçerli mi? Momentler yöntemi neden büyük örneklemlerde MLE'ye yakınsar?
\end{ozyansima}

\begin{kopru}
\textbf{Nereden geldik:} Bölüm~\ref{chp:istatistik_temelleri}'deki yeterli istatistik, Rao-Blackwell argümanının temelidir.

\textbf{Nereye gidiyoruz:} Bölüm~\ref{chp:guven_araliklari}'da nokta tahminini aralık tahminine dönüştüreceğiz; buradaki Fisher bilgisi ve asimptotik normallik, Wald güven aralığı için doğrudan kullanılacak.
\end{kopru}

%% ============================================================
%% ALISTIRMALAR
%% ============================================================

\cozumlualistirmalar

\begin{alistirma}[Al.12.1]{A}{Fisher bilgisi — Bernoulli}
$X\sim\mathrm{Ber}(p)$ için Fisher bilgisini hesapla ve Cramér-Rao alt sınırını yaz. $\bar{X}_n$'nin bu sınırı sağladığını doğrula.
\end{alistirma}
\alcozum{%
$\log f(x;p) = x\log p+(1-x)\log(1-p)$; $\partial^2/\partial p^2 = -x/p^2-(1-x)/(1-p)^2$. Beklenti $-1/(p(1-p))$; dolayısıyla $\mathcal{I}(p)=1/(p(1-p))$. Alt sınır $p(1-p)/n = \Var(\bar{X}_n)$. Eşitlik sağlanır.%
}

\begin{alistirma}[Al.12.2]{B}{Eksponansiyel UMVUE}
$X_1,\ldots,X_n\iid\Exp(\lambda)$ ($\Beklenti[X]=1/\lambda$). $T=\sum X_i$ tamamlayıcı yeterli. $g(T)=n/T$'nin yansız olduğunu göster ve UMVUE olduğunu belirt. \textit{(İpucu: $T\sim\Gama(n,1/\lambda)$.)}
\end{alistirma}
\alcozum{%
$\Beklenti[n/T]=n\cdot\Beklenti[1/T]$. $T\sim\Gama(n,1/\lambda)$; $\Beklenti[1/T]=(n-1)^{-1}\cdot\lambda$ ($n>1$ için, negatif moment formülü). Bekle: $\Beklenti[T^{-1}]=\lambda/(n-1)$. O zaman $\Beklenti[n/T]=n\lambda/(n-1)\neq\lambda$. Düzeltme: yansız olan $\hat\lambda=(n-1)/T$. Bu $T$'nin fonksiyonu ve yansız, dolayısıyla UMVUE.%
}

\alistirmalar

\begin{alistirma}[Al.12.3]{A}{Momentler — Beta}
$X\sim\Beta(\alpha,\beta)$ için momentler yöntemi tahmin edicilerini $\bar{X}_n$ ve $S^2$ cinsinden yaz.
\end{alistirma}
\alcozum{%
$\mu=\alpha/(\alpha+\beta)$, $\sigma^2=\alpha\beta/[(\alpha+\beta)^2(\alpha+\beta+1)]$. Çöz: $\hat\alpha=\bar{X}[\bar{X}(1-\bar{X})/S^2-1]$, $\hat\beta=(1-\bar{X})[\bar{X}(1-\bar{X})/S^2-1]$.%
}

\begin{alistirma}[Al.12.4]{B}{Poisson UMVUE — $e^{-\lambda}$}
$X_1,\ldots,X_n\iid\Pois(\lambda)$. $\Prob(X=0)=e^{-\lambda}$'nın UMVUE'sünün $g(T)=\left(\frac{n-1}{n}\right)^T$ ($T=\sum X_i$) olduğunu göster. \textit{(İpucu: $\Beklenti[g(T)]$ hesapla.)}
\end{alistirma}
\alcozum{%
$T\sim\Pois(n\lambda)$; $\Beklenti[(1-1/n)^T]=\sum_{t=0}^\infty(1-1/n)^t e^{-n\lambda}(n\lambda)^t/t! = e^{-n\lambda}\exp(n\lambda(1-1/n))=e^{-n\lambda}\cdot e^{(n-1)\lambda}=e^{-\lambda}$. Dolayısıyla yansız; $T$ tamamlayıcı yeterli ve $g(T)$ fonksiyonu, dolayısıyla UMVUE.%
}

\begin{alistirma}[Al.12.5]{B}{Normal-Normal Bayes}
$X_1,\ldots,X_n\mid\mu\iid\Normal(\mu,4)$ ve $\mu\sim\Normal(0,1)$. $n=9$, $\bar{x}=2$ gözlemi için sonsal dağılımı ve Bayes tahmin edicisini hesapla. $n\to\infty$ limitini yorumla.
\end{alistirma}
\alcozum{%
$1/\sigma_n^2=9/4+1/1=13/4$; $\sigma_n^2=4/13$. $\mu_n=(4/13)(9\cdot2/4+0/1)=(4/13)(9/2)=18/13\approx1{,}38$. $n\to\infty$: $\sigma_n^2\to0$, $\mu_n\to\bar{X}_n$ — MLE baskın gelir.%
}

\begin{alistirma}[Al.12.6]{C}{Cramér-Rao sınırına ulaşma koşulu}
Bir yansız tahmin edicinin Cramér-Rao sınırına ulaşmasının, skor fonksiyonunun $W$'nin doğrusal fonksiyonu olmasıyla eşdeğer olduğunu göster.
\end{alistirma}
\alcozum{%
Cauchy-Schwarz'da $[\Cov(W,U)]^2=\Var(W)\Var(U)$ eşitliği ancak $W=aU+b$ (neredeyse kesin) ile sağlanır. Dolayısıyla $\partial\log f(\mathbf{x};\theta)/\partial\theta = c(\theta)(W-\theta)$ olmalı; yani $f(\mathbf{x};\theta)=h(\mathbf{x})\exp[A(\theta)W(\mathbf{x})+B(\theta)]$ üstel aile biçiminde.%
}

\begin{alistirma}[Al.12.7]{C}{Düzgün dağılımda MLE ve yansızlık}
$X_1,\ldots,X_n\iid\Uniform[0,\theta]$. MLE'nin $X_{(n)}$ olduğunu göster, yanlılığını hesapla ve yansız tahmin ediciyi türet.
\end{alistirma}
\alcozum{%
$L(\theta)=\theta^{-n}$ ($\theta\geq X_{(n)}$) azalıyor; dolayısıyla $\MLE\theta=X_{(n)}$. $F_{X_{(n)}}(t)=(t/\theta)^n$; $\Beklenti[X_{(n)}]=n\theta/(n+1)$. Yanlılık $-\theta/(n+1)<0$. Yansız: $W=\frac{n+1}{n}X_{(n)}$.%
}

\begin{alistirma}[Al.12.8]{C}{OKH karşılaştırması}
$X_1,\ldots,X_n\iid\Normal(\mu,\sigma^2)$. $\hat\sigma^2=\frac{1}{n}\sum(X_i-\bar{X})^2$ (yanlı) ve $S^2=\frac{1}{n-1}\sum(X_i-\bar{X})^2$ (yansız) OKH'larını karşılaştır. Hangi tahmin edici her $n$ için daha küçük OKH'a sahip?
\end{alistirma}
\alcozum{%
$\hat\sigma^2=\frac{n-1}{n}S^2$. $S^2$'nin dağılımı $\sigma^2\chi^2_{n-1}/(n-1)$; $\Var(S^2)=2\sigma^4/(n-1)$. $\mathrm{OKH}(\hat\sigma^2)=\Var(\hat\sigma^2)+[\mathrm{Yanlılık}]^2=\frac{(n-1)^2}{n^2}\frac{2\sigma^4}{n-1}+\frac{\sigma^4}{n^2}=\frac{(2n-1)\sigma^4}{n^2}$. $\mathrm{OKH}(S^2)=\Var(S^2)=\frac{2\sigma^4}{n-1}$. Karşılaştır: $\frac{2n-1}{n^2}$ vs $\frac{2}{n-1}$. İlki her $n\geq2$ için daha küçük. Dolayısıyla yanlı $\hat\sigma^2$ daha küçük OKH'a sahip.%
}

\begin{alistirma}[Al.12.9]{B}{Çok parametreli Fisher bilgisi}
$X\sim\Normal(\mu,\sigma^2)$ (her ikisi bilinmiyor). (a) Fisher bilgi matrisini hesapla. (b) $\mu$ için Cramér-Rao alt sınırını yaz. (c) $\sigma^2$ için yansız tahmin edici Cramér-Rao sınırına ulaşıyor mu?
\end{alistirma}
\alcozum{%
(a) $\mathbf{I}(\mu,\sigma^2)=\begin{pmatrix}1/\sigma^2&0\\0&1/(2\sigma^4)\end{pmatrix}$ (köşegen: $\mu,\sigma^2$ bilgi ayrışır). (b) $\Var(\hat\mu)\geq\sigma^2/n$. (c) $S^2$'nin varyansı $2\sigma^4/(n-1)$; Cramér-Rao sınırı $2\sigma^4/n$. $2/(n-1)>2/n$: sınıra ulaşılmaz. Sadece Normal$(μ,σ^2$ bilinmez) için MLE $\hat\sigma^2_{\mathrm{MLE}}=\hat\sigma^2$ (yanlı) daha küçük varyansa sahip ama yanlı.%
}

\begin{alistirma}[Al.12.10]{C}{James-Stein tahmini}
$p=3$, $\mathbf{X}\sim\Normal_3(\boldsymbol\mu,I_3)$, $\mathbf{x}=(2,1,0)$ gözlendi. (a) $\|\mathbf{x}\|^2$'yi hesapla. (b) James-Stein tahmini $\hat{\boldsymbol\mu}^{JS}$'yi hesapla. (c) OKH farkı ne kadar?
\end{alistirma}
\alcozum{%
(a) $\|\mathbf{x}\|^2=4+1+0=5$. (b) $c=1-(p-2)/\|\mathbf{x}\|^2=1-1/5=0.8$. $\hat{\boldsymbol\mu}^{JS}=0.8\cdot(2,1,0)=(1.6,0.8,0)$. (c) OKH$(\mathbf{X})=p=3$. OKH$(JS)=p-(p-2)^2E[1/\|\mathbf{X}\|^2]<3$. OKH kazancı: $(p-2)^2/\|\mathbf{x}\|^2\approx1/5=0.2$ mertebesinde (bu sadece bireysel veri için tahmin; gerçek kazanç beklenti üzerinden hesaplanır).%
}

%% ============================================================
%% BÖLÜM 13 — Güven Aralıkları
%% Olasılık Teorisi ve Matematiksel İstatistik
%% ============================================================
\chapter{Güven Aralıkları}
\label{chp:guven_araliklari}

\bolumalinti{%
  Statistics is the grammar of science.
}{Karl Pearson}

\begin{kazanimlar}
Bu bölümün sonunda okuyucu:
\begin{itemize}
  \item Güven aralığı tanımını, kapsama olasılığını ve doğru yorumunu açıklar.
  \item Pivot yöntemini kullanarak güven aralığı türetir.
  \item Normal, $t$, $\chi^2$ ve $F$ dağılımlarına dayalı güven aralıkları kurar.
  \item Büyük örneklem güven aralıklarını (Wald, oran, varyans stabilizasyonu) uygular.
  \item Aralık uzunluğu ile kapsama olasılığı arasındaki dengeyi analiz eder.
  \item Homojen olmayan varyanslar için Welch güven aralığını kullanır.
\end{itemize}
\end{kazanimlar}

\begin{motivasyon}
Nokta tahmini tek bir sayı verir; ama bu sayının ne kadar güvenilir olduğunu söylemez. Güven aralığı, belirsizliği nicelleştiren bir aralık tahminidir. Bir anket şirketi "\%43 ± \%2" derken aslında güven aralığı bildirmektedir.

\medskip
\noindent\textbf{Köprü:} Bölüm~\ref{chp:nokta_tahmin}'deki MLE'nin asimptotik normalliği, büyük örneklem güven aralıklarının temelidir. Bölüm~\ref{chp:dagilim_aileleri}'deki $t$, $\chi^2$, $F$ dağılımları, küçük örneklem güven aralıklarında kullanılacak.
\end{motivasyon}

%% ============================================================
\section{Güven Aralığı Kavramı}
\label{sec:ga_kavram}
%% ============================================================

\begin{tanim}[Güven aralığı]\label{def:guven_araligi}
$X_1,\ldots,X_n\sim F_\theta$ ve $L=L(\mathbf{X})$, $U=U(\mathbf{X})$ istatistikleri olsun ($L\leq U$). $(1-\alpha)$ \textbf{kapsama olasılıklı güven aralığı} \emph{(confidence interval, CI)}:
\[
  \Prob_\theta(L \leq \theta \leq U) \geq 1-\alpha \quad \text{her } \theta\in\Theta \text{ için.}
\]
$1-\alpha$ \textbf{güven düzeyi} \emph{(confidence level)}, $\alpha$ \textbf{anlamlılık düzeyi} \emph{(significance level)} adını alır.
\end{tanim}

\begin{tanimgerekce}
Sık yapılan yorum hatası: "\%95 olasılıkla $\theta$, $[L,U]$ içinde" — bu yanlış. $\theta$ sabit; $L$ ve $U$ rasgele. Doğru yorum: Yöntemi tekrar tekrar uygulandığında aralıkların \%95'i gerçek $\theta$'yı kapsar.
\end{tanimgerekce}

\begin{hataanalizi}
"Bu özel aralık $[1.2, 2.8]$'in \%95 olasılıkla $\mu$'yü kapsadığı" ifadesi hatalıdır. Gözlemlendikten sonra aralık sabittir; $\mu$ ya içindedir ya değildir. "Bu yöntemle üretilen aralıkların \%95'i $\mu$'yü kapsar" doğru ifadedir.
\end{hataanalizi}

\subsection{Pivot Yöntemi}

\begin{tanim}[Pivot]\label{def:pivot}
$Q = Q(X_1,\ldots,X_n;\theta)$ istatistiği, dağılımı $\theta$'dan bağımsız ise \textbf{pivot} \emph{(pivotal quantity)} adını alır.
\end{tanim}

\begin{teorem}[Pivot güven aralığı]\label{thm:pivot_ga}
$Q$ pivot ve $\Prob(a\leq Q\leq b)=1-\alpha$ ise, $a\leq Q\leq b$ eşitsizliğini $\theta$ için çözmek $(1-\alpha)$ güven aralığı verir.
\end{teorem}

%% ============================================================
\section{Normal Popülasyon: Bilinen Varyans}
\label{sec:normal_bilinen}
%% ============================================================

\begin{teorem}[$z$-güven aralığı]\label{thm:z_ga}
$X_1,\ldots,X_n\iid\Normal(\mu,\sigma^2)$, $\sigma^2$ bilinen. O zaman
\[
  Z = \frac{\overline{X}_n-\mu}{\sigma/\sqrt{n}} \sim \Normal(0,1)
\]
bir pivot. $(1-\alpha)$ güven aralığı:
\[
  \left(\overline{X}_n - z_{\alpha/2}\frac{\sigma}{\sqrt{n}},\; \overline{X}_n + z_{\alpha/2}\frac{\sigma}{\sqrt{n}}\right),
\]
$z_{\alpha/2} = \Phi^{-1}(1-\alpha/2)$ ile.
\end{teorem}

\begin{proof}
$\Prob(-z_{\alpha/2}\leq Z\leq z_{\alpha/2})=1-\alpha$. İçi çözersek $\overline{X}_n - z_{\alpha/2}\sigma/\sqrt{n}\leq\mu\leq\overline{X}_n+z_{\alpha/2}\sigma/\sqrt{n}$.
\end{proof}

\begin{ornek}[$z$-aralığı]\label{ex:z_aralik}
$n=25$, $\bar{x}=12.4$, $\sigma=2$ (bilinen). $\alpha=0.05$; $z_{0.025}=1.960$. GA: $12.4\pm1.960\cdot(2/5) = 12.4\pm0.784 = [11.616, 13.184]$.
\end{ornek}

%% ============================================================
\section{Normal Popülasyon: Bilinmeyen Varyans}
\label{sec:normal_bilinmeyen}
%% ============================================================

\begin{teorem}[$t$-güven aralığı]\label{thm:t_ga}
$X_1,\ldots,X_n\iid\Normal(\mu,\sigma^2)$, $\sigma^2$ bilinmiyor. O zaman
\[
  T = \frac{\overline{X}_n-\mu}{S/\sqrt{n}} \sim t_{n-1}
\]
bir pivot. $(1-\alpha)$ güven aralığı:
\[
  \left(\overline{X}_n - t_{n-1,\alpha/2}\frac{S}{\sqrt{n}},\; \overline{X}_n + t_{n-1,\alpha/2}\frac{S}{\sqrt{n}}\right).
\]
\end{teorem}

\begin{proof}
Bölüm~\ref{chp:dagilim_aileleri}'den: Normal popülasyonda $\overline{X}_n\perp S^2$, $\overline{X}_n\sim\Normal(\mu,\sigma^2/n)$, $(n-1)S^2/\sigma^2\sim\chi^2_{n-1}$. Dolayısıyla $T = Z/\sqrt{\chi^2_{n-1}/(n-1)}\sim t_{n-1}$.
\end{proof}

\begin{ornek}[$t$-aralığı]\label{ex:t_aralik}
$n=16$, $\bar{x}=5.2$, $s=1.8$, $\sigma^2$ bilinmiyor, $\alpha=0.05$. $t_{15,0.025}=2.131$. GA: $5.2\pm2.131\cdot(1.8/4) = 5.2\pm0.959 = [4.241, 6.159]$.
\end{ornek}

\begin{kritiksorgu}
$n$ büyüdükçe $t$-aralığı $z$-aralığına yaklaşır ($t_{n-1}\to\Normal(0,1)$). Pratikte $n\geq30$ için fark ihmal edilebilir; ama normallik varsayımı hâlâ önemlidir.
\end{kritiksorgu}

%% ============================================================
\section{Varyans için Güven Aralığı}
\label{sec:varyans_ga}
%% ============================================================

\begin{teorem}[$\chi^2$-güven aralığı]\label{thm:chi2_ga}
$X_1,\ldots,X_n\iid\Normal(\mu,\sigma^2)$. O zaman
\[
  Q = \frac{(n-1)S^2}{\sigma^2} \sim \chi^2_{n-1}
\]
bir pivot. $(1-\alpha)$ güven aralığı:
\[
  \left(\frac{(n-1)S^2}{\chi^2_{n-1,\alpha/2}},\; \frac{(n-1)S^2}{\chi^2_{n-1,1-\alpha/2}}\right).
\]
\end{teorem}

\begin{proof}
$\Prob(\chi^2_{n-1,1-\alpha/2}\leq Q\leq\chi^2_{n-1,\alpha/2})=1-\alpha$; $\sigma^2$ için çöz.
\end{proof}

\begin{hataanalizi}
$\chi^2$ aralığı normallik varsayımına son derece duyarlıdır. Dağılım normalden sapıyorsa bootstrap yöntemi tercih edilir.
\end{hataanalizi}

%% ============================================================
\section{İki Örneklem Karşılaştırması}
\label{sec:iki_orneklem}
%% ============================================================

\subsection{Eşit Varyanslı İki Ortalama}

\begin{teorem}[İki örneklem $t$-aralığı]\label{thm:iki_t}
$X_1,\ldots,X_m\iid\Normal(\mu_1,\sigma^2)$ ve $Y_1,\ldots,Y_n\iid\Normal(\mu_2,\sigma^2)$ bağımsız (ortak $\sigma^2$). Havuzlanmış varyans tahmin edicisi:
\[
  S_p^2 = \frac{(m-1)S_X^2+(n-1)S_Y^2}{m+n-2}.
\]
$\mu_1-\mu_2$ için $(1-\alpha)$ güven aralığı:
\[
  (\bar{X}-\bar{Y}) \pm t_{m+n-2,\alpha/2}\,S_p\!\sqrt{\frac{1}{m}+\frac{1}{n}}.
\]
\end{teorem}

\begin{proof}
$\bar{X}-\bar{Y}\sim\Normal(\mu_1-\mu_2,\sigma^2(1/m+1/n))$ ve $(m+n-2)S_p^2/\sigma^2\sim\chi^2_{m+n-2}$ bağımsız olduğundan pivot $t_{m+n-2}$'dir.
\end{proof}

\subsection{Welch Güven Aralığı (Eşit Olmayan Varyanslarda)}

\begin{teorem}[Welch $t$-aralığı]\label{thm:welch}
$\sigma_1^2\neq\sigma_2^2$ (bilinmiyor). Satterthwaite serbestlik derecesi:
\[
  \nu = \frac{(S_X^2/m + S_Y^2/n)^2}{\dfrac{(S_X^2/m)^2}{m-1}+\dfrac{(S_Y^2/n)^2}{n-1}}.
\]
$\mu_1-\mu_2$ için yaklaşık güven aralığı:
\[
  (\bar{X}-\bar{Y}) \pm t_{\nu,\alpha/2}\!\sqrt{\frac{S_X^2}{m}+\frac{S_Y^2}{n}}.
\]
\end{teorem}

\begin{ispatiskele}
Welch-Satterthwaite yaklaşımı, $S_X^2/m+S_Y^2/n$'yi $a\chi^2_\nu$ ile eşleştirir; $a$ ve $\nu$ moment eşleştirmesiyle belirlenir. Yaklaşık (asimptotik değil); normallik gerekir.
\end{ispatiskele}

%% ============================================================
\section{Büyük Örneklem Güven Aralıkları}
\label{sec:buyuk_orneklem}
%% ============================================================

\begin{teorem}[Wald güven aralığı]\label{thm:wald_ga}
Regularity koşulları altında $\sqrt{n}(\MLE\theta-\theta)\dto\Normal(0,1/\mathcal{I}(\theta))$. Büyük örneklem $(1-\alpha)$ güven aralığı:
\[
  \MLE\theta \pm \frac{z_{\alpha/2}}{\sqrt{n\,\hat{\mathcal{I}}(\MLE\theta)}},
\]
$\hat{\mathcal{I}}$ gözlemlenen Fisher bilgisi.
\end{teorem}

\begin{ornek}[Bernoulli oranı için GA]\label{ex:oran_ga}
$X_1,\ldots,X_n\iid\mathrm{Ber}(p)$. $\MLE{p}=\bar{X}$, $\mathcal{I}(p)=1/(p(1-p))$. Wald aralığı:
\[
  \bar{X} \pm z_{\alpha/2}\sqrt{\frac{\bar{X}(1-\bar{X})}{n}}.
\]
\end{ornek}

\begin{hataanalizi}
Wald oranı güven aralığı $p\approx0$ veya $p\approx1$'de ya da küçük $n$'de kötü performans gösterir. Wilson aralığı (skorlu GA) daha sağlamdır:
\[
  \frac{\bar{X}+z_{\alpha/2}^2/(2n)}{1+z_{\alpha/2}^2/n} \pm \frac{z_{\alpha/2}\sqrt{\bar{X}(1-\bar{X})/n + z_{\alpha/2}^2/(4n^2)}}{1+z_{\alpha/2}^2/n}.
\]
\end{hataanalizi}

\subsection{Varyans Stabilizasyon Dönüşümü}

\begin{teorem}[Delta yöntemi ile GA]\label{thm:delta_ga}
$g$ türevlenebilir ve $g'(\theta)\neq0$. Delta yönteminden $\sqrt{n}(g(\MLE\theta)-g(\theta))\dto\Normal(0,[g'(\theta)]^2/\mathcal{I}(\theta))$. $g$ seçimi ile varyans sabitlenirse, GA sabit genişlikte olur.
\end{teorem}

\begin{ornek}[Poisson varyans stabilizasyonu]
$\bar{X}\sim\Pois(\lambda)$: $\Var(\bar{X})\approx\lambda/n$. $g(\lambda)=\sqrt{\lambda}$ ile delta yöntemi: $\sqrt{n}(\sqrt{\bar{X}}-\sqrt{\lambda})\dto\Normal(0,1/4)$. GA:
\[
  \sqrt{\bar{X}} \pm \frac{z_{\alpha/2}}{2\sqrt{n}} \quad\Rightarrow\quad
  \lambda: \left(\sqrt{\bar{X}} - \frac{z_{\alpha/2}}{2\sqrt{n}}\right)^{\!2} \leq \lambda \leq \left(\sqrt{\bar{X}} + \frac{z_{\alpha/2}}{2\sqrt{n}}\right)^{\!2}.
\]
\end{ornek}

%% ============================================================
\section{Güven Aralığı Uzunluğu ve Örneklem Büyüklüğü}
\label{sec:ga_uzunluk}
%% ============================================================

\begin{teorem}[Örneklem büyüklüğü belirleme]\label{thm:orneklem_buyuklugu}
$z$-aralığı ile $\mu$'yü $\pm d$ hassasiyetle tahmin etmek için gereken minimum örneklem büyüklüğü:
\[
  n \geq \left\lceil \frac{z_{\alpha/2}^2\,\sigma^2}{d^2} \right\rceil.
\]
$\sigma$ bilinmiyorsa pilot çalışma veya muhafazakâr $\sigma$ tahmini kullanılır.
\end{teorem}

\begin{proof}
GA yarı-genişliği $z_{\alpha/2}\sigma/\sqrt{n}\leq d$ eşitsizliğini $n$ için çöz.
\end{proof}

\begin{ornek}[Seçim anketi]
Bir oy oranını $\pm0.03$ hassasiyetle \%95 güvenle tahmin et. Muhafazakâr $p=0.5$ ile $\sigma^2=p(1-p)=0.25$. $n\geq\lceil 1.96^2\cdot0.25/0.03^2\rceil = \lceil 1067.1\rceil = 1068$.
\end{ornek}

\begin{kendinleifade}
Güven düzeyini $\alpha=0.05$'ten $\alpha=0.01$'e çekince aralık nasıl değişir? Örneklem büyüklüğünü 4 katına çıkarmak aralık genişliğini kaçta bire indirir?
\kendinleifadecevap{Güven düzeyi $\alpha=0.05\to0.01$: $z_{0.025}=1.96\to z_{0.005}=2.576$,
  aralık genişler (daha güvenilir = daha geniş).
  Örneklem $n\to4n$: Genişlik $\propto\sigma/\sqrt{n}\to\sigma/\sqrt{4n}=\tfrac{1}{2}\cdot\sigma/\sqrt{n}$,
  yarıya iner. Sezgi: güveni artırmak genişletir, veriyi artırmak daraltır;
  ikisi birbirini dengeler.}
\end{kendinleifade}

%% ============================================================
\section{En Kısa Güven Aralığı}
\label{sec:en_kisa_ga}
%% ============================================================

Birden fazla güven aralığı aynı kapsama olasılığına sahip olabilir. En bilgilisi en kısasıdır.

\begin{teorem}[Simetrik unimodal dağılımda en kısa GA]\label{thm:en_kisa}
Pivot $Q$'nun yoğunluğu unimodal ve simetrik ($f_Q$ genişleme merkezi $c$ olan) ise, sabit uzunluklu $(1-\alpha)$ aralık içinde kapsama olasılığını maksimize eden — eşdeğer olarak, sabit kapsama olasılığında uzunluğu minimize eden — aralık simetriktir: $\Prob(c-d\leq Q\leq c+d)=1-\alpha$.
\end{teorem}

\begin{proof}[Proof (taslak)]
Neyman-Pearson argümanı: $f_Q(q)>k$ kümesi üzerinde topla; $k$ normalleştirme sabitine göre ayarla. Unimodal simetrik $f_Q$ için bu küme simetrik bir aralık.
\end{proof}

\begin{sonuc}[Normal güven aralığı optimal]\label{cor:normal_optimal}
$\Normal(0,1)$ pivot için simetrik $z$-aralığı, verilen $1-\alpha$ değerinde en kısa olasılıklı güven aralığıdır.
\end{sonuc}

\begin{ornek}[$\chi^2$ en kısa aralığı]\label{ex:chi2_en_kisa}
$\chi^2_{19}$ asimetrik olduğundan simetrik eşit kuyruk aralığı en kısa değildir. Optimal kesim noktaları $c_L, c_U$ için:
\[
  f_{\chi^2_{19}}(c_L) = f_{\chi^2_{19}}(c_U), \quad \Prob(c_L\leq Q\leq c_U)=0.95.
\]
Sayısal çözüm gerektirir; pratikte eşit kuyruk aralığı standart olarak kullanılır.
\end{ornek}

%% ============================================================
\section{Profil Olabilirlik Güven Aralığı}
\label{sec:profil_ga}
%% ============================================================

\begin{tanim}[Profil olabilirlik]\label{def:profil_likh}
$\theta=(\psi,\lambda)$, ilgi parametresi $\psi$, tortu $\lambda$. \textbf{Profil olabilirlik}:
\[
  L_P(\psi) = \max_{\lambda} L(\psi,\lambda) = L(\psi,\hat\lambda(\psi)).
\]
\textbf{Profil olabilirlik oran GA} (Wilks teoreminden):
\[
  R_P(\psi) = 2[\ell(\hat\theta) - \ell_P(\psi)] \leq \chi^2_{1,\alpha}.
\]
\end{tanim}

\begin{teorem}[Wilks — profil olabilirlik GA]\label{thm:wilks_profil}
Düzenli model, $n\to\infty$. O zaman $\{\psi: R_P(\psi)\leq\chi^2_{1,\alpha}\}$ yaklaşık $(1-\alpha)$ güven aralığı oluşturur.
\end{teorem}

\begin{ornek}[Lognormal parametre aralığı]\label{ex:profil_lognormal}
$X_i\iid\mathrm{LogNormal}(\mu,\sigma^2)$. $\psi=e^\mu$ ilgi parametresi. Profil GA:
1. $\hat\mu,\hat\sigma^2$ = MLE. 2. Farklı $\mu$ değerleri için $\hat\sigma^2(\mu)=\frac1n\sum(\log X_i-\mu)^2$. 3. $R_P(\mu)\leq3.84$ kümesi GA.

Profil GA, Wald GA'ya göre asimetrik olabilir — gerçek GA şekline daha iyi uyar.
\end{ornek}

%% ============================================================
\section{Eşzamanlı Güven Aralıkları}
\label{sec:esz_ga}
%% ============================================================

\begin{tanim}[Eşzamanlı GA]\label{def:esz_ga}
$k$ parametre $(\theta_1,\ldots,\theta_k)$ için aralıklar $(L_j,U_j)$, $j=1,\ldots,k$. \textbf{Aile genel hata oranı} (FWER):
\[
  \Prob(\exists j: \theta_j\notin[L_j,U_j]) \leq \alpha.
\]
Bireysel $\alpha_j=\alpha/k$ (Bonferroni) bunu sağlar ama muhafazakârdır.
\end{tanim}

\begin{teorem}[Bonferroni düzeltmesi]\label{thm:bonferroni_ga}
Her aralık $\alpha/k$ düzeyinde kurulursa FWER $\leq\alpha$.
\end{teorem}

\begin{proof}
Birleşim sınırı: $\Prob(\bigcup_j A_j)\leq\sum_j\Prob(A_j)\leq k\cdot(\alpha/k)=\alpha$.
\end{proof}

\begin{teorem}[Scheffé eşzamanlı GA — regresyon]\label{thm:scheffe_ga}
Çok değişkenli regresyonda tüm doğrusal kombinasyonlar $c^\top\beta$ için eşzamanlı \%$95$ GA:
\[
  c^\top\hat\beta \pm \sqrt{p\,F_{p,n-p,\alpha}}\cdot\hat\sigma\sqrt{c^\top(X^\top X)^{-1}c}.
\]
Bu, sonsuz sayıda kombinasyonu eşzamanlı kapsar; Bonferroni'den geniş ama kapsamlıdır.
\end{teorem}

%% ============================================================
\section{Güven Bölgeleri (Çok Parametreli)}
\label{sec:guven_bolgesi}
%% ============================================================

\begin{tanim}[Güven bölgesi]\label{def:guven_bolgesi}
$p$ boyutlu parametre $\boldsymbol\theta$. Rasgele küme $C(\mathbf{X})\subseteq\mathbb{R}^p$ için $(1-\alpha)$ \textbf{güven bölgesi} \emph{(confidence region)}:
\[
  \Prob_{\boldsymbol\theta}(\boldsymbol\theta\in C(\mathbf{X})) \geq 1-\alpha \quad \text{her } \boldsymbol\theta\in\Theta \text{ için.}
\]
\end{tanim}

\begin{teorem}[Eliptik güven bölgesi]\label{thm:eliptik_bolge}
$X_1,\ldots,X_n\iid\Normal(\boldsymbol\mu,\Sigma)$, $\Sigma$ bilinen. $(\bar{\mathbf{X}}-\boldsymbol\mu)^\top (n\Sigma^{-1})(\bar{\mathbf{X}}-\boldsymbol\mu)\sim\chi^2_p$ pivottur. $(1-\alpha)$ güven bölgesi:
\[
  C(\mathbf{X}) = \{\boldsymbol\mu: (\bar{\mathbf{X}}-\boldsymbol\mu)^\top (n\Sigma^{-1})(\bar{\mathbf{X}}-\boldsymbol\mu) \leq \chi^2_{p,\alpha}\}.
\]
Bu bir $p$-boyutlu elipsoittir; $\bar{\mathbf{X}}$ merkezi, $\Sigma/(n\chi^2_{p,\alpha})$ yarı-eksen uzunlukları.
\end{teorem}

\begin{proof}
$\bar{\mathbf{X}}\sim\Normal(\boldsymbol\mu,\Sigma/n)$; standartlaştırarak $\mathbf{Z}=\Sigma^{-1/2}(\bar{\mathbf{X}}-\boldsymbol\mu)\sqrt{n}\sim\Normal(\mathbf{0},I_p)$. $\|\mathbf{Z}\|^2=(\bar{\mathbf{X}}-\boldsymbol\mu)^\top(n\Sigma^{-1})(\bar{\mathbf{X}}-\boldsymbol\mu)\sim\chi^2_p$ (standart Normal karelerinin toplamı).
\end{proof}

\begin{teorem}[Asimptotik güven bölgesi — genel]\label{thm:asimptotik_bolge}
MLE $\hat{\boldsymbol\theta}_n$ için:
\[
  n(\hat{\boldsymbol\theta}_n-\boldsymbol\theta)^\top I(\boldsymbol\theta) (\hat{\boldsymbol\theta}_n-\boldsymbol\theta) \dto \chi^2_p.
\]
Büyük örneklem güven bölgesi: $\{\boldsymbol\theta: n(\hat{\boldsymbol\theta}_n-\boldsymbol\theta)^\top \hat{I}(\hat{\boldsymbol\theta}_n)(\hat{\boldsymbol\theta}_n-\boldsymbol\theta)\leq\chi^2_{p,\alpha}\}$.
\end{teorem}

\begin{ornek}[İki parametreli Normal güven elipsi]
$X_1,\ldots,X_n\iid\Normal(\mu,\sigma^2)$, ikisi bilinmiyor; ilgi parametresi $(\mu,\sigma^2)$. Fisher bilgisi diyagonal: $I_{11}=1/\sigma^2$, $I_{22}=1/(2\sigma^4)$. Elipsin eksenler $(\mu-\hat\mu)^2/(n\hat\sigma^2)$ ve $(\sigma^2-\hat\sigma^2)^2/(2\hat\sigma^4/n)$ yönünde. İki boyutlu elipsin alanı tek boyutlu aralıkların yüzey alanıyla değil yüzey dışıyla kıyaslanır (Bonferroni daha muhafazakâr).
\end{ornek}

\begin{kritiksorgu}
Bireysel $p$ adet \%95 GA'yı oluşturmak yerine neden doğrudan güven bölgesi kurarız? Bireysel aralıkların kesişimi güven bölgesi değildir: iki \%95 aralığın kesişimi en az \%90 kapsama garantisi verir (Bonferroni). Eliptik güven bölgesi tam $1-\alpha$ garantisiyle daha verimli bilgi taşır.
\end{kritiksorgu}

%% ============================================================
\section{Tolerans Aralıkları}
\label{sec:tolerans_aralik}
%% ============================================================

\begin{tanim}[Tolerans aralığı]\label{def:tolerans_aralik}
$(1-\alpha)$ güvenle popülasyonun en az $\gamma$ oranını kapsayan aralık \textbf{tolerans aralığı} \emph{(tolerance interval)}: $L,U$ rasgeleler; gözlemlenmemiş $X_{\text{yeni}}$ için
\[
  \Prob\!\bigl[\Prob(L\leq X_{\text{yeni}}\leq U\mid\mathbf{X})\geq\gamma\bigr]\geq1-\alpha.
\]
Tolerans aralığı, güven aralığından farklıdır: bireysel gözlemleri kapsar, parametreyi değil.
\end{tanim}

\begin{teorem}[Normal tolerans aralığı]\label{thm:normal_tolerans}
$X_1,\ldots,X_n\iid\Normal(\mu,\sigma^2)$; $\mu,\sigma^2$ bilinmiyor. $(1-\alpha)$ güven ve $\gamma$ kapsama oranı için tolerans aralığı:
\[
  \bar{X} \pm k\,S,
\]
$k = k(n,\gamma,\alpha)$ tolerans çarpanı tablodan veya şu yaklaşımla: $k\approx\frac{z_{(1-\gamma)/2}}{\sqrt{1-z_\alpha^2/(2(n-1))}}\sqrt{1+\frac{1}{n}}$ (Wald-Wolfowitz yaklaşımı).
\end{teorem}

\begin{proof}[Proof — taslak]
$X_{\text{yeni}}|\mu,\sigma\sim\Normal(\mu,\sigma^2)$; $\bar X|\mu,\sigma\sim\Normal(\mu,\sigma^2/n)$. $\Prob(|\bar X-X_{\text{yeni}}|>kS)$ ifadesi $\chi^2$ ve $t$ dağılımlarını içerir; integral kapalı formda değil. $k$ tablo değerleri Monte Carlo veya yaklaşımla hesaplanır.
\end{proof}

\begin{ornek}[Tolerans aralığı — vida uzunluğu]
Makine $n=25$ vida üretti; $\bar{x}=50.2$ mm, $s=0.3$ mm. \%95 güvenle vidaların \%99'unu kapsayan tolerans aralığını bul. $k(25,0.99,0.05)\approx3.46$ (tablodan). Aralık: $50.2\pm3.46\cdot0.3=50.2\pm1.04=[49.16,\;51.24]$ mm.
\end{ornek}

\begin{hataanalizi}
Üç farklı aralık türünü karıştırmayın:
\begin{itemize}
  \item \textbf{Güven aralığı}: $\mu$ gibi bir parametre için — aralık ortalamayı kapsar.
  \item \textbf{Tahmin aralığı}: tek bir yeni gözlem için — $X_{\text{yeni}}\in[L,U]$.
  \item \textbf{Tolerans aralığı}: popülasyonun $\gamma$ oranı için — hem gelecek hem geçmiş gözlemler.
\end{itemize}
Tolerans aralığı en geniş, güven aralığı en dar. Kalite güvencesinde tolerans aralığı standart.
\end{hataanalizi}

\begin{teorem}[Tahmin aralığı — tek yeni gözlem]\label{thm:tahmin_aralik}
$X_{\text{yeni}}$ bağımsız, $\Normal(\mu,\sigma^2)$'dan. $\mu,\sigma^2$ bilinmiyor. $(1-\alpha)$ tahmin aralığı \emph{(prediction interval)}:
\[
  \bar X \pm t_{n-1,\alpha/2}\,S\sqrt{1+\frac{1}{n}}.
\]
Güven aralığından $\sqrt{1+1/n}$ faktörü kadar geniş (yeni gözlem belirsizliği ekstra).
\end{teorem}

\begin{disiplinlerarasibaglan}
\textbf{Kalite güvencesi:} Tolerans aralıkları üretim spesifikasyonu belirleme standardı (ISO 16269-6). \textbf{Tıbbi teşhis:} Referans aralıkları (sağlıklı popülasyonun \%95'ini kapsayan değerler) birer tolerans aralığıdır. \textbf{Çevre mühendisliği:} Kirletici konsantrasyonu için üst tolerans sınırı yasal düzenleme referansı.
\end{disiplinlerarasibaglan}

%% ============================================================
\section{Güven Aralıklarında Asimptotik Doğruluk: Edgeworth Düzeltmesi}
\label{sec:edgeworth_ga}
%% ============================================================

Standart Wald güven aralığı $\bar X_n \pm z_{\alpha/2}\sigma/\sqrt{n}$ adı altında kapsama olasılığı $\Phi(z_{\alpha/2})-\Phi(-z_{\alpha/2})$'ye yaklaşır; ancak bu yaklaşımın hata mertebesi $O(n^{-1/2})$'dir. Bartlett ve diğerleri, Edgeworth açılımı kullanarak daha kesin kapsama olasılığı hesaplamıştır.

\begin{teorem}[Edgeworth düzeltmesi]\label{thm:edgeworth_ga}
$X_i$ i.i.d., sonlu 4. momenti var, $\gamma_1=\mu_3/\sigma^3$ (çarpıklık). Wald güven aralığının gerçek kapsama olasılığı:
\[
  P\!\left(\theta\in\bar X_n\pm\frac{z_{\alpha/2}\sigma}{\sqrt{n}}\right)
  = (1-\alpha) + \frac{\gamma_1}{6\sqrt{n}}[z_{\alpha/2}^2-1]\cdot2\phi(z_{\alpha/2}) + O(n^{-1}).
\]
Çarpıklık terimi $\gamma_1\neq0$ ise kapsama olasılığını $(1-\alpha)$'dan saptırır.
\end{teorem}

\begin{proof}[Taslak]
$W_n=\sqrt{n}(\bar X_n-\mu)/\sigma$ için Edgeworth açılımı:
\[
  F_{W_n}(w) = \Phi(w) - \frac{\gamma_1}{6\sqrt{n}}(w^2-1)\phi(w) + O(n^{-1}).
\]
Kapsama olasılığını $F_{W_n}(z_{\alpha/2})-F_{W_n}(-z_{\alpha/2})$ olarak yazıp Edgeworth açılımını yerine koy; $\phi$ simetrik, $(w^2-1)$ anti-simetrik olmadığından fark birinci terimde sıfırlanmaz; teoremde verilen düzeltme elde edilir.
\end{proof}

\begin{ornek}[Üstel dağılımda kapsama hatası]
$X_i\sim\mathrm{Exp}(\lambda)$: $\gamma_1=2$. $n=20$, $1-\alpha=0.95$ ($z_{0.025}=1.96$).
Edgeworth düzeltmesi:
\[
  \Delta = \frac{2}{6\sqrt{20}}(1.96^2-1)\cdot2\phi(1.96) = \frac{2}{6\cdot4.47}(2.842)(0.584)\approx0.0124.
\]
Gerçek kapsama $\approx0.95+0.012=0.962$ (aşırı kapsama) — sağa çarpık dağılımda normal yaklaşım tutucu davranır.
\end{ornek}

\begin{teorem}[Bootstrap-t güven aralığının kapsama doğruluğu]\label{thm:bootstrap_kapsama}
Double bootstrap güven aralığı (Beran 1987), kapsama hatasını $O(n^{-1})$'den $O(n^{-3/2})$'ye indirger. Percentile-t bootstrap $O(n^{-1})$ doğruluk sağlar; basit yüzdelik bootstrap ise yalnızca $O(n^{-1/2})$.
\end{teorem}

\noindent\textbf{Sonuç:} Asimetrik dağılımlarda veya küçük $n$'de, double-bootstrap veya profil olabilirlik aralıkları, Wald aralığından belirgin şekilde üstündür.

%% ============================================================
\section{Adaptif Güven Aralıkları: Stein Küçültmesi ve SURE}
\label{sec:adaptif_ga}
%% ============================================================

$p\geq3$ parametreli $\boldsymbol\mu\in\mathbb{R}^p$ tahmininde, James-Stein tahminedici MLE'nin OKH'sini düşürür; bu fenomen güven bölgeleri için de paralel sonuçlar doğurur.

\begin{tanim}[SURE — Stein'ın hesapsız risk tahmincisi]
$\mathbf{X}\sim\Normal_p(\boldsymbol\mu,\sigma^2I)$, yansız risk tahmincisi:
\[
  \widehat{R}(\delta,\boldsymbol\mu) = p\sigma^2 + \sum_{i=1}^p\left[\left(\frac{\partial\delta_i}{\partial X_i}\right)^2\sigma^4 - 2\sigma^2\frac{\partial\delta_i}{\partial X_i}\right].
\]
Stein lemması ile beklenti $\boldsymbol\mu$'ya bağlı değildir; böylece risk veri gözlemlenerek tahmin edilebilir.
\end{tanim}

\begin{teorem}[James-Stein güven küresi]
$\mathbf{X}\sim\Normal_p(\boldsymbol\mu,I)$ ($p\geq3$), $\delta^{JS}=\left(1-\frac{p-2}{\|\mathbf{X}\|^2}\right)\mathbf{X}$ James-Stein tahminedici. Standart güven küresi $\{\boldsymbol\mu:\|\boldsymbol\mu-\mathbf{X}\|\leq c\}$'in kapsama olasılığı $(1-\alpha)$; James-Stein merkezli küre \emph{daha küçük} ama \emph{aynı ya da daha yüksek} kapsama sağlar (Stein 1981).
\end{teorem}

Bu, yüksek boyutlu istatistikte güven kümelerinin, doğal kayıp (Öteki yöntem: Bayes HPD bölgesi) ile tanımlandığında, sezgilerin aksine davranabileceğini gösterir.

%% ============================================================

\begin{mikroozet}
\begin{itemize}
  \item Güven aralığı: $\theta$ sabit, aralık rasgele; kapsama yorumu yönteme aittir, tek bir aralığa değil.
  \item Pivot yöntemi: dağılımı $\theta$'dan bağımsız büyüklük bul, eşitsizliği çöz.
  \item $z$/$t$/$\chi^2$ aralıkları normallik gerektirir; Welch eşitsiz varyanslarda güçlüdür.
  \item Wald büyük örneklem için evrensel; küçük oran ve uç değerlerde Wilson daha iyidir.
  \item En kısa GA: unimodal simetrik pivotta simetrik kesimler optimal.
  \item Profil olabilirlik GA: asimetrik ve parametrizasyona değişmez.
\end{itemize}
\end{mikroozet}

\begin{ozyansima}
Güven aralığının uzunluğu ne zaman belirsizliğimizi tam yansıtır? Güven aralığı ile Bayes güvenilir aralığı arasındaki felsefi fark nedir?
\end{ozyansima}

\begin{kopru}
\textbf{Nereden geldik:} Bölüm~\ref{chp:nokta_tahmin}'deki Fisher bilgisi ve asimptotik normallik Wald aralığını, Bölüm~\ref{chp:dagilim_aileleri}'deki $t/\chi^2/F$ dağılımları küçük örneklem aralıklarını mümkün kıldı.

\textbf{Nereye gidiyoruz:} Bölüm~\ref{chp:hipotez_testi}'de "bu $\theta_0$ değeri güven aralığında mı?" sorusu hipotez testine dönüşüyor. Güven aralığı ve hipotez testi dualdir.
\end{kopru}

%% ============================================================
%% ALISTIRMALAR
%% ============================================================

\cozumlualistirmalar

\begin{alistirma}[Al.13.1]{A}{$z$-aralığı}
$n=100$ gözlem; $\bar{x}=4.82$, $\sigma=1.2$ (bilinen). $\mu$ için \%90 ve \%99 güven aralıklarını hesapla ve uzunluklarını karşılaştır.
\end{alistirma}
\alcozum{%
$z_{0.05}=1.645$, $z_{0.005}=2.576$. \%90: $4.82\pm1.645\cdot0.12=[4.623,5.017]$ (uzunluk 0.394). \%99: $4.82\pm2.576\cdot0.12=[4.511,5.129]$ (uzunluk 0.617). Güven artınca uzunluk artar.%
}

\begin{alistirma}[Al.13.2]{B}{$t$-aralığı ve normallik}
$n=9$, $\bar{x}=23.1$, $s=4.5$, normal popülasyon. (a) \%95 $t$-aralığı. (b) $n=9$ yerine $n=100$ ile yaklaşık $z$-aralığını karşılaştır.
\end{alistirma}
\alcozum{%
(a) $t_{8,0.025}=2.306$; GA: $23.1\pm2.306\cdot(4.5/3)=23.1\pm3.459=[19.641,26.559]$. (b) $n=100$: $z_{0.025}=1.960$; GA: $23.1\pm1.960\cdot(4.5/10)=23.1\pm0.882$. Büyük örneklemde $t\approx z$ ve aralık çok daralır.%
}

\begin{alistirma}[Al.13.2b]{B}{Bootstrap Güven Aralığı}
$X_1,\ldots,X_{20}$ gözlemi, $\bar{x}=12.4$, $s=3.1$. (a) Nonparametrik bootstrap ile medyan için \%95 GA'nın mantığını açıklayın. (b) $B=1000$ bootstrap örnekleminde medyan standart hatası $\hat\sigma^*=0.82$ bulundu. Normal yaklaşım ile GA'yı kurun. (c) Bootstrap'ın klasik $t$-aralığına göre avantajı nedir?
\end{alistirma}
\alcozum{%
(a) Her bootstrap örnekleminden medyan $\tilde\theta^{*(b)}$ hesaplanır ($b=1,\ldots,B$). Emprik dağılım $F^*$ kullanılarak güven sınırları: $(\tilde\theta^*_{(\alpha/2)},\tilde\theta^*_{(1-\alpha/2)})$ yüzdelik yöntemi ile.
(b) Normal yaklaşım: $\tilde\theta\pm z_{\alpha/2}\hat\sigma^*=\tilde\theta\pm1.96\cdot0.82$. Medyan $\approx\bar{x}=12.4$ (simetrik ise): $12.4\pm1.607=[10.79,14.01]$.
(c) Bootstrap normallik varsaymaz; medyan, oran ve diğer düzgün olmayan istatistikler için geçerlidir; küçük örneklemde $t$-dağılımı varsayımını aşar.}

\begin{alistirma}[Al.13.2c]{B}{Bayes Güvenilir Aralığı — Hesap}
$X_1,\ldots,X_n\iid\Binom(1,p)$, $n=30$, $\sum x_i=18$. Önsel $p\sim\Beta(2,2)$. (a) Sonsal dağılımı bulun. (b) \%95 eşit-kuyruklu güvenilir aralığı Beta kantil tablosuyla kurun. (c) Bu aralığı klasik Wald aralığıyla karşılaştırın.
\end{alistirma}
\alcozum{%
(a) Sonsal: $p|X\sim\Beta(2+18,2+12)=\Beta(20,14)$.
(b) \%95 GA: $(\Beta_{0.025}(20,14), \Beta_{0.975}(20,14))$. $\Beta(20,14)$ modu $19/32=0.594$. Yaklaşık hesap: $\mu_1=20/34=0.588$, $\sigma_1^2=20\cdot14/(34^2\cdot35)=0.00696$. Normal yaklaşım: $0.588\pm1.96\cdot0.0834\approx[0.425,0.751]$.
(c) Wald: $\hat p=18/30=0.6$, $\hat\sigma=\sqrt{0.6\cdot0.4/30}=0.0894$; GA: $0.6\pm1.96\cdot0.0894=[0.425,0.775]$. Bayesçi önceden $\Beta(2,2)$ ortalamaya (0.5) doğru çeker.}

\alistirmalar

\begin{alistirma}[Al.13.3]{A}{Varyans için GA}
$n=20$, $s^2=16.4$, normal popülasyon. $\sigma^2$ için \%95 güven aralığını kur.
\end{alistirma}
\alcozum{%
$\chi^2_{19,0.025}=32.852$, $\chi^2_{19,0.975}=8.907$. GA: $[19\cdot16.4/32.852,\;19\cdot16.4/8.907]=[9.49,\;34.99]$.%
}

\begin{alistirma}[Al.13.4]{B}{Wilson oranı aralığı}
$n=20$ denemede $k=2$ başarı. $p$ için (a) Wald aralığı, (b) Wilson aralığı kur (\%95). Sonuçları yorumla.
\end{alistirma}
\alcozum{%
$\bar{p}=0.1$. (a) Wald: $0.1\pm1.96\sqrt{0.1\cdot0.9/20}=0.1\pm0.131=[-0.031,0.231]$ — olumsuz değer içeriyor, geçersiz. (b) Wilson: pay=$0.1+1.96^2/40=0.196$; payda=$1+3.84/20=1.192$; merkez$=0.164$; $\pm\approx0.155$; $[0.009,\;0.319]$ — daha sağlam.%
}

\begin{alistirma}[Al.13.5]{B}{Welch güven aralığı}
İki bağımsız grup: Grup A ($n_1=15$, $\bar{x}_1=42$, $s_1=8$), Grup B ($n_2=20$, $\bar{x}_2=37$, $s_2=12$). Varyansların eşit olmadığını varsayarak $\mu_1-\mu_2$ için \%95 Welch güven aralığını kur.
\end{alistirma}
\alcozum{%
SE $=\sqrt{64/15+144/20}=\sqrt{4.267+7.2}=\sqrt{11.467}=3.387$. $\nu=11.467^2/[(4.267^2/14)+(7.2^2/19)]=131.5/(1.298+2.737)=32.6\approx33$. $t_{33,0.025}\approx2.035$. GA: $5\pm2.035\cdot3.387=5\pm6.89=[-1.89,\;11.89]$.%
}

\begin{alistirma}[Al.13.6]{C}{Örneklem büyüklüğü optimizasyonu}
Bir üretim hattında $\sigma=5$ mm. Ortalamayı $\pm0.5$ mm hassasiyetle \%95 ile tahmin etmek için kaç ölçüm gerekir? Örneklem büyüklüğü iki katına çıkarılırsa hassasiyet nasıl değişir?
\end{alistirma}
\alcozum{%
$n\geq(1.96\cdot5/0.5)^2=384.16\Rightarrow n=385$. $n=770$'de: ME$=1.96\cdot5/\sqrt{770}=0.354$ mm, yani $1/\sqrt{2}\approx0.707$ katına düşer. Hassasiyet $\sqrt{2}$ faktörüyle artmış.%
}

\begin{alistirma}[Al.13.7]{C}{Pivot yöntemi — Uniform}
$X_1,\ldots,X_n\iid\Uniform[0,\theta]$. $Q=X_{(n)}/\theta$'nın dağılımını bul ve $\theta$ için $(1-\alpha)$ güven aralığını türet.
\end{alistirma}
\alcozum{%
$F_{X_{(n)}}(t)=(t/\theta)^n$; dolayısıyla $Q=X_{(n)}/\theta\sim\mathrm{Beta}(n,1)$ (KDF: $F_Q(q)=q^n$). $\Prob(q_{L}\leq Q\leq1)=1-\alpha$ ile $q_L=\alpha^{1/n}$. Çöz: $X_{(n)}/1\leq\theta$ her zaman; alt sınır $\Prob(Q\geq q_L)=1-q_L^n=1-\alpha$. GA: $\theta\in[X_{(n)},\;X_{(n)}/\alpha^{1/n}]$.%
}

\begin{alistirma}[Al.13.8]{B}{Eşzamanlı güven aralığı}
3 parametre $(\theta_1,\theta_2,\theta_3)$ için bireysel \%95 aralıklar kur isteniyor. (a) Bonferroni düzeltmesiyle her aralığın hangi güven düzeyinde kurulması gerektiğini söyle. (b) Eğer her aralığın uzunluğu 4 ise ve Bonferroni uygulanırsa FWER en az $1-3\cdot(1-0.9833)=0.95$ olduğunu göster.
\end{alistirma}
\alcozum{%
(a) Bonferroni: $\alpha_j=0.05/3\approx0.0167$; her aralık $\%98.33$ güven düzeyinde. (b) Her aralık hata $\Prob(A_j)=0.0167$. FWER $\leq\sum_j\Prob(A_j)=3\cdot0.0167=0.05$; yani toplam hata $\leq\alpha=0.05$. Bu doğrudan Birleşim sınırı argümanıdır.%
}

\begin{alistirma}[Al.13.9]{B}{En kısa GA karşılaştırması}
$n=10$, $s^2=9.6$, Normal popülasyon. (a) $\sigma^2$ için eşit kuyruk \%90 GA'sını hesapla ($\chi^2_{9,0.05}=16.92$, $\chi^2_{9,0.95}=3.33$). (b) Üst tek taraflı \%90 GA'sını kur. (c) Hangi GA daha kısa?
\end{alistirma}
\alcozum{%
(a) Eşit kuyruk: $[9\cdot9.6/16.92,\;9\cdot9.6/3.33]=[5.11,\;25.95]$; uzunluk $=20.84$. (b) Tek taraflı üst: $(-\infty,\;9\cdot9.6/\chi^2_{9,0.90})=(-\infty,\;9\cdot9.6/4.17)=(-\infty,\;20.72)$; tek yönlü GA daha kısa ama üst sınırı yoktur. Eşit kuyruk iki taraflı GA burada uzundur çünkü $\chi^2$ asimetriktir; en kısa iki taraflı GA sayısal optimizasyon gerektirir.%
}

\begin{alistirma}[Al.13.10]{C}{Profil olabilirlik GA}
$X_1,\ldots,X_n\iid\Normal(\mu,\sigma^2)$, ikisi de bilinmiyor. $\psi=\mu$ ilgi parametresi. (a) Profil log-olabilirliği $\ell_P(\mu)$'yü hesapla. (b) $R_P(\mu)=2[\ell(\hat\mu,\hat\sigma^2)-\ell_P(\mu)]$ ifadesini hesapla ve $n(\bar X-\mu)^2/S^2$ ile ilişkilendir.
\end{alistirma}
\alcozum{%
(a) $\ell(\mu,\sigma^2)=-\frac{n}{2}\log\sigma^2-\frac{1}{2\sigma^2}\sum(X_i-\mu)^2$. $\hat\sigma^2(\mu)=\frac1n\sum(X_i-\mu)^2$. $\ell_P(\mu)=-\frac{n}{2}\log\hat\sigma^2(\mu)-n/2$. (b) $R_P(\mu)=n\log(\hat\sigma^2(\mu)/\hat\sigma^2)=n\log(1+t^2/(n-1))$ ($t=\sqrt{n}(\bar X-\mu)/S$). $R_P\leq\chi^2_{1,\alpha}$ koşulu $|t|\leq t_{n-1,\alpha/2}$'ye indirgenir — profil olabilirlik GA, $t$-aralığıyla özdeştir.%
}

%% ============================================================
%% BÖLÜM 14 — Hipotez Testi
%% Olasılık Teorisi ve Matematiksel İstatistik
%% ============================================================
\chapter{Hipotez Testi}
\label{chp:hipotez_testi}

\bolumalinti{%
  To call in the statistician after the experiment is done may be no more than asking him to perform a post-mortem examination; he may be able to say what the experiment died of.
}{Ronald A. Fisher}

\begin{kazanimlar}
Bu bölümün sonunda okuyucu:
\begin{itemize}
  \item Hipotez testi çerçevesini (H0/H1, tip I/II hata, güç) tanımlar.
  \item Neyman-Pearson lemmasını uygulayarak en güçlü test türetir.
  \item Tek parametreli üstel ailede UMP test oluşturur.
  \item $z$, $t$, $\chi^2$, $F$ testlerini uygun durumlarda seçer ve uygular.
  \item Olabilirlik oran testini asimptotik dağılımıyla birlikte kullanır.
  \item $p$-değerini hesaplar ve doğru yorumlar.
\end{itemize}
\end{kazanimlar}

\begin{motivasyon}
Bir ilaç firması yeni ilacın plasebodan daha etkili olduğunu iddia ediyor. Bu iddiayı verilerle nasıl test ederiz? Rasgelelikten kaynaklanan farklılıkla gerçek bir etkiyi nasıl ayırt ederiz? Hipotez testi bu sorunun matematiksel yanıtıdır.

\medskip
\noindent\textbf{Köprü:} Bölüm~\ref{chp:guven_araliklari}'daki güven aralığı ve hipotez testi dual kavramlar: $\theta_0$'ın $(1-\alpha)$ GA dışında olması, $\alpha$ düzeyinde $H_0:\theta=\theta_0$'ı reddetmeye eşdeğer.
\end{motivasyon}

%% ============================================================
\section{Test Çerçevesi}
\label{sec:test_cercevesi}
%% ============================================================

\begin{tanim}[İstatistiksel test]\label{def:istatistiksel_test}
$\Theta_0\cap\Theta_1=\emptyset$, $\Theta_0\cup\Theta_1=\Theta$ ile:
\begin{align*}
  H_0 &: \theta\in\Theta_0 \quad \text{(sıfır hipotezi, null hypothesis)}\\
  H_1 &: \theta\in\Theta_1 \quad \text{(alternatif hipotez, alternative hypothesis)}
\end{align*}
\textbf{Test istatistiği} $T=T(\mathbf{X})$; \textbf{red bölgesi} $R\subseteq\mathcal{X}$: $\mathbf{x}\in R$ ise $H_0$ reddedilir.

$|\Theta_0|=|\Theta_1|=1$ ise \textbf{basit}, aksi hâlde \textbf{bileşik} hipotez.
\end{tanim}

\begin{tanim}[Hata türleri]\label{def:hata_turleri}
\[
  \alpha = \Prob_{\theta\in\Theta_0}(\mathbf{X}\in R) \quad \text{(Tip I hata oranı, anlamlılık düzeyi)}
\]
\[
  \beta = \Prob_{\theta\in\Theta_1}(\mathbf{X}\notin R) \quad \text{(Tip II hata oranı)}
\]
\textbf{Güç fonksiyonu:} $\pi(\theta) = \Prob_\theta(\mathbf{X}\in R)$.
Test $\alpha$ düzeyinde ise $\sup_{\theta\in\Theta_0}\pi(\theta)\leq\alpha$.
\end{tanim}

\begin{tanim}[$p$-değeri]\label{def:p_degeri}
$T$ test istatistiği, büyük $T$ değerlerinin $H_1$ lehine olduğu bir test için $p$\textbf{-değeri}:
\[
  p = \sup_{\theta\in\Theta_0}\Prob_\theta(T(\mathbf{X})\geq T(\mathbf{x}_{\mathrm{gözlem}})).
\]
$p\leq\alpha$ ise $H_0$ reddedilir.
\end{tanim}

\begin{tanimgerekce}
$p$-değeri "$H_0$ doğru olduğunda bu kadar veya daha uç bir gözlemi alma olasılığı"dır. $p<0.05$ "$H_0$'ın doğru olma olasılığı \%5'tir" anlamına \emph{gelmez}. Bu yaygın yorum hatasına düşmemeye dikkat.
\end{tanimgerekce}

%% ============================================================
\section{Neyman-Pearson Lemması}
\label{sec:neyman_pearson}
%% ============================================================

\begin{teorem}[Neyman-Pearson]\label{thm:neyman_pearson}
Basit hipotezler $H_0:\theta=\theta_0$ ve $H_1:\theta=\theta_1$. Olabilirlik oranı:
\[
  \Lambda(\mathbf{x}) = \frac{f(\mathbf{x};\theta_1)}{f(\mathbf{x};\theta_0)}.
\]
$k>0$ ve red bölgesi $R = \{\mathbf{x}: \Lambda(\mathbf{x})>k\}$ ile belirlenen testin anlamlılık düzeyi $\alpha^* = \Prob_{\theta_0}(\Lambda(\mathbf{X})>k)$ olsun. Bu test, aynı $\alpha\leq\alpha^*$ düzeyindeki tüm testler arasında en güçlüdür (güç fonksiyonu en yüksek).
\end{teorem}

\begin{proof}
$\phi^*$ NP testi, $\phi$ başka herhangi $\alpha$-düzeyinde test olsun. Göster: $\Beklenti_{\theta_1}[\phi^*]\geq\Beklenti_{\theta_1}[\phi]$.

Anahtar gözlem: $(\phi^*-\phi)(f(\mathbf{x};\theta_1)-kf(\mathbf{x};\theta_0))\geq0$ her $\mathbf{x}$ için (işaret kontrolü yapılır). Dolayısıyla:
\[
  0\leq\int(\phi^*-\phi)(f_1-kf_0)\,d\mathbf{x}
  = \Beklenti_{\theta_1}[\phi^*]-\Beklenti_{\theta_1}[\phi] - k(\Beklenti_{\theta_0}[\phi^*]-\Beklenti_{\theta_0}[\phi]).
\]
$\Beklenti_{\theta_0}[\phi^*]=\alpha^*\geq\Beklenti_{\theta_0}[\phi]$ olduğundan $k\geq0$ ile ikinci terim $\geq0$. Sonuç: $\Beklenti_{\theta_1}[\phi^*]\geq\Beklenti_{\theta_1}[\phi]$.
\end{proof}

\begin{ornek}[Normal NP testi]\label{ex:normal_np}
$X_1,\ldots,X_n\iid\Normal(\mu,1)$. $H_0:\mu=0$ vs $H_1:\mu=1$. NP testi:
\[
  \Lambda(\mathbf{x}) = \exp\!\left(n\bar{x} - \frac{n}{2}\right)>k \iff \bar{X}_n > c.
\]
$c = z_\alpha/\sqrt{n}$. Bu en güçlü test; güç: $\pi(1)=\Phi(\sqrt{n}-z_\alpha)$.
\end{ornek}

%% ============================================================
\section{UMP Testleri}
\label{sec:ump}
%% ============================================================

\begin{tanim}[UMP]\label{def:ump}
$\alpha$ düzeyinde test $\phi^*$, aynı düzeydeki tüm testlere göre her $\theta\in\Theta_1$ için daha yüksek güce sahipse \textbf{UMP} \emph{(uniformly most powerful)} test adını alır.
\end{tanim}

\begin{teorem}[Üstel ailede UMP]\label{thm:ustel_ump}
Tek parametreli üstel aile $f(x;\theta)=h(x)\exp[\eta(\theta)T(x)-B(\theta)]$; $\eta(\theta)$ artan. $H_0:\theta\leq\theta_0$ vs $H_1:\theta>\theta_0$ için UMP test:
\[
  \phi(\mathbf{x}) = \begin{cases}1&T(\mathbf{x})>c_\alpha\\0&T(\mathbf{x})<c_\alpha\end{cases}
\]
$c_\alpha$: $\Prob_{\theta_0}(T>c_\alpha)=\alpha$.
\end{teorem}

\begin{ispatiskele}
Karmonotonik olabilirlik oranı (MLR) özelliği: üstel ailede $\Lambda(\mathbf{x};\theta_1,\theta_0)$ yalnızca $T(\mathbf{x})$'in artan fonksiyonu. NP lemması her $(\theta_0,\theta_1)$ çifti için aynı red bölgesini ($T>c$) verir; bu nedenle "düzgün" en güçlü olur.
\end{ispatiskele}

\begin{ornek}[Poisson UMP]\label{ex:poisson_ump}
$X_1,\ldots,X_n\iid\Pois(\lambda)$. $H_0:\lambda\leq\lambda_0$ vs $H_1:\lambda>\lambda_0$. UMP: $T=\sum X_i>c_\alpha$. $T\sim\Pois(n\lambda_0)$ altında $c_\alpha$ belirlenir.
\end{ornek}

\begin{kritiksorgu}
İki yönlü alternatifler ($H_1:\theta\neq\theta_0$) için UMP test neden genellikle yoktur? Çünkü $\theta>\theta_0$ için büyük $T$ yönünde, $\theta<\theta_0$ için küçük $T$ yönünde güç artar — tek bir red bölgesi her ikisini optimize edemez.
\end{kritiksorgu}

%% ============================================================
\section{Yaygın Parametrik Testler}
\label{sec:parametrik_testler}
%% ============================================================

\subsection{Tek Örneklem $z$ ve $t$ Testleri}

\begin{teorem}[$z$-testi]\label{thm:z_testi}
$X_1,\ldots,X_n\iid\Normal(\mu,\sigma^2)$, $\sigma^2$ bilinen. $H_0:\mu=\mu_0$ için test istatistiği:
\[
  Z = \frac{\overline{X}_n-\mu_0}{\sigma/\sqrt{n}} \overset{H_0}{\sim} \Normal(0,1).
\]
\begin{itemize}
  \item $H_1:\mu>\mu_0$: red $Z>z_\alpha$.
  \item $H_1:\mu<\mu_0$: red $Z<-z_\alpha$.
  \item $H_1:\mu\neq\mu_0$: red $|Z|>z_{\alpha/2}$.
\end{itemize}
\end{teorem}

\begin{teorem}[$t$-testi]\label{thm:t_testi}
$\sigma^2$ bilinmiyor. Test istatistiği:
\[
  T = \frac{\overline{X}_n-\mu_0}{S/\sqrt{n}} \overset{H_0}{\sim} t_{n-1}.
\]
Red bölgeleri $t_{n-1}$ kantilleri ile belirlenir.
\end{teorem}

\subsection{Varyans Testi}

\begin{teorem}[$\chi^2$-testi]\label{thm:chi2_testi}
$X_1,\ldots,X_n\iid\Normal(\mu,\sigma^2)$. $H_0:\sigma^2=\sigma_0^2$ için:
\[
  Q = \frac{(n-1)S^2}{\sigma_0^2} \overset{H_0}{\sim} \chi^2_{n-1}.
\]
\end{teorem}

\subsection{İki Örneklem Testleri}

\begin{teorem}[İki örneklem $t$-testi]\label{thm:iki_t_testi}
$X_1,\ldots,X_m\iid\Normal(\mu_1,\sigma^2)$ ve $Y_1,\ldots,Y_n\iid\Normal(\mu_2,\sigma^2)$ bağımsız. $H_0:\mu_1=\mu_2$:
\[
  T = \frac{\overline{X}-\overline{Y}}{S_p\sqrt{1/m+1/n}} \overset{H_0}{\sim} t_{m+n-2}.
\]
\end{teorem}

\begin{teorem}[$F$-testi: varyans oranı]\label{thm:f_testi}
$X_1,\ldots,X_m\iid\Normal(\mu_1,\sigma_1^2)$ ve $Y_1,\ldots,Y_n\iid\Normal(\mu_2,\sigma_2^2)$ bağımsız. $H_0:\sigma_1^2=\sigma_2^2$:
\[
  F = \frac{S_X^2}{S_Y^2} \overset{H_0}{\sim} F_{m-1,n-1}.
\]
\end{teorem}

%% ============================================================
\section{Olabilirlik Oran Testi}
\label{sec:lrt}
%% ============================================================

\begin{tanim}[Olabilirlik oran istatistiği]\label{def:lrt}
\[
  \Lambda(\mathbf{x}) = \frac{\sup_{\theta\in\Theta_0}L(\theta;\mathbf{x})}{\sup_{\theta\in\Theta}L(\theta;\mathbf{x})}.
\]
$\Lambda\in[0,1]$; küçük $\Lambda$ değeri $H_0$'a karşı kanıt. Test: $\Lambda(\mathbf{x})\leq c_\alpha$ ise $H_0$ reddedilir.
\end{tanim}

\begin{teorem}[Wilks — asimptotik $\chi^2$]\label{thm:wilks}
Regularity koşulları altında $H_0:\theta\in\Theta_0$ doğruyken ($\dim\Theta-\dim\Theta_0 = r$):
\[
  -2\log\Lambda(\mathbf{X}) \dto \chi^2_r.
\]
\end{teorem}

\begin{ispatiskele}
\textbf{Fikir:} $\log L$ etrafında Taylor açılımı: kısıtlı MLE'yi kısıtsız MLE ile karşılaştır. İkinci türev Fisher bilgisini verir; Cauchy-Schwarz argümanıyla $\chi^2$ sınırı çıkar. Tam ispat: Hogg-Craig §10.4.
\end{ispatiskele}

\begin{ornek}[ORT: Normal iki parametre]\label{ex:ort_normal}
$X_1,\ldots,X_n\iid\Normal(\mu,\sigma^2)$; $H_0:\mu=\mu_0$ (bilinen $\sigma^2$ olmaksızın). $-2\log\Lambda=n\log(1+T^2/(n-1))$ ile $T=(\bar{X}-\mu_0)/(S/\sqrt{n})$; büyük $n$'de $t_{n-1}$'e ve $\chi^2_1$'e yakınsar.
\end{ornek}

%% ============================================================
\section{Çoklu Test Düzeltmesi}
\label{sec:coklu_test}
%% ============================================================

\begin{tanim}[FWER ve FDR]\label{def:coklu_hata}
$m$ hipotez testi yapılıyor.
\begin{itemize}
  \item \textbf{FWER} \emph{(family-wise error rate)}: en az bir yanlış red olasılığı.
  \item \textbf{Bonferroni düzeltmesi:} Her testi $\alpha/m$ düzeyinde yap; $\mathrm{FWER}\leq\alpha$.
  \item \textbf{FDR} \emph{(false discovery rate)}: yanlış redlerin tüm redlere oranının beklentisi.
  \item \textbf{Benjamini-Hochberg:} $p_{(1)}\leq\cdots\leq p_{(m)}$; $k^*=\max\{k:p_{(k)}\leq k\alpha/m\}$; $k^*$ kadar red.
\end{itemize}
\end{tanim}

\begin{teorem}[Bonferroni eşitsizliği]\label{thm:bonferroni}
$A_i$ = "$i$-inci testte yanlış red" olayları. Bonferroni eşitsizliği:
\[
  \Prob\!\left(\bigcup_{i=1}^m A_i\right) \leq \sum_{i=1}^m \Prob(A_i) \leq m\cdot\frac{\alpha}{m} = \alpha.
\]
\end{teorem}

%% ============================================================
\section{Güven Aralığı — Test Dualliği}
\label{sec:ga_test_dual}
%% ============================================================

\begin{teorem}[Duallik]\label{thm:duallik}
$C(\mathbf{x})$ $(1-\alpha)$ güven kümesi; $\phi(\mathbf{x};\theta_0)$ her $\theta_0$ için $\alpha$ düzeyinde test. O zaman:
\[
  \phi(\mathbf{x};\theta_0) = 1 \iff \theta_0\notin C(\mathbf{x}).
\]
Bu ilişki iki yönde de kurulabilir.
\end{teorem}

\begin{proof}
$C(\mathbf{x}) = \{\theta_0: \phi(\mathbf{x};\theta_0)=0\}$ tanımla. $\Prob_{\theta_0}(\theta_0\notin C(\mathbf{X})) = \Prob_{\theta_0}(\phi(\mathbf{X};\theta_0)=1)\leq\alpha$, dolayısıyla $\Prob_{\theta_0}(\theta_0\in C(\mathbf{X}))\geq1-\alpha$.
\end{proof}

\begin{kendinleifade}
İki yönlü $t$-testi ile $t$-güven aralığı arasındaki dualliği somutlaştır. "$H_0:\mu=5$, $\alpha=0.05$ reddediliyor" ile "$5$, \%95 GA'ya dahil değil" ifadelerinin eşdeğer olduğunu göster.
\end{kendinleifade}

%% ============================================================

\begin{mikroozet}
\begin{itemize}
  \item Test: karar kuralı; Tip I = yanlış red ($\alpha$), Tip II = yanlış kabul ($\beta$); güç = $1-\beta$.
  \item NP lemması: basit hipotezler için olabilirlik oranı en güçlü test verir.
  \item UMP: üstel ailede tek yönlü alternatif için her zaman mevcuttur.
  \item ORT: genel, asimptotik $\chi^2_r$; $p$-değeri yorumu dikkat gerektirir.
\end{itemize}
\end{mikroozet}

\begin{ozyansima}
Bir test istatistiksel olarak anlamlı ama pratik açıdan önemsiz olabilir mi? $n\to\infty$ limitinde ne olur? Büyük veri istatistiksel hipotez testini nasıl etkiler?
\end{ozyansima}

\begin{kopru}
\textbf{Nereden geldik:} Bölüm~\ref{chp:guven_araliklari}'deki güven aralıkları, bu bölümdeki testlerin çift yüzü.

\textbf{Nereye gidiyoruz:} Bölüm~\ref{chp:dogrusal_modeller}'de regresyon modellerinde $\beta$ katsayıları için $t$ ve $F$ testleri, bu bölümün doğrudan uygulamasıdır.
\end{kopru}

%% ============================================================
%% ALISTIRMALAR
%% ============================================================

\cozumlualistirmalar

\begin{alistirma}[Al.14.1]{A}{Tek örneklem $t$-testi}
$n=25$, $\bar{x}=3.8$, $s=1.2$. $H_0:\mu=4$ vs $H_1:\mu<4$. $\alpha=0.05$ düzeyinde test yap ve $p$-değerini hesapla.
\end{alistirma}
\alcozum{%
$T=(3.8-4)/(1.2/5)=-0.2/0.24=-0.833$. $t_{24}$ altında $p=\Prob(T<-0.833)\approx0.206$. $p>0.05$: $H_0$ reddedilemez.%
}

\begin{alistirma}[Al.14.2]{B}{NP lemması uygulaması}
$X\sim\Normal(\mu,1)$, tek gözlem. $H_0:\mu=0$ vs $H_1:\mu=2$. $\alpha=0.05$ için NP testini bul ve gücü hesapla.
\end{alistirma}
\alcozum{%
$\Lambda(x)=\exp(2x-2)>k\iff x>c$. $\Prob_0(X>c)=0.05\Rightarrow c=1.645$. Güç: $\Prob_2(X>1.645)=\Prob(\Normal(2,1)>1.645)=\Phi(0.355)\approx0.639$.%
}

\alistirmalar

\begin{alistirma}[Al.14.3]{A}{$\chi^2$ varyans testi}
$n=20$ gözlem; $s^2=18.5$. $H_0:\sigma^2=15$ vs $H_1:\sigma^2>15$. $\alpha=0.05$ için $Q$ hesapla ve karar ver.
\end{alistirma}
\alcozum{%
$Q=19\cdot18.5/15=23.43$. $\chi^2_{19,0.05}=30.14$. $Q<30.14$: $H_0$ reddedilemez.%
}

\begin{alistirma}[Al.14.4]{B}{Olabilirlik oran testi — Poisson}
$X_1,\ldots,X_{20}\iid\Pois(\lambda)$, $\sum x_i=38$. $H_0:\lambda=1.5$ vs $H_1:\lambda\neq1.5$ için ORT'yi hesapla ve Wilks yaklaşımıyla $p$-değerini bul.
\end{alistirma}
\alcozum{%
$\MLE\lambda=\bar{x}=1.9$. $-2\log\Lambda=-2[20(\log1.5-\log1.9)+20(1.9-1.5)]=2[20\log(1.9/1.5)-8]=2[20\cdot0.2357-8]=2[4.714-8]=-2(-3.286)=6.57$. Wilks: $\chi^2_1$; $p=\Prob(\chi^2_1>6.57)\approx0.010$. $H_0$ reddedilir.%
}

\begin{alistirma}[Al.14.5]{B}{Güç analizi}
$X_1,\ldots,X_n\iid\Normal(\mu,4)$. $H_0:\mu=0$ vs $H_1:\mu=1$; $\alpha=0.05$. Gücü \%80'e çıkarmak için kaç gözlem gerekli?
\end{alistirma}
\alcozum{%
Güç $=\Phi(\sqrt{n}\cdot1/2-z_{0.05})=\Phi(\sqrt{n}/2-1.645)$. \%80 için $\Phi^{-1}(0.80)=0.842$: $\sqrt{n}/2-1.645=0.842\Rightarrow\sqrt{n}=4.974\Rightarrow n=24.74\Rightarrow n=25$.%
}

\begin{alistirma}[Al.14.6]{C}{UMP yokluğu — iki yönlü}
$X_1,\ldots,X_n\iid\Normal(\mu,1)$. $H_0:\mu=0$ vs $H_1:\mu\neq0$. UMP testin var olmadığını göster. \textit{(İpucu: $\mu=\delta>0$ için NP ve $\mu=-\delta$ için NP farklı red bölgesi gerektirir.)}
\end{alistirma}
\alcozum{%
$H_1:\mu=\delta>0$ için NP red bölgesi $\bar{X}>c$ (sağ kuyruk). $H_1:\mu=-\delta<0$ için NP $\bar{X}<-c$ (sol kuyruk). Bu iki bölge farklı; tek bir $\phi$ her ikisi için UMP olamaz. Dolayısıyla $H_1:\mu\neq0$ için UMP test yoktur.%
}

\begin{alistirma}[Al.14.7]{C}{Bonferroni vs BH}
$m=10$ hipotez testi, $p$-değerleri: $0.001, 0.008, 0.039, 0.041, 0.042, 0.06, 0.1, 0.12, 0.3, 0.8$. $\alpha=0.05$ için (a) Bonferroni ve (b) BH yöntemiyle kaç hipotez reddedilir?
\end{alistirma}
\alcozum{%
(a) Bonferroni: $0.05/10=0.005$; sadece $p=0.001$ ve $p=0.008$ bu eşiğin altında; 2 red. (b) BH: $p_{(k)}\leq k\cdot0.05/10=0.005k$. $k=1:0.001\leq0.005$\checkmark; $k=2:0.008\leq0.010$\checkmark; $k=3:0.039\leq0.015$ — hayır. $k^*=2$; 2 red. (Bu örnekte eşit sonuç; genel olarak BH daha fazla reddeder.)%
}

%% ============================================================
%% BÖLÜM 15 — Doğrusal Modeller ve Regresyon
%% Olasılık Teorisi ve Matematiksel İstatistik
%% ============================================================
\chapter{Doğrusal Modeller ve Regresyon}
\label{chp:dogrusal_modeller}

\bolumalinti{%
  All models are wrong, but some are useful.
}{George Box}

\begin{kazanimlar}
Bu bölümün sonunda okuyucu:
\begin{itemize}
  \item Basit doğrusal regresyon modelini kurar; EKK tahmin edicilerini türetir ve yorumlar.
  \item Çok değişkenli regresyon modelini matris formunda ifade eder; $\hat{\boldsymbol{\beta}} = (\mathbf{X}^\top\mathbf{X})^{-1}\mathbf{X}^\top\mathbf{y}$ formülünü ispat eder.
  \item Gauss-Markov teoremini uygular; BLUE tahmin edicisini belirler.
  \item Normal hata varsayımı altında $t$ ve $F$ testleri yapar.
  \item Belirlilik katsayısı $R^2$ ile modeli değerlendirir.
  \item Regresyon varsayımlarını kontrol eder; çoklu doğrusallık ve aykırı değer sorunlarını tanır.
\end{itemize}
\end{kazanimlar}

\begin{motivasyon}
"Boy ile kilo arasında nasıl bir ilişki var?", "Reklam harcaması satışları nasıl etkiliyor?" — bu sorular bir değişkeni diğerlerinden tahmin etmeyi hedefler. Doğrusal regresyon bu hedefin en yaygın ve en iyi anlaşılmış matematiksel çerçevesidir.

\medskip
\noindent\textbf{Köprü:} Bölüm~\ref{chp:istatistik_temelleri}'deki yeterli istatistik ve Fisher-Neyman teoremi, regresyondaki $(\bar{X},S_{xx})$ çiftinin yeterliliğini açıklar. Bölüm~\ref{chp:dagilim_aileleri}'deki $t$, $F$, $\chi^2$ dağılımları burada doğrudan kullanılır.
\end{motivasyon}

%% ============================================================
\section{Basit Doğrusal Regresyon}
\label{sec:basit_regresyon}
%% ============================================================

\begin{tanim}[Basit doğrusal regresyon modeli]\label{def:bdr_model}
\[
  Y_i = \beta_0 + \beta_1 x_i + \varepsilon_i, \quad i=1,\ldots,n,
\]
$x_i$ sabit (rasgele değil), $\varepsilon_i\iid(0,\sigma^2)$, $\boldsymbol{\beta}=(\beta_0,\beta_1)^\top$ bilinmiyor.
\end{tanim}

\subsection{EKK Tahmin Edicileri}

\begin{teorem}[EKK tahmin edicileri]\label{thm:ekk_tah}
\[
  \hat\beta_1 = \frac{S_{xy}}{S_{xx}}, \qquad \hat\beta_0 = \bar{Y} - \hat\beta_1\bar{x},
\]
$S_{xx}=\sum(x_i-\bar{x})^2$, $S_{xy}=\sum(x_i-\bar{x})(Y_i-\bar{Y})$ ile.
\end{teorem}

\begin{proof}
$\mathrm{EKK}(\beta_0,\beta_1) = \sum_{i=1}^n(Y_i-\beta_0-\beta_1 x_i)^2$ minimizasyonu. $\partial/\partial\beta_0=0$: $\hat\beta_0=\bar{Y}-\hat\beta_1\bar{x}$. $\partial/\partial\beta_1=0$: $\hat\beta_1=S_{xy}/S_{xx}$.
\end{proof}

\begin{teorem}[EKK'nın özellikleri]\label{thm:ekk_ozell}
Normal hata varsayımı olmaksızın ($\Beklenti[\varepsilon_i]=0$, $\Var(\varepsilon_i)=\sigma^2$):
\begin{enumerate}[(a)]
  \item $\hat\beta_0$, $\hat\beta_1$ yansız: $\Beklenti[\hat\beta_j]=\beta_j$.
  \item $\Var(\hat\beta_1)=\sigma^2/S_{xx}$, $\Var(\hat\beta_0)=\sigma^2(\frac{1}{n}+\frac{\bar{x}^2}{S_{xx}})$.
  \item $S^2_e = \frac{\mathrm{SSE}}{n-2}$ yansız: $\Beklenti[S^2_e]=\sigma^2$ (burada $\mathrm{SSE}=\sum(Y_i-\hat{Y}_i)^2$).
\end{enumerate}
\end{teorem}

\begin{proof}
(a) $\hat\beta_1 = \sum c_i Y_i$ ($c_i=(x_i-\bar{x})/S_{xx}$); $\Beklenti[\hat\beta_1]=\sum c_i(\beta_0+\beta_1 x_i)=\beta_1$ ($\sum c_i=0$, $\sum c_i x_i=1$).

(b) $\Var(\hat\beta_1)=\sigma^2\sum c_i^2=\sigma^2/S_{xx}$.

(c) $\mathrm{SSE}=\mathbf{Y}^\top(I-H)\mathbf{Y}$ ($H$ şapka matrisi); $(I-H)$ projeksiyon; iz$(I-H)=n-2$. Özdeger argümanıyla $\mathrm{SSE}/\sigma^2\sim\chi^2_{n-2}$.
\end{proof}

%% ============================================================
\section{Çok Değişkenli Doğrusal Regresyon}
\label{sec:cok_regresyon}
%% ============================================================

\begin{tanim}[Matris modeli]\label{def:matris_model}
$n\times(p+1)$ tasarım matrisi $\mathbf{X}$ (birinci sütun $\mathbf{1}$), $\boldsymbol{\beta}\in\mathbb{R}^{p+1}$, $\boldsymbol{\varepsilon}\sim(\mathbf{0},\sigma^2 I_n)$:
\[
  \mathbf{Y} = \mathbf{X}\boldsymbol{\beta} + \boldsymbol{\varepsilon}.
\]
\end{tanim}

\begin{teorem}[EKK tahmini]\label{thm:cok_ekk}
$\mathbf{X}^\top\mathbf{X}$ tersinir ise EKK tahmin edicisi:
\[
  \hat{\boldsymbol{\beta}} = (\mathbf{X}^\top\mathbf{X})^{-1}\mathbf{X}^\top\mathbf{Y}.
\]
\end{teorem}

\begin{proof}
$S(\boldsymbol{\beta})=(\mathbf{Y}-\mathbf{X}\boldsymbol{\beta})^\top(\mathbf{Y}-\mathbf{X}\boldsymbol{\beta})$ minimize edilecek. $\nabla_\beta S = -2\mathbf{X}^\top(\mathbf{Y}-\mathbf{X}\boldsymbol{\beta})=\mathbf{0}$: normal denklem $\mathbf{X}^\top\mathbf{X}\hat{\boldsymbol{\beta}}=\mathbf{X}^\top\mathbf{Y}$. $\mathbf{X}^\top\mathbf{X}$ p.d. (tersinir); çözüm $\hat{\boldsymbol{\beta}}=(\mathbf{X}^\top\mathbf{X})^{-1}\mathbf{X}^\top\mathbf{Y}$.
\end{proof}

\begin{teorem}[Şapka matrisi]\label{thm:sapka_matris}
$H = \mathbf{X}(\mathbf{X}^\top\mathbf{X})^{-1}\mathbf{X}^\top$ \textbf{projeksiyon matrisi} (şapka matrisi). $\hat{\mathbf{Y}}=H\mathbf{Y}$, $\mathbf{e}=(I-H)\mathbf{Y}$ (artıklar). $H$ idempotent: $H^2=H$; $(I-H)^2=I-H$.
\end{teorem}

%% ============================================================
\section{Gauss-Markov Teoremi}
\label{sec:gauss_markov}
%% ============================================================

\begin{teorem}[Gauss-Markov]\label{thm:gauss_markov}
$\Beklenti[\boldsymbol{\varepsilon}]=\mathbf{0}$, $\Cov(\boldsymbol{\varepsilon})=\sigma^2 I$ (normallik varsayımı yok). Her $\mathbf{c}\in\mathbb{R}^{p+1}$ için $\mathbf{c}^\top\boldsymbol{\beta}$'nın tüm yansız doğrusal tahmin edicileri arasında EKK tahmini $\mathbf{c}^\top\hat{\boldsymbol{\beta}}$ en küçük varyanslıdır (\textbf{BLUE}).
\end{teorem}

\begin{proof}
$\tilde{\beta} = \mathbf{a}^\top\mathbf{Y}$ yansız: $\mathbf{a}^\top\mathbf{X}=\mathbf{c}^\top$. $\Var(\tilde\beta)=\sigma^2\|\mathbf{a}\|^2$. $\mathbf{a}=\mathbf{X}(\mathbf{X}^\top\mathbf{X})^{-1}\mathbf{c}+\mathbf{d}$ ($\mathbf{d}$ serbest vektör; yansızlıktan $\mathbf{X}^\top\mathbf{d}=\mathbf{0}$). O zaman $\|\mathbf{a}\|^2=\mathbf{c}^\top(\mathbf{X}^\top\mathbf{X})^{-1}\mathbf{c}+\|\mathbf{d}\|^2\geq\mathbf{c}^\top(\mathbf{X}^\top\mathbf{X})^{-1}\mathbf{c}$; eşitlik $\mathbf{d}=\mathbf{0}$'da.
\end{proof}

%% ============================================================
\section{Normal Hata Altında Çıkarım}
\label{sec:normal_cikarim}
%% ============================================================

\begin{teorem}[Normal regresyonda dağılımlar]\label{thm:normal_regresyon_dag}
$\boldsymbol{\varepsilon}\sim\Normal_n(\mathbf{0},\sigma^2 I)$ ise:
\begin{enumerate}[(a)]
  \item $\hat{\boldsymbol{\beta}}\sim\Normal_{p+1}(\boldsymbol{\beta},\sigma^2(\mathbf{X}^\top\mathbf{X})^{-1})$.
  \item $\mathrm{SSE}/\sigma^2\sim\chi^2_{n-p-1}$, $\hat{\boldsymbol{\beta}}\perp\mathrm{SSE}$.
  \item $T_j = (\hat\beta_j-\beta_j)/(S_e\sqrt{[(\mathbf{X}^\top\mathbf{X})^{-1}]_{jj}})\sim t_{n-p-1}$.
\end{enumerate}
\end{teorem}

\begin{proof}
(a) $\hat{\boldsymbol{\beta}}=(\mathbf{X}^\top\mathbf{X})^{-1}\mathbf{X}^\top\mathbf{Y}$ normalin doğrusal dönüşümü; $\Beklenti[\hat{\boldsymbol{\beta}}]=\boldsymbol{\beta}$, $\Cov=\sigma^2(\mathbf{X}^\top\mathbf{X})^{-1}$.

(b) $H$ ve $I-H$ ortogonal projeksiyon; $\hat{\mathbf{Y}}=H\mathbf{Y}$ ve $\mathbf{e}=(I-H)\mathbf{Y}$ bağımsız ($HG(I-H)=0$). $\mathrm{SSE}=\mathbf{e}^\top\mathbf{e}$; iz$(I-H)=n-p-1$ ile $\chi^2_{n-p-1}$.

(c) Standart $t$ argümanı.
\end{proof}

\subsection{$F$-Testi ve ANOVA Ayrışımı}

\begin{teorem}[$F$-testi — anlamsallık]\label{thm:f_test_regresyon}
$H_0:\beta_1=\cdots=\beta_p=0$ (sabit model) vs genel model.
\[
  F = \frac{(\mathrm{SST}-\mathrm{SSE})/p}{\mathrm{SSE}/(n-p-1)} \overset{H_0}{\sim} F_{p,n-p-1}.
\]
$\mathrm{SST}=\sum(Y_i-\bar{Y})^2$ toplam kareler toplamı.
\end{teorem}

\begin{tanim}[Belirlilik katsayısı]\label{def:r_kare}
\[
  R^2 = 1 - \frac{\mathrm{SSE}}{\mathrm{SST}} = \frac{\mathrm{SSR}}{\mathrm{SST}} \in [0,1].
\]
$\mathrm{SSR}=\sum(\hat{Y}_i-\bar{Y})^2$ regresyon kareler toplamı. $R^2$: bağımlı değişken varyansının model tarafından açıklanan oranı.
\end{tanim}

\begin{kritiksorgu}
$R^2$ tek başına model kalitesini yeterince ölçer mi? Değişken eklendikçe $R^2$ daima artar; düzeltilmiş $\bar{R}^2 = 1-(1-R^2)\frac{n-1}{n-p-1}$ bu sorunu kısmen çözer.
\end{kritiksorgu}

%% ============================================================
\section{Varsayım Kontrolü}
\label{sec:varsayim_kontrol}
%% ============================================================

\subsection{Çoklu Doğrusallık}

\begin{tanim}[VIF]\label{def:vif}
$j$. açıklayıcının diğer açıklayıcılara regresyonunda $R_j^2$ elde edilir. \textbf{Varyans şişirme faktörü} \emph{(VIF)}:
\[
  \mathrm{VIF}_j = \frac{1}{1-R_j^2}.
\]
$\mathrm{VIF}>10$ problematik çoklu doğrusallığa işaret eder.
\end{tanim}

\subsection{Hata Normalliği ve Eşit Varyanslılık}

\begin{itemize}
  \item \textbf{Artık grafiği:} Artık $e_i$ vs $\hat{Y}_i$; fan biçim $\Rightarrow$ heteroskedastik.
  \item \textbf{Q-Q grafiği:} Artıkların normallik kontrolü.
  \item \textbf{Breusch-Pagan testi:} $e_i^2$'yi $\mathbf{X}$'e regresse et; $\chi^2$ testi ile heteroskedastik.
\end{itemize}

\begin{hataanalizi}
Normallik sadece küçük örneklemlerde kritiktir; büyük $n$'de CLT ile tahmin ediciler zaten yaklaşık normaldir. Eşit varyanslılık ihlali standart hataları bozar — White sağlamlık standart hatası düzeltme yapar.
\end{hataanalizi}

\begin{disiplinlerarasibaglan}
\textbf{Ekonometri:} Regresyon ekonomik politika analizinin temeli. \textbf{Biyoistatistik:} Klinik deneylerde kovaryanslı F-testleri (ANCOVA). \textbf{Makine öğrenmesi:} Ridge, Lasso regresyon bu teori üzerine kurulur (düzenleme = Bayes önsel).
\end{disiplinlerarasibaglan}

%% ============================================================
\section{Ridge ve Lasso Regresyon (Giriş)}
\label{sec:ridge_lasso}
%% ============================================================

\begin{tanim}[Ridge regresyon]\label{def:ridge}
\[
  \hat{\boldsymbol{\beta}}_{\mathrm{ridge}} = \arg\min_{\boldsymbol\beta}\left[\|\mathbf{Y}-\mathbf{X}\boldsymbol\beta\|^2 + \lambda\|\boldsymbol\beta\|^2\right] = (\mathbf{X}^\top\mathbf{X}+\lambda I)^{-1}\mathbf{X}^\top\mathbf{Y}.
\]
$\lambda>0$: çözüm daima mevcuttur; çoklu doğrusallığa sağlamlık. EKK'ya kıyasla yanlı ama küçük varyanslı.
\end{tanim}

\begin{tanim}[Lasso]\label{def:lasso}
\[
  \hat{\boldsymbol{\beta}}_{\mathrm{lasso}} = \arg\min_{\boldsymbol\beta}\left[\|\mathbf{Y}-\mathbf{X}\boldsymbol\beta\|^2 + \lambda\|\boldsymbol\beta\|_1\right].
\]
$\ell_1$ cezası: bazı $\hat\beta_j=0$; değişken seçimi yapar (seyrek çözüm).
\end{tanim}

\begin{kendinleifade}
Ridge ve Lasso tahmin edicilerinin Bayes yorumunu açıkla. Ridge hangi önsel dağılıma karşılık gelir? Lasso için hangi önsel kullanılır?
\end{kendinleifade}

%% ============================================================
\section{Ağırlıklı En Küçük Kareler (WLS)}
\label{sec:wls}
%% ============================================================

Hata varyansları eşit değilse ($\Var(\varepsilon_i)=\sigma^2/w_i$, $w_i$ bilinen) klasik EKK BLUE değildir.

\begin{tanim}[WLS tahmin edicisi]\label{def:wls}
$W=\mathrm{diag}(w_1,\ldots,w_n)$ ağırlık matrisi. \textbf{Ağırlıklı EKK} \emph{(weighted least squares)}:
\[
  \hat{\boldsymbol\beta}_{\mathrm{WLS}} = (\mathbf{X}^\top W\mathbf{X})^{-1}\mathbf{X}^\top W\mathbf{Y}.
\]
\end{tanim}

\begin{teorem}[WLS Gauss-Markov]\label{thm:wls_gm}
$\Cov(\boldsymbol\varepsilon)=\sigma^2 W^{-1}$ altında $\hat{\boldsymbol\beta}_{\mathrm{WLS}}$ BLUE'dur.
\end{teorem}

\begin{proof}[Proof (taslak)]
$\tilde{\mathbf{Y}}=W^{1/2}\mathbf{Y}$, $\tilde{\mathbf{X}}=W^{1/2}\mathbf{X}$, $\tilde{\boldsymbol\varepsilon}=W^{1/2}\boldsymbol\varepsilon$. $\Cov(\tilde{\boldsymbol\varepsilon})=\sigma^2 I$. WLS dönüştürülmüş sistemde standart EKK; Gauss-Markov'u uygula.
\end{proof}

\begin{ornek}[Grupla toplanmış veri]\label{ex:wls_grup}
Her $x_i$ noktasında $n_i$ gözlem, grup ortalaması $\bar Y_i$ kullanılıyor. $\Var(\bar Y_i)=\sigma^2/n_i$; ağırlık $w_i=n_i$. WLS, başlangıç gözlemlerle EKK'ya eşdeğer.
\end{ornek}

%% ============================================================
\section{Model Seçim Kriterleri}
\label{sec:model_secim}
%% ============================================================

\begin{tanim}[AIC ve BIC]\label{def:aic_bic}
$p+1$ parametreli model, $\hat L$ maksimum olabilirlik:
\begin{align*}
\mathrm{AIC} &= -2\log\hat L + 2(p+1), \\
\mathrm{BIC} &= -2\log\hat L + (p+1)\log n.
\end{align*}
Küçük değer tercih edilir. BIC, $n$ büyüdükçe AIC'tan daha fazla cezalandırır; daha seyrek model seçer.
\end{tanim}

\begin{teorem}[Normal regresyonda AIC]\label{thm:aic_normal}
Normal regresyon modeli $\hat\sigma^2=\mathrm{SSE}/n$:
\[
  \mathrm{AIC} = n\log(\mathrm{SSE}/n) + 2(p+1) + \text{sabit}.
\]
\end{teorem}

\begin{proof}[Proof]
$\log\hat L = -\frac{n}{2}\log(2\pi\hat\sigma^2)-\frac{n}{2}$. $-2\log\hat L = n\log\hat\sigma^2+n(1+\log2\pi)$. AIC = bunun üstüne $2(p+1)$ ekle.
\end{proof}

\begin{tanim}[Çapraz doğrulama]\label{def:cv}
\textbf{Bırakılan-bir-dışarı (LOO-CV)} tahmin hatası:
\[
  \mathrm{CV} = \frac{1}{n}\sum_{i=1}^n (Y_i - \hat Y_{(-i)})^2,
\]
$\hat Y_{(-i)}$: $i$'inci gözlem dışlanarak elde edilen regresyon ile $x_i$ noktasındaki tahmin. Normal regresyonda kapalı form: $\mathrm{CV} = \frac{1}{n}\sum(e_i/(1-h_{ii}))^2$ ($h_{ii}$ şapka matrisinin köşegen elemanı).
\end{tanim}

%% ============================================================
\section{Regresyon Tanıtlayıcı İstatistikleri}
\label{sec:reg_tanimlayici}
%% ============================================================

\begin{tanim}[Leverage ve etki]\label{def:leverage}
Şapka matrisi $H$'nın köşegen elemanı $h_{ii}=[\mathbf{X}(\mathbf{X}^\top\mathbf{X})^{-1}\mathbf{X}^\top]_{ii}$ \textbf{leverage} \emph{(kaldıraç)} olarak adlandırılır.

$0\leq h_{ii}\leq1$; $\sum_i h_{ii}=p+1$. $h_{ii}$ büyükse $x_i$ uç noktada — regresyona büyük etki.
\end{tanim}

\begin{teorem}[Artık varyansı]\label{thm:artik_var}
$e_i = Y_i - \hat Y_i$ artık. $\Var(e_i) = \sigma^2(1-h_{ii})$; iç öğrenci artık:
\[
  r_i = \frac{e_i}{S_e\sqrt{1-h_{ii}}} \approx t_{n-p-2}.
\]
\end{teorem}

\begin{tanim}[Cook mesafesi]\label{def:cook}
\textbf{Cook mesafesi}: $i$'inci gözlem çıkarıldığında $\hat{\boldsymbol\beta}$'nın ne kadar değiştiği:
\[
  D_i = \frac{(\hat{\boldsymbol\beta}_{(-i)}-\hat{\boldsymbol\beta})^\top\mathbf{X}^\top\mathbf{X}(\hat{\boldsymbol\beta}_{(-i)}-\hat{\boldsymbol\beta})}{(p+1)S_e^2} = \frac{r_i^2 h_{ii}}{(p+1)(1-h_{ii})}.
\]
$D_i>1$ etkili nokta olarak kabul edilir; $D_i>4/n$ izleme kuralı.
\end{tanim}

\begin{ornek}[Leverage vs etki]\label{ex:leverage_etki}
Yüksek leverage ($h_{ii}$ büyük) mutlaka etkili nokta olmayabilir: $x_i$ uç ama modele uyuyorsa ($e_i\approx0$), $D_i$ küçük kalır. Kaldıraç fırsat, etki gerçekleşen etkidir.
\end{ornek}

%% ============================================================
\section{Genelleştirilmiş Doğrusal Modeller (GLM) Girişi}
\label{sec:glm}
%% ============================================================

\begin{kesif}
Normal hata varsayımını kaldıralım: ikili yanıt (0/1), sayma verisi (Poisson), süre (Gama). Bu yanıtlar için doğrusal model uygun değildir — bağlantı fonksiyonu (link function) devreye girer.
\end{kesif}

\begin{tanim}[GLM]\label{def:glm}
Üstel aileden yanıt $Y_i\sim f(y;\theta_i)$ ve $g(\mu_i)=\mathbf{x}_i^\top\boldsymbol\beta$ olsun; $\mu_i=\Beklenti[Y_i]$, $g$ \textbf{bağlantı fonksiyonu} \emph{(link function)}. Üç bileşen:
\begin{enumerate}
  \item Rasgele bileşen: $Y_i$'nin dağılımı (üstel aile).
  \item Sistematik bileşen: doğrusal öngörücü $\eta_i=\mathbf{x}_i^\top\boldsymbol\beta$.
  \item Bağlantı: $\eta_i=g(\mu_i)$.
\end{enumerate}
\end{tanim}

\begin{ornek}[Önemli GLM'ler]\label{ex:glm_ornekler}
\begin{itemize}
  \item \textbf{Lojistik regresyon:} $Y_i\sim\mathrm{Ber}(p_i)$; $g=\mathrm{logit}$: $\log(p_i/(1-p_i))=\mathbf{x}_i^\top\boldsymbol\beta$.
  \item \textbf{Poisson regresyon:} $Y_i\sim\Pois(\mu_i)$; $g=\log$: $\log\mu_i=\mathbf{x}_i^\top\boldsymbol\beta$.
  \item \textbf{Normal regresyon:} $g=\mathrm{identity}$: klasik doğrusal model.
\end{itemize}
\end{ornek}

\begin{teorem}[GLM'de MLE]\label{thm:glm_mle}
GLM katsayıları kapalı formla değil, yinelemeli ağırlıklı en küçük kareler (IRLS) ile bulunur:
\[
  \boldsymbol\beta^{(t+1)} = (\mathbf{X}^\top W^{(t)}\mathbf{X})^{-1}\mathbf{X}^\top W^{(t)}\mathbf{z}^{(t)},
\]
$W_{ii}^{(t)} = w_i^{(t)}$ çalışma ağırlığı, $z_i^{(t)}$ çalışma yanıt. Her adım WLS; bağlantı türevi ve dağılım varyansından $W$, $\mathbf{z}$ belirlenir.
\end{teorem}

\begin{ispatiskele}
Fisher puanlama: $\boldsymbol\beta^{(t+1)}=\boldsymbol\beta^{(t)}+I(\boldsymbol\beta^{(t)})^{-1}\mathbf{U}(\boldsymbol\beta^{(t)})$; $\mathbf{U}$ score vektörü, $I$ Fisher bilgisi. GLM yapısında bu yineleme IRLS biçimini alır.
\end{ispatiskele}

\begin{ornek}[Lojistik regresyon: yorum]
$\hat\beta_1=0.6$ ise $e^{0.6}=1.82$: $x$ bir birim artınca oranın $(\text{odds})$ 1.82 katına çıkması bekleniyor. Not: bu log-olasılık (log-odds) etkisidir; doğrudan olasılık farkı değildir — bu ayrım kritik.
\end{ornek}

\begin{teorem}[GLM asimptotik dağılım]\label{thm:glm_asimptotik}
Regularity koşulları altında $\hat{\boldsymbol\beta}_n$ MLE:
\[
  \sqrt{n}(\hat{\boldsymbol\beta}_n-\boldsymbol\beta) \dto \Normal_{p+1}(\mathbf{0},\; I(\boldsymbol\beta)^{-1}).
\]
Büyük örneklem $t$ ve $F$ testleri GLM çıkarımı için kullanılır.
\end{teorem}

\begin{disiplinlerarasibaglan}
\textbf{Tıp:} Lojistik regresyon hastalık risk faktörlerini modeller (odds ratio). \textbf{Ekoloji:} Tür sayısı için Poisson regresyon; Gamma GLM ağırlık verisi. \textbf{Sosyal bilimler:} İkili anket yanıtları (oy verme, tercihleri) için probit/logit.
\end{disiplinlerarasibaglan}

%% ============================================================
\section{Regresyon Tanılama: Artık Analizi}
\label{sec:artik_analiz}
%% ============================================================

\begin{tanim}[Artık türleri]\label{def:artik_turleri}
Ham artık $e_i=Y_i-\hat Y_i$'nin varyansı eşit değildir ($\Var(e_i)=\sigma^2(1-h_{ii})$). Standardize edilmiş türevler:
\begin{itemize}
  \item \textbf{İç öğrenci artık}: $r_i=e_i/(S_e\sqrt{1-h_{ii}})$ — yaklaşık $t_{n-p-2}$.
  \item \textbf{Dış öğrenci artık}: $r_i^*=e_i/(S_{(-i)}\sqrt{1-h_{ii}})$ — tam $t_{n-p-2}$ dağılımı.
\end{itemize}
\end{tanim}

\begin{teorem}[Dış öğrenci artığı kapalı formu]\label{thm:dis_artik}
\[
  r_i^* = r_i\,\sqrt{\frac{n-p-2}{n-p-1-r_i^2}}.
\]
$n-1$ kere regresyon çalıştırmak yerine tek formülle hesaplanır.
\end{teorem}

\begin{proof}[Proof — taslak]
$S_{(-i)}^2 = \frac{(n-p-1)S_e^2 - e_i^2/(1-h_{ii})}{n-p-2}$ (Sherman-Morrison güncelleme lemması). $S_{(-i)}$'yi $r_i^*$ formülüne yerleştirip basitleştir.
\end{proof}

\begin{tanim}[DFFITS ve DFBETAS]\label{def:dffits}
\textbf{DFFITS}: $i$'inci gözlemin kaldırılması halinde $\hat Y_i$'nin standardize değişimi:
\[
  \mathrm{DFFITS}_i = r_i^*\sqrt{\frac{h_{ii}}{1-h_{ii}}}.
\]
\textbf{DFBETAS}$_{ji}$: $i$'inci gözlemin kaldırılması halinde $j$. katsayı değişimi. Eşikler: $|\mathrm{DFFITS}|>2\sqrt{(p+1)/n}$ veya $|\mathrm{DFBETAS}|>2/\sqrt{n}$ etkili nokta.
\end{tanim}

\begin{ornek}[Artık analizi adım adım]
$n=20$, $p=2$ regresyonda 5. gözlem: $e_5=3.2$, $h_{55}=0.35$, $S_e=1.4$.
\begin{itemize}
  \item İç öğrenci: $r_5=3.2/(1.4\sqrt{0.65})=3.2/1.129=2.83$.
  \item Dış öğrenci: $r_5^*=2.83\sqrt{17/(17-8.01)}=2.83\sqrt{1.89}=3.89$.
  \item DFFITS: $3.89\sqrt{0.35/0.65}=3.89\cdot0.734=2.86$.
  \item Eşik $2\sqrt{3/20}=0.775$; bu nokta etkili.
\end{itemize}
\end{ornek}

\begin{teorem}[Bonferroni aykırı değer testi]\label{thm:aykiri_deger}
$n$ iç öğrenci artığı $r_1,\ldots,r_n$; en büyük mutlak değer $r_{\max}$. $H_0$: aykırı değer yok altında:
\[
  r_{\max} \approx t_{n-p-2,\,\alpha/(2n)}.
\]
Bonferroni düzeltmesiyle $\alpha$ düzeyinde aykırı değer testi.
\end{teorem}

%% ============================================================

\begin{mikroozet}
\begin{itemize}
  \item EKK $\hat{\boldsymbol\beta}=(\mathbf{X}^\top\mathbf{X})^{-1}\mathbf{X}^\top\mathbf{Y}$: Gauss-Markov'dan BLUE.
  \item Normal hata altında: $\hat{\boldsymbol\beta}$ normal, SSE $\chi^2$, katsayı testleri $t_{n-p-1}$, model testi $F_{p,n-p-1}$.
  \item $R^2$: açıklanan varyans oranı; artırılmış $\bar{R}^2$ değişken sayısına göre düzeltir.
  \item Ridge: $\ell_2$ ceza $\to$ her zaman çözüm, küçük varyans; Lasso: $\ell_1$ ceza $\to$ seyrek.
  \item WLS: heterojen varyans altında BLUE; dönüşümle standart EKK'ya indirgenir.
  \item AIC/BIC model seçim kriterleri; Cook mesafesi etkili nokta tespiti.
\end{itemize}
\end{mikroozet}

\begin{ozyansima}
Gauss-Markov teoremi normallik gerektirmiyor. Normallik varsayımı neden istatistiksel çıkarım için ekleniyor? Büyük örneklemde bu varsayımı zayıflatabilir miyiz?
\end{ozyansima}

\begin{kopru}
\textbf{Nereden geldik:} Bölüm~\ref{chp:dagilim_aileleri}'deki $\chi^2$, $t$, $F$ dağılımları; Bölüm~\ref{chp:nokta_tahmin}'deki BLUE kavramı.

\textbf{Nereye gidiyoruz:} Bölüm~\ref{chp:anova}'da regresyon çerçevesi grup karşılaştırmasına (ANOVA) genişletilir; $F$-testi bu bölümde türetildi, orada özel biçimine bakacağız.
\end{kopru}

%% ============================================================
%% ALISTIRMALAR
%% ============================================================

\cozumlualistirmalar

\begin{alistirma}[Al.15.1]{A}{Basit regresyon hesabı}
$(x_i,y_i)$: $(1,2),(2,4),(3,5),(4,4),(5,5)$. EKK doğrusunu bul; $R^2$'yi hesapla.
\end{alistirma}
\alcozum{%
$\bar{x}=3$, $\bar{y}=4$. $S_{xx}=10$, $S_{xy}=7$. $\hat\beta_1=0.7$, $\hat\beta_0=4-0.7\cdot3=1.9$. $\hat{Y}=1.9+0.7x$. SST$=6$, SSE$=\sum(y_i-\hat y_i)^2=0.09+0.01+0.09+0.01+0.16=\ldots$ Doğru: $\hat y$: 2.6,3.3,4,4.7,5.4. SSE=$(2-2.6)^2+(4-3.3)^2+(5-4)^2+(4-4.7)^2+(5-5.4)^2=0.36+0.49+1+0.49+0.16=2.5$. $R^2=1-2.5/6=0.583$.%
}

\begin{alistirma}[Al.15.2]{B}{Gauss-Markov — özel durum}
Basit regresyonda $\hat\beta_1=S_{xy}/S_{xx}$ ve alternatif yansız doğrusal tahmin edici $\tilde\beta_1=\sum w_i Y_i$ ($\sum w_i x_i=1$, $\sum w_i=0$) olsun. $\Var(\tilde\beta_1)\geq\Var(\hat\beta_1)$ olduğunu göster.
\end{alistirma}
\alcozum{%
$w_i = c_i + d_i$ ($c_i=(x_i-\bar x)/S_{xx}$, $d_i$ kısıtları sağlayan ek terim). Yansızlık: $\sum d_i=0$, $\sum d_i x_i=0$. $\Var(\tilde\beta_1)=\sigma^2\sum w_i^2=\sigma^2(\sum c_i^2+\sum d_i^2)\geq\sigma^2/S_{xx}=\Var(\hat\beta_1)$.%
}

\alistirmalar

\begin{alistirma}[Al.15.3]{A}{$F$-testi yorumu}
$n=30$, $p=3$, $\mathrm{SST}=200$, $\mathrm{SSE}=80$. $F$-istatistiğini hesapla ve $\alpha=0.05$'te $H_0:\beta_1=\beta_2=\beta_3=0$'ı test et.
\end{alistirma}
\alcozum{%
SSR$=120$. $F=(120/3)/(80/26)=40/3.077=13.0$. $F_{3,26,0.05}\approx2.98$. $F>2.98$: $H_0$ reddedilir.%
}

\begin{alistirma}[Al.15.4]{B}{Çoklu regresyon — matris}
$\mathbf{X}=\begin{pmatrix}1&1\\1&2\\1&3\end{pmatrix}$, $\mathbf{y}=(2,3,6)^\top$. $\hat{\boldsymbol\beta}$ ve $\hat{\mathbf{y}}$'yi hesapla.
\end{alistirma}
\alcozum{%
$\mathbf{X}^\top\mathbf{X}=\begin{pmatrix}3&6\\6&14\end{pmatrix}$; $(\mathbf{X}^\top\mathbf{X})^{-1}=\frac{1}{6}\begin{pmatrix}14&-6\\-6&3\end{pmatrix}$. $\mathbf{X}^\top\mathbf{y}=(11,28)^\top$. $\hat\beta=(14\cdot11-6\cdot28,\ -6\cdot11+3\cdot28)^\top/6=(-14/6,\;18/6)^\top=(-7/3,3)^\top$. $\hat\beta_0=-7/3$, $\hat\beta_1=3$. $\hat{\mathbf{y}}=(2/3,3,4/3+3)^\top=(2/3,3,8/3+\ldots)$— Doğrudan: $\hat y_i=-7/3+3x_i$: $2/3, 3, 16/3\approx5.33$.%
}

\begin{alistirma}[Al.15.5]{B}{Ridge vs EKK: OKH}
Tek değişkenli regresyon; $\beta$ skaler. EKK $\hat\beta$ ve Ridge $\hat\beta_\lambda=\frac{\sum x_i y_i}{\sum x_i^2+\lambda}$ karşılaştır. Ridge OKH'ının $\lambda^*$ için EKK OKH'ından küçük olduğu $\lambda$ aralığını bul.
\end{alistirma}
\alcozum{%
$S_{xx}=\sum x_i^2$. $\hat\beta=S_{xy}/S_{xx}$, $\hat\beta_\lambda=S_{xy}/(S_{xx}+\lambda)$. $\Beklenti[\hat\beta_\lambda]= \beta\cdot S_{xx}/(S_{xx}+\lambda)$; yanlılık $= -\beta\lambda/(S_{xx}+\lambda)$. Varyans$(\hat\beta_\lambda)=\sigma^2 S_{xx}/(S_{xx}+\lambda)^2$. OKH$(\hat\beta_\lambda)=\sigma^2 S_{xx}/(S_{xx}+\lambda)^2 + \beta^2\lambda^2/(S_{xx}+\lambda)^2$. EKK OKH$=\sigma^2/S_{xx}$. Ridge daha iyi: $\lambda<2\sigma^2/\beta^2$. Küçük $\beta^2$ (zayıf sinyal) geniş aralık; $\lambda^*=\sigma^2/(\beta^2)$ optimal.%
}

\begin{alistirma}[Al.15.6]{C}{Regresyon sertlikası — sonuç doğrulama}
Göster: basit regresyon $\hat\beta_1=\mathrm{Cor}(x,Y)\cdot(s_Y/s_x)$ ile $R^2=\mathrm{Cor}(x,Y)^2$ eşitlikleri geçerli ($s_x,s_Y$ örneklem standart sapmaları).
\end{alistirma}
\alcozum{%
$\hat\beta_1=S_{xy}/S_{xx}=\frac{\sum(x_i-\bar x)(Y_i-\bar Y)}{\sum(x_i-\bar x)^2}$. $r=S_{xy}/\sqrt{S_{xx}S_{YY}}$ ile $\hat\beta_1=r\sqrt{S_{YY}/S_{xx}}=r\cdot s_Y/s_x$ ($s_x=\sqrt{S_{xx}/(n-1)}$, $s_Y=\sqrt{S_{YY}/(n-1)}$). $R^2=\mathrm{SSR}/\mathrm{SST}=\hat\beta_1^2 S_{xx}/S_{YY}=r^2$.%
}

\begin{alistirma}[Al.15.7]{B}{WLS hesabı}
3 gözlem: $(x_i,Y_i,w_i)=(1,3,1),(2,5,4),(3,8,1)$. Ağırlıklı EKK ile $\hat\beta_0,\hat\beta_1$'i bul.
\end{alistirma}
\alcozum{%
$W=\mathrm{diag}(1,4,1)$. $\mathbf{X}^\top W\mathbf{X}=\begin{pmatrix}1&1&1\\1&2&3\end{pmatrix}W\begin{pmatrix}1&1\\1&2\\1&3\end{pmatrix}=\begin{pmatrix}6&9\\9&14\end{pmatrix}$. $\mathbf{X}^\top W\mathbf{Y}=\begin{pmatrix}1&4&1\\1&8&3\end{pmatrix}(3,5,8)^\top=(3+20+8,3+40+24)=(31,67)$. $(\mathbf{X}^\top W\mathbf{X})^{-1}=\frac{1}{3}\begin{pmatrix}14&-9\\-9&6\end{pmatrix}$. $\hat\beta=(14\cdot31-9\cdot67,-9\cdot31+6\cdot67)^\top/3=(434-603,-279+402)^\top/3=(-169/3,123/3)=(-56.3,41)$. Kontrol edin: ağırlıklı ortamada 2. nokta ($w=4$) büyük etkili.%
}

\begin{alistirma}[Al.15.8]{B}{AIC karşılaştırması}
Aynı veri, iki model: Model 1 ($p=1$, SSE$=80$), Model 2 ($p=3$, SSE$=60$), $n=30$. (a) Her model için AIC hesapla. (b) Hangi model tercih edilir?
\end{alistirma}
\alcozum{%
AIC $\propto n\log(\mathrm{SSE}/n)+2(p+1)$. Model 1: $30\log(80/30)+2\cdot2=30\cdot0.981+4=29.43+4=33.43$. Model 2: $30\log(60/30)+2\cdot4=30\cdot0.693+8=20.79+8=28.79$. Model 2 daha düşük AIC; tercih edilir. (SSE azalmasının fayda fazla cezadan büyük.)%
}

\begin{alistirma}[Al.15.9]{C}{Cook mesafesi yorumu}
$n=20$, $p=2$ regresyon; 3 gözlem için $(h_{ii}, r_i)$: A$(0.15, 2.1)$, B$(0.7, 0.3)$, C$(0.6, 2.5)$. Her biri için Cook mesafesini hesapla; yorumla.
\end{alistirma}
\alcozum{%
$D_i=r_i^2 h_{ii}/[(p+1)(1-h_{ii})]$. A: $D=4.41\cdot0.15/(3\cdot0.85)=0.662/2.55=0.26$ (orta). B: $D=0.09\cdot0.7/(3\cdot0.3)=0.063/0.9=0.07$ (yüksek leverage ama küçük artık — etkisiz). C: $D=6.25\cdot0.6/(3\cdot0.4)=3.75/1.2=3.125$ — $D>1$; en etkili nokta. A normal, B leverage var ama etkisiz, C hem leverage hem büyük artık: kaldırılması modeli önemli değiştirir.%
}

\begin{alistirma}[Al.15.10]{A}{Basit Regresyon — Skor Fonksiyonu}
$Y_i=\alpha+\beta x_i+\varepsilon_i$, $\varepsilon_i\iid\Normal(0,\sigma^2)$. (a) Log-olabilirlik $\ell(\alpha,\beta,\sigma^2)$'yi yaz. (b) $\partial\ell/\partial\alpha=0$, $\partial\ell/\partial\beta=0$ denklemlerini çözerek $\hat\alpha,\hat\beta$ EKK formüllerine ulaşın.
\end{alistirma}
\alcozum{%
(a) $\ell=-\frac{n}{2}\log(2\pi\sigma^2)-\frac{1}{2\sigma^2}\sum(Y_i-\alpha-\beta x_i)^2$.
(b) $\partial\ell/\partial\alpha=\frac{1}{\sigma^2}\sum(Y_i-\alpha-\beta x_i)=0\Rightarrow\bar Y=\hat\alpha+\hat\beta\bar x$. $\partial\ell/\partial\beta=\frac{1}{\sigma^2}\sum x_i(Y_i-\alpha-\beta x_i)=0\Rightarrow\hat\beta=S_{xy}/S_{xx}$, $\hat\alpha=\bar Y-\hat\beta\bar x$. EKK = MLE ($\Normal$ hatası altında).}

\begin{alistirma}[Al.15.11]{B}{Lojistik Regresyon — Newton-Raphson}
$Y\sim\mathrm{Bernoulli}(p)$, $\log\frac{p}{1-p}=\beta_0+\beta_1 x$. $n=5$ gözlem: $(x,y)$: $(0,0),(1,0),(2,1),(3,1),(4,1)$. (a) Log-olabilirliği yaz. (b) $\beta_0=0,\beta_1=1$ başlangıcından Newton adımı için skor $\mathbf{s}$ ve Hessian $\mathbf{H}$'yi numerik hesaplayın.
\end{alistirma}
\alcozum{%
(a) $\ell(\beta)=\sum_{i:y_i=1}(\beta_0+\beta_1x_i)-\sum_{i=1}^5\log(1+e^{\beta_0+\beta_1 x_i})$.
(b) $\hat p_i(0,1)$: $x=0,1,2,3,4$ için $p_i=\sigma(x_i)=1/(1+e^{-x_i})$:
$p_0=0.5,p_1=0.731,p_2=0.880,p_3=0.952,p_4=0.982$.
$s_0=\sum(y_i-p_i)=(0-0.5)+(0-0.731)+1(1-0.880)+(1-0.952)+(1-0.982)=-0.5-0.731+0.12+0.048+0.018=-1.045$.
$s_1=\sum x_i(y_i-p_i)=0\cdot(-0.5)+1(-0.731)+2(0.12)+3(0.048)+4(0.018)=-0.731+0.24+0.144+0.072=-0.275$.
Newton adımı: $\boldsymbol\beta^{(1)}=\boldsymbol\beta^{(0)}-H^{-1}\mathbf{s}$ (nümerik hesap gerektirdiğinden IRLS kullanılır).}

\begin{alistirma}[Al.15.12]{C}{Ridge Regresyon — Bias-Variance}
$\hat{\boldsymbol\beta}_{\mathrm{ridge}}=(\mathbf{X}^\top\mathbf{X}+\lambda\mathbf{I})^{-1}\mathbf{X}^\top\mathbf{Y}$. (a) $E[\hat{\boldsymbol\beta}_\mathrm{ridge}]$'i hesapla ve yanlılığını göster. (b) $\lambda\to\infty$ ve $\lambda\to0$ limitlerinde ne olur? (c) $\mathbf{X}^\top\mathbf{X}=\mathbf{I}$ basit durumunda MSE$({\hat\beta_j})$'yi $\lambda$ cinsinden ifade et ve minimize et.
\end{alistirma}
\alcozum{%
(a) $E[\hat{\boldsymbol\beta}_\mathrm{ridge}]=(\mathbf{X}^\top\mathbf{X}+\lambda\mathbf{I})^{-1}\mathbf{X}^\top\mathbf{X}\boldsymbol\beta_0\neq\boldsymbol\beta_0$: yanlı.
(b) $\lambda\to0$: OLS'ye yaklaşır; $\lambda\to\infty$: $\hat{\boldsymbol\beta}_\mathrm{ridge}\to\mathbf{0}$.
(c) $\mathbf{X}^\top\mathbf{X}=\mathbf{I}$: $\hat\beta_j=Y_j/(1+\lambda)$, $E[\hat\beta_j]=\beta_j/(1+\lambda)$.
$\mathrm{bias}^2=\beta_j^2\lambda^2/(1+\lambda)^2$, $\Var=\sigma^2/(1+\lambda)^2$.
MSE $=[\sigma^2+\beta_j^2\lambda^2]/(1+\lambda)^2$. $\partial\mathrm{MSE}/\partial\lambda=0$: $\lambda^*=\sigma^2/\beta_j^2$.}

%% ============================================================
%% BÖLÜM 16 — Varyans Analizi (ANOVA)
%% Olasılık Teorisi ve Matematiksel İstatistik
%% ============================================================
\chapter{Varyans Analizi (ANOVA)}
\label{chp:anova}

\bolumalinti{%
  The analysis of variance is not a mathematical theorem, but rather a convenient method of arranging the arithmetic.
}{Ronald A. Fisher}

\begin{kazanimlar}
Bu bölümün sonunda okuyucu:
\begin{itemize}
  \item Tek yönlü ANOVA modelini kurar; $F$-testini türetir ve uygular.
  \item Toplam kareler toplamı ayrışımını $\mathrm{SST} = \mathrm{SSA} + \mathrm{SSE}$ ispat eder.
  \item İki yönlü ANOVA'da ana etkiler ve etkileşim kavramlarını açıklar.
  \item Çoklu karşılaştırma yöntemlerini (Tukey, Bonferroni) uygular.
  \item Sabit, rasgele ve karma etkili modelleri ayırt eder.
  \item ANOVA ile çok değişkenli regresyon arasındaki ilişkiyi gösterir.
\end{itemize}
\end{kazanimlar}

\begin{motivasyon}
İki grubun ortalamasını karşılaştırmak için $t$-testi vardır. Üç veya daha fazla gruba genişletebilir miyiz? Tek tek $t$-testleri uygulamak tip I hata oranını patlatır. ANOVA gruplar arasındaki varyasyon ile grup içi varyasyonu karşılaştırarak bu sorunu çözer.

\medskip
\noindent\textbf{Köprü:} Bölüm~\ref{chp:dogrusal_modeller}'deki regresyon modeli ve $F$-testi, ANOVA'nın özel bir durumudur — $\mathbf{X}$ tasarım matrisi gösterge (dummy) değişkenlerden oluştuğunda ANOVA elde edilir.
\end{motivasyon}

%% ============================================================
\section{Tek Yönlü ANOVA}
\label{sec:tek_anova}
%% ============================================================

\begin{tanim}[Tek yönlü ANOVA modeli]\label{def:tek_anova}
$k$ grup; $i$. grupta $n_i$ gözlem:
\[
  Y_{ij} = \mu + \alpha_i + \varepsilon_{ij}, \quad i=1,\ldots,k;\; j=1,\ldots,n_i,
\]
$\varepsilon_{ij}\iid\Normal(0,\sigma^2)$, $\sum n_i\alpha_i=0$ (kısıt). $N=\sum n_i$ toplam gözlem. Hipotez:
\[
  H_0: \alpha_1=\alpha_2=\cdots=\alpha_k=0 \quad\text{(tüm grup ortalamaları eşit)}.
\]
\end{tanim}

\subsection{Kareler Toplamı Ayrışımı}

\begin{teorem}[SST ayrışımı]\label{thm:sst_ayrisim}
\[
  \underbrace{\sum_{i=1}^k\sum_{j=1}^{n_i}(Y_{ij}-\bar{Y}_{..})^2}_{\mathrm{SST}} =
  \underbrace{\sum_{i=1}^k n_i(\bar{Y}_{i.}-\bar{Y}_{..})^2}_{\mathrm{SSA}} +
  \underbrace{\sum_{i=1}^k\sum_{j=1}^{n_i}(Y_{ij}-\bar{Y}_{i.})^2}_{\mathrm{SSE}}.
\]
\end{teorem}

\begin{proof}
$Y_{ij}-\bar Y_{..} = (\bar Y_{i.}-\bar Y_{..}) + (Y_{ij}-\bar Y_{i.})$ ayrışımını kareye al ve çift çarpım teriminin $\sum_{i,j}(\bar Y_{i.}-\bar Y_{..})(Y_{ij}-\bar Y_{i.})=0$ olduğunu gör (her $i$ için $\sum_j(Y_{ij}-\bar Y_{i.})=0$).
\end{proof}

\begin{teorem}[ANOVA $F$-testi]\label{thm:anova_f}
$H_0$ altında:
\[
  F = \frac{\mathrm{SSA}/(k-1)}{\mathrm{SSE}/(N-k)} \sim F_{k-1,N-k}.
\]
$H_0$ reddedilir ise $F > F_{k-1,N-k,\alpha}$.
\end{teorem}

\begin{proof}
$\mathrm{SSE}/\sigma^2\sim\chi^2_{N-k}$ (her grup içi $\chi^2_{n_i-1}$; bağımsız toplam). $H_0$ altında $\mathrm{SSA}/\sigma^2\sim\chi^2_{k-1}$ (ortogonal projeksiyon). SSA $\perp$ SSE. $F$ tanımı gereği $F_{k-1,N-k}$.
\end{proof}

\begin{ornek}[ANOVA tablosu]\label{ex:anova_tablo}
\[
\begin{array}{lcccc}
\hline
\text{Kaynak} & \text{SD} & \text{KT} & \text{MKT} & F \\
\hline
\text{Gruplar} & k-1 & \mathrm{SSA} & \mathrm{MSA}=\mathrm{SSA}/(k-1) & \mathrm{MSA}/\mathrm{MSE} \\
\text{Hata} & N-k & \mathrm{SSE} & \mathrm{MSE}=\mathrm{SSE}/(N-k) & \\
\text{Toplam} & N-1 & \mathrm{SST} & & \\
\hline
\end{array}
\]
SD: serbestlik derecesi; KT: kareler toplamı; MKT: ortalama KT.
\end{ornek}

\begin{teorem}[Beklenti MSE ve MSA]\label{thm:mse_msa_beklenti}
\[
  \Beklenti[\mathrm{MSE}] = \sigma^2, \qquad
  \Beklenti[\mathrm{MSA}] = \sigma^2 + \frac{\sum n_i\alpha_i^2}{k-1}.
\]
\end{teorem}

\begin{proof}
MSE yansız ($\Beklenti[\mathrm{SSE}]=(N-k)\sigma^2$). MSA: $\Beklenti[\mathrm{SSA}]=\sum n_i\Var(\bar Y_{i.}-\bar Y_{..})+(k-1)\sigma^2+\sum n_i\alpha_i^2 = (k-1)\sigma^2+\sum n_i\alpha_i^2$.
\end{proof}

%% ============================================================
\section{Çoklu Karşılaştırma}
\label{sec:coklu_karsilastirma}
%% ============================================================

\begin{kesif}
$F$-testi $H_0$'ı reddedince hangi grupların farklı olduğunu söylemez. Post-hoc testler bu soruyu yanıtlar ama FWER kontrolü gereklidir.
\end{kesif}

\begin{tanim}[Tukey HSD]\label{def:tukey}
Dengeli tasarımda ($n_i=n$ herkese eşit). \textbf{Tukey dürüst anlamlı fark (HSD)}:
\[
  |\bar{Y}_{i.}-\bar{Y}_{j.}| > q_{\alpha,k,N-k}\!\sqrt{\frac{\mathrm{MSE}}{n}},
\]
$q_{\alpha,k,N-k}$ studentleştirilmiş aralık ($q$) dağılımının kritik değeri. Tüm $\binom{k}{2}$ çift karşılaştırmada FWER $=\alpha$.
\end{tanim}

\begin{teorem}[Scheffé yöntemi]\label{thm:scheffe}
Herhangi kontrast $L=\sum c_i\mu_i$ ($\sum c_i=0$) için $\%100(1-\alpha)$ güven aralığı:
\[
  \sum c_i\bar{Y}_{i.} \pm \sqrt{(k-1)F_{k-1,N-k,\alpha}\cdot\mathrm{MSE}\sum\frac{c_i^2}{n_i}}.
\]
Tüm olası kontrastları kapsayan en geniş yöntem; bireysel karşılaştırmalarda muhafazakâr.
\end{teorem}

%% ============================================================
\section{İki Yönlü ANOVA}
\label{sec:iki_anova}
%% ============================================================

\begin{tanim}[İki yönlü ANOVA modeli]\label{def:iki_anova}
$a$ düzeyli Faktör A, $b$ düzeyli Faktör B; her hücrede $n$ gözlem:
\[
  Y_{ijk} = \mu + \alpha_i + \beta_j + (\alpha\beta)_{ij} + \varepsilon_{ijk},
\]
$i=1,\ldots,a$; $j=1,\ldots,b$; $k=1,\ldots,n$; $\varepsilon_{ijk}\iid\Normal(0,\sigma^2)$.

Kısıtlar: $\sum\alpha_i=\sum\beta_j=\sum_i(\alpha\beta)_{ij}=\sum_j(\alpha\beta)_{ij}=0$.
\end{tanim}

\begin{teorem}[İki yönlü SST ayrışımı]\label{thm:iki_sst}
\[
  \mathrm{SST} = \mathrm{SSA} + \mathrm{SSB} + \mathrm{SSAB} + \mathrm{SSE}.
\]
Her terim karşılık gelen serbestlik derecesiyle $F$-testi oluşturur.
\end{teorem}

\begin{ornek}[İki yönlü ANOVA tablosu]\label{ex:iki_anova_tablo}
\[
\begin{array}{lccc}
\hline
\text{Kaynak} & \text{SD} & F \\
\hline
A & a-1 & \mathrm{MSA}/\mathrm{MSE} \\
B & b-1 & \mathrm{MSB}/\mathrm{MSE} \\
A\times B & (a-1)(b-1) & \mathrm{MSAB}/\mathrm{MSE} \\
Hata & ab(n-1) & \\
\hline
\end{array}
\]
\end{ornek}

\begin{tanim}[Etkileşim]\label{def:etkilesim}
$(\alpha\beta)_{ij}\neq0$: \textbf{etkileşim} \emph{(interaction)} mevcut. Faktör A'nın etkisi Faktör B'nin düzeyine göre değişiyor demektir. Etkileşim varken ana etkilerin yorumu anlamsız olabilir.
\end{tanim}

\begin{kritiksorgu}
Etkileşim anlamlıyken ana etkilere bakılmalı mı? Genellikle etkileşim istatistiği önce kontrol edilir; anlamlıysa ana etkiler ayrı ayrı her faktörün düzeyinde incelenir.
\end{kritiksorgu}

%% ============================================================
\section{Rasgele Etkili Model}
\label{sec:rasgele_etkili}
%% ============================================================

\begin{tanim}[Tek yönlü rasgele etki modeli]\label{def:rasgele_etki}
$\alpha_i\iid\Normal(0,\sigma_A^2)$ (sabit değil rasgele). Model:
\[
  Y_{ij} = \mu + \alpha_i + \varepsilon_{ij},
\]
$\alpha_i\perp\varepsilon_{ij}$. $\Var(Y_{ij})=\sigma_A^2+\sigma^2$; $\Cov(Y_{ij},Y_{ij'})=\sigma_A^2$ ($j\neq j'$).

\textbf{İçsınıf korelasyon katsayısı:} $\mathrm{ICC} = \sigma_A^2/(\sigma_A^2+\sigma^2)$.
\end{tanim}

\begin{teorem}[Rasgele etki modeli $F$-testi]\label{thm:rasgele_f}
$H_0:\sigma_A^2=0$ altında:
\[
  F = \frac{\mathrm{MSA}}{\mathrm{MSE}} \sim F_{k-1,N-k}.
\]
$F$-istatistiği aynı; yorumu farklı: gruplar arası varyans sıfır mı?
\end{teorem}

\begin{ispatiskele}
$H_0:\sigma_A^2=0$ altında $\Beklenti[\mathrm{MSA}]=\Beklenti[\mathrm{MSE}]=\sigma^2$; yani $F$ ortalaması $\approx 1$. $H_0$ doğruyken $\mathrm{SSA}/\sigma^2\sim\chi^2_{k-1}$ (rasgele $\alpha_i$ katkısı yoktur). Aynı $F$ dağılımı.
\end{ispatiskele}

%% ============================================================
\section{ANOVA — Regresyon Bağlantısı}
\label{sec:anova_regresyon}
%% ============================================================

\begin{teorem}[ANOVA = regresyon]\label{thm:anova_regresyon}
$k$ grup tek yönlü ANOVA, $k-1$ gösterge değişkenli regresyona eşdeğerdir:
\[
  x_{i\ell} = \mathbf{1}[i=\ell], \quad \ell=1,\ldots,k-1.
\]
ANOVA $F$-istatistiği ile regresyon $F$-istatistiği özdeştir.
\end{teorem}

\begin{proof}
Tasarım matrisi $\mathbf{X}$ (intercept + $k-1$ gösterge sütunu): $\hat{\boldsymbol\beta}=(\mu,\alpha_1-\alpha_k,\ldots,\alpha_{k-1}-\alpha_k)^\top$. $\mathrm{SSR}=\mathrm{SSA}$, $\mathrm{SSE}$ özdeş. $F$ aynı.
\end{proof}

\begin{disiplinlerarasibaglan}
\textbf{Tıp:} Klinik deneylerde tedavi grupları karşılaştırması. \textbf{Eğitim araştırmaları:} Farklı öğretim yöntemlerinin etkinliği. \textbf{Tarımsal araştırma:} Fisher'in orijinal motivasyonu — farklı gübre miktarlarının mahsule etkisi.
\end{disiplinlerarasibaglan}

%% ============================================================

\begin{mikroozet}
\begin{itemize}
  \item Tek yönlü ANOVA: SST = SSA + SSE; $F = \mathrm{MSA}/\mathrm{MSE}\sim F_{k-1,N-k}$.
  \item Tukey HSD: dengeli tasarımda tüm çift karşılaştırmalarda FWER kontrolü.
  \item İki yönlü ANOVA: etkileşim terimi ana etkilerin yorumunu etkiler; önce kontrol et.
  \item Rasgele etki: $\sigma_A^2$ varyans bileşeni; ICC; aynı $F$ farklı yorum.
\end{itemize}
\end{mikroozet}

\begin{ozyansima}
ANOVA normallik ve eşit varyans gerektirir. Bu varsayımlar bozulduğunda ne yapılır? Kruskal-Wallis testi bir alternatiftir ama gücü ne kadar?
\end{ozyansima}

\begin{kopru}
\textbf{Nereden geldik:} Bölüm~\ref{chp:dogrusal_modeller}'deki regresyon ve $F$-testi; Bölüm~\ref{chp:hipotez_testi}'deki çoklu test düzeltmesi.

\textbf{Bu kitabın sonu:} Olasılık teorisinden istatistiksel çıkarıma uzanan yolculuğu tamamladık. Bundan sonraki adımlar: Bölüm~\ref{chp:stokastik} ve~\ref{chp:asimptotik}'daki ileri konular, semiparametrik model ve Bayes hesaplamalı yöntemler.
\end{kopru}

%% ============================================================
%% ALISTIRMALAR
%% ============================================================

\cozumlualistirmalar

\begin{alistirma}[Al.16.1]{A}{Tek yönlü ANOVA hesabı}
3 grup: A$=\{4,6,8\}$, B$=\{3,5,7\}$, C$=\{8,10,12\}$. ANOVA tablosunu oluştur ve $\alpha=0.05$'te $H_0$ test et.
\end{alistirma}
\alcozum{%
$\bar Y_{A.}=6$, $\bar Y_{B.}=5$, $\bar Y_{C.}=10$; $\bar Y_{..}=7$. SSA$=3[(6-7)^2+(5-7)^2+(10-7)^2]=3[1+4+9]=42$. SSE$=[(4-6)^2+(6-6)^2+(8-6)^2]+[(3-5)^2+(5-5)^2+(7-5)^2]+[(8-10)^2+(10-10)^2+(12-10)^2]=8+8+8=24$. $F=(42/2)/(24/6)=21/4=5.25$. $F_{2,6,0.05}=5.14$. $F>5.14$: $H_0$ reddedilir.%
}

\begin{alistirma}[Al.16.2]{B}{SST ayrışımı ispatı}
Genel formel: $(Y_{ij}-\bar Y_{..}) = (\bar Y_{i.}-\bar Y_{..}) + (Y_{ij}-\bar Y_{i.})$. Her iki tarafın karesini alıp $i,j$ üzerinden toplayınca çift çarpım teriminin kaybolduğunu göster.
\end{alistirma}
\alcozum{%
$\sum_{i,j}(\bar Y_{i.}-\bar Y_{..})(Y_{ij}-\bar Y_{i.}) = \sum_i(\bar Y_{i.}-\bar Y_{..})\sum_j(Y_{ij}-\bar Y_{i.})$. Her $i$ için iç toplam $\sum_j(Y_{ij}-\bar Y_{i.})=0$. Dolayısıyla çift çarpım $=0$; SST = SSA + SSE.%
}

\alistirmalar

\begin{alistirma}[Al.16.3]{A}{Tukey HSD}
Alıştırma~\ref{Al.16.1}'deki veriler için Tukey HSD yöntemiyle (\%95) hangi grup çiftlerinin anlamlı biçimde farklı olduğunu belirle. ($q_{0.05,3,6}\approx4.34$.)
\end{alistirma}
\alcozum{%
MSE$=4$; $n=3$. HSD $=4.34\sqrt{4/3}=4.34\cdot1.155=5.01$. $|A-B|=1<5.01$, $|A-C|=4<5.01$, $|B-C|=5<5.01$ — Hiç çift anlamlı değil (görece küçük örneklem). Not: $F$-testi anlamlı bulmuştu ama Tukey daha muhafazakâr.%
}

\begin{alistirma}[Al.16.4]{B}{İki yönlü ANOVA — etkileşim}
$a=2$ (A: erkek/kadın), $b=2$ (B: ilaç/plasebo), $n=4$. Hücre ortalamaları: $\bar Y_{11}=20$, $\bar Y_{12}=15$, $\bar Y_{21}=22$, $\bar Y_{22}=25$. MSE$=10$, $N=16$. Etkileşim F-istatistiğini hesapla.
\end{alistirma}
\alcozum{%
$\bar Y_{1.}=17.5$, $\bar Y_{2.}=23.5$, $\bar Y_{.1}=21$, $\bar Y_{.2}=20$; $\bar Y_{..}=20.5$. SSAB $= 4\sum_{ij}[(Y_{ij}-\bar Y_{i.}-\bar Y_{.j}+\bar Y_{..})^2]$. Hücre etkileşim: $(\alpha\beta)_{11}=20-17.5-21+20.5=2$; $(\alpha\beta)_{12}=15-17.5-20+20.5=-2$; $(\alpha\beta)_{21}=22-23.5-21+20.5=-2$; $(\alpha\beta)_{22}=25-23.5-20+20.5=2$. SSAB$=4(4+4+4+4)=64$. MSAB$=64/1=64$. $F=64/10=6.4$. $F_{1,12,0.05}\approx4.75$: etkileşim anlamlı.%
}

\begin{alistirma}[Al.16.5]{C}{Rasgele etki: ICC tahmin}
$k=5$ sınıf, her sınıftan $n=10$ öğrenci test skoru. MSA$=120$, MSE$=30$. $\hat\sigma_A^2$ ve $\hat{\mathrm{ICC}}$'yi hesapla. Yorumla.
\end{alistirma}
\alcozum{%
$\Beklenti[\mathrm{MSA}]=\sigma^2+n\sigma_A^2$; $\hat\sigma^2=\mathrm{MSE}=30$; $\hat\sigma_A^2=(\mathrm{MSA}-\mathrm{MSE})/n=(120-30)/10=9$. $\hat{\mathrm{ICC}}=9/(9+30)=9/39\approx0.23$. Yorum: sınıf üyeliği test skoru varyansının \%23'ünü açıklıyor; sınıf etkisi önemli.%
}

\begin{alistirma}[Al.16.6]{C}{ANOVA — regresyon eşdeğerliği}
$k=3$ grup ANOVA'yı gösterge değişkenli regresyon olarak yeniden kur. Tasarım matrisini yaz; EKK tahminlerinin $\hat\mu$, $\hat\alpha_1-\hat\alpha_3$, $\hat\alpha_2-\hat\alpha_3$ olduğunu göster.
\end{alistirma}
\alcozum{%
$\mathbf{X}=[\mathbf{1},\mathbf{x}_1,\mathbf{x}_2]$; $x_{i1}=\mathbf{1}[$grup 1$]$, $x_{i2}=\mathbf{1}[$grup 2$]$. Normal denklem çözümü: intercept $\hat\beta_0=\bar Y_{3.}=\hat\mu+\hat\alpha_3$; $\hat\beta_1=\bar Y_{1.}-\bar Y_{3.}=\hat\alpha_1-\hat\alpha_3$; $\hat\beta_2=\bar Y_{2.}-\bar Y_{3.}=\hat\alpha_2-\hat\alpha_3$. $F=\mathrm{SSR}/2\div\mathrm{SSE}/(N-3)$ = ANOVA $F$.%
}

% %% ============================================================
%% BÖLÜM 17 — Parametresiz İstatistik  [OPSİYONEL]
%% Olasılık Teorisi ve Matematiksel İstatistik
%% ============================================================
\chapter{Parametresiz İstatistik}
\label{chp:parametresiz}

\bolumalinti{%
  The great advantage of the nonparametric approach is that we need not make
  any assumption about the form of the underlying population.
}{Frank Wilcoxon}

\begin{kazanimlar}
Bu bölümün sonunda okuyucu:
\begin{itemize}
  \item Parametresiz yöntemlerin motivasyonunu ve güç-sağlamlık dengesini açıklar.
  \item İşaret testi ve Wilcoxon işaret-sıra testini uygular ve asimptotik dağılımlarını türetir.
  \item Mann-Whitney $U$-testini iki bağımsız örneklem karşılaştırması için kullanır.
  \item Kruskal-Wallis ve Friedman testlerini çok grup karşılaştırmasında uygular.
  \item Kolmogorov-Smirnov testini dağılım uyum testi olarak kullanır.
  \item Bootstrap yöntemini standart hata ve güven aralığı hesabında uygular.
\end{itemize}
\end{kazanimlar}

\begin{motivasyon}
Normal dağılım varsayımı zaman zaman geçersizdir: çarpık dağılımlar, aykırı değerler, sıralı ölçek. Parametresiz yöntemler dağılım hakkında minimal varsayımla çalışır. "Sağlamlık" (robustness) modern istatistikte kritik önem taşır.

\medskip
\noindent\textbf{Köprü:} Bölüm~\ref{chp:hipotez_testi}'deki test çerçevesi burada da geçerlidir; fark, test istatistiklerinin sıralardan ya da işaretlerden inşa edilmesidir.
\end{motivasyon}

%% ============================================================
\section{İşaret Testi}
\label{sec:isaret_testi}
%% ============================================================

\begin{kesif}
Sürekli dağılım medyanı $m$ için $X_i>m$ veya $X_i<m$ ikili bilgisi kullanılabilir. $H_0:m=m_0$ altında bu ikili değerlerin dağılımı nedir?
\end{kesif}

\begin{tanim}[İşaret testi]\label{def:isaret_testi}
$X_1,\ldots,X_n\iid F$ sürekli. $H_0: F(m_0)=1/2$ (medyan $m_0$). $B = \#\{i:X_i>m_0\}$. $H_0$ altında $B\sim\Binom(n,1/2)$.

Test istatistiği $B$ (veya $\min(B,n-B)$ iki yönlü test için). Küçük $n$'de kesin Binom dağılımı; büyük $n$'de $B$ yaklaşık $\Normal(n/2,n/4)$.
\end{tanim}

\begin{teorem}[İşaret testinin asimptotik dağılımı]
$H_0$ altında $Z=(2B-n)/\sqrt{n}\dto\Normal(0,1)$.
\end{teorem}

\begin{proof}
$B=\sum_{i=1}^n\mathbf{1}[X_i>m_0]$; her terim i.i.d.\ $\mathrm{Ber}(1/2)$. MLT uygulanır.
\end{proof}

%% ============================================================
\section{Wilcoxon İşaret-Sıra Testi}
\label{sec:wilcoxon_isaret}
%% ============================================================

\begin{tanim}[Wilcoxon işaret-sıra istatistiği]\label{def:wilcoxon_isaret}
$X_1,\ldots,X_n$ simetrik sürekli dağılımdan (simetri merkezi $\theta$), $H_0:\theta=0$. $|X_1|,\ldots,|X_n|$'yi küçükten büyüğe sırala; $R_i = |X_i|$'nin sırası.

\textbf{Wilcoxon istatistiği:} $W^+ = \sum_{i=1}^n R_i\,\mathbf{1}[X_i>0]$.
\end{tanim}

\begin{teorem}[Wilcoxon asimptotik dağılımı]\label{thm:wilcoxon_asimptotik}
$H_0$ altında:
\[
  \Beklenti[W^+] = \frac{n(n+1)}{4}, \qquad \Var(W^+) = \frac{n(n+1)(2n+1)}{24}.
\]
Standartlaştırılmış $Z = (W^+ - n(n+1)/4)/\sqrt{n(n+1)(2n+1)/24} \dto \Normal(0,1)$.
\end{teorem}

\begin{proof}[Proof — beklenti]
$H_0$ altında her $X_i$ için $\Prob(X_i>0)=1/2$ (simetri). $W^+=\sum_{i=1}^n R_i\,\psi_i$ ($\psi_i=\mathbf{1}[X_i>0]$). $\psi_i\perp R_i$ (simetri) ve $\Beklenti[\psi_i]=1/2$.
\[
  \Beklenti[W^+] = \sum_{i=1}^n\Beklenti[R_i]\cdot\Beklenti[\psi_i] = \frac{1}{2}\sum_{i=1}^n i = \frac{n(n+1)}{4}.
\]
\end{proof}

\begin{ispatiskele}
Varyans hesabı: $\psi_i,\psi_j$ ve $R_i,R_j$ bağımsız olduğundan $\Cov(R_i\psi_i,R_j\psi_j)=0$ ($i\neq j$). $\Var(R_i\psi_i)=\Beklenti[R_i^2]/4$; $\sum_{i=1}^n i^2=n(n+1)(2n+1)/6$. Bölünce $n(n+1)(2n+1)/24$.
\end{ispatiskele}

%% ============================================================
\section{Mann-Whitney $U$-Testi}
\label{sec:mann_whitney}
%% ============================================================

\begin{tanim}[Mann-Whitney istatistiği]\label{def:mann_whitney}
$X_1,\ldots,X_m\iid F$ ve $Y_1,\ldots,Y_n\iid G$ bağımsız. $H_0:F=G$.

$U = \#\{(i,j): X_i < Y_j\} = \sum_{i=1}^m\sum_{j=1}^n\mathbf{1}[X_i<Y_j]$.
\end{tanim}

\begin{teorem}[Mann-Whitney asimptotik dağılımı]\label{thm:mann_whitney}
$H_0$ altında:
\[
  \Beklenti[U] = \frac{mn}{2}, \qquad \Var(U) = \frac{mn(m+n+1)}{12}.
\]
$Z = (U - mn/2)/\sqrt{mn(m+n+1)/12}\dto\Normal(0,1)$.
\end{teorem}

\begin{proof}[Proof — beklenti]
$H_0:F=G$ altında $\Prob(X_i<Y_j)=1/2$ (sürekli dağılım). $U=\sum_{ij}\mathbf{1}[X_i<Y_j]$, her biri $\mathrm{Ber}(1/2)$; $\Beklenti[U]=mn/2$.
\end{proof}

\begin{teorem}[Wilcoxon sıra-toplam — eşdeğerlik]\label{thm:wilcoxon_mann}
İki örneklem birleştirilip sıralanır; $W_X = $ X değerlerinin sıra toplamı. $U = W_X - m(m+1)/2$ (sıra toplamından Mann-Whitney elde edilir). İki test eşdeğerdir.
\end{teorem}

%% ============================================================
\section{Kruskal-Wallis ve Friedman Testleri}
\label{sec:kruskal_friedman}
%% ============================================================

\begin{tanim}[Kruskal-Wallis testi]\label{def:kruskal_wallis}
$k$ grup, grup $i$'de $n_i$ gözlem; toplam $N=\sum n_i$. Tüm gözlemler birleştirilip sıralanır; $\bar{R}_i$ = $i$. grubun ortalama sırası. Kruskal-Wallis istatistiği:
\[
  H = \frac{12}{N(N+1)}\sum_{i=1}^k n_i(\bar{R}_i - \bar{R})^2
\]
$\bar{R}=(N+1)/2$ (genel ortalama sıra).
\end{tanim}

\begin{teorem}[Kruskal-Wallis asimptotiği]\label{thm:kruskal_asimptotik}
$H_0$ (tüm gruplar aynı dağılım) altında $H\dto\chi^2_{k-1}$.
\end{teorem}

\begin{ispatiskele}
ANOVA sıra dönüşümü: sıraların $F$-istatistiği formülü uygulanınca $H$ elde edilir; $\chi^2$ asimptotiği normal yakınsama argümanından gelir.
\end{ispatiskele}

\begin{tanim}[Friedman testi]\label{def:friedman}
$n$ blok, $k$ muamele; blok içi sıralar $r_{ij}$ ($j=1,\ldots,k$). Friedman istatistiği:
\[
  F_r = \frac{12n}{k(k+1)}\sum_{j=1}^k\!\left(\bar{r}_j - \frac{k+1}{2}\right)^{\!2}
\]
$\bar{r}_j$ = $j$-inci muamelenin ortalama sırası. $H_0$ altında $F_r\dto\chi^2_{k-1}$.
\end{tanim}

%% ============================================================
\section{Kolmogorov-Smirnov Testi}
\label{sec:ks_testi}
%% ============================================================

\begin{tanim}[Ampirik dağılım fonksiyonu]\label{def:ampirik_df}
$X_1,\ldots,X_n$ gözlemi için \textbf{ampirik dağılım fonksiyonu} \emph{(empirical CDF)}:
\[
  F_n(x) = \frac{1}{n}\sum_{i=1}^n\mathbf{1}[X_i\leq x].
\]
\end{tanim}

\begin{teorem}[Glivenko-Cantelli]\label{thm:glivenko_cantelli}
$X_i\iid F$. O zaman
\[
  \sup_x |F_n(x)-F(x)| \asto 0.
\]
\end{teorem}

\begin{proof}[Proof (taslak)]
Her sabit $x$ için GBSK'dan $F_n(x)\asto F(x)$. Üst limit geçirgenliği (üstel sınır argümanı) sayılabilir yoğun $\{x_k\}$ kümesi üzerinden $\sup$'u kapsar; sol süreklilik ve sınırlılık ile genel sonuç elde edilir.
\end{proof}

\begin{tanim}[KS istatistiği]\label{def:ks_istatistik}
$H_0:F=F_0$ (belirtilmiş sürekli dağılım). \textbf{Kolmogorov-Smirnov istatistiği:}
\[
  D_n = \sup_x |F_n(x) - F_0(x)|.
\]
\end{tanim}

\begin{teorem}[Kolmogorov dağılımı]\label{thm:kolmogorov_dagilim}
$H_0$ altında $F_0$ sürekli ise:
\[
  \sqrt{n}\,D_n \dto K
\]
$K$'nın dağılım fonksiyonu: $\Prob(K\leq t) = \sum_{k=-\infty}^{\infty}(-1)^k e^{-2k^2t^2}$.
\end{teorem}

\begin{ispatiskele}
$U_i = F_0(X_i)\iid\Uniform[0,1]$ ($F_0$ sürekli). $F_n$ yerine uniform ampirik süreç $\sqrt{n}(\hat{F}_n(u)-u)$ incelenir; bu Brownian köprüye yakınsar (Donsker). Brownian köprünün sonsuz normu Kolmogorov dağılımını verir.
\end{ispatiskele}

%% ============================================================
\section{Medyan için Güven Aralığı}
\label{sec:medyan_ga}
%% ============================================================

\begin{teorem}[Medyan için dağılım-serbest GA]\label{thm:medyan_ga}
$X_1,\ldots,X_n\iid F$ sürekli; medyan $m$ bilinmiyor. Sıra istatistikleri $X_{(1)}\leq\cdots\leq X_{(n)}$.

$\Prob(X_{(r)}\leq m\leq X_{(s)}) = \sum_{k=r}^{s-1}\binom{n}{k}(1/2)^n$.

$(1-\alpha)$ güven aralığı: $r$ ve $s$'yi $\sum_{k=r}^{s-1}\binom{n}{k}2^{-n}\geq1-\alpha$ sağlayacak biçimde seç.
\end{teorem}

\begin{proof}
$\Prob(X_{(r)}\leq m) = \Prob(\text{en az } r \text{ gözlem } m\text{'den küçük}) = \Prob(B_{n,1/2}\geq r)$ ($B_{n,1/2}\sim\Binom(n,1/2)$). $\Prob(X_{(r)}\leq m\leq X_{(s)}) = \Prob(r\leq B_{n,1/2}\leq s-1)$.
\end{proof}

\begin{ornek}[Medyan GA, $n=9$, $\alpha=0.05$]
$\sum_{k=r}^{8-r}\binom{9}{k}(1/2)^9$ $r=2$, $s=8$: $\sum_{k=2}^7\binom{9}{k}/512=(9+36+84+126+126+84)/512=465/512=0.908<0.95$. $r=2,s=8$ yeterli değil. $r=1,s=9$: $\sum_{k=1}^8\binom{9}{k}/512=(512-1-1)/512=510/512=0.996\geq0.95$. GA: $[X_{(1)},X_{(9)}]$, kapsama $\geq99.6\%$.
\end{ornek}

%% ============================================================
\section{Spearman Sıra Korelasyonu}
\label{sec:spearman}
%% ============================================================

\begin{tanim}[Spearman'ın $\rho$ istatistiği]\label{def:spearman}
$(X_i,Y_i)_{i=1}^n$ çiftleri. $R_i$ = $X_i$'nin sırası, $S_i$ = $Y_i$'nin sırası. \textbf{Spearman $\rho$}:
\[
  r_S = 1 - \frac{6\sum_{i=1}^n d_i^2}{n(n^2-1)}, \quad d_i = R_i - S_i.
\]
Parametrik Pearson korelasyonunun sıra tabanlı alternatifi; aykırı değerlere sağlamlı.
\end{tanim}

\begin{teorem}[Spearman $\rho$ — bağlantı]
$r_S$ aslında sıraların Pearson korelasyonudur: $r_S = \Cor(R_i,S_i)$. $H_0:\rho=0$ (bağımsızlık) altında $T=r_S\sqrt{(n-2)/(1-r_S^2)}\dto t_{n-2}$.
\end{teorem}

\begin{proof}
Sıralar $1,\ldots,n$'nin ortalaması $(n+1)/2$; varyansı $(n^2-1)/12$. Pearson formülünü sıralara uygulayınca $r_S$ elde edilir. Büyük örneklem: $\sqrt{n-1}r_S\dto\Normal(0,1)$ ($H_0$ altında).
\end{proof}

\begin{ornek}[Spearman $\rho$ hesabı]
$(X,Y)$ çiftleri: sıralar $(R,S)$: $(1,2),(2,1),(3,4),(4,3),(5,5)$. $d_i=(-1,1,-1,1,0)$. $\sum d_i^2=4$. $r_S=1-6\cdot4/(5\cdot24)=1-0.2=0.8$.
\end{ornek}

%% ============================================================
\section{Permütasyon Testleri}
\label{sec:permutasyon_test}
%% ============================================================

\begin{tanim}[Permütasyon testi]\label{def:permutasyon_test}
$H_0$: iki grup $\{X_i\}$ ve $\{Y_j\}$ aynı dağılımdan. Test istatistiği $T$ (örn. $|\bar{X}-\bar{Y}|$).

\textbf{Permütasyon dağılımı:} Tüm $\binom{m+n}{m}$ mümkün grup atamasında $T$ değerleri hesaplanır. $p$-değeri: gözlenen $T$'den büyük ya da eşit $T$ değerlerinin oranı.
\end{tanim}

\begin{teorem}[Permütasyon testinin kesin geçerliliği]
$H_0$ altında tüm permütasyonlar eşit olasılıklı; permütasyon $p$-değeri tam anlamlılık düzeyi $\alpha$'yı kontrol eder (dağılım varsayımı gerektirmez).
\end{teorem}

\begin{ornek}[Küçük örneklem permütasyon testi]
Grup A: $\{3,7\}$, Grup B: $\{1,5\}$. $T=|\bar{A}-\bar{B}|=|5-3|=2$. Tüm $\binom{4}{2}=6$ bölme: $\{3,7\}$ vs $\{1,5\}$: $T=2$; $\{3,1\}$ vs $\{7,5\}$: $T=4$; $\{3,5\}$ vs $\{1,7\}$: $T=2$; $\{1,7\}$ vs $\{3,5\}$: $T=2$; $\{1,5\}$ vs $\{3,7\}$: $T=2$; $\{7,5\}$ vs $\{3,1\}$: $T=4$. $p=\Prob(T\geq2)=6/6=1.00$. $H_0$ reddedilemez.
\end{ornek}

%% ============================================================
\section{Asimptotik Göreli Verimlilik (ARE)}
\label{sec:are}
%% ============================================================

\begin{tanim}[ARE]\label{def:are}
$T_n^{(1)}$ ve $T_n^{(2)}$ aynı hipotez için iki test. $n_1(n_2)$: $T^{(1)}(T^{(2)})$'nin belirli güce ($1-\beta$) ulaşması için gereken örneklem büyüklükleri.
\[
  \mathrm{ARE}(T^{(1)},T^{(2)}) = \lim_{n\to\infty}\frac{n_2}{n_1}.
\]
$\mathrm{ARE}>1$: $T^{(1)}$ daha verimli; daha az gözlemle aynı güce ulaşır.
\end{tanim}

\begin{teorem}[Wilcoxon vs $t$-testi ARE]\label{thm:are_wilcoxon}
Normal dağılım altında:
\[
  \mathrm{ARE}(\text{Wilcoxon},t) = \frac{3}{\pi} \approx 0.955.
\]
Wilcoxon normallik geçerliyken $t$'nin $\%95.5$ veriminde çalışır — ama normallikten sapmalarda Wilcoxon üstün olabilir.
\end{teorem}

\begin{ispatiskele}
Pitman verimlilik: $\mathrm{ARE}=\mathrm{SH}^2_{t}/\mathrm{SH}^2_W$ lokal alternatif altında asimptotik güç eğrileri karşılaştırılarak türetilir. Normal $f$ için Pitman verimlilik $12\sigma^2[f(0)]^2 = 12\sigma^2 /(2\pi\sigma^2)=6/\pi\to3/\pi$ (tek örneklem Wilcoxon için $3/\pi\approx0.955$).
\end{ispatiskele}

\begin{teorem}[ARE — genel sonuç]
\begin{itemize}
  \item Herhangi sürekli dağılım için: $\mathrm{ARE}(\text{Wilcoxon},t)\geq108/125\approx0.864$.
  \item Çift üstel dağılımda: $\mathrm{ARE}=2.0$ (Wilcoxon iki kat verimli).
  \item Cauchy dağılımda: $\mathrm{ARE}=\infty$ ($t$-testi tutarsız; Wilcoxon tutarlı).
\end{itemize}
Sonuç: Wilcoxon normalliği hiçbir zaman çok fazla kaybetmez; ağır kuyrukta büyük kazanç.
\end{teorem}

%% ============================================================
\section{Bootstrap}
\label{sec:bootstrap}
%% ============================================================

\begin{kesif}
Tahmin edicimizin standart hatasını analitik olarak hesaplayamıyorsak ne yaparız? Bootstrap fikri: gerçek popülasyondan örnekleme yerine, elimizdeki veriyi popülasyon gibi kullan ve yeniden örnekle.
\end{kesif}

\begin{tanim}[Bootstrap]\label{def:bootstrap}
$\mathbf{x}=(x_1,\ldots,x_n)$ gözlem. Bootstrap algoritması:
\begin{enumerate}
  \item $b=1,\ldots,B$ için: $\mathbf{x}^{*b}=(x_1^*,\ldots,x_n^*)$ yerine koyarak $\mathbf{x}$'den çek.
  \item $\hat\theta^{*b} = T(\mathbf{x}^{*b})$ hesapla.
  \item Bootstrap standart hata: $\hat{\mathrm{SE}}_B = \sqrt{\frac{1}{B-1}\sum_{b=1}^B(\hat\theta^{*b}-\bar{\hat\theta}^*)^2}$.
\end{enumerate}
\end{tanim}

\begin{tanim}[Bootstrap güven aralıkları]\label{def:bootstrap_ga}
\begin{itemize}
  \item \textbf{Yüzdelik GA:} $[\hat\theta^*_{(\alpha/2)}, \hat\theta^*_{(1-\alpha/2)}]$.
  \item \textbf{BCa (bias-corrected and accelerated):} Yüzdelik yöntemini çarpıklık ve ivme için düzeltir.
\end{itemize}
\end{tanim}

\begin{teorem}[Bootstrap tutarlılığı]\label{thm:bootstrap_tutarlilik}
Geniş koşullar altında bootstrap dağılımı, gerçek örnekleme dağılımına yakınsar. Özellikle $\sqrt{n}$-tutarlı ve asimptotik normal $T(\mathbf{X})$ için bootstrap standart hata tutarlıdır.
\end{teorem}

\begin{ispatiskele}
Genelleşmiş delta yöntemi: $T(\mathbf{X})=g(\bar{X})$ biçimindeyse, $T^* - T$ dağılımının, $T - \theta$ dağılımına bootstrap yakınsaması MLT + Slütseva argümanıyla kurulur.
\end{ispatiskele}

\begin{disiplinlerarasibaglan}
\textbf{Biyoistatistik:} Küçük klinik örneklemlerde bootstrap güven aralıkları standart. \textbf{Makine öğrenmesi:} Bootstrap + bagging ile varyans azaltma (rastgele ormanlar). \textbf{Ekonometri:} Küme-sağlamlı standart hatalar bootstrap ile hesaplanır.
\end{disiplinlerarasibaglan}

%% ============================================================

\begin{mikroozet}
\begin{itemize}
  \item Parametresiz testler: minimal dağılım varsayımı; sıra veya işaret tabanlı.
  \item İşaret: medyan testi; Wilcoxon: simetri testi; Mann-Whitney: iki örneklem; KW/Friedman: çok grup.
  \item KS testi: ampirik ile teorik KDF farkı; $\sqrt{n}D_n\to$ Kolmogorov.
  \item Bootstrap: yeniden örnekleme ile standart hata ve GA; dağılım varsayımı gerektirmez.
\end{itemize}
\end{mikroozet}

\begin{ozyansima}
Parametresiz testler normallik varsayımı bozulduğunda her zaman parametrik testlerden daha iyi midir? Normallik geçerliyken ne kadar güç kaybı olur?
\end{ozyansima}

\begin{kopru}
\textbf{Nereden geldik:} Bölüm~\ref{chp:hipotez_testi}'deki test çerçevesi parametresiz yöntemleri de kapsar; temel fark test istatistiğinin seçimindedir.

\textbf{Nereye gidiyoruz:} Bölüm~\ref{chp:asimptotik}'da bootstrap tutarlılığının arkasındaki asimptotik teoriye, fonksiyonel MLT'ye ve Donsker teoreminin kesin formülasyonuna bakacağız.
\end{kopru}

%% ============================================================
%% ALISTIRMALAR
%% ============================================================

\cozumlualistirmalar

\begin{alistirma}[Al.17.1]{A}{İşaret testi}
$n=15$ hastada plasebo ile karşılaştırılan ilacın etkisi: 11 hasta iyileşti, 4 kötüleşti. $H_0:m=0$ (medyan fark sıfır) vs $H_1:m>0$. $\alpha=0.05$ için karar ver.
\end{alistirma}
\alcozum{%
$B=11\sim\Binom(15,0.5)$ ($H_0$ altında). $p=\Prob(B\geq11)=\sum_{k=11}^{15}\binom{15}{k}(0.5)^{15}=(1+15+105+455+3003\cdot0)/32768$. Doğrudan: $\Prob(B\geq11)=0.059>0.05$. $H_0$ reddedilemez (tam Binom ile $p=0.059$).%
}

\begin{alistirma}[Al.17.2]{B}{Mann-Whitney}
Grup A: $\{3,5,7,9,11\}$, Grup B: $\{2,4,8,10,12\}$. $H_0:F_A=F_B$. $U$ istatistiğini hesapla ve asimptotik $z$ ile test yap.
\end{alistirma}
\alcozum{%
$U=\#\{(A_i,B_j):A_i<B_j\}$. $3<4,8,10,12$ (4), $5<8,10,12$ (3), $7<8,10,12$ (3), $9<10,12$ (2), $11<12$ (1): $U=13$. $\Beklenti=12.5$, $\Var=5\cdot5\cdot11/12=22.9$. $Z=(13-12.5)/\sqrt{22.9}=0.1$. $H_0$ reddedilemez.%
}

\alistirmalar

\begin{alistirma}[Al.17.3]{A}{KS testi uygulaması}
$n=5$ gözlem: $\{0.2, 0.5, 0.6, 0.8, 0.9\}$. $H_0: F_0 = \Uniform[0,1]$. $D_5$'i hesapla.
\end{alistirma}
\alcozum{%
$F_5(x_i)=i/5$, $F_0(x_i)=x_i$. Farklar: $|0.2-0.2|=0$, $|0.4-0.5|=0.1$, $|0.6-0.6|=0$, $|0.8-0.8|=0$, $|1.0-0.9|=0.1$. Sol sınır farkları da hesapla: max$=0.1$. $D_5=0.1$. $\sqrt{5}\cdot0.1=0.224$ — KS tablosundan $\alpha=0.05$ kritik değeri $0.565$; $H_0$ reddedilemez.%
}

\begin{alistirma}[Al.17.4]{B}{Bootstrap standart hata}
$n=6$ gözlem: $\{2,4,5,7,9,12\}$. Medyanın bootstrap standart hatasını $B=5$ bootstrap yinelemesiyle elle hesapla (rastgele çekimler verildiğinde: $\{2,4,4,7,9,9\}$, $\{4,5,7,7,12,12\}$, $\{2,2,5,9,9,12\}$, $\{4,5,7,9,9,12\}$, $\{2,4,5,5,7,12\}$).
\end{alistirma}
\alcozum{%
Medyan(orijinal)$=6$. Bootstrap medyanları: $m_1=(4+7)/2=5.5$, $m_2=(7+7)/2=7$, $m_3=(5+9)/2=7$, $m_4=(7+9)/2=8$, $m_5=(5+5)/2=5$. $\bar{m}^*=6.5$. $\hat{SE}=\sqrt{[(5.5-6.5)^2+(7-6.5)^2+(7-6.5)^2+(8-6.5)^2+(5-6.5)^2]/4}=\sqrt{(1+0.25+0.25+2.25+2.25)/4}=\sqrt{1.5}=1.22$.%
}

\begin{alistirma}[Al.17.5]{C}{Glivenko-Cantelli hızı}
$X_1,\ldots,X_n\iid\Uniform[0,1]$. $\Beklenti[\sup_x|F_n(x)-x|]$'in üst sınırını Dvoretzky-Kiefer-Wolfowitz eşitsizliğini kullanarak ver: $\Prob(\sup_x|F_n(x)-F(x)|>\varepsilon)\leq 2e^{-2n\varepsilon^2}$.
\end{alistirma}
\alcozum{%
DKW'den $\Prob(D_n>\varepsilon)\leq2e^{-2n\varepsilon^2}$. $\Beklenti[D_n]=\int_0^\infty\Prob(D_n>t)\,dt\leq\int_0^\infty2e^{-2nt^2}\,dt=2\cdot\frac{\sqrt{\pi}}{2\sqrt{2n}}=\sqrt{\pi/(2n)}$. Dolayısıyla $\Beklenti[D_n]=O(1/\sqrt{n})$.%
}

% %% ============================================================
%% BÖLÜM 18 — Asimptotik İstatistik  [OPSİYONEL]
%% Olasılık Teorisi ve Matematiksel İstatistik
%% ============================================================
\chapter{Asimptotik İstatistik}
\label{chp:asimptotik}

\bolumalinti{%
  The virtue of the asymptotic approach is generality; its vice is the possibility
  that the approximation may be poor for the sample sizes encountered in practice.
}{Peter Bickel \& Kjell Doksum}

\begin{kazanimlar}
Bu bölümün sonunda okuyucu:
\begin{itemize}
  \item $\sqrt{n}$-tutarlılık ve asimptotik normallik kavramlarını tanımlar; delta yöntemiyle türetir.
  \item Fonksiyonel merkezi limit teoremini (Donsker) sezgisel olarak açıklar.
  \item Le Cam'ın yerel yakınsama yaklaşımını (LAN) tanımlar ve LAN ailelerinde optimum testi betimler.
  \item U-istatistiklerini tanımlar ve asimptotik dağılımlarını Hoeffding projeksiyonuyla hesaplar.
  \item M-tahmin edicilerinin (MLE, momentler, robust) genel asimptotik normalliğini uygular.
  \item Semiparametrik modellerde verimlilik sınırını kavramsal olarak açıklar.
\end{itemize}
\end{kazanimlar}

\begin{motivasyon}
Sonlu örneklem istatistiği genellikle analitik olarak çözülemez. Asimptotik teori, büyük örneklem davranışını betimleyen güçlü araçlar sunar; modern istatistiğin matematiksel temeli budur.

\medskip
\noindent\textbf{Köprü:} Bölüm~\ref{chp:merkezi_limit}'teki MLT ve delta yöntemi, burada fonksiyonel ve çok boyutlu versiyonlara genişletilir.
\end{motivasyon}

%% ============================================================
\section{$\sqrt{n}$-Tutarlılık ve Asimptotik Normallik}
\label{sec:asimptotik_normallik}
%% ============================================================

\begin{tanim}[$\sqrt{n}$-tutarlılık]\label{def:sqrtn_tutarli}
$T_n$ tahmin edici, $\sqrt{n}(T_n-\theta) = O_P(1)$ ise \textbf{$\sqrt{n}$-tutarlı} adını alır; yani her $\varepsilon>0$ için $M$ ve $n_0$ seçilebilir: $n\geq n_0 \Rightarrow \Prob(|\sqrt{n}(T_n-\theta)|>M)<\varepsilon$.
\end{tanim}

\begin{tanim}[Asimptotik normallik]\label{def:asimptotik_normal}
$T_n$ \textbf{asimptotik normal}, $\sigma^2(\theta)$ \textbf{asimptotik varyans}: $\sqrt{n}(T_n-\theta)\dto\Normal(0,\sigma^2(\theta))$.
\end{tanim}

\begin{teorem}[Çok boyutlu delta yöntemi]\label{thm:delta_cok}
$\sqrt{n}(\mathbf{T}_n-\boldsymbol{\theta})\dto\Normal_p(\mathbf{0},\Sigma)$ ve $g:\mathbb{R}^p\to\mathbb{R}^q$ $\boldsymbol{\theta}$'da türevlenebilir ($J_g$ Jacobian). O zaman:
\[
  \sqrt{n}(g(\mathbf{T}_n)-g(\boldsymbol{\theta}))\dto\Normal_q(\mathbf{0},J_g\Sigma J_g^\top).
\]
\end{teorem}

\begin{proof}
Taylor: $g(\mathbf{T}_n) - g(\boldsymbol{\theta}) = J_g(\boldsymbol{\theta})(\mathbf{T}_n-\boldsymbol{\theta}) + o(\|\mathbf{T}_n-\boldsymbol{\theta}\|)$. $\sqrt{n}\cdot o(\|\mathbf{T}_n-\boldsymbol{\theta}\|)\pto0$ (çünkü $\sqrt{n}(\mathbf{T}_n-\boldsymbol{\theta})=O_P(1)$ ve $\|\mathbf{T}_n-\boldsymbol{\theta}\|=O_P(1/\sqrt{n})$). Slütseva + sürekli harita teoremi.
\end{proof}

%% ============================================================
\section{M-Tahmin Edicileri}
\label{sec:m_tahmin}
%% ============================================================

\begin{tanim}[M-tahmin edici]\label{def:m_tahmin}
$\rho:\mathcal{X}\times\Theta\to\mathbb{R}$ kayıp fonksiyonu. \textbf{M-tahmin edici}:
\[
  \hat\theta_n = \arg\min_{\theta\in\Theta} \frac{1}{n}\sum_{i=1}^n\rho(X_i,\theta).
\]
$\rho(x,\theta)=-\log f(x;\theta)$ seçimi MLE verir.
\end{tanim}

\begin{teorem}[M-tahmin edicinin asimptotik normalliği]\label{thm:m_asimptotik}
Regularity koşulları ($\rho$ türevlenebilir, $\psi=\partial\rho/\partial\theta$) altında:
\[
  \sqrt{n}(\hat\theta_n-\theta_0)\dto\Normal\!\left(0, \frac{\Beklenti[\psi^2(X,\theta_0)]}{[\Beklenti[\dot\psi(X,\theta_0)]]^2}\right),
\]
$\dot\psi = \partial\psi/\partial\theta$.
\end{teorem}

\begin{proof}[Proof (taslak)]
Skor denklemi $\frac{1}{n}\sum\psi(X_i,\hat\theta_n)=0$'ı $\theta_0$ çevresinde Taylor:
\[
  0 = \frac{1}{n}\sum\psi(X_i,\theta_0) + (\hat\theta_n-\theta_0)\frac{1}{n}\sum\dot\psi(X_i,\theta^*).
\]
GBSK: $\frac{1}{n}\sum\dot\psi\asto\Beklenti[\dot\psi]$. MLT: $\frac{1}{\sqrt{n}}\sum\psi\dto\Normal(0,\Beklenti[\psi^2])$. Çöz: $\sqrt{n}(\hat\theta_n-\theta_0)\dto\Normal(0,\Beklenti[\psi^2]/\Beklenti[\dot\psi]^2)$.
\end{proof}

\begin{ornek}[Robust M-tahmin: Huber kaybı]\label{ex:huber}
Huber kaybı:
\[
  \rho_c(x,\theta) = \begin{cases}(x-\theta)^2/2 & |x-\theta|\leq c \\ c|x-\theta|-c^2/2 & |x-\theta|>c\end{cases}
\]
Bu MLE ve medyan arasında köprü; aykırı değerlere karşı sağlam.
\end{ornek}

%% ============================================================
\section{U-İstatistikleri}
\label{sec:u_istatistikleri}
%% ============================================================

\begin{tanim}[U-istatistiği]\label{def:u_istatistik}
$h(x_1,\ldots,x_m)$ simetrik çekirdek, $\theta = \Beklenti[h(X_1,\ldots,X_m)]$. \textbf{U-istatistiği}:
\[
  U_n = \binom{n}{m}^{-1}\sum_{1\leq i_1<\cdots<i_m\leq n}h(X_{i_1},\ldots,X_{i_m}).
\]
\end{tanim}

\begin{ornek}[U-istatistik örnekleri]
\begin{itemize}
  \item $m=1$, $h(x)=x$: $U_n=\bar{X}_n$.
  \item $m=2$, $h(x,y)=(x-y)^2/2$: $U_n=S^2$ (örneklem varyansı).
  \item $m=2$, $h(x,y)=\mathbf{1}[x+y>0]$: Wilcoxon işaret-sıra bağlantısı.
\end{itemize}
\end{ornek}

\begin{teorem}[Hoeffding projeksiyon]\label{thm:hoeffding}
$h_1(x) = \Beklenti[h(x,X_2,\ldots,X_m)] - \theta$ (birinci dereceden Hoeffding projeksiyon). O zaman:
\[
  \sqrt{n}(U_n - \theta) \dto \Normal(0, m^2\sigma_1^2)
\]
$\sigma_1^2 = \Var(h_1(X_1))$ ile.
\end{teorem}

\begin{ispatiskele}
$U_n - \theta = \frac{m}{n}\sum_{i=1}^n h_1(X_i) + R_n$ ayrışımı. $R_n = o_P(1/\sqrt{n})$ gösterilir (yüksek dereceli terimler). MLT: birinci terim $\sqrt{n}$-normalliği verir.
\end{ispatiskele}

%% ============================================================
\section{Fonksiyonel Merkezi Limit Teoremi}
\label{sec:fonksiyonel_mlt}
%% ============================================================

\begin{tanim}[Ampirik süreç]\label{def:ampirik_surec}
$X_1,\ldots,X_n\iid F$. \textbf{Ampirik süreç}:
\[
  \mathbb{G}_n(t) = \sqrt{n}(F_n(t) - F(t)), \quad t\in\mathbb{R}.
\]
\end{tanim}

\begin{teorem}[Donsker]\label{thm:donsker}
$U_i = F(X_i)\iid\Uniform[0,1]$ olsun. O zaman $\mathbb{G}_n \dto \mathbb{B}$ \emph{(fonksiyonel yakınsama)}, $\mathbb{B}$ \textbf{Brownian köprü} \emph{(Brownian bridge)}:
\begin{itemize}
  \item $\mathbb{B}(0)=\mathbb{B}(1)=0$.
  \item $\Beklenti[\mathbb{B}(s)\mathbb{B}(t)] = \min(s,t) - st$.
\end{itemize}
\end{teorem}

\begin{ispatiskele}
Yarı-metrik uzayda zayıf yakınsama. $D[0,1]$ Skorokhod uzayında sıkılık (tightness) ve finansal boyutsal yakınsama (MLT her sabit $t$'de) yeterli. Kolmogorov-Smirnov istatistiği $D_n=\|\mathbb{G}_n\|_\infty\dto\|\mathbb{B}\|_\infty$ = Kolmogorov dağılımı. Detay: van der Vaart (1998) Bölüm 19.
\end{ispatiskele}

%% ============================================================
\section{Yerel Asimptotik Normallik (LAN)}
\label{sec:lan}
%% ============================================================

\begin{tanim}[LAN koşulu]\label{def:lan}
$\{P_{\theta+h/\sqrt{n}}^{(n)}\}$ dizi model. \textbf{Yerel asimptotik normallik (LAN)} şartı: her $h\in\mathbb{R}$ için
\[
  \log\frac{dP_{\theta+h/\sqrt{n}}^{(n)}}{dP_\theta^{(n)}} = h\Delta_n(\theta) - \frac{h^2}{2}\mathcal{I}(\theta) + o_{P_\theta}(1)
\]
$\Delta_n(\theta)\dto\Normal(0,\mathcal{I}(\theta))$ altında $P_\theta^{(n)}$.
\end{tanim}

\begin{teorem}[Optimum test — LAN modeli]\label{thm:lan_optimum}
LAN koşulu altında $H_0:\theta=\theta_0$ vs $H_n:\theta=\theta_0+h/\sqrt{n}$. Yerel olarak en güçlü test: $\Delta_n(\theta_0)>z_\alpha$ (skor testi).
\end{teorem}

\begin{teorem}[Convolution teoremi]\label{thm:convolution}
LAN modeli, $\sqrt{n}$-tutarlı $T_n$. O zaman $\sqrt{n}(T_n-\theta)$'nın asimptotik dağılımı $\Normal(0,1/\mathcal{I}(\theta))*\nu$ biçimindedir ($*$ konvolüsyon; $\nu$ dağılım). Dolayısıyla asimptotik varyans $\geq 1/\mathcal{I}(\theta)$ — Cramér-Rao'nun asimptotik versiyonu.
\end{teorem}

%% ============================================================
\section{Asimptotik Göreli Verimlilik}
\label{sec:asimptotik_verimlilik}
%% ============================================================

\begin{tanim}[ARE]\label{def:are}
$T_n$ ve $T_n'$ asimptotik varyansları sırasıyla $\sigma^2$ ve $\sigma'^2$ olan iki tutarlı tahmin edici. \textbf{Asimptotik göreli verimlilik}:
\[
  \mathrm{ARE}(T_n,T_n') = \frac{\sigma'^2}{\sigma^2}.
\]
$\mathrm{ARE}>1$ ise $T_n$ daha verimli.
\end{tanim}

\begin{ornek}[Ortalama vs medyan — Normal]
$X_i\iid\Normal(\mu,\sigma^2)$. $\bar{X}_n$: asimptotik varyans $\sigma^2$. Medyan $\tilde{X}_n$: asimptotik varyans $\pi\sigma^2/2$ (Normal yoğunluğu medyanda $1/(\sigma\sqrt{2\pi})$; medyan asimptotik varyans $= 1/(4f(m)^2)$). ARE$($ medyan, ortalama$) = \sigma^2/(\pi\sigma^2/2) = 2/\pi \approx 0.637$.

Normal'de ortalama medyandan \%57 daha verimli. Ağır kuyrukta (Cauchy) ARE tersine döner.
\end{ornek}

\begin{teorem}[Pitman ARE — Wilcoxon vs $t$]\label{thm:pitman_are}
Simetrik sürekli dağılımda Wilcoxon testinin $t$-testine karşı Pitman ARE'si:
\[
  \mathrm{ARE}(W,t) = 12\sigma^2\left[\int f^2(x)\,dx\right]^2.
\]
Normal'de $\mathrm{ARE}=3/\pi\approx0.955$; Logistik'te $=1$; Cauchy'de $\to\infty$.
\end{teorem}

\begin{ispatiskele}
Her iki test asimptotik olarak normal; güç fonksiyonları $1-\beta=\Phi(c\cdot\mathrm{drift}/\sigma_T)$ biçimli. Eşit güç için gereken $n_W/n_t$ oranı $\mathrm{ARE}$'yi tanımlar. Detay: Lehmann (1975) §3.
\end{ispatiskele}

%% ============================================================

\begin{mikroozet}
\begin{itemize}
  \item M-tahmin edicileri genel asimptotik normallik sağlar; paydaki oran Cramér-Rao'ya indirgenir.
  \item U-istatistikleri: Hoeffding projeksiyonu $\sqrt{n}$-normalliği birinci dereceli projeksiyona indirir.
  \item Donsker: ampirik süreç Brownian köprüye; KS'nin asimptotiği buradan gelir.
  \item LAN + convolution: asimptotik Cramér-Rao sınırı, MLE optimalliğini kesinleştirir.
\end{itemize}
\end{mikroozet}

\begin{ozyansima}
Sonlu örneklem verimliliği ile asimptotik verimlilik çelişebilir mi? Hangi durumlarda asimptotik yaklaşım yanıltıcıdır?
\end{ozyansima}

\begin{kopru}
\textbf{Nereden geldik:} Bölüm~\ref{chp:buyuk_sayilar} ve~\ref{chp:merkezi_limit}'teki yakınsama araçları bu bölümün analitik temeli.

\textbf{Nereye gidiyoruz:} Bu kitabın kapsamı dışında kalan semiparametrik etkinlik (skor teoremi, varyans azaltma sınırları) için van der Vaart \emph{Asymptotic Statistics} (1998) kapsamlı referans.
\end{kopru}

%% ============================================================
%% ALISTIRMALAR
%% ============================================================

\cozumlualistirmalar

\begin{alistirma}[Al.18.1]{B}{Delta yöntemi — lognormal}
$X_1,\ldots,X_n\iid\mathrm{LogNormal}(\mu,\sigma^2)$ ($\log X_i\sim\Normal(\mu,\sigma^2)$). $\Beklenti[X]=e^{\mu+\sigma^2/2}$ için MLE asimptotik dağılımını delta yöntemiyle bul.
\end{alistirma}
\alcozum{%
MLE: $\hat\mu=\bar{Y}$, $\hat\sigma^2=\frac{1}{n}\sum(Y_i-\bar{Y})^2$ ($Y_i=\log X_i$). $\widehat{\Beklenti[X]}=\exp(\hat\mu+\hat\sigma^2/2)$. Delta yöntemi: $g(\mu,\sigma^2)=e^{\mu+\sigma^2/2}$; $\nabla g=(e^{\mu+\sigma^2/2}, e^{\mu+\sigma^2/2}/2)$. Asimptotik varyans: $(\Beklenti[X])^2(\sigma^2+\sigma^4/2)/n$ (Fisher bilgisi kullanılır).%
}

\begin{alistirma}[Al.18.2]{B}{U-istatistiği beklenti ve varyans}
$h(x,y) = (x-y)^2/2$ çekirdeğiyle $m=2$ U-istatistiği. (a) $\theta=\Beklenti[h]=\Var(X)$ olduğunu göster. (b) Hoeffding projeksiyon $h_1(x)$'i hesapla. (c) Asimptotik varyansı bul.
\end{alistirma}
\alcozum{%
(a) $\Beklenti[(X-Y)^2/2]=\Beklenti[(X-\mu)^2+(Y-\mu)^2]/2=\sigma^2$ ($X,Y$ bağımsız). (b) $h_1(x)=\Beklenti[h(x,X_2)]-\sigma^2=\Beklenti[(x-X_2)^2/2]-\sigma^2=(x-\mu)^2/2+\sigma^2/2-\sigma^2=(x-\mu)^2/2-\sigma^2/2$. (c) $\sigma_1^2=\Var((X-\mu)^2/2-\sigma^2/2)=\Var((X-\mu)^2)/4=(\mu_4-\sigma^4)/4$ ($\mu_4$=4. merkezi moment). Asimptotik varyans $4\sigma_1^2/n=(\mu_4-\sigma^4)/n$.%
}

\alistirmalar

\begin{alistirma}[Al.18.3]{B}{M-tahmin: medyan asimptotik varyansı}
$X_i\iid F$ sürekli, yoğunluk $f$. $\rho(x,\theta)=|x-\theta|$ seçimiyle M-tahmin edicinin $\hat\theta_n = $ medyan olduğunu ve asimptotik varyansının $1/(4f(m)^2)$ olduğunu göster.
\end{alistirma}
\alcozum{%
$\psi(x,\theta)=\mathrm{sgn}(x-\theta)$; $\Beklenti[\psi^2]=1$. $\dot\psi(x,\theta)=-2\delta(x-\theta)$; $\Beklenti[\dot\psi]=-2f(\theta)$. M-tahmin formülü: asimptotik varyans $=1/(4f(\theta_0)^2)$. $\theta_0=m$ medyan.%
}

\begin{alistirma}[Al.18.4]{C}{ARE: Double-Exponential dağılımı}
$X_i\iid\mathrm{DE}(\theta)=\frac{1}{2}e^{-|x-\theta|}$. Ortalama ve medyanın ARE'sini hesapla. Hangi tahmin edici daha verimli?
\end{alistirma}
\alcozum{%
$f(x)=\frac12 e^{-|x|}$; $\sigma^2=2$ (DE varyansı). Ortalama asimptotik varyans: $2/n$. Medyan: $f(0)=1/2$; asimptotik varyans $1/(4\cdot(1/2)^2)=1$. ARE(medyan,ortalama)$=2/1=2$. DE'de medyan ortalamadan iki kat daha verimli (ağır kuyruk etkisi).%
}

\begin{alistirma}[Al.18.5]{C}{LAN koşulunu doğrula}
$X_1,\ldots,X_n\iid\Normal(\theta,1)$. LAN koşulunun sağlandığını göster ve $\Delta_n(\theta)$'yı, $\mathcal{I}(\theta)$'yı belirle.
\end{alistirma}
\alcozum{%
$\log\frac{dP_{\theta+h/\sqrt n}^{(n)}}{dP_\theta^{(n)}}=\frac{h}{\sqrt n}\sum(X_i-\theta)-\frac{h^2}{2}$. $\Delta_n(\theta)=\frac{1}{\sqrt{n}}\sum(X_i-\theta)\dto\Normal(0,1)=\Normal(0,\mathcal{I})$ ($\mathcal{I}(\theta)=1$). Kalan terim $o_P(1)$: LAN sağlanır.%
}


% ---- Ekler ----
\appendix
%% ============================================================
%% EK A — Ölçüm Teorisi Temelleri
%% Olasılık Teorisi ve Matematiksel İstatistik
%% ============================================================
\chapter{Ölçüm Teorisi Temelleri}
\label{app:olcum_teorisi}

Bu ek, kitap boyunca kullanılan ölçüm ve integral teorisinin temel sonuçlarını özetlemektedir. İspatsız verilen lemma ve teoremlerin referansları Billingsley (1995) ve Klenke (2014)'dedir.

%% ============================================================
\section{Sigma-Cebirler ve Ölçüler}
\label{app:sigma_olcu}
%% ============================================================

\begin{tanim}[Halka ve cebir]\label{app:def:halka}
$\mathcal{R}\subseteq 2^\Omega$: eğer $A,B\in\mathcal{R}\Rightarrow A\cup B, A\setminus B\in\mathcal{R}$ ise \textbf{halka} \emph{(ring)}. Boşluk içeriyorsa \textbf{cebir} \emph{(algebra)}.
\end{tanim}

\begin{tanim}[$\sigma$-cebiri ve ölçü uzayı]
$\mathcal{F}\subseteq 2^\Omega$, $\sigma$-cebiri: $\Omega\in\mathcal{F}$; $A\in\mathcal{F}\Rightarrow A^c\in\mathcal{F}$; sayılabilir birleşim altında kapalı. $(\Omega,\mathcal{F})$ \textbf{ölçülebilir uzay} \emph{(measurable space)}. $\mu:\mathcal{F}\to[0,\infty]$ sonlu-eklenebilir ve $\sigma$-eklenebilir ise \textbf{ölçü} \emph{(measure)}.
\end{tanim}

\begin{teorem}[Caratheodory genişletme — özet]\label{app:thm:caratheodory}
$\mu_0$ halka üzerinde $\sigma$-sonlu ön-ölçü ise $\sigma(\mathcal{R})$ üzerinde yegâne ölçü $\mu$ ile $\mu\big|_\mathcal{R}=\mu_0$.
\end{teorem}

\begin{teorem}[Lebesgue ölçüsü]\label{app:thm:lebesgue_olcu}
$(\mathbb{R},\mathcal{B}(\mathbb{R}),\lambda)$ Lebesgue ölçü uzayı: $\lambda([a,b])=b-a$; Caratheodory genişletmesiyle yegâne belirlenir.
\end{teorem}

%% ============================================================
\section{Ölçülebilir Fonksiyonlar ve Lebesgue İntegrali}
\label{app:lebesgue_integral}
%% ============================================================

\begin{teorem}[Basit fonksiyon yakınsama]\label{app:thm:basit_yakin}
$f\geq0$ ölçülebilir ise $f_n\uparrow f$ n.k.\ ($f_n$ basit) olan artan basit fonksiyon dizisi vardır.
\end{teorem}

\begin{teorem}[Monoton Yakınsama Teoremi]\label{app:thm:myt}
$0\leq f_n\uparrow f$ n.k.\ ise $\int f_n\,d\mu\uparrow\int f\,d\mu$.
\end{teorem}

\begin{teorem}[Fatou Lemması]\label{app:thm:fatou}
$f_n\geq0$ ise $\int\liminf f_n\,d\mu\leq\liminf\int f_n\,d\mu$.
\end{teorem}

\begin{teorem}[Baskınlıklı Yakınsama Teoremi]\label{app:thm:byt}
$f_n\to f$ n.k., $|f_n|\leq g$, $\int g<\infty$: $\int|f_n-f|\,d\mu\to0$ ve $\int f_n\,d\mu\to\int f\,d\mu$.
\end{teorem}

\begin{teorem}[Fubini-Tonelli]\label{app:thm:fubini}
$(\Omega_1\times\Omega_2,\mathcal{F}_1\otimes\mathcal{F}_2,\mu_1\otimes\mu_2)$ çarpım uzayı. $f\geq0$ veya $\int|f|<\infty$ ise yinelenen integraller eşit:
\[
  \iint f\,d(\mu_1\otimes\mu_2) = \int\!\left(\int f\,d\mu_1\right)d\mu_2 = \int\!\left(\int f\,d\mu_2\right)d\mu_1.
\]
\end{teorem}

\begin{teorem}[Radon-Nikodym]\label{app:thm:radon_nikodym}
$\nu\ll\mu$ ($\nu$ mutlak sürekli $\mu$'ya göre). Yegâne $f\geq0$ ölçülebilir ile $\nu(A)=\int_A f\,d\mu$. $f = d\nu/d\mu$ \textbf{Radon-Nikodym türevi}.
\end{teorem}

%% ============================================================
\section{Karakteristik Fonksiyon Sonuçları}
\label{app:kf_sonuclar}
%% ============================================================

\begin{teorem}[Tersim formülü]\label{app:thm:tersim}
$\varphi_X$ karakteristik fonksiyon, $F_X$ KDF. Süreklilik noktalarında:
\[
  F_X(b) - F_X(a) = \lim_{T\to\infty}\frac{1}{2\pi}\int_{-T}^T\frac{e^{-ita}-e^{-itb}}{it}\varphi_X(t)\,dt.
\]
\end{teorem}

\begin{teorem}[Levy süreklilik teoremi]\label{app:thm:levy_sureklilk}
$X_n\dto X \iff \varphi_{X_n}(t)\to\varphi_X(t)$ her $t$ için; ve $\varphi_X$ $t=0$'da sürekli.
\end{teorem}

%% ============================================================
%% EK B — Olasılık Dağılımları Özet Tablosu
%% Olasılık Teorisi ve Matematiksel İstatistik
%% ============================================================
\chapter{Olasılık Dağılımları Özet Tablosu}
\label{app:dagilim_tablosu}

Bu ek, kitap boyunca kullanılan başlıca dağılımları sistematik biçimde özetler. Her dağılım için gösterim, yoğunluk/kütle fonksiyonu, beklenti, varyans, moment üreten fonksiyon (MÜF) ve karakteristik fonksiyon (KF) verilmiştir.

%% ============================================================
\section{Kesikli Dağılımlar}
\label{app:kesikli_dag}
%% ============================================================

\begin{center}
\renewcommand{\arraystretch}{1.6}
\small
\begin{tabular}{p{2.4cm}p{1.8cm}p{4.2cm}cc}
\hline
\textbf{Dağılım} & \textbf{Gösterim} & \textbf{KKF} $p(x)$ & $\boldsymbol{\mu}$ & $\boldsymbol{\sigma^2}$ \\
\hline
Bernoulli & $\mathrm{Ber}(p)$ & $p^x(1-p)^{1-x}$, $x\in\{0,1\}$ & $p$ & $p(1-p)$ \\
Binom & $\Binom(n,p)$ & $\binom{n}{x}p^x(1-p)^{n-x}$, $x=0,\ldots,n$ & $np$ & $np(1-p)$ \\
Poisson & $\Pois(\lambda)$ & $e^{-\lambda}\lambda^x/x!$, $x\in\mathbb{N}_0$ & $\lambda$ & $\lambda$ \\
Geometrik & $\mathrm{Geom}(p)$ & $(1-p)^{x-1}p$, $x=1,2,\ldots$ & $1/p$ & $(1-p)/p^2$ \\
Negatif Binom & $\mathrm{NB}(r,p)$ & $\binom{x-1}{r-1}p^r(1-p)^{x-r}$, $x=r,r+1,\ldots$ & $r/p$ & $r(1-p)/p^2$ \\
Hipergeometrik & $\mathrm{HG}(N,K,n)$ & $\dfrac{\binom{K}{x}\binom{N-K}{n-x}}{\binom{N}{n}}$ & $nK/N$ & $\dfrac{nK(N-K)(N-n)}{N^2(N-1)}$ \\
Uniform (kesikli) & $\mathrm{DU}\{1,\ldots,n\}$ & $1/n$, $x=1,\ldots,n$ & $(n+1)/2$ & $(n^2-1)/12$ \\
\hline
\end{tabular}
\end{center}

\subsection{MÜF ve KF — Kesikli Dağılımlar}

\begin{center}
\renewcommand{\arraystretch}{1.5}
\small
\begin{tabular}{p{2.4cm}p{4.5cm}p{4.5cm}}
\hline
\textbf{Dağılım} & \textbf{MÜF} $M(t)=\Beklenti[e^{tX}]$ & \textbf{KF} $\varphi(t)=\Beklenti[e^{itX}]$ \\
\hline
$\mathrm{Ber}(p)$ & $1-p+pe^t$ & $1-p+pe^{it}$ \\
$\Binom(n,p)$ & $(1-p+pe^t)^n$ & $(1-p+pe^{it})^n$ \\
$\Pois(\lambda)$ & $e^{\lambda(e^t-1)}$ & $e^{\lambda(e^{it}-1)}$ \\
$\mathrm{Geom}(p)$ & $pe^t/(1-(1-p)e^t)$, $t<-\log(1-p)$ & $pe^{it}/(1-(1-p)e^{it})$ \\
$\mathrm{NB}(r,p)$ & $[pe^t/(1-(1-p)e^t)]^r$ & $[pe^{it}/(1-(1-p)e^{it})]^r$ \\
\hline
\end{tabular}
\end{center}

%% ============================================================
\section{Sürekli Dağılımlar}
\label{app:surekli_dag}
%% ============================================================

\begin{center}
\renewcommand{\arraystretch}{1.7}
\small
\begin{tabular}{p{2.2cm}p{1.6cm}p{4.5cm}cc}
\hline
\textbf{Dağılım} & \textbf{Gösterim} & \textbf{PDF} $f(x)$ & $\boldsymbol{\mu}$ & $\boldsymbol{\sigma^2}$ \\
\hline
Uniform & $\Uniform[a,b]$ & $\frac{1}{b-a}$, $x\in[a,b]$ & $\frac{a+b}{2}$ & $\frac{(b-a)^2}{12}$ \\[4pt]
Normal & $\Normal(\mu,\sigma^2)$ & $\frac{1}{\sigma\sqrt{2\pi}}e^{-(x-\mu)^2/(2\sigma^2)}$, $x\in\mathbb{R}$ & $\mu$ & $\sigma^2$ \\[4pt]
Üstel & $\Exp(\lambda)$ & $\lambda e^{-\lambda x}$, $x>0$ & $1/\lambda$ & $1/\lambda^2$ \\[4pt]
Gama & $\Gama(\alpha,\beta)$ & $\frac{x^{\alpha-1}e^{-x/\beta}}{\Gamma(\alpha)\beta^\alpha}$, $x>0$ & $\alpha\beta$ & $\alpha\beta^2$ \\[4pt]
Beta & $\Beta(\alpha,\beta)$ & $\frac{x^{\alpha-1}(1-x)^{\beta-1}}{B(\alpha,\beta)}$, $x\in(0,1)$ & $\frac{\alpha}{\alpha+\beta}$ & $\frac{\alpha\beta}{(\alpha+\beta)^2(\alpha+\beta+1)}$ \\[4pt]
Cauchy & $\mathrm{Cau}(\mu,\sigma)$ & $\frac{1}{\pi\sigma[1+((x-\mu)/\sigma)^2]}$, $x\in\mathbb{R}$ & undef. & undef. \\[4pt]
$\chi^2_\nu$ & $\chi^2_\nu$ & $\Gama(\nu/2,2)$ özel durumu & $\nu$ & $2\nu$ \\[4pt]
$t_\nu$ & $t_\nu$ & $\frac{\Gamma((\nu+1)/2)}{\sqrt{\nu\pi}\,\Gamma(\nu/2)}\!\left(1+\frac{x^2}{\nu}\right)^{-(\nu+1)/2}$ & $0$ ($\nu>1$) & $\frac{\nu}{\nu-2}$ ($\nu>2$) \\[4pt]
$F_{\nu_1,\nu_2}$ & $F_{\nu_1,\nu_2}$ & Beta ailesinden & $\frac{\nu_2}{\nu_2-2}$ ($\nu_2>2$) & $\frac{2\nu_2^2(\nu_1+\nu_2-2)}{\nu_1(\nu_2-2)^2(\nu_2-4)}$ \\[4pt]
LogNormal & $\mathrm{LN}(\mu,\sigma^2)$ & $\frac{1}{x\sigma\sqrt{2\pi}}e^{-(\log x-\mu)^2/(2\sigma^2)}$, $x>0$ & $e^{\mu+\sigma^2/2}$ & $e^{2\mu+\sigma^2}(e^{\sigma^2}-1)$ \\[4pt]
Pareto & $\mathrm{Pa}(\alpha,x_m)$ & $\frac{\alpha x_m^\alpha}{x^{\alpha+1}}$, $x>x_m$ & $\frac{\alpha x_m}{\alpha-1}$ ($\alpha>1$) & $\frac{x_m^2\alpha}{(\alpha-1)^2(\alpha-2)}$ ($\alpha>2$) \\[4pt]
Weibull & $\mathrm{Wei}(k,\lambda)$ & $\frac{k}{\lambda}\!\left(\frac{x}{\lambda}\right)^{k-1}e^{-(x/\lambda)^k}$, $x>0$ & $\lambda\Gamma(1+1/k)$ & $\lambda^2[\Gamma(1+2/k)-\Gamma(1+1/k)^2]$ \\
\hline
\end{tabular}
\end{center}

\subsection{MÜF ve KF — Sürekli Dağılımlar}

\begin{center}
\renewcommand{\arraystretch}{1.5}
\small
\begin{tabular}{p{2.2cm}p{4.5cm}p{4.5cm}}
\hline
\textbf{Dağılım} & \textbf{MÜF} $M(t)$ & \textbf{KF} $\varphi(t)$ \\
\hline
$\Uniform[a,b]$ & $(e^{tb}-e^{ta})/(t(b-a))$ & $(e^{itb}-e^{ita})/(it(b-a))$ \\
$\Normal(\mu,\sigma^2)$ & $e^{\mu t+\sigma^2t^2/2}$ & $e^{i\mu t-\sigma^2t^2/2}$ \\
$\Exp(\lambda)$ & $\lambda/(\lambda-t)$, $t<\lambda$ & $\lambda/(\lambda-it)$ \\
$\Gama(\alpha,\beta)$ & $(1-\beta t)^{-\alpha}$, $t<1/\beta$ & $(1-i\beta t)^{-\alpha}$ \\
$\Beta(\alpha,\beta)$ & $1+\sum_{k=1}^\infty\frac{\prod_{r=0}^{k-1}(\alpha+r)}{B(\alpha,\beta)\,k!}t^k$ & $\vphantom{|}$Kümülatif seri \\
$\chi^2_\nu$ & $(1-2t)^{-\nu/2}$, $t<1/2$ & $(1-2it)^{-\nu/2}$ \\
$t_\nu$ & — (var olmaz) & $\frac{(\sqrt{\nu}|t|)^{\nu/2}K_{\nu/2}(\sqrt{\nu}|t|)}{2^{\nu/2-1}\Gamma(\nu/2)}$ \\
\hline
\end{tabular}
\end{center}

\noindent$K_\nu$: ikinci tür değiştirilmiş Bessel fonksiyonu.

%% ============================================================
\section{Çok Boyutlu Dağılımlar}
\label{app:cok_boyutlu_dag}
%% ============================================================

\begin{center}
\renewcommand{\arraystretch}{1.6}
\small
\begin{tabular}{p{3cm}p{2.5cm}p{7cm}}
\hline
\textbf{Dağılım} & \textbf{Gösterim} & \textbf{Açıklama} \\
\hline
Çok boyutlu Normal & $\Normal_p(\boldsymbol\mu,\Sigma)$ & $f(\mathbf{x})=(2\pi)^{-p/2}|\Sigma|^{-1/2}\exp\!\left(-\tfrac{1}{2}(\mathbf{x}-\boldsymbol\mu)^\top\Sigma^{-1}(\mathbf{x}-\boldsymbol\mu)\right)$; $\Beklenti=\boldsymbol\mu$, $\Cov=\Sigma$ \\
Dirichlet & $\mathrm{Dir}(\boldsymbol\alpha)$ & $f(\mathbf{p})\propto\prod_i p_i^{\alpha_i-1}$; Beta'nın $k$ boyutlu genellemesi; simplex üzerinde \\
Çok terimli & $\mathrm{Mult}(n,\mathbf{p})$ & $\Prob(\mathbf{X}=\mathbf{x})=\frac{n!}{x_1!\cdots x_k!}p_1^{x_1}\cdots p_k^{x_k}$; Binom'un genellemesi \\
Wishart & $W_p(n,\Sigma)$ & Normal popülasyondan $n$ örneklem; $\mathbf{S}=\sum \mathbf{X}_i\mathbf{X}_i^\top$; $\chi^2$'nin matris analogu \\
\hline
\end{tabular}
\end{center}

%% ============================================================
\section{Üstel Aile Özellikleri}
\label{app:ustel_aile}
%% ============================================================

\begin{tanim}[Tek parametreli üstel aile]
$f(x;\theta) = h(x)\exp[\eta(\theta)T(x) - B(\theta)]$: doğal parametre $\eta$, yeterli istatistik $T(x)$, bölümleme fonksiyonu $B(\theta)$.
\end{tanim}

\begin{center}
\renewcommand{\arraystretch}{1.5}
\small
\begin{tabular}{p{2.2cm}ccccc}
\hline
\textbf{Dağılım} & $h(x)$ & $\eta(\theta)$ & $T(x)$ & $B(\theta)$ & $\Beklenti[T(X)]$ \\
\hline
$\Binom(n,p)$ & $\binom{n}{x}$ & $\log\frac{p}{1-p}$ & $x$ & $n\log(1+e^\eta)$ & $np$ \\
$\Pois(\lambda)$ & $1/x!$ & $\log\lambda$ & $x$ & $e^\eta=\lambda$ & $\lambda$ \\
$\Normal(\mu,\sigma^2)$ & $(2\pi\sigma^2)^{-1/2}e^{-x^2/(2\sigma^2)}$ & $\mu/\sigma^2$ & $x$ & $\mu^2/(2\sigma^2)$ & $\mu$ \\
$\Exp(\lambda)$ & $1$ & $-\lambda$ & $x$ & $-\log\lambda$ & $1/\lambda$ \\
$\Gama(\alpha,\beta)$ & $x^{\alpha-1}/\Gamma(\alpha)$ & $-1/\beta$ & $x$ & $\alpha\log\beta$ & $\alpha\beta$ \\
\hline
\end{tabular}
\end{center}

%% ============================================================
\section{Önemli Özel Fonksiyonlar}
\label{app:ozel_fonksiyonlar}
%% ============================================================

\subsection{Gama ve Beta Fonksiyonları}

\[
  \Gamma(\alpha) = \int_0^\infty t^{\alpha-1}e^{-t}\,dt \quad (\alpha>0), \qquad
  B(\alpha,\beta) = \frac{\Gamma(\alpha)\Gamma(\beta)}{\Gamma(\alpha+\beta)} = \int_0^1 t^{\alpha-1}(1-t)^{\beta-1}\,dt.
\]

\noindent\textbf{Özel değerler:}
\[
  \Gamma(n) = (n-1)!, \quad \Gamma(1/2)=\sqrt{\pi}, \quad \Gamma(\alpha+1)=\alpha\Gamma(\alpha).
\]

\subsection{Standart Normal — Sayısal Değerler}

\[
  \Phi(z) = \int_{-\infty}^z\frac{1}{\sqrt{2\pi}}e^{-t^2/2}\,dt, \qquad \phi(z)=\Phi'(z)=\frac{1}{\sqrt{2\pi}}e^{-z^2/2}.
\]

\begin{center}
\begin{tabular}{ccccccc}
\hline
$z_\alpha$ & $0.842$ & $1.282$ & $1.645$ & $1.960$ & $2.326$ & $2.576$ \\
$\alpha$ & $0.200$ & $0.100$ & $0.050$ & $0.025$ & $0.010$ & $0.005$ \\
\hline
\end{tabular}
\end{center}

\subsection{Dağılım İlişkileri}

\begin{itemize}
  \item $X\sim\Normal(\mu,\sigma^2)$: $(X-\mu)/\sigma\sim\Normal(0,1)$.
  \item $X_1,\ldots,X_\nu\iid\Normal(0,1)$: $\sum X_i^2\sim\chi^2_\nu$.
  \item $Z\sim\Normal(0,1)$, $V\sim\chi^2_\nu$ bağımsız: $T=Z/\sqrt{V/\nu}\sim t_\nu$.
  \item $U\sim\chi^2_{\nu_1}$, $V\sim\chi^2_{\nu_2}$ bağımsız: $F=(U/\nu_1)/(V/\nu_2)\sim F_{\nu_1,\nu_2}$.
  \item $\Exp(\lambda)=\Gama(1,1/\lambda)$; $\chi^2_\nu=\Gama(\nu/2,2)$.
  \item $X\sim\Pois(\lambda)$: $\lambda\to\infty$'da CLT, $(X-\lambda)/\sqrt{\lambda}\dto\Normal(0,1)$.
  \item $X\sim\Binom(n,p)$: $n\to\infty$, $np\to\lambda$ ise $X\dto\Pois(\lambda)$.
  \item $t_\nu^2\sim F_{1,\nu}$; $F_{\nu_1,\nu_2,\alpha}=1/F_{\nu_2,\nu_1,1-\alpha}$.
  \item $X\sim F_{\nu_1,\nu_2}$: $\nu_2\to\infty$'da $\nu_1 X\dto\chi^2_{\nu_1}$.
\end{itemize}

\include{ekler/ekC_istatistik_tablolari}
%% ============================================================
%% EK D — Türkçe-İngilizce Terminoloji Sözlüğü
%% Olasılık Teorisi ve Matematiksel İstatistik
%% ============================================================
\chapter{Türkçe-İngilizce Terminoloji Sözlüğü}
\label{app:terminoloji}

Bu sözlük kitap boyunca kullanılan teknik terimlerin Türkçe-İngilizce karşılıklarını alfabetik olarak listeler. Her terimin ilk geçişi ilgili bölümde parantez içinde İngilizce karşılığıyla birlikte verilmiştir.

\section{A}

\begin{tabular}{p{5cm}p{5cm}p{4cm}}
\hline
\textbf{Türkçe} & \textbf{İngilizce} & \textbf{İlk geçiş} \\
\hline
Akış & flow & Böl.~\ref{chp:stokastik} \\
Alternatif hipotez & alternative hypothesis & Böl.~\ref{chp:hipotez_testi} \\
Ampirik dağılım fonksiyonu & empirical CDF & Böl.~\ref{chp:parametresiz} \\
Aperiyodik & aperiodic & Böl.~\ref{chp:stokastik} \\
Asimptotik göreli verimlilik & ARE & Böl.~\ref{chp:asimptotik} \\
Asimptotik normallik & asymptotic normality & Böl.~\ref{chp:asimptotik} \\
\hline
\end{tabular}

\section{B}

\begin{tabular}{p{5cm}p{5cm}p{4cm}}
\hline
\textbf{Türkçe} & \textbf{İngilizce} & \textbf{İlk geçiş} \\
\hline
Baskınlıklı yakınsama & dominated convergence & Böl.~\ref{chp:onbilgiler} \\
Beklenti & expected value & Böl.~\ref{chp:beklenti} \\
Belirlilik katsayısı & coefficient of determination ($R^2$) & Böl.~\ref{chp:dogrusal_modeller} \\
Bileşik hipotez & composite hypothesis & Böl.~\ref{chp:hipotez_testi} \\
Bootstrap & bootstrap & Böl.~\ref{chp:parametresiz} \\
Borel $\sigma$-cebiri & Borel $\sigma$-algebra & Böl.~\ref{chp:olcum_olasilik} \\
Brownian hareket & Brownian motion & Böl.~\ref{chp:stokastik} \\
\hline
\end{tabular}

\section{D}

\begin{tabular}{p{5cm}p{5cm}p{4cm}}
\hline
\textbf{Türkçe} & \textbf{İngilizce} & \textbf{İlk geçiş} \\
\hline
Dağılım fonksiyonu & distribution function (CDF) & Böl.~\ref{chp:rasgele_degiskenler} \\
Delta yöntemi & delta method & Böl.~\ref{chp:donusumler} \\
Doğrusal regresyon & linear regression & Böl.~\ref{chp:dogrusal_modeller} \\
Durağan dağılım & stationary distribution & Böl.~\ref{chp:stokastik} \\
Durma zamanı & stopping time & Böl.~\ref{chp:martingaller} \\
\hline
\end{tabular}

\section{E-G}

\begin{tabular}{p{5cm}p{5cm}p{4cm}}
\hline
\textbf{Türkçe} & \textbf{İngilizce} & \textbf{İlk geçiş} \\
\hline
En çok olabilirlik (MLE) & maximum likelihood estimator & Böl.~\ref{chp:nokta_tahmin} \\
Ergodik & ergodic & Böl.~\ref{chp:stokastik} \\
Eşit varyanslılık & homoscedasticity & Böl.~\ref{chp:dogrusal_modeller} \\
Filtrasyon & filtration & Böl.~\ref{chp:martingaller} \\
Fisher bilgisi & Fisher information & Böl.~\ref{chp:nokta_tahmin} \\
Gauss-Markov teoremi & Gauss-Markov theorem & Böl.~\ref{chp:dogrusal_modeller} \\
Geri dönüşlü & recurrent & Böl.~\ref{chp:stokastik} \\
Güven aralığı & confidence interval & Böl.~\ref{chp:guven_araliklari} \\
Güvenilir aralık & credible interval & Böl.~\ref{chp:nokta_tahmin} \\
\hline
\end{tabular}

\section{İ-K}

\begin{tabular}{p{5cm}p{5cm}p{4cm}}
\hline
\textbf{Türkçe} & \textbf{İngilizce} & \textbf{İlk geçiş} \\
\hline
İ.i.d. & i.i.d. & Böl.~\ref{chp:rasgele_degiskenler} \\
İçsınıf korelasyonu & intraclass correlation & Böl.~\ref{chp:anova} \\
İndirgeme & reduction & Böl.~\ref{chp:istatistik_temelleri} \\
Karakteristik fonksiyon & characteristic function & Böl.~\ref{chp:beklenti} \\
Konjugat önsel & conjugate prior & Böl.~\ref{chp:nokta_tahmin} \\
Koşullu beklenti & conditional expectation & Böl.~\ref{chp:beklenti} \\
Koşullu dağılım & conditional distribution & Böl.~\ref{chp:rasgele_degiskenler} \\
Kovaryans & covariance & Böl.~\ref{chp:beklenti} \\
\hline
\end{tabular}

\section{L-N}

\begin{tabular}{p{5cm}p{5cm}p{4cm}}
\hline
\textbf{Türkçe} & \textbf{İngilizce} & \textbf{İlk geçiş} \\
\hline
Lehmann-Scheffé & Lehmann-Scheffé theorem & Böl.~\ref{chp:nokta_tahmin} \\
Markov özelliği & Markov property & Böl.~\ref{chp:stokastik} \\
Markov zinciri & Markov chain & Böl.~\ref{chp:stokastik} \\
Martingal & martingale & Böl.~\ref{chp:martingaller} \\
Merkezi limit teoremi & central limit theorem & Böl.~\ref{chp:merkezi_limit} \\
Momentler yöntemi & method of moments & Böl.~\ref{chp:nokta_tahmin} \\
Monoton yakınsama & monotone convergence & Böl.~\ref{chp:onbilgiler} \\
Neredeyse kesin (n.k.) & almost surely (a.s.) & Böl.~\ref{chp:buyuk_sayilar} \\
\hline
\end{tabular}

\section{O-R}

\begin{tabular}{p{5cm}p{5cm}p{4cm}}
\hline
\textbf{Türkçe} & \textbf{İngilizce} & \textbf{İlk geçiş} \\
\hline
Olabilirlik & likelihood & Böl.~\ref{chp:istatistik_temelleri} \\
Olabilirlik oran testi & likelihood ratio test & Böl.~\ref{chp:hipotez_testi} \\
Olasılık ölçüsü & probability measure & Böl.~\ref{chp:olcum_olasilik} \\
Örneklem yolu & sample path & Böl.~\ref{chp:stokastik} \\
$p$-değeri & $p$-value & Böl.~\ref{chp:hipotez_testi} \\
Pivot & pivotal quantity & Böl.~\ref{chp:guven_araliklari} \\
Poisson süreci & Poisson process & Böl.~\ref{chp:stokastik} \\
Radon-Nikodym türevi & Radon-Nikodym derivative & Böl.~\ref{chp:beklenti} \\
Rasgele değişken & random variable & Böl.~\ref{chp:rasgele_degiskenler} \\
Rao-Blackwell & Rao-Blackwell theorem & Böl.~\ref{chp:nokta_tahmin} \\
\hline
\end{tabular}

\section{S-Z}

\begin{tabular}{p{5cm}p{5cm}p{4cm}}
\hline
\textbf{Türkçe} & \textbf{İngilizce} & \textbf{İlk geçiş} \\
\hline
Sıfır hipotezi & null hypothesis & Böl.~\ref{chp:hipotez_testi} \\
Sonsal dağılım & posterior distribution & Böl.~\ref{chp:nokta_tahmin} \\
Stokastik süreç & stochastic process & Böl.~\ref{chp:stokastik} \\
$\sigma$-cebiri & $\sigma$-algebra & Böl.~\ref{chp:olcum_olasilik} \\
Tahmin edici & estimator & Böl.~\ref{chp:nokta_tahmin} \\
Tamamlayıcı istatistik & complete statistic & Böl.~\ref{chp:istatistik_temelleri} \\
Tip I hata & type I error & Böl.~\ref{chp:hipotez_testi} \\
Tutarlı (tahmin) & consistent (estimator) & Böl.~\ref{chp:nokta_tahmin} \\
UMVUE & UMVUE & Böl.~\ref{chp:nokta_tahmin} \\
Üstel aile & exponential family & Böl.~\ref{chp:dagilim_aileleri} \\
Varyans analizi (ANOVA) & analysis of variance & Böl.~\ref{chp:anova} \\
Yansız tahmin edici & unbiased estimator & Böl.~\ref{chp:nokta_tahmin} \\
Yeterli istatistik & sufficient statistic & Böl.~\ref{chp:istatistik_temelleri} \\
Yakınsama dağılımda & convergence in distribution & Böl.~\ref{chp:buyuk_sayilar} \\
Yakınsama olasılıkta & convergence in probability & Böl.~\ref{chp:buyuk_sayilar} \\
\hline
\end{tabular}

%% ============================================================
%% EK E — Sık Kullanılan İspat Teknikleri
%% Olasılık Teorisi ve Matematiksel İstatistik
%% ============================================================
\chapter{Sık Kullanılan İspat Teknikleri}
\label{app:ispat_teknikleri}

Bu ek, olasılık teorisi ve matematiksel istatistikte sıkça başvurulan ispat tekniklerini özetler. Her teknik için kısa açıklama ve kitaptaki somut uygulamalar verilir.

%% ============================================================
\section{Sınır Geçirme: MYT, Fatou, BYT}
\label{app:sinir_gecirme}
%% ============================================================

Sınır geçirme argümanları, $\lim$ ile $\int$'ı değiştirmek için kullanılır.

\begin{itemize}
  \item \textbf{MYT (Monoton Yakınsama):} $f_n\uparrow f$, $f_n\geq0$ ise $\int f_n\uparrow\int f$.
  \item \textbf{Fatou:} $f_n\geq0$ ise $\int\liminf f_n\leq\liminf\int f_n$.
  \item \textbf{BYT (Baskınlıklı):} $f_n\to f$, $|f_n|\leq g$, $\int g<\infty$ ise $\int f_n\to\int f$.
\end{itemize}

\noindent\textbf{Uygulama:} Karakteristik fonksiyon türevi, Fisher bilgisi türetimi.

%% ============================================================
\section{$\pi$-$\lambda$ (Dynkin) Argümanı}
\label{app:pi_lambda}
%% ============================================================

Bir özelliğin tüm olaylar üzerinde geçerli olduğunu kanıtlamak için:
\begin{enumerate}
  \item Özelliğin $\pi$-sistem (kesişim altında kapalı koleksiyon) üzerinde geçerli olduğunu göster.
  \item Özelliği taşıyan kümelerin $\lambda$-sistem (monoton sınıflara kapalı) oluşturduğunu göster.
  \item Dynkin: $\pi$-sistemin ürettiği $\sigma$-cebir, $\lambda$-sistemde.
\end{enumerate}

\noindent\textbf{Uygulama:} Kolmogorov 0-1 Kanunu (Böl.~\ref{chp:buyuk_sayilar}), bağımsızlığın $\sigma$-cebiri karakterizasyonu.

%% ============================================================
\section{Kırpma Yöntemi}
\label{app:kirpma}
%% ============================================================

$X$ rasgele değişken, $c>0$: $X^{(c)} = X\mathbf{1}[|X|\leq c]$ (kırpılmış). Strateji:
\begin{enumerate}
  \item Kırpılmış $X^{(c)}$ için sonucu kanıtla (sınırlı değişken, kolay).
  \item $c\to\infty$ alarak orijinal değişkene geç (Fatou/BYT).
\end{enumerate}

\noindent\textbf{Uygulama:} Genel ZBSK (Böl.~\ref{chp:buyuk_sayilar}), MLE tutarlılık ispatı.

%% ============================================================
\section{Cauchy-Schwarz ve Jensen Kullanımı}
\label{app:cs_jensen}
%% ============================================================

\begin{itemize}
  \item \textbf{Cauchy-Schwarz:} $[\Cov(X,Y)]^2\leq\Var(X)\Var(Y)$. Cramér-Rao ispatının kalbi.
  \item \textbf{Jensen ($f$ konveks):} $f(\Beklenti[X])\leq\Beklenti[f(X)]$. Markov eşitsizliğinin üreteci; entropide bilgi eşitsizliği.
\end{itemize}

%% ============================================================
\section{Moment Üretici Fonksiyon (MÜF) Argümanı}
\label{app:muf_arguman}
%% ============================================================

$M_X(t)=\Beklenti[e^{tX}]$ bir çevresinde sonlu ise dağılımı yegâne belirler. MÜF/KF çarpım özelliği bağımsız toplamların dağılımını verir.

\noindent\textbf{Uygulama:} Binom$\to$Poisson limiti, çarpım yapısı ispatları.

%% ============================================================
\section{Slütseva Teoremi}
\label{app:slutseva}
%% ============================================================

$X_n\dto X$ ve $Y_n\pto c$ (sabit) ise:
\[
  X_n+Y_n\dto X+c, \quad X_nY_n\dto cX, \quad X_n/Y_n\dto X/c \;(c\neq0).
\]

\noindent\textbf{Uygulama:} Her asimptotik normallik ispatında ($\hat\sigma\to\sigma$), $t$-istatistiği asimptotiği.

%% ============================================================
\section{Varyans Ayrışımı}
\label{app:varyans_ayrisim}
%% ============================================================

$\Var(Y) = \Beklenti[\Var(Y\mid X)] + \Var(\Beklenti[Y\mid X])$.

\noindent\textbf{Uygulama:} Rao-Blackwell teoremi, karışık modellerde varyans bileşenleri.

%% ============================================================
\section{Simetri Argümanı}
\label{app:simetri}
%% ============================================================

Bir rasgele değişken simetrisi ve değişim argümanı:
\begin{itemize}
  \item $X\overset{d}{=}-X$ ise $\Beklenti[X]=0$ (varsa).
  \item İzotropik dağılımlar: $\mathbf{X}\overset{d}{=}O\mathbf{X}$ her ortogonal $O$ için.
\end{itemize}

\noindent\textbf{Uygulama:} Normal popülasyon teoremi (ortogonal dönüşüm invariansı), işaret testi argümanları.

%% ============================================================
\section{Karakteristik Fonksiyon İspatı Akışı}
\label{app:kf_ispat_akisi}
%% ============================================================

MLT'nin KF ispatı şu adımları izler:
\begin{enumerate}
  \item $\varphi_{\bar{X}_n}(t) = [\varphi_{X_1}(t/\sqrt{n})]^n$ (bağımsız toplam KF çarpım özelliği).
  \item Taylor: $\varphi_{X_1}(u) = 1 - \sigma^2 u^2/2 + o(u^2)$ ($u\to0$).
  \item $[1-\sigma^2t^2/(2n)+o(1/n)]^n\to e^{-\sigma^2t^2/2}$ ($n\to\infty$).
  \item Levy süreklilik teoremi: $\Normal(0,\sigma^2)$'ye yakınsama.
\end{enumerate}

Bu akış Bölüm~\ref{chp:merkezi_limit}'te tam detayla sunulmuştur.

%% ============================================================
%% EK F — Kaynakça ve İleri Okuma
%% Olasılık Teorisi ve Matematiksel İstatistik
%% ============================================================
\chapter{Kaynakça ve İleri Okuma}
\label{app:kaynaklar}

\section{Temel Kaynaklar}

Bu kitabın yazımında doğrudan yararlanılan başvurular:

\begin{description}
  \item[Akdeniz (2013)] Fikri Akdeniz. \textit{Olasılık ve İstatistik.} 18. Baskı. Çukurova Üniversitesi, Adana. — Türkçe bağlam için temel referans; örnek ve alıştırmalar.

  \item[Billingsley (1995)] Patrick Billingsley. \textit{Probability and Measure.} 3rd ed. Wiley. — Ölçüm teorisi ve olasılığın kapsamlı birleşimi; Böl.~0--2 için.

  \item[Casella-Berger (2002)] George Casella, Roger Berger. \textit{Statistical Inference.} 2nd ed. Duxbury. — İstatistik Part II'nin ana kaynağı; yeterlilik, UMP, regresyon.

  \item[Durrett (2010)] Rick Durrett. \textit{Probability: Theory and Examples.} 4th ed. Cambridge. — Martingaller, GBSK, MLT için birincil kaynak.

  \item[Feller Vol.I (1968)] William Feller. \textit{An Introduction to Probability Theory and Its Applications.} Vol.~I, 3rd ed. Wiley. — Kesikli olasılık ve limit teoremleri; sezgi için eşsiz.

  \item[Feller Vol.II (1971)] William Feller. \textit{An Introduction to Probability Theory and Its Applications.} Vol.~II, 2nd ed. Wiley. — Sürekli dağılımlar, Laplace dönüşümü, yenileme.

  \item[Gnedenko (1962)] Boris Gnedenko. \textit{The Theory of Probability.} Chelsea. — Klasik Rus ekolü; limit teoremlerinin tarihsel kökeni.

  \item[Hogg-Craig (1995)] Robert Hogg, Allen Craig. \textit{Introduction to Mathematical Statistics.} 5th ed. Prentice Hall. — İstatistik Part II için tamamlayıcı kaynak.

  \item[Klenke (2014)] Achim Klenke. \textit{Probability Theory: A Comprehensive Course.} 2nd ed. Springer. — Modern ölçüm teorisi yaklaşımı; Böl.~0 için ek referans.

  \item[Kolmogorov (1950)] Andrei Kolmogorov. \textit{Foundations of the Theory of Probability.} Chelsea. — Aksiyomatik olasılığın orijinal metni.

  \item[Lehmann-Casella (1998)] Erich Lehmann, George Casella. \textit{Theory of Point Estimation.} 2nd ed. Springer. — UMVUE, Cramér-Rao, semiparametrik etkinlik.

  \item[Lehmann-Romano (2005)] Erich Lehmann, Joseph Romano. \textit{Testing Statistical Hypotheses.} 3rd ed. Springer. — Neyman-Pearson, UMP, ORT için kapsamlı referans.

  \item[Mood-Graybill-Boes (1974)] A. Mood, F. Graybill, D. Boes. \textit{Introduction to the Theory of Statistics.} 3rd ed. McGraw-Hill.

  \item[Shiryaev Vol.I (1996)] Albert Shiryaev. \textit{Probability.} 2nd ed. Springer. — Rus geleneği; stokastik süreçler giriş.

  \item[Shiryaev Vol.II (2019)] Albert Shiryaev. \textit{Probability-2.} Springer. — Martingaller, ergodik teori.

  \item[Williams (1991)] David Williams. \textit{Probability with Martingales.} Cambridge. — Martingal teorisinin en anlaşılır sunumu; Böl.~\ref{chp:martingaller} için birincil.
\end{description}

\section{İleri Okuma}

\subsection{Olasılık Teorisi}

\begin{description}
  \item[Revuz-Yor (1999)] D. Revuz, M. Yor. \textit{Continuous Martingales and Brownian Motion.} Springer. — Stokastik integral, Itô formülü.
  \item[Kallenberg (2002)] O. Kallenberg. \textit{Foundations of Modern Probability.} 2nd ed. Springer. — Kapsamlı, modern; ileri düzey referans.
  \item[Bass (2011)] Richard Bass. \textit{Stochastic Processes.} Cambridge. — Uygulamalı stokastik süreçler.
\end{description}

\subsection{İstatistik}

\begin{description}
  \item[van der Vaart (1998)] A. W. van der Vaart. \textit{Asymptotic Statistics.} Cambridge. — Donsker, LAN, semiparametrik etkinlik; Böl.~\ref{chp:asimptotik}'ın doğal devamı.
  \item[Wainwright (2019)] M. Wainwright. \textit{High-Dimensional Statistics.} Cambridge. — Büyük boyutlu modern istatistik.
  \item[Efron-Hastie (2016)] B. Efron, T. Hastie. \textit{Computer Age Statistical Inference.} Cambridge. — Bootstrap, lasso, makine öğrenmesi köprüsü.
\end{description}

\subsection{Türkçe Kaynaklar}

\begin{description}
  \item[Akdeniz (2013)] Yukarıda.
  \item[Önder vd. (çeşitli)] Türkiye'deki çeşitli üniversitelerin olasılık-istatistik ders notları (erişim: YÖK Açık Erişim).
\end{description}

\section{Yazılım Kaynakları}

\begin{description}
  \item[R (CRAN)] \texttt{stats}, \texttt{MASS}, \texttt{survival} paketleri; parametrik ve parametresiz testler.
  \item[Python] \texttt{scipy.stats}, \texttt{statsmodels}, \texttt{numpy}; simülasyon ve büyük örneklem doğrulama.
  \item[SageMath] Sembolik hesaplama; olasılık dağılımı türevleri, MÜF hesapları.
\end{description}


\backmatter

\end{document}
